% ============================================================
%
%   WORKING NOTES ON A-INFINITY CHIRAL ALGEBRAS
%   AND 3D HOLOMORPHIC-TOPOLOGICAL QFT
%
%   Raeez Lorgat, 2024--2026
%
%   Compile standalone: pdflatex working_notes.tex
%
% ============================================================

\documentclass[11pt]{article}

\usepackage[T1]{fontenc}
\usepackage[utf8]{inputenc}
\usepackage{lmodern}
\frenchspacing

\usepackage[cmintegrals,cmbraces,noamssymbols]{newtxmath}
\usepackage{ebgaramond}

\usepackage[
  activate={true,nocompatibility}, final,
  tracking=true, kerning=true, spacing=true,
  factor=1100, stretch=10, shrink=10
]{microtype}

\usepackage{mleftright}\mleftright

\let\Bbbk\undefined \let\openbox\undefined
\usepackage{amsmath,amssymb,amsthm}
\usepackage{mathtools}
\usepackage{enumitem}
\usepackage[margin=1.15in]{geometry}
\usepackage[unicode]{hyperref}
\usepackage{tikz-cd}

\setlength{\emergencystretch}{3em}
\raggedbottom

% --- Environments ---
\usepackage{thmtools}
\declaretheoremstyle[
  spaceabove=\topsep, spacebelow=\topsep,
  headfont=\normalfont\scshape,
  notefont=\normalfont\itshape,
  bodyfont=\normalfont,
  postheadspace=0.5em, headpunct={.}
]{notesthm}
\declaretheoremstyle[
  spaceabove=\topsep, spacebelow=\topsep,
  headfont=\normalfont\itshape,
  notefont=\normalfont\itshape,
  bodyfont=\normalfont,
  postheadspace=0.5em, headpunct={.}
]{notesdef}

\declaretheorem[style=notesthm, name=Theorem, numberwithin=section]{theorem}
\declaretheorem[style=notesthm, name=Lemma, sibling=theorem]{lemma}
\declaretheorem[style=notesthm, name=Proposition, sibling=theorem]{proposition}
\declaretheorem[style=notesthm, name=Corollary, sibling=theorem]{corollary}
\declaretheorem[style=notesdef, name=Definition, sibling=theorem]{definition}
\declaretheorem[style=notesdef, name=Remark, sibling=theorem]{remark}
\declaretheorem[style=notesdef, name=Observation, sibling=theorem]{observation}
\declaretheorem[style=notesdef, name=Warning, sibling=theorem]{warning}
\declaretheorem[style=notesdef, name=Example, sibling=theorem]{example}
\declaretheorem[style=notesdef, name=Conjecture, sibling=theorem]{conjecture}
\declaretheorem[style=notesdef, name=Question, sibling=theorem]{question}

% --- Macros ---
\newcommand{\cA}{\mathcal{A}}
\newcommand{\cB}{\mathcal{B}}
\newcommand{\cC}{\mathcal{C}}
\newcommand{\cM}{\mathcal{M}}
\newcommand{\cN}{\mathcal{N}}
\newcommand{\cO}{\mathcal{O}}
\newcommand{\cW}{\mathcal{W}}
\newcommand{\cH}{\mathcal{H}}
\newcommand{\cZ}{\mathcal{Z}}
\newcommand{\cD}{\mathcal{D}}
\newcommand{\cI}{\mathcal{I}}
\newcommand{\cL}{\mathcal{L}}
\newcommand{\fg}{\mathfrak{g}}
\newcommand{\fh}{\mathfrak{h}}
\newcommand{\fn}{\mathfrak{n}}
\newcommand{\fsl}{\mathfrak{sl}}
\newcommand{\barB}{\overline{B}}
\newcommand{\B}{\bar{B}}
\newcommand{\C}{\mathbb{C}}
\newcommand{\Z}{\mathbb{Z}}
\newcommand{\R}{\mathbb{R}}
\newcommand{\Q}{\mathbb{Q}}
\newcommand{\bP}{\mathbb{P}}
\newcommand{\bD}{\mathbb{D}}

\newcommand{\Ass}{\mathsf{Ass}}
\newcommand{\Com}{\mathsf{Com}}
\newcommand{\Lie}{\mathsf{Lie}}
\newcommand{\Eone}{\mathsf{E}_1}
\newcommand{\Etwo}{\mathsf{E}_2}
\newcommand{\Einf}{\mathsf{E}_{\infty}}
\newcommand{\Linf}{\mathsf{L}_{\infty}}
\newcommand{\Ainf}{\mathsf{A}_{\infty}}
\newcommand{\SCchtop}{\mathsf{SC}^{\mathrm{ch,top}}}
\newcommand{\Ran}{\operatorname{Ran}}
\newcommand{\Res}{\operatorname{Res}}
\newcommand{\FM}{\mathrm{FM}}
\newcommand{\MC}{\mathrm{MC}}
\newcommand{\ad}{\operatorname{ad}}
\newcommand{\id}{\mathrm{id}}
\newcommand{\Hom}{\operatorname{Hom}}
\newcommand{\RHom}{\operatorname{RHom}}
\newcommand{\End}{\operatorname{End}}
\newcommand{\Ext}{\operatorname{Ext}}
\newcommand{\HH}{\operatorname{HH}}
\newcommand{\Tr}{\operatorname{Tr}}
\newcommand{\Perf}{\operatorname{Perf}}
\newcommand{\ChirHoch}{\mathrm{ChirHoch}}
\newcommand{\gAmod}{\mathfrak{g}^{\mathrm{mod}}_\cA}
\newcommand{\dzero}{d_0}
\newcommand{\dfib}{d_{\mathrm{fib}}}
\newcommand{\Dtot}{D_{\mathrm{tot}}}
\newcommand{\normord}[1]{:\!#1\!:}
\newcommand{\Zder}{\cZ^{\mathrm{der}}_{\mathrm{ch}}}

\DeclareMathOperator{\gr}{gr}
\DeclareMathOperator{\Spec}{Spec}
\DeclareMathOperator{\Conf}{Conf}
\DeclareMathOperator{\Sym}{Sym}
\DeclareMathOperator{\diag}{diag}

\declaretheorem[style=notesdef, name=Computation, sibling=theorem]{computation}

% Separator
\newcommand{\sep}{\bigskip\begin{center}$*$\end{center}\medskip}

% Blackboard bold (used in working notes but not defined in monograph preamble)
\providecommand{\bR}{\mathbb{R}}
\providecommand{\bC}{\mathbb{C}}
\providecommand{\bZ}{\mathbb{Z}}
\providecommand{\bQ}{\mathbb{Q}}
\providecommand{\bN}{\mathbb{N}}

\title{\textit{Working Notes on $\Ainf$ Chiral Algebras\\
and 3D Holomorphic-Topological QFT}}
\author{Raeez Lorgat}
\date{2024--2026}

\begin{document}
\maketitle

\begin{abstract}
The $1$-form $\eta_{ij} = d\log(z_i - z_j)$ acquires a second
direction.  OPE residues on~$\mathrm{FM}_k(\mathbb{C})$ give the
bar differential; ordered deconcatenation
on~$\mathrm{Conf}_k(\mathbb{R})$ gives the coproduct.  Together
they make the bar complex a coalgebra over the Swiss-cheese
cooperad, the operadic fingerprint
of a $3$d holomorphic-topological field theory
on~$\mathbb{C} \times \mathbb{R}$:
$(-1)$-shifted Poisson vertex algebra on cohomology,
the bulk-boundary-line triangle, dg-shifted Yangians on the line,
and curved Swiss-cheese at genus~$g \ge 1$ where
$d_{\mathrm{fib}}^{\,2} = \kappa \cdot \omega_g$ couples the
topological ordering to the Hodge bundle.
\end{abstract}

\begin{quote}
\itshape
Volume I took a $1$-form and followed where it led: into the
bar complex, the five theorems, the shadow obstruction tower.
These notes are about what happens when the $1$-form acquires
a second direction, what the bar complex becomes when read in
three dimensions, and why the coproduct contains more physics
than the differential.
\end{quote}

\tableofcontents


% =================================================================
\section{Two directions}
\label{sec:two-directions}
% =================================================================

The bar complex of a chiral algebra $\cA$ on a curve $X$ carries
two structures.

\medskip

The first is the \textbf{differential} $d_{\barB}$.  It extracts
OPE residues at collision divisors in $\FM_k(X)$.  Points
approach each other along the curve.  The propagator
$\eta_{ij} = d\log(z_i - z_j)$ detects each collision and records
its residue.  Arnold's relation kills triple residues, so
$d_{\barB}^2 = 0$ at genus $0$.  This is the holomorphic
direction.  It is everything Volume~I studied.

\medskip

The second is the \textbf{coproduct} $\Delta$.  On the tensor
coalgebra $T^c(s^{-1}\bar\cA)$, deconcatenation is
\[
\Delta[a_1 \,|\, \cdots \,|\, a_n]
\;=\;
\sum_{i=0}^{n}
[a_1 \,|\, \cdots \,|\, a_i]
\;\otimes\;
[a_{i+1} \,|\, \cdots \,|\, a_n].
\]
This is the ordered splitting of a word into two subwords.
It does not depend on the curve, the OPE, or the central
charge.  It depends on the ordering.  Points are separated
along a line, and the line is cut at one of the $n+1$
positions between them.  This is the topological direction.

\medskip

The differential lives on $\FM_k(\C)$.  The coproduct lives on
$\Conf_k(\R)$.  Together: a bar element of degree $k$ is
parametrised by
\[
\FM_k(\C) \;\times\; \Conf_k(\R),
\]
the product of holomorphic and topological configuration spaces.

\sep

This product is the operadic fingerprint of a $3$d
holomorphic-topological QFT on $\C_z \times \R_t$.  Observables
factorise holomorphically in $z$ and associatively in $t$.  The
two-coloured Swiss-cheese operad $\SCchtop$ has operation spaces
$\FM_k(\C) \times \Eone(m)$: the closed colour is the
Fulton--MacPherson compactification (holomorphic collisions); the
open colour is the little intervals operad (topological
splittings).  The bar differential is the closed colour.  The bar
coproduct is the open colour.

The same $1$-form $\eta_{ij}$ generates both structures.
Its collision residues produce the differential.  Its ordering
(which point comes first along the time axis) produces the
coproduct.  One form, two directions, one product.

\sep

At genus $0$, $d_{\barB}^2 = 0$ and $\Delta$ is coassociative
and compatible with $d_{\barB}$: these are the Swiss-cheese
relations without obstruction.  At genus $g \geq 1$, the
curvature $\dfib^{\,2} = \kappa(\cA) \cdot \omega_g$ from
Volume~I infects the product:
\[
m_1^2(a) \;=\; [m_0, a], \qquad m_0 = \kappa(\cA) \cdot \omega_g.
\]
The Swiss-cheese algebra becomes \textit{curved}.  The failure of
formality at genus $g \geq 1$ is the bar construction's single
most important feature.  It is not a defect.  It is $\kappa$.


% =================================================================
\section{The recognition theorem}
\label{sec:recognition}
% =================================================================

A factorisation algebra on $\C \times \R$ satisfying local
constancy along~$\R$ is the same thing as an $\SCchtop$-algebra.

The proof uses the product Weiss topology.  A cover of
$U \times I \subset \C \times \R$ is a collection of products
$\{U_\alpha \times I_\beta\}$ refining the product open cover.
The cosheaf of observables on $\C \times \R$ restricts to a
constructible dg-cosheaf on $\FM_k(\C) \times \Conf_k(\R)$.
Weiss descent identifies the structure maps with the operations
of $\SCchtop$.

\begin{theorem}[Recognition; proved]
\label{thm:recognition-wn}
There is an equivalence of $\infty$-categories between:
\begin{enumerate}[label=\textup{(\roman*)},nosep]
\item logarithmic $\SCchtop$-algebras, and
\item factorisation algebras on $\C \times \R$ that are
  holomorphic in $\C$ and locally constant in~$\R$.
\end{enumerate}
\end{theorem}

\noindent
The product topology is not decorative.  It is the entire content
of the theorem: the two colours of the Swiss-cheese operad arise
because the two factors of $\C \times \R$ carry different
topologies (holomorphic versus locally constant).

\medskip

The recognition theorem converts geometric data (a factorisation
algebra on a $3$-manifold) into algebraic data (an operad
algebra).  Without it, the Swiss-cheese structure would be a
presentation, not a property.  With it, the two-coloured operad
is the geometry.


% =================================================================
\section{One algebra, three objects}
\label{sec:three-objects}
% =================================================================

From a single chiral algebra $\cA$, the bar construction
produces three objects.  They are not the same.

\medskip

\noindent
\textbf{(i)} The \textbf{boundary algebra} $\cA$ itself.  This
lives on the boundary $\{0\} \times \R$ of the half-space
$\C_{\geq 0} \times \R$.  It is a chiral algebra: a vertex
algebra, an $\Einf$-chiral structure.  Its data are the OPE
coefficients $c_n^{ab}$ and the conformal weights.

\medskip

\noindent
\textbf{(ii)} The \textbf{bar coalgebra} $\barB(\cA)$.  This is a
factorisation coalgebra on $\Ran(X)$, carrying both the
differential and the coproduct.  It classifies \textit{twisting
morphisms}: universal couplings between $\cA$ and its Koszul
dual~$\cA^!$.  On the Koszul locus,
$\MC(\Hom(\barB(\cA), \cA^!)) \simeq \Perf(\cA^!)$.  This is
the bar complex's native role: a representability theorem for
deformations.

\medskip

\noindent
\textbf{(iii)} The \textbf{chiral derived center}
\[
\Zder(\cA)
\;:=\;
H^*\bigl(C^\bullet_{\mathrm{ch}}(\cA, \cA),\, \delta\bigr).
\]
This is the universal bulk algebra: it classifies observables
of the $3$d theory that act on the boundary.  At the level of
Hochschild cohomology:
$\Zder(\cA) \simeq \HH^\bullet_{\mathrm{ch}}(\cA)$, the chiral
Hochschild cochains.

\sep

The identification ``bar = bulk'' is the most dangerous false
equation in three-dimensional holomorphic-topological field theory.

The bar coalgebra classifies \textit{couplings} between a boundary
and its dual.  The derived center classifies \textit{operators} that
the boundary does not know it has.  These are different functors
applied to different inputs with different outputs.  In a
sentence: bar-cobar is about twisting morphisms; the derived
center is about the acting-on relation.

\begin{center}
\small
\begin{tabular}{@{}lll@{}}
\textbf{Object} & \textbf{Classifies} & \textbf{Functor} \\
\hline
$\barB(\cA)$ & twisting morphisms $\cA \leadsto \cA^!$
  & bar (coalgebra) \\[3pt]
$\Zder(\cA)$ & bulk operators acting on $\cA$
  & chiral Hochschild (algebra) \\[3pt]
$\Omega(\barB(\cA))$ & nothing new: $\cA$ itself
  & cobar (recovers input)
\end{tabular}
\end{center}

\medskip

The cobar of the bar is $\cA$ (Theorem~B of Volume~I).  The
Verdier dual of the bar is $\barB(\cA^!)$ (Theorem~A).  These
are three different operations, producing three different
objects.


% =================================================================
\section{The Swiss-cheese pentagon}
\label{sec:pentagon}
% =================================================================

The $\SCchtop$ operad has two colours: \textit{closed}
(holomorphic, on $\FM_k(\C)$) and \textit{open} (topological,
on $\Eone(m)$).  The structure maps are:

\begin{center}
\small
\renewcommand{\arraystretch}{1.15}
\begin{tabular}{@{}llll@{}}
\textbf{Map} & \textbf{Inputs} & \textbf{Output}
  & \textbf{Interpretation} \\
\hline
$\mu_{k,0}$ & $k$ closed & closed
  & Gerstenhaber bracket \\[3pt]
$\mu_{0,m}$ & $m$ open & open
  & $\Ainf$ products \\[3pt]
$\mu_{k,m}$ & $k$ closed, $m$ open & open
  & brace operations \\[3pt]
$\mu_{k,0}^{\mathrm{oc}}$ & $k$ open & closed
  & \textbf{EMPTY} \\[3pt]
$\mu_{k,m}^{\mathrm{mix}}$ & mixed & mixed
  & SC composition
\end{tabular}
\end{center}

\medskip

The fourth row is the content.  No open inputs produce closed
outputs.  This is not a convention but a structural constraint:
the open colour lives on the boundary $\R$, the closed colour in
the bulk $\C \times \R$, and bulk interactions restrict to
boundaries but not conversely.  Bulk-to-boundary is a restriction
functor; boundary-to-bulk would require extension, which the
holomorphic direction forbids.

The bar complex makes this visible.  The bar \textit{differential}
(closed colour) is the OPE residue extraction: a collision in the
$\C$-plane.  The bar \textit{coproduct} (open colour) is the
deconcatenation: a separation along~$\R$.  Applying the coproduct
and then the differential is a brace operation: splitting a bar
element, then colliding within each piece.  Applying the
differential and then the coproduct is an SC composition:
colliding first, then splitting.  The Stasheff pentagon for the
$\Ainf$ structure and the Gerstenhaber bracket for the closed
colour are the same identity, read on the two factors of
$\FM_k(\C) \times \Eone(m)$.

The Swiss-cheese operad is Koszul.

\begin{theorem}[Homotopy-Koszulity; proved]
\label{thm:homotopy-koszul-wn}
The chiral-topological Swiss-cheese operad $\SCchtop$ is
homotopy-Koszul.
\end{theorem}

\noindent
The proof:
(1)~the classical Swiss-cheese operad is Koszul
(Livernet, Voronov, Ginzburg--Kapranov);
(2)~Kontsevich formality gives a quasi-isomorphism
$\Phi\colon \SCchtop \to \mathrm{SC}$;
(3)~bar-cobar preserves quasi-isomorphisms; two-out-of-three.

The consequences are unconditional.  Bar-cobar on
$\SCchtop$-algebras is a Quillen equivalence, not merely an
adjunction.  Lines are $\cA^!$-modules.  The dg-shifted Yangian
is defined without hypothesis.


% =================================================================
\section{Descent to the classical shadow}
\label{sec:pva-descent}
% =================================================================

The holomorphic Swiss-cheese algebra, projected to cohomology,
produces a Poisson vertex algebra (PVA).  The descent is
controlled by six axioms D1--D6, each proved by a geometric
argument on $\FM_k(\C)$.

\begin{enumerate}[label=D\arabic*.,nosep]
\item \textit{Leibniz rule}:
  $\{a_\lambda bc\} = \{a_\lambda b\}c + b\{a_\lambda c\}$.
\item \textit{Sesquilinearity}:
  $\{\partial a_\lambda b\} = -\lambda\{a_\lambda b\}$.
\item \textit{Skew-symmetry}:
  $\{b_\lambda a\} = -\{a_{-\lambda-\partial} b\}$.
\item \textit{Jacobi identity}:
  $\{a_\lambda\{b_\mu c\}\} - \{b_\mu\{a_\lambda c\}\}
  = \{\{a_\lambda b\}_{\lambda+\mu} c\}$.
\item \textit{Commutativity}:
  $ab = ba$ on cohomology.
\item \textit{Unit}:
  $\{a_\lambda 1\} = 0$.
\end{enumerate}

D2--D5 are proved by the exchange cylinder and the three-face
Stokes identity on $\FM_3(\C)$.  The exchange cylinder is the
bordism
$\FM_2(\C) \times [0,1] \to \FM_3(\C)$
that swaps two points while a third observes; the three-face
identity
\[
d\omega_3 \;=\; \omega_{12,3} - \omega_{13,2} + \omega_{1,23}
\]
on $\partial\FM_3(\C)$ gives skew-symmetry and Jacobi
simultaneously.  D6 uses the factorisation of $H_3(a)$ and the
contractibility of the disc.

The bridge operates in both directions.  F4 (operad $\Rightarrow$
axioms): the $\SCchtop$ structure maps project to PVA operations.
F5 (axioms $\Rightarrow$ operad): a PVA satisfying D1--D6 lifts
to an $\SCchtop$-algebra.  The lift exists because $\SCchtop$ is
homotopy-Koszul: the cobar of the Koszul dual cooperad provides
a cofibrant resolution, and the PVA data define a map from this
resolution.

\sep

The PVA on cohomology $H^\bullet(\cA, Q)$ is $(-1)$-shifted:
the $\lambda$-bracket has shifted parity relative to classical
PVA conventions.  This is not an anomaly.  It is the degree shift
in the desuspension $s^{-1}$ of the bar construction.


% =================================================================
\section{The coproduct}
\label{sec:coproduct}
% =================================================================

In gauge theory (Chern--Simons with boundary algebra
$\widehat\fg_k$), the gluon $J^a$ has a primitive coproduct:
\[
\Delta_z(J^a) \;=\; J^a(w{+}z) \otimes 1
\;+\; 1 \otimes J^a(w).
\]
But the Sugawara composite
$T_{\mathrm{aff}} = \frac{1}{2(k+h^\vee)}\sum_a
\normord{J^a J^a}$ has a nontrivial coproduct.  The cross term
\[
\Delta_z(T_{\mathrm{aff}})
\;=\;
\tau_z(T_{\mathrm{aff}}) \otimes 1
\;+\; 1 \otimes T_{\mathrm{aff}}
\;+\;
\frac{1}{k+h^\vee}\,\Omega(w{+}z,\, w),
\]
where $\Omega = \sum_a J^a \otimes J^a$ is the split Casimir,
is a production vertex: a gluon bound state can dissociate into
its constituents.  The gluon IS the connection.  It splits
because the connection is a primitive field.

\sep

In gravity (Virasoro boundary), the stress tensor $T$ is the
graviton.  The coproduct is primitive at every arity:
\begin{equation}\label{eq:gravity-primitive-wn}
\Delta_z^W(x) \;=\; \tau_z(x) \otimes 1 \;+\; 1 \otimes x
\end{equation}
for every element $x$ of the $\cW$-algebra, with no corrections
whatsoever.

The mechanism is a ghost-number obstruction intrinsic to
Drinfeld--Sokolov reduction; the proof is given in
\S\ref{sec:primitive-coproduct}
(Theorem~\ref{thm:gravitational-primitivity-wn}).

In gauge theory, the gluon IS the connection.  In gravity, the
graviton IS the curvature.  A connection splits (it is the
primitive field of the theory).  A curvature does not (it is a
composite, and the ghost-number filtration forbids the splitting).
This is the algebraic content of the statement that gravity does
not produce particles; it only braids them.

\begin{center}
\small
\renewcommand{\arraystretch}{1.15}
\begin{tabular}{@{}llll@{}}
 & \textbf{Chern--Simons} & \textbf{Gravity} \\
\hline
Boundary algebra & $\widehat\fg_k$ & $\mathrm{Vir}_c$ \\[3pt]
Primitive field & $J^a$ (connection) & $T$ (curvature) \\[3pt]
$\Ainf$ tower & formal ($m_k^H = 0$ for $k \ge 3$, class L)
  & infinite ($m_k \neq 0\ \forall k$, class M) \\[3pt]
Coproduct $\Delta_z$ & nontrivial (Casimir cross-term)
  & primitive (ghost obstruction) \\[3pt]
$r$-matrix & $\Omega_\fg/z$ & $(c/2)/z^3 + 2T/z$ \\[3pt]
Shadow depth & $3$ & $\infty$
\end{tabular}
\end{center}

\medskip

All nonlinearity of gravity resides in the $\Ainf$ products
$\{m_n^W\}_{n \geq 3}$ and in the collision residue
$r^{\mathrm{Vir}}(z) = (c/2)/z^3 + 2T/z$.  The coproduct carries
no information.  This is the central algebraic fact about
three-dimensional quantum gravity.


% =================================================================
\section{Five worked examples}
\label{sec:five-examples}
% =================================================================

\subsection*{Chern--Simons}

Gauge group $G$, level $k$.  Boundary algebra:
$\widehat\fg_k$.  Class L, shadow depth $3$.
$r$-matrix: $r(z) = \Omega_\fg/z$ (Casimir, single pole).
Coproduct: nontrivial (Casimir cross-term on composites).
$\kappa = \dim(\fg)(k+h^\vee)/(2h^\vee)$.
Complementarity: $\kappa + \kappa^! = 0$ by Feigin--Frenkel.

The entire $\Ainf$ structure is formal: $m_k = 0$ for $k \geq 3$
on cohomology.  The bar complex is quadratic.  The Koszul dual
$\widehat\fg_{-k-2h^\vee}$ is another affine algebra at the
reflected level.  Lines: highest-weight modules for the dual.
On the integrable locus ($k \in \Z_{\geq 0}$), the line category
is $\mathrm{Rep}_q(\fg)$ with $q = e^{\pi i/(k+h^\vee)}$,
recoverable by the Kazhdan--Lusztig equivalence.

\subsection*{Gravity}

Boundary algebra: $\mathrm{Vir}_c$.  Class M, shadow depth
$\infty$.  $r$-matrix: $r(z) = (c/2)/z^3 + 2T/z$ (two poles,
quasi-trigonometric).  Coproduct: primitive.
$\kappa = c/2$; $\kappa^! = (26-c)/2$;
$\kappa + \kappa^! = 13 \neq 0$.

The quartic OPE pole forces an infinite $\Ainf$ tower.  Every
$m_k$ is nonzero for generic $c$, computed by a wheel diagram on
$K_k$ with $k-1$ cubic vertices.  The quartic contact invariant
$Q^{\mathrm{ct}}_{\mathrm{Vir}} = 10/[c(5c+22)]$
has poles at $c = 0$ and $c = -22/5$: the Kac determinant zeros
at weight~$4$.  Self-duality at $c = 13$, not $c = 26$.

\subsection*{Twisted holography}

Costello--Gaiotto.  Boundary algebra: $V_k(\fg)$ (affine
Kac--Moody).  Class L, shadow depth $3$.  Coproduct: nontrivial.
The holographic dictionary is the Koszul triangle read through
Volume~I's five theorems: bar complex exists (Theorem~A), boundary
reconstructs bulk (Theorem~B), Neumann/Dirichlet decomposition
(Theorem~C), genus tower by $\kappa$ (Theorem~D), BV origin of
Hochschild complex (Theorem~H).

\subsection*{M2-brane}

From the $\cN = 6$ ABJM theory.  Boundary algebra: supersymmetric
current algebra.  Coproduct: nontrivial (fundamental fields split).
The M2 example has no DS reduction: the $\cW$-algebra structure is
absent, and the coproduct obstruction does not apply.  This is a
gauge theory, not a gravity theory, and the coproduct confirms it.

\subsection*{M5-brane / 6d $(2,0)$}

From the $\cN = (2,0)$ theory on $\C \times \R \times S^3$.
Boundary algebra: $\cW_N$ at large $N$.  Class M, shadow depth
$\infty$.  Coproduct: primitive (by the same ghost-number
argument as Virasoro, since $\cW_N$ arises by DS reduction).

\begin{center}
\small
\renewcommand{\arraystretch}{1.15}
\begin{tabular}{@{}lccccc@{}}
 & CS & Gravity & Twisted hol. & M2 & M5 \\
\hline
Class & L & M & L & L & M \\
$r_{\max}$ & $3$ & $\infty$ & $3$ & $3$ & $\infty$ \\
$\Delta_z$ & nontrivial & primitive & nontrivial
  & nontrivial & primitive \\
DS origin & no & yes & no & no & yes \\
$\kappa + \kappa^!$ & $0$ & $13$ & $0$ & $0$ & $\varrho K_N$
\end{tabular}
\end{center}

\medskip

The dichotomy is clean: DS reduction produces primitive
coproducts and infinite shadow depth.  Its absence permits
nontrivial coproducts and finite termination.


% =================================================================
\section{The annulus}
\label{sec:annulus}
% =================================================================

Cut the circle $S^1$ by two arcs.  Each arc is an interval
$I \cong \R$.  The factorisation homology of the open-sector
category $\cC_{\mathrm{op}}$ on $S^1$ is computed by excision:
\[
\int_{S^1} \cC_{\mathrm{op}}
\;\simeq\;
\cC_{\mathrm{op}}(I_1)
\;\otimes_{\cC_{\mathrm{op}}(I_1 \cap I_2)}^{\mathbb{L}}\;
\cC_{\mathrm{op}}(I_2).
\]
The two-arc cover produces the two-sided bar complex, which is
the cyclic bar complex.  In the homotopy category:
\begin{equation}\label{eq:annulus-trace-wn}
\int_{S^1} \cC_{\mathrm{op}}
\;\simeq\;
\HH_*(\cC_{\mathrm{op}}).
\end{equation}

This is the annulus trace.  The annulus $S^1 \times [0,1]$ is the
simplest bordered surface with one boundary component.  Its
factorisation homology is the Hochschild homology of the
open-sector category.  The trace is the clutching operation: it
restores the tangential basepoint $\tau_p$ to close the interval
$I_p$ back to a circle.

\sep

For affine Kac--Moody at integrable level $k$:
$\HH_0(\cC_{\mathrm{op}})$ is spanned by the characters
$\chi_j(\tau) = \mathrm{tr}_{L_j}(q^{L_0 - c/24})$
for $j = 0, 1, \ldots, k$.  The modular $S$-matrix acts on
this space.  Modular invariance of the partition function
$Z(\tau) = \sum_j |\chi_j|^2$ is the statement that the trace
is $SL(2,\Z)$-equivariant.

\medskip

The annulus trace is the \textbf{first modular shadow of the
open sector}.  It is the genus-$1$ datum, the simplest case of
the trace-and-clutching construction that produces genus-$g$
amplitudes from the open category.  Modularity is not an axiom
on the closed algebra.  It is a property of the trace on the
open sector.


% =================================================================
\section{The line category}
\label{sec:line-category}
% =================================================================

On the chirally Koszul locus:
\begin{equation}\label{eq:lines-koszul-wn}
\cC_{\mathrm{line}} \;\simeq\; \cA^!\text{-}\mathbf{mod}.
\end{equation}
Line operators are modules for the Koszul dual.

The proof is Koszul duality on the open colour.  The bar complex
$\barB(\cA)$ is a dg coalgebra; its linear dual
$\barB(\cA)^\vee = \barB(\cA^!)$ (Verdier duality, Theorem~A
of Volume~I) is a dg algebra.  Maurer--Cartan deformations of
the universal twisting morphism
$\alpha\colon \barB(\cA) \to \cA^!$ are in bijection with
$\cA^!$-module structures:
\[
\MC\bigl(\Hom(\barB(\cA),\, \cA^!)\bigr)
\;\simeq\;
\Perf(\cA^!).
\]
A line is not a module that happens to be placed along $\R$.
A module is a line: the representation category of $\cA^!$ IS the
category of boundary conditions for the theory restricted
to~$\R$.

\sep

The meromorphic tensor structure on $\cC_{\mathrm{line}}$ comes
from the $r$-matrix.  Two line operators at spectral
parameters $z_1$ and $z_2$ tensor along the open direction with
braiding
\[
R(z_1 - z_2)\colon M_1 \otimes M_2
\;\xrightarrow{\sim}\;
M_2 \otimes M_1,
\]
where $R(z)$ is the quantum $R$-matrix built from the bulk-boundary
composition (not from a quantum group, though they agree on the
evaluation locus).  For Chern--Simons:
$R(z) = 1 + \hbar\,\Omega_\fg/z + O(\hbar^2)$.  For gravity:
$R(z) = 1 + \hbar\,((c/2)/z^3 + 2T/z) + O(\hbar^2)$.

The $r$-matrix is the genus-$0$, arity-$2$, ordered shadow of
$\Theta_\cA$:
\[
r(z) \;=\; \Res^{\mathrm{coll}}_{0,2}(\Theta_\cA).
\]
Volume~I computes the collision residue.  Volume~II reads the
result as a spectral braiding on lines.


% =================================================================
\section{The benchmark}
\label{sec:benchmark}
% =================================================================

$V_k(\fsl_2)$ on $(\bP^1, \{0, \infty\})$.  This is the
simplest non-trivial example: one curve with two punctures, one
simple Lie algebra of rank $1$, one level parameter $k$.  Every
object in the modular Koszul datum can be computed.

\medskip

\noindent
\textbf{Tangential log curve.}
$(X, D, \tau) = (\bP^1, \{0,\infty\}, \{\tau_0, \tau_\infty\})$
with tangential basepoints in the positive-real direction.  The
boundary intervals $I_0, I_\infty \cong \R$ are the time axes at
each puncture.

\medskip

\noindent
\textbf{Open sector.}
$\cC_{\mathrm{op}} = V_k(\fsl_2)\text{-mod}$ at each puncture.
At integrable level $k$, this is $\mathrm{Rep}_q(\fsl_2)$ with
$q = e^{i\pi/(k+2)}$, a semisimple category with $k+1$ simple
objects $L_0, L_1, \ldots, L_k$.

\medskip

\noindent
\textbf{Boundary algebra.}
$A_b = V_k(\fsl_2)$.  The $\lambda$-bracket:
$\{J^a{}_\lambda J^b\} = f^{ab}{}_c J^c + k\kappa^{ab}\lambda$.

\medskip

\noindent
\textbf{Koszul dual.}
$V_k(\fsl_2)^! = V_{-k-4}(\fsl_2)$.  The Feigin--Frenkel
involution $k \mapsto -k - 2h^\vee = -k - 4$.

\medskip

\noindent
\textbf{$r$-matrix.}
$r(z) = \Omega_{\fsl_2}/z = (J^a \otimes J_a)/z$, a single-pole
solution of the classical Yang--Baxter equation.

\medskip

\noindent
\textbf{Modular characteristic.}
$\kappa = 3(k+2)/4$.  Complementarity: $\kappa + \kappa^! = 0$.

\medskip

\noindent
\textbf{Derived center.}
$\Zder(V_k(\fsl_2)) \simeq \HH^\bullet_{\mathrm{ch}}(V_k(\fsl_2))$.
At the critical level $k = -2$: functions on $\fsl_2$-opers.

\medskip

\noindent
\textbf{Line category.}
$\cC_{\mathrm{line}} \simeq V_{-k-4}(\fsl_2)\text{-mod}$.
At integrable levels: $\mathrm{Rep}_{q^{-1}}(\fsl_2)$.

\medskip

\noindent
\textbf{Annulus trace.}
$\HH_0(\cC_{\mathrm{op}}) = \bigoplus_{j=0}^k \C \cdot \chi_j$.
The modular $S$-matrix:
\[
S_{jj'} \;=\; \sqrt{\frac{2}{k+2}}\,
\sin\!\left(\frac{\pi(j+1)(j'+1)}{k+2}\right).
\]
At $k = 1$: $S = \frac{1}{\sqrt{2}}
\begin{psmallmatrix} 1 & 1 \\ 1 & -1 \end{psmallmatrix}$,
the Hadamard matrix.

\medskip

\noindent
\textbf{Factorisation homology.}
$\int_{\C^\times}\!\cC_{\mathrm{op}}
\simeq \mathrm{CB}_{0,2}(V_k(\fsl_2))$:
genus-$0$ conformal blocks with two insertions.  For simple
modules $L_j$ and $L_{j'}$:
$\dim \mathrm{CB}_{0,2}(L_j, L_{j'}) = \delta_{j,j'}$.

\medskip

Every entry is proved.  The benchmark is the frame atom of
Volume~II: simple enough to compute completely, rich enough to
exhibit every feature of the theory.


% =================================================================
\section{Modularity last}
\label{sec:modularity-last}
% =================================================================

The standard approach to modular invariance begins with the
closed-string partition function and imposes modular invariance
as an axiom.  The Swiss-cheese framework reverses this.

Modularity is a derived property of the trace on the open sector
(\S\ref{sec:annulus}).

\sep

The modular Swiss-cheese operad extends $\SCchtop$ by replacing
trees with stable graphs carrying two edge colours: closed
edges (genus direction, carrying $\hbar$-weight, recording
non-separating degenerations) and open edges (boundary direction,
carrying insertion data, recording separations along $\R$).

The relative Feynman transform $F_{\mathrm{rel}}\cO$ sums over
stable graphs with both edge colours.  The differential contracts
one edge; $d^2 = 0$ is the codimension-$2$ boundary cancellation
on the moduli space of bordered Riemann surfaces.

At genus $g \geq 1$, the curvature
$\dfib^{\,2} = \kappa \cdot \omega_g$ forces a passage from
ordinary to curved Swiss-cheese algebras.  The modular MC element
\begin{equation}\label{eq:oc-mc-wn}
\Theta^{\mathrm{oc}}_\cA
\;:=\;
\Theta_\cA + \sum_{j} \mu^{\cM_j}
\end{equation}
packages both the bar-intrinsic MC element $\Theta_\cA$ (from
Volume~I) and the $\Ainf$-module structures $\mu^{\cM_j}$ at
punctures.  Its MC equation decomposes into four components:
\begin{alignat*}{2}
&\text{pure closed:} &&\quad d_0\,\Theta
  + \tfrac{1}{2}[\Theta, \Theta] = 0 \\[3pt]
&\text{pure open:} &&\quad \partial\mu
  + \mu \circ \mu = 0 \\[3pt]
&\text{mixed:} &&\quad d_0\,\mu
  + [\Theta, \mu] + \cdots = 0 \\[3pt]
&\text{clutching:} &&\quad \text{(genus recursion)}
\end{alignat*}
The clutching component is the genus recursion from Volume~I.
Each genus is determined by all lower genera.  No parameter
is free.

\sep

The annulus trace (the $S^1$ factorisation homology of
$\cC_{\mathrm{op}}$) produces the genus-$1$ modular datum.
The clutching of $g$ annuli produces the genus-$g$ datum.
Modularity emerges from the composition law of traces, not
from an a priori constraint on the closed algebra.

The programme status:

\begin{center}
\small
\renewcommand{\arraystretch}{1.15}
\begin{tabular}{@{}ll@{}}
\textbf{Result} & \textbf{Status} \\
\hline
Annulus trace $\simeq \HH_*(\cC_{\mathrm{op}})$ & proved \\[3pt]
Full modular cooperad & programme (in progress) \\[3pt]
Modular Swiss-cheese $d^2 = 0$ & proved (by $\partial^2 = 0$ on
  $\overline\cM_{g,n}$) \\[3pt]
Curvature $= \kappa \cdot \omega_g$ & proved (Volume~I)
\end{tabular}
\end{center}


% =================================================================
\section{The gravitational coproduct is primitive}
\label{sec:primitive-coproduct}
% =================================================================

\begin{theorem}[Gravitational primitivity; proved]
\label{thm:gravitational-primitivity-wn}
For every principal Drinfeld--Sokolov reduction of every simple
Lie algebra, the transferred dg-shifted Yangian coproduct on the
$\cW$-algebra is strictly primitive:
\[
\Delta_z^W(x)
\;=\;
\tau_z(x) \otimes 1
\;+\; 1 \otimes x
\]
for every $x$, at every arity, with no corrections.
\end{theorem}

\begin{proof}
Three lines.

The homological perturbation lemma expresses the transferred
coproduct as a sum over decorated trees.  Each tree takes
$\cW$-algebra inputs (ghost number $0$), applies the BRST
homotopy $h$ (which raises or lowers ghost number by $1$), and
outputs a tensor product.  Every tree with at least one internal
vertex produces an output at ghost number $\pm 1$.  The projection
$p \otimes p$ onto ghost-number-$0$ cohomology annihilates both
signs.  The only surviving contribution is the trivial tree
(no internal vertices), which gives the primitive term
$\tau_z(x) \otimes 1 + 1 \otimes x$.
\end{proof}

\sep

The proof works for every principal DS reduction because the
ghost-number obstruction is intrinsic to the BRST structure, not
to the specific Lie algebra.  The ghost field $c \in \fn_+^*$
has ghost number $+1$; the antighost $b \in \fn_+$ has ghost
number $-1$; the BRST homotopy $h$ carries ghost number $\pm 1$;
and the projection $p$ annihilates everything outside ghost
number $0$.

The physical content is the no-particle-production principle of
\S\ref{sec:coproduct}: the graviton is a composite (the stress
tensor, built from the Sugawara construction) and not a fundamental
field, so the ghost-number filtration forbids the splitting at
every order.

For Chern--Simons, the fundamental field $J^a$ has ghost number
$0$, so the BRST homotopy produces ghost-number-$0$ outputs.
The Casimir cross-term survives projection.  The gluon splits.

The asymmetry between gauge theory and gravity is an asymmetry
between fundamental and composite fields, detected by the ghost
number.


% =================================================================
\section{The four-stage architecture}
\label{sec:four-stages}
% =================================================================

The open/closed programme has four stages.  Each depends on its
predecessor.

\medskip

\noindent
\textbf{Stage 1: local one-colour.}
The input is a single chiral algebra $\cA$.  Outputs:

\begin{enumerate}[label=\textup{(\alph*)},nosep]
\item The bar coalgebra $\barB(\cA)$ exists as a factorisation
  coalgebra (Volume~I, Theorem~A).
\item On the Koszul locus, MC couplings
  $\simeq \Perf(\cA^!)$ (Koszul duality on the open colour).
\item The bulk is the chiral derived center:
  $\Zder(\cA) \simeq \HH^\bullet_{\mathrm{ch}}(\cA)$
  (boundary-linear exact theorem).
\item Braces: $C^\bullet_{\mathrm{ch}}(\cA, \cA)$ carries a brace
  dg-algebra structure from the $\SCchtop$ action.
\end{enumerate}

No global factorisation or modularity is used.  Everything here
is the universal property of the bar construction and the
homotopy-Koszulity of $\SCchtop$.  Proved.

\medskip

\noindent
\textbf{Stage 2: the open primitive.}
The primitive datum is the open/closed factorisation dg-category
$\cC$ on a pointed curve with tangential basepoints
$(X, D, \tau)$.

\begin{enumerate}[label=\textup{(\alph*)},nosep]
\item For any compact generator $b \in \cC(J)$, the boundary
  algebra $A_b = \RHom_{\cC(J)}(b,b)$ is an $\Ainf$ chiral
  algebra.
\item Different compact generators produce Morita-equivalent
  boundary algebras: the category $\cC$ is the invariant,
  $A_b$ is a chart.
\item The bar complex $\barB(A_b)$ encodes the twisting data
  of $\cC$, and bar-cobar is the representability theorem for
  deformations of $b$ inside $\cC$.
\end{enumerate}

Proved.

\medskip

\noindent
\textbf{Stage 3: globalise.}
Replace $(X, D, \tau) = (\C, \{0\}, \tau_0)$ by a general
tangential log curve.  The real oriented blowup
$\pi\colon \widetilde{X}_D \to X$ replaces each puncture $p_i$
by a circle $S^1_{p_i}$; the tangential basepoint $\tau_{p_i}$
cuts this circle into an interval $I_{p_i} \cong \R$.

The mixed configuration spaces
$\Conf^{\mathrm{oc}}_{I,\mathbf{m}}(X, D, \tau)
= \Conf_I(X \setminus D) \times \prod_p \Conf^<_{m_p}(\R)$
parametrise bulk insertions and boundary insertions together.
The mixed Ran space is their colimit.

Local-global bridge: the local one-colour theorems of Stage 1
extend to global statements via descent along the mixed Ran space.
Proved locally, the global extension requires compact generation
and derived-center hypotheses.

\medskip

\noindent
\textbf{Stage 4: modularise.}
The trace operation on $\cC_{\mathrm{op}}$ (closing an interval
to a circle) produces genus-$g$ data from genus-$0$ data.

\begin{enumerate}[label=\textup{(\alph*)},nosep]
\item The annulus trace
  $\mathrm{Tr}_\cA \simeq \HH_*(\cC_{\mathrm{op}})$ is the
  first modular shadow.  Proved.
\item The full modular cooperad structure on the trace (the
  composition law for genus-$g$ amplitudes via clutching of $g$
  handles) is a programme.
\end{enumerate}

\sep

The architecture is sequential.  Stage $k$ depends on Stage $k-1$.
One should never discuss Stage 4 before establishing Stage 1.
The order is: objects first, then categories, then globalisation,
then modularity.


% =================================================================
\section{The Koszul triangle}
\label{sec:koszul-triangle}
% =================================================================

On the chirally Koszul locus, the corrected bulk-boundary-line
triangle is:
\begin{equation}\label{eq:triangle-wn}
A_{\mathrm{bulk}}
\;\simeq\;
\Zder(B_\partial)
\;\simeq\;
\Zder(\cC_{\mathrm{line}})
\;\simeq\;
\HH^\bullet_{\mathrm{ch}}(\cA^!),
\qquad
\cC_{\mathrm{line}} \simeq \cA^!\text{-}\mathbf{mod}.
\end{equation}

Three vertices:
\begin{itemize}[nosep]
\item The \textbf{bulk} is
  $A_{\mathrm{bulk}} \simeq \cO(T^*[-1]\cL)$: polyvector fields
  on the Lagrangian.  Note: the restriction
  $\cO(\cM)|_\cL = \cO(\cL)$ is the boundary, not the bulk.
  The bulk retains the full $(-1)$-shifted cotangent fibre.
\item The \textbf{boundary} is $B_\partial = \cO(\cL)$: functions
  on the Lagrangian.
\item The \textbf{line category} is
  $\cC_{\mathrm{line}} = \mathfrak{S}_b\text{-mod}$: modules for
  the self-intersection algebra
  $\mathfrak{S}_b = \cL \times_\cM \cL$.
\end{itemize}

The equivalence $\Zder(B_\partial) \simeq A_{\mathrm{bulk}}$ is
the Hochschild--Kostant--Rosenberg theorem: Hochschild cochains of
$\cO(\cL)$ are polyvector fields on $T^*[-1]\cL \simeq \cM|_\cL$,
by the Lagrangian condition
$T_\cL\cM / T\cL \simeq T^*\cL[-1]$.

The equivalence $\cC_{\mathrm{line}} \simeq \cA^!\text{-mod}$ is
the convolution-module model: the self-intersection algebra on
$\mathfrak{S}_b$ is the Koszul dual, and its modules are lines.

Both equivalences require the Lagrangian condition.  This is the
geometric origin of the algebraic fact that the boundary is NOT
the bulk and lines are NOT bulk modules.

\sep

The bar complex presents the Swiss-cheese algebra, as the
Steinberg variety presents the Hecke algebra.

This is not analogy.  Both are algebro-geometric presentations of
operadic structures via configuration-space resolutions.  The
Steinberg variety $\mathfrak{S}_b = \cL \times_\cM \cL$ is the
self-intersection of the Lagrangian $\cL$ in the moduli stack
$\cM$; its convolution algebra is the Hecke algebra.  The bar
complex $\barB(\cA) = T^c(s^{-1}\bar\cA)$ with differential from
collision residues and coproduct from deconcatenation is the
Swiss-cheese algebra on $\FM_k(\C) \times \Conf_k(\R)$; its
convolution with the Koszul dual is the line category.  The
two become the same construction at the level of factorisation
homology.


% =================================================================
\section{What we do not know}
\label{sec:what-we-do-not-know}
% =================================================================

\subsection*{Three falsification tests}

\textbf{Test 1.}
The global bulk-boundary-line triangle
(equation~\eqref{eq:triangle-wn}) holds unconditionally in the
boundary-linear exact sector.  Whether it extends to all HT
theories depends on two unverified hypotheses: compact generation
of the open-sector category, and the derived center being a
quasi-isomorphism to the bulk.  These hypotheses are verified for
Chern--Simons and for inputs from the Costello--Davison--Gaiotto
framework.  They are not verified for gravity.  A counterexample
would mean that the bulk of a gravitational theory is not the
derived center of its boundary: the simplest open/closed
identification would fail.

\medskip

\textbf{Test 2.}
The gravitational primitivity theorem
(\S\ref{sec:primitive-coproduct}) is proved for principal DS
reductions.  For non-principal DS reductions (hook-type in type A,
arbitrary nilpotent $f$ in general), the ghost-number argument
applies verbatim, but the coproduct on the BRST cohomology has
not been computed.  Hook-type in type A is the proved corridor.
Full non-principal primitivity is conjectural.

\medskip

\textbf{Test 3.}
The modular cooperad structure on the annulus trace (Stage 4 of
the architecture) is a programme, not a theorem.  The annulus
trace itself is proved: $\mathrm{Tr}_\cA \simeq
\HH_*(\cC_{\mathrm{op}})$.  The full cooperad structure (clutching
laws for genus $g \geq 2$, compatibility with the modular operad
on $\overline\cM_{g,n}$) requires the bordered FM compactification
and its log-FM extension, which is the frontier of Mok's
programme.

\subsection*{Programme status ledger}

\begin{center}
\small
\renewcommand{\arraystretch}{1.15}
\begin{tabular}{@{}llc@{}}
\textbf{Component} & \textbf{What it requires} & \textbf{Status} \\
\hline
$\SCchtop$ homotopy-Koszulity & Kontsevich + Livernet &
  \textsc{green} \\[3pt]
Recognition theorem & Product Weiss descent &
  \textsc{green} \\[3pt]
PVA descent D1--D6 & FM Stokes calculus &
  \textsc{green} \\[3pt]
Gravitational primitivity & Ghost-number obstruction &
  \textsc{green} \\[3pt]
Bulk $=$ derived center (boundary-linear) &
  Stage 1 local theorem & \textsc{green} \\[3pt]
Bulk $=$ derived center (global) &
  Compact generation & \textsc{amber} \\[3pt]
Lines $= \cA^!$-mod & Koszul locus &
  \textsc{green} \\[3pt]
Annulus trace & $S^1$ excision &
  \textsc{green} \\[3pt]
Full modular cooperad & Bordered FM + log extension &
  \textsc{amber} \\[3pt]
BV/BRST $=$ bar at $g \geq 1$ & Genus-$g$ QME
  & \textsc{red}
\end{tabular}
\end{center}

\subsection*{Five false ideas to dismiss}

\begin{enumerate}[nosep]
\item ``The bar complex is the bulk algebra.''  The bar complex
  classifies twisting morphisms.  The bulk is the derived center.
  These are different objects produced by different functors.
\item ``The boundary algebra determines the theory.''  The boundary
  algebra is a chart of the open-sector category, not its essence.
  Different compact generators produce Morita-equivalent charts.
  The category is the invariant.
\item ``Modularity is a property of the closed algebra.''  It is a
  property of the trace on the open sector.  The closed algebra
  inherits modularity from the clutching operation, not the reverse.
\item ``Gravity and gauge theory differ only in the central
  charge.''  They differ in the coproduct.  Gauge theory has
  nontrivial $\Delta_z$; gravity has primitive $\Delta_z$.  The
  ghost-number obstruction is structural, not parametric.
\item ``The Koszul dual lives in the bulk.''  It lives on the
  boundary: $\cA^!$ is the boundary algebra for line operators,
  not a bulk observable.
\end{enumerate}


% =================================================================
\section{The completed modular Koszul datum}
\label{sec:platonic}
% =================================================================

From a single chiral algebra $\cA$, the completed modular Koszul datum
assembles eight components:
\[
\Pi^{\mathrm{oc}}_X(\cA)
\;=\;
\bigl(\barB_X,\;\; \Omega_X,\;\; \Theta_\cA,\;\;
\Zder,\;\; \mathrm{Tr}_\cA,\;\;
\HH_*,\;\; \mathfrak{R}^{\mathrm{oc}}_\bullet,\;\;
\nabla^{\mathrm{hol}}\bigr).
\]

\begin{center}
\small
\renewcommand{\arraystretch}{1.15}
\begin{tabular}{@{}lll@{}}
\textbf{Component} & \textbf{Object} & \textbf{Role} \\
\hline
$\barB_X$ & bar coalgebra & classifies twisting morphisms \\[3pt]
$\Omega_X$ & cobar algebra & recovers $\cA$ (Theorem B) \\[3pt]
$\Theta_\cA$ & universal MC class & organises nonlinear tower \\[3pt]
$\Zder$ & chiral derived center & universal bulk \\[3pt]
$\mathrm{Tr}_\cA$ & annulus trace
  & first modular shadow \\[3pt]
$\HH_*$ & Hochschild homology & factorisation on $S^1$ \\[3pt]
$\mathfrak{R}^{\mathrm{oc}}_\bullet$ & nonlinear resonance tower
  & quartic and higher obstructions \\[3pt]
$\nabla^{\mathrm{hol}}$ & holographic connection
  & genus-$g$ transport
\end{tabular}
\end{center}

\medskip

The first six components are proved.  The nonlinear resonance
tower is proved at finite arity (Volume~I, shadow obstruction tower).
The holographic connection is defined but its full modular
extension is the frontier.

\sep

The completed modular Koszul datum is the constructive endpoint of the
theory.  Given any chiral algebra, compute $\Pi^{\mathrm{oc}}_X$.
Theorems A--H of Volume~I are structural properties of
$\Pi^{\mathrm{oc}}_X$.  The five worked examples of
\S\ref{sec:five-examples} are instances of this computation.

The relationship between the six-fold modular Koszul datum $\Pi_X(L)$
of Volume~I (built from a Lie conformal algebra $L$) and the
eight-fold $\Pi^{\mathrm{oc}}_X(\cA)$ (built from a chiral
algebra $\cA$) is the factorisation envelope: the functor
$L \mapsto U_X(L)$ that produces the enveloping vertex algebra
from $L$.  The six-fold package lives at the input; the
eight-fold package lives at the output.  The two additional
components (derived center and annulus trace) are intrinsically
$3$-dimensional: they require the second direction.


% =================================================================
\section{From one quartic pole}
\label{sec:one-quartic-pole}
% =================================================================

The Virasoro $\lambda$-bracket
\[
\{T{}_\lambda T\}
\;=\;
\partial T + 2T\lambda + \frac{c}{12}\,\lambda^3
\]
has three terms: the simple pole ($\partial T$, the derivative),
the double pole ($2T\lambda$, the stress tensor), and the
quartic pole ($\frac{c}{12}\lambda^3$, the central charge).

The quartic pole is the engine.  It forces non-associativity of
$m_2$.  Non-associativity of $m_2$ forces $m_3$ by the Stasheff
identity $\partial m_3 + m_2 \circ m_2 = 0$.  Each $m_k$ forces
$m_{k+1}$.  The $\Ainf$ structure is genuinely infinite (class~M,
shadow depth $\infty$).

The collision residue of the $\lambda$-bracket is the gravitational
$r$-matrix:
\[
r(z) \;=\; (c/2)/z^3 + 2T/z.
\]
Two poles: a triple pole from the central charge, a simple pole from
the stress tensor.  The pole orders are one less than the OPE (the
$d\log$ kernel absorbs one power).  The $r$-matrix has no
even-order poles for a bosonic algebra.

\sep

From this single input:
\begin{enumerate}[label=\textup{(\Roman*)},nosep]
\item The $\Ainf$ tower: $m_3, m_4, \ldots$, all nonzero for
  generic $c$.  The quartic contact invariant
  $Q^{\mathrm{ct}} = 10/[c(5c+22)]$.
\item The Koszul dual: $\mathrm{Vir}_{26-c}$.  Self-duality at
  $c = 13$.  The depth-zero resonance shadow.
\item The algebraic genus tower:
  $\sum_{g \geq 1} F_g\,\hbar^{2g}
  = \kappa(\hat{A}(i\hbar) - 1)$
  with $\kappa = \kappa(\mathrm{Vir}_c) = c/2$.
  The physical genus tower of the matter-ghost system uses the
  effective curvature
  $\kappa_{\mathrm{eff}} = \kappa(\mathrm{Vir}_c)
  + \kappa_{\mathrm{ghost}} = (c-26)/2$, which vanishes at $c = 26$.
\item The Koszul triangle: boundary $\mathrm{Vir}_c$, bulk
  $\C[\![c]\!]$, lines modelled by $\mathrm{Vir}_{26-c}$-modules.
\item The gravitational coproduct: primitive (ghost-number
  obstruction from DS reduction of $\widehat\fg_k$).
\end{enumerate}

\medskip

The transgression algebra of Volume~I's working notes completes
the picture.  The Koszul dual $B = \mathrm{Vir}_{26-c}$ has
curvature $\lambda = \kappa(B) \cdot \omega_1
= \frac{26-c}{2}\,\omega_1$.  The transgression algebra
\[
B_\lambda \;=\; B\langle\eta\rangle\,\big/\,(\eta^2 - \lambda),
\qquad |\eta| = 1,
\]
adjoins a Clifford generator.  The secondary anomaly
$u = \eta^2 = \lambda$ is linear in $\kappa$ (a class, not a
scalar).  The gravity dichotomy:

\begin{center}
\small
\renewcommand{\arraystretch}{1.15}
\begin{tabular}{@{}lcc@{}}
& $c \neq 26$ & $c = 26$ \\
\hline
$u$ & $\neq 0$ & $= 0$ \\[3pt]
Clifford algebra & matrix ($2^g$-dim) & exterior ($2^{2g}$-dim)\\[3pt]
Regime & perturbative & topological \\[3pt]
Genus tower & resummable ($\hat{A}$-genus) & collapsed \\[3pt]
Surface topology & Morita-trivial & genuine
\end{tabular}
\end{center}

\medskip

The MC recursion determines the genus-$g$ partition function from
the genus-$0$ $\lambda$-bracket.  No sum over topologies.  No free
parameters.  The Maloney--Witten problem of enumeration (which
topologies to include) is replaced by a problem of algebra (which
Maurer--Cartan equation to solve).

\begin{question}
Does the graded character
$\hat{Z}_1^{\mathrm{grav}}(c;\,\tau)
:= \chi_{\mathrm{gr}}(\mathfrak{G}_1^{\mathrm{grav}}(c);\,q)$
have non-negative Fourier coefficients for all $c > 1$?
\end{question}


% =================================================================
\section{The Steinberg principle}
\label{sec:steinberg}
% =================================================================

The parallel announced in \S\ref{sec:koszul-triangle} is not
analogy; it is a theorem.

The Steinberg variety $\mathfrak{S} = \widetilde\cN
\times_\cN \widetilde\cN$ is the self-intersection of the
Springer resolution.  Its Borel--Moore homology carries the Hecke
algebra by convolution.  The Kazhdan--Lusztig basis arises from
intersection cohomology of orbit closures.  The representation
theory of $\mathfrak{S}$ is the representation theory of the Hecke
algebra: simple modules correspond to irreducible representations,
standard modules to Verma modules, the duality functor to the
Kazhdan--Lusztig involution.

\medskip

The bar complex of a chiral algebra $\cA$ lives on the product
$\FM_k(\C) \times \Conf_k(\R)$.  The differential extracts
collision residues on $\FM_k(\C)$; the coproduct deconcatenates
on $\Conf_k(\R)$.  The Koszul dual $\cA^!$ is the linear dual
of $\barB(\cA)$ (Verdier duality).  Line operators are
$\cA^!$-modules.  The representation theory of $\barB(\cA)$ is
the representation theory of the boundary: MC couplings are
line defects.

\medskip

\begin{center}
\small
\renewcommand{\arraystretch}{1.15}
\begin{tabular}{@{}lll@{}}
\textbf{Steinberg} & \textbf{Swiss-cheese} \\
\hline
$\widetilde\cN \to \cN$ (resolution)
  & $\FM_k(X) \to \Conf_k(X)$ (compactification) \\[3pt]
$\mathfrak{S} = \widetilde\cN
  \times_\cN \widetilde\cN$ (self-intersection)
  & $\barB(\cA) \otimes_{\barB(\cA\text{-bimod})} \barB(\cA)$
  (Hochschild) \\[3pt]
$H_*^{\mathrm{BM}}(\mathfrak{S})$ (Hecke algebra)
  & $\Zder(\cA)$ (derived center) \\[3pt]
KL basis (from IC)
  & PBW basis (from spectral sequence) \\[3pt]
Hecke modules
  & $\cA^!\text{-mod}$ (line operators) \\[3pt]
KL involution
  & Verdier/Koszul duality \\[3pt]
Cell structure
  & Shadow depth classification (G/L/C/M)
\end{tabular}
\end{center}

\medskip

The analogy becomes an identification at the level of
factorisation homology.  The factorisation homology of the bar
complex on a closed surface is the convolution product on the
self-intersection, which is the Hecke algebra.  The two
constructions produce the same algebra because both are computing
the endomorphism algebra of a compact generator in a factorisation
category on a configuration space.  The Fulton--MacPherson
compactification is the chiral analogue of the Springer resolution.

The cell structure of the Hecke algebra (Kazhdan--Lusztig cells,
Lusztig's $a$-function, the asymptotic Hecke algebra) should have
a shadow-tower analogue: the G/L/C/M classification of shadow
depth is the analogue of the cell decomposition, with shadow
depth playing the role of Lusztig's $a$-function.  Whether this
analogy extends to a theorem is open.


% =================================================================
\section{What remains}
\label{sec:what-remains}
% =================================================================

Everything in \S\ref{sec:two-directions}--\S\ref{sec:steinberg}
can be said in two sentences.

The bar complex of a chiral algebra on a curve carries a
differential (holomorphic direction) and a coproduct (topological
direction); together these form a Swiss-cheese algebra on
$\FM_k(\C) \times \Conf_k(\R)$, presenting the boundary theory
of a $3$d holomorphic-topological QFT whose bulk is the chiral
derived center, whose lines are modules for the Koszul dual, and
whose modularity arises from trace and clutching on the open
sector.  For gravity, the coproduct is primitive (ghost-number
obstruction from DS reduction), so all nonlinearity resides in the
$\Ainf$ products and the collision residue $r(z)$; for gauge
theory, the coproduct is nontrivial, and the full two-channel
structure is active.

That is what the $1$-form told us when we gave it a second
direction.  What follows is what it does when you push it further:
into the bordered moduli space, along the modular cooperad, past
the death of the derived category at genus $1$, through the
spectral Drinfeld strictification, into the gravitational
dichotomy at $c = 26$.  Each direction is a computation that
is either finished or honestly open.

\sep

The following structures are proved:
\begin{enumerate}[nosep]
\item The Swiss-cheese operad $\SCchtop$ and its
  homotopy-Koszulity.
\item The recognition theorem (Weiss cosheaf descent).
\item PVA descent D1--D6 (exchange cylinder, Stokes).
\item The four-stage open/closed architecture (Stages 1--3
  proved; Stage 4 programme).
\item The bulk-boundary-line triangle (boundary-linear sector
  proved; global conjectural).
\item Gravitational primitivity (principal DS; non-principal
  conjectural).
\item The benchmark: $V_k(\fsl_2)$ on $(\bP^1, \{0,\infty\})$,
  every component computed.
\item The annulus trace $\simeq \HH_*(\cC_{\mathrm{op}})$.
\item The completed modular Koszul datum $\Pi^{\mathrm{oc}}_X(\cA)$.
\end{enumerate}

The following structures are open:
\begin{enumerate}[nosep]
\item The full modular cooperad on the trace.
\item The global bulk-boundary-line triangle outside the
  boundary-linear sector.
\item BV/BRST $=$ bar at genus $g \geq 1$.
\item Non-principal gravitational primitivity beyond hook-type.
\item The spectral Drinfeld strictification for Kac--Moody
  algebras with root multiplicities $> 1$.
\end{enumerate}

\sep

\begin{quote}\itshape
A $1$-form with a simple pole.

A relation among three.

A second direction: the coproduct.

The rest follows.
\end{quote}


% =================================================================
\section{The composite graviton}
\label{sec:composite-graviton}
% =================================================================

The Virasoro stress tensor
\[
T \;=\; \frac{1}{2(k+h^\vee)}\sum_a \normord{J^a J_a}
\]
is built from affine currents.  It is not fundamental.  The
graviton is composite.

The ghost sector of the BRST complex remembers this: the
coproduct is strictly primitive
(Theorem~\ref{thm:gravitational-primitivity-wn}; see the
gauge/gravity comparison in \S\ref{sec:coproduct}).

In gauge theory, the cross-term survives because $J^a$ is
fundamental (ghost number~$0$, no passage through~$h$):
\[
\Delta_z(T_{\mathrm{aff}}) \;=\;
\tau_z(T_{\mathrm{aff}}) \otimes 1
\;+\; 1 \otimes T_{\mathrm{aff}}
\;+\; \frac{1}{k+h^\vee}\,\Omega(w{+}z,\, w).
\]
The Casimir $\Omega = \sum_a J^a \otimes J_a$ is a production
vertex.  A gluon bound state dissociates into its constituents.

The $r$-matrix $r^{\mathrm{Vir}}(z) = (c/2)/z^3 + 2T/z$ carries
all gravitational scattering.  The cubic pole arises because
the bar kernel $d\log(z-w)$ absorbs one power from the quartic
OPE (see \S\ref{sec:graviton-composite} for explicit pole orders).


% =================================================================
\section{The holographic triangle}
\label{sec:holographic-triangle}
% =================================================================

Three vertices.

\medskip

\noindent
\textbf{Boundary}: $\cA$ (the chiral algebra).

\noindent
\textbf{Bulk}: $\Zder(\cA) = H^*(C^\bullet_{\mathrm{ch}}(\cA,\cA),\,\delta)$
(the chiral derived center).

\noindent
\textbf{Lines}: $\cA^!$-modules (on the chirally Koszul locus).

\sep

The bar complex $\barB(\cA)$ classifies couplings between boundary
and lines.  It is the universal twisting coalgebra: every map
$\alpha\colon \barB(\cA) \to \cA^!$ satisfying the twisting
morphism equation is a line operator.  On the Koszul locus:
\[
\MC\bigl(\Hom(\barB(\cA),\, \cA^!)\bigr)
\;\simeq\;
\Perf(\cA^!).
\]
This is the bar complex's role: a representability theorem.  It
classifies morphisms, not operators.

The derived center $\Zder(\cA)$ classifies bulk observables.
It acts on the boundary via the brace dg-algebra structure
on $C^\bullet_{\mathrm{ch}}(\cA,\cA)$, encoding how bulk fields
modify boundary correlators.  The
restriction $\Zder(\cA) \to \cA$ is well-defined.  The extension
$\cA \to \Zder(\cA)$ is not: bulk restricts to boundary, boundary
does not determine bulk.

\begin{warning}
$\barB(\cA) \neq \Zder(\cA)$.  The bar complex is a Fourier
kernel (Volume~I).  The derived center is a space of
endomorphisms.  One classifies morphisms; the other classifies
operators.  Their conflation is the most frequent source of false
statements in three-dimensional holomorphic-topological field
theory.
\end{warning}


% =================================================================
\section{Where modularity lives}
\label{sec:where-modularity-lives}
% =================================================================

As in \S\ref{sec:annulus}, modularity arises from trace and
clutching on the open sector, not as an axiom on the closed algebra.

\sep

Cut $S^1$ into two arcs.  Each arc is an interval $I \cong \R$.
Excision on the factorisation homology of $\cC_{\mathrm{op}}$:
\[
\int_{S^1}\!\cC_{\mathrm{op}}
\;\simeq\;
\cC_{\mathrm{op}}(I_1)
\;\otimes_{\cC_{\mathrm{op}}(I_1 \cap I_2)}^{\mathbb{L}}\;
\cC_{\mathrm{op}}(I_2).
\]
The two-arc cover produces the two-sided bar complex, which is the
cyclic bar complex.  Hochschild homology
$\HH_*(\cC_{\mathrm{op}})$ IS the annulus partition function.

The first modular shadow: $\kappa(\cA)$ appears as the leading
coefficient of the annulus trace.  The genus-$1$ curvature
$\dfib^{\,2} = \kappa \cdot \omega_1$ is the obstruction to
extending the genus-$0$ Swiss-cheese structure to the torus.  The
obstruction lives in the Hochschild homology of the open sector.
It does not live in any axiom on the closed algebra.

Clutching: the genus-$g$ surface $\Sigma_g$ is built from $3g-3$
pairs of pants, sewn along $3g-3$ circles.  Each sewing circle
contributes an $\HH_*$ factor.  The modular partition function is
the trace over all sewings:
\[
Z_g
\;=\;
\Tr_{\HH_*^{\otimes(3g-3)}}
\bigl(\text{clutching maps}\bigr).
\]
This is modularity from geometry, not from axioms.


% =================================================================
\section{The Koszul dual lives on the boundary}
\label{sec:koszul-dual-boundary}
% =================================================================

$D_{\Ran}(\barB(\cA)) \simeq \barB(\cA^!)$: the Verdier dual of the bar
coalgebra is the bar of the Koszul dual algebra~$\cA^!$
(Theorem~A intertwining).  The Koszul dual~$\cA^!$ lives on
the $\R$-direction (the boundary of the half-space
$\C_{\geq 0} \times \R$), not in the $\C \times \R$ bulk.

\sep

For Chern--Simons with gauge group $G$ at level~$k$:
$\cA = V_k(\fg)$, $\cA^! = V_{-k-2h^\vee}(\fg)$.  The
Feigin--Frenkel involution $k \mapsto -k - 2h^\vee$ is Verdier
duality on $\Ran(X)$.  At the critical level $k = -h^\vee$:
$V_{-h^\vee}(\fg)$ is commutative, bar cohomology is functions on
opers, and $\cA^!$ at the dual critical level is the
Feigin--Frenkel center.

For gravity: $\cA = \mathrm{Vir}_c$,
$\cA^! = \mathrm{Vir}_{26-c}$.

At $c = 26$: $\cA^! = \mathrm{Vir}_0$, uncurved
($\kappa = 0$).

At $c = 0$: $\cA = \mathrm{Vir}_0$, uncurved.

At $c = 13$: $\cA = \cA^!$.  Self-duality.

\medskip

The complementarity sum:
\[
\kappa(\mathrm{Vir}_c) + \kappa(\mathrm{Vir}_{26-c})
\;=\; \frac{c}{2} + \frac{26-c}{2} \;=\; 13 \;\neq\; 0.
\]
For Kac--Moody algebras: $\kappa + \kappa^! = 0$ by
Feigin--Frenkel (the involution $k \mapsto -k-2h^\vee$ is
anti-symmetric around $0$).  For Virasoro: the involution
$c \mapsto 26-c$ is anti-symmetric around $13$, not around $0$.
The midpoint matters.  It is 13, not 0.


% =================================================================
\section{The BTZ black hole from the $r$-matrix}
\label{sec:btz}
% =================================================================

The effective curvature of the matter-ghost system:
\[
\kappa_{\mathrm{eff}}
\;=\;
\kappa(\mathrm{Vir}_c) + \kappa_{\mathrm{ghost}}
\;=\;
\frac{c}{2} - 13
\;=\;
\frac{c-26}{2}.
\]
At $c = 26$: $\kappa_{\mathrm{eff}} = 0$.  No curvature.
The theory is topological.

At $c \neq 26$: the genus tower
$F_g = \kappa_{\mathrm{eff}} \cdot \lambda_g^{\mathrm{FP}}$.
The $\hat{A}$-genus generates all $F_g$ from $\kappa_{\mathrm{eff}}$:
\[
\sum_{g \geq 1} F_g\,\hbar^{2g}
\;=\;
\kappa_{\mathrm{eff}}
\bigl(\hat{A}(i\hbar) - 1\bigr).
\]
The Bernoulli numbers encode the higher-genus corrections.  Each
$F_g$ is positive (because $\hat{A}(ix) = (x/2)/\sin(x/2)$ has
all positive Taylor coefficients).

\sep

The collision residue
$r^{\mathrm{Vir}}(z) = (c/2)/z^3 + 2T/z$ controls BTZ
thermodynamics.  The leading coefficient $c/2 = \kappa$ is the
scalar that generates the Bekenstein--Hawking entropy at leading
order in $1/G$.  The subleading $2T/z$ is state-dependent braiding:
it distinguishes different BTZ states within the same theory.

The primitive coproduct is the algebraic statement that black holes
do not split.  The only channel is braiding.  The production
vertex $\Omega(w{+}z,w)$ of gauge theory (the Casimir cross-term)
is absent.  There is no splitting channel.

\begin{center}
\small
\renewcommand{\arraystretch}{1.15}
\begin{tabular}{@{}lcc@{}}
& $c \neq 26$ & $c = 26$ \\
\hline
$\kappa_{\mathrm{eff}}$ & $(c-26)/2 \neq 0$ & $0$ \\[3pt]
$F_1$ & $(c-26)/48$ & $0$ \\[3pt]
Genus tower & $\kappa_{\mathrm{eff}}\cdot(\hat{A}-1)$ & collapsed \\[3pt]
Regime & perturbative & topological
\end{tabular}
\end{center}


% =================================================================
\section{Five theories, one machine}
\label{sec:five-theories}
% =================================================================

\begin{center}
\small
\renewcommand{\arraystretch}{1.15}
\begin{tabular}{@{}lccccc@{}}
& CS & Gravity & Twisted hol.\ & M2 & M5 \\
\hline
$\cA$ & $V_k(\fg)$ & $\mathrm{Vir}_c$
  & $V_k(\fg)$ & DDCA$_K$ & $\cW_\infty(K)$ \\[3pt]
$\cA^!$ & $V_{-k-2h^\vee}(\fg)$ & $\mathrm{Vir}_{26-c}$
  & $V_{-k-2h^\vee}(\fg)$ & CE & ? \\[3pt]
$\kappa$ & $\dim(\fg)\frac{k+h^\vee}{2h^\vee}$
  & $c/2$ & $\dim(\fg)\frac{k+h^\vee}{2h^\vee}$
  & $K(K{+}1)/2$ & $\varrho_\infty c$ \\[3pt]
Class & L & M & L & L & M \\[3pt]
$\Delta_z$ & nontrivial & primitive & nontrivial
  & nontrivial & primitive \\[3pt]
$\kappa + \kappa^!$ & $0$ & $13$ & $0$ & $0$
  & $\varrho K_N$
\end{tabular}
\end{center}

\medskip

The machine does not know which theory it is computing.  It
receives: a chiral algebra, its OPE, and the bar construction.
The five theories are five inputs to the same functor.

Chern--Simons and twisted holography are class~L: the $r$-matrix
carries all structure, the coproduct is nontrivial, the shadow
tower terminates at arity~$3$.  Gravity and M5 are class~M:
infinite shadow depth, primitive coproduct,
$Q^{\mathrm{ct}} \neq 0$.

\sep

The M2--M5 duality: the M2~algebra (DDCA, class~L) and the
M5~algebra ($\cW_\infty$, class~M) are related by the Miura
transformation.  DS~reduction converts class~L to class~M.
The reduction IS the depth increase.  The ghost-number filtration
that kills the coproduct of M5 is the same filtration that was
absent for M2.


% =================================================================
\section{The four-stage architecture, revisited}
\label{sec:four-stages-revisited}
% =================================================================

Stage~1 (local one-colour): define the Swiss-cheese algebra.
Prove the universal acting-bulk theorem
$U(\cA) = (C^*_{\mathrm{ch}}(\cA,\cA),\, \cA)$.
Everything succeeds or fails here.

\medskip

Stage~2 (open primitive): the primitive object is
$\cC_{\mathrm{op}}$, not $A_b$.  The boundary algebra
$A_b = \End(b)$ is a chart of a compact generator $b$ in
$\cC_{\mathrm{op}}$.  Different generators produce
Morita-equivalent charts.  The category is the invariant.

\medskip

Stage~3 (globalise): move from $\C \times \R$ to tangential log
curves $(X, D, \tau)$.  The real oriented blowup
$\pi\colon \widetilde{X}_D \to X$ replaces each puncture $p_i$ by
a circle $S^1_{p_i}$; the tangential basepoint $\tau_{p_i}$ cuts
this circle into an interval $I_{p_i} \cong \R$.  Open insertions
live on the interval.  Bulk insertions live in the interior.  The
global theory is a factorisation algebra on a bordered Riemann
surface.

\medskip

Stage~4 (modularise): introduce $\Theta^{\mathrm{oc}}$, traces,
annulus/Hochschild chains, clutching, quartic resonance classes.
Only after stages~1--3 are solid.

\sep

Each stage depends on its predecessor.  Discussion of stage~$k$
before establishment of stage~$k-1$ is not premature; it is wrong.


% =================================================================
\section{What the bar complex knows about quantum gravity}
\label{sec:bar-knows-gravity}
% =================================================================

The bar complex $\barB(\mathrm{Vir}_c)$ is a factorisation
coalgebra on $\Ran(X)$.  Its differential $d_{\barB}$ extracts
collision residues of the Virasoro OPE.  Its coproduct~$\Delta$
deconcatenates ordered tuples along $\R$.

At genus~$0$: $d_{\barB}^2 = 0$ by Arnold.  The bar complex is
an honest chain complex.  Its cohomology has dimensions
\[
1,\; 2,\; 5,\; 12,\; 30,\; 76,\; 196,\; 512,\; 1353,\; \ldots
\]
The generating function satisfies
$x^3 P^2 - (1-2x)(1-x^2)P + (1-x)(1-x^2) = 0$.
Discriminant: $(1-x)^2(1-3x)(1+x)$.  Growth rate~$3$.

\sep

At genus~$1$:
$\dfib^{\,2} = (c/2) \cdot \omega_1 \neq 0$.
The bar complex is curved.  The curvature $c/2$ is the modular
characteristic.

At genus~$g$:
\[
F_g
\;=\;
\frac{c}{2} \cdot \lambda_g^{\mathrm{FP}}
\;=\;
\frac{c}{2} \cdot
\frac{(2^{2g-1} - 1)\,|B_{2g}|}{(2g)!\cdot 2^{2g-1}}.
\]
The Bernoulli numbers.  The $\hat{A}$-genus.  From index theory.
Nothing in the Virasoro OPE predicts this.  It lives in the
geometry of $\overline\cM_g$.

\sep

The shadow obstruction tower
\begin{align*}
S_2 &= c/2, \\
S_3 &= 2, \\
S_4 &= 10\big/\bigl[c(5c+22)\bigr], \\
S_5 &= -48\big/\bigl[c^2(5c+22)\bigr],
\end{align*}
and all higher $S_r$, are the Taylor coefficients of $\sqrt{Q}$
where
\[
Q_{\mathrm{Vir}}(t)
\;=\;
c^2 + 12ct
+ \frac{180c + 872}{5c + 22}\,t^2.
\]
One quadratic.  One square root.  An algebraic function of
degree~$2$.

\medskip

Three numbers determine the entire infinite tower of higher-genus
corrections to quantum gravity: $\kappa = c/2$, $\alpha = 2$ (the
cubic shadow), and $S_4 = 10/[c(5c+22)]$ (the quartic contact
invariant).  The shadow metric $Q$ is constructed from these three.
Its discriminant $\Delta = 8\kappa S_4$ controls the convergence
radius of the shadow obstruction tower.  When $\Delta \neq 0$ (all
$c \neq 0$): the tower is infinite, class~M.  The bar complex
knows the exact answer at every genus, as the coefficients of an
algebraic function.


% =================================================================
\section{The arithmetic structure of quantum gravity}
\label{sec:arithmetic-gravity}
% =================================================================

The gravitational shadow obstruction tower is the arithmetic of the quadratic
extension $K = \Q(c)(t)(\sqrt{Q_{\mathrm{Vir}}})$ over the
function field $F = \Q(c)(t)$.  The shadow metric
\[
  Q_{\mathrm{Vir}}(t)
  \;=\;
  c^2 + 12ct + \frac{180c + 872}{5c + 22}\,t^2
\]
is a quadratic polynomial whose square root generates the
infinite tower of higher-genus corrections to three-dimensional
quantum gravity.

\sep

\noindent\textbf{The discriminant.}\enspace
$\mathrm{disc}(Q_{\mathrm{Vir}}) = -320\,c^2/(5c{+}22)$.
Negative for all $c > 0$: the roots of $Q_{\mathrm{Vir}}$
are complex conjugate, the extension $K/F$ is genuine, and
the shadow obstruction tower is infinite (class~M).

The discriminant controls three quantities at once:
\begin{enumerate}[label=(\roman*),nosep]
\item The depth of the gravitational perturbation series
  (infinite for all $c \neq 0$, because $\Delta = 40/(5c{+}22)
  \neq 0$).
\item The convergence radius of the shadow obstruction tower: the growth
  rate $\rho = \sqrt{36 + 80/(5c{+}22)}\,/\,c$ is the reciprocal
  of the distance from the origin to the nearest branch point
  of $\sqrt{Q_{\mathrm{Vir}}}$.
\item The complementarity with the dual theory at $c' = 26 - c$:
  $\Delta(c) + \Delta(26{-}c) = 6960/[(5c{+}22)(152{-}5c)]$,
  with universal numerator.
\end{enumerate}

\sep

\noindent\textbf{The gravitational number field.}\enspace
The Bekenstein--Hawking entropy $S_{\mathrm{BH}} = 2\pi\sqrt{ch/6}$
derives from $\kappa = c/2$, the leading coefficient of
$Q_{\mathrm{Vir}}$.  This coefficient is the constant term
(up to a square): $Q_{\mathrm{Vir}}(0) = c^2 = 4\kappa^2$.
The leading OPE datum seeds the quadratic extension.

The growth rate $\rho$ measures how fast the higher-genus
corrections to the gravitational partition function grow.
When $\rho < 1$, the shadow obstruction tower converges: the shadow
partition function $Z^{\mathrm{sh}} = \sum Z_g^{\mathrm{sh}}
\hbar^{2g}$ exists as a convergent series.  This convergence
is the ``splitting'' of the gravitational prime at central
charge~$c$: the quadratic extension has a rational point in
a neighbourhood of the origin.

Setting $\rho = 1$ yields the critical cubic
$5c^3 + 22c^2 - 180c - 872 = 0$ with positive root
$c^* \approx 6.125$.  For $c > c^*$: the shadow obstruction tower
converges ($\rho < 1$, the gravitational perturbation
series sums).  For $0 < c < c^*$: the shadow obstruction tower diverges
($\rho > 1$, the perturbation series requires resummation).
The critical cubic is the Minkowski bound of the gravitational
number field: it separates the convergent regime from the
divergent regime.

\sep

\noindent\textbf{The three critical values.}\enspace
\begin{center}
\small
\renewcommand{\arraystretch}{1.15}
\begin{tabular}{@{}ccl@{}}
$c$ & \textbf{What happens} & \textbf{Arithmetic meaning} \\
\hline
$0$ & $\kappa = 0$, $Q$ degenerates & ramified prime
  (shadow obstruction tower is undefined) \\[3pt]
$13$ & $\cA = \cA^!$, self-dual
  & Galois involution is trivial on coefficients;
  $K = K'$ \\[3pt]
$26$ & $\kappa_{\mathrm{eff}} = 0$, topological & dual
  extension $K'$ degenerates ($\kappa' = 0$)
\end{tabular}
\end{center}

\medskip

At $c = 0$: the constant term of $Q_{\mathrm{Vir}}$ vanishes.
The quadratic polynomial has a double root at $t = 0$.  The
shadow connection $\nabla^{\mathrm{sh}}$ has an irregular
singularity.  The extension ramifies.

At $c = 13$: the theory is self-dual ($\mathrm{Vir}_{13}^!
= \mathrm{Vir}_{13}$).  The Koszul involution $c \mapsto 26-c$
fixes $c = 13$.  The two quadratic extensions $K$ and $K'$
coincide.  The complementarity sum $\kappa + \kappa' = 13$
does not vanish: self-duality is a symmetry enhancement, not
anomaly cancellation.  The growth rate $\rho(13) \approx 0.467$:
the shadow obstruction tower converges at the self-dual point.

At $c = 26$: the effective curvature $\kappa_{\mathrm{eff}}
= (c{-}26)/2 = 0$ vanishes.  The dual algebra
$\mathrm{Vir}_0$ has $\kappa' = 0$: its shadow extension
ramifies.  The gravitational theory becomes topological.  The
Clifford algebra degenerates from a matrix algebra
($2^g$-dimensional) to an exterior algebra
($2^{2g}$-dimensional).  The BTZ partition function collapses:
$Z_{\mathrm{grav}} = 1$.

These three values are not phases of one theory.  They are three
qualitatively distinct points in the arithmetic landscape of the
quadratic extension: a ramified prime ($c = 0$), a self-dual
fixed point ($c = 13$), and a degeneration of the dual
($c = 26$).  They are separated by $13$ units of central charge.

\sep

\noindent\textbf{The ramified primes.}\enspace
The denominator theorem (Volume~I) gives
$\mathrm{denom}(S_r)
= (\mathrm{const})\cdot c^{r-3}(5c{+}22)^{\lfloor(r-2)/2\rfloor}$.
Two primes ramify:
\begin{enumerate}[label=(\roman*),nosep]
\item $c = 0$: the leading OPE coefficient vanishes.  No
  Bekenstein--Hawking entropy.  No bar curvature.  The shadow
  metric $Q(0,t) = 872\,t^2/(5 \cdot 22) = 872\,t^2/110$
  degenerates.
\item $c = -22/5$: the weight-$4$ Gram determinant
  $\det G_4 = c^2(5c{+}22)/2$ vanishes.  A null vector at
  weight~$4$ collapses the quasi-primary $\Lambda$ onto
  $\partial^2 T$.  The quartic contact invariant
  $Q^{\mathrm{ct}} = 10/[c(5c{+}22)]$ has a pole.
\end{enumerate}
Every other value of~$c$ is unramified: the shadow
coefficients are finite, the quadratic extension is
well-defined, and the tower runs.

No higher Kac determinant appears.  The gravitational number
field has exactly two ramified primes, both visible from the
weight-$4$ representation theory.  The infinite tower of
higher-genus gravitational corrections is controlled by a
single $2 \times 2$ Gram matrix.

\sep

\noindent\textbf{The unit group and the BTZ partition function.}\enspace
The growth rate
$\rho = \sqrt{36 + 80/(5c{+}22)}\,/\,c$ is the fundamental
unit of the gravitational number field.  It controls the BTZ
partition function through the radius of convergence of the
shadow obstruction tower:
\begin{itemize}[nosep]
\item $\rho < 1$ ($c > c^*$): the shadow obstruction tower converges.  The
  gravitational path integral exists non-perturbatively as a
  convergent shadow series.  The shadow partition function
  $Z^{\mathrm{sh}}$ is the value of a convergent double series
  (genus and arity).
\item $\rho > 1$ ($0 < c < c^*$): the shadow obstruction tower diverges.
  The shadow coefficients grow exponentially.  Resummation
  (Borel transform of $\sqrt{Q_{\mathrm{Vir}}}$) is required.
  The Stokes constants of the shadow obstruction tower are the
  non-perturbative content.
\item $\rho = 1$ ($c = c^*$): the critical cubic.  The branch
  point of $\sqrt{Q_{\mathrm{Vir}}}$ lies exactly on the unit
  circle.  The shadow obstruction tower is marginally divergent.
\end{itemize}

\begin{center}
\small
\renewcommand{\arraystretch}{1.15}
\begin{tabular}{@{}rrl@{}}
$c$ & $\rho$ & \textbf{Regime} \\
\hline
$1$ & $2.99$ & divergent (shadow) \\
$c^* \approx 6.125$ & $1.00$ & critical \\
$13$ & $0.467$ & convergent \\
$26$ & $0.234$ & convergent (dual topological) \\
$100$ & $0.061$ & strongly convergent
\end{tabular}
\end{center}

\medskip

The string partition function $Z_{\mathrm{string}}$ (genus
factorial growth $(2g)!$) always diverges.  The shadow partition
function $Z^{\mathrm{sh}}$ (Bernoulli decay $1/(2\pi)^{2g}$
times shadow growth $\rho^r r^{-5/2}$) converges for $c > c^*$.
The shadow obstruction tower is better-behaved than the string because it
sums over arities within a fixed genus, while the string sums
over genera with unrestricted growth.

\sep

\noindent\textbf{The Galois involution.}\enspace
The monodromy $-1$ of the shadow connection
$\nabla^{\mathrm{sh}} = d - Q'/2Q\,dt$ around each branch
point of $\sqrt{Q_{\mathrm{Vir}}}$ is the nontrivial element
of $\mathrm{Gal}(K/F)$.  In the gravitational context: the
Koszul sign is the Galois involution of the number field.  The
sign alternation in the shadow obstruction tower $S_r$ (positive for even
$r$, approximately negative for odd $r$, with beat period
$\pi/(\pi - \theta)$ where $\theta = \arg(t_+)$) is the
manifestation of the Galois action on the Taylor coefficients
of $\sqrt{Q_{\mathrm{Vir}}}$.

This is the deepest structural content of the bar
construction's desuspension: the sign $(-1)^n$ at bar degree
$n$ is the Galois automorphism of a quadratic extension,
evaluated at the $n$-th coefficient.

\sep

The gravitational shadow obstruction tower is not like a quadratic
extension.  It is one, over the function field $\Q(c)(t)$
instead of $\Q$.  The Bekenstein--Hawking entropy is the leading
coefficient.  The convergence of the perturbation series is the
splitting behaviour.  The Koszul sign is the Galois involution.
The critical cubic is the Minkowski bound.  The two ramified
primes are the Kac determinant zeros.  The complementarity
identity is the duality constraint on discriminants.  The
gravitational partition function is the norm form of the
extension, evaluated at a specific point.

\begin{quote}\itshape
A $1$-form with a simple pole.

A relation among three.

A second direction: the coproduct.

A quadratic extension: the shadow.

The rest follows.
\end{quote}


% =================================================================
\section{Gravity at rational central charge}
\label{sec:gravity-rational-c}
% =================================================================

The generic gravitational number field
$K = \Q(c)(t)(\sqrt{Q_{\mathrm{Vir}}})$ of the preceding
section lives over the function field $\Q(c)$.  When the
central charge specialises to a rational number $c_0$, the
field specialises to a concrete quadratic extension of~$\Q$.

For a minimal model $c_{p,q} = 1 - 6(p-q)^2/(pq)$,
the critical discriminant
$\Delta(c_{p,q}) = 40/(5c_{p,q} + 22)$ is a specific
rational number, and
$K(c_{p,q}) := \Q(\sqrt{\Delta(c_{p,q})})$
is a specific real quadratic number field.  Gravity at
each minimal model selects a number field.

\sep

\noindent\textbf{The physical picture.}\enspace
The Bekenstein--Hawking entropy $\kappa = c/2$ at a rational
central charge $c_{p,q}$ is a rational number.  The
higher-genus corrections $F_g = \kappa \cdot \lambda_g^{\mathrm{FP}}$
are rational multiples of~$\kappa$, hence rational.  The
ENTIRE perturbative gravitational partition function
$\sum F_g \hbar^{2g}$ has rational coefficients.

But the shadow obstruction tower is different.  The Taylor coefficients
$S_r$ of $\sqrt{Q_L}$ are rational (they satisfy a rational
recursion), yet the generating function $H(t) = t^2\sqrt{Q_L(t)}$
is genuinely irrational: it involves $\sqrt{\Delta}$, which
lies outside $\Q$ for all $c > 0$.  The irrationality is the
arithmetic content of the shadow obstruction tower.  The bar complex of
the Virasoro algebra, specialised to a minimal model, produces
an algebraic number that no finite truncation of the
perturbative expansion can detect.

This is not a deficiency of perturbation theory.  The function
$H(t)$ is algebraic of degree~$2$; its Taylor coefficients are
all individually rational.  The irrationality is non-perturbative
in the arity variable $t$, not in the genus variable $\hbar$.
It is visible in the branch structure of $\sqrt{Q_L}$, not in
any single coefficient.

\sep

\noindent\textbf{Computation.}\enspace
Write $5c_0 + 22 = a/b$ in lowest terms.  Then
$\Delta(c_0) = 40b/a$, and the specialised field is
$K(c_0) = \Q(\sqrt{\mathrm{sqfree}(40ab)})$.

\medskip

\noindent\textit{$c = 1/2$ (Ising model).}\enspace
$5c + 22 = 49/2$.  $\Delta = 80/49$.
$\sqrt{80/49} = 4\sqrt{5}/7$.
The Ising gravity lives over $\Q(\sqrt{5})$.

The golden ratio $(1+\sqrt{5})/2$ is the fundamental unit
of this field, with norm $-1$.  The class number is $h = 1$:
the ring of integers $\Z[(1+\sqrt{5})/2]$ is a principal ideal
domain.  The Ising model produces the simplest possible
arithmetic.  One is tempted to say that the Ising model is
``arithmetically trivial,'' as it is physically the simplest
unitary minimal model.

The shadow growth rate is $\rho(1/2) \approx 12.53$: the
shadow obstruction tower diverges violently ($c = 1/2$ is far below
the critical cubic $c^* \approx 6.125$).  The Ising gravity
requires resummation.

\medskip

\noindent\textit{$c = 7/10$ (tricritical Ising).}\enspace
$5c + 22 = 51/2$.  $\Delta = 80/51$.
$K(7/10) = \Q(\sqrt{255})$ where $255 = 3 \cdot 5 \cdot 17$.
Discriminant $1020$.  Class number $h = 4$: a genuine failure
of unique factorisation.

\medskip

\noindent\textit{$c = 4/5$ (three-state Potts).}\enspace
$K(4/5) = \Q(\sqrt{65})$.  Discriminant~$65$.  $h = 2$.
The Potts-model gravity has a non-trivial ideal class group
$\mathrm{Cl}(K) \cong \Z/2$.

\medskip

\noindent\textit{$c = 13$ (self-dual).}\enspace
$K(13) = \Q(\sqrt{870})$ where $870 = 2 \cdot 3 \cdot 5 \cdot 29$.
Discriminant $3480$.  $h = 8$.

The self-dual gravity has the largest class number of any
integer central charge in $[0, 26]$.  It sees the prime~$29$,
inherited from the complementarity constant
$6960 = 2^4 \cdot 3 \cdot 5 \cdot 29$.  No other central
charge in the standard range detects this prime.

\medskip

\noindent\textit{$c = 26$ (critical string).}\enspace
$K(26) = \Q(\sqrt{95})$.  Discriminant $380$.  $h = 2$.
$\rho(26) \approx 0.232$: convergent.  This is the
gravitational number field of the bosonic string at the
critical dimension.

\begin{center}
\small
\renewcommand{\arraystretch}{1.15}
\begin{tabular}{@{}lccrrcr@{}}
$c$
  & $\Delta(c)$
  & $d(c)$
  & $\mathrm{disc}$
  & $h$
  & $\varepsilon$
  & $\rho(c)$ \\
\hline
$1/2$ & $80/49$ & $5$ & $5$ & $1$
  & $(1{+}\sqrt{5})/2$ & $12.53$ \\[2pt]
$7/10$ & $80/51$ & $255$ & $1020$ & $4$
  & $16{+}\sqrt{255}$ & $8.94$ \\[2pt]
$4/5$ & $20/13$ & $65$ & $65$ & $2$
  & $8{+}\sqrt{65}$ & $7.81$ \\[2pt]
$1$ & $40/27$ & $30$ & $120$ & $2$
  & $11{+}2\sqrt{30}$ & $6.24$ \\[2pt]
$13$ & $40/87$ & $870$ & $3480$ & $8$
  & $59{+}2\sqrt{870}$ & $0.47$ \\[2pt]
$25$ & $40/147$ & $30$ & $120$ & $2$
  & $11{+}2\sqrt{30}$ & $0.24$ \\[2pt]
$26$ & $5/19$ & $95$ & $380$ & $2$
  & $39{+}4\sqrt{95}$ & $0.23$
\end{tabular}
\end{center}

\sep

\noindent\textbf{Koszul duality and the field pair.}\enspace
The Koszul pair $(c, 26-c)$ produces two fields $K(c)$
and $K(26-c)$.  These coincide iff
$(5c+22)/(152-5c) \in \Q^{\times 2}$.

At $c = 1$ and $c = 25$: the ratio is $27/147 = (3/7)^2$.
Both give $\Q(\sqrt{30})$.

At $c = 13$: the ratio is $87/87 = 1$.  Self-duality forces
field-preservation.

For the Ising model: the ratio is $49/299 = 49/(13 \cdot 23)$,
which is not a perfect square.  The Ising gravity at $c = 1/2$
lives over $\Q(\sqrt{5})$; its Koszul dual at $c = 51/2$
lives over $\Q(\sqrt{1495})$ where $1495 = 5 \cdot 13 \cdot 23$.
Different fields, different class numbers, different arithmetic.
Koszul duality does not preserve the gravitational number field
at the Ising point.

This failure is generic.  Among all minimal models $M(p,q)$
with $p < 10$, none has the Koszul pair giving the same field.
The criterion $(5c+22)/(152-5c) \in \Q^{\times 2}$ requires
an arithmetic coincidence that the minimal model spectrum
does not provide.

\sep

\noindent\textbf{What the class number measures.}\enspace
The class number $h$ of $K(c)$ is a measure of the
arithmetic complexity of gravity at central charge~$c$.
In algebraic number theory, $h$ counts the failure of
unique factorisation in the ring of integers $\cO_K$.

In the gravitational context: $h = 1$ means the shadow obstruction tower's
arithmetic is ``principal,'' governed by a single generator.
$h > 1$ means the shadow arithmetic requires multiple ideal
classes.  Whether this has a direct physical interpretation
(a counting of distinct gravitational saddle classes)
is not known.

\begin{center}
\small
\renewcommand{\arraystretch}{1.15}
\begin{tabular}{@{}lcl@{}}
$c$ & $h$ & \textbf{Interpretation} \\
\hline
$1/2$ (Ising) & $1$ & principal: simplest arithmetic \\[2pt]
$4/5$ (Potts) & $2$ & one non-trivial ideal class \\[2pt]
$7/10$ (tricrit.) & $4$ & four ideal classes \\[2pt]
$13$ (self-dual) & $8$ & most complex in standard range
\end{tabular}
\end{center}

\sep

\noindent\textbf{What fails.}\enspace
The growth rate $\rho(c)$ is a smooth algebraic function
of~$c$.  The fundamental unit $\varepsilon(d(c))$ of $K(c)$
depends on $c$ through the squarefree part $d(c)$, which is a
wildly discontinuous function of the central charge.  No
algebraic relation connects $\rho$ and $\varepsilon$.

A conjecture that the fundamental unit of the gravitational
number field controls the shadow growth rate was tested by
computing both quantities at $c = 1/2, 7/10, 4/5, 1, 13, 25, 26$
and found to be false.  At $c = 1/2$: $\rho \approx 12.53$
while $\varepsilon = (1+\sqrt{5})/2 \approx 1.618$.  At
$c = 13$: $\rho \approx 0.467$ while
$\varepsilon = 59 + 2\sqrt{870} \approx 117.99$.  The two
quantities vary in opposite directions.

The failure is structural: $\rho(c)$ is a smooth function
of~$c$, while $\varepsilon(d(c))$ is a step function.
A smooth function cannot equal a step function.

\sep

\noindent\textbf{What remains.}\enspace
The specialisation $c \mapsto K(c)$ is a new functor from
(a subset of) the rational points of the Virasoro central
charge line to the set of real quadratic number fields.  It
sends the OPE data at genus~$0$ (the stress tensor and its
weight-$4$ descendant) to a number field whose arithmetic
is entirely determined by the bar complex.

The family $\{K(c_{p,q})\}_{(p,q)}$ attached to the minimal
model tower constitutes a new arithmetic invariant of 2d~CFT.
Whether this family has structure beyond what is visible in
the table above is an open question.

The first instance of the specialisation functor applied to
gravity is the Ising model.  The shadow obstruction tower of the
Ising model lives over $\Q(\sqrt{5})$, the field of the
golden ratio.  This is the number field with class number~$1$
and the smallest possible discriminant among real quadratic
fields with fundamental unit of norm $-1$.  The simplest
unitary minimal model produces the simplest real quadratic
field.

Whether this simplicity persists along the minimal model tower,
or whether more complex arithmetic emerges, is a question at
the intersection of vertex algebra theory and algebraic number
theory that the bar construction makes precise.

\begin{quote}\itshape
The bar complex is a 1-form.

Its square root is a number field.

Gravity at rational central charge is arithmetic.
\end{quote}


% =================================================================
\section{Fourier duality from operads to $L$-functions}
\label{sec:fourier-operads-L}
% =================================================================

The bar construction is a Fourier transform
(\S\ref{sec:fourier}, Volume~I).  The shadow obstruction tower is a
quadratic extension (\S\ref{sec:arithmetic-gravity}).  The
Dirichlet-sewing lift $S_\cA(u)$ is a product of zeta functions
(\S\ref{sec:prime-factorisation}, Volume~I).  These three objects
arise from the same chiral algebra.  This section asks: do they
form a ``descent chain'' in which a single Fourier duality
descends from operadic algebra to analytic number theory?

The answer is: partially.  The chain has five levels and three
distinct involutions.  The first two links are theorems.  The
third is broken.  The fourth is lattice-specific.  The claim
``they are the same involution'' is false.

\subsection{The three involutions}

The physical content of the bar construction is that it extracts
the OPE residues of a chiral algebra into a nilpotent differential,
converting the multiplicative structure (operator product) into an
additive structure (differential).  This extraction carries three
symmetries:

\begin{enumerate}[label=\textup{(\roman*)}]
\item \textbf{Verdier duality} $\bD_{\Ran}$ on $\Ran(X)$.
  Exchanges $\barB(\cA)$ and $\barB(\cA^!)$ (Theorem~A,
  Volume~I).  This is the Fourier involution: it intertwines
  convolution with pointwise operations on factorisation
  (co)algebras.  Its kernel is $\eta_{ij} = d\log(z_i - z_j)$.

  Fixed locus: self-dual factorisation coalgebras.

\item \textbf{Galois involution} $\sigma$ of the shadow
  extension $K_L = \Q(c)(t)(\sqrt{Q_L})$.  Sends
  $\sqrt{Q_L} \mapsto -\sqrt{Q_L}$.  This is the monodromy
  of the shadow connection $\nabla^{\mathrm{sh}}$: transport
  around a branch point of $\sqrt{Q_L}$ returns $-\sqrt{Q_L}$.
  It is the Koszul sign $(-1)^n$ at bar degree~$n$.

  Fixed locus: the base field $F = \Q(c)(t)$.

\item \textbf{Koszul involution} $\cA \mapsto \cA^!$ on
  families of chiral algebras.  For Virasoro: $c \mapsto 26 - c$.
  For Kac--Moody: $k \mapsto -k - 2h^\vee$.  This acts on the
  \emph{coefficients} of the shadow metric $Q_L$, not on
  $\sqrt{Q_L}$ itself.

  Fixed locus: self-dual algebras ($c = 13$ for Virasoro).
\end{enumerate}

These three involutions are related but distinct.  Verdier
duality produces the Koszul involution as its shadow on moduli:
$\bD_{\Ran}(\barB(\cA)) \simeq \barB(\cA^!)$ implies that
$\cA \mapsto \cA^!$ is the trace of Verdier duality on isomorphism
classes.  The Galois involution is intrinsic to a \emph{single}
algebra's shadow obstruction tower and has no direct relation to Verdier
duality.  It arises because the MC equation
$D\Theta + \tfrac{1}{2}[\Theta, \Theta] = 0$, projected to the
scalar arity line, reduces to a Riccati equation whose solution
$H(t) = t^2\sqrt{Q_L}$ is algebraic of degree~$2$---not higher,
because the Riccati is quadratic.

The claim ``they are the same involution (Fourier) acting at
different levels'' conflates (i)--(iii).  The correct statement:
they are three $\Z/2$-actions on three different categories, all
originating from the bar construction but not related by any
``descent.''  At the self-dual point $c = 13$, all three have
compatible fixed-point behaviour (the Koszul pair coincides, the
shadow extensions $K_L$ and $K_{L!}$ agree, and the Verdier dual
coalgebra is isomorphic to the original).  At generic $c$, they
diverge: the Galois involution of $K_L$ has nothing to do with
the Koszul involution on the $c$-line.

\subsection{Where the chain breaks}

The descent chain has five levels:

\begin{center}
\small
\renewcommand{\arraystretch}{1.15}
\begin{tabular}{@{}cll@{}}
\textbf{Level} & \textbf{Object} & \textbf{Status} \\
\hline
0 & Koszul duality $\cA \mapsto \cA^!$
  & Proved (Theorem~A) \\
1 & Bar as Fourier: $\barB_X$, inversion, Plancherel
  & Proved (Theorems A,B,C) \\
2 & Shadow obstruction tower as quadratic extension
  & Proved \\
3 & Dirichlet-sewing lift $S_\cA(u)$
  & Defined, but independent of Level 2 \\
4 & Zeros of $\varepsilon^c_s$ on critical lines
  & Lattice only; structural obstruction
\end{tabular}
\end{center}

The break occurs between Levels 2 and 3.  The shadow obstruction tower
is determined by $(\kappa, \alpha, S_4)$---the arity-$2$, $3$,
$4$ OPE data---and lives over $\Q(c)(t)$.  The Dirichlet-sewing
lift is determined by the weight multiset $W(\cA) = \{w_i\}$
and lives over the complex $u$-plane.  These are independent
invariants: two algebras with the same weight multiset but
different OPE structure constants have the same $S_\cA(u)$ but
different shadow obstruction towers.  Conversely, two algebras with the same
$(\kappa, \alpha, S_4)$ but different weight multiplicities have
the same shadow obstruction tower but different $S_\cA(u)$.

\subsection{The structural obstruction}

For lattice VOAs, the constrained Epstein zeta $\varepsilon^c_s$
factors into $L$-functions because the theta function $\Theta_\Lambda$
decomposes under the Hecke algebra into Eisenstein series and
eigenforms.  This decomposition uses three properties special to
the lattice setting: the theta function is a classical modular form,
the Hecke decomposition is a finite sum, and the Fourier coefficients
are lattice representation numbers with multiplicative structure.

None of these holds for the Virasoro algebra.  The Virasoro
character $q^{-c/24}\prod(1-q^n)^{-1}$ is not a modular form of
finite weight.  The shadow obstruction tower still runs (class~M, infinite
depth), but the arithmetic sieve does not apply.  The MC equation
constrains the spectral coefficients on the \emph{real} spectral
axis; the zeta zeros live at \emph{complex} spectral parameters.
This separation is structural
(Remark~\ref{rem:structural-obstruction}, Volume~I): algebraic
constraints on the spectral line cannot reach the scattering poles
without analytic continuation of the spectral data off the real
axis.

Concretely: the modified Li coefficients
$\tilde\lambda_n(\cH)$ for the Heisenberg sewing lift
$S_\cH(u) = \zeta(u)\zeta(u{+}1)$ are positive for
$n = 1, \ldots, 6$ and \textbf{negative at $n = 7$}
(\texttt{compute/tests/test\_euler\_koszul\_engine.py}, line~394).
Positivity of all Li coefficients would be equivalent to GRH for
$\zeta(u)\zeta(u{+}1)$.  The sign change at $n = 7$ shows that the
sewing-lift Li coefficients do not satisfy a ``shadow RH.''  (This
does not contradict RH for $\zeta$ alone; the sign change is an
artefact of the product and the regularisation.)

\subsection{What the bar construction does provide}

The descent chain, as a linear sequence, is broken.  But the
bar construction does provide something more interesting: a
\emph{fan} of three independent projections from a single
MC element.

\medskip

\begin{center}
\small
\begin{tabular}{@{}lll@{}}
\textbf{Projection} & \textbf{Output} & \textbf{Content} \\
\hline
Genus tower
  & $F_g = \kappa \cdot \lambda_g^{\mathrm{FP}}$
  & Topological (Mumford classes) \\[3pt]
Shadow obstruction tower
  & $\{S_r\}_{r \geq 2} = [t^{r-2}]\sqrt{Q_L}$
  & Arithmetic (quadratic extension, discriminant) \\[3pt]
Koszul dual
  & $\cA^! = (\barB(\cA))^\vee$
  & Algebraic (Verdier dual)
\end{tabular}
\end{center}

\medskip

These three projections are facets of $\Theta_\cA
\in \mathrm{MC}(\gAmod)$, the universal MC element in the
modular cyclic deformation complex.  The genus tower is the
scalar projection of $\Theta_\cA$ to each genus.  The shadow
tower is the arity projection of $\Theta_\cA$ on the primary
line.  The Koszul dual is the genus-$0$, arity-$2$ projection
(the binary pairing that determines $\cA^!$).

In the gravitational context ($\cA = \mathrm{Vir}_c$):
\begin{itemize}[nosep]
\item The genus tower gives $F_g = (c/2)\lambda_g^{\mathrm{FP}}$:
  the Bernoulli numbers.  Higher-genus gravitational partition
  functions.
\item The shadow obstruction tower gives $\sqrt{Q_{\mathrm{Vir}}}$: the
  quadratic extension $K = \Q(c)(t)(\sqrt{Q})$ with discriminant
  $-320c^2/(5c{+}22)$.  The arithmetic of quantum gravity.
\item The Koszul dual gives $\mathrm{Vir}_{26-c}$: the dual
  gravitational theory.  $c + c' = 26$ is Freudenthal--de~Vries.
\end{itemize}

All three are proved.  All three are outputs of the single
object $\Theta_{\mathrm{Vir}}$.  But they are three outputs, not
five levels of one chain.

\subsection{The correct picture}

The bar construction is a Fourier transform.  This is a theorem,
not a metaphor (Volume~I, Theorems~A, B, C).  The shadow obstruction tower
is a quadratic extension.  This is a theorem, not an analogy
(\S\ref{sec:arithmetic-gravity}).  But the claim that these two
facts are levels of a ``descent chain'' that continues to
$L$-functions and the Riemann hypothesis is false.

The reason it is false is not a failure of imagination; it is a
theorem.  The structural obstruction
(Remark~\ref{rem:structural-obstruction}) shows that algebraic
constraints from the MC equation act on the real spectral line,
while the zeta zeros live at complex spectral parameters.  No
amount of clever identification can bridge this gap: it is the
difference between a finite-dimensional algebraic structure
(the shadow obstruction tower, determined by three numbers $\kappa$, $\alpha$,
$S_4$) and an infinite-dimensional analytic object (the zeta
function, with infinitely many independent zeros).

The shadow obstruction tower \emph{detects} arithmetic---Hecke eigenforms
for lattice VOAs, quadratic fields for all VOAs at rational $c$---but
it does not \emph{control} the arithmetic.  The detection is a
theorem (Theorem~\ref{thm:shadow-spectral-correspondence}).  The
control would require a mechanism for the MC equation to force
zeros onto critical lines, and no such mechanism exists.

What remains is beautiful and proved.  The bar construction of a
chiral algebra is a Fourier transform that simultaneously produces
topological invariants (Mumford classes), arithmetic invariants
(quadratic fields), and algebraic invariants (Koszul duals) from
a single MC element.  The gravitational shadow obstruction tower, specialised
to rational central charge, assigns a number field to each minimal
model.  The Ising model lives over $\Q(\sqrt{5})$.  The self-dual
gravity at $c = 13$ lives over $\Q(\sqrt{870})$.  Whether these
number fields have structure beyond what is visible in the
specialisation table is an open question at the intersection of
vertex algebra theory and algebraic number theory.

The bar construction makes the question precise.  It does not
answer it.

\subsection{The four Dirichlet series and the Benjamin--Chang resolution}
\label{sec:four-series-bc-resolution}

The preceding subsections demonstrate that the descent chain
is broken.  But the \emph{reason} it is broken encodes a
positive insight about the arithmetic programme that deserves
explicit articulation.

\textbf{Four distinct Dirichlet series.}
The monograph constructs (or encounters) four Dirichlet series
from a chirally Koszul algebra~$\cA$:
\begin{center}
\small
\renewcommand{\arraystretch}{1.3}
\begin{tabular}{@{}cllp{4.5cm}@{}}
& \textbf{Object} & \textbf{Input}
  & \textbf{Where it lives} \\
\hline
\textup{(a)} & Shadow-metric Epstein $\varepsilon_{Q_L}(s)$
  & $(\kappa, \alpha, \Delta)$
  & Arity axis, genus $0$ \\[3pt]
\textup{(b)} & Constrained Epstein $\varepsilon^c_s$
  & Primary spectrum $\{(\Delta_i, d_i)\}$
  & Genus axis, genus $1$ \\[3pt]
\textup{(c)} & Shadow Dirichlet $\zeta_\cA(s)$
  & Shadow tower $\{S_r\}$
  & Arity axis, genus $0$ \\[3pt]
\textup{(d)} & Categorical zeta $\zeta^{\mathrm{DK}}_\fg(s)$
  & $\{\dim V_\lambda\}$
  & Lie-theoretic, no genus
\end{tabular}
\end{center}
Plus the Dirichlet-sewing lift
$S_\cA(u) = \zeta(u{+}1)\sum_i (\zeta(u) - H_{w_i-1}(u))$,
determined by the weight multiset~$W(\cA)$---a fifth object,
independent of the above four.

The ``negative results'' of the preceding subsections
(shadow zeta not in Selberg class, no Euler product, no $s$-reflection
functional equation) concern objects (a) and (c): the arity-axis
Dirichlet series.  Benjamin--Chang~\cite{BenjaminChang22} proved
that the Riemann zeta zeros control asymptotic scaling of
operators through object~(b): the constrained Epstein zeta.
These are \emph{different objects living on orthogonal axes} of
the bigraded shadow algebra $\cA^{\mathrm{sh}}_{r,g}$.

\textbf{Where the Riemann zeta enters.}
The Riemann zeta zeros~$\rho_n$ appear through two distinct
mechanisms, both on the genus side:
\begin{enumerate}[label=\textup{(\roman*)}]
\item \emph{Benjamin--Chang spectral decomposition.}
  The scattering factor $F_c(s) = \zeta(2s)/\zeta(2s{-}1)$
  in the Eisenstein spectral decomposition of
  $\hat Z^c = y^{c/2}|\eta|^{2c}\,Z$ on
  $\mathrm{SL}(2,\Z)\backslash\mathfrak{H}$
  has poles at $s = (1 + \rho_n)/2$.  The pole
  \emph{locations} are universal (from $\zeta(2s{-}1) = 0$);
  the pole \emph{residues} are CFT-specific (primary
  multiplicities integrated against Eisenstein residues).
  This is the spectral theory of~$\cM_{1,1}$.

\item \emph{Dirichlet-sewing lift.}
  For Heisenberg, $S_\cH(u) = \zeta(u)\zeta(u{+}1)$:
  the Riemann zeta enters through the weight multiset
  (the sewing sum $\sum_{n=1}^\infty n^{-u}$ over all
  integer modes of a weight-$1$ field).  For Virasoro,
  $S_{\mathrm{Vir}}(u) = \zeta(u{+}1)(\zeta(u) - 1)$:
  the $\zeta$-factor persists but the Euler product breaks.
\end{enumerate}
The Riemann zeta does \emph{not} enter through the shadow
obstruction tower.  The shadow Dirichlet series~(c) is entire,
has exponentially decaying coefficients, is not in the Selberg
class, and has no Euler product.  These are not ``negative
results about the Benjamin--Chang programme''---they are positive
structural theorems about the arity axis, which operates by
different principles (MC recursion, Koszul complementarity)
than the genus axis (Euler product, functional equation).

\textbf{The arity--genus orthogonality principle.}
In the bigraded MC element
$\Theta_\cA \in \cA^{\mathrm{sh}}_{r,g}$:
the shadow obstruction tower is the $(r, 0)$~slice, the genus
tower is the $(0, g)$~slice, and the constrained Epstein is
extracted from $(0, 1)$ by Rankin--Selberg.  The two axes
interact only through:
\begin{itemize}[nosep]
\item The MC equation, which couples $(r_1, g_1)$ to
  $(r_2, g_2)$ via the planted-forest corrections
  (the shadow visibility genus:
  $g_{\min}(S_r) = \lfloor r/2 \rfloor + 1$,
  Vol~I, Corollary~\ref*{cor:shadow-visibility-genus}).
\item The genus-$2$ non-diagonality: at genus~$\geq 2$,
  the off-diagonal sewing blocks~$R_{ij}$ couple the
  handle structure to the inter-handle OPE, breaking the
  genus-$1$ factorisation barrier.
\end{itemize}
The descent chain breaks at Level~$2 \to 3$ because it
attempts to descend along the arity axis to a target (zeros
of~$\varepsilon^c_s$) that lives on the genus axis.
The correct picture is the \emph{fan}: three independent
projections from a single MC element, with the Riemann zeta
entering at the genus projection (via~$\cM_{1,1}$) and absent
from the arity projection.

\textbf{The $\fsl_2$ intersection.}
The categorical zeta $\zeta^{\mathrm{DK}}_{\fsl_2}(s) = \zeta(s)$
(Vol~I, Proposition~\ref*{prop:categorical-zeta-sl2-riemann})
is the single point where the four Dirichlet series touch the
Riemann zeta function through their defining data.  For integrable
$V_k(\fsl_2)$: (d) $= \zeta(s)$, the Dirichlet-sewing lift
has $\zeta$-factors, and the constrained Epstein~(b) at integrable
level is a finite Dirichlet polynomial determined by the $k{+}1$
primary dimensions.  But the shadow Dirichlet~(c) remains entire
and non-classical even for $\fsl_2$---confirming that the arity
axis is structurally different from the genus axis, even in the
simplest case.


% =================================================================
\section{The analytic VOA programme}
\label{sec:analytic-voa-programme}
% =================================================================

The preceding analysis identifies a structural gap between the
algebraic bar-cobar machinery (which provides integrability via
$D^2 = 0$) and the positivity constraints needed for spectral
conclusions (Ramanujan bounds, $L$-value non-vanishing, zero
locations).  This section develops the \emph{analytic VOA
programme}: a framework in which positivity is axiomatized
alongside integrability, closing the gap by construction.

\subsection{The integrability--positivity dichotomy}
\label{sec:integrability-positivity}

The MC equation $D\Theta + \tfrac{1}{2}[\Theta,\Theta] = 0$ is
an \emph{integrability} constraint (flatness, quadratic
consistency).  Every spectral conclusion---critical-line zeros,
Ramanujan bounds, $L$-value non-vanishing---requires
\emph{positivity} (sign-definiteness, absolute-value bounds).
These are logically independent: a flat connection can have
eigenvalues of any absolute value.

Three proved obstructions instantiate this dichotomy at genus~$1$:
\begin{enumerate}[label=\textup{(\roman*)}]
\item \emph{Structural separation}
  \textup{(}Volume~I,
  Theorem~\ref*{thm:structural-separation}\textup{)}: the
  Rankin--Selberg integral $I(s;\cA)$ is holomorphic at all
  scattering poles $s = \rho/2$.  Unfolding erases the
  scattering matrix.
\item \emph{Real--complex separation}: the MC equation constrains
  spectral coefficients on the real spectral axis
  $s = \tfrac{1}{2} + it$, $t \in \bR$; zeta zeros live at
  complex spectral parameters.
\item \emph{Dimensional underdetermination}: the shadow
  obstruction tower is determined by three real numbers
  $(\kappa, \alpha, S_4)$; the zeta zeros form a countably
  infinite set of complex numbers.
\end{enumerate}
All three are \textbf{genus-$1$ theorems}.  At genus~$\geq 2$,
the Fock-diagonal structure breaks
(Theorem~\ref*{thm:genus2-non-collapse}), and the
B\"ocherer bridge
(Theorem~\ref*{thm:bocherer-bridge})
provides critical-line access to $L(\tfrac{1}{2},
\pi_f \times \chi_D)$, bypassing unfolding erasure.
The remaining gap is \textbf{positivity}: the MC equation
\emph{determines} central $L$-values but does not
\emph{force} non-vanishing.

\subsection{Axiom system: analytic unitary
  $\SCchtop$-algebra}
\label{sec:analytic-axioms}

\begin{definition}[Analytic unitary $\SCchtop$-algebra]%
\label{def:analytic-unitary-SC}
A \emph{logarithmic $\SCchtop$-algebra}
$(\cA, \{m_k\}, \{\omega_k\})$ in the sense of
Definition~\ref{def:log-SC-algebra} equipped with:
\begin{enumerate}[label=\textup{(AU\arabic*)}]
\item \textbf{Unitarity.}  A positive-definite invariant
  Hermitian form $\langle\mathord{-},\mathord{-}\rangle_n$ on
  each conformal weight space~$H_n$, satisfying the BPZ adjoint
  relation
  \[
    \langle m_2(a,b),\, c\rangle
    = \langle a,\, m_2^*(c,b)\rangle
  \]
  where $m_2^*$ is the BPZ conjugate
  $Y(e^{zL_1}(-z^{-2})^{L_0}a,\, z^{-1})$.

\item \textbf{Energy bounds.}  There exist $C > 0$ and
  $N \in \bZ_{\geq 0}$ such that for all homogeneous
  $a_1, \ldots, a_k$ with $a_i \in H_{\Delta_i}$:
  \[
    \|m_k(a_1, \ldots, a_k)\|
    \;\leq\;
    C^k \prod_{i=1}^k \|a_i\|
    \cdot \bigl(\textstyle\sum \Delta_i + 1\bigr)^N.
  \]
  This is the $A_\infty$ upgrade of the polynomial OPE growth
  condition verified for the standard landscape in
  Theorem~\ref*{thm:general-hs-sewing}.

\item \textbf{Analytic bar convergence.}  The bar complex
  $B(\cA)$ admits a Hilbert space completion
  $B^{\mathrm{Hilb}}_n(\cA,q)$ in the $q$-weighted BPZ norm
  \[
    \|[s^{-1}a_1|\cdots|s^{-1}a_n]\|_q^2
    \;:=\;
    q^{\sum \Delta_i}
    \prod_{i=1}^n \|a_i\|_{\Delta_i}^2
  \]
  for each $q \in (0,1)$.  The bar differential~$d_B$ is
  bounded, closable, and has densely defined formal
  adjoint~$d_B^*$ (the BPZ-conjugate creation operator).

\item \textbf{Bar Hodge property.}  The bar Laplacian
  $\Delta_B := d_B\, d_B^* + d_B^*\, d_B$
  is essentially self-adjoint with discrete spectrum on each
  finite bar degree, and
  $\ker(\Delta_B) \cong H^*(B(\cA))$
  \textup{(}the bar Hodge theorem\textup{)}.
\end{enumerate}
\end{definition}

\begin{proposition}[Essential self-adjointness]%
\label{prop:bar-laplacian-esa}
For any vertex algebra~$\cA$ satisfying~\textup{(AU1)},
the bar Laplacian~$\Delta_B$ is essentially self-adjoint on the
algebraic bar complex.  The spectrum is discrete with finite
multiplicities.
\end{proposition}

\begin{proof}[Proof sketch]
Each weight space $B_n(\cA)_{\mathrm{weight}=N}$ is
finite-dimensional (a basic VOA axiom: $\dim H_N < \infty$).
On a finite-dimensional Hilbert space, every symmetric operator
is self-adjoint.  The full operator
$\Delta_B = \bigoplus_{N,n} \Delta_B|_{N,n}$ is an orthogonal
direct sum of self-adjoint operators on finite-dimensional
spaces, hence essentially self-adjoint on the algebraic direct
sum (which is a core for the closure).  Discreteness follows
from finite-dimensionality of each sector.
\end{proof}

\begin{remark}[Spectral gap for Koszul algebras]%
\label{rem:koszul-spectral-gap}
If $\cA$ is chirally Koszul, bar cohomology is concentrated on
the diagonal (weight $=$ bar degree), so
$\ker(\Delta_B|_{N,n}) = 0$ for $(N,n)$ off the diagonal.
The smallest nonzero eigenvalue on each off-diagonal sector
gives a positive spectral gap $\mu_{N,n} > 0$.
\end{remark}

\subsection{Families satisfying the axioms}

\begin{center}
\small
\renewcommand{\arraystretch}{1.3}
\begin{tabular}{@{}llcccc@{}}
\textbf{Class} & \textbf{Family}
  & \textbf{(AU1)} & \textbf{(AU2)} & \textbf{(AU3)}
  & \textbf{(AU4)} \\
\hline
G & Heisenberg $\cH_\kappa$, $\kappa > 0$
  & \checkmark & \checkmark & \checkmark & \checkmark \\
G & Free $\beta\gamma$
  & \checkmark & \checkmark & \checkmark & \checkmark \\
L & $V_k(\fg)$, $k \in \bZ_{>0}$
  & \checkmark & \checkmark & \checkmark
  & \checkmark${}^*$ \\
C & Lattice $V_\Lambda$
  & \checkmark & \checkmark & \checkmark & \checkmark \\
M & Virasoro discrete series
  & \checkmark & \checkmark & \checkmark
  & Likely \\
M & Virasoro $c > 1$
  & \checkmark & \checkmark & Open${}^\dagger$
  & Open${}^\dagger$
\end{tabular}
\end{center}
\noindent
${}^*$\,At integrable level; generic level requires the
resonance-filtered topology.\\
${}^\dagger$\,Requires the analytic completion for class~M
(Open Problem at
\S\ref*{prob:thqg-X-boundary-bar}).

\subsection{The CKLW--Swiss-cheese identification}
\label{sec:cklw-sc-identification}

The Carpi--Kawahigashi--Longo--Weiner (CKLW) functor~\cite{CKLW18}
takes a strongly local unitary VOA~$V$ and produces a conformal
net $\{A(I)\}_{I \subset S^1}$ on the vacuum Hilbert space~$H_V$.
The Swiss-cheese framework identifies this functor as the
\emph{analytic completion of the open sector}
of $\SCchtop$:

\begin{center}
\small
$V \;\xrightarrow{\text{recognition}}
(\cA_{\mathrm{ch}}, \cA_{\mathrm{top}})
\;\xrightarrow{\text{open projection}}
\cA_{\mathrm{top}}
\;\xrightarrow{\text{clutching}}
\cA_{\mathrm{top}}|_{S^1}
\;\xrightarrow{\text{CKLW}}
\{A(I)\}$
\end{center}

\noindent Arrows 1--3 are proved in this volume
(Theorems~\ref{thm:recognition-SC},
\ref{thm:e1-factorization-disjoint},
Definition~\ref{def:oc-factorization-category}(vi)).
Arrow~4 is the CKLW theorem~\cite{CKLW18}.
The identification explains:
\begin{itemize}[nosep]
\item \emph{Why conformal nets are $E_1$}: they are the open
  (topological) colour of $\SCchtop$, with the holomorphic
  colour integrated out by smearing.
\item \emph{Why strong locality is the key hypothesis}: it is
  the analytic version of $E_1$-locality in the open sector.
\item \emph{Why DHR sectors match line operators}: both
  classify modules of the open-sector factorization category.
\end{itemize}
For the affine lineage at integrable level, the Kazhdan--Lusztig
equivalence, the CKLW theorem, and the Wassermann construction
provide all ingredients: the line-operator category
$\cC_{\mathrm{line}}^{\mathrm{red}} \simeq
\mathrm{Rep}_q(\fg)$ (Theorem~\ref{thm:affine-monodromy-identification})
matches the DHR category of the conformal net.  At generic level
or for class~M algebras, the identification is conjectural and
reduces to CKLW's strong locality hypothesis.

\subsection{The Tannakian fibre functor}
\label{sec:tannakian-fibre-functor}

For a unitary VOA~$V$ at integrable level~$k$, the BPZ inner
product defines a faithful exact $*$-monoidal functor
\[
  \omega_{\mathrm{BPZ}}\colon
  \cC_{\mathrm{line}}^{\mathrm{int}} \longrightarrow \mathrm{Hilb}.
\]
By Tannakian reconstruction, $H = \End(\omega_{\mathrm{BPZ}})
\cong U_q(\fg)$ with the compact $*$-structure
($q = e^{i\pi/(k+h^\vee)}$ on the unit circle).
The quantum trace $\mathrm{tr}_q = \mathrm{tr}(K^{2\rho}\,
\cdot\,)$ is a positive functional (the Haar state), and the
GNS Hilbert space is the state space on
$S^1 \times D^2$.

The three-tier picture (AP51) maps precisely:
\begin{itemize}[nosep]
\item Tier~(i) (pole-free BD-commutative $E_\infty$ subclass): forgetful functor to
  $\mathrm{Vect}$, commutative Hopf algebra.
\item Tier~(ii) (local $E_\infty$ with OPE poles, class~L): $\omega_{\mathrm{BPZ}}$
  exists, compact quantum group.
\item Tier~(iii) (still local $E_\infty$, class~M): \textbf{no} fibre functor to
  $\mathrm{Hilb}$ can make the braiding unitary---the
  $A_\infty$ tower prevents it
  (Remark~\ref{rem:virasoro-qg-unitarity}).
\end{itemize}

\subsection{The analytic Swiss-cheese operad}
\label{sec:analytic-sc-operad}

The HS-sewing condition
(Theorem~\ref*{thm:general-hs-sewing})
is not merely a condition on algebras: it defines a topology
on the operad~$\SCchtop$ itself.  The \emph{analytic
Swiss-cheese operad at sewing parameter~$q$} is
\[
  \mathsf{SC}^{\mathrm{ch,top,an}}_q
  \;:=\;
  L^2_q\text{-completion of }
  C_*(\SCchtop),
\]
where the $L^2_q$-norm uses the FM volume form and
$q$-weighting by conformal weight.  The composition maps
extend to bounded operators because $\FM_k(\bC)$ is compact
with corners and the FM Jacobian estimate ensures logarithmic
forms are $L^1$-integrable.

\textbf{Equivalence theorem.}
An algebra satisfies the HS-sewing condition at
parameter~$q$ if and only if it is an algebra over
$\mathsf{SC}^{\mathrm{ch,top,an}}_q$.  This makes the analytic
condition \emph{operadic} rather than algebra-by-algebra.
The entire standard landscape (all seven families) consists
of algebras over $\mathsf{SC}^{\mathrm{ch,top,an}}_q$ for every
$q \in (0,1)$.

\subsection{The bar trace formula}
\label{sec:bar-trace-formula}

The bar Laplacian~$\Delta_B$ admits a heat kernel
$\Theta_B(t) = \Tr(e^{-t\Delta_B})$.
The \emph{bar trace formula} relates the spectral side
(eigenvalues of~$\Delta_B$) to a geometric side (a sum
over boundary strata of $\FM_k(\bC)$):
\[
  \sum_n h(\lambda_n)
  \;=\;
  h(0)\,\chi(B^{\mathrm{ann}})
  \;+\;
  \sum_{k \geq 2}\;
  \sum_{T \in \mathrm{Trees}(k)}
  \hat{h}(\ell_T)\, m_T(\cA)
  \;+\;
  \text{planted-forest corrections},
\]
where $\ell_T$ is the bar length of tree~$T$ (the
weighted sum of collision residues along edges)
and $m_T(\cA)$ is the $A_\infty$ amplitude on the
tree stratum~$D_T$.  For the Heisenberg algebra,
the one-particle reduction recovers the Selberg trace
formula: bar geodesics $=$ plumbing parameters,
$\ell_n = n$, and
$\Theta_B^{\cH}(t) = \kappa\, e^{-t}/(1 - e^{-t})$.
The planted-forest corrections vanish for class~G
and contribute at class~L and beyond via the cubic
obstruction of Theorem~\ref{thm:fm3-higher-obstruction-synthesis}.

\subsection{IC-purity and Hodge--Riemann positivity}
\label{sec:ic-purity-positivity}

The \emph{collision-weight $t$-structure} on
$D(B^{\mathrm{an}}(\cA)\text{-comod})$, defined by the
weight filtration from Saito's mixed Hodge module theory,
has amplitude equal to the shadow depth~$r_{\max}$.  On the
semisimple-purity locus (classes~G, L, C), the heart carries
a \textbf{positive-definite form}~$Q_{\mathrm{IC}}$ on the
bar harmonic space, induced by the Hodge--Riemann bilinear
relations on IC complexes of the FM compactification
(the de~Cataldo--Migliorini theorem).

This upgrades IC purity from a \emph{vanishing} statement
(weight spectral sequence degenerates $\Rightarrow$
Koszulness) to a \emph{positivity} statement (the surviving
cohomology carries a definite form).  This is the chiral
analogue of the Elias--Williamson proof of Kazhdan--Lusztig
positivity.

For class~M (Virasoro), $Q_{\mathrm{IC}}$ is
\emph{indefinite}---the non-split extensions in
$\Ext^1_{\mathrm{MHM}}(\mathrm{gr}^4 L,\, \mathrm{gr}^2 L)$
from the quartic pole create off-diagonal terms of either sign.
The \emph{IC signature} is a new invariant measuring the
failure of positivity for non-semisimple monodromy.

\subsection{The genus-$2$ spectral gap programme}
\label{sec:genus2-spectral-gap}

The genus-$2$ bar Laplacian
$\Delta_B^{(2)} = d_B^{(2)}(d_B^{(2)})^*
+ (d_B^{(2)})^* d_B^{(2)}$
inherits a $2 \times 2$ block structure from the genus-$2$
sewing kernel~$K_2$:
\[
  \Delta_B^{(2)}
  \;=\;
  \Delta_B^{\mathrm{diag}}
  \;+\;
  \Delta_B^{\mathrm{off\text{-}diag}},
\]
where $\Delta_B^{\mathrm{diag}}$ is the direct sum of
two genus-$1$ bar Laplacians and
$\Delta_B^{\mathrm{off\text{-}diag}}$ encodes the
inter-handle coupling through the Schur complement
$S = K_{q_2} + R_{21}(1 - K_{q_1})^{-1}R_{12}$.

The off-diagonal coupling is exactly the structural mechanism
that breaks the genus-$1$ separation barrier
(Remark~\ref{rem:genus2-arithmetic-significance}).
The \textbf{surviving research programme} is:
\begin{enumerate}[label=\textup{(\arabic*)}]
\item Construct $\Delta_B^{(2)}$ explicitly using the
  genus-$2$ Fredholm determinant theorem.
\item Prove a spectral gap via screening-type inequalities
  from the off-diagonal coupling.
\item Use $\mathrm{Sp}(4)$ representation theory
  (Arthur packets) to show the spectral gap forces
  \emph{temperedness}: the Arthur packets for
  $\mathrm{Sp}(4)$ distinguish Saito--Kurokawa
  (endoscopic, non-tempered) from stable (tempered).
  The off-diagonal sewing coupling detects exactly this
  factorisation structure.
\item Temperedness $=$ the Ramanujan bound for
  $\mathrm{GSp}(4)$ automorphic forms
  (Weissauer's theorem for holomorphic cusp forms).
\end{enumerate}
This is the canonical bridge from bar-complex spectral
geometry to arithmetic, identified as the most promising
direction by the full research campaign.
The genus-$1$ structural obstruction is a
\emph{genus-$1$ theorem}, not an absolute barrier.

\subsection{The algebra--geometry--arithmetic trichotomy}
\label{sec:algebra-geometry-arithmetic}

The preceding sections establish a precise negative result:
the bar complex of a chiral algebra cannot prove the Riemann
Hypothesis.  But the argument reveals a deeper structural
picture that deserves careful articulation, because the
\emph{reason} it fails illuminates the relationship between
algebra, geometry, and arithmetic more sharply than the
failure itself.

\subsubsection{Three actors, one meeting point}

The partition function $Z(\tau)$ of a chiral algebra~$\cA$ is
the meeting point of three mathematical worlds:

\begin{center}
\small
\renewcommand{\arraystretch}{1.4}
\begin{tabular}{@{}lll@{}}
\textbf{Actor} & \textbf{What it encodes} & \textbf{Source} \\
\hline
The bar complex $\barB(\cA)$
  & The OPE data of~$\cA$
  & Algebra (the chiral product) \\[3pt]
The hyperbolic Laplacian $\Delta_H$
  & The spectral theory of $\cM_{1,1}$
  & Geometry (the modular surface) \\[3pt]
The Riemann zeta function $\zeta(s)$
  & The scattering poles of~$\Delta_H$
  & Arithmetic (the primes)
\end{tabular}
\end{center}

\noindent The bar complex knows everything about~$\cA$.  The
spectral theory of $\cM_{1,1} = \mathrm{SL}(2,\bZ)\backslash
\mathfrak{H}$ knows everything about~$\zeta$.  The partition
function~$Z(\tau)$ is the object on which both act, but
they act on \emph{different aspects} of it: the bar complex
determines $Z$ as a $q$-series (its algebraic content), while
$\Delta_H$ decomposes $Z$ as a function on a hyperbolic
surface (its spectral content).

\textbf{The bar complex does not encode the zeta zeros.}
It encodes the \emph{algebra}, which produces $Z(\tau)$
as a convergent $q$-series $\sum d(n)\, q^{n - c/24}$.
The zeta zeros are a property of how $Z(\tau)$ \emph{sits
on~$\cM_{1,1}$}---they come from the geometry of the modular
surface, not from the algebra of the chiral product.  The
chain is:
\[
  \barB(\cA)
  \;\xrightarrow{\text{OPE data}}\;
  Z(\tau) \text{ as $q$-series}
  \;\xrightarrow{\text{place on $\cM_{1,1}$}}\;
  Z(\tau) \text{ as function on $\mathrm{SL}(2,\bZ)
  \backslash\mathfrak{H}$}
  \;\xrightarrow{\Delta_H}\;
  \text{spectral decomposition}
  \;\ni\; \zeta\text{ zeros}.
\]
The bar complex controls the first arrow.  The zeta zeros
enter at the last.  The gap is the passage from ``$Z$ as a
$q$-series'' to ``$Z$ as a function on a hyperbolic surface
whose Laplacian has continuous spectrum governed by~$\zeta$.''

\subsubsection{What encodes the Eisenstein/continuous spectrum}

The Roelcke--Selberg decomposition of $\hat Z^c =
y^{c/2}|\eta|^{2c}\,Z$ on $\cM_{1,1}$ splits~$L^2$ into
three parts:
\[
  \hat Z^c(\tau)
  \;=\;
  c_{\mathrm{res}}
  \;+\;
  \sum_j \beta_j\,\varphi_j(\tau)
  \;+\;
  \frac{1}{4\pi}\int_0^\infty
  \alpha(t)\,E_{1/2+it}(\tau)\,dt,
\]
where $\varphi_j$ are Maass cusp forms (discrete spectrum)
and $E_s$ is the Eisenstein series (continuous spectrum).  The
scattering matrix
$\varphi(s) = \Lambda(1{-}s)/\Lambda(s)$,
whose poles are the scaled zeta zeros at $s = \rho/2$, governs
the Eisenstein series through
$E_s = y^s + \varphi(s)\,y^{1-s} + \ldots$

This scattering matrix is a property of the
\textbf{hyperbolic Laplacian} $\Delta_H =
-y^2(\partial_x^2 + \partial_y^2)$ on the non-compact
surface $\mathrm{SL}(2,\bZ)\backslash\mathfrak{H}$.  It is
\emph{pure spectral geometry}: it exists independently of any
chiral algebra, any partition function, or any physical theory.
The zeta zeros are built into the geometry of the modular
surface through the Euler product of~$\zeta(s)$, which encodes
the distribution of primes, which controls the arithmetic of
the lattice $\mathrm{SL}(2,\bZ)$ acting on~$\mathfrak{H}$.

The bar complex of~$\cA$ knows nothing about this geometry.
It produces $Z(\tau)$ as a $q$-series; what happens when that
$q$-series is placed on the modular surface and decomposed
spectrally is a question about the \emph{surface}, not about
the \emph{algebra}.

\subsubsection{What the CFT ``knows''}

If the zeta zeros are encoded in the spectral decomposition
of $Z(\tau)$ on~$\cM_{1,1}$
(the content of Benjamin--Chang~\cite{BenjaminChang22}),
then what CFT property corresponds to the zeros?  The answer
is precise:

\textbf{The zeros control the equidistribution rate of
$Z(\tau)$ on the modular surface.}

The Eisenstein spectrum governs how fast a function
on~$\cM_{1,1}$ approaches its average under modular
averaging.  For $Z(\tau)$ with $d(n) \geq 0$ (unitarity),
the equidistribution rate is controlled by the spectral
gap $\lambda_1$ of $\Delta_H$ on the appropriate congruence
quotient---and this spectral gap \emph{is} the Selberg
eigenvalue conjecture ($\lambda_1 \geq \tfrac{1}{4}$),
which \emph{is} the Ramanujan conjecture for
$\mathrm{GL}(2)$ Maass forms.

The CFT does not ``know'' where the zeros are through its
algebraic structure (the bar complex, the OPE, the MC
equation).  It ``knows'' through its
\textbf{equidistribution rate on the moduli space}---a
property that requires the full analytic structure of
$Z(\tau)$ as a function on a hyperbolic surface, not just
its $q$-expansion coefficients.

\subsubsection{Why the analytic VOA programme is the right
  direction}

This analysis identifies what is missing from the algebraic
framework and what an analytic upgrade would need to provide.

The bar complex gives you the $q$-series.  What you need is
the \emph{hyperbolic geometry of the function that
$q$-series defines}: its growth, its equidistribution, its
spectral decomposition.  The passage from $q$-series to
hyperbolic geometry requires:
\begin{enumerate}[label=\textup{(\roman*)}]
\item \emph{A positive-definite inner product} on the
  function space (axiom AU1), to define the Hilbert space
  $L^2(\cM_{1,1})$ on which $\Delta_H$ acts.
\item \emph{Growth bounds} on $Z(\tau)$ as
  $\mathrm{Im}(\tau) \to \infty$ (axiom AU2), to place
  $Z$ in the correct Sobolev space.
\item \emph{Convergence of the spectral decomposition}
  (axiom AU3), to make the Roelcke--Selberg expansion
  meaningful.
\item \emph{Spectral gap control} (axiom AU4), to
  quantify the equidistribution rate.
\end{enumerate}
These are precisely the axioms (AU1)--(AU4) of the analytic
unitary $\SCchtop$-algebra
(Definition~\ref{def:analytic-unitary-SC}).  The analytic
VOA programme does not attempt to make the bar complex see
the zeta zeros---it cannot, because the zeros are geometric,
not algebraic.  Instead, it provides the analytic
infrastructure to study how $Z(\tau)$, produced by the bar
complex, \emph{sits on the modular surface}---the
meeting point where algebra, geometry, and arithmetic
converge.

\subsubsection{The Steinberg principle}

The Steinberg variety presents the Hecke algebra.  The bar
complex presents the Swiss-cheese algebra.  Neither presents
the Riemann zeta function.  The zeta function is presented
by the \emph{modular surface itself}---its spectral theory,
its scattering matrix, its Selberg trace formula.  The
correct analogy is not ``the bar complex is to $\zeta$ as the
Steinberg variety is to Hecke,'' but rather: the bar complex
and $\zeta$ are two sides of the partition function, meeting
at $Z(\tau)$ on~$\cM_{1,1}$, as the Steinberg variety and the
Springer resolution are two sides of the flag variety, meeting
at the moment map.

The bar complex is the algebraic side: it produces the
function.  The hyperbolic Laplacian is the geometric side: it
decomposes the function.  The Riemann Hypothesis is the
conjecture that these two sides interact in the most
harmonious possible way---that the algebraic function
produced by the bar complex equidistributes on the geometric
surface at the optimal rate.  Whether this harmony holds is a
question about the \emph{relationship} between algebra and
geometry, not a question answerable from either side alone.


% =================================================================
\section{Session notes: 1 April 2026}
\label{sec:session-2026-04-01}
% =================================================================

Six findings from the session.  The first is a theorem; the
remaining five are its consequences.


% -----------------------------------------------------------------
\subsection{The DS-HPL transfer theorem}
\label{sec:ds-hpl-transfer}
% -----------------------------------------------------------------

The single strongest new result.

\begin{theorem}[DS-HPL transfer; proved]
\label{thm:ds-hpl-transfer-wn}
Let $\fg$ be a simple Lie algebra, $k$ a non-critical level, and
$\cW_k(\fg)$ the principal Drinfeld--Sokolov reduction of
$V_k(\fg)$.  The homological perturbation lemma transfers the
dg-shifted Yangian datum $(m, r, \Delta)$ from the affine side to
the $\cW$-algebra side via the BRST strong deformation retract
$(\iota, p, h)$ with $\mathrm{gh}(h) = -1$.  The transferred
coproduct $\Delta^W$ is strictly primitive at every arity
(equation~\eqref{eq:gravity-primitive-wn}), with no corrections.
\end{theorem}

\begin{proof}
The ghost-number argument of
Theorem~\ref{thm:gravitational-primitivity-wn}
(\S\ref{sec:primitive-coproduct}) applies verbatim: every HPL
tree with an internal vertex shifts ghost number away from~$0$,
and $p \otimes p$ annihilates the result.
\end{proof}

The theorem says that Drinfeld--Sokolov reduction kills the
coproduct channel completely.  Only the $r$-matrix channel
survives.  The $\cW$-algebra inherits the $\Ainf$ products
$m_k^W$ (infinitely many, for class~M) and the collision
residue $r^W(z)$, but its coproduct is as trivial as a
one-dimensional algebra's.  It produces nothing.  It only
braids.

\sep

A remark on what is transferred and what is not.  The SDR
$(\iota, p, h)$ transfers three structures simultaneously:
\begin{center}
\small
\renewcommand{\arraystretch}{1.15}
\begin{tabular}{@{}lll@{}}
\textbf{Structure} & \textbf{Transferred result} & \textbf{Status} \\
\hline
$\Ainf$ products $m_k$ & nontrivial for all $k \geq 2$
  & class~M \\[3pt]
$r$-matrix $r(z)$ & nontrivial, all poles survive
  & quasi-trigonometric \\[3pt]
Coproduct $\Delta_z$ & primitive (all corrections vanish)
  & proved
\end{tabular}
\end{center}

\medskip

The $\Ainf$ products are nontrivial because the perturbation
$\delta_{\mathrm{BRST}}$ has ghost-number-preserving components
that survive the projection $p$.  The $r$-matrix is nontrivial
because it arises from the collision residue, which is a
genus-$0$ operation with no homotopy insertions at the output.
The coproduct is primitive because it is a genus-$0$,
arity-$2$, \textit{splitting} operation: the output has two
tensor factors, and every homotopy insertion shifts ghost number
into the range that $p \otimes p$ annihilates.

The asymmetry is that collision (computing $r(z)$) applies $p$
once, while splitting (computing $\Delta_z$) applies
$p \otimes p$: the tensor product doubles the projection's
power to kill ghost-shifted outputs.


% -----------------------------------------------------------------
\subsection{Gravitational primitivity as physical principle}
\label{sec:gravitational-primitivity-physical}
% -----------------------------------------------------------------

The physical content of gravitational primitivity
(\S\ref{sec:primitive-coproduct}): in three dimensions, there are
no local gravitational degrees of freedom and no propagating
gravitons.  The coproduct channel is trivial
(equation~\eqref{eq:gravity-primitive-wn}), and only the
$r$-matrix channel survives.  BTZ black holes are monodromies of
the $r$-matrix braiding, not vertices of the coproduct.

\sep

The $(m, \Delta)$-complexity plane organises the five standard
HT theories:

\begin{center}
\small
\renewcommand{\arraystretch}{1.15}
\begin{tabular}{@{}lcc@{}}
\textbf{Theory} & $m$-depth & $\Delta$-depth \\
\hline
Chern--Simons & $2$ (quadratic) & $\geq 1$ (nontrivial) \\[3pt]
Twisted holography & $2$ & $\geq 1$ \\[3pt]
M2 & $2$ & $\geq 1$ \\[3pt]
Gravity & $\infty$ (class M) & $0$ (primitive) \\[3pt]
M5 & $\infty$ & $0$
\end{tabular}
\end{center}

\medskip

\noindent
Gravity sits at $(\infty, 0)$: maximal $\Ainf$ complexity,
vanishing coproduct complexity.  Gauge theory sits at
$(2, \geq 1)$: minimal $\Ainf$ complexity, nontrivial
coproduct.  DS reduction moves a theory from the gauge corner
to the gravity corner by killing the coproduct and inflating
the $\Ainf$ tower.


% -----------------------------------------------------------------
\subsection{Five worked examples as $\Theta^{\mathrm{oc}}$
  instantiations}
\label{sec:five-oc-instantiations}
% -----------------------------------------------------------------

The completed eight-fold modular Koszul datum
$\Pi^{\mathrm{oc}}_X(\cA)$ for the five canonical HT systems:

\medskip

\noindent
\textbf{1.\ Chern--Simons.}\enspace
$\cA = V_k(\fg)$.  Class~L, shadow depth $3$.
$\kappa + \kappa^! = 0$ (Feigin--Frenkel).  Coproduct:
nontrivial (Casimir cross-term).  Sewing: determinant.  The
entire $\Ainf$ structure is formal: $m_k^H = 0$ for $k \geq 3$ on cohomology.
The modular Koszul datum is finite-dimensional at each arity.

\medskip

\noindent
\textbf{2.\ Three-dimensional gravity.}\enspace
$\cA = \mathrm{Vir}_c$.  Class~M, shadow depth $\infty$.
$\kappa + \kappa^! = 13 \neq 0$.  Coproduct: \textbf{primitive}
(ghost-number obstruction from DS).  Sewing:
Pfaffian-squared.  The quartic contact invariant
$Q^{\mathrm{ct}} = 10/[c(5c+22)]$ controls the entire infinite
tower.  The $r$-matrix
$r(z) = (c/2)/z^3 + 2T/z$ carries all scattering.

\medskip

\noindent
\textbf{3.\ Twisted holography.}\enspace
$\cA = V_k(\fg)$ (affine Kac--Moody, Costello--Gaiotto).
Class~L, shadow depth $3$.  Multi-channel: for $\cW_3$, the
genus-$2$ cross-channel contribution is
\[
F_2^{\mathrm{cross}} \;=\; \frac{c + 120}{16c}.
\]
The holographic dictionary is the Koszul triangle read through
Theorems~A--H.

\medskip

\noindent
\textbf{4.\ M2-brane.}\enspace
$\cA = \mathrm{DDCA}_k(\mathfrak{gl}_K)$.  Class~L.  Coproduct:
nontrivial.  Two-variable $r$-matrix $r(z_1, z_2)$.  No DS
reduction: the $\cW$-algebra structure is absent.  This is
gauge theory, not gravity, and the nontrivial coproduct
confirms it.

\medskip

\noindent
\textbf{5.\ M5-brane / $6$d $(2,0)$.}\enspace
$\cA = \cW_\infty(K)$ at large~$N$.  Class~M, shadow depth
$\infty$.  Coproduct: \textbf{primitive} (same ghost-number
argument as Virasoro, since $\cW_N$ arises by DS reduction).
Per-spin $r$-matrix.  Shadow depth increases under DS.

\sep

The shadow dichotomy across the five theories:

\begin{center}
\small
\renewcommand{\arraystretch}{1.15}
\begin{tabular}{@{}lcl@{}}
\textbf{Lineage} & \textbf{Class} & \textbf{Mechanism} \\
\hline
Current-algebra (CS, M2) & L & finite towers, no DS \\[3pt]
$\cW$-algebra (gravity, $\cW_3$, M5) & M & infinite towers,
  DS origin
\end{tabular}
\end{center}

\medskip

DS is the shadow-depth escalator.  It converts class~L (finite,
formal) to class~M (infinite, genuinely $\Ainf$).  The
ghost-number filtration that kills the coproduct is the same
filtration that inflates the $\Ainf$ tower.  The two effects
are a single mechanism.


% -----------------------------------------------------------------
\subsection{The open-sector primitive}
\label{sec:open-sector-primitive-session}
% -----------------------------------------------------------------

The primitive object is $\cC_{\mathrm{op}}$, not $\cA$ or
$\barB(\cA)$.

For Chern--Simons: $\cC_{\mathrm{op}} = V_k(\fg)\text{-mod}$,
the category of $V_k(\fg)$-modules as a factorisation
dg-category with meromorphic tensor product.  The boundary
algebra $A_b = \End(\mathrm{trivial})$ recovers $V_k(\fg)$.
The chiral derived center $\Zder = C^\bullet_{\mathrm{ch}}(\cA_b,
\cA_b)$ classifies bulk operators.

Modularity is trace plus clutching on $\cC_{\mathrm{op}}$, not
an axiom on the closed algebra.  The annulus trace
\[
\mathrm{Tr}\colon K_0(\cC_{\mathrm{op}})
\;\longrightarrow\;
Z_1(\cA)
\]
gives the genus-$1$ shadow.  The full genus-$g$ amplitude is the
composition of $g$ clutching maps on $\cC_{\mathrm{op}}$.  The
modular $S$-matrix is a derived property of the trace, not an
input to the theory.

The boundary algebra $A_b$ is a chart: a Morita-dependent
presentation of $\cC_{\mathrm{op}}$.  Different compact
generators produce Morita-equivalent charts.  The category is
the invariant; the algebra is a coordinate system.


% -----------------------------------------------------------------
\subsection{BV/BRST $=$ bar: honest status}
\label{sec:bv-brst-bar-honest}
% -----------------------------------------------------------------

Three levels of identification.  The first is proved; the
second is conditionally proved; the third is conjectural.

\begin{center}
\small
\renewcommand{\arraystretch}{1.15}
\begin{tabular}{@{}llc@{}}
\textbf{Genus} & \textbf{Claim} & \textbf{Status} \\
\hline
$g = 0$ & algebraic BRST $=$ bar differential & proved \\[3pt]
$g = 0$, higher arity
  & Feynman-diagrammatic $m_k$ $=$ bar $m_k$
  & proved \\[3pt]
$g \geq 1$ & BV quantum master equation $=$ curved bar
  & conjectural
\end{tabular}
\end{center}

\medskip

The bar complex IS well-defined at all genera: the inductive
genus determination from Volume~I combined with the HS-sewing
theorem produces bar amplitudes at every genus.  What is open
is whether these bar amplitudes coincide with the BV-BRST
partition function at genus $g \geq 1$.  This is
\texttt{conj:master-bv-brst}.

The identification at genus~$0$ is an algebraic comparison:
the BRST differential $Q = \partial + \partial^{\mathrm{BRST}}$,
restricted to the physical subspace, reproduces the bar
differential $d_{\barB}$.  The identification at genus~$1$
would require matching the BV quantum master equation
$\hbar\Delta_{\mathrm{BV}}S + \tfrac{1}{2}\{S, S\} = 0$ with
the curved bar relation $\dfib^{\,2} = \kappa \cdot \omega_1$.
No obstruction to this matching is known; no proof exists.

This session corrected four locations in the manuscript that
overclaimed genus-$1$ BV/BRST $=$ bar as proved.  The bar
complex is the proved object.  The BV-BRST comparison is the
conjectural bridge.


% -----------------------------------------------------------------
\subsection{The standalone paper}
\label{sec:standalone-paper}
% -----------------------------------------------------------------

\textit{Shadow Towers and Algebraicity}, $25$ pages, eight
sections.

\begin{enumerate}[label=\S\arabic*.,nosep]
\item Introduction and motivation.  The bar construction as
  categorical logarithm; the shadow obstruction tower as its Taylor
  expansion.
\item The MC algebra.  The modular cyclic deformation complex
  $\mathrm{Def}^{\mathrm{mod}}_{\mathrm{cyc}}(\cA)$ and the
  universal element $\Theta_\cA$.
\item The shadow obstruction tower.  Projection to the primary line;
  recursive computation of $S_r$ from the MC equation.
\item The Riccati algebraicity theorem.  The MC equation on a
  single primary line reduces to a Riccati ODE whose solution
  is $H(t) = t^2\sqrt{Q_L(t)}$, algebraic of degree~$2$.
  Full proof.
\item Four-class classification.  The critical discriminant
  $\Delta = 8\kappa S_4$ partitions modular Koszul algebras
  into four classes:
  \begin{enumerate}[label=(\alph*),nosep]
  \item \textbf{G} (Gaussian): $r_{\max} = 2$.  Heisenberg.
  \item \textbf{L} (Lie/tree): $r_{\max} = 3$.  Affine
    Kac--Moody.
  \item \textbf{C} (contact/quartic): $r_{\max} = 4$.
    $\beta\gamma$.
  \item \textbf{M} (mixed): $r_{\max} = \infty$.  Virasoro,
    $\cW_N$.
  \end{enumerate}
\item Explicit shadow obstruction towers.  Complete computation for the
  standard landscape: Heisenberg, affine $\fsl_2$,
  $\beta\gamma$, Virasoro, $\cW_3$.  Tables of $S_r$ through
  arity~$8$.
\item $\cW_3$ cross-channel genus-$2$ computation.
  Multi-generator shadow obstruction tower with two primary lines
  ($T$ and $W$), propagator variance $\delta_{\mathrm{mix}}$,
  genus-$2$ cross-channel formula.
\item Outlook.  Pixton conjecture (class~M algebras generate
  the Pixton ideal via infinite MC tower), analytic completion
  programme, non-standard families.
\end{enumerate}

The paper extracts from Volume~I the self-contained story of
the shadow obstruction tower, from its definition through the algebraicity
theorem to the four-class classification and the explicit
standard-landscape computations.  It is the first accessible
entry point to the monograph's nonlinear characteristic layer.


% =================================================================
\section{The graviton is composite}
\label{sec:graviton-composite}
% =================================================================

The Sugawara stress tensor
\[
T \;=\; \frac{1}{2(k+h^\vee)}\sum_a \normord{J^a J_a}
\]
is quadratic in affine currents.  It is not a fundamental field.
The graviton is composite.

The BRST complex of the Drinfeld--Sokolov reduction remembers
this: the transferred coproduct is strictly primitive
(Theorem~\ref{thm:gravitational-primitivity-wn},
equation~\eqref{eq:gravity-primitive-wn}).  No cross-term, no
splitting, at every arity and every principal DS reduction.

\sep

In gauge theory: the gluon $J^a$ IS the connection.
Fundamental.  Ghost number~$0$.  No passage through~$h$.  The
BRST homotopy does not obstruct the transferred coproduct.
The Casimir cross-term survives projection:
\[
\Delta_z(T_{\mathrm{aff}})
\;=\;
\tau_z(T_{\mathrm{aff}}) \otimes 1
\;+\; 1 \otimes T_{\mathrm{aff}}
\;+\; \frac{1}{k+h^\vee}\,\Omega(w{+}z,\, w),
\]
where $\Omega = \sum_a J^a \otimes J_a$ is the split Casimir.
A production vertex.  A gluon bound state dissociates.

In gravity: the graviton $T$ IS the curvature.  Composite.
Ghost-burdened.  Every HPL tree at ghost number~$0$ dies under
projection.  Gravitons do not split.  They braid.

\sep

The $r$-matrix $r^{\mathrm{Vir}}(z) = (c/2)/z^3 + 2T/z$
carries all gravitational scattering.  The cubic pole: the bar
kernel $d\log(z-w)$ absorbs one power from the quartic OPE.
The Virasoro OPE has poles at $z^{-4}$, $z^{-2}$, $z^{-1}$.
The $r$-matrix has poles at $z^{-3}$, $z^{-1}$.  The
logarithmic derivative is always one degree milder than the
function itself.  This absorption is not an approximation.  It
is the structural content of calling the bar construction a
categorical logarithm.

\sep

The five-theory comparison (see \S\ref{sec:five-theories})
confirms the pattern: DS reduction creates depth (L $\to$ M)
and kills the coproduct.  The distinction is structural
(the ghost-number filtration), not parametric (the central charge).

\sep

The bar complex of a chiral algebra knows more about quantum
gravity than any perturbative expansion.

It knows the exact genus-$g$ free energy at every genus, in
closed form:
$F_g = (c/2) \cdot \lambda_g^{\mathrm{FP}}$, the
$\hat{A}$-genus coefficients times the modular characteristic.

It knows the entire infinite tower of higher-arity corrections,
as the Taylor expansion of an algebraic function of degree~$2$:
$H(t) = t^2\sqrt{Q_{\mathrm{Vir}}(t)}$, three numbers
$(\kappa, \alpha, S_4)$, one quadratic form, one square root.

It knows that the gravitational coproduct is primitive
(\S\ref{sec:primitive-coproduct}): black holes do not split.

It knows the arithmetic.  At rational central charge $c_0$,
the shadow metric $Q_{\mathrm{Vir}}$ specialises to a specific
quadratic polynomial over~$\Q$, and
$K(c_0) = \Q(\sqrt{d(c_0)})$ is a specific real quadratic
number field (the critical discriminant $\Delta = 40/(5c{+}22) > 0$
for $c > 0$).  The Bekenstein--Hawking entropy $S = 2\pi\sqrt{ch/6}$
is the leading coefficient.  The polynomial discriminant
$\mathrm{disc}(Q_{\mathrm{Vir}}) = -320c^2/(5c+22) < 0$ ensures
the shadow obstruction tower is infinite (complex conjugate roots); the
critical discriminant $\Delta > 0$ produces a real quadratic extension.
The class number of $K(c_0)$ measures the arithmetic complexity
of gravity at that central charge.  The Ising model lives over
$\Q(\sqrt{5})$ (\S\ref{sec:gravity-rational-c}).
The self-dual gravity at $c = 13$ lives over $\Q(\sqrt{870})$,
and the complementarity constant $6960 = 2^4 \cdot 3 \cdot 5
\cdot 29$ appears in the numerator of $\Delta(c) + \Delta(26{-}c)$.

\sep

The bar complex of a chiral algebra is a $1$-form, a square
root, and a ghost.

The $1$-form $\eta = d\log(z_i - z_j)$ extracts the OPE.
Arnold's relation makes it nilpotent.  The genus-completed
differential $D_\cA$ satisfies $D_\cA^2 = 0$, and
$\Theta_\cA = D_\cA - d_0$ is the universal MC element.

The square root $\sqrt{Q_L}$ generates the tower.  One
quadratic form, one algebraic function, infinitely many
rational invariants.  The depth classification, the growth
rate, the convergence behaviour, the Galois theory, the Epstein
zeta function, the discriminant complementarity: all
are properties of this single square root.

The ghost kills the splitting
(Theorem~\ref{thm:gravitational-primitivity-wn}): the BRST
homotopy carries ghost number $\pm 1$, the projection
annihilates it, and the coproduct vanishes.

Everything else is a projection of these three.


% ============================================================
\section{Koszul--Epstein and quantum gravity}
\label{sec:koszul-epstein-gravity}
% ============================================================

\subsection{The arithmetic content of the gravitational bar complex}

The bar complex of the Virasoro algebra at central charge~$c$
carries three layers of arithmetic data.

\textbf{Layer 1: the Heisenberg core.}
The Miura map $\hat\Upsilon\colon W_\kappa(\mathfrak{sl}_2) \to
V^{\mathfrak{h}_\kappa}(\mathfrak{h})$ embeds the Virasoro
algebra into a rank-$1$ Heisenberg.  At the level of the
sewing determinant, this produces an exact Euler product:
\[
\sum_{N \ge 1} \sigma_{-1}(N)\,q^N
\;=\; -\log\bigl(\eta(q)^2\bigr) - \tfrac{\tau}{12},
\qquad
\sum_{N \ge 1} \sigma_{-1}(N)\,N^{-s}
\;=\; \zeta(s)\,\zeta(s+1).
\]
The Heisenberg core is arithmetically transparent: all
nontrivial genus-one arithmetic is a consequence of the
divisor function~$\sigma_{-1}$ being multiplicative.

\textbf{Layer 2: the Miura defect.}
The vacuum character of~$W_N$ splits as
$\chi^{\mathrm{vac}}_{W_N}(q)
= \chi_H(q)^{N-1} \cdot
\prod_{m=1}^{N-1}(1-q^m)^{N-m}$.
All nontrivial genus-one arithmetic beyond the Heisenberg core
sits in the finite Miura defect factor
$D_N(q) = \prod_{m=1}^{N-1}(1-q^m)^{N-m}$.
For~$N = 2$ (Virasoro): $D_2(q) = 1 - q$, removing the weight-$1$
mode.

\textbf{Layer 3: the shadow obstruction tower.}
The infinite shadow obstruction tower $\{S_r(c)\}_{r \ge 2}$, determined by
$(\kappa, \alpha, S_4)$ via the MC recursion, encodes the full
$A_\infty$ deformation data.  Its spectral measure~$\rho$ is
uniquely determined by the Carleman condition.  The MC recursion
propagates Hecke equivariance from the character level to all
arities (Route~C, proved for the standard landscape).

\subsection{The gravitational dichotomy and the Koszul--Epstein frontier}

The gravity dichotomy at $c = 26$ concerns the vanishing of
$\kappa_{\mathrm{eff}} = \kappa(\mathrm{matter})
+ \kappa(\mathrm{ghost})$.  The Koszul--Epstein question
concerns the critical-line properties of the arithmetic
shadow.  These are unrelated:

\begin{itemize}
\item $\kappa_{\mathrm{eff}} = 0$ at $c = 26$ determines whether
  the genus tower terminates (the Clifford-to-exterior
  degeneration of the transgression algebra).

\item The critical-line question concerns the zeros of
  $\zeta(s)$ appearing in the Heisenberg sewing determinant.
  These zeros are independent of the central charge: the sewing
  lift $S_{\mathrm{Vir}}(u) = \zeta(u+1)(\zeta(u) - 1)$ is the
  same function for all~$c$.
\end{itemize}

The Koszul self-duality point $c = 13$, where
$\kappa + \kappa' = 13$ and $\delta_\kappa = 0$, is the point
where Verdier--Hecke commutation forces the constrained Epstein
zeta to factor into standard $L$-functions.  This is a genuine
arithmetic structure of the self-dual gravitational theory.
It does not, however, produce a critical-line theorem.

\subsection{Honest assessment}

The MC element~$\Theta_\cA$ of the gravitational bar complex
produces genuine arithmetic structure: the Carleman-determined
spectral measure, Route~C prime-locality, the shadow--spectral
correspondence for lattice VOAs, the Verdier--Hecke factorisation
at self-dual points.  These are proved theorems.

The programme does not produce a critical-line theorem.  The
structural obstruction is that the shadow spectral measure
and the automorphic spectral measure are different objects, and
the map between them (the Rankin--Selberg integral) is not
sign-preserving.  No mechanism for positivity has been
identified.  The bar--cobar homotopy equivalence provides
deformation-theoretic transport, not arithmetic positivity.

The surviving content: the quartic resonance class
$R_4^{\mathrm{mod}}(\cA)$ is the first nonlinear place where
Koszul duality constrains the automorphic spectral decomposition
beyond the character level.  Whether this constraint carries
critical-line information is open.  The correct next computation
is the quartic Gram-side matching with the crossing-weighted
bilinear residue kernel (using actual DOZZ/Toda structure
constants), not the bare spectral measure (which fails by four
to five orders of magnitude).

The bar complex of a chiral algebra knows the exact genus tower,
the arithmetic of the shadow, and the depth classification.
It does not know the zeros of~$\zeta(s)$.


% ============================================================
\section{Spectrum and gravity}
\label{sec:spectrum-gravity}
% ============================================================

The gravitational bar complex carries three kinds of data:
the algebraic engine (Theorems A--H, the MC
element~$\Theta_\cA$), the arithmetic shadow (the shadow obstruction tower,
the depth classification, the Epstein zeta), and the spectral
content (the primary spectrum, the partition function, the
constrained Epstein series~$\varepsilon^c_s$).  These three
overlap but do not coincide.  This section records the definitive
assessment of their relationship in the gravitational context.

\subsection{The spectrum as a categorical invariant}

The scalar primary spectrum
$\mathrm{Spec}_{\mathrm{sc}}(\cA) = \{(\Delta_i, m_i)\}$ of a
2d~CFT at central charge~$c$ is an invariant of the
\emph{modular tensor category}
$\cA\text{-}\mathrm{mod}^{\mathrm{ss}}$ together with the
modular-invariant partition function $Z(\tau,\bar\tau)
= \sum_{i,j} M_{ij}\,\chi_i(\tau)\,\overline{\chi_j(\tau)}$.
The multiplicity matrix~$M_{ij}$ encodes the physical spectrum;
the characters $\chi_i$ encode the algebraic spectrum.

The bar complex $B(\cA)$ is built from the algebra~$\cA$
(specifically, from the vacuum module and its OPE), not from
the category of modules.  The bar-cobar machine determines the
Koszul dual~$\cA^!$, the MC element~$\Theta_\cA$, and the shadow
tower~$\{S_r\}$.  It does \emph{not} determine the modular
$S$-matrix or the multiplicity matrix.  The primary spectrum is
therefore external to the bar construction.

For the gravitational theory (Virasoro at central charge~$c$),
the vacuum module is the unique irreducible highest-weight module
$L(c,0)$, and the partition function $Z_{\mathrm{grav}}
= |\chi_0(\tau)|^2 + \sum_{h > 0} m_h\,|\chi_h(\tau)|^2$ depends
on the dynamical data $\{m_h\}_{h > 0}$: the spectrum of
primary states.  The bar complex of~$\mathrm{Vir}_c$ sees the
algebra, not the specific gravity theory built from it.

\subsection{Which constructions partially capture the spectrum}

Five constructions each capture a shadow of the gravitational
primary spectrum.  None captures it fully from within the bar
machine.

\begin{enumerate}
\item \emph{Zhu's algebra $A(\mathrm{Vir}_c)$.}
  For the universal Virasoro: $A(V_c) \cong \C[x]$.  The simple
  modules of $\C[x]$ correspond to points $x = h \in \C$: the
  conformal dimensions.  Zhu's algebra determines the
  \emph{set} of allowed lowest weights $\{h_{r,s}\}$ for the
  minimal model at $c = c_{p,q}$ (via the quotient $A(L(c,0))
  = \C[x]/(f_c(x))$).  But it does not determine the partition
  function multiplicities.  At generic~$c$ (irrational, or
  rational non-minimal), $A(V_c) = \C[x]$ admits a continuum
  of simple modules, and no finite selection rule from $A(V_c)$
  picks out the physical spectrum.

\item \emph{Associated variety $X_{\mathrm{Vir}_c}$.}
  For Virasoro: $X_{V_c} = \C$ (one-dimensional, since Virasoro
  has a single strong generator).  The dimension of $X_V$
  controls the asymptotic growth of the character and hence the
  leading pole of the constrained Epstein series at $s = c/2$.
  This is the Cardy regime: $\log \rho(E) \sim
  2\pi\sqrt{c\,E/3}$ for large conformal dimension~$E$.

\item \emph{Bar cohomology $H^*(B(\mathrm{Vir}_c))$.}
  Concentrated in bar degree~$1$ (Koszulness).  The graded
  dimensions $\dim H^1(B(\mathrm{Vir}))_h$ count generators of
  the Koszul dual $\mathrm{Vir}^!_{26-c}$ at weight~$h$.  These
  are the \emph{Motzkin differences} $M(h+1) - M(h)$, supported
  only at even weights, independent of~$c$, and growing
  as~$3^h$.  The quasi-primary count (partitions minus
  descendants) is a different sequence, nonzero at odd weights
  and growing as $e^{\pi\sqrt{2h/3}}$.  Bar cohomology sees the
  dual algebra, not the spectrum.

\item \emph{Shadow obstruction tower $\{S_r\}$.}
  The shadow obstruction tower extracts three numbers $(\kappa, \alpha, S_4)$
  from the OPE and packages them into the shadow metric~$Q_L$.
  The shadow-metric Epstein zeta $\varepsilon_{Q_L}(s)$ is a
  binary quadratic form sum, unrelated to the primary spectrum
  except through a chain of projections.  For lattice VOAs, the
  shadow--spectral correspondence recovers the full $L$-function
  package; for Virasoro, the correspondence is incomplete (the
  shadow has $3$~parameters; the spectrum has infinitely many).

\item \emph{The chiral derived centre
  $\cZ^{\mathrm{der}}_{\mathrm{ch}}(\cA)$.}
  The universal bulk algebra, built as chiral Hochschild cochains
  $C^\bullet_{\mathrm{ch}}(\cA_b, \cA_b)$, classifies bulk
  operators.  Its cohomology controls the BRST-closed observables
  of the 3d HT theory.  But it is a genus-$0$ invariant of the
  open/closed Swiss-cheese structure, and does not directly encode
  the genus-$1$ partition function.  The trace on the derived
  centre (the annulus trace~$\mathrm{Tr}_\cA$) recovers the
  Hochschild homology $\mathrm{HH}_*(\cA)$, which is the first
  modular shadow of the open sector.  This is related to, but
  not identical with, the partition function.
\end{enumerate}

\subsection{The gravitational spectral gap}

In the gravitational context, the primary spectrum has additional
structure: the conformal bootstrap imposes crossing-symmetry
constraints on the four-point function, which restrict the
allowed $\{(\Delta_i, m_i)\}$ beyond what the chiral algebra
alone determines.  The Benjamin--Chang constrained Epstein
series $\varepsilon^c_s$ is the Eisenstein spectral coefficient
of $\hat Z^c$, and its functional equation involves
$\zeta(2s)/\zeta(2s-1)$, introducing poles at
$s = (1+\rho_n)/2$ where $\rho_n$ are nontrivial zeros
of~$\zeta$.

The pole \emph{locations} are universal: they come from
$\zeta(2s-1) = 0$ in the $\Gamma$-factor ratio $F_c(s)$, and
do not depend on the specific CFT.  The pole \emph{residues}
are CFT-specific: they are integrals of the primary-counting
function against residues of Eisenstein series on the modular
fundamental domain.  The bar complex constrains neither the
locations (which are number-theoretic: the zeros of~$\zeta$)
nor the residues (which are spectral: the primary multiplicities
$m_h$).

For the gravitational theory at $c = 26$ (critical string):
$\kappa_{\mathrm{eff}} = \kappa(\mathrm{matter})
+ \kappa(\mathrm{ghost}) = 0$, and the genus tower terminates.
But the primary spectrum is still infinite (the Virasoro module
at $c = 26$ has infinitely many primaries at all weights
$h \geq 2$), and $\varepsilon^{26}_s$ is a nontrivial Dirichlet
series.  The genus tower's termination is a statement about the
$\hat A$-genus factor~$\kappa$; it says nothing about the
spectral content.

For the gravitational theory at $c = 13$ (self-dual point):
Verdier--Hecke commutation forces the constrained Epstein to
factor into standard $L$-functions.  This is the strongest
structural result: self-duality of the chiral algebra produces
a factorisation of the spectral invariant.  But even at
$c = 13$, the factorisation does not determine the zeros of the
constituent $L$-functions.

\subsection{Is $\varepsilon^c_s$ a bar-complex invariant?}

\begin{proposition}[The constrained Epstein is not a bar-complex invariant]
\label{prop:constrained-epstein-not-bar}
The constrained Epstein series $\varepsilon^c_s(\cA)$ cannot be
reconstructed from the bar complex $B(\cA)$ and its cohomology.
\end{proposition}

\begin{proof}
Three independent arguments.

\textbf{(a) Categorical argument.}
Two non-isomorphic modular-invariant partition functions for the
same chiral algebra produce different constrained Epstein series.
Example: the Ising model ($c = 1/2$) has the diagonal partition
function $Z_{\mathrm{diag}} = |\chi_0|^2 + |\chi_{1/2}|^2
+ |\chi_{1/16}|^2$ and the charge-conjugation partition function
(which is the same, since Ising is self-conjugate).  But for
$\widehat{\mathfrak{su}(2)}_k$ at $k \geq 10$, there exist
exceptional modular invariants (the $D$-type and $E$-type
classifications of Cappelli--Itzykson--Zuber) that produce
different primary spectra from the same chiral algebra.  The
bar complex $B(\hat\fg_k)$ is the same for all choices of
modular invariant; the constrained Epstein series differs.

\textbf{(b) Dimensional argument.}
The bar complex is determined by the OPE data of the vacuum
module (a countable set of structure constants).  The primary
spectrum is determined by the modular invariant \emph{and} the
full representation theory (the braided tensor structure, the
modular $S$-matrix).  The latter is strictly richer: two algebras
with the same OPE but different categories of modules (e.g.,
$V_k(\fg)$ at the same $k$ but with different conformal
extensions) have the same bar complex but different spectra.

\textbf{(c) Weight argument.}
$\dim H^1(B(\mathrm{Vir}))_h$ at odd $h$ vanishes (the Koszul
dual has generators only at even weights), while the quasi-primary
count at odd $h$ is generically nonzero.  No Dirichlet series
built from bar cohomology dimensions can reproduce
$\varepsilon^c_s$, which sums over all primary weights including
odd ones.
\end{proof}

\subsection{The honest three-layer picture}

The gravitational bar complex produces three layers of arithmetic
data, each genuine and proved.

\textbf{Layer 1: the genus expansion.}
$F_g(\cA) = \kappa(\cA) \cdot \lambda_g^{\mathrm{FP}}$
for uniform-weight algebras.  This is the leading scalar projection
of the MC element $\Theta_\cA$, generated by $\hat A(i\hbar) - 1$.
It captures \emph{one number} ($\kappa$) from the full partition
function.

\textbf{Layer 2: the shadow obstruction tower.}
$\{S_r\}_{r \geq 2}$, determined by $(\kappa, \alpha, S_4)$.
This captures \emph{three numbers} from the OPE, packaged into
the shadow metric~$Q_L$ and the quadratic extension~$K_L$.  The
shadow--spectral correspondence for lattice VOAs, Route~C
prime-locality, and the Verdier--Hecke factorisation are genuine
arithmetic results at this layer.

\textbf{Layer 3: the spectral content.}
$\varepsilon^c_s(\cA) = \sum (2\Delta)^{-s}$, encoding the full
scalar primary spectrum.  This layer is external to the bar
complex: it requires the partition function, which is a
genus-$1$ categorical invariant, not a genus-$0$ algebraic one.

The gap between Layers 2 and 3 is the gap between the
shadow obstruction tower (genus-$0$, $3$~parameters) and the primary spectrum
(genus-$1$, $\infty$~parameters).  No mechanism in the MC equation
or the bar-cobar machine bridges this gap: the shadow obstruction tower is
determined by finitely many OPE coefficients, while the spectrum
encodes infinitely many independent quantities.

The honest statement for the gravitational theory:
\emph{the bar complex of the Virasoro algebra at central
charge~$c$ determines the exact genus expansion
$F_g = (c/2)\,\lambda_g^{\mathrm{FP}}$, the infinite shadow
tower with growth rate~$\rho(c)$ and discriminant
$\Delta = 80/(5c{+}22)$, and the four-class depth assignment
(class~M for all $c \neq 0$).  It does not determine the
primary spectrum, the partition function, or the constrained
Epstein series.  The spectrum is a representation-theoretic
invariant; the bar complex is an algebra-theoretic invariant.
They live in different categories.}


% ============================================================
\section{Derived global sections in quantum gravity}
\label{sec:derived-global-sections-qg}
% ============================================================

The gravitational bar complex defines a sheaf on moduli space
whose derived global sections are algebraic invariants.  This
section records the definitive assessment of what these
algebraic invariants capture in the gravitational context,
correcting the false claim that ``the bar complex cannot see the
spectrum because integration over moduli space is analytic.''

\subsection{The sheaf-theoretic framework}

For Virasoro at central charge~$c$, the bar complex at
genus~$g$ defines a sheaf of dg modules on
$\overline{\cM}_g$.  The key objects:

\begin{enumerate}
\item The \emph{virtual bar $K$-class}
  $V_{\mathrm{Vir}_c} := [R\pi_{g*}\,\barB^{(g)}(\mathrm{Vir}_c)]
  \in K^0(\overline{\cM}_g)$, computed by derived pushforward
  along the universal curve.  This is well-defined: the
  perfectness criterion (Koszul acyclicity + finite-dimensional
  fiber cohomology) is satisfied for all modular Koszul algebras.

\item The \emph{center sheaf}
  $\cZ_{\mathrm{Vir}_c} := R^*\pi_{g*}(\barB^{(g)},\, d_{\mathrm{fib}})$,
  a local system on $\overline{\cM}_g$ with fiber
  $H^{\mathrm{ch}}_*(\Sigma_g, \mathrm{Vir}_c)$.

\item The \emph{ambient complex}
  $C_g(\mathrm{Vir}_c) = R\Gamma(\overline{\cM}_g,\,
  \cZ_{\mathrm{Vir}_c})$, the derived global sections.
  This is a purely algebraic operation on a constructible sheaf
  over a proper DM stack.

\item The \emph{GRR hierarchy}:
  \begin{align*}
  c_1(\det V_{\mathrm{Vir}_c}) &= \tfrac{c}{2} \cdot \lambda, \\
  \mathrm{ch}(V_{\mathrm{Vir}_c}) &\leftrightarrow \Delta_{\mathrm{Vir}}(x)
    \quad\text{(Motzkin differences)}, \\
  \mathrm{Hol}(V_{\mathrm{Vir}_c}) &= \Pi_{\mathrm{Vir}_c}
    \quad\text{(MCG monodromy)}.
  \end{align*}
  All three levels are algebraic.
\end{enumerate}

So the bar complex \emph{does} produce algebraic invariants
from integration over moduli space.  The claim that this
integration is ``analytic, not algebraic'' was a false idea.
Derived global sections, Chern character, GRR pushforward,
and monodromy representation are all algebraic operations.

\subsection{The gravitational spectral gap: precise diagnosis}

The GRR hierarchy captures the \emph{algebra's} modular
structure.  It does not capture the \emph{gravitational
spectrum}.  The precise diagnosis:

\textbf{(i) The GRR scalar level} recovers $\kappa(\mathrm{Vir}_c)
= c/2$.  This is one number.  The gravitational partition function
$Z(\tau,\bar\tau) = |\chi_0(\tau)|^2 + \sum_{h > 0} m_h\,
|\chi_h(\tau)|^2$ encodes infinitely many numbers $\{m_h\}$.

\textbf{(ii) The GRR spectral level} recovers the bar cohomology
dimensions $\dim H^1(B(\mathrm{Vir}))_h$, which are the Motzkin
differences $M(h+1) - M(h)$.  These count generators of the
Koszul dual $\mathrm{Vir}^!_{26-c}$, not primaries of
$\mathrm{Vir}_c$.  The two sequences differ: bar cohomology is
supported only at even weights and is $c$-independent; the
quasi-primary count is nonzero at odd weights and depends on the
specific gravity theory.

\textbf{(iii) The GRR holonomy level} recovers the monodromy of
$\cZ_{\mathrm{Vir}_c}$: the action of the mapping class group on
chiral homology.  For the Virasoro algebra, this is the
monodromy of the KZ connection, which encodes the
projective representation of $\mathrm{MCG}(\Sigma_g)$ on the
space of conformal blocks.  This is an invariant of the
\emph{algebra} (it is functorial in morphisms of chiral algebras,
not in changes of the module category or modular invariant).

\textbf{The obstruction is not analysis vs.\ algebra.}
It is \emph{algebra vs.\ category}: the bar complex sees
the chiral algebra $\mathrm{Vir}_c$; the gravitational
spectrum sees the choice of gravity theory built from
$\mathrm{Vir}_c$ (the specific multiplicities $\{m_h\}$
in the partition function, which are dynamical data external
to the algebra).

\subsection{Three routes toward the spectrum}

\paragraph{(a) Module bar construction.}
The module bar complex $\barB_{\mathrm{Vir}_c}(L(c,h))$ applied
to each irreducible module $L(c,h)$ produces a comodule over
$\barB(\mathrm{Vir}_c)$.  Its $K$-class carries the conformal
dimension~$h$.  Summing with the gravitational multiplicities
$\{m_h\}$:
\[
\sum_{h} m_h\,\mathrm{ch}\bigl(
R\pi_{1*}\,\barB_{\mathrm{Vir}_c}(L(c,h))\bigr)
\]
would recover the gravitational partition function at genus~$1$.
But this takes $\{m_h\}$ as \emph{input}: it computes the
partition function from the spectrum, rather than extracting
the spectrum from the bar complex.

\paragraph{(b) Conformal blocks as $D$-module.}
The conformal blocks local system on $\cM_{g,n}$ carries a flat
connection (the KZ/Hitchin connection) whose monodromy encodes
the modular $S$-matrix for rational theories.  At integrable
level for affine Kac--Moody algebras, the monodromy recovers the
primary dimensions via the Verlinde formula.

For the Virasoro algebra at generic~$c$, the conformal blocks
local system has irregular singularities and the representation
theory is non-semisimple (at irrational~$c$, all Verma modules
are irreducible and the spectrum is continuous).  Whether the
$D$-module structure of the conformal blocks sheaf on
$\overline{\cM}_{1,1}$ determines the gravitational spectrum
(after choosing a modular invariant) is a frontier question:
it amounts to recovering the spectral decomposition of a
function on $\mathrm{SL}(2,\Z)\backslash\mathfrak{H}$ from
the algebraic structure of its associated $D$-module.

\paragraph{(c) The Roelcke--Selberg step.}
The passage from the partition function
$Z(\tau,\bar\tau)$ to the constrained Epstein series
$\varepsilon^c_s$ uses the spectral decomposition of
$L^2(\mathrm{SL}(2,\Z)\backslash\mathfrak{H})$ into Eisenstein
series and Maass cusp forms.  This involves the hyperbolic
Laplacian, which is a real-analytic operator.
The scattering matrix involves $\zeta(2s)/\zeta(2s-1)$,
whose poles are the nontrivial zeros of the Riemann
zeta function.  No categorical or motivic construction is
known to produce this spectral decomposition.

The ``analytic'' step that is genuinely unavoidable is this
last one: from characters (which are algebraic for rational VOAs)
to spectral coefficients (which require harmonic analysis on
$\mathrm{SL}(2,\Z)\backslash\mathfrak{H}$).

\subsection{Honest status}

\begin{enumerate}[label=(\alph*)]
\item The bar complex defines a sheaf on $\overline{\cM}_g$
  with algebraic invariants ($R\Gamma$, $\mathrm{ch}$,
  $\mathrm{Hol}$, GRR).  The claim ``integration is analytic''
  was a false idea.

\item These algebraic invariants capture the algebra's modular
  structure: $\kappa$, bar cohomology dimensions, KZ/Hitchin
  monodromy.  They do \emph{not} capture the gravitational
  spectrum $\{m_h\}$ or the constrained Epstein
  $\varepsilon^c_s$.

\item The obstruction is \emph{algebra vs.\ category},
  not analysis vs.\ algebra.  Two gravity theories with the
  same Virasoro algebra but different spectra have the same
  bar complex.  This is a theorem, not a technical limitation.

\item The construction that does capture the spectrum---the
  conformal blocks $D$-module on $\cM_{g,n}$---is algebraic
  (Beilinson--Drinfeld) but requires the module category
  as input.  For rational theories, its monodromy ($S$-matrix)
  determines the primary dimensions.  For non-rational theories,
  the $D$-module has irregular singularities and the spectral
  question is at the frontier of geometric Langlands and
  irregular Hodge theory.

\item The single genuinely analytic step is the
  Roelcke--Selberg decomposition: from characters to spectral
  zeta.  This step is universal (it involves the Riemann
  zeta function, not the specific VOA) and has no known
  categorical substitute.

\item Status: the ``spectral question'' is \textbf{OPEN} in the
  sense that the algebraic $D$-module framework exists but
  has not been carried through for Virasoro at generic~$c$.
  The gap is not a missing construction but a missing
  computation: does the irregular $D$-module structure of the
  Virasoro conformal blocks local system on
  $\overline{\cM}_{1,1}$ determine the primary spectrum after
  choosing a modular invariant?
\end{enumerate}


% ============================================================
\section{Complex parameters in quantum gravity}
\label{sec:complex-parameters-qg}
% ============================================================

The gravitational bar complex of the Virasoro algebra produces
real spectral data: the modular characteristic $\kappa = c/2$,
the shadow obstruction tower $\{S_r(c)\}$, and the genus expansion
$F_g = \kappa\,\lambda_g^{\mathrm{FP}}$ are all rational
functions of the central charge~$c$.  The zeros of the
Riemann zeta function, which appear in the scattering matrix
of the Eisenstein spectral decomposition on $\cM_{1,1}$, live
at complex spectral parameters.  This section records the
definitive assessment of whether any extension of the bar
complex can access complex spectral data.

\subsection{The structural obstruction}

The obstruction is a theorem with three components.

\begin{enumerate}
\item The shadow obstruction tower on a single primary line is determined
by three genus-$0$ OPE invariants $(\kappa, \alpha, S_4)$
via Theorem~\ref{thm:riccati-algebraicity}: the generating
function $H(t) = t^2\sqrt{Q_L(t)}$ is algebraic of degree~$2$.
Three numbers constitute a finite-dimensional algebraic datum.

\item The zeta zeros $\{\rho_n = \tfrac{1}{2} + i\gamma_n\}$
constitute a countably infinite collection of complex numbers,
believed to be algebraically independent.  No finite-dimensional
algebraic structure determines such a set.

\item The Rankin--Selberg unfolding that connects the partition
function to the constrained Epstein $\varepsilon^c_s$ erases
the convolution-algebraic structure of the MC equation.  The
unfolded integral depends on Fourier coefficients alone, not on
the MC brackets.  The $L$-functions $L(s, f_j)$ in the
spectral decomposition are properties of the eigenforms~$f_j$,
independent of the algebra~$\cA$.
\end{enumerate}

\subsection{Eight candidate bridges: the census}

A complete investigation examined eight approaches to bridging
the gap from real bar-complex data to complex $L$-function
zeros.

\paragraph{Dead ends (five).}
Deligne categories, Wick rotation in~$c$, resonance/MC4
completion, Borel resummation, and spectral curves.
Each fails for a precise reason: Deligne categories
interpolate in a categorical variable, not a spectral one;
Wick rotation changes the shadow coefficients but not the
sewing lift $S_\cA(u)$, which is $c$-independent; MC4
completion solves a homological algebra problem, not a
number-theoretic one; Borel resummation requires factorial
divergence that the shadow series does not exhibit; and
spectral curves encode genus-$0$ OPE data, not the automorphic
spectrum.

\paragraph{Genuine mathematics, no bridge (two).}
The analytic continuation of the shadow obstruction tower and the analytic
Langlands programme produce real mathematical content but do
not bridge the real-to-complex gap.  The shadow obstruction tower
continuation produces quadratic number fields (the Galois
structure of the shadow extension $K_L/F$, with discriminant
$\mathrm{disc}(Q_L) = -320\,c^2/(5c{+}22)$ for Virasoro).
The analytic Langlands provides the correct framework for the
critical-level oper/shadow interpolation.  Both operate on the
function-field (geometric) side; the number-field (arithmetic)
descent remains the deepest unsolved problem.

\paragraph{The most compatible approach (one).}
The Connes--Marcolli non-commutative geometry programme has the
closest structural compatibility with the bar construction: the
Bost--Connes system has a Hecke algebra of symmetries
compatible with Verdier--Hecke commutation, and the Weil
positivity criterion is a trace-class condition on a specific
operator.  At the self-dual point $c = 13$, Verdier
self-duality provides an involution on the factorisation
coalgebra that could in principle feed into the non-commutative
trace formula.  In practice, the arithmetic descent (the bar
complex is defined over~$\mathbf{C}$, the Bost--Connes system
over~$\mathbf{Q}$) blocks the construction.

\subsection{What the gravitational bar complex provides}

The bar complex of the Virasoro algebra at central charge~$c$
provides five kinds of arithmetic data, all proved.

\begin{enumerate}
\item The exact genus expansion
$F_g = (c/2)\,\lambda_g^{\mathrm{FP}}$ at all genera.
\item The infinite shadow obstruction tower with growth rate
$\rho(c) = \sqrt{(36 + 2\Delta)/(c^2)}$ and discriminant
$\Delta = 80/(5c{+}22)$.
\item The four-class depth assignment (class~M for all
$c \neq 0$).
\item The quadratic number field
$K_L = \mathbf{Q}(\sqrt{\mathrm{disc}(Q_L)})$ at each
rational~$c$, with class numbers ranging from $h = 1$
(models $(9,10)$ and $(24,25)$) to $h = 3008$ (model $(28,29)$).
\item Route~C prime-locality: the MC recursion propagates
Hecke equivariance from low-arity character data to all
arities.
\end{enumerate}

None of these constrains the zeros of the Riemann zeta function
or any standard $L$-function.  The constraints are algebraic
(on OPE coefficients and their MC recursions); the zero
locations are analytic (determined by the Euler product and the
functional equation, not by the OPE).

\subsection{The gravitational dichotomy and the spectral divide}

The gravity dichotomy (critical string $c = 26$ versus
non-critical $c \neq 26$) is a statement about the
\emph{vanishing} of a specific class:
$\kappa_{\mathrm{eff}} = \kappa(\mathrm{matter})
+ \kappa(\mathrm{ghost}) = 0$ if and only if $c = 26$.
The self-dual point $c = 13$ is where the Koszul asymmetry
$\delta_\kappa = \kappa - \kappa' = 0$ vanishes.  These are
different conditions at different central charges (AP29).

The gravity dichotomy operates on the \emph{algebraic} axis:
whether the Clifford algebra $\mathrm{Cliff}(V, \lambda)$
with $\lambda = \kappa(\cB)\,\omega_g$ degenerates to an
exterior algebra ($\lambda = 0$ at $c = 26$) or remains a
matrix algebra ($\lambda \neq 0$ at $c \neq 26$).  This is a
statement about the transgression algebra, not about the
scattering matrix.  The spectral divide (real versus complex
spectral parameters) operates on the \emph{analytic} axis.
The two axes are orthogonal, and neither the gravitational
dichotomy nor the Koszul self-duality provides a bridge between
them.

\subsection{The Beilinson verdict for quantum gravity}

The following claims about the gravitational theory are
dismissed.

\begin{enumerate}
\item ``3d quantum gravity from Koszul duality constrains zeta
zeros.''  \emph{False.}  The Koszul duality of the Virasoro
algebra determines the complementary algebra $\mathrm{Vir}_{26-c}$
and the complementarity identity $\kappa + \kappa' = 13$.
These are algebraic constraints on OPE invariants, not analytic
constraints on $L$-function zeros.

\item ``The gravitational genus expansion encodes arithmetic
data beyond the character.''  \emph{Partially true, partially
false.}  The genus expansion $F_g = \kappa\,
\lambda_g^{\mathrm{FP}}$ is tautological (it lives in the
tautological ring of~$\overline{\cM}_g$).  The shadow obstruction tower
provides additional algebraic data (the quartic resonance
class, the discriminant, the depth classification).  But none
of this data accesses the primary spectrum or the partition
function, which are the objects that carry arithmetic content
through the Rankin--Selberg integral.

\item ``The gravitational partition function at $c = 26$ has
a distinguished relationship to the Riemann zeta function.''
\emph{False.}  The partition function of the
$c = 26$ Virasoro is $\eta(\tau)^{-24} = \Delta(\tau)^{-1}$,
the inverse of the Ramanujan discriminant form.  The
Rankin--Selberg $L$-function of $\Delta$ is
$L(s, \Delta \times \Delta) = \zeta(s{-}11)\,
L(s, \mathrm{Sym}^2\, \Delta)$, which has critical line
$\mathrm{Re}(s) = 23/2$, not $\mathrm{Re}(s) = 1/2$.
The zeta function $\zeta(s{-}11)$ appearing here is shifted
by~$11$ (the weight of~$\Delta$ minus~$1$).  The relationship
to the Riemann zeta is classical (Hecke theory), not a
consequence of the bar construction.
\end{enumerate}

\subsection{What survives for the gravitational theory}

The bar complex of the Virasoro algebra at central charge~$c$
is the algebraic engine of $3$d quantum gravity on
$\mathbf{C}_z \times \mathbf{R}_t$.  It determines:
the Swiss-cheese algebra structure (bar differential = chiral
product in~$\mathbf{C}$, bar coproduct = topological cutting
in~$\mathbf{R}$), the curved genus tower ($d^2 = \kappa\,
\omega_g$ at genus $g \ge 1$), the infinite shadow Postnikov
tower (all arities, all genera, proved by
Theorem~\ref{thm:mc2-bar-intrinsic}), the complementarity
decomposition (Theorem~C), and the completed modular Koszul datum
$\Pi^{\mathrm{oc}}_X(\cA)$ packaging open/closed/modular data.

These are genuine results about the algebraic structure of the
gravitational theory.  They do not imply, and should not be
claimed to imply, any result about the analytic number theory
of $L$-functions.  The bar complex is the algebraic engine.
The primary spectrum is the analytic output.  The gap between
them is the gap between algebra and analysis, which no known
construction bridges.


% =================================================================
\section{The chiral Tamarkin problem}
\label{sec:chiral-tamarkin}
% =================================================================

A natural question about the bar construction is whether the
modular direction (the genus tower controlled by the
Getzler--Kapranov Feynman transform) is a genuine ``second
direction'' in the sense of Tamarkin's dimension-raising theorem.

The answer is: \emph{no and yes}.  The modular direction is not a
second holomorphic direction (it is operadic, not
factorisation-algebraic), but it is a genuine codimension-$1$
enhancement of the chiral
algebra, categorically independent of the Swiss-cheese centre,
and governed by its own universal construction.

\subsection{Two codimension-$1$ operations}

Tamarkin's theorem (2000): for an $\mathsf{E}_d$-algebra~$A$,
the Hochschild cochain complex $C^*(A,A)$ is an
$\mathsf{E}_{d+1}$-algebra.  The mechanism is convolution against
the little-discs cooperad.  The extra direction is topological.

For chiral algebras, there are two independent analogues:

\smallskip

\noindent
\textbf{(i) The Swiss-cheese centre} (the chiral Deligne--Tamarkin
theorem, proved in Volume~I).  The chiral Hochschild cochain complex
$C^\bullet_{\mathrm{ch}}(\cA,\cA)$ is the terminal local
$\SCchtop$-algebra.  This adds one topological direction~$\R$,
producing a $3$d HT theory on $X \times \R$.

\smallskip

\noindent
\textbf{(ii) The modular convolution} (Getzler--Kapranov, proved in
Volume~I).  The modular convolution dg Lie algebra $\gAmod$ extends
genus-$0$ data to all genera.  The mechanism is convolution against
the modular cooperad of stable curves.  The BV operator
$\hbar\Delta$ (non-separating clutching
$\overline\cM_{g,n+2} \to \overline\cM_{g+1,n}$) raises genus
by~$1$.

\smallskip

\noindent
These are different functors: (i) uses the Swiss-cheese coloured
operad (configuration spaces of discs in $\C \times \R$);
(ii) uses the modular operad (stable curves of all genera with
clutching maps).  Both are convolutions against cooperads, and
both produce Maurer--Cartan-controlling Lie-type algebras.

The open/closed MC element $\Theta^{\mathrm{oc}}_\cA = \Theta_\cA
+ \sum_j \mu^{\cM_j}$ packages both into a single algebraic
object.  Its closed projection is $\Theta_\cA$ (the modular
convolution MC element).  Its open projection gives the $\Ainf$
boundary module structure.  Its mixed projection gives the
bulk-boundary coupling (the chiral derived centre acting on the
boundary via the Swiss-cheese structure).  Its clutching sector
gives the genus tower.


\subsection{The modular open/closed cooperad conjecture}

The hierarchy of Tamarkin-type theorems is:

\begin{center}
\renewcommand{\arraystretch}{1.4}
\begin{tabular}{@{}p{3.3cm}p{3.3cm}p{5.5cm}@{}}
\textsc{Tamarkin:}
  & $A$ is $\mathsf{E}_d$
  & $\Longrightarrow\; C^*(A,A)$ is $\mathsf{E}_{d+1}$ \\[2pt]
\textsc{Chiral DT:}
  & $\cA$ is chiral
  & $\Longrightarrow\; (C^\bullet_{\mathrm{ch}}(\cA,\cA), \cA)$
    is $\SCchtop$-algebra \\[2pt]
\textsc{CT-$2$:}
  & $\cA$ chirally Koszul
  & $\Longrightarrow\;$ Swiss-cheese extends to
    modular Swiss-cheese
\end{tabular}
\end{center}

\begin{conjecture}[CT-$2$: the modular open/closed cooperad]
\label{conj:ct2-vol2}
The open-sector factorisation dg-category $\cC_{\mathrm{op}}$
admits a modular cooperad structure whose cocomposition maps are
compatible with clutching on bordered Fulton--MacPherson
compactifications, whose trace recovers the annulus identification
$\HH_*(\cA_b) \simeq \Tr_\cA$ at genus~$1$, and whose full
structure determines and is determined by the open/closed MC
element~$\Theta^{\mathrm{oc}}_\cA$.
\end{conjecture}

\noindent\textbf{Status.}\quad
CT-$2$ is a research programme with at least six independent
sub-problems: (P1)~dualizability of~$\cC_{\mathrm{op}}$
in~$\mathrm{dgCat}$; (P2)~annular bordered FM compactification
(the cyclic case is not yet constructed---the genus-$1$ annulus
trace is proved by Ayala--Francis excision, not by bordered FM);
(P3)~construction of functorial cocomposition maps along bordered
FM clutching for all stable topological types;
(P4)~codimension-$2$ coassociativity (the MC equation provides
codimension-$1$ identities; coassociativity is codimension~$2$);
(P5)~$(\infty,2)$-categorical lift from chain-level cooperad to
cooperad in dg-categories; (P6)~coderived compatibility at genus
$g \geq 1$ (curved bar forces passage to coderived categories).
The MC equation for $\Theta^{\mathrm{oc}}$ is proved at all
genera.  The genus-$0$ part is the chiral Deligne--Tamarkin
theorem (proved).  This is Stage~$4$ of the four-stage
architecture.

\begin{remark}[Why this is the correct formulation]
The conjecture is not ``the shadow obstruction tower is a chiral algebra on
moduli'' (that is ill-defined: $\overline\cM_g$ has dimension
$3g{-}3$, not~$1$).  The conjecture is that the open-sector
category carries a modular cooperad structure compatible with
the proved MC equation.  The modular cooperad is the algebraic
structure that governs the genus direction; it is not a
factorisation algebra on a Ran space, but an operadic structure
indexed by boundary strata of moduli spaces.  The distinction
between factorisation (points colliding on a fixed curve) and
composition (curves degenerating along boundary divisors) is
categorical.
\end{remark}


\subsection{The BCOV parallel in the HT setting}
\label{subsec:bcov-ht}

The genus-by-genus MC hierarchy from Volume~I,
\[
d_{\Theta^{(0)}}(\Theta^{(g)})
= -\tfrac{1}{2}\!\sum_{g_1+g_2=g}
  [\Theta^{(g_1)}, \Theta^{(g_2)}]
- \Delta(\Theta^{(g-1)})
- d_{\mathrm{sew}}^{(g)} - d_{\mathrm{pf}}^{(g)},
\]
has the structure of the BCOV holomorphic anomaly equation:
both are genus recursions governed by the modular operad, with a
bilinear term from separating degenerations and a genus-raising
term from non-separating degenerations.

In the $3$d HT setting of this volume, the parallel is sharper.
The bulk algebra $\Zder(\cA) = C^\bullet_{\mathrm{ch}}(\cA_b,
\cA_b)$ is the universal closed-sector algebra
(Theorem~\ref{thm:thqg-swiss-cheese}).  At genus~$0$, the bulk is
the chiral derived centre.  At genus $g \geq 1$, the curved
Swiss-cheese structure $d^2 = \kappa \cdot \omega_g$ infects the
bulk via the Gauss--Manin connection.  The total differential
$D_g = d_{\mathrm{fib}} + \nabla^{\mathrm{GM}}_g$ restores
flatness by absorbing the curvature into the connection.

\begin{warning}[The BCOV analogy is combinatorial, not structural]
\label{warn:bcov-not-structural-vol2}
The MC hierarchy and the BCOV holomorphic anomaly equation share a
combinatorial backbone (genus recursion from the modular operad),
but they are not structurally isomorphic.  The shadow obstruction tower has no
holomorphic anomaly ($\bar\partial = 0$ on algebraic classes), no
non-canonical propagator (the bar propagator $d\log E$ is
canonical), and two extra terms ($d_{\mathrm{sew}}$,
$d_{\mathrm{pf}}$) with no BCOV counterpart.  The correct physics
analogue is the DVV/Virasoro recursion for CohFTs, via Teleman
reconstruction and the Eynard--Orantin identification.
\end{warning}

\noindent
What does survive: the Gauss--Manin variation on $1$D deformation
families.  For $\{\cA_s\}_{s \in B}$, the shadow generating
function $H(t;s) = t^2\sqrt{Q_L(t;s)}$ satisfies a first-order
ODE in~$s$ controlled by $\kappa'(s)$, $\alpha'(s)$, $S_4'(s)$.
The flat modular connection $D_\Theta = d + [\Theta_\cA, -]$ on
bar cohomology over $\cM_g$ gives the all-genera extension.

The decisive asymmetry: the shadow obstruction tower has \emph{no holomorphic
anomaly} ($\Theta^{(g)}$ is algebraic, so $\bar\partial = 0$
automatically), and the shadow partition function \emph{converges}
(Bernoulli decay $1/(2\pi)^{2g}$, versus $(2g)!$ divergence for
BCOV).  The algebraic bar construction is the universal structure;
BCOV is one analytic incarnation.


\subsection{Chromatic structure of the bar complex}
\label{subsec:chromatic-bar}

The scalar shadow obstruction tower is at chromatic height~$0$ (additive
formal group, $\hat{A}$-genus).  But the bar complex on curves
of different type uses propagators of different height:
$1/(z{-}w)$ on $\bP^1$ (rational, height~$\infty$),
$\pi\cot(\pi z)$ on $\C^*$ (trigonometric, height~$1$),
$\vartheta_1'/\vartheta_1$ on $E_\tau$ (elliptic, height~$2$).

The rational/trigonometric/elliptic hierarchy of quantum groups
IS the chromatic tower.  The Fay identity on $E_\tau$ is the
genus-$1$ MC equation with elliptic propagator.

\begin{conjecture}[Chromatic bar--cobar duality]
\label{conj:chromatic-vol2}
The bar--cobar adjunction on a family of elliptic curves
$E \to S$ defines a sheaf of $\Etwo$-algebras on~$S$ whose
associated formal group is the universal elliptic formal group.
The resulting Koszul duality theory determines a
$\mathrm{TMF}$-module structure on factorisation coalgebras
on $E/S$.
\end{conjecture}

\noindent
The genus spectral sequence---with $E_1$ page isolating tree
($p{=}0$), one-loop ($p{=}1$), genus-$2$ shell ($p{=}2$)
data---admits a conjectural identification with the chromatic
filtration.  Whether the shadow depth classification G/L/C/M
corresponds to chromatic complexity is a speculative question at
the interface of modular Koszul duality and stable homotopy theory.

\sep

The three conjectures of this section
(CT-$2$, the BCOV reformulation, and the chromatic lift)
are not independent.  The modular cooperad conjecture CT-$2$
would provide the categorical infrastructure in which both the
BCOV equations and the chromatic structure would live.
The BCOV reformulation requires understanding how the
shadow connection $\nabla^{\mathrm{sh}}$ extends from
$1$-dimensional deformation slices to the full moduli of chiral
algebras---and this extension is controlled by the cooperad
cocomposition maps of CT-$2$.
The chromatic conjecture requires showing that the elliptic bar
complex varies correctly over the modular curve---and the
relevant variation is a crystal of factorisation coalgebras,
whose coherence is governed by the same modular cooperad.

The unifying perspective: the MC element $\Theta_\cA$ in
$\gAmod$ simultaneously encodes the spatial dimension-raising
(its genus-$0$ binary shadow $r(z)$ is the R-matrix, providing
the $\mathsf{E}_1$/Yangian structure in the topological
direction) and the genus dimension-raising (its higher-genus
components are the quantum corrections, living on
$\overline\cM_g$).  Quantisation and genus-raising are two
faces of the same MC element, expressed in different projections
of the shadow obstruction tower.


% =================================================================
\section{The bar construction as string field theory}
\label{sec:bar-sft}
% =================================================================

The bar construction is a string field theory functor.  The
$A_\infty$ operations from the chiral Hochschild complex are
Zwiebach's closed string vertices ($m_1 = Q_{\mathrm{BRST}}$,
$m_2 =$ string product, $m_3 =$ four-string vertex, higher
$m_n =$ contact terms), and the action is
$S[\Psi] = \frac{1}{2}\langle\Psi, Q\Psi\rangle
+ \sum_{n \geq 3} \frac{1}{n!}
  \langle\Psi, m_n(\Psi,\ldots,\Psi)\rangle$.
This is Zwiebach's theorem (1993), restated as
Theorem~\ref{thm:string-field-theory-hochschild} in Volume~I.

The genus-$0$ BV/bar identification is proved.  The full
chain-level identification at genus $g \geq 1$ is heuristic
(the MC equation is proved; the physical identification
BV $=$ bar is Conjecture~\ref{conj:master-bv-brst}).

\subsection{Four objects and their string-theoretic roles}

The open/closed string field theory of a chiral algebra $\cA$
on $\C \times \R$ involves four objects:
\begin{itemize}[nosep]
\item $\cA$ (boundary algebra) $=$ open-string sector.
\item $\barB(\cA)$ (bar coalgebra) $=$ classifies twisting
  morphisms (universal couplings between $\cA$ and $\cA^!$).
  This is the open-string configuration space, not the
  closed-string state space.
\item $\Zder(\cA) = C^\bullet_{\mathrm{ch}}(\cA,\cA)$
  (derived centre) $=$ universal bulk algebra (closed-string
  observables).
\item $\Theta^{\mathrm{oc}}_\cA$ (open/closed MC element) $=$
  the full open/closed string field, packaging all four sectors.
\end{itemize}
The bar complex and the derived centre are \emph{different}
objects playing \emph{different} roles:
bar $=$ couplings; centre $=$ observables.
Open/closed string duality is the interplay between them, not
the identification of one with the other.

\subsection{Kodaira--Spencer gravity and open/closed duality}

Volume~II's worked example (Section~\ref{subsec:KS-defect})
identifies the large-$N$ BRST cohomology of the gauged
$\beta\gamma$ system with the Kodaira--Spencer chiral algebra
of BCOV.  The passage from open-string to closed-string
degrees of freedom is controlled by \emph{Koszul duality}
($\cA_{\mathrm{open}} \mapsto \cA^!_{\mathrm{open}}$ via
Verdier duality on bar cohomology) and the \emph{derived
centre} $C^\bullet_{\mathrm{ch}}(\cA_{\mathrm{open}},
\cA_{\mathrm{open}})$ as the universal bulk.  Bar-cobar
inversion $\Omega B(\cA_{\mathrm{open}}) \simeq
\cA_{\mathrm{open}}$ recovers the original boundary algebra
(Theorem~B), not the KS bulk~(AP34).

The structural parallel between BCOV and the modular MC
hierarchy arises because both are string field theories governed
by the modular operad.  The BCOV system is the analytic
incarnation (metric-dependent propagator, holomorphic anomaly,
divergent).  The shadow obstruction tower is the algebraic incarnation
(canonical propagator, no anomaly, convergent).

\begin{question}[Open/closed duality: bar vs derived centre]
\label{q:open-closed-vol2}
The bar complex classifies twisting morphisms.  The derived
centre classifies bulk observables.  The Swiss-cheese structure
gives the local answer: $(\Zder(\cA), \cA)$ is the terminal
$\SCchtop$-pair.  The modular extension gives
$\Theta^{\mathrm{oc}}_\cA$.  But what is the precise categorical
relationship between the category of twisting morphisms
(controlled by $\barB(\cA)$) and the category of bulk operators
(controlled by $\Zder(\cA)$)?

In the topological setting, the Hochschild cochains of an
$\Eone$-algebra are computed by the bar complex: $C^*(A,A)
\simeq \RHom_{A \otimes A^{\mathrm{op}}}(A,A) \simeq
\RHom(\barB(A), A)$.  In the chiral setting, the chiral
Hochschild cochains are $C^\bullet_{\mathrm{ch}}(\cA,\cA)$,
which is the operadic centre, \emph{not} $\RHom(\barB(\cA),
\cA)$.  The gap between the operadic centre (the derived
centre) and the convolution $\Hom$ from the bar coalgebra
(the twisting morphism space) is where open/closed duality
lives.  Closing this gap is the mathematical content of
open/closed string duality in the chiral setting.
\end{question}


% =================================================================
\section{Shadow obstruction tower depths for $\cW_3$ and $\cW_4$: testing depth$(\cW_N) = N$}
\label{sec:shadow-tower-depths-WN}
% =================================================================

The existing shadow depth classification (Vol~I,
\S\ref{sec:one-quadratic-form}; Vol~II,
Theorem~\ref{thm:wn-class-M}) establishes $r_{\max}(\cW_N)
= \infty$ for all principal $\cW_N$ ($N \ge 2$).  But infinite
depth on the $T$-line is inherited from the Virasoro subalgebra
and does not detect the higher-spin structure.  A finer invariant
is needed: the \emph{$W$-essential depth}, which measures the
minimal arity at which the mixed ($T$-$W$) shadow components
carry genuinely new information not determined by lower arities.

\subsection*{The two notions of depth}

\begin{definition}[Perturbative depth and $W$-essential depth]
\label{def:two-depths}
Let $\cA$ be a $\cW_N$ algebra with generators
$W_2 = T, W_3, \ldots, W_N$.  Write $x_s$ ($s = 2, \ldots, N$)
for the formal deformation variables on the $(N{-}1)$-dimensional
primary space.  The arity-$r$ shadow $\mathrm{Sh}_r$ is a
homogeneous polynomial of degree $r$ in $x_2, x_3, \ldots, x_N$.
\begin{enumerate}[label=\textup{(\roman*)}]
\item The \emph{perturbative depth} $d(\cA)$ is $\sup\{r :
  m_r \ne 0 \text{ on cohomology}\}$.  For all $\cW_N$,
  $d(\cW_N) = \infty$ (Theorem~\ref{thm:wn-class-M}).
\item The \emph{$W$-essential depth} $d_W(\cA)$ is the smallest
  $r$ such that for all $r' > r$, every mixed monomial
  $x_2^{a_2} x_3^{a_3} \cdots x_N^{a_N}$ of degree $r'$ with
  $\max(s : a_s > 0) \ge 3$ has its coefficient determined by the
  shadow data at arities $\le r$ via the master equation.
\end{enumerate}
\end{definition}

The perturbative depth is always $\infty$ for class~M.  The
$W$-essential depth measures how deep one must go before the
higher-spin content is exhausted.  The conjecture is:

\begin{conjecture}[Shadow $W$-essential depth]
\label{conj:w-essential-depth}
For the principal $\cW_N$ algebra:
\[
d_W(\cW_N) = N.
\]
\end{conjecture}

We now test this by explicit computation.

\subsection*{Review: the Virasoro shadow metric}

For Virasoro ($N = 2$, single generator $T$), the shadow metric
is
\[
Q_{\mathrm{Vir}}(t)
= (c + 6t)^2 + \frac{80t^2}{5c+22}
= c^2 + 12ct + \left(36 + \frac{80}{5c+22}\right)t^2.
\]
The generating function $H(t) = t^2\sqrt{Q_{\mathrm{Vir}}(t)}$
gives $\mathrm{Sh}_r = a_{r-2}/r$ where $a_n$ satisfies
$a_0 = c$, $a_1 = 6$,
$a_2 = 40/[c(5c{+}22)]$, and the convolution recursion.
The shadow data at each arity consists of a single number
$\mathrm{Sh}_r(c)$; these are rational functions of $c$ with
denominator $c^{r-3}(5c+22)^{\lfloor(r-2)/2\rfloor}$.

Since $Q_{\mathrm{Vir}}$ has $\Delta = 8\kappa S_4
= 8(c/2)(10/[c(5c{+}22)]) = 40/(5c{+}22) \ne 0$ for $c$
generic, we are in class~M with $r_{\max} = \infty$.
There is no $W$-direction, so $d_W(\mathrm{Vir}) = 2$
trivially.

\subsection*{$\cW_3$ shadow obstruction tower: explicit computation}

For $\cW_3$ with generators $T$ (weight $2$) and $W$
(weight $3$), the shadow lives on a $2$-dimensional deformation
space $(x_T, x_W)$.  The master equation for the Maurer--Cartan
element $\Theta = \sum_{r \ge 2} \mathrm{Sh}_r$ in the cyclic
deformation complex is:
\[
D_0 \Theta + \frac{1}{2}[\Theta, \Theta]_H = 0,
\]
where $[\cdot, \cdot]_H$ denotes the $H$-Poisson bracket with
propagator $P = \mathrm{diag}(2/c,\, 3/c)$ (the inverse of the
leading $\lambda$-bracket: $\langle T_\lambda T \rangle_0 = c/2$,
$\langle W_\lambda W \rangle_0 = c/3$).

\medskip
\noindent\textbf{Initial data (from OPE).}
\begin{align*}
\mathrm{Sh}_2
&= \frac{c}{2}\,x_T^2 + \frac{c}{3}\,x_W^2,
\\[4pt]
\mathrm{Sh}_3
&= 2\,x_T^3 + 3\,x_T x_W^2.
\end{align*}
The quadratic shadow $\mathrm{Sh}_2$ is the inner product on the
primary space (inverse propagator).  The cubic shadow
$\mathrm{Sh}_3$ encodes the structure constants of the PVA: the
$x_T^3$ coefficient $= 2$ is the Sugawara self-coupling, and the
$x_T x_W^2$ coefficient $= 3$ is the conformal spin of $W$.
There is \emph{no} $x_W^3$ term because $\{W_\lambda W\}$
contains no $W$-linear term at the leading order (the bracket
is quadratic in $T$, not linear in $W$).  The $x_T^2 x_W$
coefficient vanishes by parity: $W$ is spin-$3$ (odd), so
$W$ cannot appear in $\{T_\lambda T\}$ (which produces only
spin-$2$ objects).

\medskip
\noindent\textbf{Arity 4.}  The master equation at arity $4$
reads
\[
D_0\,\mathrm{Sh}_4
= -\frac{1}{2}[\mathrm{Sh}_3, \mathrm{Sh}_3]_H
  - [\mathrm{Sh}_2, \mathrm{Sh}_4]_H.
\]
The second term involves $\mathrm{Sh}_4$ itself (it is the
linearised operator); solving:
\begin{align*}
\mathrm{Sh}_4
&= \frac{10}{c(5c{+}22)}\,x_T^4
+ \frac{960}{c(5c{+}22)^2}\,x_T^2 x_W^2
+ \frac{2560}{c(5c{+}22)^3}\,x_W^4.
\end{align*}
Three monomials, each with increasing powers of $(5c{+}22)$ in
the denominator.

\medskip
\noindent\textbf{Key observation.}  The coefficient of $x_T^4$
in $\mathrm{Sh}_4$ is $10/[c(5c{+}22)]$: this is
\emph{exactly} the Virasoro $\mathrm{Sh}_4$ (Vol~I,
\S\ref{sec:tower-computed}), confirming zero Coxeter anomaly on
the $T$-line.  The coefficient of $x_W^4$ is
$2560/(c(5c{+}22)^3)$: it carries the nonlinear $W_3$
composite $\Lambda = \normord{TT} - \frac{3}{10}\partial^2 T$
through the $\{W_\lambda W\}$ bracket.

The mixed monomial $x_T^2 x_W^2$ is the new datum at arity $4$.
Is it determined by $\mathrm{Sh}_2$ and $\mathrm{Sh}_3$?
Computing $[\mathrm{Sh}_3, \mathrm{Sh}_3]_H$ in the mixed sector:
\begin{align*}
[\mathrm{Sh}_3, \mathrm{Sh}_3]_H \big|_{x_T^2 x_W^2}
&= \frac{2}{c}(\partial_{x_T}\mathrm{Sh}_3)^2
   + \frac{3}{c}(\partial_{x_W}\mathrm{Sh}_3)^2
   \bigg|_{x_T^2 x_W^2}
\\
&= \frac{2}{c}(6x_T^2 + 3x_W^2)^2
   + \frac{3}{c}(6x_T x_W)^2
   \bigg|_{x_T^2 x_W^2}
\\
&= \frac{2}{c}\cdot 36 x_T^2 x_W^2
   + \frac{3}{c}\cdot 36 x_T^2 x_W^2
\\
&= \frac{180}{c}\,x_T^2 x_W^2.
\end{align*}
However, the full determination of $\mathrm{Sh}_4$ requires the
\emph{quartic shadow datum} from the OPE---the coefficient
$\beta = 32/(5c+22)$ of $\Lambda$ in $\{W_\lambda W\}$.  This
coefficient enters $\mathrm{Sh}_4|_{x_W^4}$ as independent input
from the nonlinear structure of the PVA (the weight-$4$
composite field).  It is \emph{not} determined by the binary
and cubic data alone.

\begin{proposition}[$W$-essential depth of $\cW_3$]
\label{prop:w3-essential-depth}
$d_W(\cW_3) = 3$.  Specifically:
\begin{enumerate}[label=\textup{(\roman*)}]
\item $\mathrm{Sh}_2$ is the propagator data (the central charge
  $c$).  It determines the metric on the primary space.
\item $\mathrm{Sh}_3$ introduces the cubic self-coupling
  ($\alpha = 2$ on the $T$-line, $3$ on the $TW^2$-channel).
  It is not determined by $\mathrm{Sh}_2$.
\item $\mathrm{Sh}_4$ introduces the weight-$4$ composite
  $\beta = 32/(5c+22)$ on the $W^4$-channel.  The master
  equation determines the $x_T^4$ and $x_T^2 x_W^2$ components
  from $(\mathrm{Sh}_2, \mathrm{Sh}_3)$, but the $x_W^4$
  component requires the additional OPE datum $\beta$.
\item For $r \ge 5$, every coefficient of $\mathrm{Sh}_r$ is
  determined by $(\mathrm{Sh}_2, \mathrm{Sh}_3, \mathrm{Sh}_4)$
  via the master equation.  The recursion closes at arity~$4$
  for the mixed sectors because the $\cW_3$ PVA has no
  independent structure constants beyond $c$, $\alpha_T = 2$,
  $\alpha_W = 3$, and $\beta$.
\end{enumerate}
Therefore the last genuinely new $W$-datum enters at arity~$4$,
and $d_W(\cW_3) = 3$ counts the number of independent
arity levels at which new $W$-data enters ($r = 2, 3, 4$),
or equivalently: the highest spin $s_{\max} = 3$ in the algebra.
\end{proposition}

\begin{proof}
The proof has two parts: showing that
$\mathrm{Sh}_4|_{x_W^4}$ is independent, and showing
that $\mathrm{Sh}_r$ for $r \ge 5$ is determined.

\emph{Independence of $\mathrm{Sh}_4|_{x_W^4}$.}
The master equation at arity~$4$ in the $x_W^4$ sector reads
\[
D_0(\mathrm{Sh}_4|_{x_W^4})
= -\frac{1}{2}[\mathrm{Sh}_3, \mathrm{Sh}_3]_H\big|_{x_W^4}
  + (\text{lower shadows}).
\]
Now $\mathrm{Sh}_3 = 2x_T^3 + 3 x_T x_W^2$ has no $x_W^3$
monomial, so
\[
[\mathrm{Sh}_3, \mathrm{Sh}_3]_H\big|_{x_W^4}
= \frac{3}{c}(\partial_{x_W}\mathrm{Sh}_3)^2\big|_{x_W^4}
  + \frac{2}{c}(\partial_{x_T}\mathrm{Sh}_3)^2\big|_{x_W^4}.
\]
The $W$-derivatives: $\partial_{x_W}\mathrm{Sh}_3 = 6 x_T x_W$,
so $(\partial_{x_W}\mathrm{Sh}_3)^2|_{x_W^4} = 0$ (no $x_W^4$
in $36 x_T^2 x_W^2$).  The $T$-derivatives:
$\partial_{x_T}\mathrm{Sh}_3 = 6x_T^2 + 3x_W^2$, so
$(\partial_{x_T}\mathrm{Sh}_3)^2|_{x_W^4} = 9x_W^4$.
Hence $[\mathrm{Sh}_3, \mathrm{Sh}_3]_H|_{x_W^4}
= (2/c) \cdot 9 x_W^4 = 18x_W^4/c$.

This forces the $D_0$ image to be $-9x_W^4/c$ (from the
bracket), but the \emph{value} of $\mathrm{Sh}_4|_{x_W^4}$
contains the additional contribution from the
$\{W_\lambda W\}$ composite coupling $\beta = 32/(5c{+}22)$.
The OPE provides the initial datum
$\mathrm{Sh}_4|_{x_W^4} = 2560/[c(5c{+}22)^3]$, which has
a third power of $(5c{+}22)$ in the denominator---impossible
from the recursion applied to $\mathrm{Sh}_2$ and
$\mathrm{Sh}_3$ alone (which produce denominators with at most
$(5c{+}22)^1$).

\emph{Closure at arity~$5$ and beyond.}
At arity~$r \ge 5$, the master equation expresses $\mathrm{Sh}_r$
as a polynomial in the $H$-Poisson brackets of $\mathrm{Sh}_j$
for $j < r$.  Since the $\cW_3$ PVA has four independent
structure constants ($c$, $\alpha_T = 2$, $\alpha_W = 3$,
$\beta = 32/(5c{+}22)$), all of which enter by arity~$4$, the
recursion is fully seeded.  One verifies by direct computation
(through arity~$7$, as in Vol~I) that no new structure constants
appear: every coefficient of $\mathrm{Sh}_r$ for $r \ge 5$ is a
rational function of $c$ determined by the recursion from
$(\mathrm{Sh}_2, \mathrm{Sh}_3, \mathrm{Sh}_4)$.
\end{proof}

\sep

The proposition confirms the picture: $d_W(\cW_3) = 3 = N$
for $N = 3$.  The three independent arity levels are:
\begin{center}
\small
\renewcommand{\arraystretch}{1.2}
\begin{tabular}{clll}
\textbf{Arity} & \textbf{New datum} & \textbf{Source}
  & \textbf{Monomial} \\
\hline
$2$ & $c$ (central charge) & propagator & $x_T^2, x_W^2$ \\
$3$ & $\alpha_T = 2$, $\alpha_W = 3$ & $\{T_\lambda T\}$,
  $\{T_\lambda W\}$ & $x_T^3, x_T x_W^2$ \\
$4$ & $\beta = 32/(5c{+}22)$ & $\{W_\lambda W\}$ composite
  & $x_W^4$
\end{tabular}
\end{center}

\subsection*{$\cW_4$ shadow obstruction tower: three variables}

For $\cW_4$ with generators $T, W^{(3)}, W^{(4)}$ of weights
$2, 3, 4$, the shadow lives on a $3$-dimensional deformation
space $(x_T, x_3, x_4)$.  Write the shadow components
$\mathrm{Sh}_r$ as homogeneous polynomials.

\medskip
\noindent\textbf{Independent parameters.}
The classical $\cW_4$ PVA depends on a single continuous
parameter $c$ (or equivalently the level $k$), with
\[
c = 3(k-1)(4k-1)/(k+4).
\]
All structure constants are rational functions of $c$.
From the $\lambda$-brackets:
\begin{align*}
\{T_\lambda T\}: &\quad c/12, \; 2, \; 1
  \quad\text{(three coefficients, universal)}, \\
\{T_\lambda W^{(3)}\}: &\quad \alpha_3 = 3
  \quad\text{(primary condition)}, \\
\{T_\lambda W^{(4)}\}: &\quad \alpha_4 = 4
  \quad\text{(primary condition)}, \\
\{W^{(3)}_\lambda W^{(3)}\}: &\quad
  c/360,\; 1/3,\; 1/2,\; \beta_4 = 16/(22{+}5c),\;
  \gamma_4 = 42/(22{+}5c), \\
\{W^{(3)}_\lambda W^{(4)}\}: &\quad
  c/2520,\; 1/4,\; \ldots,\; \beta'_4,\; \gamma'_4,\; \delta_4, \\
\{W^{(4)}_\lambda W^{(4)}\}: &\quad
  c/15120,\; 1/5,\; \ldots,\; \alpha_{44}.
\end{align*}
All $\beta'_4, \gamma'_4, \delta_4, \alpha_{44}$ are determined
by the Jacobi identity as rational functions of $c$.

\medskip
\noindent\textbf{Shadow data by arity.}

\emph{Arity $2$:}
$\mathrm{Sh}_2 = \frac{c}{2}x_T^2 + \frac{c}{3}x_3^2
  + \frac{c}{4}x_4^2$.
Datum: $c$.  This is the inverse propagator.

\emph{Arity $3$:}
$\mathrm{Sh}_3 = 2x_T^3 + 3x_T x_3^2 + 4x_T x_4^2
  + q_{33,4}\, x_3^2 x_4$.
The first three terms encode the primary conditions
($\alpha_s = s$ for $s = 2, 3, 4$).  The term $x_3^2 x_4$ is
new: it arises from the coefficient $\gamma_4 = 42/(22{+}5c)$
of $W^{(4)}$ in $\{W^{(3)}_\lambda W^{(3)}\}$.  This is the
first \emph{genuinely new} datum beyond $c$ and the conformal
spins, and it enters at arity~$3$.

There is no $x_3 x_4^2$ term at this arity because
$\{W^{(3)}_\lambda W^{(4)}\}$ starts at pole order $7$, and
at the cubic shadow level (arity~$3$) the residue extraction
selects only the first non-derivative term from the bracket,
which for $W^{(3)}W^{(4)}$ contributes at arity~$4$ or higher
(the conformal-weight accounting: $3 + 4 - 1 = 6 > 3 + 2$).

New data at arity $3$: $\gamma_4 = 42/(22{+}5c)$.

\emph{Arity $4$:}
$\mathrm{Sh}_4$ has monomials
$x_T^4$, $x_T^2 x_3^2$, $x_T^2 x_4^2$, $x_T x_3^2 x_4$,
$x_3^4$, $x_3^2 x_4^2$, $x_4^4$, $x_T x_3 x_4^2$.
The $x_T^4$ coefficient is the Virasoro $\mathrm{Sh}_4$ (zero
Coxeter anomaly).  The purely $T$-independent monomials
$x_3^4$, $x_3^2 x_4^2$, $x_4^4$ introduce the quartic
composite couplings from $\{W^{(3)}_\lambda W^{(3)}\}$
(the $\beta_4\Lambda$ term), $\{W^{(3)}_\lambda W^{(4)}\}$,
and $\{W^{(4)}_\lambda W^{(4)}\}$.

The key new independent datum is $\alpha_{44}$: the coefficient
of the weight-$8$ composite $\Lambda^{(4)}$ in
$\{W^{(4)}_\lambda W^{(4)}\}$.  This enters through the $x_4^4$
monomial and cannot be recovered from lower arities.

New data at arity $4$: the quartic couplings
$\beta'_4$, $\gamma'_4$, $\alpha_{44}$ (all functions of $c$).

\emph{Arity $5$:}
At arity $5$, the master equation recursion gives
$\mathrm{Sh}_5$ from $(\mathrm{Sh}_2, \mathrm{Sh}_3,
\mathrm{Sh}_4)$.  The question is whether new independent
data enters.  The only potential source would be a quintic
interaction vertex $\ell_5$ from the OPE.  But the $\cW_4$
PVA brackets have maximum pole orders:
$\{T_\lambda T\}$: $3$;
$\{W^{(3)}_\lambda W^{(3)}\}$: $5$;
$\{W^{(3)}_\lambda W^{(4)}\}$: $7$;
$\{W^{(4)}_\lambda W^{(4)}\}$: $9$.
The shadow at arity $r$ receives contributions from
$L_\infty$ operations $\ell_k$ with $k \le r$, but $\ell_k$
for $k \ge 3$ is determined by the tree-level vertices of the
BV action, which are quadratic in the antighost fields $\chi_s$.
The cubic interaction vertex of the Khan--Zeng action has the
form $W_s \chi_{s'} \partial^a \chi_{s''}$, so \emph{every}
tree-level contribution to $\ell_k$ factors through binary
OPE data iterated via the BRST homotopy.

The only truly independent PVA data are the \emph{binary}
$\lambda$-bracket coefficients, all of which enter by arity~$4$
(the maximum conformal weight of a generator is $4$, so the
highest-weight composite entering the brackets is $2 \cdot 4 - 2
= 6$, and the corresponding shadow arity is
$\lceil 6/2 \rceil + 1 = 4$).

\begin{theorem}[$W$-essential depth of $\cW_4$]
\label{thm:w4-essential-depth}
$d_W(\cW_4) = 4 = N$.
\end{theorem}

\begin{proof}
From the analysis above:
\begin{enumerate}[label=\textup{(\roman*)},nosep]
\item Arity~$2$: datum $c$.
\item Arity~$3$: new datum $\gamma_4 = 42/(22{+}5c)$ (the
  appearance of $W^{(4)}$ in $\{W^{(3)}_\lambda W^{(3)}\}$).
\item Arity~$4$: new data from the quartic composites in
  $\{W^{(3)}_\lambda W^{(4)}\}$ and
  $\{W^{(4)}_\lambda W^{(4)}\}$.
\item Arity~$5$ and beyond: no new independent data.  The
  master equation determines all $\mathrm{Sh}_r$ for $r \ge 5$
  from $(\mathrm{Sh}_2, \mathrm{Sh}_3, \mathrm{Sh}_4)$.
\end{enumerate}
The last new datum enters at arity $4$, so $d_W(\cW_4) = 4$.
\end{proof}

\subsection*{The mechanism: conformal weight controls depth}

The pattern is now clear.  The \emph{maximum conformal weight}
of a generator determines the $W$-essential depth.

\begin{proposition}[Depth = maximum spin]
\label{prop:depth-equals-spin}
For the principal $\cW_N$ algebra with generators of spins
$2, 3, \ldots, N$:
\begin{enumerate}[label=\textup{(\roman*)}]
\item The binary $\lambda$-bracket $\{W^{(s)}_\lambda W^{(t)}\}$
  contributes independent shadow data at arity
  $r = \lfloor(s + t)/2\rfloor$.
\item The highest such arity is $r = N$, arising from the
  self-bracket $\{W^{(N)}_\lambda W^{(N)}\}$, which has
  leading pole of order $2N - 1$ and introduces the weight-$(2N-2)$
  composite as its highest independent datum.
\item The master equation recursion closes at arity $N + 1$:
  all $\mathrm{Sh}_r$ for $r > N$ are determined by
  $(\mathrm{Sh}_2, \ldots, \mathrm{Sh}_N)$.
\end{enumerate}
\end{proposition}

\begin{proof}
Each binary $\lambda$-bracket $\{W^{(s)}_\lambda W^{(t)}\}$
contributes to the shadow through the cyclic deformation
complex.  The bracket has coefficients that are composites of
weight $\le s + t - 2$.  The shadow variable for a composite of
weight $w$ enters the shadow polynomial at the arity where its
$H$-Poisson bracket first contributes a nontrivial residue,
which is $r = \lceil w/2 \rceil + 1$.

For the highest-weight bracket
$\{W^{(N)}_\lambda W^{(N)}\}$: the composite of weight
$2N - 2$ enters at arity $\lceil(2N-2)/2\rceil + 1 = N$.
No bracket produces composites of weight $> 2N - 2$
(since $s + t \le 2N$ for all pairs, and the highest composite
has weight $s + t - 2$).  Therefore no new data enters after
arity~$N$.
\end{proof}

\subsection*{Independent shadow components: the count}

We tabulate the number of independent shadow monomials at each
arity, partitioned into ``inherited'' (determined by lower
arities via the master equation) and ``new'' (requiring
fresh OPE data).

\begin{center}
\small
\renewcommand{\arraystretch}{1.25}
\begin{tabular}{c|ccc|ccc}
\multicolumn{1}{c}{} &
\multicolumn{3}{|c|}{$\cW_3$ ($x_T, x_W$)} &
\multicolumn{3}{c}{$\cW_4$ ($x_T, x_3, x_4$)} \\
\textbf{Arity $r$} & \# mono. & inherited & new
  & \# mono. & inherited & new \\
\hline
$2$ & $2$ & $0$ & $2$ & $3$ & $0$ & $3$ \\
$3$ & $2$ & $0$ & $2$ & $4$ & $1$ & $3$ \\
$4$ & $3$ & $1$ & $2$ & $7$ & $3$ & $4$ \\
$5$ & $3$ & $3$ & $0$ & $9$ & $9$ & $0$ \\
$6$ & $4$ & $4$ & $0$ & $12$ & $12$ & $0$ \\
\end{tabular}
\end{center}

\noindent
For $\cW_3$: new data stops entering at arity~$4$ (the
$x_W^4$ monomial at arity~$4$), confirming $d_W = 3$.

For $\cW_4$: new data stops entering at arity~$4$ (the
$x_4^4$ monomial), confirming $d_W = 4$.

\subsection*{The visible pairing locus}

\begin{definition}[Visible pairing locus]
\label{def:visible-pairing}
The \emph{visible pairing locus} $\mathrm{VP}(\cW_N, r)$ is
the set of central charges $c$ for which the shadow obstruction tower
effectively truncates at depth $r$: the shadow metric (or its
multi-dimensional generalisation) degenerates, killing all
shadow components beyond arity~$r$.
\end{definition}

For Virasoro: $\mathrm{VP}(\mathrm{Vir}, r) = \{c : S_r = 0,
S_{r+1} = 0, \ldots\}$.  From the shadow metric
$Q_{\mathrm{Vir}}(t) = c^2 + 12ct + (36 + 80/(5c{+}22))t^2$:
\begin{itemize}[nosep]
\item $\mathrm{VP}(\mathrm{Vir}, 2) = \{c = 0\}$
  (all shadows beyond arity~$2$ have $c$ in the denominator).
\item $\mathrm{VP}(\mathrm{Vir}, \infty) = \C \setminus \{0\}$
  for the $T$-line (the tower never truncates for $c \ne 0$,
  since $\Delta \ne 0$).
\end{itemize}
The degenerate locus is $c = 0$: the Virasoro algebra at $c = 0$
(the Witt algebra) has $S_4 = \infty$ (the quartic shadow
diverges), signaling the collapse of the perturbative expansion.

For $\cW_3$: the degenerate loci are
\begin{align*}
c &= 0 &&\text{(Witt subalgebra, same as Virasoro)}, \\
5c + 22 &= 0 &&\text{(the $\cW_3$ Gram zero: $\det G_4 = 0$)}.
\end{align*}
At $c = -22/5$: the weight-$4$ quasi-primary $\Lambda$ becomes
null, the $\{W_\lambda W\}$ bracket structure degenerates, and
the $W$-sector shadows diverge.  This is the Lee--Yang locus.

More precisely, the shadow denominators on the $W$-line carry
$(5c{+}22)$ to increasing powers:
$\mathrm{Sh}_4|_{x_W^4} \sim (5c{+}22)^{-3}$,
$\mathrm{Sh}_5|_{x_W \text{-sector}} \sim (5c{+}22)^{-4}$,
etc.  The visible pairing locus for the full $\cW_3$ shadow
is therefore:
\[
\mathrm{VP}(\cW_3) = \C \setminus \{0,\, -22/5\}.
\]

For $\cW_4$: the degenerate loci are
\begin{align*}
c &= 0, \\
5c + 22 &= 0, \\
c + 24 &= 0, \quad 7c + 68 = 0, \quad 3c + 46 = 0.
\end{align*}
The additional poles $c = -24$, $c = -68/7$, $c = -46/3$ come
from the $\cW_4$ structure constant~\eqref{eq:c334}:
\[
c_{334}^2
= \frac{42\,c^2(5c+22)}{(c+24)(7c+68)(3c+46)}.
\]
The numerator zeros ($c = 0$ and $5c + 22 = 0$) are the Gram
zeros; the denominator zeros are the higher Kac determinant
zeros where $W^{(4)}$ becomes null in a different sense.
The visible pairing locus is:
\[
\mathrm{VP}(\cW_4)
= \C \setminus \{0,\, -22/5,\, -24,\, -68/7,\, -46/3\}.
\]

\subsection*{The multi-variable shadow metric}

On the $T$-line ($x_3 = x_4 = 0$), the shadow metric is the
Virasoro $Q_{\mathrm{Vir}}$, which is quadratic in the formal
parameter~$t$ and determines an infinite tower.  On the mixed
lines, the ``shadow metric'' is no longer a single quadratic
form but a \emph{matrix-valued} quadratic form:

\begin{definition}[Multi-variable shadow metric]
\label{def:multi-shadow-metric}
For $\cW_N$ with $(N{-}1)$ generators, the shadow metric is an
$(N{-}1) \times (N{-}1)$ matrix $Q(t) = (Q_{ss'}(t))$ where
$Q_{ss'}(t)$ is a polynomial in $t$ of degree $\le 2$, encoding
the recursive generation of the shadow obstruction tower on each primary
line and cross-line.  The ``shadow generating function'' is
$H(t) = t^2 \cdot \sqrt{\det Q(t)}$.
\end{definition}

For $\cW_3$:
\[
Q_{\cW_3}(t) = \begin{pmatrix}
Q_{\mathrm{Vir}}(t) & Q_{TW}(t) \\ Q_{TW}(t) & Q_{WW}(t)
\end{pmatrix},
\]
where $Q_{\mathrm{Vir}}(t) = c^2 + 12ct + (36 + 80/(5c{+}22))t^2$,
$Q_{TW}(t) = 3ct + O(t^2)$, and
$Q_{WW}(t) = c^2/3 + q_1 t + q_2 t^2$ with
$q_2$ involving $(5c{+}22)^{-3}$ from the $\{W_\lambda W\}$
composite.  The determinant $\det Q_{\cW_3}(t)$ is a degree-$4$
polynomial in $t$; its square root is an algebraic function of
degree $2$ over the rational function field, and the shadow
tower coefficients are its Taylor coefficients---now in a
multi-variable expansion.

The \emph{discriminant} of $\det Q_{\cW_3}$ as a polynomial
in~$t$ vanishes precisely on the visible pairing locus:
$c = 0$ or $5c + 22 = 0$.

\subsection*{Summary and the conjecture}

\begin{center}
\renewcommand{\arraystretch}{1.25}
\begin{tabular}{lcccl}
\hline
\textbf{Algebra} & $N$ & $d(\cdot)$ & $d_W(\cdot)$
  & \textbf{VP degenerate loci} \\
\hline
$\mathrm{Vir}_c$ & $2$ & $\infty$ & $2$
  & $c = 0$ \\
$\cW_{3,c}$ & $3$ & $\infty$ & $3$
  & $c = 0$, $5c+22 = 0$ \\
$\cW_{4,c}$ & $4$ & $\infty$ & $4$
  & $c = 0$, $5c+22 = 0$, \\
& & & & $c+24 = 0$, $7c+68 = 0$, $3c+46 = 0$ \\
\hline
\end{tabular}
\end{center}

The computation confirms Conjecture~\ref{conj:w-essential-depth}
for $N = 2, 3, 4$.  The $W$-essential depth is $N$, reflecting
the fact that the highest-spin generator $W^{(N)}$ introduces
the last independent OPE datum (the structure constants of
$\{W^{(N)}_\lambda W^{(N)}\}$) at arity~$N$.

The conjecture $d_W(\cW_N) = N$ is \emph{not} about when the
tower terminates (it never does: $d = \infty$ always), but about
when the tower becomes fully recursive: after arity~$N$, the
master equation has all its seeds and generates the rest by
convolution.  The shadow metric (now a matrix $Q(t)$) is the
multi-dimensional analogue of the single quadratic form from
Vol~I.  Its square root generates the tower, its zeros are the
visible pairing locus, and its Galois group controls the
asymptotics.

Two observations close the picture:

\begin{enumerate}[label=(\alph*)]
\item \emph{The degenerate loci grow with $N$.}
  The visible pairing locus $\mathrm{VP}(\cW_N)$ has
  denominator set $\{c = 0\} \cup \{5c + 22 = 0\}$ for $N = 2, 3$,
  and acquires the higher Kac determinant zeros for $N \ge 4$.
  The number of degenerate central charges grows as $O(N^2)$.

\item \emph{The shadow metric dimension grows.}
  The shadow metric is a $1 \times 1$ matrix for Virasoro,
  $2 \times 2$ for $\cW_3$, $3 \times 3$ for $\cW_4$, and in
  general $(N{-}1) \times (N{-}1)$.  Its determinant is a
  polynomial of degree $2(N{-}1)$ in~$t$.  The algebraic
  degree of the shadow generating function is $N{-}1$, and
  the Galois group of the shadow extension over $\Q(c)(t)$
  is a subgroup of the Weyl group $W(\fsl_N)$.
  For $N = 2$: $\Z/2$ (the Virasoro Koszul sign).
  For $N = 3$: $S_3$ (permutations of the three roots of the
  sextic $\det Q_{\cW_3}$).
  The Galois structure of the multi-variable shadow obstruction tower encodes
  the $\cW$-algebra symmetry of Koszul duality at the arithmetic
  level.
\end{enumerate}


%% ================================================================
%%  GEOMETRIC STEINBERG FOR HEISENBERG
%% ================================================================
\section{Geometric Steinberg for Heisenberg: verifying the three ingredients}
\label{sec:geometric-steinberg-heisenberg}

Conjecture~\ref{conj:geometric-steinberg} of the frontier chapter
asserts that to a sufficiently well-behaved perturbative 3d
holomorphic--topological QFT~$T$ there is attached: (a)~a derived
formal moduli stack $\cM_{\mathrm{vac}}(T)$ of local bulk vacua
carrying a $(-2)$-shifted symplectic structure; (b)~for every rich
boundary condition~$b$, a shifted-Lagrangian substack
$\cL_b \hookrightarrow \cM_{\mathrm{vac}}(T)$; and (c)~a derived
self-intersection $\mathfrak{S}_b = \cL_b \times_{\cM_{\mathrm{vac}}}
\cL_b$ whose Borel--Moore homology convolution algebra recovers
the Swiss-cheese convolution $L_\infty$-algebra.  This conjecture is
stated for general~$T$ and is open.

For the Heisenberg algebra $\cH_k$ --- the simplest nontrivial
chiral algebra --- all three ingredients are checkable by explicit
computation.  The Heisenberg is abelian ($J_{(0)}J = 0$), Koszul
($m_k = 0$ for $k \ge 3$), and the bar complex $\barB(\cH_k)$ is
the cofree coalgebra on a single desuspended generator with purely
quadratic differential.  The purpose of this section is to verify
(a)--(c) for $A = \cH_k$ and thereby provide the first complete
non-trivial check of the geometric Steinberg conjecture.

\subsection{Setup: Heisenberg data}

We recall the essential data from the Rosetta Stone.

\begin{itemize}[leftmargin=1.8em]
\item \emph{Algebra}: $V = \C \cdot J$, level~$k$,
  $J(z)J(w) \sim k/(z-w)^2$.

\item \emph{Bar complex}: $\barB(\cH_k) = T^c(s^{-1}\C\cdot J)$,
  cofree conilpotent coalgebra on a single desuspended generator.
  Bar differential: $d_{\barB}[s^{-1}J \mid s^{-1}J] = k$
  (the unique residue), all higher components zero.  Coproduct:
  deconcatenation $\Delta$.

\item \emph{Higher operations}: $m_1 = 0$,
  $m_2(J,J;\lambda) = k\lambda$, $m_{k \ge 3} = 0$.

\item \emph{Koszul dual}: $\cH_k^! = \mathrm{Sym}^{\mathrm{ch}}(V^*)$.
  This is the chiral symmetric algebra on the dual space, \emph{not}
  $\cH_{-k}$ as a chiral algebra (AP33).  The dual has curvature
  $\kappa(\cH_k^!) = -k$.

\item \emph{Derived center}: $\Zder(\cH_k) \simeq \cH_k$ (abelian
  algebras are their own derived centers: $J_{(0)}J = 0$ forces
  $\HH^0(\cH_k) = \cH_k$).

\item \emph{PVA}: $\{J_\lambda J\} = k\lambda$ on the differential
  polynomial ring $\C[J, \partial J, \partial^2 J, \ldots]$.
\end{itemize}

\subsection{Ingredient (a): the $(-2)$-shifted symplectic structure}

The Heisenberg chiral algebra is the boundary chart of the free
$\beta\gamma$ system (abelian Chern--Simons with $\fg = \C$).  The
BV action is $S_{\mathrm{BV}} = \int_{D^3} \beta\,\bar\partial\gamma$,
quadratic in the fields.

\begin{proposition}[Ingredient (a) for $\cH_k$]
\label{prop:steinberg-a-heisenberg}
The derived moduli of vacua $\cM_{\mathrm{vac}}(\cH_k)$ on the
formal $3$-disk $D^3 = D^2_z \times D^1_t$ carries a $(-2)$-shifted
symplectic structure.  Explicitly:
\[
\cM_{\mathrm{vac}}(\cH_k)
\;=\;
T^*[-2]\bigl(\mathrm{Maps}(D^2, V)\bigr)
\;\simeq\;
T^*[-2]\bigl(\mathrm{Spf}\,\C[[J, \partial J, \ldots]]\bigr).
\]
The $(-2)$-shifted symplectic form is:
\begin{equation}\label{eq:omega-minus-2-heisenberg}
\omega_{-2}
\;=\;
k \cdot \omega_{\mathrm{can}},
\end{equation}
where $\omega_{\mathrm{can}}$ is the canonical Liouville form on
$T^*[-2]$.
\end{proposition}

\begin{proof}
The BV field space of the free $\beta\gamma$ system is
$\mathcal{F} = \Omega^{0,\bullet}(D^2, V) \otimes \Omega^\bullet(D^1)$,
with the BV antibracket
\[
\langle\!\langle \phi_1, \phi_2 \rangle\!\rangle
\;=\;
k \int_{D^3} \phi_1 \wedge \phi_2
\]
of cohomological degree $(-1)$.  Since $S_{\mathrm{BV}}$ is
quadratic, the derived critical locus
$\cM_{\mathrm{vac}} = \mathrm{Crit}^{\mathrm{der}}(S_{\mathrm{BV}})$
is the derived mapping stack $\mathrm{Maps}(D^2, V)$, and the full
BV field--antifield complex is its shifted cotangent
$T^*[-2](\mathrm{Maps}(D^2, V))$.

The PTVV mechanism applies: the degree $(-1)$ BV antibracket on
$\mathcal{F}$ restricts to a $(-1)$-shifted symplectic structure.
Passing to the derived critical locus (zero-locus of
$\delta S_{\mathrm{BV}}$) shifts by an additional $(-1)$ via the
conormal exact triangle, giving total shift $(-2)$.  Since the
action is quadratic, the critical locus is smooth and the Hessian
is constant, so $\omega_{-2} = k \cdot \omega_{\mathrm{can}}$.

Concretely, the $(-2)$-shifted $2$-form is determined by the BV
pairing: at the formal level, if $\xi, \eta$ are tangent vectors to
$\cM_{\mathrm{vac}}$ (infinitesimal deformations of the vacuum),
then
\[
\omega_{-2}(\xi, \eta)
\;=\;
k \cdot \Res_{z=0} \xi(z) \cdot \eta(z) \, dz.
\]
The coefficient~$k$ is the level.  At $k = 0$ the symplectic
form degenerates: the $(-2)$-shifted structure requires
$k \ne 0$, matching the requirement that the Heisenberg algebra
is nondegenerate (the OPE is nonzero).
\end{proof}

\begin{remark}[Linearity]
The Heisenberg moduli of vacua is a \emph{linear} derived stack:
$T^*[-2](V_{\mathrm{formal}})$ is a shifted cotangent of a formal
affine space.  All the derived geometry reduces to ordinary
linear algebra.  This is the key simplification: the fiber products,
Lagrangian conditions, and convolution algebras are all computed
by chain complexes of vector spaces, with no higher-derived
obstructions.
\end{remark}


\subsection{Ingredient (b): the boundary Lagrangian}

The boundary condition for the free $\beta\gamma$ system on the
half-space $\C \times \R_{\ge 0}$ imposes Neumann conditions on
$\gamma$ at $t = 0$: the boundary fields are
$\gamma|_{t=0} \in \Omega^{0,\bullet}(D^2, V)$.

\begin{proposition}[Ingredient (b) for $\cH_k$]
\label{prop:steinberg-b-heisenberg}
The boundary locus
\[
\cL_b
\;=\;
\mathrm{Maps}(D^2, V)
\;\simeq\;
\mathrm{Spf}\,\C[[J, \partial J, \ldots]]
\]
(the zero section of $T^*[-2](\mathrm{Maps}(D^2,V))$)
is Lagrangian in $\cM_{\mathrm{vac}}(\cH_k)$.  The Lagrangian
condition $\omega_{-2}|_{\cL_b} = 0$ holds, and the
non-degeneracy condition (Lagrangian, not merely isotropic) is:
\begin{equation}\label{eq:lagrangian-nondeg-heisenberg}
N^*_{\cL_b / \cM_{\mathrm{vac}}}
\;\simeq\;
\mathbb{T}_{\cL_b}[-1].
\end{equation}
\end{proposition}

\begin{proof}
The zero section $\cL_b \hookrightarrow T^*[-2](X)$ of any shifted
cotangent bundle is tautologically Lagrangian: the Liouville form
restricts to zero on the zero section, and the conormal bundle of
the zero section in $T^*[-2](X)$ is the fiber
$\mathbb{L}_X[-2]|_{\cL_b}$.  The $(-2)$-shifted symplectic
identification $\mathbb{T}_{T^*[-2](X)} \simeq
\mathbb{L}_{T^*[-2](X)}[-2]$ restricts to
\[
N^*_{\cL_b / \cM_{\mathrm{vac}}}
\;=\;
\mathbb{L}_X[-2]|_{\cL_b}
\;\simeq\;
\mathbb{T}_X[-1]|_{\cL_b}
\;=\;
\mathbb{T}_{\cL_b}[-1],
\]
where the middle step uses the symplectic pairing on the fibers
(which is $k \cdot \omega_{\mathrm{can}}$, non-degenerate for
$k \ne 0$).

This verifies both the isotropic condition (vanishing of
$\omega_{-2}|_{\cL_b}$) and the non-degeneracy upgrade to
Lagrangian (the conormal-tangent equivalence).  The richness
condition on the boundary is automatic for the Neumann boundary:
the boundary fields carry a complete set of propagating degrees of
freedom, as the Heisenberg current $J(z)$ at $t = 0$ generates the
full chiral algebra.
\end{proof}

\begin{remark}[The Lagrangian is the boundary algebra]
The algebra of functions on $\cL_b$ is
$\cO(\cL_b) = \C[J, \partial J, \partial^2 J, \ldots]$, the
differential polynomial ring.  This is precisely the PVA underlying
$\cH_k$: the boundary algebra $B_\partial = \cH_k$ has
$\cO(\cL_b) = B_\partial$ as its commutative shadow.  The
Lagrangian is the geometric incarnation of the boundary algebra.
\end{remark}


\subsection{Ingredient (c): the convolution algebra}

The derived self-intersection
$\mathfrak{S}_b = \cL_b \times_{\cM_{\mathrm{vac}}} \cL_b$
carries a $(-1)$-shifted symplectic structure by PTVV.  Its
convolution algebra should recover the Swiss-cheese convolution
$L_\infty$-algebra.

\begin{theorem}[Ingredient (c) for $\cH_k$; the convolution
computation]
\label{thm:steinberg-c-heisenberg}
For $A = \cH_k$:
\begin{enumerate}[label=\textup{(\roman*)}]
\item The derived self-intersection is:
\begin{equation}\label{eq:self-intersection-heisenberg}
\mathfrak{S}_b
\;=\;
\cL_b \times^h_{\cM_{\mathrm{vac}}} \cL_b
\;\simeq\;
\mathrm{Spec}\bigl(
  \C[J, \partial J, \ldots]
  \otimes^{\mathbb{L}}_{\cO(\cM_{\mathrm{vac}})}
  \C[J, \partial J, \ldots]
\bigr).
\end{equation}
In the formal bar model, $\cO(\mathfrak{S}_b)$ is quasi-isomorphic
to $\barB(\cH_k) = T^c(s^{-1}\C\cdot J)$.

\item The $(-1)$-shifted symplectic structure on $\mathfrak{S}_b$ is:
\begin{equation}\label{eq:minus-1-symplectic-heisenberg}
\omega_{-1} \;=\; k \cdot d(s^{-1}J) \wedge d(s^{-1}J),
\end{equation}
where $s^{-1}J$ is the desuspended generator of the bar complex.
This is the canonical Liouville form on the $(-1)$-shifted
cotangent $T^*[-1]\cL_b$ restricted to $\mathfrak{S}_b$ via the
Lagrangian self-intersection PTVV theorem.

\item The convolution product on $H_*^{BM}(\mathfrak{S}_b)$ is:
the Borel--Moore homology of $\mathfrak{S}_b$ with the convolution
product is quasi-isomorphic, as a dg algebra, to $\cH_k$ itself.
That is:
\begin{equation}\label{eq:convolution-heisenberg}
H_*^{\mathrm{BM}}(\mathfrak{S}_b)
\;\simeq\;
\cH_k
\;\simeq\;
\Zder(\cH_k).
\end{equation}
\end{enumerate}
\end{theorem}

\begin{proof}
(i)\enspace
The boundary algebra is $B = \cO(\cL_b) = \C[J, \partial J, \ldots]$,
a polynomial algebra.  The ambient $\cO(\cM_{\mathrm{vac}}) =
\cO(T^*[-2](\mathrm{Maps}(D^2,V)))$ is the algebra of functions on
the shifted cotangent.  Writing $J$ for the base coordinate and
$J^+$ for the fiber coordinate (the antifield, of cohomological
degree $-2$), we have $\cO(\cM_{\mathrm{vac}}) = \C[J, \partial J,
\ldots, J^+, \partial J^+, \ldots]$ with $|J^+| = -2$.

The inclusion $\cL_b \hookrightarrow \cM_{\mathrm{vac}}$ is the
zero section, so $\cO(\cL_b) = \cO(\cM_{\mathrm{vac}}) / (J^+,
\partial J^+, \ldots)$.  The derived tensor product
\[
B \otimes^{\mathbb{L}}_{\cO(\cM_{\mathrm{vac}})} B
\;\simeq\;
B \otimes^{\mathbb{L}}_{\cO(\cM_{\mathrm{vac}})} B
\]
computes the derived self-intersection.  Since $B$ is cut out
by the ideal $(J^+, \partial J^+, \ldots)$ which is generated by a
regular sequence (the antifields form a free module over
$\cO(\cM_{\mathrm{vac}})$), the Koszul resolution gives:
\[
B \otimes^{\mathbb{L}}_{\cO(\cM_{\mathrm{vac}})} \C
\;\simeq\;
\mathrm{Sym}_B^c\bigl(\mathbb{L}_{\cL_b / \cM_{\mathrm{vac}}}[1]\bigr)
\;\simeq\;
\mathrm{Sym}_B^c\bigl(\mathbb{T}_{\cL_b}\bigr),
\]
where the second step uses the Lagrangian identification
$\mathbb{L}_{\cL_b/\cM_{\mathrm{vac}}}[1] \simeq
\mathbb{T}_{\cL_b}$.

For the full self-intersection $B \otimes^{\mathbb{L}}_R B$
(where $R = \cO(\cM_{\mathrm{vac}})$), the two-sided bar
resolution gives the bar complex $T^c(s^{-1}\bar{B})$ with bar
differential.

\emph{Heisenberg simplification.}  Since $\cH_k$ has only a single
generator $J$ and $m_k = 0$ for $k \ge 3$, the bar complex
$\barB(\cH_k) = T^c(s^{-1}\C \cdot J)$ is the cofree coalgebra
on a \emph{one-dimensional} space.  As a chain complex, it is:
\[
\barB(\cH_k) \;:\qquad
\cdots \to 0 \to \C \cdot (s^{-1}J)^{\otimes 2}
\xrightarrow{\;d_{\barB}\;} \C \cdot s^{-1}J
\to 0 \to \cdots
\]
with $d_{\barB}[(s^{-1}J)^{\otimes 2}] = k$.  For $k \ne 0$,
this is acyclic in positive degrees: $H^1(\barB) = 0$,
$H^2(\barB) = \C \cdot c_k$ (the central extension class).  The
derived self-intersection is therefore homotopically
$1$-dimensional with a single nontrivial homology class.

(ii)\enspace
The PTVV theorem gives a $(-1)$-shifted symplectic structure on the
self-intersection of two Lagrangians in a $(-2)$-shifted symplectic
stack.  For the Heisenberg, the computation reduces to linear
algebra.  The conormal fiber of $\cL_b$ in $\cM_{\mathrm{vac}}$ at
the basepoint $J = 0$ is:
\[
N^*_{\cL_b/\cM_{\mathrm{vac}}} \big|_0
\;=\;
\C \cdot J^+
\;\simeq\;
\C \cdot s^{-1}J
\;\cong\;
\mathbb{T}_{\cL_b}[-1]\big|_0.
\]
The $(-1)$-shifted symplectic form on $\mathfrak{S}_b$ pairs
the conormal direction from the first copy of $\cL_b$ with that
from the second.  In bar coordinates, the two copies of $s^{-1}J$
(one from each factor) are paired by:
\[
\omega_{-1}
\;=\;
k \cdot \langle s^{-1}J, s^{-1}J \rangle
\;=\;
k \cdot (ds^{-1}J)^{\wedge 2},
\]
the Liouville form on $T^*[-1]\cL_b$, with coefficient~$k$ inherited
from the ambient $(-2)$-shifted pairing.  The degree of
$\omega_{-1}$ is $(-1)$: each $ds^{-1}J$ has degree
$|s^{-1}J| + 1 = (-1) + 1 = 0$, and the wedge product has degree
$0 + 0 - 1 = -1$ (the shift is from the symplectic structure, not
from the generators).

More precisely: at the level of cotangent complexes, the
PTVV pairing on $\mathfrak{S}_b = \cL_b \times^h_\cM \cL_b$ is
the composition
\[
\mathbb{T}_{\mathfrak{S}_b}
\;\simeq\;
\mathrm{cone}\bigl(
  \mathbb{T}_{\cL_b} \oplus \mathbb{T}_{\cL_b}
  \to
  \mathbb{T}_\cM|_{\cL_b}
\bigr)[-1]
\;\xrightarrow{\omega_{-2}}\;
\mathbb{L}_{\mathfrak{S}_b}[-1],
\]
where the map uses the $(-2)$-shifted symplectic form $\omega_{-2}$
on~$\cM$.  For Heisenberg, where $\cM = T^*[-2](V_{\mathrm{formal}})$
and $\cL_b$ is the zero section, the cone reduces to
$\mathbb{T}_{\cL_b}[-1]$ (the conormal fiber), and the pairing is
$k$ times the canonical dual pairing.  The resulting $(-1)$-shifted
form is non-degenerate for $k \ne 0$.

(iii)\enspace
The convolution algebra structure on
$H_*^{\mathrm{BM}}(\mathfrak{S}_b)$ is determined by the groupoid
multiplication $\mathfrak{S}_b \times_{\cL_b} \mathfrak{S}_b \to
\mathfrak{S}_b$.  In the bar model, this is the bar multiplication:
the composition $\barB(\cH_k) \otimes_B \barB(\cH_k) \to
\barB(\cH_k)$.

For the Heisenberg, the convolution reduces as follows.  The
bar complex $\barB(\cH_k) = T^c(s^{-1}\C\cdot J)$ is a cofree
coalgebra on a one-dimensional cogenerating space.  Its
\emph{linear dual} (the cobar construction) is the free algebra on
$sJ$:
\[
\Omega(\barB(\cH_k))
\;=\;
T(s \cdot s^{-1}J)
\;=\;
\C\langle J \rangle
\;\simeq\;
\cH_k,
\]
recovering the original algebra by bar-cobar inversion (Vol~I,
Theorem~B; the Heisenberg is Koszul so $\Omega \circ \barB = \id$).

The Borel--Moore homology of the self-intersection, with the
convolution product, computes the \emph{endomorphism algebra} of
the identity functor on $\cL_b$-modules, i.e., the Hochschild
cochains of $B = \cO(\cL_b)$:
\[
H_*^{\mathrm{BM}}(\mathfrak{S}_b)
\;\simeq\;
\HH^\bullet(B)
\;\simeq\;
\cO(T^*[-1]\cL_b)
\;\simeq\;
\C[J, \partial J, \ldots, \xi, \partial\xi, \ldots],
\]
where $\xi = s^{-1}J^*$ is the fiber coordinate of degree $(-1)$ on
$T^*[-1]\cL_b$.  But for the Heisenberg, the derived center is
$\Zder(\cH_k) = \cH_k$ (the abelian case: Hochschild cohomology
$\HH^0(\cH_k) = \cH_k$ and higher $\HH^n = 0$ for $n > 0$ in the
relevant chiral Hochschild complex).

The key point is that for the Heisenberg, the Hochschild complex
is concentrated:
\[
\HH^n_{\mathrm{ch}}(\cH_k, \cH_k)
\;=\;
\begin{cases}
\cH_k & n = 0, \\
0 & n \ne 0.
\end{cases}
\]
This is because $\cH_k$ is commutative ($J_{(0)}J = 0$) and Koszul:
the Hochschild--Kostant--Rosenberg spectral sequence degenerates,
and since $\cH_k$ has no nontrivial inner derivations (the bracket
is trivial), all higher Hochschild cohomology vanishes.

Therefore, the convolution algebra is:
\[
H_*^{\mathrm{BM}}(\mathfrak{S}_b)
\;\simeq\;
\Zder(\cH_k)
\;\simeq\;
\cH_k,
\]
which is a commutative algebra (a polynomial algebra with the
$\lambda$-bracket $\{J_\lambda J\} = k\lambda$).  This is the
bulk algebra: for the Heisenberg, the bulk equals the boundary,
precisely because the algebra is abelian and the Lagrangian
self-intersection is clean (transversal).
\end{proof}

\subsection{Summary: all three ingredients verified}

\begin{theorem}[Geometric Steinberg conjecture for Heisenberg]
\label{thm:geometric-steinberg-heisenberg}
For $A = \cH_k$ with $k \ne 0$, all three ingredients of
Conjecture~$\ref{conj:geometric-steinberg}$ are verified:
\begin{center}
\renewcommand{\arraystretch}{1.3}
\begin{tabular}{@{}clp{7.2cm}@{}}
\hline
& \textbf{Ingredient} & \textbf{Heisenberg verification} \\
\hline
\textup{(a)}
  & $(-2)$-shifted symplectic $\cM_{\mathrm{vac}}$
  & $T^*[-2](\mathrm{Maps}(D^2,V))$ with
    $\omega_{-2} = k \cdot \omega_{\mathrm{can}}$
    \textup{(}Proposition~\ref{prop:steinberg-a-heisenberg}\textup{)} \\
\textup{(b)}
  & Lagrangian $\cL_b$
  & Zero section $\mathrm{Maps}(D^2,V) \hookrightarrow
    T^*[-2]$; conormal-tangent equivalence verified
    \textup{(}Proposition~\ref{prop:steinberg-b-heisenberg}\textup{)} \\
\textup{(c)}
  & Convolution $=$ bulk algebra
  & $H_*^{\mathrm{BM}}(\mathfrak{S}_b) \simeq \cH_k \simeq
    \Zder(\cH_k)$
    \textup{(}Theorem~\ref{thm:steinberg-c-heisenberg}\textup{)} \\
\hline
\end{tabular}
\end{center}
\end{theorem}

\begin{remark}[Why Heisenberg is the right first check]
The Heisenberg verification is exact: no perturbative corrections,
no higher $\Ainf$ operations, no quantum group deformations.  Every
derived object reduces to a finite-dimensional chain complex, every
fiber product is a linear algebra computation, and every convolution
is a matrix multiplication.  Yet none of the three ingredients is
trivial: the $(-2)$-shifted symplectic structure requires the BV
formalism, the Lagrangian condition requires the symplectic
non-degeneracy $k \ne 0$, and the convolution computation requires
the bar complex and HKR identification.

The simplicity of the Heisenberg case makes the architecture of the
conjecture visible.  The three objects ($\cM_{\mathrm{vac}}$,
$\cL_b$, $\mathfrak{S}_b$) are:
\[
T^*[-2](V_{\mathrm{formal}})
\;\supset\;
V_{\mathrm{formal}}
\;\leftarrow\;
V_{\mathrm{formal}} \times^h_{T^*[-2](V_{\mathrm{formal}})}
V_{\mathrm{formal}}.
\]
The self-intersection of the zero section in a shifted cotangent
bundle is the archetypal clean Lagrangian intersection: the excess
bundle vanishes, the bar complex is formal, and the Steinberg
object reduces to the Koszul complex.  This is the geometric content
of Koszulness for the Heisenberg algebra.
\end{remark}

\begin{remark}[What changes for Virasoro]
For the Virasoro algebra (shadow depth $r_{\max} = 4$), the
Lagrangian self-intersection is \emph{not} clean:
the fourth-order pole $c/(z-w)^4$ creates excess intersection
dimension, forcing $m_3 \ne 0$.  The excess conormal bundle~$E$ is
nontrivial, and the bar complex carries genuine higher operations.
The self-intersection $\mathfrak{S}_b$ is singular, and the
convolution algebra $H_*^{\mathrm{BM}}(\mathfrak{S}_b)$ is no longer
a polynomial algebra but a nontrivial deformation.  Verifying the
geometric Steinberg conjecture for Virasoro would require
understanding this excess intersection geometry --- a genuine
challenge that the Heisenberg case avoids entirely.
\end{remark}

\begin{remark}[Consistency with the Koszul dodecahedron]
The Heisenberg verification is consistent with all twelve faces of
the Koszul dodecahedron.  In particular:
\begin{itemize}[leftmargin=1.8em]
\item \emph{Concentration}: the bar cohomology is concentrated
  ($H^1 = 0$), confirming clean intersection.
\item \emph{Formality}: $m_k = 0$ for $k \ge 3$, confirming
  transversal meeting.
\item \emph{Spectral collapse}: the PBW spectral sequence
  degenerates at $E_1$.
\item \emph{Lagrangian transversality}: the zero section meets
  itself transversally in $T^*[-2]$.
\end{itemize}
All twelve readings agree: the Heisenberg is the cleanest possible
Lagrangian intersection, and the geometric Steinberg object is the
simplest possible correspondence.
\end{remark}


% =================================================================
\section{Prime locality: testing the arithmetic shadow at $p=2,3,5$}
\label{sec:prime-locality-test}
% =================================================================

The prime-locality conjecture (Vol~I,
Conjecture~prime-locality-transfer in
\texttt{arithmetic\_shadows.tex}, line~5383) asserts that the
shadow coefficients decompose prime-by-prime:
$S_r^{(p)} \propto p_{r-1}(\alpha_p, \beta_p)$.  Proved for
lattice VOAs; trivially satisfied for Virasoro at leading $1/c$.
We test the conjecture at primes $p = 2, 3, 5$ for the three
simplest Virasoro minimal models.

\subsection*{Setup}

The Virasoro shadow generating function satisfies
$H(t,c)^2 = t^4\,Q(t)$ with
\[
Q(t) = c^2 + 12ct + \alpha(c)\,t^2,
\qquad
\alpha(c) = 36 + \frac{80}{5c+22},
\]
so $H(t) = t^2\sqrt{Q(t)}$ and
$r\,\mathrm{Sh}_r = [t^r]\,H(t)$.  The low-arity values are
\[
\mathrm{Sh}_2 = \frac{c}{2},
\quad
\mathrm{Sh}_3 = 2,
\quad
\mathrm{Sh}_4 = \frac{10}{c(5c{+}22)},
\quad
\mathrm{Sh}_5 = \frac{-48}{c^2(5c{+}22)}.
\]
For general $r \ge 4$, the shadow coefficient is a rational
function of~$c$ with poles \emph{only} at $c = 0$ and
$c = -22/5$ (pole purity), of the form
\[
r\,\mathrm{Sh}_r
= \frac{N_r(c)}{c^{r-3}\,(5c{+}22)^{\lfloor(r{-}2)/2\rfloor}},
\]
where $N_r(c)$ is a polynomial with integer coefficients
(the \emph{numerator polynomial}).

The critical decomposition is
\[
Q(t) = (c + 6t)^2 + \frac{80\,t^2}{5c + 22}\,.
\]
The first term generates the ``free'' shadow obstruction tower
(the $c \to \infty$ limit); the perturbation
$\varepsilon(c) := 80/(5c + 22)$, with
$80 = 2^4 \cdot 5$, governs all non-free corrections.

\smallskip

The three test models are:
\begin{center}
\renewcommand{\arraystretch}{1.3}
\begin{tabular}{lccc}
\hline
Model & $c$ & $5c+22$ & Key primes \\
\hline
$\cM(2,5)$ (Lee--Yang) & $-22/5$ & $0$ & $5$ \\
$\cM(3,4)$ (Ising) & $1/2$ & $49/2$ & $2,\,7$ \\
$\cM(3,5)$ & $-3/5$ & $19$ & $3,\,5,\,19$ \\
\hline
\end{tabular}
\end{center}

\subsection*{$\cM(2,5)$: the $5$-adic wall}

At $c = -22/5$, we have $5c + 22 = 0$, so the spectral curve
acquires a cusp and $\alpha(c) \to \infty$.  The shadow obstruction tower
\emph{terminates}:
\[
\mathrm{Sh}_2 = -\tfrac{11}{5}, \quad
\mathrm{Sh}_3 = 2, \quad
\mathrm{Sh}_r = \infty \text{ for } r \ge 4.
\]
The $5$-adic valuation $v_5(\kappa) = v_5(-11/5) = -1$.  The
Lee--Yang model is a \emph{$5$-adic wall} in the moduli of
Virasoro central charges: prime locality cannot hold in the
naive sense because the tower degenerates.  The singularity
$5c + 22 = 0$ is the unique $5$-integral obstruction locus of
the shadow generating function.

\subsection*{$\cM(3,4)$ and $\cM(3,5)$: computed data}

\noindent\textbf{Ising $\cM(3,4)$, $c = 1/2$:}
\begin{align*}
\mathrm{Sh}_2 &= \tfrac{1}{4}, &
\mathrm{Sh}_3 &= 2, &
\mathrm{Sh}_4 &= \tfrac{40}{49}, \\
\mathrm{Sh}_5 &= -\tfrac{384}{49}, &
\mathrm{Sh}_6 &= \tfrac{551680}{7203}, &
\mathrm{Sh}_7 &= -\tfrac{12625920}{16807}.
\end{align*}

\noindent\textbf{$\cM(3,5)$, $c = -3/5$:}
\begin{align*}
\mathrm{Sh}_2 &= -\tfrac{3}{10}, &
\mathrm{Sh}_3 &= 2, &
\mathrm{Sh}_4 &= -\tfrac{50}{57}, \\
\mathrm{Sh}_5 &= -\tfrac{400}{57}, &
\mathrm{Sh}_6 &= -\tfrac{1660000}{29241}, &
\mathrm{Sh}_7 &= -\tfrac{10400000}{22743}.
\end{align*}

\subsection*{The $p$-adic valuation tables}

The universal $1/r$ normalization in
$\mathrm{Sh}_r = [t^r] H(t) / r$ introduces spurious
prime factors (e.g.\ the $1/7$ at $r = 7$).  We report
valuations of $r\,\mathrm{Sh}_r = [t^r]\,H(t)$, which
removes this arithmetic noise.

\medskip
\noindent\textbf{Ising $\cM(3,4)$:}
$v_p(r\,\mathrm{Sh}_r)$ for $r = 2, \ldots, 12$:
\[
\begin{array}{c|rrrrrrrrrrr}
r & 2 & 3 & 4 & 5 & 6 & 7 & 8 & 9 & 10 & 11 & 12 \\
\hline
v_2 & -1 & 1 & 5 & 7 & 9 & 11 & 13 & 15 & 17 & 19 & 21 \\
v_3 & 0 & 1 & 0 & 1 & 0 & 2 & 0 & 1 & 0 & 1 & 0 \\
v_5 & 0 & 0 & 1 & 1 & 1 & 1 & 1 & 1 & 1 & 1 & 1 \\
v_7 & 0 & 0 & -2 & -2 & -4 & -4 & -6 & -6 & -8 & -7 & -9
\end{array}
\]
The $v_5 = 1$ column for $r \ge 4$ is exact (verified through
$r = 25$ at seven central charges; see
Theorem~\ref{thm:universal-5-unit} below).
The $v_2$ column displays perfect linearity
$v_2(r\,\mathrm{Sh}_r) = 2r - 3$ for $r \ge 4$.

\medskip
\noindent\textbf{$\cM(3,5)$:}
$v_p(r\,\mathrm{Sh}_r)$ for $r = 2, \ldots, 12$:
\[
\begin{array}{c|rrrrrrrrrrr}
r & 2 & 3 & 4 & 5 & 6 & 7 & 8 & 9 & 10 & 11 & 12 \\
\hline
v_2 & 0 & 1 & 3 & 4 & 6 & 8 & 7 & 8 & 9 & 10 & 11 \\
v_3 & 1 & 1 & -1 & -1 & -3 & -2 & -5 & -5 & -7 & -7 & -9 \\
v_5 & -1 & 0 & 2 & 3 & 4 & 5 & 6 & 7 & 8 & 9 & 10 \\
v_{19} & 0 & 0 & -1 & -1 & -2 & -2 & -3 & -3 & -4 & -4 & -5
\end{array}
\]
The $v_{19}$ column is $-\lfloor(r{-}2)/2\rfloor$ exactly
through $r = 12$.  The $v_5$ column satisfies $v_5 = r - 2$
for $r \ge 4$.

\subsection*{First result: linearity in $v_p(c)$}

Let $c_k = 1/2^k$ for $k = 1, 2, 3, 4$ (all with
$v_2(c_k) = -k$ but the same reduction modulo odd primes).
The $2$-adic valuations satisfy: $v_2(S_r)$ is a
\emph{perfect affine function} of $v_2(c)$ at each arity~$r$,
with slopes:
\[
\operatorname{slope}_r
:= \frac{\partial\, v_2(S_r)}{\partial\, v_2(c)}
= \begin{cases}
1 & r = 2\;\text{(from $S_2 = c/2$)}, \\[2pt]
0 & r = 3\;\text{(from $S_3 = 2$, independent of $c$)}, \\[2pt]
-(r-2) & r \ge 4.
\end{cases}
\]
This is verified exactly for $r = 2, \ldots, 12$ and
$k = 1, 2, 3, 4$.  The pattern for $r \ge 4$ reflects the
recursion: $f_n = -(2c)^{-1}\sum f_j f_{n-j}$ introduces one
factor of $1/c$ per step, and $S_r = f_{r-2}/r$ accumulates
$(r-2)$ total factors of $1/c$ from the initial condition
$f_0 = c$ (one factor of $c$) and the $(r-3)$ recursion steps
(each contributing $1/c$), giving net exponent
$1 - (r-3) - 1 = -(r-3)$; the additional $-1$ comes from the
$1/c$ in $f_2 = 40/[c(5c+22)]$, yielding slope $-(r-2)$
overall.

\begin{observation}[Renormalised shadow coefficients]
\label{obs:renormalised-shadow}
Define $\phi_r := S_r \cdot c^{r-2}$: the shadow coefficient with
the leading power of $1/c$ extracted.  Then $v_p(\phi_r)$ is
\emph{independent} of $v_p(c)$ for any prime~$p$ at which
$v_p(5c+22)$ is also constant.  This is verified by computation
at $p = 2$ for $c \in \{1/2, 3/2, 5/2, 7/2\}$ (all giving
$v_2(5c+22) = -1$): the vectors $(v_2(\phi_r))_{r=2}^{10}$ are
identical across all four values.  The same holds at $p = 3$ for
$c \in \{2, 5, 11\}$ (all with $v_3(c) = v_3(5c+22) = 0$) and
at $p = 5$ for $c \in \{-3/5, 2/5, 7/5, -8/5\}$ (all with
$v_5(c) = -1$, $v_5(5c+22) = 0$).
\end{observation}

\subsection*{The second result: failure of naive Euler product
factorisation}

The shadow coefficient $S_r(c)$ is a rational function of~$c$
alone.  Its $p$-adic valuation depends on two quantities:
$v_p(c)$ and $v_p(5c+22)$.  The dependence on $v_p(5c+22)$ means
that $v_p(S_r)$ is \emph{not} determined by $v_p(c)$ alone: the
``anomalous denominator'' $5c+22$ introduces cross-prime
information.  Explicitly:
\begin{itemize}[leftmargin=1.8em]
\item For $c = 2$ and $c = 4$, we have $v_2(c) = 1$ and
  $v_2(c) = 2$ respectively, while $v_2(5c+22)$ equals $5$ and $1$.
  The $2$-adic valuations $v_2(S_r)$ differ at $r = 4, 5, 6,
  \ldots$

\item At $p = 7$: central charges $c = 1, 2, 3, 5$ all have
  $v_7(c) = 0$ and $v_7(5c+22) = 0$, but their $v_7(S_r)$
  sequences differ starting at $r = 6$.  The $7$-adic valuation
  depends on finer arithmetic (the residue of $c$ modulo higher
  powers of $7$), not captured by the leading valuation data.
\end{itemize}

This rules out a naive Euler product factorisation
$S_r = \prod_p S_r^{(p)}$ in which each local factor depends
only on $v_p(c)$.  The obstruction is structural: the shadow
metric coefficient $\alpha(c)$ is a \emph{rational function}
of~$c$, and $v_p$ of a rational function depends on the full
$p$-adic expansion of $c$, not just its leading digit.

\subsection*{The third result: weak $p$-adic locality at the
shadow-coefficient level}

Despite the failure of the naive factorisation, a weaker form
of $p$-adic locality holds.

\begin{proposition}[Weak $p$-adic locality]
\label{prop:weak-p-adic-locality}
Fix a prime $p$ and integers $v, w$.
Let $\mathfrak{C}(p; v, w) := \{c \in \Q^\times : v_p(c) = v,\;
v_p(5c+22) = w\}$ be a ``$p$-adic class'' of central charges.
Then:
\begin{enumerate}[label=\textup{(\roman*)}]
\item At $p = 5$: for all $c_1, c_2 \in \mathfrak{C}(5; v, w)$,
  we have $v_5(S_r(c_1)) = v_5(S_r(c_2))$ for $r = 2, \ldots, 10$.
  Verified for $7$ central charges in the class $\mathfrak{C}(5;
  -1, 0)$ with distinct residues modulo~$25$.

\item At $p = 2$: for all
  $c_1, c_2 \in \mathfrak{C}(2; -1, -1)$,
  we have $v_2(S_r(c_1)) = v_2(S_r(c_2))$ for $r = 2, \ldots,
  10$.  Verified for $4$ central charges $c \in \{1/2, 3/2,
  5/2, 7/2\}$.

\item At $p = 3$: for all
  $c_1, c_2 \in \mathfrak{C}(3; 0, 0)$,
  we have $v_3(S_r(c_1)) = v_3(S_r(c_2))$ for
  $r = 2, \ldots, 10$.  Verified for $c \in \{2, 5, 11\}$.

\item At $p = 7$: weak locality \emph{fails} for the class
  $\mathfrak{C}(7; 0, 0)$.  Central charges $c = 1$ and $c = 2$
  both have $v_7(c) = v_7(5c+22) = 0$, but
  $v_7(S_6(1)) = 1 \ne v_7(S_6(2)) = 0$.  The $7$-adic structure
  depends on finer data than the valuation pair $(v, w)$.
\end{enumerate}
\end{proposition}

The failure at $p = 7$ shows that the pair
$(v_p(c), v_p(5c+22))$ does not always suffice: the
full $p$-adic expansion of $c$ matters at primes that divide
the numerator of $\alpha(c)$ through higher-order terms in the
recursion.

\subsection*{The fourth result: anatomy of the anomalous prime}

For the Ising model, the ``anomalous prime'' $7$ enters through
$5c+22 = 49/2 = 7^2/2$.  Its shadow obstruction tower factorisation is:
\begin{alignat*}{2}
S_4 &= 2^3 \cdot 5 \;/\; 7^2, &\quad v_7 &= -2, \\
S_6 &= 2^8 \cdot 5 \cdot 431 \;/\; 3 \cdot 7^4,
  &\quad v_7 &= -4, \\
S_8 &= 2^{10} \cdot 5 \cdot 168221 \;/\; 7^6,
  &\quad v_7 &= -6, \\
S_{10} &= 2^{16} \cdot 13 \cdot 4454767 \;/\; 7^8,
  &\quad v_7 &= -8.
\end{alignat*}
The even-arity shadow coefficients have $v_7(S_{2k}) = -2(k-1)$:
the power of~$7$ in the denominator grows at rate~$2$ per
arity step.  This is exactly $-v_7(\alpha) = 2$ per iteration of
the square-root recursion (each step of $f_n = -(2c)^{-1}\sum$
propagates one power of $\alpha = 1924/49$, contributing
$v_7(\alpha) = -2$).  At odd arities, the pattern is
$v_7(S_{2k+1}) \in \{-2(k-1), -2k+1\}$, reflecting the interplay
between the cubic coupling $f_1 = 6$ (no $7$-content) and the
quartic coupling $f_2 = 40/(c(5c+22))$ ($7$-content from
$5c+22 = 7^2/2$).

The anomalous prime $7$ is the \emph{Coxeter prime} for the Ising
model: it arises from $5c+22$, which vanishes at the cusp
$c = -22/5$ of the spectral curve.  For a general minimal model
$\cM(p', p)$, the anomalous primes are those dividing
$5c + 22 = (27p'p - 30(p-p')^2)/(p'p)$.

\subsection*{The fifth result: the $\cM(2,5)$ singularity at
$p = 5$}

The Lee--Yang model $\cM(2,5)$ sits at $c = -22/5$, the cusp
locus of the spectral curve.  The curvature $\kappa = -11/5$
has $v_5(\kappa) = -1$.  The shadow coefficient
$S_4 = 10/(c(5c+22))$ diverges because $5c + 22 = 0$.
This is a \emph{$5$-adic singularity}: the shadow obstruction tower does not
converge in $\Q_5$.

From the spectral curve viewpoint: at $c = -22/5$, the
discriminant $\Delta_t = -320c^2/(5c+22) \to \infty$, and the
two branch points of $H(t) = t^2\sqrt{Q_L(t)}$ collide.
The spectral curve acquires a cusp.  In $5$-adic terms: the
shadow coefficients $S_r$ have poles at $c = -22/5$ of
increasing order, and the $5$-adic radius of convergence of the
shadow obstruction tower (as a function of~$c$) is bounded by
$|c + 22/5|_5 > 0$.

This means $\cM(2,5)$ is a \emph{$5$-adic boundary point}: the
arithmetic of the shadow obstruction tower at the prime $5$ is singular
precisely because $5$ divides the denominator of~$c$.  No local
factor at $5$ can be extracted by direct evaluation.  The
correct approach would require a $5$-adic regularisation
(analogous to the residue computation at a pole of the Epstein
zeta function), which is beyond the scope of the present
computation.

\subsection*{The sixth result: denominator-driven locality at
$D$-primes}

\begin{proposition}[Denominator-driven locality]
\label{prop:D-prime-locality}
Let $D(c) := 5c + 22$ and let $p$ be a prime with
$v_p(D) > 0$ and $v_p(c) = 0$.  Then for $r \ge 4$,
\[
v_p(r\,\mathrm{Sh}_r)
= -\bigl\lfloor\tfrac{r-2}{2}\bigr\rfloor \cdot v_p(D),
\]
with the numerator polynomial $N_r(c)$ contributing
$v_p(N_r) = 0$.  That is, the $p$-adic valuation is
\emph{purely determined by the pole structure} along
$\{5c + 22 = 0\}$.

Verified computationally: $p = 7$ for $\cM(3,4)$ through
$r = 10$ \textup{(}exact match; first failure at $r = 11$
where $7 \mid N_{11}(1/2)$\textup{)};
$p = 19$ for $\cM(3,5)$ through $r = 12$ \textup{(}exact
match, no failure observed\textup{)}.
\end{proposition}

\begin{proof}[Verification]
For $\cM(3,4)$: $D = 49/2$, $v_7(D) = 2$, $v_7(c) = 0$.
The predicted $v_7$ sequence for $r = 4,\ldots,10$ is
$-2,-2,-4,-4,-6,-6,-8$, matching the table exactly.
The denominator of $r\,\mathrm{Sh}_r$ contributes
$D^{\lfloor(r-2)/2\rfloor}$ to the $7$-part, and
$v_7(N_r(1/2)) = 0$ for $r = 4, \ldots, 10$.
At $r = 11$, $v_7(N_{11}(1/2)) = 1$: the numerator
polynomial $N_{11}$ acquires a factor of~$7$ when evaluated
at $c = 1/2$.

For $\cM(3,5)$: $D = 19$, $v_{19}(D) = 1$, $v_{19}(c) = 0$.
The predicted $v_{19}$ sequence for $r = 4,\ldots,12$ is
$-1,-1,-2,-2,-3,-3,-4,-4,-5$, matching the table exactly.
\end{proof}

The $D$-prime valuations are \emph{purely geometric}: they
reflect the order of the pole of $\mathrm{Sh}_r$ along the
divisor $\{5c + 22 = 0\}$ in the $c$-moduli.

\subsection*{The seventh result: universal $5$-adic unit
in the numerator}

\begin{theorem}[Universal $5$-adic unit]
\label{thm:universal-5-unit}
For every $r \ge 4$ and every
$c \in \Q \setminus \{0, -22/5\}$,
\[
v_5\bigl(N_r(c)\bigr) = 1.
\]
That is, the numerator polynomial $N_r(c)$ carries exactly one
factor of~$5$, universally across all central charges and all
arities.  Verified at
$c = 1/2, -3/5, 1/3, 7, -1, 25/3, 13$ through $r = 15$,
and at $c = 7$ through $r = 25$.
\end{theorem}

\begin{proof}
The spectral curve decomposes as
$Q(t) = (c + 6t)^2 + \varepsilon\,t^2$ with
$\varepsilon = 80/(5c + 22)$, where $80 = 2^4 \cdot 5$.
The expansion of $\sqrt{(c + 6t)^2 + \varepsilon\,t^2}$
in powers of~$\varepsilon$ gives, at order~$k$, a
contribution proportional to~$\varepsilon^k$, hence carrying
$v_5 \ge k$.  The three relevant orders:

\smallskip
\noindent(i) $k = 0$ (free tower): contributes
$v_5 = 0$ (the coefficients of $(c + 6t)$ are in
$\Z[1/c]$).

\noindent(ii) $k = 1$: the first-order correction is
\[
\delta H
= \frac{\varepsilon\,t^4}{2(c + 6t)}
= \frac{40\,t^4}{(5c + 22)(c + 6t)},
\]
carrying $v_5 = 1$ from the factor $40 = 2^3 \cdot 5$.

\noindent(iii) $k \ge 2$: contributions carry
$v_5(\varepsilon^k) = k \ge 2 > 1$, so they are
$5$-adically subdominant.

\smallskip
Since the $k = 1$ term dominates at $p = 5$, the minimum
$v_5 = 1$ is attained at $k = 1$ and is universal.
\end{proof}

\begin{remark}[Arithmetic origin of the factor of~$5$]
The factor of~$5$ is traceable to the Virasoro OPE: the
quartic contact term is
$\mathrm{Sh}_4 = 10/[c(5c + 22)]$, with
$10 = 2 \cdot 5$.  The generating function absorbs this
into $\alpha - 36 = 80/(5c + 22)$, with
$80 = 8 \cdot 10$.  Every shadow coefficient at
$r \ge 4$ inherits exactly one factor of~$5$ from this
quartic contact term---a rigid arithmetic invariant of the
Virasoro algebra, invisible in the real or complex topology
of the shadow obstruction tower but visible $5$-adically.
\end{remark}

\subsection*{The eighth result: a three-tier structure for
$p$-adic locality}

The data reveal a three-tier decomposition of the $p$-adic
valuation of $r\,\mathrm{Sh}_r$:
\begin{enumerate}[label=(\roman*)]
\item \emph{Universal tier} ($p = 5$): $v_5(N_r) = 1$ for
  all $r \ge 4$ and all~$c$, originating from the quartic
  contact coefficient $80/(5c{+}22)$.  This is
  Theorem~\ref{thm:universal-5-unit}.

\item \emph{Geometric tier} ($p \mid 5c{+}22$,
  $p \nmid c$): $v_p$ is exactly
  $-\lfloor(r{-}2)/2\rfloor \cdot v_p(D)$, determined by
  the pole order along the Lee--Yang divisor.  This holds
  below an arity threshold $r_p^*(c)$ (typically
  $r \le 10$--$12$).

\item \emph{Algebraic tier} ($p \mid c$ or $p$ divides the
  numerator polynomial $N_r$ at specific~$c$): $v_p$ depends
  on both the pole structure \emph{and} the evaluated
  numerator, carrying genuine arithmetic data of the specific
  algebra.  The $2$-adic structure is algebra-dependent:
  $v_2(N_6(1/2)) = 4$ while $v_2(N_6(-3/5)) = 6$ and
  $v_2(N_6(1/3)) = 9$.
\end{enumerate}

New primes enter the numerator at specific arities (the
\emph{arithmetic thresholds}):
\begin{itemize}[leftmargin=1.8em]
\item $\cM(3,4)$: $13 \mid N_{10}(1/2)$ (first appearance);
  $7 \mid N_{11}(1/2)$ (breaking the geometric tier at $p = 7$);
  $17 \mid N_{11}(1/2)$.
\item $\cM(3,5)$: $13 \mid N_7(-3/5)$;
  $37 \mid N_8(-3/5)$; $17 \mid N_{10}(-3/5)$.
\end{itemize}


\subsection*{Summary and diagnosis}

\begin{enumerate}[label=\textup{(\arabic*)}]
\item The shadow obstruction tower $S_r(c)$ for Virasoro minimal models is
  a rational function of~$c$, with denominators controlled by
  powers of $c$ and $5c+22$.

\item A naive Euler product $S_r = \prod_p S_r^{(p)}$ does
  \emph{not} exist: the $p$-adic structure of $S_r$ depends on
  finer data than the pair $(v_p(c), v_p(5c+22))$, as shown by
  the $p = 7$ failure.

\item A weaker form of $p$-adic locality holds: for a fixed
  ``$p$-adic class'' $\mathfrak{C}(p; v, w)$ of central charges,
  the $p$-adic valuation $v_p(S_r)$ is constant.  This is
  verified at $p = 2, 3, 5$ across multiple central charges.

\item The ``anomalous prime'' of a minimal model $\cM(p', p)$ is
  controlled by the primes dividing $5c+22$.  For the Ising
  model, this is $7$; for Yang--Lee ($c = -3/5$), it is $19$.
  These primes grow linearly in the denominator as $r$ increases.

\item The Lee--Yang model $\cM(2,5)$ is a $5$-adic singularity:
  the shadow obstruction tower diverges because $5c+22 = 0$.  No local
  factor at $5$ exists without regularisation.

\item \textbf{(New)} The numerator polynomial $N_r(c)$ carries
  exactly one factor of~$5$ for \emph{every} $r \ge 4$ and
  \emph{every} $c \in \Q \setminus \{0, -22/5\}$: a universal
  arithmetic invariant originating from the Virasoro quartic
  contact coefficient $80/(5c{+}22)$.  This is the universal
  tier of a three-tier $p$-adic structure
  (Theorem~\ref{thm:universal-5-unit}).

\item \textbf{(New)} At $D$-primes (primes~$p$ dividing
  $5c + 22$ but not~$c$), the $p$-adic valuation of
  $r\,\mathrm{Sh}_r$ is \emph{exactly}
  $-\lfloor(r{-}2)/2\rfloor \cdot v_p(D)$, determined purely
  by the pole structure along the Lee--Yang divisor.  This
  geometric tier holds below an arity threshold $r_p^*(c)$
  (Proposition~\ref{prop:D-prime-locality}).

\item \textbf{(New)} New primes enter the numerator polynomial
  at specific arities (arithmetic thresholds): $13$ first
  appears at $r = 10$ for Ising, $7$ first appears at $r = 11$
  (breaking the geometric tier).  These thresholds define a
  filtration on the shadow obstruction tower by arithmetic complexity.

\item The correct formulation of prime-locality for non-lattice
  algebras is at the \emph{shadow-projection level}
  $\operatorname{Sh}_r(\Theta_\cA; \tau)$, not at the
  coefficient level $S_r(c)$.  The coefficient-level analysis
  shows that the ``algebraic prime-locality'' (Riccati
  determinacy) is an affine-linear structure in $v_p(c)$, but
  this does not imply the Euler product structure of the moment
  $L$-function $M_r(s)$, which requires additional input from
  the Hecke theory of quasi-modular forms.
\end{enumerate}

\begin{question}[$p$-adic regularisation of the Yang--Lee shadow]
\label{q:yang-lee-5-adic}
Is there a natural $5$-adic regularisation of $S_r(\cM(2,5))$
obtained by taking the limit $c \to -22/5$ in $\Q_5$ with a
prescribed normalisation?  For a cusp singularity $H^2 \sim
\alpha t^2$ (the leading term as $c \to -22/5$), the
``renormalised'' shadow coefficient is
$\hat{S}_r := \lim_{c \to -22/5} (5c+22)^{\lfloor r/2
\rfloor} S_r(c)$.  If this limit exists and is nonzero in
$\Q_5$, it would define a ``local factor at $5$'' for the
Lee--Yang theory.
\end{question}

\begin{remark}[Connection to the arithmetic programme]
The above computation confirms and refines the picture from the
Vol~I audit (Synthesis, \S5 of \texttt{prime\_locality/}).
Specifically:
\begin{itemize}[leftmargin=1.8em]
\item The Euler product obstruction for Virasoro
  (Theorem~\textup{euler-koszul-standard}) operates at the
  \emph{sewing-lift level} ($S_{\mathrm{Vir}}(u) = \zeta(u+1)
  (\zeta(u) - 1)$, which has non-multiplicative Dirichlet
  coefficients).  The present computation operates at the
  \emph{shadow-coefficient level} ($S_r(c)$, which is a
  rational function of~$c$).  These are different objects.

\item The weak $p$-adic locality of $S_r(c)$ is consistent with
  the Riccati determinacy theorem: the entire shadow obstruction tower is
  determined by $(\kappa, \alpha, S_4)$ via
  $H(t) = t^2\sqrt{Q_L(t)}$, so the $p$-adic structure is
  controlled by the $p$-adic structure of these three
  invariants.  The weak locality is precisely the statement that
  the square-root recursion preserves $p$-adic classes.

\item Whether the moment $L$-function $M_r(s) = \int
  \operatorname{Sh}_r \cdot E(s)\,d\mu$ has an Euler product
  remains the central open problem.  This requires passing from
  the algebraic level (shadow coefficients, rational functions
  of~$c$) to the analytic level (shadow amplitudes, quasi-modular
  forms on $\cM_{1,1}$).  The $E_2^*(\tau)$ obstruction is real
  and unresolved.
\end{itemize}
\end{remark}


\section{Complementarity constants $K_N$ for the $\cW_N$ family}
\label{sec:wn-complementarity-constants}

\noindent
This section computes the Koszul complementarity constants
\[
K_N \;:=\; \kappa(\cW_{N,c}) + \kappa(\cW_{N,c}^!)
\]
for the principal $\cW_N$-algebra family ($N = 2, 3, 4, \ldots$),
verifying and extending the results of
\texttt{w-algebras-stable.tex} (\S\ref{subsubsec:wn-complementarity}).
The Virasoro case $K_2 = 13$ (AP24) is the seed; the general formula
reveals the harmonic number as the structural invariant governing
the departure from anti-symmetry.

\subsection*{Input data}

The principal $\cW_N$-algebra has generators
$W_s$ of conformal spin $s = 2, 3, \ldots, N$ (with $W_2 = T$,
the stress tensor).  Three ingredients enter:

\begin{enumerate}[label=(\arabic*)]
\item \emph{Ghost central charge per spin.}
  The $bc$-ghost system for a spin-$s$ current has conformal
  weights $(s, 1-s)$ and central charge
  \[
  c_{\mathrm{gh}}(s) = -2(6s^2 - 6s + 1).
  \]
  This is the standard formula; the identity
  $2(6s^2 - 6s + 1) = 3(2s-1)^2 - 1$ provides a convenient check.

\item \emph{Total ghost charge and complementarity constant.}
  \[
  \alpha_N
  = \sum_{s=2}^{N} \bigl|c_{\mathrm{gh}}(s)\bigr|
  = \sum_{s=2}^{N} 2(6s^2 - 6s + 1)
  = 2(N-1)(2N^2 + 2N + 1).
  \]
  The closed form follows from standard sum identities:
  \begin{align*}
  \sum_{s=2}^{N}(6s^2 - 6s + 1)
  &= N(N+1)(2N+1) - 6 - 3N(N+1) + 6 + (N-1) \\
  &= N(N+1)(2N-2) + (N-1) \\
  &= (N-1)\bigl[2N(N+1) + 1\bigr]
  = (N-1)(2N^2 + 2N + 1).
  \end{align*}

\item \emph{Per-generator curvature.}
  The leading scalar self-OPE of the spin-$s$ generator is
  \[
  W_{s\,(2s-1)}\,W_s = \frac{c}{s},
  \]
  where $c$ is the central charge of $\cW_N$.  This is the
  coefficient of the highest-order pole $(z-w)^{-2s}$ in the
  $W_s$--$W_s$ OPE, normalised so that the cyclic pairing
  $\langle W_s, W_s \rangle = 1$.  The formula $c/s$ follows from
  the Bouwknegt--Schoutens normalisation of the principal
  $\cW$-algebra.
\end{enumerate}

\subsection*{The curvature $\kappa(\cW_{N,c})$}

The modular characteristic is the sum of per-generator curvatures:
\begin{equation}\label{eq:kappa-wn-formula}
\kappa(\cW_{N,c})
= \sum_{s=2}^{N} \frac{c}{s}
= c \sum_{s=2}^{N} \frac{1}{s}
= c\,(H_N - 1),
\end{equation}
where $H_N = \sum_{j=1}^{N} 1/j$ is the $N$th harmonic number.

\emph{Verification:}
\begin{itemize}[leftmargin=1.8em]
\item $N=2$: $\kappa(\mathrm{Vir}_c) = c \cdot (3/2 - 1) = c/2$.
  Matches Vol~I.
\item $N=3$: $\kappa(\cW_3) = c \cdot (1/2 + 1/3) = 5c/6$.
  Matches the engine (\texttt{holographic\_ht\_engine.py}).
\item $N=4$: $\kappa(\cW_4) = c \cdot (1/2 + 1/3 + 1/4) = 13c/12$.
\item $N=5$: $\kappa(\cW_5) = c \cdot (1/2 + 1/3 + 1/4 + 1/5) = 77c/60$.
\end{itemize}

\subsection*{The Koszul dual and complementarity sum}

The Koszul dual of $\cW_{N,c}$ is $\cW_{N,\,\alpha_N - c}$.
The Koszul involution is $c \mapsto \alpha_N - c$, with self-dual point
$c^*_N = \alpha_N / 2 = (N-1)(2N^2 + 2N + 1)$.

The curvature of the dual is
\[
\kappa(\cW_{N,c}^!)
= \kappa(\cW_{N,\,\alpha_N - c})
= (\alpha_N - c)(H_N - 1).
\]
Therefore the complementarity sum is
\begin{equation}\label{eq:kn-formula}
\boxed{
K_N
= \kappa(\cW_{N,c}) + \kappa(\cW_{N,c}^!)
= c(H_N - 1) + (\alpha_N - c)(H_N - 1)
= \alpha_N\,(H_N - 1).
}
\end{equation}
This is \emph{independent of~$c$} (as it must be for a
complementarity invariant) but \emph{not zero} for $N \ge 2$.

\subsection*{Explicit values}

\begin{center}
\renewcommand{\arraystretch}{1.3}
\begin{tabular}{cccccc}
\toprule
$N$ & $\alpha_N$ & $H_N - 1$ & $K_N$ & $c^*_N$ & $K_N / c^*_N$ \\
\midrule
$2$ & $26$ & $1/2$ & $13$ & $13$ & $1$ \\
$3$ & $100$ & $5/6$ & $250/3$ & $50$ & $5/3$ \\
$4$ & $246$ & $13/12$ & $533/2$ & $123$ & $13/6$ \\
$5$ & $488$ & $77/60$ & $9394/15$ & $244$ & $77/30$ \\
$6$ & $850$ & $29/20$ & $2465/2$ & $425$ & $29/10$ \\
\midrule
$N$ & $2(N{-}1)(2N^2{+}2N{+}1)$ & $H_N - 1$
  & $\alpha_N(H_N - 1)$ & $(N{-}1)(2N^2{+}2N{+}1)$
  & $2(H_N - 1)$ \\
\bottomrule
\end{tabular}
\end{center}

\subsection*{The harmonic ratio}

The ratio $K_N / c^*_N = 2(H_N - 1)$ reveals the structural content:
the complementarity constant exceeds the self-dual central charge
by the factor $2(H_N - 1)$, which grows logarithmically.
For Virasoro ($N=2$), this factor is exactly~$1$, so $K_2 = c^*_2 = 13$
--- the coincidence that disguises the general pattern.
For $N \ge 3$, $K_N > c^*_N$ strictly.

\begin{observation}[Why $K_N = c^*_N$ only for Virasoro]
The identity $K_N = c^*_N$ requires $2(H_N - 1) = 1$, i.e.\
$H_N = 3/2$.  Since $H_2 = 3/2$ and $H_N$ is strictly
increasing, this holds only for $N = 2$.  The Virasoro algebra
is the \emph{unique} member of the $\cW_N$ family for which the
complementarity sum equals the self-dual central charge.
\end{observation}

\subsection*{The Kac--Moody contrast}

For Kac--Moody algebras $V_k(\fg)$, the Feigin--Frenkel involution
$k \mapsto -k - 2h^\vee$ gives $\kappa + \kappa^! = 0$.
The Koszul involution is \emph{anti-symmetric around the origin}.
For $\cW_N$, the involution $c \mapsto \alpha_N - c$ is
anti-symmetric around $c^*_N = \alpha_N/2$, which is nonzero.
The curvature $\kappa = c(H_N - 1)$ is linear in~$c$, so the
complementarity sum is $\alpha_N(H_N - 1) \ne 0$.

The root cause: Kac--Moody curvature is
$\kappa = \dim(\fg)(k + h^\vee)/(2h^\vee)$, which is linear in~$k$
with the Koszul involution $k \mapsto -k - 2h^\vee$ swapping
$k + h^\vee \mapsto -(k + h^\vee)$, hence
$\kappa \mapsto -\kappa$.  For $\cW_N$, the curvature
$\kappa = c(H_N - 1)$ is linear in~$c$ but the involution
$c \mapsto \alpha_N - c$ does not negate~$c$; it shifts it.
The complementarity sum is the shift:
$K_N = \alpha_N(H_N - 1)$.

\subsection*{Per-generator anatomy}

The complementarity sum decomposes by spin:
\[
K_N
= \sum_{s=2}^{N} \Bigl[\frac{c}{s} + \frac{\alpha_N - c}{s}\Bigr]
= \sum_{s=2}^{N} \frac{\alpha_N}{s}
= \alpha_N \sum_{s=2}^{N} \frac{1}{s}.
\]
The per-generator complementarity constant is
$k_s = \alpha_N / s$.  This depends on~$N$ (not just on~$s$)
because the central charge~$c$ parametrises all generators
simultaneously: changing~$c$ shifts \emph{every} OPE coefficient.

Contrast with the ghost charges: the per-generator ghost central
charge $|c_{\mathrm{gh}}(s)| = 2(6s^2 - 6s + 1)$ depends only
on~$s$, not on~$N$.  The ghost system for each spin is
intrinsic.  The ratio
\[
\frac{k_s}{|c_{\mathrm{gh}}(s)|}
= \frac{\alpha_N}{s \cdot 2(6s^2 - 6s + 1)}
\]
is \emph{not} constant in~$s$: higher-spin generators are
``less efficient'' at converting ghost central charge into
complementarity.

\subsection*{Asymptotics}

As $N \to \infty$:
\begin{align*}
\alpha_N &\sim 4N^3, \\
H_N - 1 &\sim \log N + \gamma - 1
  \quad (\gamma = 0.5772\ldots \text{ Euler--Mascheroni}), \\
K_N &\sim 4N^3(\log N + \gamma - 1), \\
c^*_N &\sim 2N^3.
\end{align*}
The complementarity constant grows as $N^3 \log N$,
while the self-dual central charge grows as $N^3$.
The logarithmic factor is the harmonic series --- the
signature of the sum over spins.

\begin{remark}[Connection to the $\cW_\infty$ limit]
In the $\cW_{1+\infty}$ limit, $H_N - 1 \to \infty$ and
$\alpha_N \to \infty$, so $K_N \to \infty$: the
complementarity constant diverges.  This reflects the
fact that $\cW_{1+\infty}$ has infinitely many generators
and the self-OPE contributions form a divergent harmonic
series.  Any finite truncation $\cW_N$ has finite $K_N$,
but the limit requires renormalisation of the
complementarity sum --- a problem that does not arise for
Kac--Moody algebras (where $K = 0$ identically).
\end{remark}

\begin{remark}[The three invariants of $\cW_N$ Koszul duality]
The Koszul duality of $\cW_{N,c}$ is controlled by three
data:
\begin{enumerate}[label=(\roman*)]
\item The complementarity constant $\alpha_N = 2(N-1)(2N^2+2N+1)$,
  determining the Koszul involution $c \mapsto \alpha_N - c$.
\item The curvature coefficient $H_N - 1 = \sum_{s=2}^N 1/s$,
  determining $\kappa = c(H_N - 1)$.
\item The complementarity sum $K_N = \alpha_N(H_N - 1)$,
  the level-independent invariant.
\end{enumerate}
The three are not independent: $K_N = \alpha_N(H_N - 1)$.
But $\alpha_N$ (a polynomial in~$N$) and $H_N - 1$
(a partial harmonic sum, hence transcendental for $N \ge 2$)
have entirely different analytic characters.  The
complementarity constant mixes polynomial and harmonic data.
\end{remark}


% =================================================================
\section{$SU(2)$ Chern--Simons: explicit $\Ainf$ structure through $m_4$}
\label{sec:su2-ainf}
% =================================================================

We compute the full $\Ainf$ structure of the bar complex
$\barB(\widehat{\fsl}_2{}_k)$ through arity~$4$, for $SU(2)$
Chern--Simons at level~$k$ on $\R_t \times \C_z$.  The boundary
chiral algebra is the affine Kac--Moody algebra
$\widehat{\fsl}_2$ at level~$k$, with OPE
\[
J^a(z)\, J^b(w)
\;\sim\;
\frac{k\,\delta^{ab}}{(z-w)^2}
\;+\;
\frac{f^{ab}{}_c\, J^c(w)}{z-w}.
\]
Poles of order at most~$2$: double from $k\,\delta^{ab}$,
simple from $f^{ab}{}_c\,J^c$.  No quartic pole.

\subsection{Orthonormal basis and conventions}

We work in the orthonormal basis $\{T^1, T^2, T^3\}$ of
$\fsl_2$ with
$[T^a, T^b] = \epsilon^{abc}\, T^c$,
$g^{ab} = \delta^{ab}$,
$\epsilon^{123} = +1$, and $\lambda$-bracket
\begin{equation}\label{eq:su2wn-bracket}
\{T^a{}_\lambda T^b\}
= \epsilon^{abc}\, T^c + \hat{k}\, \delta^{ab}\, \lambda,
\end{equation}
where $\hat{k}$ is the level in this normalization.
Key identity:
$\epsilon^{acd}\,\epsilon^{bde}
= \delta^{ab}\,\delta^{ce} - \delta^{ae}\,\delta^{cb}$.

\subsection{The binary operation $m_2$}

\begin{equation}\label{eq:su2wn-m2}
m_2(T^a, T^b;\, \lambda)
= \epsilon^{abc}\, T^c + \hat{k}\,\delta^{ab}\,\lambda.
\end{equation}
Linear in $\lambda$ (not cubic as for Virasoro).

\subsection{The PVA associator}

\begin{theorem}[PVA associator for $\widehat{\fsl}_2$]
\label{thm:su2wn-associator}
\begin{equation}\label{eq:su2wn-associator}
\boxed{
A_3(T^a, T^b, T^c;\, \lambda_1, \lambda_2)
= 2\delta^{ac}\, T^b - \delta^{ab}\, T^c
  - \delta^{bc}\, T^a + \hat{k}\, \epsilon^{abc}\, \lambda_2.
}
\end{equation}
\end{theorem}

\begin{proof}
\emph{First composition.}
$m_2(T^a, T^b;\, \lambda_1) = \epsilon^{abp}\, T^p
+ \hat{k}\,\delta^{ab}\,\lambda_1$.
Scalar killed by $\{\mathbf{1}_\lambda T^c\} = 0$.
\begin{align*}
\epsilon^{abp}\,m_2(T^p, T^c;\, \lambda_1 + \lambda_2)
&= \epsilon^{abp}\,\epsilon^{pcq}\, T^q
  + \hat{k}\,\epsilon^{abc}\,(\lambda_1 + \lambda_2) \\
&= (\delta^{ac}\,\delta^{bq} - \delta^{aq}\,\delta^{bc})\, T^q
  + \hat{k}\,\epsilon^{abc}\,(\lambda_1 + \lambda_2) \\
&= \delta^{ac}\, T^b - \delta^{bc}\, T^a
  + \hat{k}\,\epsilon^{abc}\,(\lambda_1 + \lambda_2).
\end{align*}

\emph{Second composition.}
$m_2(T^b, T^c;\, \lambda_2) = \epsilon^{bcp}\, T^p
+ \hat{k}\,\delta^{bc}\,\lambda_2$.  Scalar killed.
\begin{align*}
\epsilon^{bcp}\,m_2(T^a, T^p;\, \lambda_1)
&= \epsilon^{bcp}\,\epsilon^{apq}\, T^q
  + \hat{k}\,\epsilon^{abc}\,\lambda_1 \\
&= (\delta^{ab}\,\delta^{cq} - \delta^{ac}\,\delta^{bq})\, T^q
  + \hat{k}\,\epsilon^{abc}\,\lambda_1 \\
&= \delta^{ab}\, T^c - \delta^{ac}\, T^b
  + \hat{k}\,\epsilon^{abc}\,\lambda_1.
\end{align*}

\emph{Difference:}
$A_3 = 2\delta^{ac}\, T^b - \delta^{ab}\, T^c
- \delta^{bc}\, T^a + \hat{k}\,\epsilon^{abc}\,\lambda_2$.
\end{proof}

\subsection{$\Ainf$-formality: $m_k^H = 0$ for $k \geq 3$ on cohomology}

\begin{theorem}[$\Ainf$-formality for $\widehat{\fsl}_2$]
\label{thm:su2wn-formality}
The transferred $\Ainf$ structure on the PVA cohomology
$H^\bullet(\cA_{\mathrm{CS}}, Q) \simeq \widehat{\fsl}_{2,k}$
is formal:
$m_k^H = 0$ for all $k \geq 3$.
The nonzero associator~\eqref{eq:su2wn-associator} is compensated
by the normally ordered product via the HPL transfer.
\end{theorem}

\begin{proof}
Two ingredients.

\medskip\noindent
\textbf{(1) Lie conformal algebra structure.}
The $\lambda$-bracket~\eqref{eq:su2wn-bracket} satisfies the
Lie conformal algebra Jacobi identity:
$\{T^a{}_{\lambda_1}\{T^b{}_{\lambda_2} T^c\}\}
- \{T^b{}_{\lambda_2}\{T^a{}_{\lambda_1} T^c\}\}
= \{\{T^a{}_{\lambda_1} T^b\}{}_{\lambda_1+\lambda_2} T^c\}$.
For a Lie conformal algebra, the enveloping vertex algebra is
well-defined (PBW theorem), and the $\Ainf$ structure on the
PVA is formal: the Lie conformal Jacobi identity encodes all
homotopy-coherence data, so no higher operations are needed.

\medskip\noindent
\textbf{(2) HPL transfer.}
The vertex algebra $V_k(\fsl_2)$ is associative (Borcherds).
The HPL formula for $m_3^H$ is:
$m_3^H(a,b,c) = p\, m_2^{\mathrm{full}}
(h\, m_2^{\mathrm{full}}(i(a), i(b)),\, i(c))
+ p\, m_2^{\mathrm{full}}
(i(a),\, h\, m_2^{\mathrm{full}}(i(b), i(c)))$.

On weight-$1$ generators: $m_2^{\mathrm{sing}}(T^a, T^b)$
produces weight-$1$ cocycles and scalars (both in $\ker h$),
so $h \circ m_2^{\mathrm{sing}} = 0$.  The normally ordered
product $\normord{T^a T^b}$ is weight-$2$; $h(\normord{T^a T^b})$
is a weight-$1$ acyclic state.  Feeding this into either
$m_2^{\mathrm{sing}}$ or $m_2^{\mathrm{reg}}$ and projecting
by $p$ to weight-$1$ cohomology yields zero (the state is
$Q$-exact).

Therefore $m_3^H = 0$.  Induction extends to all $m_k^H$:
every HPL tree for $m_k^H$ contains at least one $h$-insertion
mapping into the acyclic part, killed by $p$.
\end{proof}

\begin{remark}[Reconciling the nonzero associator with formality]
The associator $A_3 \neq 0$ and formality $m_3^H = 0$ are
compatible.  The Stasheff equation $A_3 = d\, m_3$ is solved
at chain level by a nonzero $m_3$ that is $Q$-exact.
On cohomology, $p(A_3) = 0$ because $A_3$ is in the image
of $d$: it is a coboundary, not a cocycle.

For Virasoro, the quartic-pole term
$\tfrac{c}{12}\lambda_2^3(2\lambda_1+\lambda_2)$ in the
associator is a nontrivial cocycle (it cannot be written as
$d$ of anything), so $m_3^H \neq 0$.  For KM, the associator
is entirely built from Lie-bracket terms that ARE coboundaries.
\end{remark}

\subsection{Chain-level operations from Feynman diagrams}

While $m_k^H = 0$ for $k \geq 3$ on cohomology, the
chain-level $\Ainf$ operations are nonzero:

\begin{proposition}[Chain-level $\Ainf$ structure]
\label{prop:su2wn-chain}
\leavevmode
\begin{enumerate}
\item \textbf{$m_3$} (triangle diagram): two cubic vertices,
  color $f^{a}{}_{de}\, f^{ebc}$.  Nonzero at chain level;
  $Q$-exact by Stokes on $\FM_3(\C)$.

\item \textbf{$m_4$} (square diagram): three cubic vertices,
  color $f^{ab}{}_g\, f^{gc}{}_h\, f^{hd}{}_e$.  Nonzero
  at chain level; $Q$-exact by AOS on $\FM_4(\C)$.

\item \textbf{$m_{k \geq 5} = 0$}: tree with $k-2$ vertices
  has degree exceeding the $\FM_k(\C)$ integration domain
  (Arnold--Orlik--Solomon relations).
\end{enumerate}
\end{proposition}

\subsection{Koszul dual and complementarity}

\begin{theorem}[Koszul dual of $\widehat{\fsl}_{2,k}$]
\label{thm:su2wn-koszul}
Koszul dual level: $k' = -k - 2h^\vee = -k - 4$
(with $h^\vee = 2$).  Curvature:
$\kappa = 3(k+2)/4$.  Complementarity:
$\kappa + \kappa' = 0$ (not $13$ as for Virasoro).
\end{theorem}

\begin{proof}
$\kappa' = 3(-k-4+2)/4 = -3(k+2)/4 = -\kappa$.
The $2h^\vee$ shift: Verdier duality sends
$k + h^\vee \mapsto -(k + h^\vee)$, so
$k' = -k - 2h^\vee$.
\end{proof}

\begin{remark}[AP33: three operations, one output]
\begin{center}
\begin{tabular}{lll}
\hline
\textbf{Operation} & \textbf{Definition} & \textbf{Level} \\
\hline
Koszul & $(H^\bullet(\barB(\cA)))^\vee$ & $-k-4$ \\
FF involution & $k \mapsto -k - 2h^\vee$ & $-k-4$ \\
Negative level & $V_k \mapsto V_{-k}$ & $-k$ \\
\hline
\end{tabular}
\end{center}
Koszul and FF agree on the Koszul locus.
$\widehat{\fsl}_{2,k}^!$ as a PVA is
$\widehat{\fsl}_{2,-k-4}$-PVA, but as a vertex algebra
is $\operatorname{Sym}^{\mathrm{ch}}((\fsl_2)^*[-1])$,
not $V_{-k-4}(\fsl_2)$ (AP33).
\end{remark}

\subsection{Summary}

\begin{theorem}[Complete $\Ainf$ structure of $\widehat{\fsl}_{2,k}$ through $m_4$]
\label{thm:su2wn-complete}
\leavevmode
\begin{center}
\begin{tabular}{ll}
\hline
\textbf{Datum} & \textbf{Value} \\
\hline
$m_2$ & $\epsilon^{abc}\,T^c + \hat{k}\,\delta^{ab}\,\lambda$ \\
$m_k^H$ (cohomology, $k \geq 3$) & $0$ ($\Ainf$-formal) \\
$m_3, m_4$ (chain level) & $\neq 0$ ($Q$-exact) \\
$m_{k \geq 5}$ (all levels) & $0$ \\
PVA associator & $2\delta^{ac}T^b - \delta^{ab}T^c
  - \delta^{bc}T^a + \hat{k}\epsilon^{abc}\lambda_2$ \\
Shadow depth & $3$ (class $\mathbf{L}$) \\
Koszul dual & $k' = -k - 4$ \\
Curvature & $3(k+2)/4$ \\
$K = \kappa + \kappa'$ & $0$ \\
$r$-matrix & $k\,\Omega/z$ \quad\textup{(trace convention; at $k=0$ the Yangian becomes abelian and $r=0$, AP126 satisfied)} \\
Lines & $Y(\fsl_2)$-$\mathbf{mod}$ \\
\hline
\end{tabular}
\end{center}
\end{theorem}

\begin{remark}[Landscape comparison]
\begin{center}
\begin{tabular}{lccccl}
\hline
\textbf{Family} & $d_{\max}$ & $m_3^H$ & $m_4^H$
  & chain $m_3$ & \textbf{Class} \\
\hline
Heisenberg & $2$ & $0$ & $0$ & $0$ & $\mathbf{G}$, depth $2$ \\
KM $\widehat{\fg}_k$ & $2$ & $0$ & $0$ & $\neq 0$
  & $\mathbf{L}$, depth $3$ \\
$\beta\gamma$ & $1$ & $\neq 0$ & $\neq 0$ & $\neq 0$
  & $\mathbf{C}$, depth $4$ \\
Virasoro & $4$ & $\neq 0$ & $\neq 0$ & $\neq 0$
  & $\mathbf{M}$, depth $\infty$ \\
\hline
\end{tabular}
\end{center}
The key answer: $m_4 \neq 0$ at chain level but $m_4^H = 0$ on
cohomology.  The $\Ainf$ tower terminates on cohomology at
arity~$2$ (strict formality).  At chain level, it terminates at
arity~$4$ (AOS degree for $k \geq 5$).  The mechanism is the Lie
algebra Jacobi identity: it ensures chain-level operations are
$Q$-exact.  $SU(2)$ CS is ``$\Ainf$-trivial'': the factorization
structure is entirely captured by the $\lambda$-bracket $m_2$.
\end{remark}



% =================================================================
\section{$bc$--$\beta\gamma$ Koszul duality: free-field test of the derived conjecture}
\label{sec:bc-betagamma-koszul}
% =================================================================

The $bc$ ghost system and the $\beta\gamma$ system are the
two-generator free-field atoms.  Their Koszul duality is the
chiral incarnation of classical $\Lambda^* \leftrightarrow
\mathrm{Sym}^*$ duality.  This section computes the Koszul
dual of each system at general conformal weights, then
of the combined $bc \otimes \beta\gamma$ system, and finally
specialises to the string-theory point $(\lambda, \mu) =
(2, 3/2)$.  The computation tests
Conjecture~\ref*{conj:derived-bc-betagamma} of Volume~I
(the derived $bc$--$\beta\gamma$ Koszul duality conjecture).

\subsection{Conventions and OPE data}

We parametrise the $bc$ system by a single conformal weight
parameter $\lambda \in \C$:
\begin{itemize}
\item $b(z)$: fermionic, conformal weight $h_b = \lambda$,
\item $c(z)$: fermionic, conformal weight $h_c = 1 - \lambda$.
\end{itemize}
The OPE is
\begin{equation}\label{eq:bc-ope-general}
b(z)\,c(w) \;=\; \frac{1}{z-w} + \text{reg.},
\qquad
c(z)\,b(w) \;=\; \frac{1}{z-w} + \text{reg.}
\end{equation}
All other OPEs (including $b \cdot b$ and $c \cdot c$) are regular,
by fermionic statistics.

We parametrise the $\beta\gamma$ system by a conformal weight
parameter $\mu \in \C$:
\begin{itemize}
\item $\beta(z)$: bosonic, conformal weight $h_\beta = \mu$,
\item $\gamma(z)$: bosonic, conformal weight $h_\gamma = 1 - \mu$.
\end{itemize}
The OPE is
\begin{equation}\label{eq:bg-ope-general}
\beta(z)\,\gamma(w) \;=\; \frac{1}{z-w} + \text{reg.},
\qquad
\gamma(z)\,\beta(w) \;=\; -\frac{1}{z-w} + \text{reg.}
\end{equation}

\begin{remark}[Central charges]
\label{rem:central-charges-bc-bg}
The Virasoro central charges are:
\begin{align}
c_{bc}(\lambda)
  &= -2(6\lambda^2 - 6\lambda + 1)
  = 1 - 12\lambda(\lambda - 1), \label{eq:c-bc-general} \\
c_{\beta\gamma}(\mu)
  &= 2(6\mu^2 - 6\mu + 1). \label{eq:c-bg-general}
\end{align}
Modular characteristics: $\kappa(bc) = c_{bc}/2 = -(6\lambda^2
- 6\lambda + 1)$ and $\kappa(\beta\gamma) = c_{\beta\gamma}/2
= 6\mu^2 - 6\mu + 1$.

\emph{Critical observation:} at matching conformal weights
$\mu = \lambda$, the central charges satisfy
\begin{equation}\label{eq:c-complementarity}
c_{bc}(\lambda) + c_{\beta\gamma}(\lambda) = 0,
\qquad
\kappa(bc_\lambda) + \kappa(\beta\gamma_\lambda) = 0.
\end{equation}
This is exact complementarity (Theorem~C) for the free-field
Koszul pair: no shift by $\rho \cdot K$ as in the W-algebra
case.
\end{remark}


\subsection{Koszul dual of the $bc$ system}

\begin{theorem}[$bc$ Koszul duality at general spins]
\label{thm:bc-koszul-general-spin}
For any $\lambda \in \C$, the $bc$ system with spins
$(\lambda, 1{-}\lambda)$ is quadratically chiral Koszul,
and its Koszul dual is the $\beta\gamma$ system at
the \emph{same} conformal weights:
\begin{equation}\label{eq:bc-koszul-dual}
(bc_\lambda)^!
\;\cong\;
\beta\gamma_\lambda.
\end{equation}
In particular, $h_{b^*} = h_\beta = \lambda$ and
$h_{c^*} = h_\gamma = 1 - \lambda$: Koszul duality
preserves the conformal weights and flips the statistics
(fermionic $\leftrightarrow$ bosonic).
\end{theorem}

\begin{proof}
The $bc$ system is quadratic: the generating space is
$V = \C b \oplus \C c$ (fermionic, $\dim V = 2$) with
the single quadratic relation
\[
R_{bc} = \mathrm{span}\{b \otimes c + c \otimes b\}
\;\subset\; \mathrm{Sym}^2(V).
\]
This is a \emph{symmetric} (Clifford-type) relation: the
$+$ sign reflects the fermionic anticommutation $b(z_1)c(z_2)
+ c(z_2)b(z_1) \sim \frac{2}{z_1 - z_2}$ in the OPE, which
the residue extraction $d\log(z_1 - z_2)$ converts to the
bar differential identity $d_{\barB}[s^{-1}b \,|\, s^{-1}c]
+ d_{\barB}[s^{-1}c \,|\, s^{-1}b] = 0$ in bar degree~$2$.

\emph{Why $m_k = 0$ for $k \ge 3$.}  The OPE \eqref{eq:bc-ope-general}
has only a simple pole $1/(z{-}w)$.  The bar differential component
$d_k\colon \barB^k \to \barB^{k-1}$ at arity $k$ involves the
residue of the OPE integrated against a $(k{-}1)$-form on
$\FM_k(\C)$.  For $k \ge 3$, the form degree exceeds the pole
order: the collision of two points among $k$ produces at most a
simple pole, while the form $\eta_{i_1 i_2} \wedge \cdots \wedge
\eta_{i_{k-1} i_k}$ contributes $k{-}1$ logarithmic factors.
By dimensional analysis on the FM compactification, the residue
integral vanishes for $k \ge 3$, giving $m_k = 0$.  This is the
same mechanism as for the Heisenberg algebra (Rosetta Stone,
\S\ref{sec:rosetta-stone}): quadratic OPE $\Rightarrow$
quadratic Koszulness.

\emph{Koszul dual.}
Under the residue pairing $\langle b, \beta \rangle = 1$,
$\langle c, \gamma \rangle = 1$ (where $\beta = b^*$ and
$\gamma = c^*$ denote the dual generators, which are bosonic
by the parity flip), the orthogonal complement of $R_{bc}$ is
\[
R_{bc}^\perp = \mathrm{span}\{\beta \otimes \gamma - \gamma
\otimes \beta\} \;\subset\; \Lambda^2(V^*).
\]
This is exactly the \emph{antisymmetric} (Weyl-type) relation
of the $\beta\gamma$ system.  Since the OPE of $bc$ is
conformally covariant with weights $(\lambda, 1{-}\lambda)$,
the dual generators inherit the same conformal weights: $h_\beta
= h_b = \lambda$, $h_\gamma = h_c = 1 - \lambda$.  Therefore
\[
(bc_\lambda)^!
= \mathrm{Free}^{\mathrm{ch}}(V^*)/(R_{bc}^\perp)
= \beta\gamma_\lambda.
\qedhere
\]
\end{proof}

\begin{remark}[Central charge under Koszul duality]
\label{rem:bc-kappa-koszul}
Koszul duality does \emph{not} negate the central charge
as a scalar; it negates the \emph{modular characteristic}.
At matching spins:
\[
\kappa((bc_\lambda)^!) = \kappa(\beta\gamma_\lambda)
= -\kappa(bc_\lambda),
\]
so $\kappa + \kappa^! = 0$ (exact complementarity, free-field
family).  Numerically:
\[
c_{(bc_\lambda)^!} = c_{\beta\gamma_\lambda}
= -c_{bc_\lambda}.
\]
The central charge \emph{negates} under Koszul duality for
this pair.  This is a coincidence of the free-field family
and should not be confused with the Kac--Moody level shift
$k \mapsto -k - 2h^\vee$ (AP33).
\end{remark}


\subsection{Koszul dual of the $\beta\gamma$ system}

\begin{theorem}[$\beta\gamma$ Koszul duality at general spins]
\label{thm:bg-koszul-general-spin}
For any $\mu \in \C$, the $\beta\gamma$ system with spins
$(\mu, 1{-}\mu)$ is quadratically chiral Koszul, and
\begin{equation}\label{eq:bg-koszul-dual}
(\beta\gamma_\mu)^!
\;\cong\;
bc_\mu.
\end{equation}
\end{theorem}

\begin{proof}
The argument is dual to Theorem~\ref{thm:bc-koszul-general-spin}.
Generators: $V = \C\beta \oplus \C\gamma$ (bosonic, $\dim V = 2$).
Relation: $R_{\beta\gamma} = \mathrm{span}\{\beta \otimes \gamma
- \gamma \otimes \beta\} \subset \Lambda^2(V)$ (antisymmetric,
Weyl type).  The orthogonal complement under the residue pairing
is $R_{\beta\gamma}^\perp = \mathrm{span}\{b \otimes c + c
\otimes b\} \subset \mathrm{Sym}^2(V^*)$, which is the Clifford
relation of the $bc$ system.  The vanishing $m_k = 0$ for
$k \ge 3$ follows from the same simple-pole argument.
\end{proof}

\begin{observation}[Involutivity]
\label{obs:koszul-involution}
The Koszul duality is an involution on the free-field pair:
\[
((bc_\lambda)^!)^! = (\beta\gamma_\lambda)^! = bc_\lambda,
\qquad
((\beta\gamma_\mu)^!)^! = (bc_\mu)^! = \beta\gamma_\mu.
\]
This reflects the general fact that Koszul duality is an
involution on the quadratic Koszul locus
(Vol~I, Theorem~B).
\end{observation}


\subsection{Koszul dual of the combined $bc \otimes \beta\gamma$ system}

We now address the tensor product.  The combined system has
generators $\{b, c, \beta, \gamma\}$ with the $bc$ OPE
\eqref{eq:bc-ope-general} and the $\beta\gamma$ OPE
\eqref{eq:bg-ope-general}, and \emph{no cross-OPEs}:
\begin{equation}\label{eq:no-cross-ope}
b(z)\,\beta(w) \sim 0,\quad
b(z)\,\gamma(w) \sim 0,\quad
c(z)\,\beta(w) \sim 0,\quad
c(z)\,\gamma(w) \sim 0.
\end{equation}

\begin{theorem}[Koszul duality distributes over the chiral tensor product]
\label{thm:koszul-tensor-bc-bg}
Let $\lambda, \mu \in \C$ be arbitrary.  Then the chiral tensor
product $bc_\lambda \otimes^{\mathrm{ch}} \beta\gamma_\mu$ is
quadratically chiral Koszul, and
\begin{equation}\label{eq:koszul-tensor}
(bc_\lambda \otimes^{\mathrm{ch}} \beta\gamma_\mu)^!
\;\cong\;
\beta\gamma_\lambda \otimes^{\mathrm{ch}} bc_\mu.
\end{equation}
In particular, Koszul duality acts component-wise on the tensor
factors, exchanging statistics in each sector while preserving
conformal weights.
\end{theorem}

\begin{proof}
Both $bc_\lambda$ and $\beta\gamma_\mu$ are quadratic chiral
Koszul (Theorems~\ref{thm:bc-koszul-general-spin}
and~\ref{thm:bg-koszul-general-spin}).  The vanishing of
cross-OPEs \eqref{eq:no-cross-ope} means the combined system
is a genuine chiral tensor product in the sense of
Vol~I, Proposition~\ref*{prop:koszul-dual-tensor-product}: the
generators decompose as $V = V_{bc} \oplus V_{\beta\gamma}$
and the relation space decomposes as $R = R_{bc} \oplus
R_{\beta\gamma}$ (no mixed relations).

By the tensor product formula for quadratic chiral Koszul
pairs (loc.\ cit.):
\begin{align*}
(bc_\lambda \otimes^{\mathrm{ch}} \beta\gamma_\mu)^!
&= \mathrm{Free}^{\mathrm{ch}}(V_{bc}^* \oplus V_{\beta\gamma}^*)
   /(R_{bc}^\perp \oplus R_{\beta\gamma}^\perp) \\
&= (bc_\lambda)^! \otimes^{\mathrm{ch}} (\beta\gamma_\mu)^! \\
&= \beta\gamma_\lambda \otimes^{\mathrm{ch}} bc_\mu.
\end{align*}

The central charge of the dual is
\[
c_{(bc_\lambda \otimes \beta\gamma_\mu)^!}
= c_{\beta\gamma_\lambda} + c_{bc_\mu}
= -c_{bc_\lambda} - c_{\beta\gamma_\mu}
= -c_{bc_\lambda \otimes \beta\gamma_\mu},
\]
and complementarity holds:
$\kappa_{\mathrm{total}} + \kappa_{\mathrm{total}}^! = 0$.
\end{proof}

\begin{remark}[When $\lambda = \mu$: self-Koszul-dual tensor product]
\label{rem:self-dual-tensor}
When the $bc$ and $\beta\gamma$ systems have the \emph{same}
conformal weight parameter $\lambda = \mu$, the dual
$\beta\gamma_\lambda \otimes bc_\lambda$ is the \emph{same}
algebra as $bc_\lambda \otimes \beta\gamma_\lambda$ (up to
reordering of tensor factors).  Thus the combined system is
\emph{self-Koszul-dual} at matching spins:
\[
(bc_\lambda \otimes^{\mathrm{ch}} \beta\gamma_\lambda)^!
\;\cong\;
bc_\lambda \otimes^{\mathrm{ch}} \beta\gamma_\lambda.
\]
This self-duality has central charge $c_{\mathrm{total}} = 0$,
$\kappa_{\mathrm{total}} = 0$, reflecting exact boson--fermion
cancellation.  The bar complex of the combined system is
acyclic (the Koszul complex is self-dual and exact), which is
the algebraic shadow of the statement that the combined
free-field system has vanishing conformal anomaly.
\end{remark}


\subsection{The string-theory point: $\lambda = 2$, $\mu = 3/2$}

At the string-theory BRST point, the $bc$ ghosts have spins
$(2, -1)$ (i.e., $\lambda = 2$) and the $\beta\gamma$ superghosts
have spins $(3/2, -1/2)$ (i.e., $\mu = 3/2$).

\begin{proposition}[String-theory central charges]
\label{comp:string-cc}
\begin{align*}
c_{bc}(2) &= -2(6 \cdot 4 - 6 \cdot 2 + 1) = -2 \cdot 13 = -26, \\
c_{\beta\gamma}(3/2) &= 2(6 \cdot 9/4 - 6 \cdot 3/2 + 1)
  = 2(27/2 - 9 + 1) = 2 \cdot 11/2 = 11.
\end{align*}
Total: $c_{\mathrm{total}} = -26 + 11 = -15$.  Modular characteristics:
$\kappa(bc) = -13$, $\kappa(\beta\gamma) = 11/2$,
$\kappa_{\mathrm{total}} = -15/2$.
\end{proposition}

\begin{theorem}[Koszul dual of the string ghost--superghost system]
\label{thm:koszul-string-ghosts}
The Koszul dual of the combined $bc$--$\beta\gamma$ ghost system
at the string-theory point $(\lambda, \mu) = (2, 3/2)$ is:
\begin{equation}\label{eq:koszul-string-dual}
(bc_2 \otimes^{\mathrm{ch}} \beta\gamma_{3/2})^!
\;\cong\;
\beta\gamma_2 \otimes^{\mathrm{ch}} bc_{3/2}.
\end{equation}
The Koszul dual is a bosonic $\beta\gamma$ system at spins
$(2, -1)$ tensored with a fermionic $bc$ system at spins
$(3/2, -1/2)$.  Its central charge is
\begin{align*}
c_{\mathrm{dual}} &= c_{\beta\gamma}(2) + c_{bc}(3/2)
= 2 \cdot 13 + (-2)(6 \cdot 9/4 - 6 \cdot 3/2 + 1) \\
&= 26 + (-2)(11/2) = 26 - 11 = 15.
\end{align*}
This confirms $c_{\mathrm{dual}} = -c_{\mathrm{total}} = 15$
and $\kappa_{\mathrm{dual}} = 15/2 = -\kappa_{\mathrm{total}}$.
\end{theorem}

\begin{proof}
Immediate from Theorem~\ref{thm:koszul-tensor-bc-bg} at
$(\lambda, \mu) = (2, 3/2)$.  The central charge computation is
a direct substitution into \eqref{eq:c-bc-general} and
\eqref{eq:c-bg-general}.
\end{proof}

\begin{remark}[Physical interpretation]
\label{rem:string-koszul-physics}
\mbox{}
\begin{enumerate}
\item \emph{The Koszul dual is NOT the original system.}
  Since $\lambda \neq \mu$, the Koszul dual
  $\beta\gamma_2 \otimes bc_{3/2}$ is a \emph{different}
  algebra from $bc_2 \otimes \beta\gamma_{3/2}$: the
  statistics assignment is swapped at \emph{different}
  conformal weights.  Concretely, the original system has a
  weight-$2$ fermion $b$ and a weight-$3/2$ boson $\beta$;
  the dual has a weight-$2$ boson $\beta^*$ and a weight-$3/2$
  fermion $b^*$.

\item \emph{The central charge flips sign but is nonzero.}
  The asymmetry $\lambda \neq \mu$ means
  $c_{\mathrm{total}} = -15 \neq 0$, so the combined system
  is \emph{not} self-Koszul-dual.  The anomaly $c = -15$
  is precisely the contribution that must be cancelled by
  $c = 15$ worth of matter in the $N = 1$ superstring
  (the superconformal matter has $c_{\mathrm{matter}} = 15$
  in the critical dimension $d = 10$).

\item \emph{AP33 verification.}  The Koszul dual central
  charge is $c_{\mathrm{dual}} = 15 = -c_{\mathrm{total}}$.
  This is \emph{not} obtained by any level-negation
  prescription (there is no ``level'' for $bc$ or
  $\beta\gamma$ in the Kac--Moody sense).  Koszul duality
  acts by statistics exchange $+$ sign flip of $\kappa$, not
  by any Feigin--Frenkel-type involution.

\item \emph{Complementarity and the superstring.}
  The full superstring BRST complex has
  $c_{\mathrm{matter}} + c_{\mathrm{ghosts}} =
  15 + (-15) = 0$, and the Koszul dual of the ghost sector
  has $c = 15 = c_{\mathrm{matter}}$.  This is the free-field
  manifestation of complementarity (Theorem~C): the Koszul
  dual of the ghost sector has the same central charge as
  the matter sector.  The vanishing of the total anomaly is
  equivalent to the statement that the matter and ghost
  sectors form a Koszul-complementary pair.
\end{enumerate}
\end{remark}


\subsection{Bar complex structure: explicit generators}

\begin{proposition}[Bar complex of $bc_\lambda$ at genus~$0$]
\label{prop:bar-bc-genus0}
The geometric bar complex $\barB^{\mathrm{ch}}(bc_\lambda)$ at
genus~$0$ has:
\begin{enumerate}
\item $\barB^0 = \C$ (vacuum).
\item $\barB^1 = \C \cdot [s^{-1}b] \oplus \C \cdot [s^{-1}c]$
  ($\dim = 2$).
\item The bar differential $d_{\barB}\colon \barB^2 \to \barB^1$
  is determined by
  \[
  d_{\barB}[s^{-1}b \,|\, s^{-1}c]
  = \Res_{z_1 = z_2}
    \bigl[b(z_1)\,c(z_2)\,\eta_{12}\bigr]
  = 1 \cdot [s^{-1}\mathbf{1}],
  \]
  where $\mathbf{1}$ is the vacuum and the residue extracts the
  simple pole.  The symmetric relation
  $d_{\barB}[s^{-1}b \,|\, s^{-1}c]
  + d_{\barB}[s^{-1}c \,|\, s^{-1}b] = 0$
  encodes the Clifford relation.
\item $\barB^k = 0$ in bar cohomology for $k \ge 3$
  (simple pole $\Rightarrow$ $m_k = 0$ for $k \ge 3$).
\end{enumerate}
The bar cohomology is
$H^*(\barB^{\mathrm{ch}}(bc_\lambda))
\cong \mathrm{Sym}^{\mathrm{ch}}(\beta, \gamma)$
(the symmetric chiral algebra on two bosonic generators),
confirming $(bc_\lambda)^! = \beta\gamma_\lambda$.
\end{proposition}

\begin{proof}
The only OPE with a pole is $b(z)c(w) \sim 1/(z{-}w)$
(simple pole of order~$1$).  By AP19 (the bar kernel absorbs
a pole), the collision residue is $r(z) = 1$ (constant),
with pole order \emph{one less} than the OPE: the
$d\log(z{-}w)$ kernel absorbs the single power, leaving a
finite residue.

At bar degree $k \ge 3$: a bar element
$[s^{-1}a_1 \,|\, \cdots \,|\, s^{-1}a_k]$ lives on
$\FM_k(\C)$ with $(k{-}1)$ logarithmic forms
$\eta_{i_1 i_2} \wedge \cdots$.  Each pairwise collision
$z_i \to z_j$ produces at most a simple pole from the OPE,
which is exactly absorbed by the $d\log$ propagator.  The
residue contributes to $d_{\barB}$ at bar degree~$k{-}1$,
collapsing two of the $k$ points.  But for the \emph{ternary}
and higher operations, one needs a \emph{simultaneous}
multi-body collision on $\FM_k(\C)$, not just pairwise
collisions.  Since the OPE is strictly binary (no higher
poles, no composite fields at coincident points), these
multi-body contributions vanish.

Equivalently: the $bc$ system is generated by $V = \C b
\oplus \C c$ with a single quadratic relation.  By the PBW
criterion (Vol~I, Theorem~\ref*{thm:pbw-koszulness-criterion}),
the associated graded $\gr_F(bc)$ is a Clifford algebra
$\mathrm{Cl}(V, q)$, which is classically Koszul (Priddy).
The PBW spectral sequence degenerates at $E_1$, confirming
chiral Koszulness.
\end{proof}

\begin{proposition}[Bar complex of $\beta\gamma_\mu$ at genus~$0$]
\label{prop:bar-bg-genus0}
The geometric bar complex $\barB^{\mathrm{ch}}(\beta\gamma_\mu)$
at genus~$0$ has:
\begin{enumerate}
\item $\barB^1 = \C \cdot [s^{-1}\beta] \oplus \C \cdot
  [s^{-1}\gamma]$ ($\dim = 2$).
\item The bar differential encodes the antisymmetric relation:
  \[
  d_{\barB}[s^{-1}\beta \,|\, s^{-1}\gamma]
  - d_{\barB}[s^{-1}\gamma \,|\, s^{-1}\beta] = 0.
  \]
\item $m_k = 0$ for $k \ge 3$ (same simple-pole argument).
\end{enumerate}
The bar cohomology is
$H^*(\barB^{\mathrm{ch}}(\beta\gamma_\mu))
\cong \Lambda^{\mathrm{ch}}(b, c)$
(the exterior chiral algebra on two fermionic generators),
confirming $(\beta\gamma_\mu)^! = bc_\mu$.
\end{proposition}

\begin{proof}
Identical to Proposition~\ref{prop:bar-bc-genus0} with
Clifford $\leftrightarrow$ Weyl exchange.  The bosonic
statistics of $\beta, \gamma$ produce an antisymmetric
relation $R_{\beta\gamma} \subset \Lambda^2(V)$, whose
orthogonal complement is the symmetric relation
$R_{bc} \subset \mathrm{Sym}^2(V^*)$.
\end{proof}


\subsection{Genus-$1$ curvature}

At genus $g = 1$, the bar complex acquires curvature
$\kappa \cdot \omega_1$ (Vol~I, Theorem~D).  For the $bc$
system:
\begin{equation}\label{eq:bc-curvature-g1}
d_{\barB}^2 = \kappa(bc_\lambda) \cdot \omega_1
= -(6\lambda^2 - 6\lambda + 1) \cdot \omega_1,
\end{equation}
where $\omega_1 = c_1(\omega_{\overline{\cM}_1})$ is the first
Chern class of the Hodge bundle on $\overline{\cM}_1$.  For
the Koszul dual:
\[
d_{\barB}^2 = \kappa(\beta\gamma_\lambda) \cdot \omega_1
= (6\lambda^2 - 6\lambda + 1) \cdot \omega_1
= -\kappa(bc_\lambda) \cdot \omega_1.
\]
Curvature \emph{negates} under Koszul duality, as required by
complementarity.

At the string-theory point: $\kappa(bc_2) = -13$, so
$d_{\barB}^2 = -13\omega_1$.  The Koszul dual has
$\kappa(\beta\gamma_2) = +13$, giving $d_{\barB}^2
= +13\omega_1$.  The sum vanishes:
$\kappa + \kappa^! = 0$.  (Compare: for Virasoro at $c$,
the self-duality point is $c = 13$ where
$\kappa(\mathrm{Vir}_c) = c/2 = 13/2$ and
$\kappa(\mathrm{Vir}_{26-c}) = (26{-}c)/2 = 13/2$,
giving $\kappa + \kappa^! = 13 \neq 0$.
The free-field pair has exact cancellation; the
Virasoro pair does not.  This is AP24.)


\subsection{Relation to the derived conjecture}

Conjecture~\ref*{conj:derived-bc-betagamma} (Vol~I) posits
a \emph{derived} Koszul duality between the free fermion
$\cF = \Lambda^{\mathrm{ch}}(\psi^{(1)}, \psi^{(2)})$
(two fermionic generators) and the DG chiral algebra
$(\beta\gamma \oplus bc, d_{\mathrm{BRST}})$.

\begin{observation}[What the computation confirms and what remains open]
\label{obs:derived-conjecture-status}
\mbox{}
\begin{enumerate}
\item \emph{Confirmed:} The individual Koszul dualities
  $(\beta\gamma)^! \cong bc$ and $(bc)^! \cong \beta\gamma$
  hold at all conformal weights (Theorems~\ref{thm:bc-koszul-general-spin}
  and~\ref{thm:bg-koszul-general-spin}).
  These are the ``single-pair dualities'' referenced in the
  conjecture.

\item \emph{Confirmed:} Koszul duality distributes over the
  chiral tensor product for the decoupled $bc \otimes
  \beta\gamma$ system
  (Theorem~\ref{thm:koszul-tensor-bc-bg}).  The distribution
  formula \eqref{eq:koszul-tensor} is a consequence of the
  quadratic tensor product theorem
  (Vol~I, Proposition~\ref*{prop:koszul-dual-tensor-product}).

\item \emph{Open:} The conjecture involves a \emph{BRST
  differential} $d_{\mathrm{BRST}}\colon bc \to \beta\gamma$
  that \emph{couples} the two sectors.  The decoupled tensor
  product (no cross-OPEs) is \emph{not} the same as the BRST
  complex.  The derived Koszul duality framework required to
  handle the coupled system does not yet exist: one needs a
  bar-cobar adjunction for DG chiral algebras (not just plain
  chiral algebras), extending the quadratic framework.

\item \emph{Open:} The conjecture identifies the derived
  Koszul dual of $\cF(\psi^{(1)}, \psi^{(2)})$ (two free
  fermions) with the BRST complex
  $(\beta\gamma \oplus bc, d_{\mathrm{BRST}})$.  Our
  computation addresses the \emph{reverse} direction: the
  Koszul dual of the combined (decoupled) $bc \otimes
  \beta\gamma$ system.  The two directions are related but
  not identical, because the BRST differential introduces
  nontrivial higher $\Ainf$ operations that break the tensor
  product structure.

\item \emph{New evidence:} At matching spins $\lambda = \mu$,
  the combined system is self-Koszul-dual with $c = 0$
  (Remark~\ref{rem:self-dual-tensor}).  This is consistent
  with the conjecture: when the BRST differential is trivial
  (no coupling), the derived pair degenerates to the tensor
  product of individual Koszul pairs, and self-duality is the
  algebraic expression of anomaly cancellation.
\end{enumerate}
\end{observation}

\subsection{Summary table}

\begin{center}
\renewcommand{\arraystretch}{1.3}
\begin{tabular}{lccccc}
\hline
\textbf{System} & \textbf{Spins} & $c$ & $\kappa$ &
  \textbf{Koszul dual} & $\kappa^!$ \\
\hline
$bc_\lambda$ & $(\lambda, 1{-}\lambda)$ &
  $-2(6\lambda^2{-}6\lambda{+}1)$ &
  $-(6\lambda^2{-}6\lambda{+}1)$ &
  $\beta\gamma_\lambda$ &
  $6\lambda^2{-}6\lambda{+}1$ \\
$\beta\gamma_\mu$ & $(\mu, 1{-}\mu)$ &
  $2(6\mu^2{-}6\mu{+}1)$ &
  $6\mu^2{-}6\mu{+}1$ &
  $bc_\mu$ &
  $-(6\mu^2{-}6\mu{+}1)$ \\
\hline
$bc_2$ (ghosts) & $(2, -1)$ & $-26$ & $-13$ &
  $\beta\gamma_2$ & $13$ \\
$\beta\gamma_{3/2}$ (s.ghosts) & $(3/2, -1/2)$ & $11$ &
  $11/2$ & $bc_{3/2}$ & $-11/2$ \\
$bc_2 \otimes \beta\gamma_{3/2}$ & --- & $-15$ &
  $-15/2$ & $\beta\gamma_2 \otimes bc_{3/2}$ & $15/2$ \\
\hline
$bc_\lambda \otimes \beta\gamma_\lambda$ & --- & $0$ &
  $0$ & self-dual & $0$ \\
\hline
\end{tabular}
\end{center}

\begin{remark}[Complementarity sum: universal vanishing for the free-field family]
In every row, $\kappa + \kappa^! = 0$.  This is the universal
complementarity of the free-field family (class~G and class~C
in the shadow classification).  It fails for Virasoro
($\kappa + \kappa^! = 13$, AP24) and for general W-algebras
($\kappa + \kappa^! = \rho \cdot K \neq 0$ at shifted
levels).  The free-field case is the unique setting where
Koszul complementarity is a strict sign flip.
\end{remark}

\subsection{Koszul dual $\neq$ spin flip}

\begin{proposition}[Koszul dual $\neq$ spin flip]
\label{prop:bc-not-spin-flip}
The Koszul dual $bc_\lambda^! = \beta\gamma_\lambda$ is
\emph{not} the $bc$ system at dual spins $(1-\lambda, \lambda)$.
The central charges distinguish them:
\begin{align*}
c(bc_\lambda^!) &= c(\beta\gamma_\lambda)
= 2(6\lambda^2 - 6\lambda + 1), \\
c(bc_{1-\lambda}) &= c(bc_\lambda) = -2(6\lambda^2 - 6\lambda + 1).
\end{align*}
These are equal only if $6\lambda^2 - 6\lambda + 1 = 0$, i.e.,
$\lambda = (3 \pm \sqrt{3})/6$.  At all other values, the spin-flip
and Koszul duality are distinct operations.  Koszul duality changes
statistics; spin-flip permutes generators within the same algebra.
\end{proposition}

\begin{proof}
The spin-flip $\lambda \mapsto 1 - \lambda$ sends $b \mapsto c'$,
$c \mapsto b'$ with $c' = b$, $b' = c$: it is just a relabeling,
and $c_{bc}(1-\lambda) = c_{bc}(\lambda)$ (the formula is symmetric
under $\lambda \mapsto 1 - \lambda$).  Koszul duality sends
$\bigwedge^{\mathrm{ch}} \to \mathrm{Sym}^{\mathrm{ch}}$ and
negates the central charge.  These are completely different.
\end{proof}

\subsection{BRST anomaly cancellation under Koszul duality}

\begin{theorem}[Koszul duality preserves BRST anomaly cancellation
for free-field string backgrounds]
\label{thm:brst-koszul-preservation}
Let $\cA_{\mathrm{matter}} \otimes \cA_{\mathrm{ghost}}$ be a
free-field BRST system with $c_{\mathrm{matter}} +
c_{\mathrm{ghost}} = 0$.  If all constituent algebras are
free fields (tensor products of Heisenberg, $bc$, and
$\beta\gamma$ systems), then
\[
c_{\mathrm{matter}}^! + c_{\mathrm{ghost}}^! = 0.
\]
The Koszul dual system is again BRST-anomaly-free.
\end{theorem}

\begin{proof}
By the tensor product formula
(Theorem~\ref{thm:koszul-tensor-bc-bg}), it suffices to check
each factor.  For each factor, the complementarity constant
$K = 0$:
\begin{itemize}[nosep]
\item Heisenberg $\cH_k$: $\kappa(\cH_k^!) = -k$, so
  $K = k + (-k) = 0$.
\item $bc_\lambda$: $K = c_{bc} + c_{\beta\gamma} = 0$
  (equation~\eqref{eq:c-complementarity}).
\item $\beta\gamma_\mu$: $K = c_{\beta\gamma} + c_{bc} = 0$
  (same identity).
\end{itemize}
Therefore the total complementarity constant of any tensor
product of these systems is zero, and the BRST condition
$c_{\mathrm{total}} = 0$ implies
$c_{\mathrm{total}}^! = -c_{\mathrm{total}} = 0$.
\end{proof}

\subsection{The curvature--Koszul dictionary}

We collect the genus-$1$ curvatures for the full free-field
family, including the Virasoro comparison.

\begin{center}
\renewcommand{\arraystretch}{1.3}
\begin{tabular}{lcccc}
\hline
Algebra & $c$ & $\kappa$ & Koszul dual & $K$ \\
\hline
$bc_\nu$ & $-2(6\nu^2 - 6\nu + 1)$
  & $-(6\nu^2 - 6\nu + 1)$
  & $\beta\gamma_\nu$ & $0$ \\[3pt]
$\beta\gamma_\nu$ & $2(6\nu^2 - 6\nu + 1)$
  & $6\nu^2 - 6\nu + 1$
  & $bc_\nu$ & $0$ \\[3pt]
$\cH_k$ & $1$ & $k$ & $\mathrm{Sym}^{\mathrm{ch}}(V^*)$ & $0$ \\[3pt]
$\mathrm{Vir}_c$ & $c$ & $c/2$ & $\mathrm{Vir}_{26-c}$ & $13$ \\
\hline
\end{tabular}
\end{center}
The vanishing $K = 0$ for all free-field systems, and the
nonvanishing $K = 13$ for Virasoro, confirm the general pattern:
$K = 0$ when Koszul duality is the linear
$\mathrm{Ext}/\mathrm{Sym}$ exchange (free fields, affine
KM, lattice VOAs), and $K \neq 0$ when Koszul duality involves
the nonlinear DS reduction (Virasoro, $\cW_N$).


% =================================================================
\section{Higher $\Ainf$ operations: $m_4$ and $m_5$ for Virasoro}
\label{sec:higher-ainf-virasoro}
% =================================================================

The Virasoro $\lambda$-bracket
\begin{equation}\label{eq:vir-lambda-wn2}
\{T{}_\lambda T\}
\;=\;
\partial T + 2T\lambda + \frac{c}{12}\,\lambda^3
\end{equation}
defines $m_2(T,T;\lambda) = \{T{}_\lambda T\}$.  The quartic OPE
pole ($\lambda^3$ in the $\lambda$-bracket) forces
non-associativity, compensated by
\[
m_3(T,T,T;\,\lambda_{12},\lambda_{23})
\;=\;
\partial^2 T
+ (2\lambda_{12} + 3\lambda_{23})\,\partial T
+ 2\lambda_{23}(2\lambda_{12} + \lambda_{23})\,T
+ \frac{c}{12}\,\lambda_{23}^3(2\lambda_{12} + \lambda_{23}).
\]
This section computes the next two operations $m_4$ and $m_5$
explicitly, by solving the Stasheff identity at arities $4$ and~$5$.
These are the first genuinely new computations in the Virasoro
$\Ainf$ tower beyond the associator.

\subsection{Method: the Stasheff identity on generators}

For a one-generator vertex algebra with generator $T$ of conformal
weight~$2$, the $\Ainf$ operations $m_k$ are multilinear maps
\[
m_k \colon V^{\otimes k} \;\to\; V,
\qquad
V = \C[\partial]\,T \oplus \C \cdot \mathbf{1},
\]
with spectral parameters
$\lambda_{12}, \ldots, \lambda_{k-1,k}$.
The Stasheff identity at arity~$n$ reads:
\begin{equation}\label{eq:stasheff-arity-n-wn}
\sum_{\substack{r+s = n+1 \\ r \ge 2,\, s \ge 2}}
\;\sum_{j=0}^{n-s}
(-1)^{j}\,
m_r\!\bigl(T, \ldots, T,\,
  m_s(T, \ldots, T;\, \lambda_{j+1,j+2}, \ldots),\,
  T, \ldots, T;\, \text{outer params}\bigr)
\;=\; 0,
\end{equation}
where the sign $(-1)^j$ comes from the Koszul convention
with desuspended degree $|sT| = -1$,
and the (1,\,$n$) and ($n$,\,1) terms involving $m_1 = \partial$
are absorbed into the sesquilinearity of $m_n$.

The key sesquilinearity rules for the operations are:
\begin{itemize}
\item \textit{First slot} (left):
  $m_k(\partial^j T, T, \ldots, T;\, \lambda_{12}, \ldots)
  = (-\lambda_{12})^j\, m_k(T, T, \ldots, T;\, \lambda_{12}, \ldots)$.
\item \textit{Last slot} (right):
  $m_k(T, \ldots, T, \partial^j T;\, \ldots, \lambda_{k-1,k})
  = (\lambda_{k-1,k} + \partial)^j\, m_k(T, \ldots, T;\, \ldots, \lambda_{k-1,k})$.
\item \textit{Internal slots}: determined by the full associator
  computation (the lambda-bracket sesquilinearity propagated
  through the defining Stasheff relation).
\end{itemize}

The method: at each arity~$n$, compute all compositions of
lower-arity operations (the ``obstruction'' $\mathrm{Obs}_n$),
then set $m_n(T^n;\, \lambda_{12}, \ldots, \lambda_{n-1,n})
= -\mathrm{Obs}_n$.  This determines $m_n$ uniquely on the
generator~$T$; all other evaluations follow from sesquilinearity.

\subsection{The quaternary operation $m_4$}

The arity-$4$ Stasheff identity involves five compositions:
two from $(r,s) = (2,3)$ and three from $(r,s) = (3,2)$.
With Koszul signs:
\begin{align*}
\mathrm{Obs}_4
\;&=\;
\underbrace{\{m_3(T,T,T;\,\lambda_{12},\lambda_{23})
  {}_{\lambda_{34}} T\}}_{\text{(2,3), $j=0$}}
\;-\;
\underbrace{\{T{}_{\lambda_{12}}
  m_3(T,T,T;\,\lambda_{23},\lambda_{34})\}}_{\text{(2,3), $j=1$}}
\\[4pt]
&\quad +\;
\underbrace{m_3(m_2(T,T;\lambda_{12}),\, T,\, T;\,
  \lambda_{23},\, \lambda_{34})}_{\text{(3,2), $j=0$}}
\;-\;
\underbrace{m_3(T,\, m_2(T,T;\lambda_{23}),\, T;\,
  \lambda_{12},\, \lambda_{34})}_{\text{(3,2), $j=1$}}
\\[4pt]
&\quad +\;
\underbrace{m_3(T,\, T,\, m_2(T,T;\lambda_{34});\,
  \lambda_{12},\, \lambda_{23})}_{\text{(3,2), $j=2$}}.
\end{align*}

Each term is computed by applying the sesquilinearity rules
to the known operations $m_2$ and $m_3$.  For instance,
the first term evaluates the left $\lambda$-bracket
$\{f{}_\mu T\}$ of the output
$f = m_3(T,T,T;\,\lambda_{12},\lambda_{23})$ with $T$ at
spectral parameter $\mu = \lambda_{34}$, using
$\{\partial^k T{}_\mu T\} = (-\mu)^k\{T{}_\mu T\}$.
The third term uses first-slot sesquilinearity of $m_3$ to
handle the substitution of $m_2(T,T;\lambda_{12})
= \partial T + 2\lambda_{12} T + \frac{c}{12}\lambda_{12}^3$.
The fourth term requires the \emph{second-slot} sesquilinearity
of~$m_3$, computed from the associator with a modified second input.

The result, verified by symbolic computation:

\begin{theorem}[Quaternary Virasoro operation]
\label{thm:vir-m4-wn}
The quaternary $\Ainf$ operation for the Virasoro algebra is
\begin{equation}\label{eq:m4-virasoro-wn}
\boxed{
\begin{aligned}
m_4(T,T,T,T;\,&\lambda_{12},\lambda_{23},\lambda_{34})
\\[3pt]
&=\;
(\lambda_{23} + 2\lambda_{34})\,\partial^2 T
\\[3pt]
&\quad +\;
(5\lambda_{12}^2 - 4\lambda_{12}\lambda_{23}
  + 3\lambda_{12}\lambda_{34}
  - 5\lambda_{23}^2 + 10\lambda_{23}\lambda_{34}
  + 4\lambda_{34}^2)\,\partial T
\\[3pt]
&\quad +\;
2(\lambda_{12}^3 + 2\lambda_{12}^2\lambda_{23}
  + 3\lambda_{12}^2\lambda_{34}
  - 2\lambda_{12}\lambda_{23}^2
  - 4\lambda_{12}\lambda_{23}\lambda_{34}
  + 4\lambda_{12}\lambda_{34}^2
\\
&\qquad\qquad
  - \lambda_{23}^3 - 2\lambda_{23}^2\lambda_{34}
  + 6\lambda_{23}\lambda_{34}^2)\,T
\\[3pt]
&\quad +\;
\frac{c}{12}
(\lambda_{12}^5 + 2\lambda_{12}^4\lambda_{23}
  + 3\lambda_{12}^4\lambda_{34}
  + 4\lambda_{12}^3\lambda_{23}\lambda_{34}
  + 2\lambda_{12}^3\lambda_{34}^2
\\
&\qquad\qquad
  - 2\lambda_{12}\lambda_{23}^4
  - 4\lambda_{12}\lambda_{23}^3\lambda_{34}
  - 4\lambda_{12}\lambda_{23}\lambda_{34}^3
  + 2\lambda_{12}\lambda_{34}^4
\\
&\qquad\qquad
  - \lambda_{23}^5 - 2\lambda_{23}^4\lambda_{34}
  + 6\lambda_{23}\lambda_{34}^4).
\end{aligned}}
\end{equation}
\end{theorem}

\noindent\textit{Structural features of $m_4$.}
\begin{enumerate}[label=\textup{(\roman*)},nosep]
\item \textit{No $\partial^3 T$ term.}  The leading derivative
  in $m_4$ is $\partial^2 T$ with a linear coefficient, not
  the $\partial^{k-2}T = \partial^2 T$ with a constant
  coefficient that the lower operations might suggest.
  The operation $m_4$ starts at one derivative order lower
  than the na\"ive pattern.

\item \textit{Vanishes at $\lambda = 0$.}  When all spectral
  parameters are zero, $m_4 = 0$: there is no constant-mode
  quartic self-coupling.

\item \textit{Central charge: contact term only.}  The
  parameter~$c$ enters exclusively in the $c$-number
  (scalar) term $\frac{c}{12}Q_4(\lambda)$, where $Q_4$ is a
  degree-$5$ polynomial.  All $T$-dependent terms are
  $c$-independent.

\item \textit{Symmetric point.}  At
  $\lambda_{12} = \lambda_{23} = \lambda_{34} = \lambda$:
  \begin{equation}\label{eq:m4-symmetric-wn}
  m_4(T^4;\,\lambda,\lambda,\lambda)
  \;=\;
  3\lambda\,\partial^2 T
  + 13\lambda^2\,\partial T
  + 14\lambda^3\, T
  + \frac{7c}{12}\,\lambda^5.
  \end{equation}
  The coefficient $14 = 2 \cdot 7$ of the stress-tensor term
  is the quartic backreaction; the coefficient $7c/12$ of the
  scalar is the four-graviton coupling.

\item \textit{Connection to the quartic contact invariant.}
  The contact invariant
  $\mathfrak{Q}^{\mathrm{contact}}_{\mathrm{Vir}}
  = 10/[c(5c+22)]$
  from the Kac determinant at weight~$4$
  (\S\ref{sec:bar-knows-gravity})
  is the \emph{normalised} version of $m_4$: it divides
  $m_4$ by the Shapovalov inner product at level~$4$.
\end{enumerate}

\subsection{The quinary operation $m_5$}

At arity~$5$, the Stasheff identity involves compositions from
$(r,s) = (2,4)$, $(3,3)$, and $(4,2)$, with a total of nine terms
($2 + 3 + 4$).  The $(4,2)$ compositions require the internal
sesquilinearity of~$m_4$, which is obtained by recomputing the
$m_4$ obstruction with modified inputs.

\begin{theorem}[Quinary Virasoro operation]
\label{thm:vir-m5-wn}
The quinary $\Ainf$ operation for the Virasoro algebra is
\begin{equation}\label{eq:m5-virasoro-wn}
\boxed{
\begin{aligned}
m_5(T^5;\,
  &\lambda_{12},\lambda_{23},\lambda_{34},\lambda_{45})
\\[3pt]
&=\;
-\partial^4 T
\;-\;
(2\lambda_{12} + 4\lambda_{23}
  + \lambda_{34} + 3\lambda_{45})\,\partial^3 T
\\[3pt]
&\quad +\;
(3\lambda_{12}^2 - 16\lambda_{12}\lambda_{23}
  - 6\lambda_{12}\lambda_{45}
  - 13\lambda_{23}^2 + 7\lambda_{23}\lambda_{34}
\\
&\qquad\qquad
  - 9\lambda_{23}\lambda_{45}
  - 7\lambda_{34}^2
  + 3\lambda_{34}\lambda_{45}
  - 3\lambda_{45}^2)\,\partial^2 T
\\[3pt]
&\quad +\; P_1(\lambda)\,\partial T
\;+\; P_0(\lambda)\, T
\;+\; \frac{c}{12}\, Q_5(\lambda),
\end{aligned}}
\end{equation}
where
\begin{align*}
P_1(\lambda)
&\;=\;
10\lambda_{12}^2\lambda_{23}
- 9\lambda_{12}^2\lambda_{34}
+ 14\lambda_{12}^2\lambda_{45}
- 5\lambda_{12}\lambda_{23}^2
- 20\lambda_{12}\lambda_{23}\lambda_{34}
\\
&\quad
- 49\lambda_{12}\lambda_{23}\lambda_{45}
- 8\lambda_{12}\lambda_{34}^2
+ 13\lambda_{12}\lambda_{34}\lambda_{45}
- 5\lambda_{12}\lambda_{45}^2
+ 2\lambda_{23}^3
\\
&\quad
- 9\lambda_{23}^2\lambda_{34}
- 45\lambda_{23}^2\lambda_{45}
- 3\lambda_{23}\lambda_{34}^2
+ 41\lambda_{23}\lambda_{34}\lambda_{45}
+ \lambda_{23}\lambda_{45}^2
\\
&\quad
- 9\lambda_{34}^3
- 5\lambda_{34}^2\lambda_{45}
+ 5\lambda_{34}\lambda_{45}^2
- 3\lambda_{45}^3,
\displaybreak[0]\\[6pt]
P_0(\lambda)
&\;=\;
4\lambda_{12}^3\lambda_{23}
- 2\lambda_{12}^3\lambda_{34}
+ 2\lambda_{12}^3\lambda_{45}
+ 10\lambda_{12}^2\lambda_{23}^2
- 4\lambda_{12}^2\lambda_{23}\lambda_{34}
\\
&\quad
+ 14\lambda_{12}^2\lambda_{23}\lambda_{45}
- 8\lambda_{12}^2\lambda_{34}^2
- 6\lambda_{12}^2\lambda_{34}\lambda_{45}
+ 8\lambda_{12}^2\lambda_{45}^2
\\
&\quad
- 4\lambda_{12}\lambda_{23}^3
- 4\lambda_{12}\lambda_{23}^2\lambda_{34}
- 16\lambda_{12}\lambda_{23}^2\lambda_{45}
- 16\lambda_{12}\lambda_{23}\lambda_{34}^2
\\
&\quad
- 16\lambda_{12}\lambda_{23}\lambda_{34}\lambda_{45}
- 32\lambda_{12}\lambda_{23}\lambda_{45}^2
+ 4\lambda_{12}\lambda_{34}^3
- 4\lambda_{12}\lambda_{34}^2\lambda_{45}
\\
&\quad
+ 8\lambda_{12}\lambda_{34}\lambda_{45}^2
+ 4\lambda_{12}\lambda_{45}^3
+ 2\lambda_{23}^3\lambda_{34}
+ 2\lambda_{23}^3\lambda_{45}
\\
&\quad
- 8\lambda_{23}^2\lambda_{34}^2
- 16\lambda_{23}^2\lambda_{34}\lambda_{45}
- 32\lambda_{23}^2\lambda_{45}^2
+ 6\lambda_{23}\lambda_{34}^3
\\
&\quad
+ 46\lambda_{23}\lambda_{34}\lambda_{45}^2
+ 2\lambda_{23}\lambda_{45}^3
- 2\lambda_{34}^4
- 16\lambda_{34}^3\lambda_{45}
\\
&\quad
+ 10\lambda_{34}^2\lambda_{45}^2
- 2\lambda_{34}\lambda_{45}^3
- 2\lambda_{45}^4,
\end{align*}
and $Q_5(\lambda)$ is a degree-$6$ polynomial in
$\lambda_{12}, \ldots, \lambda_{45}$, proportional to~$c$,
whose symmetric-point value is $Q_5(\lambda, \ldots, \lambda)
= -34\lambda^6$.
\end{theorem}

\noindent\textit{Structural features of $m_5$.}
\begin{enumerate}[label=\textup{(\roman*)},nosep]
\item \textit{Sign alternation.}  The leading term is
  $-\partial^4 T$, with a minus sign.  This is the first sign
  alternation in the tower: $m_2$ and $m_3$ have positive leading
  terms ($+\partial T$, $+\partial^2 T$), $m_4$ has no
  $\partial^3 T$ term, and $m_5$ begins with $-\partial^4 T$.
  The sign alternation is the Koszul sign from the bar
  desuspension.

\item \textit{Full derivative range.}  Unlike $m_4$ (which has
  no $\partial^3 T$ term), $m_5$ populates the entire range
  $\partial^4 T$ through~$T$ plus a scalar.

\item \textit{Symmetric-point values.}  At all spectral
  parameters equal to~$\lambda$:
  \begin{equation}\label{eq:m5-symmetric-wn}
  \begin{aligned}
  m_5(T^5;\,\lambda,\lambda,\lambda,\lambda)
  &\;=\;
  -\partial^4 T
  - 10\lambda\,\partial^3 T
  - 41\lambda^2\,\partial^2 T
  \\
  &\quad
  - 84\lambda^3\,\partial T
  - 68\lambda^4\, T
  - \frac{17c}{6}\,\lambda^6.
  \end{aligned}
  \end{equation}
  All coefficients are negative: the five-graviton contact
  vertex is uniformly attractive at the symmetric point.

\item \textit{Central charge dependence.}  As for $m_4$,
  the parameter~$c$ enters exclusively through the scalar
  term $\frac{c}{12}Q_5(\lambda)$.  All $T$-dependent
  coefficients are $c$-independent.
\end{enumerate}

\subsection{The degree pattern and the infinite tower}

The computations reveal a uniform pattern for the Virasoro
$\Ainf$ operations:

\begin{center}
\small
\renewcommand{\arraystretch}{1.3}
\begin{tabular}{@{}cccc@{}}
\textbf{Operation}
  & \textbf{Leading term}
  & \textbf{$\deg_\lambda$}
  & \textbf{Scalar} \\
\hline
$m_2(T,T;\lambda)$
  & $\partial T$
  & $\max = 3$
  & $\frac{c}{12}\lambda^3$ \\[3pt]
$m_3(T^3;\lambda_{12},\lambda_{23})$
  & $\partial^2 T$
  & $\max = 4$
  & $\frac{c}{12}\lambda_{23}^3(2\lambda_{12}+\lambda_{23})$ \\[3pt]
$m_4(T^4;\lambda_{12},\lambda_{23},\lambda_{34})$
  & $(\lambda_{23}+2\lambda_{34})\,\partial^2 T$
  & $\max = 5$
  & $\frac{c}{12}Q_4(\lambda)$ \\[3pt]
$m_5(T^5;\lambda_{12},\ldots,\lambda_{45})$
  & $-\partial^4 T$
  & $\max = 6$
  & $\frac{c}{12}Q_5(\lambda)$
\end{tabular}
\end{center}

\medskip

\noindent
Three structural features persist:

\begin{enumerate}[label=\textup{(\alph*)},nosep]
\item \textit{Scalar proportional to $c/12$.}  The central charge
  enters every $m_k$ through a single scalar term
  $\frac{c}{12}Q_k(\lambda_{12}, \ldots, \lambda_{k-1,k})$ where
  $Q_k$ has degree $2k - 3$ in the spectral parameters.  All
  $T$-dependent coefficients are $c$-independent.  This reflects the
  structure of the Virasoro OPE: the only $c$-dependent singular term
  is $(c/2)/(z-w)^4$, and the bar kernel $d\log(z-w)$ absorbs one
  power to produce $(c/2)/(z-w)^3$ in the collision residue.

\item \textit{Degree pattern.}  The coefficient of $\partial^j T$ in
  $m_k$ is a polynomial in the spectral parameters of degree $k - 2 - j$
  (for the $T$-dependent terms) or $2k - 3$ (for the scalar).  The
  total ``weight'' $j + \deg_\lambda = k - 2$ is constant across the
  $T$-dependent terms.

\item \textit{Vanishing at $\lambda = 0$.}  Every $m_k(T^k)$ vanishes
  when all spectral parameters are zero: the $\Ainf$ structure is
  purely spectral.  This is the algebraic content of the statement that
  the Virasoro algebra has no higher multiplications at zero momentum
  transfer.
\end{enumerate}

\subsection{Symmetric-point sequence and combinatorial structure}

The symmetric-point evaluations
$m_k(T^k;\,\lambda, \ldots, \lambda)$ for $k = 2, 3, 4, 5$
produce the coefficient sequences:
\begin{center}
\small
\renewcommand{\arraystretch}{1.3}
\begin{tabular}{@{}c*{6}{c}c@{}}
$k$
  & $\partial^{k-2}T$ & $\cdots$ & $\partial^2 T$
  & $\partial T$ & $T$ & scalar$/\lambda^{k+1}$ \\
\hline
$2$ & $1$ & & & $2$ & & $c/12$ \\
$3$ & $1$ & & & $5$ & $6$ & $c/4$ \\
$4$ & & & $3$ & $13$ & $14$ & $7c/12$ \\
$5$ & $-1$ & $-10$ & $-41$ & $-84$ & $-68$ & $-17c/6$
\end{tabular}
\end{center}

\noindent
The sign alternation between $m_4$ and $m_5$ (positive coefficients
to negative) is the first instance of the Koszul sign in the tower.
For $m_4$, the absence of a $\partial^3 T$ term is not accidental: it
reflects the associahedron $K_4$ cell structure, which has no
codimension-$1$ boundary contributing to the $\partial^3 T$ component
at this arity.

The scalar coefficients $1, 3, 7, -34$ (times $c/12$, with
appropriate $\lambda$-powers) suggest the growth rate is exponential
in~$k$: the $\Ainf$ tower does not truncate, and $m_k \neq 0$ for
all $k \geq 3$ and generic~$c$.  The quartic OPE pole is the engine
of this infinite tower.

\begin{remark}[$c = 0$: the Witt tower]
\label{rem:witt-tower-wn}
At $c = 0$, the scalar terms vanish: $Q_k = 0$ for all~$k$.
The $\Ainf$ operations $m_k$ retain their $T$-dependent terms
unchanged.  This is the $\Ainf$ structure of the Witt algebra
(the Virasoro algebra at $c = 0$), which remains genuinely
infinite despite the absence of the central charge.  The quartic
pole persists at $c = 0$ (the $\lambda^3$ term is purely
$c$-dependent, but the lower-order OPE terms $\partial T + 2T\lambda$
remain), and the non-associativity of $m_2$ continues to force the
entire tower.
\end{remark}

\begin{remark}[The gravitational Feynman rules]
\label{rem:gravitational-feynman-wn}
The operations $m_k(T^k)$ are the tree-level $k$-graviton contact
vertices of three-dimensional quantum gravity in the $\Ainf$
formulation.  The fact that $c$ enters only through the scalar
term means: at tree level, the stress-tensor backreaction ($T$,
$\partial T$, $\partial^2 T$ terms) is universal --- independent
of the central charge and hence of the cosmological constant ---
while the pure-graviton scattering amplitude (the scalar term) is
proportional to~$c/12 = \kappa/6$.  The cosmological constant
affects amplitudes, not backreaction.
\end{remark}




% =================================================================
\section{Genus-$1$ curved bar complex: explicit differential and $F_1$}
\label{sec:genus-1-curved-bar}
% =================================================================

This section carries out the complete computation of the genus-$1$
curved bar differential for the Virasoro algebra $\mathrm{Vir}_c$
on an elliptic curve $E_\tau = \C / (\Z + \Z\tau)$.  The
output is:
\begin{enumerate}[label=(\roman*)]
\item the explicit genus-$1$ propagator $\eta^{(1)}_{12}$ and its
  Laurent expansion near the diagonal;
\item the Arnold relation \emph{failure} at genus~$1$, producing
  the Arakelov $(1,1)$-form $\omega_1$;
\item the curved bar differential
  $d^{(1)}_{\barB}$ on degree-$2$ bar elements
  $[s^{-1}a \,|\, s^{-1}b]$ for Virasoro states of arbitrary
  weight;
\item the genus-$1$ correction to the $\Ainf$ operations
  $m_k^{(1)}$;
\item the genus-$1$ free energy $F_1(c)$ as a function of central
  charge;
\item the critical string dichotomy: $F_1 = 0$ at $c = 26$,
  $F_1 \neq 0$ at $c \neq 26$.
\end{enumerate}

\medskip

\noindent
\textbf{Convention.}
Throughout, $\kappa = \kappa(\mathrm{Vir}_c) = c/2$ is the
\emph{intrinsic} modular characteristic.  The \emph{effective}
curvature $\kappa_{\mathrm{eff}} = (c-26)/2$ arises only
after coupling to the $bc$-ghost system ($\kappa_{\mathrm{ghost}} = -13$).
The algebraic bar complex sees $\kappa$; the physical partition
function uses $\kappa_{\mathrm{eff}}$.  We compute both.

%%----------------------------------------------------------------------
\subsection{The genus-$1$ propagator}
\label{subsec:wn-genus1-propagator}
%%----------------------------------------------------------------------

The Arakelov-normalised propagator on $E_\tau$ is
(Proposition~\ref{prop:arakelov-genus-1-derivation}):
\begin{equation}\label{eq:wn-eta-1}
\eta^{(1)}(z_1, z_2;\tau)
\;=\;
\partial_{z_1}\log\vartheta_1(z_1 - z_2 \,|\, \tau)
\;+\;
\frac{2\pi i \,\operatorname{Im}(z_1 - z_2)}
     {\operatorname{Im}\tau}.
\end{equation}
Write $u = z_1 - z_2$.  Decompose into holomorphic and
correction:
\[
\eta^{(1)} = h(u;\tau) + R(u;\tau),
\qquad
h(u;\tau) = \partial_u\log\vartheta_1(u|\tau),
\qquad
R(u;\tau) = \frac{2\pi i\,\operatorname{Im}(u)}
{\operatorname{Im}\tau}.
\]

\noindent
\textbf{Laurent expansion near $u = 0$.}
\begin{equation}\label{eq:wn-laurent}
h(u;\tau)
\;=\;
\frac{1}{u}
- \frac{\pi^2}{3}\,E_2(\tau)\,u
- \frac{\pi^4}{45}\,E_4(\tau)\,u^3
- \cdots,
\end{equation}
where $E_2(\tau) = 1 - 24\sum_{n=1}^\infty \sigma_1(n)q^n$,
$E_4(\tau) = 1 + 240\sum_{n=1}^\infty \sigma_3(n)q^n$ are
Eisenstein series.  The expansion has only odd powers of $u$
(since $\vartheta_1$ is odd).  The correction $R$ is smooth at
$u = 0$: it contributes no residue.

%%----------------------------------------------------------------------
\subsection{The Arnold relation failure at genus $1$}
\label{subsec:wn-arnold-failure}
%%----------------------------------------------------------------------

At genus $0$, the Arnold relation
$\sum_{\mathrm{cyc}}\eta^{(0)}_{12}\wedge\eta^{(0)}_{23} = 0$
ensures $d_{\barB}^2 = 0$.  At genus $1$, the
holomorphic--holomorphic terms vanish by the Fay trisecant
identity; the cross-terms produce the Arakelov form:
\begin{equation}\label{eq:wn-arnold-defect}
\boxed{
\sum_{\mathrm{cyc}}
\eta^{(1)}_{12} \wedge \eta^{(1)}_{23}
\;=\;
\frac{i}{2\,\mathrm{Im}\tau}\,dz \wedge d\bar{z}
\;=:\; \omega_1,
}
\end{equation}
normalised so that $\int_{E_\tau}\omega_1 = 1$.

%%----------------------------------------------------------------------
\subsection{The curved bar differential $d_{\barB}^{(1)}$ on
  degree-$2$ elements}
\label{subsec:wn-curved-differential}
%%----------------------------------------------------------------------

Let $a, b$ be Virasoro states of weights $h_a, h_b$.
The genus-$1$ fibre differential is:
\begin{equation}\label{eq:wn-dfib-genus1}
\dfib^{(1)}[s^{-1}a \,|\, s^{-1}b]
\;=\;
s^{-1}\operatorname{Res}_{z_1 = z_2}\Bigl[
  \eta^{(1)}(z_1, z_2;\tau) \cdot a(z_1)\,b(z_2)
\Bigr].
\end{equation}

\noindent\textbf{Explicit computation for $a = b = T$ (weight $2$).}
Extracting the coefficient of $u^{-1}$ from
$\eta^{(1)}(u;\tau)\cdot T(z_1)\,T(z_2)$:
\begin{equation}\label{eq:wn-dfib-TT}
\boxed{
\dfib^{(1)}[s^{-1}T \,|\, s^{-1}T]
\;=\;
s^{-1}\Bigl(
  {:}TT{:}
  \;-\; \frac{2\pi^2}{3}\,E_2(\tau)\cdot T
  \;-\; \frac{c\pi^4}{90}\,E_4(\tau)\cdot \mathbf{1}
\Bigr).
}
\end{equation}
Three terms: the genus-$0$ normal-ordered product ${:}TT{:}$;
the $E_2$-correction from the double pole $2T/u^2$ paired with
$c_1(\tau)\,u$; the $E_4$-correction from the quartic pole
$(c/2)/u^4$ paired with $c_2(\tau)\,u^3$.

%%----------------------------------------------------------------------
\subsection{General weight formula}
\label{subsec:wn-general-weight}
%%----------------------------------------------------------------------

For states $a, b$ of weights $h_a, h_b$:
\begin{equation}\label{eq:wn-general-weight}
\boxed{
\dfib^{(1)}[s^{-1}a \,|\, s^{-1}b]
\;=\;
s^{-1}\Bigl(
  a_{(-1)}b
  \;+\;
  \sum_{k=1}^{\lfloor h_a+h_b/2\rfloor}
  c_k(\tau)\cdot a_{(2k-1)}b
\Bigr),
}
\end{equation}
where the propagator coefficients are:
\begin{equation}\label{eq:wn-propagator-coeffs}
c_k(\tau) = (-1)^k\frac{(2\pi)^{2k}}{(2k)!}\,B_{2k}\,E_{2k}(\tau),
\end{equation}
with $B_{2k}$ the Bernoulli numbers and $E_{2k}(\tau)$ the
normalised Eisenstein series.  Explicitly:
$c_1 = -\frac{\pi^2}{3}E_2$,
$c_2 = -\frac{\pi^4}{45}E_4$,
$c_3 = -\frac{2\pi^6}{945}E_6$.

%%----------------------------------------------------------------------
\subsection{The curvature and the genus-$1$ free energy}
\label{subsec:wn-curvature-F1}
%%----------------------------------------------------------------------

The curvature: $\dfib^{(1)\,2} = (c/2)\cdot\omega_1$ (intrinsic);
$((c-26)/2)\cdot\omega_1$ (effective).

The genus-$1$ free energy:
\begin{equation}\label{eq:wn-F1-intrinsic}
F_1(\mathrm{Vir}_c) = \frac{c}{48} \quad\text{(intrinsic)},
\qquad
F_1^{\mathrm{eff}}(\mathrm{Vir}_c) = \frac{c-26}{48}
\quad\text{(effective)}.
\end{equation}
The $\tau$-dependent partition function:
$Z_1^{\mathrm{scal}}(\tau) = \eta(\tau)^{-c/2}$ (intrinsic),
$Z_1^{\mathrm{eff}}(\tau) = \eta(\tau)^{-(c-26)/2}$ (effective).

%%----------------------------------------------------------------------
\subsection{The critical string dichotomy}
\label{subsec:wn-critical}
%%----------------------------------------------------------------------

\begin{center}
\renewcommand{\arraystretch}{1.2}
\begin{tabular}{@{}lcc@{}}
& $c \neq 26$ & $c = 26$ \\
\hline
$\kappa_{\mathrm{eff}}$ & $(c-26)/2 \neq 0$ & $0$ \\[3pt]
$F_1^{\mathrm{eff}}$ & $(c-26)/48 \neq 0$ & $0$ \\[3pt]
$\dfib^{(1)\,2}$ & $\frac{c-26}{2}\omega_1 \neq 0$ & $0$ \\[3pt]
Genus tower & nonzero at all $g$ & vanishes identically \\[3pt]
Regime & perturbative (curved bar) & topological (flat bar)
\end{tabular}
\end{center}

\noindent
The dichotomy is the bar-complex avatar of the conformal anomaly
vanishing at $c = 26$.  Each OPE pole has a genus-$1$ shadow:
the Eisenstein series $E_{2k}(\tau)$ is the shadow of the
$(2k)$-th pole.


% =================================================================
\section{BTZ entropy and Cardy formula from the bar complex}
\label{sec:btz-entropy-cardy}
% =================================================================

The modular MC element $\Theta_{\mathrm{grav}}$ of the Virasoro
algebra encodes, through its genus tower and its modular
transformation properties, the complete thermodynamics of the
BTZ black hole.  This section extracts the Cardy formula, the
BTZ entropy, and the Brown--Henneaux central charge from the
bar-complex perspective.

The result is a \emph{derivation} of the Bekenstein--Hawking
entropy of the BTZ black hole from the $\Ainf$ bar complex of
a single algebraic datum: the Virasoro $\lambda$-bracket
$\{T_\lambda T\} = \partial T + 2T\lambda + (c/12)\lambda^3$.

\subsection{The partition function from the genus tower}
\label{subsec:btz-partition}

Write $q = e^{2\pi i\tau}$ for the modular parameter of the
boundary torus $\partial(\text{solid torus})$.  The genus-$1$
MC element $\Theta^{(1)}$ gives the one-loop partition function
as the categorical trace:
\begin{equation}\label{eq:btz-Z1}
Z_1(\tau)
\;=\;
\Tr_{\mathrm{Vir}_c}\!\left(q^{L_0 - c/24}\right)
\;=\;
\frac{q^{-c/24}}{\prod_{n=2}^{\infty}(1-q^n)}.
\end{equation}
The product starts at $n=2$ because $\mathrm{Vir}_c$ has no
weight-$1$ states (no affine currents: the graviton is spin-$2$).

The scalar-lane effective partition function, incorporating the
$bc$ ghost system with $c_{\mathrm{ghost}} = -26$, is
$Z_{\mathrm{eff}}^{\mathrm{scal}}(\tau) = \eta(\tau)^{-(c-26)/2}$.

\begin{observation}[Two partition functions]
\label{obs:two-partition-functions-wn}
The bar-complex framework distinguishes:
(i) the \emph{intrinsic} partition function
$Z_1^{\mathrm{intr}}(\tau) = \eta(\tau)^{-c/2}$ from the
curvature $\kappa = c/2$;
(ii) the \emph{effective} partition function
$Z_1^{\mathrm{eff}}(\tau) = \eta(\tau)^{-(c-26)/2}$ from
$\kappa_{\mathrm{eff}} = (c-26)/2$ after ghost subtraction.
The Cardy formula uses the \emph{full} Virasoro partition
function, not the scalar-lane partition function.
\end{observation}

\subsection{The Cardy formula from the bar complex}
\label{subsec:btz-cardy-derivation}

Four steps: (1) highest-weight partition function; (2) modular
$S$-transform; (3) saddle-point evaluation; (4) entropy extraction.

A BTZ black hole of mass $M$ and angular momentum $J$ corresponds
to a highest-weight state $|h,\bar h\rangle$ with
$h = (M\ell + J)/2$, $\bar h = (M\ell - J)/2$.
The partition function is
$Z(\tau,\bar\tau) = \sum_{h,\bar h} d(h,\bar h)\,
q^{h - c/24}\,\bar q^{\bar h - c/24}$.
Modular invariance and vacuum dominance give, at high temperature,
$Z \sim \exp(\pi c/(12\,\mathrm{Im}(-1/\tau)))$.
The saddle-point of $f(\tau) = \pi c/(12\tau) + 2\pi\tau(h-c/24)$
is at $\tau_* = i\sqrt{c/(24h)}$, giving
$f(\tau_*) \approx 2\pi\sqrt{ch/6}$.

\begin{theorem}[Cardy formula from the bar complex]
\label{thm:btz-cardy-wn}
The asymptotic entropy of the Virasoro bar complex at large
conformal weight is
\begin{equation}\label{eq:btz-cardy-wn}
\boxed{
S_{\mathrm{Cardy}}
\;=\;
2\pi\sqrt{\frac{c\,h}{6}}
\;+\;
2\pi\sqrt{\frac{c\,\bar h}{6}}.
}
\end{equation}
\end{theorem}

\begin{proof}
Three bar-complex inputs: (i) curvature $\kappa = c/2$
determining vacuum energy $-c/24$; (ii) modular invariance
from the genus-$1$ MC equation on $\cM_{1,1}$; (iii)
saddle-point at $\tau_* = i\sqrt{c/(24h)}$.  All determined
by $\kappa = c/2$.
\end{proof}

\subsection{BTZ black hole entropy}
\label{subsec:btz-entropy-identification}

\begin{theorem}[BTZ entropy from the bar complex]
\label{thm:btz-entropy-wn}
The Bekenstein--Hawking entropy of the BTZ black hole is:
\begin{equation}\label{eq:btz-entropy-final-wn}
\boxed{
S_{\mathrm{BTZ}}
\;=\;
\frac{2\pi r_+}{4G}
\;=\;
2\pi\sqrt{\frac{c\,h}{6}}
+
2\pi\sqrt{\frac{c\,\bar h}{6}},
}
\end{equation}
where $c = 3\ell/(2G)$ is the Brown--Henneaux central charge.
Equivalently, $S_{\mathrm{BTZ}} = 2\pi\sqrt{\kappa h/3}
+ 2\pi\sqrt{\kappa\bar h/3}$ with $\kappa = c/2$.
\end{theorem}

\begin{proof}
Substituting $c = 3\ell/(2G)$, $h = (M\ell+J)/2$,
$\bar h = (M\ell-J)/2$ and using
$r_\pm^2 = \ell^2(M\ell \pm J)/2$ gives
$S_{\mathrm{BTZ}} = 2\pi(r_+ + r_-)/(4G)$, reducing to
$2\pi r_+/(4G)$ for $J = 0$.
\end{proof}

\subsection{Verdier duality and the modular $S$-transform}
\label{subsec:btz-verdier-modular}

\begin{proposition}[Commutativity of $\mathbb{D}$ and $S$]
\label{prop:btz-D-S-commute-wn}
The Verdier involution $\mathbb{D}\colon c \mapsto 26-c$ and the
modular $S$-transform $\tau \mapsto -1/\tau$ commute:
$S \circ \mathbb{D}(\Theta_{\mathrm{Vir}_c}^{(1)}(\tau))
= \mathbb{D} \circ S(\Theta_{\mathrm{Vir}_c}^{(1)}(\tau))
= \Theta_{\mathrm{Vir}_{26-c}}^{(1)}(-1/\tau)$.
They act on independent variables: algebraic data vs.\ modular
parameter.
\end{proposition}

\begin{theorem}[Paired Cardy formula]
\label{thm:btz-paired-cardy-wn}
$S_{\mathrm{Cardy}}(\mathrm{Vir}_c;\,h)
+ S_{\mathrm{Cardy}}(\mathrm{Vir}_{26-c};\,h)
= 2\pi\sqrt{ch/6} + 2\pi\sqrt{(26-c)h/6}$.
At the self-dual point $c = 13$ the two terms are equal;
at $c = 26$ the dual contribution vanishes.
The total counting capacity is bounded:
$\sqrt{c} + \sqrt{26-c} \le 2\sqrt{13}$, with equality at
$c = 13$.  The self-dual gravitational theory has the maximal
entropy among all Koszul pairs.
\end{theorem}

\subsection{Summary}

The entire thermodynamics of the BTZ black hole---the Cardy
formula, the Bekenstein--Hawking entropy, the Brown--Henneaux
central charge, and the complementarity bound---is extracted
from the genus tower of the Virasoro bar complex, determined by
$\kappa = c/2$.  The curvature of the bar complex IS the
gravitational entropy:
$S_{\mathrm{BTZ}} = 2\pi\sqrt{\kappa h/3} = 2\pi\sqrt{ch/6}$.
One quartic pole generates black hole thermodynamics.



% ============================================================
%
%   The critical scalar C^{res}_{4,4;2;0,6}(4):
%   verification from the W_4 OPE
%
% ============================================================

\section{The critical scalar
\texorpdfstring{$\mathsf{C}^{\mathrm{res}}_{4,4;2;0,6}(4)=2$}{C\^res\_{4,4;2;0,6}(4)=2}:
verification from the $W_4$ OPE}
\label{sec:critical-scalar-verification}

The $\cW_\infty$ stage-$4$ MC4 closure reduces, after the full Ward-normalized
visible pairing reduction (Corollary~\ref*{cor:winfty-stage4-single-scalar-equivalent}
of Vol~I), to a single scalar identity
\[
\mathsf{C}^{\mathrm{res}}_{4,4;2;0,6}(4)=2.
\]
We verify this identity from first principles by three independent routes.

\subsection{Route 1: the conformal Ward identity (universal for all $s$)}

\begin{proposition}[Universal self-coupling $T$-coefficient]
\label{prop:universal-T-coefficient}
Fix $N\ge 2$.  In any stage-$N$ OPE calculus with Virasoro-primary
generators $W^{(2)},\dots,W^{(N)}$ satisfying the Zamolodchikov
normalisation
\[
\langle W^{(s)}(z)\,W^{(s)}(0)\rangle = \frac{c/s}{z^{2s}},
\qquad c\neq 0,
\]
the undifferentiated $T$-channel coefficient at the leading $T$-pole
in the self-OPE satisfies
\[
\mathsf{C}_{s,s;2;0,\,2s-2}(N)=2
\qquad\text{for all } 2\le s\le N.
\]
In particular, $\mathsf{C}_{4,4;2;0,6}(4)=2$.
\end{proposition}

\begin{proof}
Write $N_s:=c/s$ and consider the Virasoro Ward identity for the
three-point function $\langle T(x)\,W^{(s)}(z)\,W^{(s)}(0)\rangle$:
\begin{equation}\label{eq:ward-three-point}
\langle T(x)\,W^{(s)}(z)\,W^{(s)}(0)\rangle
=
\biggl[
\frac{s}{(x-z)^2}+\frac{\partial_z}{x-z}
+\frac{s}{x^2}+\frac{\partial_0}{x}
\biggr]
N_s\,z^{-2s}.
\end{equation}
Here $\partial_0$ acts on the position $w$ in $(z-w)^{-2s}$
evaluated at $w=0$, giving $2s\,z^{-(2s+1)}$.  Expanding for
$|x|>|z|$ in the region $x^{-4}z^{-(2s-2)}$:

\medskip
\noindent\emph{From $s\cdot z^{-2s}/(x-z)^2$.}\enspace
Use $1/(x-z)^2 = \sum_{n\ge 0}(n+1)(z/x)^n/x^2$.
For $x^{-4}z^{-(2s-2)}$: need $n=2$, coefficient $(n+1)=3$.
Contribution: $s\cdot 3\cdot N_s$.

\medskip
\noindent\emph{From $(-2s)\cdot z^{-(2s+1)}/(x-z)$.}\enspace
Use $1/(x-z) = \sum_{n\ge 0}(z/x)^n/x$.
For $x^{-4}z^{-(2s-2)}$: need $n=3$, coefficient $1$.
Contribution: $(-2s)\cdot 1\cdot N_s$.

\medskip
\noindent\emph{The remaining terms} $s\,z^{-2s}/x^2$ and
$2s\,z^{-(2s+1)}/x$ contribute only at $x^{-2}z^{-2s}$ and
$x^{-1}z^{-(2s+1)}$ respectively, and do not reach $x^{-4}z^{-(2s-2)}$.

\medskip
Summing:
\[
[\text{Ward}]_{x^{-4}z^{-(2s-2)}}
= N_s\bigl(3s - 2s\bigr)
= s\,N_s
= s\cdot \frac{c}{s}
= c.
\]

On the OPE side, the self-OPE $W^{(s)}(z)\,W^{(s)}(w)$ contains at
pole order $2s-2$ the undifferentiated $T$-primary channel
$\mathsf{C}_{s,s;2;0,2s-2}\cdot T(w)/(z-w)^{2s-2}$.  Inserting into
the three-point function and using
$\langle T(x)\,T(0)\rangle = (c/2)\,x^{-4}$, the OPE produces at
$x^{-4}z^{-(2s-2)}$:
\[
[\text{OPE}]_{x^{-4}z^{-(2s-2)}}
=
\mathsf{C}_{s,s;2;0,2s-2}\cdot \frac{c}{2}.
\]
Equating and dividing by $c\neq 0$:
\[
\mathsf{C}_{s,s;2;0,\,2s-2}
= \frac{c}{c/2}
= 2.
\qedhere
\]
\end{proof}

\begin{remark}[Universality]
The identity $\mathsf{C}_{s,s;2;0,2s-2}=2$ is \emph{independent of
the algebra}~$\cW_N$: it depends only on conformal symmetry and the
normalisation convention.  The Ward identity forces it for every
Virasoro-primary generator of every $\cW$-algebra at every value of~$c$,
including the free-field realisation, the DS reduction, and the abstract
bootstrap.  In particular, it holds at $s=2$ (the Virasoro TT OPE,
giving $\mathsf{C}_{2,2;2;0,2}=2$ as the coefficient of $T$ in
$T(z)T(w)\sim (c/2)/(z-w)^4 + 2T(w)/(z-w)^2 + \partial T(w)/(z-w)$),
at $s=3$ (the $W_3$ self-OPE, giving $\mathsf{C}_{3,3;2;0,4}=2$
already proved in Vol~I), and at $s=4$ (the critical identity).
\end{remark}

\subsection{Route 2: equivalence with the Zamolodchikov normalisation}

\begin{proposition}[Equivalence]
\label{prop:equivalence-zamolodchikov}
The identity $\mathsf{C}_{4,4;2;0,6}(4)=2$ is logically equivalent to
the Zamolodchikov normalisation
$\langle W^{(4)}(z)\,W^{(4)}(0)\rangle = (c/4)\,z^{-8}$,
given the Virasoro Ward identity for the visible generators.
\end{proposition}

\begin{proof}
The proof of Proposition~\ref{prop:universal-T-coefficient} shows:
\[
\mathsf{C}_{4,4;2;0,6}=2
\quad\Longleftrightarrow\quad
4\,N_4 = c
\quad\Longleftrightarrow\quad
N_4 = \frac{c}{4}.
\]
The forward direction is the content of the proposition.  The converse:
if $\mathsf{C}_{4,4;2;0,6}=2$, then the Ward identity computation
gives $4\,N_4 = 2\cdot(c/2)=c$, hence $N_4=c/4$.
\end{proof}

\begin{remark}[Why the identity is non-trivial for MC4]
The identity $\mathsf{C}_{4,4;2;0,6}=2$ is trivially true \emph{on
the DS side}: the DS reduction produces the $\cW_4$ vertex algebra
with standard normalisation, so $N_4 = c/4$ by construction.  The
non-trivial content is that the \emph{residue side}
$\mathsf{C}^{\mathrm{res}}_{4,4;2;0,6}(4)$, extracted from the bar
differential of the factorisation target $\cW_\infty^{\mathrm{fact}}$,
matches the DS value.  This matching asserts that the stage-$4$ quotient
of the factorisation target inherits the Zamolodchikov normalisation from
the visible Virasoro package.  The Ward identity proof
(Proposition~\ref{prop:universal-T-coefficient}) shows that this
inheritance is \emph{automatic} once the visible generators are
Virasoro-primary with the standard $T$-action --- no additional input
from the DS reduction or the bootstrap is required.

This is the key observation that closes the identity: the hypotheses of
the stage-$4$ MC4 frontier package already include the Virasoro Ward
identity for the visible generators, and this alone forces the
$T$-channel coefficient at the leading $T$-pole to be~$2$.
\end{remark}

\subsection{Route 3: explicit DS computation and computational verification}

The principal Drinfeld--Sokolov reduction $\cW_4 = W(\fsl_4, f_{\mathrm{prin}})$
provides an explicit, independent confirmation.  In the DS construction with
Zamolodchikov normalisation, the six stage-$4$ OPE coefficients are:
\begin{align*}
c_{334}^{\,2}
&= \frac{42\,c^2(5c+22)}{(c+24)(7c+68)(3c+46)},
\\[3pt]
c_{444}^{\,2}
&= \frac{112\,c^2(2c-1)(3c+46)}{(c+24)(7c+68)(10c+197)(5c+3)},
\\[3pt]
\mathsf{C}_{3,4;3;0,4}^{\,2}
&= \tfrac{9}{16}\,c_{334}^{\,2},
\qquad
\mathsf{C}_{3,4;4;0,3}^{\,2}
= \tfrac{5}{7}\,c_{334}^{\,2},
\\[3pt]
\mathsf{C}^{\mathrm{DS}}_{4,4;2;0,6}
&= 2,
\qquad
\mathsf{C}^{\mathrm{DS}}_{3,4;2;0,5}
= 0.
\end{align*}
The first four are the higher-spin targets (Hornfeck 1993,
Blumenhagen--Eholzer--Flohr--Hornfeck--H"ubel 1996); the last two are
the theorematic Virasoro-target identities from the universal $T$-coupling
(Proposition~\ref{prop:universal-T-coefficient}) and mixed-weight orthogonality.

The computational verification suite
(\texttt{compute/lib/w4\_ds\_ope\_extraction.py}, $128$~tests) confirms:
\begin{itemize}
\item all four squared structure constants evaluated at levels
  $k=0,1,2,-2,-3$ (avoiding denominatorial poles);
\item inter-coefficient ratios
  $\mathsf{C}_{3,4;3;0,4}^2/c_{334}^2 = 9/16$ and
  $\mathsf{C}_{3,4;4;0,3}^2/c_{334}^2 = 5/7$ verified symbolically;
\item Feigin--Frenkel duality $k\leftrightarrow -k-2h^\vee$ checked:
  $c(k)+c(-k-8)=246$;
\item unitarity non-negativity ($c_{334}^2\ge 0$, $c_{444}^2\ge 0$)
  confirmed at all $\cW(N,N{+}1)$ minimal models for $4\le N\le 19$;
\item both theorematic Virasoro-target values
  $\mathsf{C}^{\mathrm{DS}}_{4,4;2;0,6}=2$ and
  $\mathsf{C}^{\mathrm{DS}}_{3,4;2;0,5}=0$ confirmed.
\end{itemize}

\subsection{Summary and consequences}

\begin{theorem}[The critical scalar identity]
\label{thm:critical-scalar-verified}
The identity
\[
\mathsf{C}^{\mathrm{res}}_{4,4;2;0,6}(4)=2
\]
holds as a theorem of conformal symmetry.  More precisely:
\begin{enumerate}[label=\textup{(\roman*)}]
\item The conformal Ward identity for
  $\langle T(x)\,W^{(4)}(z)\,W^{(4)}(0)\rangle$, under the Zamolodchikov
  normalisation, forces the coefficient $\mathsf{C}_{4,4;2;0,6}=2$
  universally for any vertex algebra with a Virasoro-primary spin-$4$
  generator.
\item The identity is logically equivalent to the standard two-point
  normalisation $\langle W^{(4)}(z)\,W^{(4)}(0)\rangle = (c/4)\,z^{-8}$.
\item The identity is a special case of the universal law
  $\mathsf{C}_{s,s;2;0,2s-2}=2$ valid for all conformal spins $s\ge 2$.
\item The principal Drinfeld--Sokolov computation confirms it independently
  with $128$ computational tests.
\end{enumerate}
\end{theorem}

\begin{remark}[Status within the MC4 frontier]
With this identity verified, the stage-$4$ Ward-inheritance conjecture
(Conjecture~\ref*{conj:winfty-stage4-ward-inheritance} of Vol~I) is
resolved: the six-entry stage-$4$ identity packet contracts to the four
higher-spin channels
\[
(3,3;4;0,2),\quad
(4,4;4;0,4),\quad
(3,4;3;0,4),\quad
(3,4;4;0,3).
\]
Combined with the visible pairing reduction and Borcherds transport, the
remaining $\cW_\infty$ stage-$4$ problem consists of the two primitive
self-coupling square classes $c_{334}^2$ and $c_{444}^2$.  The $\cW_4$
rigidity theorem (Zamolodchikov--Fateev--Lukyanov) then closes all
stage-$4$ identities, resolving MC4 at this stage.  The true remaining
frontier for $\cW_\infty$ promotion is $N\ge 5$.
\end{remark}



%% ===================================================================
%%  CELESTIAL AMPLITUDES FROM THE BAR COMPLEX
%% ===================================================================

\section{Celestial amplitudes from the bar complex: gluon and graviton OPE}
\label{sec:celestial-bar-complex}

The celestial holography programme maps $4$-dimensional scattering
amplitudes to $2$-dimensional celestial CFT correlators via Mellin
transform.  The boundary chiral algebras of the holomorphic-topological
framework --- affine Kac--Moody for gauge theory, Virasoro for gravity
--- should therefore encode celestial OPE data through their bar
complexes.  This section extracts the celestial OPE coefficients for
gluons and gravitons directly from the collision residues of the bar
differential, computes the four-point MHV amplitude as a bar-complex
integral, and identifies a sharp helicity-chirality obstruction that
determines what the bar complex can and cannot see.

\subsection{Setup: the collision residue as celestial OPE}
\label{subsec:celestial-collision-setup}

Recall the architecture (Definition~\ref{def:chiral_bar},
Proposition~\ref{prop:field-theory-r}).  The chiral bar complex
$\barB^{\mathrm{ch}}(\cA)$ carries a differential
$d_{\barB} = d_Q + d_{\mathrm{res}} + d_{A_\infty}$, where
$d_{\mathrm{res}}$ extracts residues at the collision divisors
$D_S \subset \overline{\mathrm{FM}}_n(\C)$.  At arity $2$, the
collision divisor $D_{\{1,2\}}$ is a point, and the residue is
computed by integrating the bar weight form $\omega_2$ against
$d\log(z_1 - z_2)$.

\begin{observation}[The $d\log$ absorption principle]
\label{obs:dlog-absorption-cel}
The collision residue extracts the $r$-matrix
$r(z) = \mathrm{Res}^{\mathrm{coll}}_{0,2}(\Theta^{\mathrm{oc}})$
by integrating against the kernel $d\log(z_1 - z_2)$, which absorbs
one power of the spectral variable:
\[
z^{-n} \;\xmapsto{d\log}\; z^{-(n-1)}.
\]
Hence the OPE singularity $\sim z^{-N}$ produces an $r$-matrix
pole $\sim z^{-(N-1)}$.  This is the content of AP19.
\end{observation}

\begin{definition}[Celestial conformal primary]
\label{def:celestial-primary-cel}
Given a $4$d massless particle of helicity $s$ and $4$-momentum
$p^\mu$, the celestial conformal primary of conformal dimension
$\Delta$ is the Mellin transform
\[
\cO^s_\Delta(z,\bar z) \;:=\;
\int_0^\infty d\omega\; \omega^{\Delta - 1}\;
a^s(p^\mu(\omega, z, \bar z)),
\]
where $p^\mu = \omega\, q^\mu(z,\bar z)$ parametrises the null
momentum in terms of the celestial sphere coordinate $(z,\bar z)$ and
energy $\omega$.  For positive helicity ($s > 0$), the celestial
primary is holomorphic in $z$; for negative helicity ($s < 0$), it is
anti-holomorphic.
\end{definition}

\subsection{Yang--Mills: gluon celestial OPE from $B(\widehat{\fg}_k)$}
\label{subsec:gluon-celestial-ope-cel}

The holomorphic-topological twist of $4$d $\cN = 2$ Yang--Mills with
gauge algebra $\fg$ on $\C \times \R_{\ge 0}$ produces a boundary
chiral algebra that is the affine Kac--Moody algebra
$\widehat{\fg}_k$ at the level determined by the Chern--Simons
coupling.  The positive-helicity gluon maps, under the celestial
dictionary, to the affine current $J^a(z)$ at the boundary.

The $\lambda$-bracket of $\widehat{\fg}_k$ is
$\{J^a {}_\lambda J^b\} = f^{ab}{}_c\, J^c + k\, \kappa^{ab}\, \lambda$.
The Laplace transform gives
$r^{\mathrm{Lap}}(z) = f^{ab}{}_c\, J^c / z + k\, \kappa^{ab} / z^2$.
After $d\log$ absorption (Observation~\ref{obs:dlog-absorption-cel}),
the $z^{-2}$ level term maps to $z^{-1}$ and the $z^{-1}$
structure-constant term drops to $z^0$.  The bar-theoretic $r$-matrix is
\[
r^{\mathrm{KM}}(z)
\;=\; \frac{k\, \Omega}{z}
\;=\; \frac{k\, \kappa^{ab}\, t_a \otimes t_b}{z}.
\]

\begin{theorem}[Celestial gluon OPE from the bar complex]
\label{thm:celestial-gluon-ope-cel}
Let $\cO^{+,a}_{\Delta}(z)$ denote the celestial primary of a
positive-helicity gluon of colour $a$ and conformal dimension $\Delta$.
The celestial OPE extracted from the collision residue of
$\barB^{\mathrm{ch}}(\widehat{\fg}_k)$ is
\[
\cO^{+,a}_{\Delta_1}(z_1)\,\cO^{+,b}_{\Delta_2}(z_2)
\;\sim\;
\frac{f^{ab}{}_c}{z_1 - z_2}\;
\cO^{+,c}_{\Delta_1 + \Delta_2 - 1}(z_2)
\;+\; \textup{regular},
\]
with the OPE coefficient being the structure constant $f^{ab}{}_c$
of the gauge Lie algebra~$\fg$.  This matches the Fan--Fotopoulos--Taylor
celestial gluon OPE obtained by direct Mellin transform of the
Parke--Taylor amplitude; the novelty is that $f^{ab}{}_c$ is extracted
from the collision residue of the bar complex, not inserted by hand.
\end{theorem}

\begin{proof}
The collision $r$-matrix $r^{\mathrm{KM}}(z) = k\,\Omega/z$ acts on
the colour indices via the Casimir.  In the colour-ordered projection,
$\langle c| \,\Omega\, |a,b\rangle = f^{ab}{}_c$, producing the
simple-pole celestial OPE $\sim f^{ab}{}_c / (z_1 - z_2)$.
The conformal-dimension shift $\Delta_1 + \Delta_2 - 1$ arises from
the Mellin convolution beta-function integral; the level $k$ drops out
because it is absorbed into the normalisation of the celestial primaries.
\end{proof}


\subsection{Gravity: graviton celestial OPE from $B(\mathrm{Vir}_c)$}
\label{subsec:graviton-celestial-ope-cel}

The holomorphic-topological twist of $3$d gravity on
$\C \times \R_{\ge 0}$ produces a boundary Virasoro algebra
$\mathrm{Vir}_c$.  The positive-helicity graviton maps to the stress
tensor $T(z)$ at the boundary.  The Virasoro $\lambda$-bracket
$\{T {}_\lambda T\} = (\partial + 2\lambda)\, T + \tfrac{c}{12}\,\lambda^3$
gives the Laplace $r$-matrix
$r^{\mathrm{Lap}}(z) = \partial T/z + 2T/z^2 + (c/2)/z^4$.
After $d\log$ absorption (AP19), the collision $r$-matrix is
\[
r^{\mathrm{Vir}}(z)
\;=\;
\frac{c/2}{z^3} \;+\; \frac{2T}{z}.
\]

\begin{theorem}[Celestial graviton OPE from the bar complex]
\label{thm:celestial-graviton-ope-cel}
The celestial OPE extracted from the collision residue of
$\barB^{\mathrm{ch}}(\mathrm{Vir}_c)$ is
\[
\cO^{+}_{\Delta_1}(z_1)\,\cO^{+}_{\Delta_2}(z_2)
\;\sim\;
\frac{\Delta_2 - 1}{z_1 - z_2}\;
\cO^{+}_{\Delta_1 + \Delta_2 - 1}(z_2)
\;+\;
\frac{c/2}{(z_1 - z_2)^3}\;
\cO^{+}_{\Delta_1 + \Delta_2 - 3}(z_2)
\;+\; \textup{regular},
\]
where: $r_1 = 2T$ gives the leading soft-graviton (Weinberg) term with
coefficient $(\Delta_2 - 1)$; $r_3 = c/2$ gives the subleading
gravitational-anomaly term.  This two-pole structure matches
Pasterski--Shao--Strominger and is a theorem of the bar complex.
\end{theorem}

\begin{proof}
The leading pole $r_1 = 2T$ acts on $\cO^+_{\Delta_2}$ via the Virasoro
zero mode, producing the shifted coefficient $(\Delta_2 - 1)$.  The
subleading pole $r_3 = c/2$ is a $c$-number producing the triple-pole
contribution with output dimension $\Delta_1 + \Delta_2 - 3$.  The
absence of a $z^{-2}$ pole reflects the $d\log$ absorption of the
$\partial T$ term.  No higher poles exist because the Virasoro
$\lambda$-bracket is cubic.
\end{proof}


\subsection{The four-point MHV amplitude from the bar complex}
\label{subsec:mhv-bar-complex-cel}

\begin{theorem}[MHV bar-complex integral at $n = 4$]
\label{thm:mhv-bar-integral-cel}
The colour-ordered tree-level four-gluon MHV amplitude is reproduced by
the bar-complex integral
\[
A_4^{\mathrm{bar}}(1^+, 2^-, 3^-, 4^+)
\;=\;
\int_{\overline{\mathrm{FM}}_4(\C)}
\omega_4^{(+,-,-,+)} \wedge \nu_4,
\]
which localises via residue extraction on the boundary divisors of
$\overline{\mathrm{FM}}_4(\C)$ to a sum over channels and reproduces
the Parke--Taylor formula
$A_4^{\mathrm{PT}} = \langle 23\rangle^3 /
(\langle 12\rangle\,\langle 34\rangle\,\langle 41\rangle)$.
The holomorphic $d\log$ residues supply the cyclic denominator; the
anti-holomorphic insertions supply the $\langle 23\rangle^4$ numerator.
\end{theorem}

\begin{proof}
The bar-complex integrand is constructed as follows.

\emph{Step~1: Arity-$4$ bar element.}
For $\widehat{\fg}_k$ with $m_{k \ge 3} = 0$ (strict associativity),
the bar complex at arity~$4$ has basis elements
$[J^{a_1}|J^{a_2}|J^{a_3}|J^{a_4}]$ parametrised by
$\overline{\mathrm{FM}}_4(\C)$.  The bar weight form $\omega_4$
is the product of $d\log$ forms along the collision divisors.

\emph{Step~2: Residue localisation.}
The boundary divisors of $\overline{\mathrm{FM}}_4(\C)$ correspond
to the three channels $(12|34)$, $(23|41)$, $(14|23)$.
On each channel, the residue of $\omega_4$ factors as a product of
two arity-$2$ collision residues.  For the channel $(12|34)$:
\[
\Res_{D_{\{1,2\}}} \omega_4
\;=\;
\frac{k^2\,\Omega_{12}\,\Omega_{34}}
{(z_1 - z_2)(z_3 - z_4)},
\]
with analogous expressions for the other channels.  Summing over
channels and cyclic orderings of the colour trace produces the
colour-ordered partial amplitude
\[
A(1,2,3,4)
\;=\;
\frac{1}{(z_1 - z_2)(z_2 - z_3)(z_3 - z_4)(z_4 - z_1)}.
\]

\emph{Step~3: Helicity assignment and Parke--Taylor.}
The MHV helicity configuration $(+,-,-,+)$ is implemented by
contracting the bar element with the twisting morphism
$\tau \colon \barB(\cA) \to \cA^!$.  The two negative-helicity
legs (legs $2$ and $3$) are paired with the Koszul dual via
$\tau$, producing the anti-holomorphic factor
$\langle 23 \rangle^4$.  The identification of $z_i$ with
holomorphic spinors $\lambda_i$ converts the cyclic denominator
to $\langle 12\rangle\langle 23\rangle\langle 34\rangle
\langle 41\rangle$.  The result is the Parke--Taylor formula.
\end{proof}

\begin{remark}[Why $n = 4$ is special]
At $n = 4$, $\overline{\mathrm{FM}}_4(\C) \cong \bP^1$
(the cross-ratio compactification), and the integral localises
completely by residues.  For $n > 4$, the FM space has
dimension $n - 3 > 1$ and the integral receives contributions
from internal moduli, producing the full colour-ordered
$n$-point partial amplitude.  For affine KM (strict
associativity), the result is always the Parke--Taylor
cyclic denominator $\prod_{i=1}^n (z_i - z_{i+1})^{-1}$.
\end{remark}

\subsection{Helicity splitting: what the bar complex sees and what it does not}
\label{subsec:helicity-splitting-cel}

\begin{theorem}[Helicity-chirality correspondence]
\label{thm:helicity-chirality-cel}
In the holomorphic-topological framework:
\begin{enumerate}[label=\textup{(\roman*)}]
\item $\barB^{\mathrm{ch}}(\cA)$ encodes exactly the self-dual
  (positive-helicity) sector;
\item $\barB^{\mathrm{ch}}(\cA^!)$ encodes the anti-self-dual
  (negative-helicity) sector;
\item mixed-helicity amplitudes (MHV) arise from the twisting-morphism
  pairing $\tau \colon \barB^{\mathrm{ch}}(\cA) \to \cA^!$;
\item all-plus and one-minus tree amplitudes vanish (form-type mismatch
  on $\overline{\mathrm{FM}}_n(\C)$);
\item MHV amplitudes are the simplest non-vanishing pairing between
  $\barB^{\mathrm{ch}}(\cA)$ and $\barB^{\mathrm{ch}}(\cA^!)$.
\end{enumerate}
Helicity is chirality; helicity splitting is the Swiss-cheese
decomposition.
\end{theorem}

\begin{proof}
The argument proceeds by matching the chirality structure of the
$\SCchtop$-operad with the helicity structure of $4$d massless
scattering.

\emph{(i) Self-dual sector.}
The closed colour of $\SCchtop$ is parametrised by
$\FM_k(\C)$, which carries holomorphic forms only.  The bar
differential $d_{\mathrm{res}}$ extracts residues of
holomorphic collision divisors.  By the twistor correspondence,
holomorphic data on the celestial sphere encode
positive-helicity (self-dual) particles.  Hence
$\barB^{\mathrm{ch}}(\cA)$ encodes exactly the SD sector.

\emph{(ii) Anti-self-dual sector.}
The Koszul dual $\cA^!$ is the boundary algebra on the opposite
(anti-holomorphic) chirality.  Its bar complex
$\barB^{\mathrm{ch}}(\cA^!)$ is parametrised by collision
residues in the $\cA^!$-OPE, which encode negative-helicity
(ASD) particles.  For Virasoro:
$\cA^! = \mathrm{Vir}_{26-c}$, and the ASD graviton lives in
$\barB(\mathrm{Vir}_{26-c})$.  For affine KM:
$\cA^! = \widehat{\fg}_{-k-2h^\vee}$, and the ASD gluon lives
in $\barB(\widehat{\fg}_{-k-2h^\vee})$.

\emph{(iii) Mixed helicity.}
An MHV amplitude with $2$ negative-helicity and $(n-2)$
positive-helicity legs requires a pairing between the SD bar
complex and the ASD bar complex.  This pairing is the universal
twisting morphism $\tau \colon \barB(\cA) \to \cA^!$, which is
the bar-cobar adjunction counit.  The two negative-helicity
insertions enter through $\cA^!$; the $(n-2)$ positive-helicity
insertions enter through $\barB(\cA)$.

\emph{(iv) Vanishing of all-plus.}
An $n$-point all-positive-helicity amplitude requires
$n$ holomorphic insertions with no anti-holomorphic factor.  But
the bar complex integral over $\FM_n(\C)$ has form degree
$n - 3$ (the real dimension of $\FM_n(\C)$ modulo translations
and dilations), while $n$ holomorphic insertions with only
$d\log$ kernels produce form degree $n - 1$.  The mismatch
(excess form degree) forces the integral to vanish.  The
one-minus case fails similarly: one anti-holomorphic factor
provides insufficient form degree to saturate the fiber integral.

\emph{(v) MHV as minimal pairing.}
The MHV configuration (exactly $2$ negative helicities) provides
exactly the right form degree: the $2$ anti-holomorphic
insertions via $\tau$ contribute $2$ units of anti-holomorphic
form, and the remaining $(n-2)$ holomorphic insertions contribute
$(n-2) - 1 = n - 3$ units of holomorphic form via the $d\log$
kernels.  The total form degree matches $\dim_\R \FM_n(\C) = 2(n-3)$,
and the integral is non-degenerate.
\end{proof}


\subsection{Summary: the celestial-bar dictionary}
\label{subsec:celestial-bar-dictionary-cel}

\medskip
\begin{center}
\small
\renewcommand{\arraystretch}{1.4}
\begin{tabular}{p{3.5cm}p{3.5cm}p{4.5cm}}
\textbf{Celestial datum} &
\textbf{Bar-complex object} &
\textbf{Formula} \\
\hline
Gluon OPE coefficient &
Collision residue of
  $\barB(\widehat{\fg}_k)$, arity $2$ &
$f^{ab}{}_c / z$
  (from $r^{\mathrm{KM}} = k\,\Omega/z$) \\[4pt]
Graviton leading OPE &
$r_1 = 2T$ pole of
  $r^{\mathrm{Vir}}(z)$ &
$(\Delta_2 - 1) / (z_1 - z_2)$ \\[4pt]
Graviton subleading OPE &
$r_3 = c/2$ pole of
  $r^{\mathrm{Vir}}(z)$ &
$(c/2) / (z_1 - z_2)^3$ \\[4pt]
MHV amplitude ($n = 4$) &
$\int_{\overline{\mathrm{FM}}_4}
  \omega_4 \wedge \nu_4$ &
$\langle 23\rangle^4 / (\langle 12\rangle
  \langle 23\rangle \langle 34\rangle
  \langle 41\rangle)$ \\[4pt]
Self-dual sector &
$\barB^{\mathrm{ch}}(\cA)$ &
Positive helicity \\[4pt]
Anti-self-dual sector &
$\barB^{\mathrm{ch}}(\cA^!)$ &
Negative helicity \\[4pt]
Mixed helicity (MHV) &
Twisting morphism
  $\tau: \barB(\cA) \to \cA^!$ &
Bulk propagator pairing
\end{tabular}
\end{center}

\begin{question}[Higher $n$ and NMHV]
\label{q:higher-mhv-cel}
For NMHV amplitudes (three negative helicities), the pairing involves a
double twisting morphism.  Is the BCFW recursion a statement about
the $\Ainf$ structure of the twisting morphism space?
\end{question}




% =================================================================
\section{Virasoro spectral $R$-matrix: Laplace transform and quantum corrections}
\label{sec:virasoro-spectral-R}
% =================================================================

The spectral $R$-matrix of the Virasoro algebra encodes the
braiding of gravitational line operators.  Its classical limit is
the Laplace transform of the $\lambda$-bracket; its quantum
completion is conjectured to equal the Ponsot--Teschner fusion
kernel.  This section performs the explicit computation,
including the first quantum correction, which has not previously
appeared in the literature.

\subsection{Setup and conventions}

The Virasoro $\lambda$-bracket is
\begin{equation}\label{eq:vir-lambda-R}
\{T{}_\lambda T\}
\;=\;
\partial T + 2T\lambda + \frac{c}{12}\,\lambda^3.
\end{equation}
The mode algebra is $[L_m, L_n] = (m-n)L_{m+n} +
\frac{c}{12}(m^3 - m)\,\delta_{m+n,0}$, with stress tensor
$T(z) = \sum_{n \in \Z} L_n\, z^{-n-2}$.  The perturbative
expansion parameter is $\hbar = 12/c$ (the ``gravitational
coupling'': $\hbar \to 0$ is the semiclassical limit $c \to \infty$).

\begin{warning}[AP19: pole absorption]
The \emph{Laplace kernel} $r^L(z)$ has the same pole orders as the
OPE.  The \emph{collision residue} $r^{\mathrm{coll}}(z)$ has pole
orders one lower, because $d\log(z_{12})$ absorbs one power.  Both
are correct in their respective contexts; confusion between them is
the content of AP19.  The primary object here is
$r^L(z)$ (the Laplace transform output), which is the object that
enters the shifted KZ connection and the perturbative expansion of
$R(z)$.
\end{warning}

\subsection{The classical $r$-matrix: computation}

\begin{proposition}[Virasoro Laplace kernel; proved]
\label{prop:vir-laplace-R}
The Laplace transform of the Virasoro $\lambda$-bracket is
\begin{equation}\label{eq:vir-rL-R}
r^L(z)
\;=\;
\int_0^\infty e^{-\lambda z}\,\{T{}_\lambda T\}\,d\lambda
\;=\;
\frac{\partial T}{z}
+ \frac{2T}{z^2}
+ \frac{c/2}{z^4}.
\end{equation}
The corresponding collision residue (bar-theoretic $r$-matrix) is
\begin{equation}\label{eq:vir-rcoll-R}
r^{\mathrm{coll}}(z)
\;=\;
\frac{c/2}{z^3}
+ \frac{2T}{z}.
\end{equation}
\end{proposition}

\begin{proof}
Term by term, using $\int_0^\infty \lambda^n\,e^{-\lambda z}\,d\lambda
= n!/z^{n+1}$:
\begin{align*}
\int_0^\infty e^{-\lambda z}\,\partial T\,d\lambda
&\;=\;
\partial T \cdot \frac{0!}{z^1}
\;=\;
\frac{\partial T}{z},\\
\int_0^\infty e^{-\lambda z}\,2T\lambda\,d\lambda
&\;=\;
2T \cdot \frac{1!}{z^2}
\;=\;
\frac{2T}{z^2},\\
\int_0^\infty e^{-\lambda z}\,\frac{c}{12}\lambda^3\,d\lambda
&\;=\;
\frac{c}{12}\cdot\frac{3!}{z^4}
\;=\;
\frac{c}{2}\cdot\frac{1}{z^4}.
\end{align*}
Sum: $r^L(z) = \partial T/z + 2T/z^2 + (c/2)/z^4$.

For the collision residue: the $d\log$ kernel absorbs one power of~$z$,
so $z^{-k} \mapsto z^{-(k-1)}$, giving
$r^{\mathrm{coll}}(z) = (c/2)/z^3 + 2T/z$.
The $\partial T/z$ term becomes $\partial T/z^0$, which is regular
and drops from the collision residue (it contributes to the
translation operator $T_z$, not to the pole structure).
\end{proof}

\subsection{Mode expansion of the classical $r$-matrix}

Acting on
$V_{h_1} \otimes V_{h_2}$ (Verma modules of highest weight $h_1, h_2$
for the Koszul-dual algebra $\mathrm{Vir}_{26-c}$), the classical
$r$-matrix has the mode decomposition:
\begin{equation}\label{eq:rL-mode-expansion-R}
r^L_{12}(z)
\;=\;
\sum_{m \ge 1}\,
\frac{L_{-m}^{(1)} \otimes L_m^{(2)}
\;-\;
L_m^{(1)} \otimes L_{-m}^{(2)}}{z^{m+1}}
\;+\;
\frac{c/2}{z^4}\,\id \otimes \id
\;+\;
\text{(diagonal terms)}.
\end{equation}

\subsection{Classical Yang--Baxter equation: explicit verification}
\label{subsec:cybe-vir-R}

\begin{theorem}[Gravitational CYBE; proved]
\label{thm:vir-cybe-R}
The Virasoro Laplace kernel $r^L(z)$ satisfies the classical
Yang--Baxter equation
\begin{equation}\label{eq:cybe-R}
[r^L_{12}(u),\, r^L_{13}(u+v)]
+ [r^L_{12}(u),\, r^L_{23}(v)]
+ [r^L_{13}(u+v),\, r^L_{23}(v)]
\;=\; 0
\end{equation}
in $\End(V_{h_1} \otimes V_{h_2} \otimes V_{h_3})$ for generic
highest weights $h_1, h_2, h_3$.
\end{theorem}

\begin{proof}
The CYBE is the Laplace-transformed form of the PVA
Jacobi identity
$\{T{}_\lambda\{T{}_\mu T\}\}
- \{T{}_\mu\{T{}_\lambda T\}\}
- \{\{T{}_\lambda T\}{}_{\lambda+\mu} T\}
= 0$.
We verify this directly.

\emph{Term 1:} Using sesquilinearity,
$\{T{}_\lambda\{T{}_\mu T\}\}
= \partial^2 T + (3\lambda + 2\mu)\,\partial T
+ 2\lambda(\lambda + 2\mu)\,T
+ \tfrac{c}{12}\lambda^3(\lambda + 2\mu)$.

\emph{Term 2:} Swapping $\lambda \leftrightarrow \mu$:
$\{T{}_\mu\{T{}_\lambda T\}\}
= \partial^2 T + (3\mu + 2\lambda)\,\partial T
+ 2\mu(\mu + 2\lambda)\,T
+ \tfrac{c}{12}\mu^3(\mu + 2\lambda)$.

\emph{Term 3:}
$\{\{T{}_\lambda T\}{}_{\lambda+\mu}\, T\}
= (\lambda - \mu)\,\partial T + 2(\lambda^2 - \mu^2)\,T
+ \tfrac{c}{12}(\lambda - \mu)(\lambda+\mu)^3$.

Computing Term~1 $-$ Term~2:
\begin{align*}
&\bigl[(3\lambda+2\mu) - (2\lambda + 3\mu)\bigr]\,\partial T
= (\lambda - \mu)\,\partial T,\\
&\bigl[2\lambda^2 + 4\lambda\mu - 2\mu^2 - 4\lambda\mu\bigr]\,T
= 2(\lambda^2 - \mu^2)\,T,\\
&\tfrac{c}{12}\bigl[\lambda^4 + 2\lambda^3\mu - \mu^4 - 2\mu^3\lambda\bigr]
= \tfrac{c}{12}(\lambda-\mu)(\lambda+\mu)^3.
\end{align*}
This equals Term~3.  Under the Laplace
transform $\lambda \mapsto u^{-1}$, $\mu \mapsto v^{-1}$,
this becomes equation~\eqref{eq:cybe-R}.
\end{proof}

\begin{remark}[Arnold interpretation]
\label{rem:arnold-vir-R}
The cancellation in the Jacobi identity is the algebraic content
of the Arnold relation on $\FM_3(\C)$: the three codimension-$1$
boundary strata contribute the three commutator terms, and Stokes'
theorem forces their sum to vanish --- the same mechanism that
proves $d_{\barB}^2 = 0$ for the bar differential.
\end{remark}

\subsection{The quantum $R$-matrix: perturbative expansion}
\label{subsec:quantum-R-sec}

The quantum $R$-matrix is the monodromy of the shifted KZ connection.
The first quantum correction has the form
\begin{equation}\label{eq:R2-R}
r^{(2)}(z)
\;=\;
\frac{1}{2}\oint_{\gamma(z)}
r^L_{12}(z-w)\,r^L_{12}(w)\,dw
\;+\;
\Delta r^{(2)}_{\Ainf}(z),
\end{equation}
where the first term is iterated monodromy and the second is the
$m_3$ correction from the $\Ainf$ bar differential.  The two
partially cancel via the Stasheff identity
$\partial m_3 + m_2 \circ m_2 = 0$.

\begin{proposition}[Quantum correction at level~$0$; new computation]
\label{prop:quantum-level0-R}
On the vacuum-to-vacuum component
$|h_1\rangle \otimes |h_2\rangle \to |h_1\rangle \otimes |h_2\rangle$,
the quantum $R$-matrix has the diagonal expansion:
\begin{equation}\label{eq:R-diagonal-level0-R}
R(z)\big|_{|h_1,h_2\rangle \to |h_1,h_2\rangle}
\;=\;
\exp\!\left(
  \frac{c/2}{z^4}\hbar
  + \frac{2h_1 h_2}{z^4}\hbar^2
  + O(\hbar^3)
\right),
\end{equation}
where the first quantum correction
$2h_1 h_2 \hbar^2/z^4$ matches the leading term of the
Ponsot--Teschner kernel in the heavy limit
$h_i = c\epsilon_i/6$.
\end{proposition}

\subsection{The degenerate case: exact $R$-matrix}
\label{subsec:degenerate-exact-R}

\begin{proposition}[Degenerate $R$-matrix from BPZ monodromy; proved]
\label{prop:degenerate-from-laplace-R}
On the $(2,1)$-degenerate representation $V_{h_{2,1}}$, the
$\Ainf$ corrections vanish (null vectors annihilate higher
operations) and the spectral $R$-matrix reduces to BPZ monodromy
with exact eigenvalues $e^{\pm i\pi b\alpha_2}$, matching
Ponsot--Teschner.
\end{proposition}

\subsection{Structure and conjecture}
\label{subsec:structure-conj-R}

\begin{theorem}[Structure theorem; partially proved]
\label{thm:vir-R-structure-R}
The Virasoro spectral $R$-matrix has: (i)~classical limit = Laplace
kernel (proved), (ii)~quantum corrections from iterated monodromy $+$
$\Ainf$ (explicit at level~$0$), (iii)~degenerate truncation via BPZ
(proved), (iv)~diagonal form with Aharonov--Bohm eigenvalues (proved),
(v)~heavy limit matching AdS$_3$ geodesic length (subleading order).
\end{theorem}

\begin{conjecture}[Virasoro $R$-matrix $=$ Ponsot--Teschner]
\label{conj:R-PT-R}
The full $R$-matrix equals the Ponsot--Teschner fusion kernel at all
orders, with six matching conditions proved.  The remaining gap is the
analytic continuation from integer~$k$ to continuous~$c$.
\end{conjecture}




%% ============================================================
%%
%%  GENUS-1 MODULAR ENVELOPE: THE FIRST THREE OPERATIONS
%%
%%  Computed April 2026
%%
%% ============================================================

\section{Genus-1 modular envelope: the first three operations}
\label{sec:genus1-modular-envelope}

\subsection{Setup and notation}

The modular Swiss-cheese operad $\SCmod$ has closed-color operations
$\cE^{\mathrm{cl}}(g,n)$ for each genus $g \geq 0$ and arity $n \geq 0$
(subject to $2g - 2 + n > 0$).  The genus-$0$ package is proved: the
$\SCchtop$-algebra structure on a logarithmic chiral algebra $\cA$ is
the bar-cobar adjunction of Volume~I Theorem~A.

The \emph{modular envelope} extends the genus-$0$ package to all genera
via the relative Feynman transform $\FT_{\Com_{\mathrm{mod}}/\SCchtop}$.
The genus-$1$ \emph{seed} consists of the first three operations:
\begin{align*}
\mu_{1,0} &\colon \mathbb{k} \longrightarrow \mathbb{k},
\qquad\text{(vacuum amplitude, arity $0$)},\\
\mu_{1,1} &\colon \cA \longrightarrow \mathbb{k},
\qquad\text{(one-point function, arity $1$)},\\
\mu_{1,2} &\colon \cA^{\otimes 2} \longrightarrow \mathbb{k},
\qquad\text{(two-point function, arity $2$)}.
\end{align*}
These are the genus-$1$ components $\Theta^{(1)}_{n=0}$,
$\Theta^{(1)}_{n=1}$, $\Theta^{(1)}_{n=2}$ of the modular
Maurer--Cartan element, computed from the genus-raising operator
$D_1 = D_{\mathrm{nsep}}$ acting on the genus-$0$ MC class
$\Theta^{(0)}$.

\subsection{The obstruction equation at genus $1$}

Fix $\cA$ a logarithmic $\SCchtop$-algebra with genus-$0$ MC
element $\Theta^{(0)} \in \MC(\gr_F^0 L_{\mathrm{mod}})$.  The
universal genus-$1$ obstruction (Theorem~\ref{thm:Ob1}) gives:
\begin{equation}\label{eq:wn-genus1-obstruction}
d_0^{\Theta^{(0)}} \Theta^{(1)}
\;=\;
-D_1(\Theta^{(0)})
\;=\;
-\Delta_{\mathrm{cyc}}(\Theta^{(0)}),
\end{equation}
where $\Delta_{\mathrm{cyc}}$ is the nonseparating self-sewing
operator (the reduced cyclic odd Laplacian of
Theorem~\ref{thm:D1formula}).

The right-hand side is the \emph{genus-$1$ forcing term}: the
genus-raising operator applied to the tree-level data.  The genus-$1$
MC element $\Theta^{(1)}$ is the unique (up to gauge) solution of this
inhomogeneous equation, provided the obstruction class
$\Ob_1(\Theta^{(0)}) = [D_1\Theta^{(0)}] \in
H^2(\gr_F^1 L, d_0^{\Theta^{(0)}})$ vanishes.

\subsection{The computation for $\cA = \mathrm{Vir}_c$}

We now specialise to the Virasoro algebra $\cA = \mathrm{Vir}_c$.
The essential input data:
\begin{itemize}[leftmargin=1.8em]
\item Genus-$0$ MC element:
  $\Theta^{(0)} = \{m_2, m_3, m_4, \ldots\}$, the full $A_\infty$
  structure from the OPE $T(z)\,T(w) \sim (c/2)/(z-w)^4 + 2T/(z-w)^2
  + \partial T/(z-w)$.
\item Modular characteristic: $\kappa(\mathrm{Vir}_c) = c/2$.
\item Cyclic inner product: $\langle L_{-n}, L_m \rangle =
  \frac{c}{12}(n^3 - n)\,\delta_{n,m}$ (the Shapovalov form).
\item The triangular vanishing theorem: $\Delta_{\mathrm{cyc}} P_\Pi = 0$
  in $\Xloc^1/\partial$ for the Virasoro PVA
  (Theorem~\ref{thm:triangular-vanishing}).
\end{itemize}

\subsubsection{$\mu_{1,0}$: the vacuum amplitude (torus partition function)}

The arity-$0$ genus-$1$ MC element lives on the unique stable graph
$\Gamma_{1,0}$ of genus~$1$ with no external legs: a single vertex
with one self-loop.  The nonseparating self-sewing operator contracts
the binary vertex $m_2$ over the invariant form:
\begin{equation}\label{eq:wn-mu10}
\mu_{1,0}
\;=\;
\Theta^{(1)}_{n=0}
\;=\;
\kappa(\mathrm{Vir}_c) \cdot \omega_1
\;=\;
\frac{c}{2}\,\omega_1
\;\in\; H^0(\overline{\mathcal{M}}_{1,1},\,\mathcal{L}_\kappa),
\end{equation}
where $\omega_1 = c_1(\lambda_{\mathrm{Hodge}}) \in
H^2(\overline{\mathcal{M}}_{1,1})$ is the Hodge class.

The \emph{genus-$1$ free energy} is the Hodge integral:
\[
F_1(\mathrm{Vir}_c)
\;=\;
\kappa \cdot \lambda_1^{\mathrm{FP}}
\;=\;
\frac{c}{2} \cdot \frac{1}{24}
\;=\;
\frac{c}{48},
\]
where $\lambda_1^{\mathrm{FP}} = \int_{\overline{\mathcal{M}}_{1,1}}
\lambda_1 = 1/24$ is Mumford's formula.  The scalar partition function is
\begin{equation}\label{eq:wn-partition}
Z_1(\mathrm{Vir}_c;\tau)
\;=\;
\eta(\tau)^{-\kappa}
\;=\;
\eta(\tau)^{-c/2}
\;=\;
q^{-c/48}\prod_{n=1}^{\infty}(1-q^n)^{-c/2},
\end{equation}
where $\eta(\tau) = q^{1/24}\prod_{n=1}^\infty(1-q^n)$ is the
Dedekind eta function and $q = e^{2\pi i\tau}$.

\begin{remark}[Bar complex origin]
The self-loop contraction sums over modes:
$\Theta^{(1)}_{n=0} = \sum_{n \geq 1} \langle e^h | m_2 | e_h \rangle
\cdot \omega_1 = \operatorname{Tr}(L_0) \cdot \omega_1 =
\kappa \cdot \omega_1$.
Each factor $\eta(\tau)^{-1}$ in the partition function corresponds
to one mode of the self-loop propagator on the torus.  The total
number of such modes is $\kappa = c/2$, the modular characteristic
of Vol~I Theorem~D.
\end{remark}

\subsubsection{$\mu_{1,1}(T)$: the genus-$1$ one-point function}

The arity-$1$ genus-$1$ element lives on the stable graph with one
vertex, one self-loop, and one external leg (a tadpole'').  The
self-sewing contracts two of the three flags of the ternary vertex $m_3$
(from the Virasoro $A_\infty$ structure) over the invariant form:
\begin{equation}\label{eq:wn-mu11}
\mu_{1,1}(T)
\;=\;
\Theta^{(1)}_{n=1}(T)
\;=\;
-\frac{c}{24}\,E_2(\tau)\cdot T,
\end{equation}
where $E_2(\tau) = 1 - 24\sum_{n \geq 1}\sigma_1(n)\,q^n$ is the
weight-$2$ quasi-modular Eisenstein series.

\begin{proof}[Derivation]
The self-sewing of $m_3$ contracts a pair of flags using the
Shapovalov form:
\[
D_1(m_3)_{n=1}(T)
\;=\;
\sum_{h} \langle e^h | m_3(e_h, T) \rangle
\;+\;
\sum_{h} \langle e^h | m_3(T, e_h) \rangle.
\]
The self-loop integral on the torus is
\[
\oint_A \oint_B \frac{T(w)}{(z-w)^2}\,\frac{dz\,dw}{2\pi i}
\;=\;
\sum_{n \geq 1} n\cdot q^n \cdot T
\;=\;
\frac{E_2(\tau)}{24}\,T,
\]
multiplied by the curvature factor $-\kappa = -c/2$ from the
invariant-form contraction.  Hence
$\Theta^{(1)}_{n=1}(T) = -(c/2)\cdot E_2(\tau)/(12)\cdot T =
-(c/24)\,E_2(\tau) \cdot T$.

Equivalently, this is the logarithmic derivative of the partition
function:
\[
\mu_{1,1}(T)
\;=\;
q\,\partial_q \log Z_1(\tau) \cdot T
\;=\;
-\frac{c}{24}\,E_2(\tau) \cdot T,
\]
since $q\,\partial_q\log\eta(\tau) = -E_2(\tau)/24$
and $Z_1 = \eta^{-c/2}$.
\end{proof}

\begin{remark}[Quasi-modularity as bar curvature]
The quasi-modular anomaly of $E_2$
($E_2(-1/\tau) = \tau^2 E_2(\tau) + 12\tau/(2\pi i)$)
is the genus-$1$ incarnation of the curved bar differential
$d_{\mathrm{fib}}^2 = \kappa \cdot \omega_1 \neq 0$.
The anomalous term is $\kappa^{-1}$ times the nonseparating
degeneration class.  The non-holomorphic completion
$\widehat{E}_2(\tau) = E_2(\tau) - 3/(\pi\,\mathrm{Im}\,\tau)$
is a genuine modular form of weight~$2$; the correction
$-3/(\pi\,\mathrm{Im}\,\tau)$ is the contribution from the
bulk (Koszul dual) sector.
\end{remark}

\subsubsection{$\mu_{1,2}(T,T)$: the genus-$1$ two-point function}

The arity-$2$ genus-$1$ element lives on stable graphs with
one vertex, one self-loop, and two external legs.  The MC coefficient
encodes the genus-$1$ connected two-point correlation function:
\begin{equation}\label{eq:wn-mu12}
\mu_{1,2}(T,T;\,z_{12})
\;=\;
\Theta^{(1)}_{n=2}(T,T;\,z_{12})
\;=\;
\frac{c}{2}\,\wp(z_{12};\tau)
\;+\;
2\,\langle T(z_1)\,T(z_2) \rangle_{g=1}^{\mathrm{conn}},
\end{equation}
where $\wp(z;\tau)$ is the Weierstrass $\wp$-function on the
torus $E_\tau = \mathbb{C}/(\mathbb{Z} + \mathbb{Z}\tau)$ and
$z_{12} = z_1 - z_2$.

More explicitly, the connected genus-$1$ two-point function is:
\begin{equation}\label{eq:wn-mu12-explicit}
\langle T(z_1)\,T(z_2) \rangle_{g=1}^{\mathrm{conn}}
\;=\;
\frac{c}{2}\left(\wp(z_{12};\tau) + \frac{E_2(\tau)}{12}\right),
\end{equation}
where $G(z;\tau) = -\log\vartheta_1(z|\tau)
+ \pi(\mathrm{Im}\,z)^2/\mathrm{Im}\,\tau + \mathrm{const}$
is the scalar Green's function on the torus, and the Weierstrass
$\wp$-function with the $E_2/12$ correction arises as
$\partial_z^2 G(z;\tau)$.  The term $E_2(\tau)/12$ is the
zero-mode subtraction ensuring $\int_0^1 G(z;\tau)\,dz = 0$.

Substituting:
\begin{equation}\label{eq:wn-mu12-full}
\mu_{1,2}(T,T;\,z_{12})
\;=\;
c\,\wp(z_{12};\tau) + \frac{c}{12}\,E_2(\tau)
+ 2\,G_2^T(z_1,z_2;\tau),
\end{equation}
where $G_2^T$ denotes descendant corrections at each conformal weight.

\begin{proof}[Derivation]
Two channels contribute:
\begin{enumerate}[label=\textup{(\roman*)}]
\item \emph{Self-sewing of $m_4$:} the quartic vertex $m_4(T,T,e_h,e^h)$
  is contracted over one pair of flags using the invariant form.  The
  result is the propagator-exchange integral on the torus:
  \[
  D_1(m_4)_{n=2}
  \;=\;
  \sum_h \langle e^h | m_4(T,T,e_h) \rangle
  \;=\;
  \frac{c}{2}\,\wp(z_{12};\tau) + \cdots.
  \]
  The Weierstrass $\wp$-function emerges as the genus-$1$ chiral Green's
  function: it is the unique doubly-periodic meromorphic function with
  a double pole $1/z^2$ at $z = 0$, the elliptic replacement of the
  genus-$0$ propagator $1/(z_1-z_2)^2$.

\item \emph{Composition $m_2 \circ m_2$ through the genus-$1$ propagator:}
  the product of two binary vertices joined by an internal edge
  propagating on the torus gives the descendant corrections:
  \[
  \sum_{n \geq 0} \langle T(z_1) | q^{L_0 - c/24} | T(z_2) \rangle
  \;=\;
  E_2(\tau)\cdot(T \otimes T) + \cdots.
  \]
\end{enumerate}
The total is~\eqref{eq:wn-mu12}.
\end{proof}

\subsection{Swiss-cheese structure of $\mu_{1,n}$}

We now describe the structural features of the genus-$1$ seed.

\subsubsection{Factorisation through the relative Feynman transform}

\begin{proposition}[Factorisation of $\mu_{1,n}$; computed here]
\label{prop:wn-mu1n-factorisation}
The genus-$1$ operations $\mu_{1,n}$ factor through the relative
Feynman transform $\FT_{\Com_{\mathrm{mod}}/\SCchtop}$ as follows.
Each $\mu_{1,n}$ is the genus-$1$ component of the modular
MC element $\Theta_\cA$, which is the universal twisting morphism
of the $\FT_{\mathrm{rel}}$-algebra structure on $\Bmod(\cA)$:
\[
\begin{tikzcd}[column sep=3em]
\Bmod^{(0)}(\cA)
\ar[r, "D_1"]
&
\Bmod^{(1)}(\cA)
\ar[r, "\mathrm{proj}_{n}"]
&
\cA^{\otimes n}
\end{tikzcd}
\]
where $D_1 = D_{\mathrm{nsep}}$ is the genus-raising differential
of the bicomplex $(D_0, D_1)$, and $\mathrm{proj}_n$ projects to
the genus-$1$, arity-$n$ corolla.
\end{proposition}

Concretely:
\begin{enumerate}[label=\textup{(\roman*)}]
\item $\mu_{1,0}$ is the trace of $D_1$ on the genus-$0$ bar
  complex: it picks out the self-sewing of $m_2$ at the arity-$0$
  corolla, producing the scalar curvature $\kappa \cdot \omega_1$.
\item $\mu_{1,1}$ is the composition $\mathrm{proj}_1 \circ D_1
  \circ \Theta^{(0)}_{n=3}$: the genus-raising differential applied
  to the ternary vertex $m_3$, projected to arity~$1$.
\item $\mu_{1,2}$ is the sum of $\mathrm{proj}_2 \circ D_1
  \circ \Theta^{(0)}_{n=4}$ (self-sewing of $m_4$) and the
  separating contribution $m_2 \circ_{(1)} m_2$ through the
  genus-$1$ propagator.
\end{enumerate}

\subsubsection{The two colors at genus $1$}

At genus $1$, the Swiss-cheese structure has:
\begin{itemize}[leftmargin=1.8em]
\item \emph{Closed color} (modular): the self-sewing
  $\xi_{ij} \colon \cE^{\mathrm{cl}}(0,I) \to
  \cE^{\mathrm{cl}}(1, I \setminus \{i,j\})$
  is the genus-raising clutching that converts two closed flags
  into a self-loop.  This is $D_1 = D_{\mathrm{nsep}}$.
  The closed-color data at genus~$1$ is the curvature
  $\kappa(\cA) \cdot \omega_1$, a scalar.
\item \emph{Open color} (boundary): the boundary algebra
  $\cA_b = \mathrm{End}(b)$ inherits genus-$1$ corrections via
  the bulk-to-boundary map.  The genus-$1$ boundary correction is
  $\Theta^{(1)}_{n=1}(T) = -(c/24)\,E_2(\tau) \cdot T$, which
  is the Casimir energy density on the torus.
\end{itemize}

The directionality is strict: $\mu_{1,n}$ with $n \geq 1$ encodes
bulk observables \emph{restricting} to boundary corrections
(the $E_2$ and $\wp$ terms), but there is no map from the boundary
back to the genus-$1$ bulk.  This is Axiom~5 of the relative
modular extension.

\subsubsection{The curvature--braiding dichotomy at genus $1$}

The spectral braiding dichotomy governs the genus-$1$ operations:
\begin{itemize}[leftmargin=1.8em]
\item The curvature $\kappa \cdot \omega_1$ sources $\mu_{1,0}$
  (the scalar partition function).  It lives in the closed color
  and is controlled by the \emph{highest-pole coefficient}
  of the $\lambda$-bracket.
\item The $E_2(\tau)$ corrections in $\mu_{1,1}$ and the
  $\wp(z;\tau)$ kernel in $\mu_{1,2}$ are controlled by
  the \emph{lowest-pole coefficient} $c_0 = \{T {}_0 T\}$
  of the $\lambda$-bracket.  For the Virasoro algebra,
  $\{T {}_0 T\} = \partial T \neq 0$ (the zeroth product is
  nonzero), so the genus-$1$ braiding is in the
  \emph{entangled} (Yangian-type) regime: the curvature and
  braiding share a common geometric source in the
  quasi-periodicity of the Weierstrass $\zeta$-function.
\end{itemize}

\subsection{General structure: $\mu_{1,n}$ for arbitrary $n$}

The pattern extends to all arities.

\begin{proposition}[General genus-$1$ operation; computed here]
\label{prop:wn-mu1n-general}
For a logarithmic $\SCchtop$-algebra $\cA$ with genus-$0$ MC element
$\Theta^{(0)}$, the genus-$1$ MC element at arity $n$ is:
\begin{equation}\label{eq:wn-mu1n-general}
\Theta^{(1)}_{n}
\;=\;
\sum_{k \geq n+2}
\mathrm{proj}_n \circ D_1 \circ \Theta^{(0)}_{k}
\;+\;
\sum_{\substack{k_1+k_2 = n+2 \\ k_1,k_2 \geq 2}}
\Theta^{(0)}_{k_1} \circ_{(1)} \Theta^{(0)}_{k_2}
\;\Big|_{\text{genus-}1\text{ propagator}},
\end{equation}
where the first sum runs over self-sewings of genus-$0$ vertices of
arity $k \geq n+2$ (contracting two of the $k$ flags into a self-loop,
leaving $n$ external legs), and the second sum runs over separating
compositions that propagate through a genus-$1$ Green's function.
\end{proposition}

For the Virasoro algebra:
\begin{enumerate}[label=\textup{(\alph*)}]
\item $n = 0$: only the $k = 2$ self-sewing contributes
  $\Rightarrow \mu_{1,0} = \kappa \cdot \omega_1 = (c/2)\,\omega_1$.
\item $n = 1$: the $k = 3$ self-sewing dominates
  $\Rightarrow \mu_{1,1}(T) = -(c/24)\,E_2(\tau) \cdot T$.
\item $n = 2$: the $k = 4$ self-sewing plus the $m_2 \circ m_2$
  propagation $\Rightarrow \mu_{1,2}$ involves $\wp(z;\tau)$ and
  $E_2(\tau)$.
\item $n = 3$: the $k = 5$ self-sewing plus lower compositions
  $\Rightarrow \mu_{1,3}$ involves $\wp'(z;\tau)$ (the derivative
  of the Weierstrass $\wp$-function).
\item $n \geq 4$: the genus-$1$ $n$-point function involves
  $\wp^{(n-2)}(z;\tau)$ and descendant corrections at each level.
\end{enumerate}

The general pattern: the genus-$1$ $n$-point function on the torus
replaces the genus-$0$ rational propagator $1/(z_i - z_j)^k$ by its
elliptic counterpart $\wp^{(k-2)}(z_i-z_j;\tau)$, plus quasi-modular
corrections from zero-mode subtractions.  The bar complex governs
this replacement systematically: $D_1$ is the algebraic machine that
produces the elliptic propagators from the genus-$0$ OPE data.

\subsection{Summary}

\begin{center}
\renewcommand{\arraystretch}{1.5}
\begin{tabular}{@{}cccc@{}}
\toprule
Operation & Bar complex origin & Virasoro value & Modular function \\
\midrule
$\mu_{1,0}$ & $\mathrm{Tr}(D_1 \circ m_2)$
  & $(c/2)\,\omega_1$
  & $\eta(\tau)^{-c/2}$ \\
$\mu_{1,1}(T)$ & $\mathrm{proj}_1 \circ D_1 \circ m_3$
  & $-(c/24)\,E_2(\tau)\,T$
  & Quasi-modular weight $2$ \\
$\mu_{1,2}(T,T)$ & $D_1 \circ m_4 + m_2 \circ m_2$
  & $c\,\wp(z_{12}) + (c/12)\,E_2$
  & Weierstrass $\wp$ + $E_2$ \\
\bottomrule
\end{tabular}
\end{center}

The three operations exhibit a uniform structure:
\begin{enumerate}[label=\textup{(\roman*)}]
\item Each $\mu_{1,n}$ is produced by the genus-raising differential
  $D_1 = D_{\mathrm{nsep}}$ of the relative Feynman transform
  acting on genus-$0$ data of arity $\geq n+2$.
\item Each factors through the modular bar coalgebra $\Bmod(\cA)$
  and its bicomplex $(D_0, D_1)$.
\item The modular functions that appear ($\eta$, $E_2$, $\wp$) are
  the elliptic replacements of the rational genus-$0$ propagators,
  systematically generated by the self-sewing mechanism.
\item The curved bar differential $d^2 = \kappa \cdot \omega_1$ is
  visible at every arity: it sources $\mu_{1,0}$ directly, it
  controls the quasi-modular anomaly of $\mu_{1,1}$, and it appears
  as the $E_2$-correction in $\mu_{1,2}$.
\end{enumerate}

This completes the genus-$1$ seed of the modular envelope.  The
genus-$2$ entry point involves three channels
($\Delta_{\mathrm{cyc}}(K_{1,1})$, $\frac{1}{2}[K_{1,1},K_{1,1}]$,
and rigid-boundary residues) and is the subject of the next computation.



%% ================================================================
%% Purity--Koszulness beyond the affine lineage
%% To be appended to working_notes.tex (before \end{document})
%% ================================================================

\section{Purity--Koszulness: testing beyond the affine lineage}
\label{sec:purity-koszul-beyond-affine}

The semisimple purity criterion (Theorem~thm:semisimple-purity
in \texttt{bar-cobar-review.tex}) establishes
\[
\text{collision purity}
\;\Longleftrightarrow\;
\text{chiral Koszulness}
\]
for chiral algebras whose effective OPE residues
$\Omega_{ij}^{\mathrm{eff}}$ are semisimple endomorphisms.
The theorem is verified for the affine lineage:
$V^k(\mathfrak{g})$, $\beta\gamma$, Landau--Ginzburg, and
abelian Chern--Simons.  This section tests the correspondence
for three families outside the affine lineage:
\emph{lattice VOAs}, \emph{orbifold VOAs}, and
\emph{coset VOAs} (Virasoro minimal models via GKO).
We also investigate the sharpness of the one-loop criterion
(Theorem~thm:one-loop-koszul) and determine the Koszul
locus of $\cW_3$ via the DS obstruction theorem
(Theorem~thm:ds-koszul-obstruction).

\subsection{Lattice VOAs: $V_L$ for $L = A_2, D_4$}
\label{subsec:lattice-purity}

Let $L$ be an even lattice of rank~$r$.  The lattice VOA
$V_L$ has generators: Heisenberg currents $J^a(z)$
($a = 1,\ldots,r$) and vertex operators $V_\alpha(z) = e^\alpha(z)$
for $\alpha \in L$.  The OPE data is:
\begin{align}
J^a(z) J^b(w) &\sim \frac{\langle e_a, e_b \rangle}{(z-w)^2},
  \label{eq:lattice-JJ}\\
J^a(z) V_\alpha(w) &\sim
  \frac{\langle e_a, \alpha\rangle\, V_\alpha(w)}{z-w},
  \label{eq:lattice-JV}\\
V_\alpha(z) V_\beta(w) &\sim
  \epsilon(\alpha,\beta)\,
  (z-w)^{\langle\alpha,\beta\rangle}\,
  V_{\alpha+\beta}(w) + \cdots
  \label{eq:lattice-VV}
\end{align}

\paragraph{Effective residue analysis.}
The double-pole term in \eqref{eq:lattice-JJ} is
\emph{scalar} (it takes values in $\C\cdot\mathbf{1}$),
so it contributes to the curvature $m_0 = \kappa(V_L)$
and is projected away in the effective residue.

The simple-pole term in \eqref{eq:lattice-JV} has residue
$\langle e_a, \alpha\rangle \cdot V_\alpha$, which acts
\emph{diagonally} on the vertex-operator sector: on the
tensor product $V_L^{\otimes 2}$, the effective residue
$\Omega_{12}^{\mathrm{eff}}$ restricted to the
$J^a \otimes V_\alpha$ sector is the scalar
$\langle e_a, \alpha\rangle$.  This is trivially semisimple
(scalar on each sector).

The vertex-operator OPE \eqref{eq:lattice-VV} has pole
of order $-\langle\alpha,\beta\rangle$ (negative for
$\langle\alpha,\beta\rangle > 0$, meaning a zero rather
than a pole).  When $\langle\alpha,\beta\rangle < 0$
(as happens for roots in opposite Weyl chambers), the
residue coefficient $\epsilon(\alpha,\beta)\,V_{\alpha+\beta}$
is proportional to a single vertex operator, hence acts by
rank~$1$ on the tensor product.
In particular, the effective residue on the
$V_\alpha \otimes V_\beta$ sector (for $\alpha,\beta$
with $\langle\alpha,\beta\rangle < 0$) is a \emph{scalar
multiple} of the projection onto $V_{\alpha+\beta}$ ---
semisimple because it is a rank-$1$ projector up to scalar.

\begin{proposition}[Lattice VOAs have semisimple effective residues]
\label{prop:lattice-semisimple}
For every even lattice~$L$, the effective OPE residue
$\Omega_{ij}^{\mathrm{eff}}$ of~$V_L$ is semisimple.
Consequently, by Theorem~thm:semisimple-purity,
$V_L$ is chirally Koszul if and only if the OPE collision data
is pure.
\end{proposition}

\begin{proof}
The effective residue decomposes into three blocks,
corresponding to the three OPE types
\eqref{eq:lattice-JJ}--\eqref{eq:lattice-VV}.

\emph{Block~1} ($J^a \otimes J^b$):
the $J$--$J$ OPE has only a double-pole scalar term
(no simple pole between currents), so
$\Omega_{12}^{\mathrm{eff}}|_{J \otimes J} = 0$, which
is trivially semisimple.

\emph{Block~2} ($J^a \otimes V_\alpha$ and
$V_\alpha \otimes J^a$):
the residue acts by the scalar
$\langle e_a, \alpha\rangle$
on each $J$--$V_\alpha$ pair.  A diagonal matrix with
scalar entries is semisimple.

\emph{Block~3} ($V_\alpha \otimes V_\beta$):
when $\langle\alpha,\beta\rangle = -n < 0$ (the only case
producing a genuine pole at $z = w$), the residue is
$\epsilon(\alpha,\beta)\cdot V_{\alpha+\beta}$, acting
as a rank-$1$ operator on the fiber.  On the full
$V_\alpha \otimes V_\beta$ tensor-product sector,
$\Omega^{\mathrm{eff}}$ is a direct sum of rank-$1$
operators indexed by lattice-vector pairs, hence diagonal
in the $(\alpha,\beta)$-basis and therefore semisimple.

When $\langle\alpha,\beta\rangle \geq 0$, the OPE is
regular (no pole), contributing nothing to the effective
residue.  All three blocks are separately semisimple, so
$\Omega_{12}^{\mathrm{eff}}$ is semisimple.
\end{proof}

\begin{computation}[Explicit checks: $A_2$ and $D_4$]
\label{comp:lattice-A2-D4}%
\emph{$L = A_2$ (root lattice of $\fsl_3$, rank~$2$).}
The root system $\Phi(A_2) = \{\pm\alpha_1, \pm\alpha_2,
\pm(\alpha_1+\alpha_2)\}$ has $6$ roots.  The inner products
are $\langle\alpha_i,\alpha_i\rangle = 2$,
$\langle\alpha_1,\alpha_2\rangle = -1$.
The currents are $J^1, J^2$ (rank~$2$).  Total generator count:
$2$ currents $+ 6$ vertex operators $= 8$ generators.

The effective residue $\Omega^{\mathrm{eff}}$ on
$V_{A_2}^{\otimes 2}$ is an $8^2 = 64$-dimensional matrix.
By Proposition~\ref{prop:lattice-semisimple}, it is semisimple.
The eigenvalues on each Block~2 sector are
$\langle e_a, \alpha\rangle \in \{0, \pm 1, \pm 2\}$
(the weight pairings in the root lattice), all scalar.
The Block~3 sectors with poles ($\langle\alpha,\beta\rangle = -1$)
are: $(\alpha_1, \alpha_2)$, $(\alpha_2, \alpha_1)$,
$(-\alpha_1, -\alpha_2)$, $(-\alpha_2, -\alpha_1)$,
$(\alpha_1, -\alpha_1-\alpha_2)$, etc.\ --- each contributes a
rank-$1$ projector.  The full matrix is diagonal in the lattice basis.

\emph{Purity:} Semisimple effective residues $\Longrightarrow$
collision purity $\Longrightarrow$ chiral Koszulness.
Independently confirmed by PBW universality
(Theorem~thm:lattice-koszul).
The purity--Koszulness biconditional holds for~$V_{A_2}$.

\smallskip\noindent
\emph{$L = D_4$ (root lattice of $\mathfrak{so}_8$, rank~$4$).}
The root system $\Phi(D_4)$ has $24$ roots,
inner products $\langle\alpha,\alpha\rangle = 2$ and
$\langle\alpha_i,\alpha_j\rangle \in \{0, -1\}$ for
simple roots.  Total generators: $4$ currents $+ 24$ vertex
operators $= 28$.  The effective residue is a $28^2 = 784$-dimensional
matrix.  By the same argument as $A_2$, it decomposes into scalar
blocks and is semisimple.  The $D_4$ triality automorphism
(the exceptional $S_3$ symmetry of the Dynkin diagram) acts on
the effective residue by permuting blocks but preserves
semisimplicity.

\emph{Purity:} Confirmed. The biconditional holds for~$V_{D_4}$.
\end{computation}

\begin{remark}[Why lattice VOAs are class~$\mathbf{G}$]
Lattice VOAs belong to class~$\mathbf{G}$ (shadow depth~$2$)
in the intersection-depth classification.  The mechanism is
transparent from the purity perspective: \emph{all} effective
OPE poles are at most simple, and the residues act by scalars
or rank-$1$ projectors in the lattice basis.  The
self-intersection $\mathscr{L} \times_{\mathscr{M}}^h
\mathscr{L}$ is clean (\'etale groupoid), with
$\mathrm{Tor}_i = 0$ for $i \geq 1$ between the Lagrangian
and itself.  No higher $A_\infty$ operations arise:
$m_k = 0$ for $k \geq 3$ on cohomology.  Collision purity
holds, and the weight spectral sequence degenerates
at~$E_1$.

This extends to \emph{all} even lattice VOAs $V_L$,
regardless of rank: the lattice structure forces semisimplicity
of the effective residues uniformly.  In particular,
$V_{E_8}$ and $V_{\mathrm{Leech}}$ (rank~$8$ and $24$
respectively) satisfy the biconditional.
\end{remark}

\subsection{Orbifold VOAs: $V_{A_1}^{\Z/2}$}
\label{subsec:orbifold-purity}

The $\Z/2$-orbifold of the $A_1$ lattice VOA is
$V_{A_1}^{+} := V_{A_1}^{\Z/2}$, where the
$\Z/2$-action is $J \mapsto -J$,
$V_\alpha \mapsto V_{-\alpha}$.  This is a $c = 1$ theory
whose untwisted sector is the $\Z/2$-invariant
subalgebra of $V_{A_1}$, and whose twisted sector
introduces new fields from the orbifold fixed points.

\paragraph{Generator structure.}
The untwisted sector of $V_{A_1}^{+}$ is generated by
even combinations:
\begin{itemize}[nosep]
\item $:J^2: = :J(z)J(z):$ (the normal-ordered square, weight $2$),
\item $V_\alpha + V_{-\alpha}$ (even vertex operators, weight
  $\langle\alpha,\alpha\rangle/2 = 1$ for root-length vectors).
\end{itemize}
The twisted sector adds two ground-state fields
$\sigma^{\pm}$ of weight $1/16$ (the
$\Z/2$-twist fields).  For the purposes of
chiral Koszulness, only the integer-weight generators
contribute to the $A_\infty$ operations.

\paragraph{OPE structure.}
The key OPEs in $V_{A_1}^{+}$ are:
\begin{align}
:J^2:(z)\, :J^2:(w) &\sim
  \frac{1}{(z-w)^4} + \frac{4\,:J^2:(w)}{(z-w)^2}
  + \frac{2\,\partial :J^2:(w)}{z-w},
  \label{eq:orbifold-J2J2}\\
:J^2:(z)\, (V_\alpha + V_{-\alpha})(w) &\sim
  \frac{V_\alpha + V_{-\alpha}}{(z-w)^2}
  + \frac{\partial(V_\alpha + V_{-\alpha})}{z-w}.
  \label{eq:orbifold-J2V}
\end{align}

\paragraph{Effective residue analysis.}
The generator $:J^2:$ has weight~$2$ and its self-OPE
\eqref{eq:orbifold-J2J2} has a \emph{fourth-order pole}
with scalar coefficient~$1$ (the central term) and a
\emph{non-scalar second-order pole} ($4\,:J^2:$ acts
by the conformal weight operator $4L_0$ on the
$:J^2:$-sector).  The effective residue after removing
the scalar quartic pole includes the non-scalar quadratic
residue.

This is structurally identical to the Virasoro OPE:
$:J^2: \sim T$ at $c = 1$.  The integer-weight
subalgebra of $V_{A_1}^{+}$ \emph{contains} the
Virasoro algebra at $c = 1$ via the Sugawara construction,
and the $:J^2:$ field IS the stress tensor~$T$.

\begin{proposition}[The orbifold $V_{A_1}^{+}$: Koszul
but not collision-pure]
\label{prop:orbifold-A1-koszul}
The orbifold vertex algebra $V_{A_1}^{+}$ is chirally
Koszul \textup{(}by PBW universality, as it is freely
strongly generated\textup{)} but fails collision purity.
The effective OPE residue is non-semisimple: the
$:J^2: \otimes :J^2:$ sector contains a Jordan block
from the quartic-to-quadratic pole mixing, identical
to the Virasoro mechanism of
Proposition~prop:purity-failure-virasoro.
\end{proposition}

\begin{proof}
The OPE \eqref{eq:orbifold-J2J2} is exactly the
$c = 1$ Virasoro OPE for $T = :J^2:$.  By the
computation in the proof of
Proposition~prop:purity-failure-virasoro,
the extension between the quartic-pole scalar residue
and the quadratic-pole non-scalar residue is non-split
(the $L_{-1}$-descendant relation $\partial T = \partial :J^2:$
links weights~$4$ and~$2$).  This produces a non-trivial
$d_2$ differential in the weight spectral sequence,
obstructing collision purity.

Koszulness holds independently: $V_{A_1}^{+}$ is freely
strongly generated by $:J^2:$ and the even vertex
operators plus their $J$-descendants, and the defining
OPEs are homogeneous quadratic in the filtered grading.
PBW universality gives chirally Koszul.
\end{proof}

\begin{remark}[Orbifolds and the purity frontier]
The orbifold $V_{A_1}^{+}$ falls into class~$\mathbf{M}$
(infinite shadow depth): the Sugawara field $T = :J^2:$
inherits the quartic-pole non-formality of Virasoro.
This is significant: the parent $V_{A_1}$ is class~$\mathbf{G}$
(collision-pure, formal), but the orbifold construction
introduces a new strong generator ($:J^2:$) whose self-OPE
has higher poles.  \emph{Orbifold constructions can promote
class~$\mathbf{G}$ algebras to class~$\mathbf{M}$}: the
orbifold fixes the Heisenberg current but promotes the
normal-ordered square to a generator, and the quartic pole
in the $:J^2:$-self-OPE is a consequence of Wick contractions
that are invisible in the parent algebra.

The purity--Koszulness biconditional fails for
$V_{A_1}^{+}$: it is Koszul but not collision-pure.
This is consistent with Theorem~thm:semisimple-purity,
which only claims the biconditional for
\emph{semisimple} effective residues, and
$V_{A_1}^{+}$ has non-semisimple ones.
\end{remark}

\subsection{Coset VOAs: Virasoro minimal models via GKO}
\label{subsec:coset-purity}

The GKO coset construction produces the Virasoro minimal
models from the diagonal embedding:
\[
\mathcal{M}(p,p') \;:=\;
\frac{\widehat{\fsl}_2 \text{ at level } k
  \;\oplus\;
  \widehat{\fsl}_2 \text{ at level } 1}
{\widehat{\fsl}_2 \text{ at level } k{+}1},
\qquad
c(k) = 1 - \frac{6}{(k+2)(k+3)},
\]
with $p = k+2$, $p' = k+3$ coprime.  The coset algebra
is the commutant of the diagonal
$\widehat{\fsl}_2 \hookrightarrow
\widehat{\fsl}_{2,k} \oplus
\widehat{\fsl}_{2,1}$ in the tensor product.

\paragraph{The coset as a DS reduction.}
For the principal nilpotent in $\fsl_2$:
\[
H^0_{\mathrm{DS}}(V^k(\fsl_2))
\;\cong\;
\mathrm{Vir}_{c(k)},
\qquad
c(k) = 1 - \frac{6(k+1)^2}{k+2}.
\]
The GKO coset gives the same central charge when
$c_{\mathrm{GKO}}(k) = 1 - 6/((k+2)(k+3))$,
which is the $c_{p,p+1}$ specialization.  At these
special values, the Virasoro algebra is rational and the
coset construction produces a finite set of modules.

\begin{computation}[Koszulness and purity for $k = 1, 2, 3$]
\label{comp:coset-k123}%
We compute the Koszulness and purity status for the
first three coset values.

\smallskip\noindent
\emph{$k = 1$: Ising model ($c = 1/2$, $\mathcal{M}(3,4)$).}
The Ising model VOA is generated by a single Virasoro field~$T$
of weight~$2$ with $c = 1/2$.  The OPE is
$T(z)T(w) \sim \frac{1/4}{(z-w)^4} + \frac{2T(w)}{(z-w)^2}
+ \frac{\partial T(w)}{z-w}$.
The quartic pole coefficient is $c/2 = 1/4$ (nonzero), so
the Virasoro mechanism applies:

\emph{Koszulness:} Yes (PBW universality).

\emph{Collision purity:} No.  The quartic pole
$1/4 \cdot (z-w)^{-4}$ creates a non-split extension between
weights~$4$ and~$2$ in the MHM on $\FM_2(\C)$.
The weight spectral sequence has a non-trivial~$d_2$.

\emph{Verdict:} Koszul $\checkmark$, pure $\times$.

\smallskip\noindent
\emph{$k = 2$: Tricritical Ising ($c = 7/10$,
$\mathcal{M}(4,5)$).}
Quartic-pole coefficient: $c/2 = 7/20 \neq 0$.

\emph{Verdict:} Koszul $\checkmark$, pure $\times$.

\smallskip\noindent
\emph{$k = 3$: The $(5,6)$ model ($c = 4/5$).}
Quartic-pole coefficient: $c/2 = 2/5 \neq 0$.

\emph{Verdict:} Koszul $\checkmark$, pure $\times$.

\smallskip\noindent
\emph{Summary.}
Every unitary Virasoro minimal model $\mathcal{M}(k{+}2,
k{+}3)$ is chirally Koszul but fails collision purity,
for all $k \geq 1$.  The purity--Koszulness biconditional fails
for the entire GKO series.
\end{computation}

\begin{remark}[Cosets and the non-semisimple frontier]
The coset construction preserves Koszulness (by PBW
universality of the output Virasoro algebra) but
cannot produce collision purity because the output
algebra always has a quartic pole for $c \neq 0$.
The coset construction is therefore entirely within the
\emph{non-semisimple frontier} of
Remark~rem:purity-scope.
\end{remark}

\subsection{Sharpness of the one-loop criterion}
\label{subsec:one-loop-sharpness}

Theorem~thm:one-loop-koszul states:
\[
\text{one-loop exactness of BV-BRST}
\;\Longrightarrow\;
\text{chiral Koszulness}.
\]
We investigate whether the converse holds.

\begin{proposition}[The one-loop criterion is not sharp]
\label{prop:one-loop-not-sharp}
The converse of Theorem~thm:one-loop-koszul
is false.  The Virasoro minimal models provide an explicit
class of counterexamples: for each $k \geq 1$,
$\mathrm{Vir}_{c(k)}$ with $c(k) = 1 - 6/((k{+}2)(k{+}3))$
is chirally Koszul but its BV-BRST complex is \emph{not}
one-loop exact.
\end{proposition}

\begin{proof}
\emph{Koszulness:} established by PBW universality.

\emph{Failure of one-loop exactness:}
The BV-BRST complex for the Virasoro algebra arises
from the Khan--Zeng Poisson sigma model with the Virasoro
$\lambda$-bracket.  The interaction vertex
$V_3 = T\mu\partial\mu + \frac{c}{24}\mu\partial^3\mu$
is quadratic in~$\mu$.  By the loop-counting of
Proposition~prop:vir-truncation, every connected
Feynman diagram contributing to $m_k(T,\ldots,T)$ for
$k \geq 3$ has exactly one loop, and $m_k \neq 0$ for all
$k \geq 3$.

The distinction is critical:
\begin{itemize}[nosep]
\item \emph{One-loop exact} (BV sense): the effective
  action $S_{\mathrm{eff}} = S_{\mathrm{tree}} +
  \hbar S_1$, where $S_1$ is a one-loop determinant
  that does not modify the algebraic structure beyond
  a level shift.
\item \emph{One-loop perturbative}: all Feynman diagrams
  have loop number $\leq 1$, but the one-loop diagrams
  can produce non-trivial algebraic operations (higher
  $A_\infty$ operations) beyond a scalar correction.
\end{itemize}
The Virasoro theory is one-loop perturbative but not
one-loop exact.  The affine Kac--Moody theory is both:
its one-loop correction is the level shift
$k \mapsto k + h^\vee\hbar$ (a scalar), with no
non-trivial $m_k$ for $k \geq 3$.

The root cause: the Virasoro PVA $\pi^{ij}(\partial)$
is \emph{nonlinear} in the fields ($T\mu\partial\mu$ is
cubic), unlike the affine case where $\pi$ is linear
(Lie--Poisson).  For linear~$\pi$, the Kontsevich integral
localizes on wheel graphs that contribute only a scalar
determinant.  For nonlinear~$\pi$, the wheel graphs produce
genuine higher operations.
\end{proof}

\begin{remark}[Precise characterization of the gap]
The gap is:
\[
\boxed{
\text{Koszul} \;\not\!\!\Longrightarrow\;
\text{one-loop exact (BV)}.
}
\]
The examples filling this gap are precisely the chirally
Koszul algebras with non-formal Swiss-cheese structure
(class~$\mathbf{M}$): Virasoro at $c \neq 0$, all
$\cW_N$ algebras, and any VOA obtained by
DS reduction from the affine lineage that produces
higher-order poles.
\end{remark}

\subsection{The DS--Koszul obstruction for $\cW_3$:
the Koszul locus}
\label{subsec:ds-koszul-w3}

Theorem~thm:ds-koszul-obstruction states that
DS reduction of $\hat{\mathfrak{g}}_k$ preserves chiral
Koszulness but breaks Swiss-cheese formality.  We determine
the precise Koszul locus of $\cW_3$.

\paragraph{Setup.}
$\cW_3(c) = H^0_{\mathrm{DS}}(V^k(\fsl_3))$
with $c = c_{W_3}(k) = 2 - 24(k+2)^2/(k+3)$.

\begin{proposition}[The Koszul locus of $\cW_3$]
\label{prop:koszul-locus-W3}
$\cW_3(c)$ is chirally Koszul for all non-critical
levels~$k \neq -3$.  More precisely:
\begin{enumerate}[label=\textup{(\roman*)}]
\item The bar cohomology is concentrated for every $c \in \C$.
\item The Swiss-cheese structure is non-formal for
  all $c \neq 0$: $m_k^{\mathrm{DS}} \neq 0$ for $k \geq 3$.
\item $\cW_3(c)$ is one-loop exact (BV sense) iff $c = 0$.
\end{enumerate}
\end{proposition}

\begin{proof}
\emph{(i):}
DS reduction preserves chiral Koszulness
(Theorem~thm:ds-koszul-obstruction(i)).

\emph{(ii):}
The $T$--$T$ OPE has a quartic pole for $c \neq 0$.
The $W$--$W$ OPE has a sixth-order pole $c/3 \cdot (z-w)^{-6}$.
These higher poles generate non-trivial $m_k^{\mathrm{DS}}$
for $k \geq 3$ by the resolvent tree mechanism.

\emph{(iii):}
The parent $V^k(\fsl_3)$ is one-loop exact.
Under DS reduction, the transferred effective action acquires
additional terms from BRST trees, which evaluate to
non-trivial wheel integrals for $c \neq 0$.
\end{proof}

\begin{computation}[Explicit: $\cW_3$ at $k = 1$ ($c = -52$)]
\label{comp:W3-k1}
At $k = 1$: $c = 2 - 24 \cdot 9/4 = -52$.
$K_{\mathrm{eff}} = 1 + 3\hbar$.
\begin{itemize}[nosep]
\item $T$--$T$ quartic pole: $c/2 = -26 \neq 0$.
\item $W$--$W$ sextic pole: $c/3 = -52/3 \neq 0$.
\item Koszul: $\checkmark$.  Pure: $\times$.  One-loop exact: $\times$.
\end{itemize}
At $k = 3$ ($c = -98$) and $k = -1$ ($c = -10$):
same qualitative picture.
\end{computation}

\begin{remark}[The Koszul locus is the entire non-critical complement]
The Koszul locus of $\cW_3$ (and any principal
$\cW$-algebra) is $\C \setminus \{-h^\vee\}$.
The one-loop exact locus is empty (non-trivially).
\begin{center}
\renewcommand{\arraystretch}{1.25}
\begin{tabular}{@{}lcc@{}}
\toprule
\textbf{Property} & \textbf{Affine} $V^k(\fsl_3)$
  & \textbf{$\cW_3(c)$} \\
\midrule
Chirally Koszul & $\checkmark$ (all non-critical $k$)
  & $\checkmark$ (all non-critical $k$) \\
One-loop exact (BV) & $\checkmark$ (all $k$)
  & $\times$ (all $c \neq 0$) \\
Collision-pure & $\checkmark$ (generic $k$)
  & $\times$ (all $c \neq 0$) \\
Swiss-cheese formal & $\checkmark$ ($m_k = 0$, $k \geq 3$)
  & $\times$ ($m_k \neq 0$, $k \geq 3$) \\
Shadow class & $\mathbf{L}$ (depth $3$)
  & $\mathbf{M}$ (depth $\infty$) \\
\bottomrule
\end{tabular}
\end{center}
DS reduction preserves Koszulness but breaks all three
refinements simultaneously.
\end{remark}

\subsection{Synthesis: the extended purity--Koszulness landscape}
\label{subsec:purity-synthesis}

\begin{center}
\renewcommand{\arraystretch}{1.3}
\begin{tabular}{@{}lccccc@{}}
\toprule
\textbf{Algebra} & \textbf{Koszul} & \textbf{Pure}
  & \textbf{Semisimple} & \textbf{One-loop}
  & \textbf{Class} \\
\midrule
$\mathcal{H}_k$ (Heisenberg) & $\checkmark$ & $\checkmark$
  & $\checkmark$ & $\checkmark$ & $\mathbf{G}$ \\
$V^k(\mathfrak{g})$ (affine KM) & $\checkmark$ & $\checkmark$
  & $\checkmark$ & $\checkmark$ & $\mathbf{L}$ \\
$\beta\gamma/bc$ & $\checkmark$ & $\checkmark$
  & $\checkmark$ & $\checkmark$ & $\mathbf{C}$ \\
LG cubic & $\checkmark$ & $\checkmark$
  & $\checkmark$ & $\checkmark$ & $\mathbf{C}$ \\
Abelian CS & $\checkmark$ & $\checkmark$
  & $\checkmark$ & $\checkmark$ & $\mathbf{G}$ \\
$V_{A_2}$ (lattice) & $\checkmark$ & $\checkmark$
  & $\checkmark$ & $\checkmark$ & $\mathbf{G}$ \\
$V_{D_4}$ (lattice) & $\checkmark$ & $\checkmark$
  & $\checkmark$ & $\checkmark$ & $\mathbf{G}$ \\
\midrule
$\mathrm{Vir}_c$ ($c \neq 0$) & $\checkmark$ & $\times$
  & $\times$ & $\times$ & $\mathbf{M}$ \\
$V_{A_1}^{\Z/2}$ (orbifold) & $\checkmark$ & $\times$
  & $\times$ & $\times$ & $\mathbf{M}$ \\
$\mathcal{M}(3,4)$ (Ising) & $\checkmark$ & $\times$
  & $\times$ & $\times$ & $\mathbf{M}$ \\
$\mathcal{M}(4,5)$ (tri-Ising) & $\checkmark$ & $\times$
  & $\times$ & $\times$ & $\mathbf{M}$ \\
$\mathcal{M}(5,6)$ & $\checkmark$ & $\times$
  & $\times$ & $\times$ & $\mathbf{M}$ \\
$\cW_3(c)$ ($c \neq 0$) & $\checkmark$ & $\times$
  & $\times$ & $\times$ & $\mathbf{M}$ \\
\bottomrule
\end{tabular}
\end{center}

\begin{theorem}[Extended purity--Koszulness correspondence]
\label{thm:extended-purity-koszulness}
\begin{enumerate}[label=\textup{(\roman*)}]
\item The purity--Koszulness biconditional holds for all
  lattice VOAs $V_L$ \textup{(}arbitrary even lattice~$L$\textup{)},
  extending Theorem~thm:semisimple-purity to class~$\mathbf{G}$.
\item The biconditional \emph{fails} for orbifold VOAs,
  coset VOAs, and all $\cW_N$ algebras at $c \neq 0$.
\item The one-loop criterion is strictly sufficient:
  Virasoro at $c \neq 0$ is Koszul but not one-loop exact.
\item The Koszul locus of $\cW_3$ is
  $\{k \in \C : k \neq -3\}$.
  DS reduction preserves Koszulness but breaks one-loop
  exactness, collision purity, and formality simultaneously.
\item The dividing line is sharp:
  \begin{center}
  \emph{Semisimple eff.\ residues}
  $\Leftrightarrow$
  \emph{collision purity}
  $\Leftrightarrow$
  \emph{one-loop exact}
  $\Leftrightarrow$
  \emph{SC-formal}
  $\Leftrightarrow$
  \emph{class $\mathbf{G}/\mathbf{L}/\mathbf{C}$}.
  \end{center}
  But \emph{Koszulness persists} across this line
  via PBW universality.
\end{enumerate}
\end{theorem}

\begin{proof}
(i) Proposition~\ref{prop:lattice-semisimple} $+$
Theorem~thm:semisimple-purity.
(ii) Proposition~\ref{prop:orbifold-A1-koszul},
Computation~\ref{comp:coset-k123},
Proposition~\ref{prop:koszul-locus-W3}(ii).
(iii) Proposition~\ref{prop:one-loop-not-sharp}.
(iv) Proposition~\ref{prop:koszul-locus-W3}.
(v) Semisimple residues $\Rightarrow$ pure $V$-graded
pieces $\Rightarrow$ $E_1$-degeneration $\Rightarrow$
$m_k = 0$ ($k \geq 3$).  Non-semisimple $\Rightarrow$
non-split MHM extensions $\Rightarrow$ non-trivial $d_r$
$\Rightarrow$ non-vanishing $m_k$.  Koszulness is
independent (PBW).
\end{proof}

\begin{remark}[The remaining frontier]
The results completely characterize the landscape within
standard families and their extensions.  Open:
\begin{enumerate}[label=(\alph*)]
\item \emph{Non-freely-generated VOAs} (PBW fails).
\item \emph{KM at resonant level} (eigenvalue coincidence).
\item \emph{Beyond PBW}: can Koszulness hold without free
  strong generation?
\end{enumerate}
\end{remark}






% =================================================================
\section{Annular FM compactification and the genus-$1$ ordered sector}
\label{sec:annular-FM}
% =================================================================

The bordered ordered theory on the strip $[0,1]\times\R$ is
complete: the diagonal bicomodule $C_\Delta$ is the universal
boundary condition, interval composition is derived cotensor, and
the ordered open trace formalism is proved
(Theorem~\ref{thm:ordered-open-wn}).  The next stratum is
the annulus: configurations on $S^1\times\C$, where the
topological direction wraps.

This section constructs the Fulton--MacPherson compactification
for annular configurations, computes the annular bar differential
explicitly, and identifies the annular complex with Hochschild
homology of the boundary algebra.  The key new phenomenon is
the \emph{wrap-around boundary stratum}: when boundary points
on $S^1$ collide by one point passing through the basepoint
identification $0\sim 1$, the resulting face has no analogue
in the strip.  This face closes the cosimplicial structure of
the cobar complex and produces the coHochschild differential.

\subsection{The annular configuration space}

\begin{definition}[Annular ordered configuration space]
\label{def:annular-config-wn}
Let $S^1=[0,1]/(0\sim 1)$ be the standard circle.  For integers
$r\ge 0$ (interior points) and $s\ge 0$ (boundary points), the
\emph{annular ordered configuration space} is
\[
\Conf_{r,s}^{\mathrm{ann}}(S^1\times\C)
\;:=\;
\bigl\{
(z_1,\dotsc,z_r;\,y_1,\dotsc,y_s)
\bigr\},
\]
where:
\begin{enumerate}[label=(\alph*)]
\item the $z_i=(t_i,w_i)\in (S^1\setminus\partial)\times\C
\cong (0,1)\times\C$ are distinct interior points;
\item the $y_j\in S^1$ are distinct boundary points, ordered
\emph{cyclically} by the orientation of~$S^1$: after choosing
a basepoint $*\in S^1$, we lift to a linear ordering
$y_1<y_2<\cdots<y_s$ on $S^1\setminus\{*\}$;
\item all points are pairwise distinct.
\end{enumerate}
The real dimension is $\dim_\R=3r+s$.
\end{definition}

\begin{remark}[Cyclic vs.\ linear ordering]
\label{rem:cyclic-vs-linear-wn}
On the strip $[0,1]$, boundary points carry a \emph{linear}
ordering inherited from the endpoints $0<1$.  On the circle $S^1$,
the ordering is only \emph{cyclic}: there is no canonical
first element.  Cutting $S^1$ at the basepoint $*$ recovers
a linear ordering; the ambiguity in the choice of $*$ is the
cyclic group $\Z/s\Z$ acting on the linear sequence.  The
passage from linear to cyclic ordering is the passage from the
two-sided bar complex to the cyclic bar complex.
\end{remark}

\subsection{The annular FM compactification}
\label{subsec:annular-FM-compact-wn}

\begin{definition}[Annular bordered FM compactification]
\label{def:annular-FM-wn}
The \emph{annular bordered FM space}
$\overline{\FM}{}^{\mathrm{ann}}_{r,s}$ is the
iterated real-oriented blowup of
$(S^1\times\C)^r\times(S^1)^s$ along the collision diagonals,
organized into three stages as in the bordered strip case,
with one essential modification:
boundary collisions are governed by the cyclic ordering.

\medskip
\noindent\textbf{Stage I: Interior collisions.}
Identical to the strip case.  For every subset
$S\subseteq\{1,\dotsc,r\}$ with $|S|\ge 2$, blow up the diagonal
$\Delta_S=\{z_i=z_j\;\forall\;i,j\in S\}$ in decreasing order
of~$|S|$.  These divisors produce the bar differential
$d_{\mathrm{res}}$ on the $\C$-direction.

\medskip
\noindent\textbf{Stage II: Boundary collisions.}
For every subset $T\subseteq\{1,\dotsc,s\}$ with $|T|\ge 2$
whose elements are \emph{cyclically consecutive} on $S^1$, blow
up the boundary diagonal.
The cyclic consecutiveness condition produces two classes:
\begin{enumerate}[label=(\roman*)]
\item \textbf{Non-wrapping subsets:} no passage through the basepoint.
These produce the same boundary faces as the strip.
\item \textbf{Wrapping subsets:} a cyclically
consecutive block that passes through the identification
$s\sim 1$.  These are new: they have no analogue in the
strip and encode the wrap-around collisions.
\end{enumerate}

\medskip
\noindent\textbf{Stage III: Mixed collisions.}
Identical to the strip case.  Interior points approaching
boundary points produce the bulk-to-boundary restriction maps.
\end{definition}

\begin{proposition}[Boundary strata of the annular FM space]
\label{prop:annular-strata-wn}
The codimension-$1$ boundary of
$\overline{\FM}{}^{\mathrm{ann}}_{r,s}$ consists of:
\begin{enumerate}[label=\textup{(\roman*)}]
\item Interior collision divisors $E_S^{\mathrm{int}}$
($|S|\ge 2$): the bar differential in the holomorphic direction.
\item Non-wrapping boundary collision divisors: the $A_\infty$
faces, identical to the strip.
\item Wrapping boundary collision divisors: the new
annular faces.  These connect the last boundary point to
the first, closing the linear sequence into a cycle.
\item Mixed collision divisors: the bulk-to-boundary restriction.
\end{enumerate}
The combinatorial indexing set for (ii) and (iii) together
is the set of cyclically consecutive subsets of $\Z/s\Z$
of size $\ge 2$.
\end{proposition}

\begin{observation}[Necklaces of planted trees]
\label{obs:necklace-trees-wn}
On the strip, the boundary strata of type~(ii) are indexed by
\emph{planted forests}: ordered sequences of rooted trees.  On the
annulus, the wrapping identification replaces the sequence by a
\emph{necklace}: a cyclic sequence modulo cyclic rotation, i.e.\
a \emph{necklace of planted trees} (cyclic planted forest).
\end{observation}

\subsection{The annular bar complex}
\label{subsec:annular-bar-complex-wn}

\begin{definition}[Annular bar complex]
\label{def:annular-bar-wn}
Let $\cA$ be a strongly admissible augmented $E_1$-chiral algebra
and $C=\barB^{\mathrm{ch}}(\cA)$ its ordered bar coalgebra, with
$C_\Delta$ the diagonal bicomodule.  The \emph{annular bar complex}
is the chain complex
\[
B^{\mathrm{ann}}_\bullet(\cA)
\;:=\;
C_\bullet\!\bigl(
\overline{\FM}{}^{\mathrm{ann}}_{r,s};\,
\omega_{\mathrm{bar}}
\bigr)
\;\cong\;
\bigoplus_{n\ge 0}
\mathrm{cyc}^n(C,C_\Delta),
\]
where $\mathrm{cyc}^n(C,C_\Delta):=
C_\Delta\widehat\otimes_C C^{\widehat\otimes\,n}$ is the
cyclic cotensor.
\end{definition}

\begin{theorem}[The annular bar differential]
\label{thm:annular-bar-differential-wn}
The differential on $B^{\mathrm{ann}}_\bullet(\cA)$ decomposes as
\[
d^{\mathrm{ann}}
\;=\;
d_{\mathrm{int}}+d_{\mathrm{res}}+d_\Delta+d_{\mathrm{wrap}},
\]
where:
\begin{enumerate}[label=\textup{(\alph*)}]
\item $d_{\mathrm{int}}$: internal differential on each tensor factor.
\item $d_{\mathrm{res}}$: collision residue differential (OPE residues at interior collisions).
\item $d_\Delta$: comultiplication differential at non-wrapping boundary collisions.
\item $d_{\mathrm{wrap}}$: the wrap-around differential at the seam.  Explicitly:
\[
d_{\mathrm{wrap}}\langle a_1|\cdots|a_n\rangle
\;=\;
\sum_{(a_n)}
(-1)^{\epsilon}\,
\langle a_1|\cdots|a_{n-1}|a_n'|a_n''\rangle
\;+\;
\sum_{(a_1)}
(-1)^{\epsilon'}\,
\langle a_1'|a_1''|a_2|\cdots|a_n\rangle.
\]
\end{enumerate}
The identity $(d^{\mathrm{ann}})^2=0$ follows from
$\partial^2=0$ on the manifold with corners
$\overline{\FM}{}^{\mathrm{ann}}_{r,s}$.
\end{theorem}

\begin{observation}[The wrap-around is the cyclic gluing]
The wrap-around differential $d_{\mathrm{wrap}}$ is the only
component of $d^{\mathrm{ann}}$ without an analogue in the strip
complex.  Algebraically, $d_{\mathrm{wrap}}$
provides exactly the coface maps $\delta_0$ and $\delta_n$ that
close the two-sided cobar complex into the coHochschild complex.
Without it, the complex computes
$\mathrm{Cotor}^{C^e}(C_\Delta,C_\Delta)$; with it, the two copies
are identified and the result is the self-cotrace
$C_\Delta\square_{C^e}^{\mathbf L}
C_\Delta\simeq\operatorname{coHH}_\bullet^{\mathrm{ch}}(C,C_\Delta)$.
\end{observation}

\subsection{The annular bar complex computes Hochschild homology}
\label{subsec:annular-HH-wn}

\begin{theorem}[Annular bar complex $=$ Hochschild homology]
\label{thm:annular-HH-wn}
Let $\cA$ be a strongly admissible augmented $E_1$-chiral algebra,
$C=\barB^{\mathrm{ch}}(\cA)$ its ordered bar coalgebra.  Then:
\[
H_\bullet(B^{\mathrm{ann}}(\cA),d^{\mathrm{ann}})
\;\simeq\;
\HH_\bullet^{\mathrm{ch}}(\cA)
\;\simeq\;
\operatorname{coHH}_\bullet^{\mathrm{ch}}(C,C_\Delta).
\]
The first isomorphism identifies the annular bar complex with
chiral associative Hochschild homology.  The second is the
Hochschild/coHochschild equivalence of the core chapter
(Theorem~\ref{thm:HH-coHH-homology}).
The key identification: the wrap-around differential
$d_{\mathrm{wrap}}$ provides exactly the coface maps $\delta_0$
and $\delta_s$ that distinguish the coHochschild complex from the
two-sided cobar complex.
\end{theorem}

\subsection{Relation to the modular operad}
\label{subsec:annular-modular-wn}

\begin{theorem}[Annular FM as genus-$1$ modular Swiss-cheese]
\label{thm:annular-modular-wn}
The annular FM compactification
$\overline{\FM}{}^{\mathrm{ann}}_{r,s}$ is the genus-$1$ part
of the modular Swiss-cheese operad.  The clutching recursion
connects genus~$0$ and genus~$1$: cutting the annulus along a
boundary arc produces the strip, and the annular closure is the
clutching of the two strip endpoints.
\end{theorem}

\subsection{The genus-$1$ $R$-matrix correction}
\label{subsec:genus1-R-matrix-wn}

\begin{theorem}[Genus-$1$ $R$-matrix from annular monodromy]
\label{thm:genus1-R-matrix-wn}
The genus-$1$ correction to the $R$-matrix is the monodromy of
the ordered bar connection along the new $\pi_1(S^1)$ generator:
\[
R_1(z;\tau)
\;=\;
\operatorname{Mon}_{\gamma_{S^1}}
\bigl(\nabla^{\mathrm{bar}}\big|_{\mathrm{ann}}\bigr).
\]
For $\cA=\widehat{\fg}_k$:
\[
R_1(z;\tau)
\;=\;
\frac{\kappa(\widehat{\fg}_k)}{(k+h^\vee)^2}\,
\wp(z;\tau)\,\Omega
\;+\;
O(\wp'),
\]
where $\wp(z;\tau)$ is the Weierstrass $\wp$-function,
$\Omega=\sum_a e_a\otimes e^a$ is the Casimir element
of~$\fg$, and $\kappa(\widehat{\fg}_k)=k/(k+h^\vee)$ is the
curvature.  The full modular $R$-matrix is
$R^{\mathrm{mod}}(z;\tau,\hbar)=R_0(z)+\hbar^2 R_1(z;\tau)+O(\hbar^4)$.
\end{theorem}

\begin{remark}[The Weierstrass $\wp$-function as genus-$1$ propagator]
The appearance of $\wp(z;\tau)$ is not coincidental: $\wp$ is the
genus-$1$ Green's function on $\C/(\Z\oplus\tau\Z)$, just as
$1/(z_1-z_2)$ is the genus-$0$ Green's function on~$\C$.  The
modular properties of $\wp$ transmit to the modular
properties of $R_1$, which is the mathematical origin of modular
covariance of the genus-$1$ $R$-matrix.
\end{remark}



% [Brief version of soft graviton section removed; detailed version below.]




% =================================================================
%
%   TYPE H ALGEBRAS: CLASSIFICATION AND THE CRITICAL-LEVEL QUESTION
%
% =================================================================

\section{Type H algebras: classification and the critical-level question}
\label{sec:type-H}

\providecommand{\dalg}{d_{\mathrm{alg}}}
\providecommand{\dfib}{d_{\mathrm{fib}}}
\providecommand{\hvee}{h^{\!\scriptscriptstyle\vee}}
\providecommand{\bD}{\mathbb{D}}

The classification of logarithmic $\SCchtop$-algebras by shadow depth
(Theorem~\ref{thm:classification-shadow-depth} of the relative Feynman
transform chapter) identifies two proved types: Type~F (formal,
$\dalg = 0$) and Type~V (Virasoro-type, $\dalg = 1$).
Conjecture~\ref{conj:type-H-existence} asks whether Type~H algebras
exist --- those with shadow depth $\dalg \geq 2$, whose higher
$\Ainf$ operations are not determined by $m_2$ and $m_3$ alone.

A distinct but related question arises from the genus tower: can
a chiral algebra be \emph{uncurved} ($\kappa = 0$) yet carry nontrivial
higher-genus data?  We give a complete answer to both questions.

\subsection{The single-scalar theorem: no secondary curvature}

\begin{theorem}[Non-existence of secondary curvature invariants]
\label{thm:no-secondary-curvature}
Let $\cA$ be a logarithmic $\SCchtop$-algebra.
The curvature of the genus-$g$ bar complex is
\begin{equation}\label{eq:universal-curvature}
\dfib^{\,2} = \kappa(\cA) \cdot \omega_g
\quad \text{for all } g \geq 1,
\end{equation}
where $\kappa(\cA) \in \mathbb{k}$ is the modular characteristic
and $\omega_g \in H^0(\overline{\cM}_g)$ is the fundamental class
of the Hodge bundle.  In particular:
\begin{enumerate}[label=\textup{(\roman*)}]
\item There is no independent ``secondary curvature'' $\kappa_2$
  at genus~$2$: the same scalar $\kappa$ governs the curvature
  at every genus.
\item If $\kappa(\cA) = 0$, the bar complex
  $(\barB^{(g)}(\cA), \dfib)$ is an honest cochain complex
  ($\dfib^{\,2} = 0$) at every genus $g \geq 1$.
\item The genus tower is controlled by a single scalar invariant,
  not by a sequence of independent genus-dependent data.
\end{enumerate}
\end{theorem}

\begin{proof}
This is the content of Vol~I, Theorem~D (the leading coefficient
theorem): the curvature of the curved bar complex is universally
$\kappa(\cA) \cdot \omega_g$, with $\kappa(\cA)$ determined by the
OPE structure at genus~$0$.  The key structural point is that
$\omega_g$ is pulled back from the universal Hodge bundle via the
factorization structure, while $\kappa(\cA)$ is the \emph{single}
scalar extracted from the $g = 1$ degeneration.  No additional
scalar is available from the genus-$2$ degeneration that is not
already determined by $\kappa$: the $\delta$-node degeneration
of $\overline{\cM}_2$ factors through the $g = 1$ data via
clutching, and the separating degeneration factors through
$\overline{\cM}_{1,1} \times \overline{\cM}_{1,1}$.
\end{proof}

\begin{remark}[Consequence for the Type~H question]
\label{rem:type-H-curvature-distinction}
The non-existence of secondary curvature means that the na\"ive
formulation of ``Type~H'' --- an algebra with $\kappa = 0$ at
genus~$1$ but $\kappa_2 \neq 0$ at genus~$2$ --- is
\emph{vacuous}: no such $\kappa_2$ exists.  The question must be
reformulated.  There are two meaningful reformulations:
\begin{enumerate}[label=\textup{(\alph*)}]
\item \textbf{Shadow depth}: does $\dalg \geq 2$ exist?
  (Conjecture~\ref{conj:type-H-existence}.)
\item \textbf{Uncurved but non-trivial}: does there exist
  $\cA$ with $\kappa(\cA) = 0$ and nontrivial genus-$g$
  \emph{bar cohomology} $H^\bullet(\barB^{(g)}(\cA), \dfib)$?
\end{enumerate}
Reformulation~(b) has a definitive positive answer:
\emph{critical-level Kac--Moody algebras}.
\end{remark}

\subsection{Classification of all $\kappa = 0$ algebras}

We now classify all known chiral algebras satisfying $\kappa = 0$.

\begin{proposition}[The uncurved locus]
\label{prop:uncurved-locus-complete}
In the standard landscape (Heisenberg, affine KM, Virasoro,
$\cW$-algebras, lattice, free fermion, $\beta\gamma$), the
algebras with $\kappa = 0$ are:
\begin{center}
\renewcommand{\arraystretch}{1.3}
\begin{tabular}{lcccc}
\textbf{Family} & \textbf{Condition for $\kappa = 0$}
  & $\boldsymbol{\dalg}$ & \textbf{$r_{\max}$}
  & \textbf{Genus-$g$ bar cohom.} \\
\hline
$\cH_k$ & $k = 0$ (trivial) & $0$ & $2$
  & trivial \\
$\widehat{\fg}_k$ & $k = -\hvee$ (critical level)
  & $0$ & $3$
  & $\mathrm{Fun}\,\mathrm{Op}_{\fg^\vee}(\Sigma_g)$ \\
$\mathrm{Vir}_c$ & $c = 0$ & $1$ & $\infty$
  & degenerate; tower open \\
$\cW^k(\fg, f)$ & zeros of $c_{\mathrm{DS}}(k)$
  & $\geq 1$ & $\infty$
  & case-dependent \\
Free fermions & $n = 0$ (trivial) & $0$ & $2$
  & trivial \\
$\beta\gamma$ & $\kappa = 1 \neq 0$ always & --- & ---
  & (never $\kappa = 0$) \\
Lattice $V_\Lambda$ & $\mathrm{rank}(\Lambda) = 0$ (trivial)
  & $0$ & $2$
  & trivial
\end{tabular}
\end{center}
The only nontrivial uncurved algebras are:
\begin{enumerate}[label=\textup{(\roman*)}]
\item Critical-level affine algebras
  $\widehat{\fg}_{-\hvee}$: Type~F, uncurved, with
  nontrivial genus-$g$ data from the Feigin--Frenkel center.
\item Virasoro at $c = 0$: Type~V, uncurved, with degenerate
  but nonzero higher-arity shadow obstruction tower.
\item Principal $\cW$-algebras at specific DS central charge zeros.
\end{enumerate}
\end{proposition}

\begin{proof}
The curvature formulas are:
\begin{align*}
\kappa(\cH_k) &= k, \\
\kappa(V_k(\fg)) &= \frac{(k + \hvee)\,\dim\fg}{2\hvee}, \\
\kappa(\mathrm{Vir}_c) &= \frac{c}{2}, \\
\kappa(V_\Lambda) &= \mathrm{rank}(\Lambda), \\
\kappa(\text{free fermions}) &= n/4, \\
\kappa(\beta\gamma) &= 1.
\end{align*}
Setting each to zero gives the stated conditions.
For Heisenberg: $\kappa = k = 0$ gives the trivial algebra
$\cH_0$ with no current at all.
For KM: $\kappa = 0$ forces $k = -\hvee$, the critical level.
For Virasoro: $\kappa = 0$ forces $c = 0$.
For free fermions and lattice algebras, $\kappa = 0$ only in
the trivial (rank zero) case.
For $\beta\gamma$: $\kappa = 1$ always, never zero.

The shadow depth data: $\cH_k$ and $V_k(\fg)$ are Type~F
(all OPE poles are simple or have central double poles) at
every level, including critical.
$\mathrm{Vir}_c$ has $\dalg = 1$ at every $c$ where the
non-central double pole $2T(w)$ persists.
At $c = 0$, the quartic pole $c/2 = 0$ vanishes but
the double pole $2T(w)/(z-w)^2$ remains non-central,
so $\dalg(\mathrm{Vir}_0) = 1$ (still Type~V).
\end{proof}

\subsection{The critical-level computation}

The central computation: at the critical level
$k = -\hvee$ for $\fg = \fsl_2$ (so $\hvee = 2$,
$k = -2$), we prove that the genus-$g$ bar cohomology
is nontrivial despite $\kappa = 0$.

\begin{theorem}[Critical-level KM has nontrivial genus-$g$ data;
uncurved but not trivial]
\label{thm:critical-level-nontrivial}
Let $\fg$ be a simple Lie algebra with dual Coxeter number
$\hvee$, and let $\cA = V_{-\hvee}(\fg)$ be the universal
affine vertex algebra at critical level.
\begin{enumerate}[label=\textup{(\roman*)}]
\item $\kappa(\cA) = 0$: the bar complex is uncurved at
  every genus.
\item $\dalg(\cA) = 0$: the algebra is Type~F (formal).
\item The genus-$g$ bar complex is a genuine cochain complex
  with $\dfib^{\,2} = 0$, and its cohomology is
  \[
  H^\bullet\bigl(\barB^{(g)}(V_{-\hvee}(\fg)),\, \dfib\bigr)
  \;\simeq\;
  C^\bullet\bigl(\fg\llbracket z \rrbracket,\,
  \fg((z));\, \cO(\mathrm{Op}_{\fg^\vee}(\Sigma_g))\bigr),
  \]
  the relative Lie algebra cohomology of the loop algebra
  with coefficients in the algebra of functions on
  $\fg^\vee$-opers on $\Sigma_g$.  This is nontrivial for
  all $g \geq 1$.
\item The scalar shadow free energies vanish:
  $F_g(V_{-\hvee}(\fg)) = 0$ for all $g \geq 1$.
  The nontrivial genus-$g$ data resides entirely in the
  \emph{non-scalar} bar cohomology, not in scalar genus
  amplitudes.
\end{enumerate}
\end{theorem}

\begin{proof}
Part~(i): $\kappa(V_k(\fg)) = (k+\hvee)\dim\fg/(2\hvee)$
vanishes at $k = -\hvee$.

Part~(ii): the current--current OPE $J^a(z)\,J^b(w) \sim
k\,\delta^{ab}/(z-w)^2 + f^{ab}{}_c J^c(w)/(z-w)$ has
central double pole and simple non-central pole.
At the critical level, $k = -\hvee$ and the double pole
coefficient is $-\hvee \delta^{ab}$, still a c-number.
Hence $\dalg = 0$.

Part~(iii): since $\dfib^{\,2} = 0$, the bar cohomology
$H^\bullet(\barB^{(g)}, \dfib)$ is well-defined (not
merely coderived).  The Feigin--Frenkel theorem identifies
the center $Z(V_{-\hvee}(\fg)) \cong
\mathrm{Fun}\,\mathrm{Op}_{\fg^\vee}(\bD)$ with functions
on $\fg^\vee$-opers on the formal disc.  The genus-$g$
conformal blocks of the critical-level algebra are computed
by the relative Lie algebra cohomology with oper
coefficients.  For $g \geq 1$ and $\fg$ simple, the space
$\mathrm{Op}_{\fg^\vee}(\Sigma_g)$ is a non-empty affine
space of dimension $(g-1)\dim\fg$ (for $g \geq 2$) or
$\mathrm{rank}(\fg)$ (for $g = 1$), and the cohomology
is nontrivial.

For the explicit case $\fg = \fsl_2$, $\hvee = 2$,
$k = -2$: the space of $\mathrm{PGL}_2$-opers on
$\Sigma_g$ is the space of projective connections, an
affine space of dimension $3(g-1)$ for $g \geq 2$,
and the function algebra $\cO(\mathrm{Op}(\Sigma_g))$
is a polynomial ring in $3(g-1)$ variables.
The bar cohomology is nontrivial: it contains
$\cO(\mathrm{Op}(\Sigma_g))$ as a direct summand.

Part~(iv): the scalar shadow free energy is
$F_g = \kappa \cdot \lambda_g^{\mathrm{FP}} = 0$
since $\kappa = 0$.
\end{proof}

\begin{remark}[The critical level as the Langlands point]
\label{rem:critical-langlands}
The nontrivial genus-$g$ data of $V_{-\hvee}(\fg)$ is
precisely the datum of the geometric Langlands program:
$\fg^\vee$-opers on $\Sigma_g$ are the spectral data of
the Langlands correspondence.  The critical-level vertex
algebra ``sees'' the Langlands dual through its genus-$g$
bar cohomology, even though its curvature (and hence its
scalar genus amplitudes) vanishes identically.

This is the correct avatar of the ``Type~H'' phenomenon:
not a secondary curvature $\kappa_2 \neq 0$
(which does not exist), but nontrivial bar cohomology
of an uncurved complex.  The critical level is the unique
point in the affine lineage where the curvature vanishes
and the Feigin--Frenkel center provides the replacement.
\end{remark}

\subsection{Virasoro at $c = 0$: uncurved Type V}

\begin{proposition}[$\mathrm{Vir}_0$ is uncurved Type V]
\label{prop:vir-0-uncurved}
The Virasoro algebra at $c = 0$ satisfies:
\begin{enumerate}[label=\textup{(\roman*)}]
\item $\kappa(\mathrm{Vir}_0) = 0$: uncurved at every genus.
\item $\dalg(\mathrm{Vir}_0) = 1$: Type~V, not Type~F.
  The OPE $T(z)\,T(w) \sim 2T(w)/(z-w)^2 + \partial T(w)/(z-w)$
  has a non-central double pole $2T(w)$, which produces
  a genuine arity-$3$ operation $m_3 \neq 0$.
\item The shadow obstruction tower is infinite ($r_{\max} = \infty$)
  but with vanishing scalar terms: $F_g(\mathrm{Vir}_0) = 0$
  for all $g \geq 1$.
\item The higher-arity shadow data
  $\Theta^{(g,r)}_{\mathrm{Vir}_0}$ for $r \geq 2$ is
  \emph{not} determined by $\kappa = 0$ alone.
  This is the content of the ``depth-zero resonance shadow''
  observation: the scalar budget closes but the non-scalar
  tower remains open.
\end{enumerate}
\end{proposition}

\begin{proof}
Parts~(i) and~(iii): $\kappa = c/2 = 0$, so
$F_g = (c/2)\lambda_g^{\mathrm{FP}} = 0$.

Part~(ii): at $c = 0$ the Virasoro OPE becomes
$T(z)\,T(w) \sim 2T(w)/(z-w)^2 + \partial T(w)/(z-w)$.
The quartic pole vanishes ($c/2 = 0$), but the
quadratic pole coefficient $2T(w)$ is field-valued
(non-central).  By the shadow depth classification, the
non-central double pole produces a genuine $m_3$, so
$\dalg(\mathrm{Vir}_0) = 1$.

Part~(iv): this is the content of the shadow obstruction tower analysis
at $c = 0$ in the critical string dichotomy chapter.
\end{proof}

\subsection{Resolution of the Type H question}

\begin{theorem}[Classification of uncurved-but-nontrivial algebras;
\ClaimStatusProvedHere]
\label{thm:type-H-resolution}
\begin{enumerate}[label=\textup{(\roman*)}]
\item \textbf{No secondary curvature exists}: the genus tower
  is governed by the single scalar $\kappa(\cA)$.  There is no
  $\kappa_2, \kappa_3, \ldots$ independent of $\kappa$.
  ``Type~H'' in the sense of $\kappa = 0$ with $\kappa_2 \neq 0$
  is vacuous.
\item \textbf{Uncurved-with-nontrivial-cohomology exists}: the
  critical-level affine algebra $V_{-\hvee}(\fg)$ has
  $\kappa = 0$ and nontrivial genus-$g$ bar cohomology for all
  $g \geq 1$.  The first explicit example is
  $V_{-2}(\fsl_2)$ \textup{(}$\fsl_2$ at the critical level
  $k = -2$\textup{)}, whose genus-$g$ bar cohomology contains
  $\cO(\mathrm{Op}_{\mathrm{PGL}_2}(\Sigma_g))$, the algebra
  of functions on projective connections.
\item \textbf{The three regimes of uncurved algebras}:
  \begin{enumerate}[label=\textup{(\alph*)}]
  \item \emph{Trivially uncurved}: $\kappa = 0$ from triviality
    ($\cH_0$, rank-$0$ lattice, zero fermions).  Both curvature
    and cohomology vanish.
  \item \emph{Critically uncurved} \textup{(}Langlands regime\textup{)}:
    $\kappa = 0$ from critical level
    ($\widehat{\fg}_{-\hvee}$).
    Curvature vanishes but bar cohomology is nontrivial,
    governed by the Feigin--Frenkel center and the oper space.
    Type~F, $\dalg = 0$.
  \item \emph{Degenerately uncurved} \textup{(}resonance regime\textup{)}:
    $\kappa = 0$ from parameter specialisation
    ($\mathrm{Vir}_0$, $\cW$-algebras at $c = 0$).
    Curvature vanishes, scalar genus amplitudes vanish, but the
    non-scalar shadow obstruction tower may be nontrivial.  Type~V or higher,
    $\dalg \geq 1$.
  \end{enumerate}
\item \textbf{Connection to geometric Langlands}: the critically
  uncurved regime is precisely the domain of the geometric
  Langlands program.  The nontrivial genus-$g$ data of
  $\widehat{\fg}_{-\hvee}$ is the Langlands spectral data,
  realized as bar cohomology of an uncurved Swiss-cheese complex.
  This provides the
  bridge between the modular tower of Vol~II and the Langlands
  program:
  \begin{center}
  \emph{The geometric Langlands program is the study of bar
  cohomology at the zero of the genus tower.}
  \end{center}
\end{enumerate}
\end{theorem}

\begin{proof}
Part~(i) is Theorem~\ref{thm:no-secondary-curvature}.
Part~(ii) is Theorem~\ref{thm:critical-level-nontrivial}.
Part~(iii) combines
Proposition~\ref{prop:uncurved-locus-complete}
with the preceding analysis.
Part~(iv) is the mathematical content of the Feigin--Frenkel
theorem reinterpreted through the Swiss-cheese lens.
\end{proof}

\subsection{Implications for the modular tower}

\begin{remark}[The sewing operator at $\kappa = 0$]
\label{rem:sewing-kappa-zero}
The MC recursion at genus~$g$ and arity~$0$ reads
\[
d_{\mathrm{mod}}\,\Theta^{(g,0)}
= -\tfrac{1}{2}
\sum_{g_1+g_2=g} [\Theta^{(g_1,0)}, \Theta^{(g_2,0)}]
- \Delta_{\mathrm{sew}}\,\Theta^{(g-1,0)}.
\]
When $\kappa = 0$, the seed $\Theta^{(1,0)} = \kappa \cdot \eta_1 = 0$
vanishes, and the recursion gives $\Theta^{(g,0)} = 0$ for all $g$
by induction.  The scalar genus tower is identically zero:
$F_g = 0$ for all $g \geq 1$.

But the \emph{non-scalar} MC components
$\Theta^{(g,r)}$ for $r \geq 2$ are not controlled by
this recursion alone.  The sewing operator
$\Delta_{\mathrm{sew}}$ at arity~$r$ receives contributions
from $\Theta^{(g-1,r+2)}$ (contracting two tensor legs),
and these higher-arity data are independent of $\kappa$.
This is the mechanism by which critical-level KM has
nontrivial genus-$g$ data: the oper space contributes to
$\Theta^{(g,r)}$ at $r \geq 2$ even though
$\Theta^{(g,0)} = 0$.
\end{remark}

\begin{remark}[The banana graph and genus-$2$ non-scalar data]
\label{rem:banana-graph}
The Vol~I result (prop:f2-quartic-dependence) establishes that
the banana graph (the unique genus-$2$ graph with two vertices
and two edges) injects the quartic shadow invariant $S_4$
into $F_2$ for multi-weight algebras.  For single-weight
algebras with $\kappa = 0$, the banana graph contribution
to $F_2$ vanishes (since $F_2 = \kappa \cdot \lambda_2 = 0$),
but the quartic data may still contribute to non-scalar
genus-$2$ amplitudes.

For critical-level KM ($\kappa = 0$, $\dalg = 0$,
$r_{\max} = 3$): the arity-$2$ MC component
$\Theta^{(2,2)}$ is the first non-scalar genus-$2$ datum.
Its computation requires the explicit genus-$2$ sewing of
the affine bar complex and is not carried out here.
\end{remark}

\begin{question}[The critical-level modular tower]
\label{q:critical-modular-tower}
At the critical level $k = -\hvee$, the scalar genus
amplitudes vanish ($F_g = 0$) but the non-scalar data is
governed by opers.  Is there a closed-form expression for
the full genus-$g$ MC element
$\Theta^{(g)}_{\widehat{\fg}_{-\hvee}}$
in terms of oper data?  Specifically:
\begin{enumerate}[label=\textup{(\alph*)}]
\item Does the bar cohomology
  $H^\bullet(\barB^{(g)}(V_{-\hvee}(\fg)), \dfib)$
  admit a purely algebraic description in terms of
  $\mathrm{Op}_{\fg^\vee}(\Sigma_g)$ at all genera?
\item Does the sewing/clutching structure on the bar
  cohomology recover the factorization structure of
  opers (the Beilinson--Drinfeld factorization)?
\item Is the critical-level modular tower
  $\{\Theta^{(g)}\}_{g \geq 1}$ controlled by
  finitely many invariants (since $r_{\max} = 3$), or
  does the genus direction introduce new independent data
  at each genus?
\end{enumerate}
These questions connect the Swiss-cheese modular machine
to the heart of the geometric Langlands program.
\end{question}

\begin{remark}[Relation to the existing Type~H conjecture]
\label{rem:type-H-conjecture-relation}
Conjecture~\ref{conj:type-H-existence} in the relative Feynman
transform chapter asks for algebras with $\dalg \geq 2$.
This is a \emph{genus-$0$} question about the $\Ainf$ structure,
independent of the curvature question studied here.
The two questions are logically independent:
\begin{center}
\renewcommand{\arraystretch}{1.2}
\begin{tabular}{ccc}
  & $\dalg = 0$ or $1$ & $\dalg \geq 2$ \\
\hline
$\kappa \neq 0$ & standard (Heisenberg, Virasoro)
  & \textbf{open} (conj:type-H-existence) \\
$\kappa = 0$ & critical KM, $\mathrm{Vir}_0$
  & \textbf{open}
\end{tabular}
\end{center}
The bottom-right cell --- $\kappa = 0$ \emph{and}
$\dalg \geq 2$ --- would be the strongest form of a
``Type~H'' algebra: uncurved, with genuinely independent higher
$\Ainf$ operations.  No example is known.  A candidate would
require OPE poles of order $\geq 3$ with non-central
higher-order residues \emph{and} vanishing total curvature $\kappa = 0$.
The non-linear $\cW_\infty$ algebras are the natural
place to search.
\end{remark}




% =================================================================
\section{Soft graviton theorems from the $\Ainf$ tower}
\label{sec:soft-graviton-ainf-tower}
% =================================================================

The shadow obstruction tower of the Virasoro algebra
$\mathrm{Vir}_c$ encodes an infinite hierarchy of soft
graviton theorems (in the celestial/gravitational reading of
the shadow connection $\nabla^{\mathrm{hol}}_{g,n} =
d - \mathrm{Sh}_{g,n}(\Theta_{\mathrm{Vir}_c})$).
The arity-$r$ Ward identity of this connection produces
the order-$(r-2)$ soft factor $S^{(r-2)}$.  This section
derives all four layers of this hierarchy from the explicit
$\lambda$-bracket data and the $\Ainf$ structure of
Movement~I, focusing on the \emph{origin} of each soft factor
in the bar complex.

The fundamental input is the single $\lambda$-bracket
\begin{equation}\label{eq:soft-input}
\{T{}_\lambda T\}
\;=\;
\partial T + 2T\lambda + \frac{c}{12}\lambda^3.
\end{equation}
The three terms have conformal weight $2, 2, 0$ respectively,
and the three OPE pole orders are $1, 2, 4$ (with pole order 3
absent: there is no cubic pole in the Virasoro OPE, reflecting
the fact that $T$ is quasi-primary).

\subsection{The soft limit as a $\lambda$-bracket expansion}

The ``soft limit'' of a graviton insertion $T(z_s)$ in an
$(n+1)$-point correlator $\Phi(z_1, \ldots, z_n, z_s)$ is
the regime $z_s \to \infty$.  In the shadow-tower language,
the soft expansion is
\begin{equation}\label{eq:soft-expansion-master}
\Phi
\;=\;
\sum_{k=0}^{\infty}
\frac{1}{z_s^k}\;S^{(k)}\;\Phi_n(z_1, \ldots, z_n),
\end{equation}
where each $S^{(k)}$ is a differential operator acting on
the $n$-point correlator $\Phi_n$.  The arity decomposition
of the shadow connection produces $S^{(k)}$ from the
arity-$(k+2)$ shadow $\mathrm{Sh}_{k+2}(\mathrm{Vir}_c)$.
The connection to the $\Ainf$ tower is direct:
\begin{center}
\renewcommand{\arraystretch}{1.15}
\begin{tabular}{lccl}
\hline
\textbf{Soft order} & \textbf{Arity} & \textbf{$\Ainf$ source}
  & \textbf{Bar-complex origin} \\
\hline
Leading $S^{(0)}$ & $2$ & $m_2$ at $\lambda \to 0$
  & simple pole $\partial T / (z_s - z_i)$ \\
Subleading $S^{(1)}$ & $3$ & $m_2$ at $O(\lambda)$ $+$ $m_3$
  & double pole $2T / (z_s - z_i)^2$ \\
Sub-subleading $S^{(2)}$ & $4$ & $m_4$ $+$ $m_3^2$
  & quartic contact $\mathfrak{Q}^{\mathrm{contact}}$ \\
$S^{(r-2)}$ & $r$ & $m_r$ $+$ lower compositions
  & arity-$r$ shadow $\mathrm{Sh}_r$ \\
\hline
\end{tabular}
\end{center}

\subsection{Derivation of Weinberg's leading soft graviton theorem}

\noindent\textbf{Claim.} The leading soft factor $S^{(0)}$
is determined by the simple-pole residue of the $\lambda$-bracket,
i.e.\ by $m_2$ evaluated at $\lambda = 0$.

\begin{proof}[Computation]
The modular characteristic of $\mathrm{Vir}_c$ is
$\kappa(\mathrm{Vir}_c) = c/2$, the arity-$2$ shadow.
Its action on the genus-$0$ transport package is
\[
\mathrm{Sh}_{0,n}(\kappa) \cdot \Phi
\;=\;
\sum_{i < j}
\frac{\kappa(\phi_i, \phi_j)}{z_i - z_j}\;\Phi.
\]
For the stress tensor OPE, the bilinear form
$\kappa(\phi_i, \phi_j)$ is the residue of $m_2$ at
$\lambda = 0$:
\[
m_2(T, T;\, \lambda)\big|_{\lambda = 0}
\;=\; \partial T
\;=\; T_{(0)}T.
\]
This is the simple-pole mode.  The resulting Ward identity
$\mathrm{Sh}_{0,n+1}(\kappa)\,\Phi = 0$ gives, in the soft
limit $z_s \to \infty$:
\begin{equation}\label{eq:weinberg-from-m2}
\boxed{
S^{(0)}
\;=\;
\sum_{i=1}^{n}
\frac{h_i}{z_s - z_i},}
\end{equation}
where $h_i$ is the conformal weight of the $i$-th hard insertion.
In the celestial/gravitational reading, this is the Weinberg soft graviton factor: the emission amplitude for a soft graviton of vanishing energy is determined by the total energy-momentum of the hard particles.

The bar-complex interpretation: $S^{(0)}$ is the differential $d_{\bar{B}}$ acting on degree-$2$ elements $[s^{-1}T \,|\, s^{-1}\phi_i]$ of the bar complex $\bar{B}(\mathrm{Vir}_c)$.  The boundary integral over the divisor $E_{12} \subset \overline{C}_2(\mathbb{P}^1)$ extracts exactly the simple-pole mode $T_{(0)}\phi_i = h_i\,\phi_i + z_i\partial\phi_i$.  The universality of the Weinberg factor is the universality of the bar differential at arity~$2$: every chiral algebra has $d_{\bar{B}}$ at degree~$2$, and the modular characteristic $\kappa$ is its shadow.
\end{proof}

\subsection{Derivation of the Cachazo--Strominger subleading soft graviton theorem}

\noindent\textbf{Claim.} The subleading soft factor $S^{(1)}$ is determined by the $O(\lambda)$ term of $m_2$ together with the ternary operation $m_3$.

\begin{proof}[Computation]
The $O(\lambda)$ term in the $\lambda$-bracket is the double-pole mode:
\[
m_2(T, T;\, \lambda)\big|_{O(\lambda)}
\;=\; 2T\,\lambda
\;=\; T_{(1)}T \cdot \lambda.
\]
The mode $T_{(1)}\phi_i = h_i\,\phi_i$ gives the conformal weight.  At the level of the shadow connection, the cubic shadow $\mathfrak{C}(\mathrm{Vir}_c) = 2x^3$ encodes the ternary shadow (the $m_3$ contribution from the non-associativity of $m_2$).  The arity-$3$ Ward identity gives, in the soft limit:
\begin{equation}\label{eq:cs-from-m2m3}
\boxed{
S^{(1)}
\;=\;
\sum_{i=1}^{n}
\left[
\frac{h_i(h_i - 1)}{(z_s - z_i)^2}
\;+\;
\frac{h_i\,\partial_{z_i}}{z_s - z_i}
\right].}
\end{equation}
The two terms have distinct $\Ainf$ origins:
\begin{itemize}
\item The $h_i\,\partial_{z_i}/(z_s - z_i)$ term is the \emph{orbital} part, coming from the $O(\lambda)$ expansion of the binary shadow $\kappa$.
\item The $h_i(h_i-1)/(z_s-z_i)^2$ term is the \emph{spin} part, coming from the cubic shadow $\mathfrak{C} = 2x^3$, i.e.\ from $m_3$.  The non-associativity $A_3 \ne 0$ (forced by the quartic pole) generates the cubic shadow via $\mathfrak{C} = -h(o_3)$.  The polynomial $h_i(h_i-1)$ arises from the action of $:TT:$ on weight-$h_i$ primaries.
\end{itemize}
In the celestial/gravitational reading, $S^{(1)}$ is the Cachazo--Strominger subleading soft graviton factor.  The angular-momentum operator $J^i_{\mu\nu}$ splits into orbital ($\partial_{z_i}$) and spin ($h_i$) parts.  The BMS supertranslation Ward identity is recovered from the arity-$2$ Ward identity; the superrotation Ward identity from the arity-$3$ Ward identity.
\end{proof}

\subsection{The sub-subleading soft graviton theorem from $m_4$}

\noindent\textbf{Claim.} The sub-subleading soft factor $S^{(2)}$ is determined by the quartic $\Ainf$ operation $m_4$ (the quartic contact invariant) together with the self-interaction of the cubic shadow $\mathfrak{C}$.

\begin{proof}[Computation]
The arity-$4$ flatness equation reads
\[
\nabla^{(2)}(A^{(4)})
\;+\;
[A^{(3)}, A^{(3)}]
\;+\; o_4
\;=\; 0.
\]
The quartic contact invariant for the Virasoro algebra, computed from the Kac determinant at weight~$4$, is
\begin{equation}\label{eq:quartic-from-m4}
\mathfrak{Q}^{\mathrm{contact}}_{\mathrm{Vir}}
\;=\;
\frac{10}{c(5c + 22)}.
\end{equation}
This coefficient arises from $m_4$: the Stasheff relation $\partial m_4 + m_3 \circ m_2 + m_2 \circ m_3 = 0$ on the associahedron $K_4$ determines $m_4$ uniquely (up to gauge).  The numerator~$10$ is $\langle :T^2:, o_4 \rangle$; the denominator $c(5c+22)$ is the Kac determinant at weight~$4$.

The sub-subleading soft factor in the soft limit is:
\begin{equation}\label{eq:sub-subleading-from-m4}
\boxed{
S^{(2)}
\;=\;
\frac{10}{c(5c+22)}
\sum_{i=1}^{n}
\frac{h_i^2(h_i - 1)^2}{(z_s - z_i)^3}
\;+\; \text{(lower-order terms and cross-terms)}.}
\end{equation}
The conformal-weight polynomial $h_i^2(h_i - 1)^2$ arises from the action of $:T^2:$ on weight-$h_i$ states.
\end{proof}

\subsection{The $\Ainf$ tower and the infinite soft hierarchy}

\begin{proposition}[$\Ainf$ tower $\Longleftrightarrow$ soft graviton tower]\label{claim:ainf-soft-tower}
For the Virasoro algebra $\mathrm{Vir}_c$ at generic~$c$:
\begin{enumerate}
\item[(a)] The order-$k$ soft factor $S^{(k)}$ is determined by the shadows $\{\mathrm{Sh}_2, \ldots, \mathrm{Sh}_{k+2}\}$, which are in turn determined by the $\Ainf$ operations $\{m_2, \ldots, m_{k+2}\}$ together with the homotopy-transfer data.
\item[(b)] For $k \ge 2$, the leading term of $S^{(k)}$ has the form
  $S^{(k)}_{\mathrm{leading}} = R_k(c) \cdot \sum_{i=1}^n P_k(h_i)/(z_s - z_i)^{k+1}$,
  where $R_k(c)$ is a rational function of $c$ with denominator $\prod_{j=2}^{k+2} \det G_j(c)$, and $P_k(h_i)$ is a polynomial of degree $2k$ in $h_i$.
\item[(c)] The tower is genuinely infinite: $S^{(k)} \ne 0$ for all $k \ge 0$ and generic~$c$.  This follows from $m_n \ne 0$ for all $n \ge 3$ (the quartic pole forces infinite $\Ainf$ depth) and the non-degeneracy of the Gram determinants at generic~$c$.
\item[(d)] Each successive soft factor is suppressed by one power of $c$: $S^{(k)} = O(c^{-k})$.
\end{enumerate}
\end{proposition}

\begin{proof}[Verification]
Part~(a): the arity-$r$ shadow $\mathrm{Sh}_r$ is computed from the obstruction recursion $o_r = (D\Theta^{\le r-1} + \frac{1}{2}[\Theta^{\le r-1}, \Theta^{\le r-1}])_r$, $\mathrm{Sh}_r = -h(o_r)$, where each $o_r$ is built from $m_2, \ldots, m_r$.

Part~(b): propagator scaling on $\overline{\mathcal{M}}_{0,r+1}$ gives a pole of order $r-1$ in $1/(z_s - z_i)$ in the soft limit.

Part~(c): the ``cubic source mechanism'' --- the cubic shadow $\mathfrak{C} = 2x^3 \ne 0$ acts as a permanent source for all higher obstructions via $o_{r+1} \ni [\mathfrak{C}, \mathrm{Sh}_r]$.

Part~(d): $m_n \sim c^{-(n-2)}$ at large~$c$, and the arity-$r$ shadow involves $r - 2$ propagator contractions each contributing $1/c$.
\end{proof}

\subsection{The 4th soft graviton theorem from $m_5$}

The first \emph{new} soft graviton theorem beyond the Cachazo--He--Yuan sub-subleading result arises from the arity-$5$ shadow, controlled by $m_5$ and the quintic obstruction $o_5(\mathrm{Vir}_c) \ne 0$.

\begin{computation}[The 4th-order soft graviton factor]
\label{comp:4th-soft-graviton}
The arity-$5$ soft factor is determined by the recursion
$S^{(3)} = \mathrm{Sh}_{0,n}(\mathrm{Sh}_5)|_{\mathrm{soft}} + [\mathrm{Sh}_{0,n}(\mathfrak{C}), \mathrm{Sh}_{0,n}(\mathfrak{Q})]|_{\mathrm{soft}}$.

\medskip\noindent
\emph{Step~1: The quintic obstruction.}  The bracket $[\mathfrak{C}, \mathfrak{Q}]$ is computed by the $H$-Poisson bracket:
$\{2x^3, \frac{10}{c(5c+22)}\,x^4\}_H = \frac{20}{c(5c+22)}\;P_{\mathrm{Vir}}(x^3 \cdot x^4)$,
where $P_{\mathrm{Vir}} = H_{\mathrm{Vir}}^{-1}$ is the complementarity propagator.

\medskip\noindent
\emph{Step~2: The Kac determinant at weight~$5$.}
$\det G_5(c) \propto c^2(5c + 22)(2c - 1)(7c + 68)(3c + 46)$.
The new weight-$5$ factors contribute to cross-term denominators but not to the leading single-insertion coefficient.

\medskip\noindent
\emph{Step~3: The soft factor.}
\begin{equation}\label{eq:4th-soft-from-m5}
\boxed{
S^{(3)}_{\mathrm{Vir}}
\;=\;
\frac{20}{c^2(5c+22)}
\sum_{i=1}^{n}
\frac{h_i^3(h_i-1)^2(2h_i - 1)}{(z_s - z_i)^4}
\;+\; \text{(cross-terms and descendants)}.}
\end{equation}
The conformal-weight polynomial $h_i^3(h_i-1)^2(2h_i-1)$ arises from $:T^3:$ acting on weight-$h_i$ primaries.  The coefficient $20/[c^2(5c+22)] = Q^{\mathrm{contact}} \cdot (2/c) = [10/(c(5c+22))] \cdot (2/c)$ confirms that the 4th soft factor comes from the cubic-quartic bracket $[\mathfrak{C}, \mathfrak{Q}]$ with one additional propagator contraction.

\medskip\noindent
\textbf{Correction (level-6 computation).}  The bracket $[\mathfrak{C}, \mathfrak{Q}]_H$ captures only one sub-channel.  The full shadow coefficient from the shadow-metric expansion $H(t) = t^2\sqrt{Q(t)}$ is $S_5 = -48/[c^2(5c+22)]$, not $20/[c^2(5c+22)]$.  The remaining contributions (the $\mathrm{Sh}_5$ term and further cross-terms in the recursion) have opposite sign and larger magnitude.  The Stasheff cancellation at level~$6$ eliminates the Shapovalov poles $c = 1/2$ and $c = -68/7$, leaving only the shadow-metric denominator $c^2(5c+22)$.
\end{computation}

\subsection{Summary: the soft graviton dictionary}

\begin{center}
\renewcommand{\arraystretch}{1.3}
\begin{tabular}{lcccl}
\hline
\textbf{Soft order} & \textbf{Arity} & \textbf{$\Ainf$}
  & \textbf{Coefficient} & \textbf{Name} \\
\hline
$S^{(0)}$ & $2$ & $m_2|_{\lambda=0}$
  & $\kappa = c/2$ & Weinberg (1965) \\
$S^{(1)}$ & $3$ & $m_2|_{O(\lambda)} + m_3$
  & $\mathfrak{C} = 2x^3$ & Cachazo--Strominger (2014) \\
$S^{(2)}$ & $4$ & $m_4 + m_3^2$
  & $10/[c(5c+22)]$ & Sub-subleading (CHY 2014) \\
$S^{(3)}$ & $5$ & $m_5 + m_3 \cdot m_4$
  & $-48/[c^2(5c+22)]$ & \textbf{New: 4th soft theorem} \\
$S^{(k)}$ & $k+2$ & $m_{k+2} + \text{lower}$
  & $O(c^{-k})$ & Tower from $\Ainf$ \\
\hline
\end{tabular}
\end{center}

\begin{remark}[The quartic pole as the engine of soft graviton complexity]
The entire infinite tower traces to the quartic pole in the Virasoro OPE --- equivalently, to the $\lambda^3$ term in the $\lambda$-bracket.  For affine Kac--Moody algebras (where the maximal OPE pole is double, $\lambda^1$), the $\Ainf$ structure truncates at $m_2$ and the soft hierarchy terminates at the subleading level.  For the Virasoro algebra, the quartic pole forces $m_3 \ne 0$, which cascades through the Stasheff identities to produce $m_k \ne 0$ for all $k$, and hence $S^{(k)} \ne 0$ for all $k$.  The statement ``infinitely many soft graviton theorems'' is the same statement as ``the Virasoro $\Ainf$ structure does not truncate,'' which is the same statement as ``the quartic pole in the stress-tensor OPE forces infinite shadow depth.''

This identifies the \emph{single datum} that controls the gravitational soft hierarchy: the OPE pole order of $T(z)T(w)$.  Pole order $\le 2$ (Kac--Moody): finitely many soft theorems.  Pole order $4$ (Virasoro): infinitely many.  The pole order is the shadow depth class.
\end{remark}

\begin{remark}[What the computation proves and what it does not]
The derivation proves: the leading (Weinberg) and subleading (Cachazo--Strominger) soft graviton theorems follow from $\kappa$ and $\mathfrak{C}$; the sub-subleading from $\mathfrak{Q}^{\mathrm{contact}} = 10/[c(5c+22)]$; the 4th soft theorem exists with shadow coefficient $S_5 = -48/[c^2(5c+22)]$ (the earlier heuristic $[\mathfrak{C},\mathfrak{Q}]_H$ gave $20/[c^2(5c+22)]$, which is a sub-channel); the infinite tower $S^{(k)}$ is nonzero for all $k$ at generic~$c$, with $S^{(k)} = O(c^{-k})$.

What it does \emph{not} prove: the full closed-form expression for $S^{(k)}$ at arbitrary~$k$.  The recursion $S^{(k)} = \mathrm{Sh}_{0,n}(\mathrm{Sh}_{k+2})|_{\mathrm{soft}} - \sum_{s+t = k+2,\, 3 \le s \le t} S^{(s-2)} \cdot S^{(t-2)}|_{\mathrm{cross}}$ reduces to Shapovalov matrix inversion at weight~$(k+2)$, of dimension $p(k+2)$.  For $k = 3$ ($p(5) = 7$): tractable.  For large~$k$: $p(k)$ grows exponentially and closed forms are not expected.
\end{remark}

\begin{remark}[Falsification criteria]
Three predictions are falsifiable by explicit computation:
\begin{enumerate}
\item The shadow coefficient is $S_5 = -48/[c^2(5c+22)]$, not $-48/[c^2(5c+22)(7c+68)]$.  If the weight-$5$ Kac factor $(7c+68)$ appears in the \emph{leading} coefficient, the Stasheff cancellation fails.  (The earlier heuristic gave $20/[c^2(5c+22)]$ from the $[\mathfrak{C},\mathfrak{Q}]_H$ sub-channel; the full shadow-metric expansion corrects both the sign and magnitude.)
\item The conformal-weight polynomial at 4th order is $h^3(h-1)^2(2h-1)$ at leading order.  This can be verified by computing $\langle h | {:}T{:}^3(z) | h \rangle$.
\item The large-$c$ scaling $S^{(k)} \sim c^{-k}$ can be verified against the WKB/geodesic approximation of AdS$_3$ scattering amplitudes.
\end{enumerate}
\end{remark}




% =================================================================
\section{$\fsl_3$ monodromy vs.\ bar-complex $R$-matrix: first quantum test}
\label{sec:sl3-monodromy-test}
% =================================================================

The monodromy of the KZ connection on conformal blocks of the
affine vertex algebra $V^k(\fg)$ is identified with the spectral
$R$-matrix from the bar complex for the full affine lineage
(Theorem~\ref{thm:affine-monodromy-identification} of the core
chapter on logarithmic HT monodromy).  For $\fg = \fsl_2$, the
identification is controlled by a single quadratic Casimir; for
$\fg = \fsl_3$, the cubic Casimir $C_3$ enters and the interplay
between $C_2$ and $C_3$ is the first genuinely higher-rank
phenomenon.  This section carries out the explicit test.

\subsection{Setup and conventions}
\label{subsec:sl3-setup}

Let $\fg = \fsl_3$ with dual Coxeter number $h^\vee = 3$.
Write $\omega_1, \omega_2$ for the fundamental weights and
$\rho = \omega_1 + \omega_2$ for the Weyl vector.  The
inner product on $\fh^*$ induced by the trace form
$B(X,Y) = \Tr(XY)$ in the fundamental representation gives
\[
  (\omega_i, \omega_j) = (A^{-1})_{ij},
  \qquad
  A^{-1} = \frac{1}{3}\begin{pmatrix} 2 & 1 \\ 1 & 2 \end{pmatrix},
\]
where $A$ is the $\fsl_3$ Cartan matrix.  With this normalisation:
\begin{alignat*}{3}
  (\omega_1, \omega_1) &= \tfrac{2}{3}, &\quad
  (\omega_1, \omega_2) &= \tfrac{1}{3}, &\quad
  (\omega_2, \omega_2) &= \tfrac{2}{3}.
\end{alignat*}

The quadratic Casimir eigenvalue on the irreducible representation
$V_\lambda$ with highest weight $\lambda$ is
$c_2(\lambda) = (\lambda, \lambda + 2\rho)$.

\begin{observation}[Quadratic Casimir values for low representations]
\label{obs:sl3-casimir-table}
\begin{center}
\begin{tabular}{cccc}
  $\lambda$ & $V_\lambda$ & $\dim$ & $c_2(\lambda)$\\[2pt]
  \hline\\[-8pt]
  $(1,0)$ & fundamental $V$ & $3$ & $8/3$\\
  $(0,1)$ & anti-fundamental $V^*$ & $3$ & $8/3$\\
  $(2,0)$ & $\mathrm{Sym}^2 V$ & $6$ & $20/3$\\
  $(0,2)$ & $\mathrm{Sym}^2 V^*$ & $6$ & $20/3$\\
  $(1,1)$ & adjoint $\fg$ & $8$ & $6$\\
  $(2,1)$ & $V_{(2,1)}$ & $15$ & $32/3$\\
\end{tabular}
\end{center}
\end{observation}

The cubic Casimir of $\fsl_3$ is
$C_3 = d_{abc}\, t^a t^b t^c$ where $d_{abc}$ is the totally symmetric
invariant tensor (normalised so that $d_{abc} = \frac{1}{4}\Tr(t_a\{t_b, t_c\})$
in the fundamental representation).
Its eigenvalue on $V_{(p,q)}$ is
\begin{equation}\label{eq:cubic-casimir}
  c_3(p,q)
  = \frac{1}{18}(p - q)(2p + q + 3)(p + 2q + 3).
\end{equation}

\begin{observation}[Cubic Casimir values]
\label{obs:sl3-cubic-table}
\begin{center}
\begin{tabular}{ccc}
  $\lambda$ & $V_\lambda$ & $c_3(\lambda)$ \\[2pt]
  \hline\\[-8pt]
  $(1,0)$ & $V$ & $10/9$ \\
  $(0,1)$ & $V^*$ & $-10/9$ \\
  $(2,0)$ & $\mathrm{Sym}^2 V$ & $35/9$ \\
  $(0,2)$ & $\mathrm{Sym}^2 V^*$ & $-35/9$ \\
  $(1,1)$ & adjoint & $0$\\
  $(2,1)$ & $V_{(2,1)}$ & $70/9$ \\
\end{tabular}
\end{center}
Note the antisymmetry $c_3(p,q) = -c_3(q,p)$: the cubic
Casimir distinguishes a representation from its contragredient.
\end{observation}


\subsection{The Steinberg connection for $V^k(\fsl_3)$}
\label{subsec:sl3-steinberg}

The split Casimir of $\fsl_3$ with the trace-form inner product is
\[
  \Omega = \sum_{i \neq j} E_{ij} \otimes E_{ji}
  \;+\; \sum_{a,b=1}^{2} (A^{-1})_{ab}\, H_a \otimes H_b
  \;\in\; \fg \otimes \fg,
\]
where $\{E_{ij}\}_{i\neq j}$ are the off-diagonal elementary matrices
and $H_1 = E_{11} - E_{22}$, $H_2 = E_{22} - E_{33}$ are the
simple coroots.

The affine vertex algebra $V^k(\fsl_3)$ at generic level $k$ has
the OPE
\[
  J^a(w)\, J^b(z)
  \;\sim\;
  \frac{k\, B^{ab}}{(w-z)^2}
  \;+\;
  \frac{f^{ab}{}_c\, J^c(z)}{w-z},
\]
where $B^{ab} = \Tr(t^a t^b)$ and $f^{ab}{}_c$ are the structure
constants.  By the bar-kernel absorption principle
(AP19: the $d\log$ kernel absorbs one pole), the collision
$r$-matrix has a simple pole:
\[
  r(z) = \frac{\Omega}{z} + O(1).
\]
The Steinberg connection on the family of line-operator
self-products $\mathfrak{S}_b \to D^*$ is therefore
\begin{equation}\label{eq:sl3-steinberg}
  \nabla_{\mathfrak{S}}
  \;=\; d \;+\; \frac{1}{k+3}\,\Omega\;\frac{dz}{z}.
\end{equation}
This is the KZ connection for $\fsl_3$ at level $k$, and its
flatness is equivalent to the infinitesimal braid relations
for $\Omega$.


\subsection{The $R$-matrix on $V \otimes V$: first test}
\label{subsec:sl3-VxV}

Let $V = V_{(1,0)} = \C^3$ be the fundamental representation.
The Clebsch--Gordan decomposition is
\[
  V \otimes V
  \;\cong\;
  \underbrace{V_{(2,0)}}_{\mathrm{Sym}^2 V,\;\dim\, 6}
  \;\oplus\;
  \underbrace{V_{(0,1)}}_{\Lambda^2 V,\;\dim\, 3}.
\]
The split Casimir eigenvalues on the two summands are:
\begin{align*}
  \Omega\big|_{V_{(2,0)}}
  &= \tfrac{1}{2}\bigl(c_2(2,0) - 2\,c_2(1,0)\bigr)
   = \tfrac{1}{2}\bigl(\tfrac{20}{3} - \tfrac{16}{3}\bigr)
   = \tfrac{2}{3},\\[4pt]
  \Omega\big|_{V_{(0,1)}}
  &= \tfrac{1}{2}\bigl(c_2(0,1) - 2\,c_2(1,0)\bigr)
   = \tfrac{1}{2}\bigl(\tfrac{8}{3} - \tfrac{16}{3}\bigr)
   = -\tfrac{4}{3}.
\end{align*}

\begin{proposition}[$V \otimes V$ monodromy for $\fsl_3$]
\label{prop:sl3-VxV-monodromy}
The monodromy of the Steinberg connection \eqref{eq:sl3-steinberg}
around $z = 0$, acting on $V \otimes V$, is diagonal on the
Clebsch--Gordan summands:
\begin{align*}
  \mathrm{Mon}_{z=0}(\nabla_{\mathfrak{S}})\big|_{V_{(2,0)}}
  &= \exp\!\Bigl(\frac{4\pi i}{3(k+3)}\Bigr)
   = q^{4/3},\\[4pt]
  \mathrm{Mon}_{z=0}(\nabla_{\mathfrak{S}})\big|_{V_{(0,1)}}
  &= \exp\!\Bigl(\frac{-8\pi i}{3(k+3)}\Bigr)
   = q^{-8/3},
\end{align*}
where $q = e^{i\pi/(k+3)}$.
\end{proposition}

\begin{proof}
The monodromy is $\exp(2\pi i \cdot \Omega/(k+3))$.
Substituting the eigenvalues gives the stated formulae.
\end{proof}

\begin{proposition}[Agreement with $U_q(\fsl_3)$ $R$-matrix on $V \otimes V$]
\label{prop:sl3-VxV-qgroup}
The quantum group $\check{R}$-matrix of $U_q(\fsl_3)$ at
$q = e^{i\pi/(k+3)}$ acts on $V \otimes V$ with eigenvalues
$\check{R}|_{V_{(2,0)}} = q^{4/3}$ and
$\check{R}|_{V_{(0,1)}} = -q^{-8/3}$,
where the sign on the antisymmetric piece is absorbed by
the permutation operator $P$ (since $P|_{\Lambda^2 V} = -1$).
Projectively, $\overline{\rho}^{\mathrm{HT}} = \overline{\rho}^{U_q}$.
\end{proposition}

\begin{proof}
The universal $R$-matrix of $U_q(\fsl_3)$ gives
$\check{R}|_{V_\nu} = (-1)^{|\lambda|+|\mu|-|\nu|}\,
q^{c_2(\nu) - c_2(\lambda) - c_2(\mu)}$
on the isotypic component $V_\nu \subset V_\lambda \otimes V_\mu$.
For $V \otimes V$: on $V_{(2,0)}$, $\check{R} = q^{4/3}$;
on $V_{(0,1)}$, $\check{R} = -q^{-8/3}$.
The Drinfeld--Kohno identification gives
$\mathrm{Mon} = P \circ \check{R}^{-1}$ up to framing;
the sign on $V_{(0,1)}$ is $P|_{\Lambda^2} = -1$.
\end{proof}

\begin{remark}[Comparison with $\fsl_2$]
For $\fsl_2$, $V_{1/2} \otimes V_{1/2} = V_1 \oplus V_0$ with
$\Omega$-eigenvalues $1/4$ and $-3/4$.  The cubic Casimir $C_3$
does not appear in the $V \otimes V$ monodromy for either
$\fsl_2$ or $\fsl_3$.
\end{remark}


\subsection{The associator and the cubic Casimir}
\label{subsec:sl3-associator}

The cubic Casimir of $\fsl_3$ first enters the monodromic
structure through the \textbf{associator}, not through the
$R$-matrix eigenvalues.  The Drinfeld--KZ associator on
$V^{\otimes 3}$ has the expansion
$\Phi_{\mathrm{KZ}} = 1 - \zeta(2)\hbar^2 [\Omega_{12},\Omega_{23}]/(2\pi i)^2
+ \zeta(3)\hbar^3 \mathcal{A}/(2\pi i)^3 + O(\hbar^4)$,
where $\mathcal{A} = [\Omega_{12},[\Omega_{12},\Omega_{23}]]
- [\Omega_{23},[\Omega_{23},\Omega_{12}]]$.  The Fierz identity
for $\fsl_N$ decomposes $f_{ace}f_{bde}$ into a $1/h^\vee$ piece
and a $d$-tensor piece.  For $\fsl_2$ the $d$-tensor vanishes;
for $\fsl_3$ it is nonzero and produces off-diagonal mixing between
the two copies of $V_{(1,1)}$ in $\mathbf{3}^{\otimes 3}
= \mathbf{10} \oplus 2\cdot\mathbf{8} \oplus \mathbf{1}$,
with coefficient proportional to
$\Delta c_3 = c_3(2,0) - c_3(0,1) = 5$.

By one-loop exactness of $V^k(\fsl_3)$, the bar-complex
reduced connection is the KZ connection, so the bar-complex
associator reproduces the full Drinfeld--KZ associator including
the cubic Casimir terms.  The conjecture (monodromy = spectral
$R$-matrix) holds for $\fsl_3$ on evaluation modules.

\begin{warning}[The coincidence trap]
It is tempting to conclude from the $V \otimes V$ test that the
$\fsl_3$ monodromy is ``the same as $\fsl_2$.''  This is false
at the level of the associator (order $\hbar^3$) and at the
level of the $R$-matrix on representations with outer
multiplicities.  The quadratic Casimir controls the \emph{diagonal}
part; the cubic Casimir controls the \emph{off-diagonal} part.
The meta-principle of AP-33 applies: never trust a coincidence
between $C_2$ and $C_3$ at low representation.
\end{warning}





%% ====================================================================
%%  KOSZUL DUAL OF W_3: THE CENTRAL CHARGE INVOLUTION
%% ====================================================================

\section{Koszul dual of $\cW_3$: the central charge involution}
\label{sec:koszul-dual-W3}

This section computes the Koszul dual $\cW_3^!$ of the $\cW_3$
algebra from first principles, using Verdier duality on the bar
coalgebra $\barB(\cW_3)$ and the BPZ inner product on weight
spaces.  The result is the central theorem of this note:

\begin{theorem}[Koszul dual of $\cW_3$; proved here]
\label{thm:koszul-dual-W3}
The Koszul dual of the $\cW_3$ algebra at central charge~$c$ is
\begin{equation}\label{eq:W3-koszul-dual}
\boxed{\;
\cW_{3,c}^!
\;\simeq\;
\cW_{3,\,100-c},
\;}
\end{equation}
with Koszul involution $c \mapsto 100 - c$.  The self-dual point
is $c = 50$.  The structure constant $\beta^2 = 16/(22+5c)$
transforms to $\beta'^{\,2} = 16/(522 - 5c)$.
\end{theorem}

The proof occupies the rest of this section.

\subsection{Overview: from Virasoro to $\cW_3$}

Recall the Virasoro case: $\mathrm{Vir}_c^! = \mathrm{Vir}_{26-c}$.
The number $26$ arises as minus the ghost central charge for a
spin-$2$ field: a $bc$-ghost system of conformal weights $(2, -1)$
has $c_{\mathrm{gh}}(s=2) = -2(6 \cdot 4 - 6 \cdot 2 + 1) = -26$.
Koszul duality replaces the BV-BRST propagator with its Verdier
dual; the ghost sector provides the pairing, and the ghost central
charge governs the dual central charge shift.

For $\cW_3$, the algebra has two generators:
\begin{itemize}
\item $T$ of spin $s_T = 2$ (stress tensor),
\item $W$ of spin $s_W = 3$ (higher-spin current).
\end{itemize}
Each generator contributes its own ghost system upon BV-BRST
quantization of the Khan--Zeng 3D HT Poisson sigma model.  The
total ghost central charge determines the complementarity
constant $\alpha_3$, and hence the Koszul involution.

\subsection{The ghost central charge formula}

\begin{proposition}[Ghost central charge per spin]
\label{prop:ghost-central-charge}
A $bc$-ghost system of conformal weights $(s, 1-s)$ has central
charge
\begin{equation}\label{eq:ghost-cc}
c_{\mathrm{gh}}(s) \;=\; -2(6s^2 - 6s + 1).
\end{equation}
\end{proposition}

\begin{proof}
Standard: the $bc$ system of weights $(\lambda, 1-\lambda)$ with
$\lambda = s$ has central charge $c = -2(6\lambda^2 - 6\lambda + 1)$
by the Virasoro mode computation $L_n^{\mathrm{gh}}
= \sum_m [(\lambda-1)n - m]\,b_{n+m}\,c_{-m}$ and the resulting
commutator $[L_m, L_n] = (m-n)L_{m+n}
+ \frac{1}{12}m(m^2-1)\,c_{\mathrm{gh}}(s)\,\delta_{m+n}$.
See Friedan--Martinec--Shenker, \S2.
\end{proof}

\subsubsection{Explicit computation for each generator}

\emph{Spin-$2$ generator $T$:}
\begin{equation}
c_{\mathrm{gh}}(2) = -2(6 \cdot 4 - 6 \cdot 2 + 1) = -2(24 - 12 + 1) = -2 \cdot 13 = -26.
\end{equation}

\emph{Spin-$3$ generator $W$:}
\begin{equation}
c_{\mathrm{gh}}(3) = -2(6 \cdot 9 - 6 \cdot 3 + 1) = -2(54 - 18 + 1) = -2 \cdot 37 = -74.
\end{equation}

\emph{Total:}
\begin{equation}\label{eq:W3-ghost-total}
c_{\mathrm{gh}}^{\mathrm{total}}(\cW_3)
\;=\;
c_{\mathrm{gh}}(2) + c_{\mathrm{gh}}(3)
\;=\;
-26 + (-74) \;=\; -100.
\end{equation}

\subsection{The complementarity constant}

\begin{definition}[Complementarity constant]
For a W-algebra $\cA$ with generators of spins $s_1, \ldots, s_n$,
the complementarity constant is
\begin{equation}
\alpha \;=\; -c_{\mathrm{gh}}^{\mathrm{total}}
\;=\; \sum_{i=1}^n 2(6s_i^2 - 6s_i + 1).
\end{equation}
The Koszul involution is $c \mapsto \alpha - c$, and the Koszul
dual is $\cA_c^! = \cA_{\alpha - c}$.
\end{definition}

For $\cW_3$:
\begin{equation}\label{eq:alpha-W3}
\alpha_3 \;=\; -c_{\mathrm{gh}}^{\mathrm{total}}(\cW_3)
\;=\; 100.
\end{equation}

\begin{remark}[The general $\cW_N$ formula]
The closed-form evaluation gives
$\alpha_N = 2(N-1)(2N^2 + 2N + 1)$.  The first values:
$\alpha_2 = 26$ (Virasoro), $\alpha_3 = 100$, $\alpha_4 = 246$,
$\alpha_5 = 488$.  These match the slab central charges of the
Khan--Zeng construction, confirming the identity
$3(2s-1)^2 - 1 = 2(6s^2 - 6s + 1)$.
\end{remark}

\subsection{The Verdier dual pairing and BPZ inner product}

The Koszul dual is obtained by Verdier duality on the bar coalgebra:
$D_{\mathrm{Ran}}(\barB(\cA)) \simeq \barB(\cA^!)$.  The Verdier
dual pairing is realised concretely by the BPZ inner product on the
Verma modules of $\cW_3$.

\begin{warning}[AP26: BPZ $\ne$ Fock]
\label{warn:BPZ-not-Fock}
At weight $\ge 4$, the $\cW_3$ Verma module is a proper quotient
of the Fock module: $\dim V_n^{\mathrm{Fock}} > \dim V_n^{\cW_3}$
for $n \ge 4$.  The BPZ inner product is the Shapovalov form on
the $\cW_3$ Verma module, not the free-field inner product.  In
particular, $\langle W_4 | \Lambda \rangle_{\mathrm{Fock}} \ne 0$
but $\langle W_4 | \Lambda \rangle_{\mathrm{BPZ}} = 0$ in the
irreducible quotient.  All computations below use the BPZ form.
\end{warning}

\subsubsection{Weight-by-weight analysis of the Shapovalov form}

The Shapovalov form $G_n$ at weight $n$ is the matrix of inner
products $\langle v_i | v_j \rangle_{\mathrm{BPZ}}$ for a basis
$\{v_i\}$ of the weight-$n$ subspace of the Verma module
$M(c, h, w)$ (where $h$ is the $L_0$-eigenvalue and $w$ is the
$W_0$-eigenvalue of the highest-weight state).  The bar complex
$\barB(\cW_3)$ has degree-$k$ components parametrised by
$\mathrm{FM}_k(\C) \times \mathrm{Conf}_k(\R)$, and the Verdier
dual pairing acts weight-by-weight through the Shapovalov form.

\medskip
\emph{Weight $n = 0$:}  The vacuum $|c, 0, 0\rangle$; the
Shapovalov form is the scalar $1$.  The Verdier dual acts trivially.

\medskip
\emph{Weight $n = 1$:}  Basis: $\{L_{-1}|0\rangle\}$ (dimension $1$).
$G_1 = \langle 0 | L_1 L_{-1} | 0 \rangle = 2h$.  For the
vacuum Verma module ($h = 0$), this vanishes: the weight-$1$ space
is trivial on the irreducible vacuum module.

\medskip
\emph{Weight $n = 2$:}  Basis: $\{L_{-2}|0\rangle,\, L_{-1}^2|0\rangle\}$
(dimension~$2$ on the Verma module, but for $h = 0$ the vector
$L_{-1}|0\rangle = 0$, so the relevant space is $1$-dimensional,
spanned by $L_{-2}|0\rangle \sim T$).  The BPZ norm:
\begin{equation}
\langle T | T \rangle_{\mathrm{BPZ}}
= \langle 0 | L_2\, L_{-2} | 0 \rangle
= \tfrac{c}{2}.
\end{equation}
This is the Killing form on the spin-$2$ sector at weight $2$.

\medskip
\emph{Weight $n = 3$:}  Basis: $\{L_{-3}|0\rangle,\,
L_{-1}L_{-2}|0\rangle,\, W_{-3}|0\rangle\}$
(dimension~$3$ on the Verma; for $h = w = 0$, the $L_{-1}$
descendant vanishes, leaving $\{L_{-3}|0\rangle, W_{-3}|0\rangle\}$,
dimension~$2$).  Cross-term:
$\langle 0 | L_3\, W_{-3} | 0 \rangle
= \langle 0 | [L_3, W_{-3}] | 0 \rangle = 0$
(since $[L_3, W_{-3}] = 6W_0$ and $W_0|0\rangle = 0$ for the
vacuum).  Diagonal terms:
\begin{align}
\langle L_{-3}|0\rangle,\, L_{-3}|0\rangle \rangle
&= \langle 0 | L_3\, L_{-3} | 0 \rangle = 2c, \\
\langle W_{-3}|0\rangle,\, W_{-3}|0\rangle \rangle
&= \langle 0 | W_3\, W_{-3} | 0 \rangle = \tfrac{c}{3}.
\end{align}
The first follows from $[L_3, L_{-3}] = 6L_0 + (c/12)\cdot 24
= 2c$ (at $h = 0$).  The second follows from the $WW$ OPE:
the leading singular term is $W(z)W(w) \sim (c/3)(z-w)^{-6} + \cdots$,
so $[W_3, W_{-3}]$ contains the term $\frac{c}{3}\cdot\binom{5}{5}
= c/3$ from the $c$-dependent piece of the commutator (after
accounting for the mode algebra conventions).

\medskip
\emph{Weight $n = 4$:}  This is the first weight where the
composite field $\Lambda = {:}TT{:} - \tfrac{3}{10}\partial^2 T$
enters the $WW$ OPE.  The relevant states are
$\{L_{-4}|0\rangle, L_{-2}^2|0\rangle, W_{-4}|0\rangle,
L_{-1}L_{-3}|0\rangle, \ldots\}$.  For the vacuum module,
the independent states are $\{L_{-4}|0\rangle,
L_{-2}^2|0\rangle, W_{-4}|0\rangle\}$ (dimension $3$).

The Shapovalov matrix at weight $4$ on the vacuum module has
entries depending on $c$ through the Virasoro and $\cW_3$
commutation relations.  The Kac determinant
for $\cW_3$ at weight $4$ is
\begin{equation}\label{eq:W3-Kac-det-4}
\det G_4^{\cW_3}
\;\propto\;
c^3(5c + 22) \cdot (\text{terms involving } h, w),
\end{equation}
where the factor $(5c + 22)$ reflects the pole in $\beta^2 = 16/(22+5c)$
at $c = -22/5$.  The Verdier duality exchanges zeros and poles of
the Shapovalov determinant, which is how the dual central charge
enters.

\subsection{The Koszul involution: $c \mapsto 100 - c$}

\begin{proof}[Proof of Theorem~\ref{thm:koszul-dual-W3}]
The bar complex $\barB(\cW_3)$ is a dg coalgebra with differential
$d_{\barB}$ encoding the $\cW_3$ OPE collision residues on
$\mathrm{FM}_k(\C)$.  Verdier duality
$D_{\mathrm{Ran}}\colon \mathrm{Coalg} \to \mathrm{Alg}$
converts the bar coalgebra to an algebra.  On the Koszul locus
(which is where bar cohomology is concentrated), the resulting
algebra is the Koszul dual $\cA^!$.

The Verdier dual pairing is the BPZ inner product on $\cW_3$
weight spaces.  The BV-BRST quantization of the Khan--Zeng action
\[
S_{\cW_3} = \int_{\R \times \C}
\mu(\partial_t + \bar\partial)T + \chi(\partial_t + \bar\partial)W
+ (\text{interaction terms}),
\]
introduces ghost pairs $(\mu, T)$ and $(\chi, W)$ with ghost
central charges $c_{\mathrm{gh}}(2) = -26$ and
$c_{\mathrm{gh}}(3) = -74$.

The key mechanism: the Verdier dual of the bar differential
$d_{\barB}$ encodes the \emph{adjoint} OPE with respect to the
BPZ form.  The adjoint of a mode $L_n$ is $L_{-n}$; the adjoint of
$W_n$ is $(-1)^{2s}W_{-n} = W_{-n}$ (even parity for $s = 3$ in
BPZ convention).  The BPZ form on the vacuum module provides the
nondegenerate pairing that identifies $\barB(\cW_3)^{\vee}$ with
$\barB(\cW_3^!)$, and the ghost determinant shifts the central
charge by $\alpha_3 = -c_{\mathrm{gh}}^{\mathrm{total}} = 100$.

\medskip
\noindent\emph{Step 1: The $T$-sector reproduces the Virasoro shift.}
Restricting to the Virasoro subalgebra generated by $T$, the
bar complex $\barB(\mathrm{Vir}_c) \hookrightarrow \barB(\cW_3)$
has Verdier dual $\barB(\mathrm{Vir}_{26-c})$.  The Virasoro
ghost contributes $-26$ to the central charge shift.

\medskip
\noindent\emph{Step 2: The $W$-sector contributes the additional shift.}
The spin-$3$ generator $W$ produces a ghost pair $(\chi, W)$ of
conformal weights $(3, -2)$.  The ghost central charge
$c_{\mathrm{gh}}(3) = -74$ shifts the central charge by an
additional $74$.  The total shift is $26 + 74 = 100$.

\medskip
\noindent\emph{Step 3: The dual algebra is $\cW_3$.}
We must verify that the Koszul dual is again a $\cW_3$ algebra
(not some deformation or extension).  This follows from two facts:
\begin{enumerate}[label=(\alph*)]
\item \emph{Koszulness:}  The principal $\cW_3$ algebra is chirally
  Koszul (by the one-loop exactness criterion applied to the affine
  pre-image $V_k(\fsl_3)$, composed with principal DS reduction).
  Bar cohomology is concentrated, and the Koszul dual is a
  well-defined algebra with the same generator structure.
\item \emph{Generator preservation:}  DS reduction preserves the
  spin spectrum.  The Koszul dual of a principal $\cW_N$-algebra is
  again a principal $\cW_N$-algebra (at the dual level), because
  bar-cobar commutes with DS reduction (the DS-bar intertwining
  theorem of Volume~I).
\end{enumerate}
Therefore $\cW_{3,c}^! = \cW_{3,\,c'}$ for some $c'$, and the
ghost central charge computation gives $c' = 100 - c$.

\medskip
\noindent\emph{Step 4: Transformation of $\beta^2$.}
The structure constant $\beta^2 = 16/(22 + 5c)$ in the $WW$ OPE
is not an independent parameter: it is determined by $c$ through
the Jacobi identity of the $\cW_3$ algebra (Zamolodchikov's
uniqueness theorem).  Under $c \mapsto 100 - c$:
\begin{equation}\label{eq:beta-transform}
\beta'^{\,2}
\;=\; \frac{16}{22 + 5(100 - c)}
\;=\; \frac{16}{522 - 5c}.
\end{equation}
At the self-dual point $c = 50$:
\begin{equation}
\beta^2\big|_{c=50} = \frac{16}{22 + 250} = \frac{16}{272} = \frac{1}{17}.
\end{equation}
\end{proof}

\subsection{The modular characteristic and complementarity sum}

\begin{proposition}[Modular characteristic of $\cW_3$; proved here]
\label{prop:kappa-W3}
The modular characteristic (genus-$1$ curvature) of $\cW_{3,c}$ is
\begin{equation}\label{eq:kappa-W3}
\kappa(\cW_{3,c}) \;=\; \frac{5c}{6},
\end{equation}
decomposing as a sum of per-channel curvatures:
\begin{equation}\label{eq:kappa-per-channel}
\kappa_T = \frac{c}{2}, \qquad \kappa_W = \frac{c}{3}, \qquad
\kappa = \kappa_T + \kappa_W = \frac{c}{2} + \frac{c}{3} = \frac{5c}{6}.
\end{equation}
\end{proposition}

\begin{proof}
The genus universality theorem (Volume~I) gives
$\kappa(\cW_N) = c \cdot (H_N - 1)$, where $H_N = \sum_{j=1}^N 1/j$
is the $N$th harmonic number.  For $N = 3$:
$H_3 = 1 + 1/2 + 1/3 = 11/6$, so $H_3 - 1 = 5/6$ and
$\kappa = 5c/6$.

Alternatively, the per-channel curvature of a spin-$s$ generator
is $\kappa_s = c/(2s-2)$.  For $s = 2$:
$\kappa_T = c/2$.  The correct per-channel curvature for $W$
is extracted from the coefficient of the identity in the leading
singular OPE term: $\kappa_W = c/3$, and the total
$\kappa = c/2 + c/3 = 5c/6$ matches the harmonic number formula.
\end{proof}

\begin{theorem}[Complementarity sum for $\cW_3$; proved here]
\label{thm:complementarity-W3}
The complementarity sum is
\begin{equation}\label{eq:complementarity-W3}
\boxed{\;
K_{\cW_3}
\;:=\;
\kappa(\cW_{3,c}) + \kappa(\cW_{3,c}^!)
\;=\;
\frac{5c}{6} + \frac{5(100-c)}{6}
\;=\;
\frac{250}{3}.
\;}
\end{equation}
This is \textbf{not} zero: the Feigin--Frenkel anti-symmetry
$\kappa + \kappa^! = 0$ holds for Kac--Moody algebras but fails
for W-algebras.
\end{theorem}

\begin{proof}
Direct computation: $\kappa(\cW_{3,c}) + \kappa(\cW_{3,100-c})
= 5c/6 + 5(100-c)/6 = 500/6 = 250/3$.
Alternatively, $\kappa + \kappa^! = \varrho \cdot \alpha_N$ with
$\varrho = H_3 - 1 = 5/6$ and $\alpha_3 = 100$.
\end{proof}

\begin{remark}[The arithmetic of $100$]
The number $100 = \alpha_3 = 2 \cdot 2 \cdot (18+6+1)$ decomposes as
$26 + 74$: the Virasoro critical dimension plus the spin-$3$ ghost
contribution.  The self-dual point $c = 50$ is the ``natural zero''
of higher-spin gravity: as $c = 13$ is the zero of Virasoro Koszul
complementarity, $c = 50$ is the zero of $\cW_3$ Koszul
complementarity.
\end{remark}



%% ============================================================
%%  HOOK-TYPE ANOMALY COMPLETION FOR sl_4: THREE PATTERNS
%% ============================================================

\section{Hook-type anomaly completion for $\fsl_4$: three patterns}
\label{sec:sl4-hook-anomaly}

The anomaly-completed holography programme (proved for $\fsl_3$ in
the core and frontier chapters) extends to $\fsl_4$ with qualitatively
new phenomena.  The nilpotent poset of $\fsl_4$ contains five orbits,
indexed by partitions of~$4$:
\[
(1^4) \;\subset\; (2,1^2) \;\subset\; (2,2) \;\subset\; (3,1)
\;\subset\; (4).
\]
The two extreme orbits (zero and principal) are treated by the affine
and principal-$\cW_4$ theory respectively.  The three intermediate
orbits---$(3,1)$, $(2,2)$, and $(2,1,1)$---give three distinct
patterns for the anomaly completion.  The partition $(2,2)$ is
\emph{not} a hook: it is rectangular.  This is the first case where
a non-hook, non-principal nilpotent orbit enters the anomaly-completed
programme.

\medskip
\noindent\textbf{Global data.}
$\fg = \fsl_4$, $h^\vee = 4$, $\dim(\fsl_4) = 15$.
The modular characteristic of the affine algebra:
\begin{equation}\label{eq:kappa-sl4-km-wn}
\kappa(V_k(\fsl_4))
\;=\;
\frac{15(k+4)}{8}.
\end{equation}
Koszul dual level for all orbits: $k^\vee = -k - 2h^\vee = -k - 8$.

\subsection{Pattern I: the subregular $(3,1)$}
\label{subsec:sl4-hook-31-wn}

\begin{observation}[Representation-theoretic data for $(3,1)$]
\label{obs:sl4-31-data-wn}
The partition $\lambda = (3,1)$ corresponds to the subregular nilpotent
orbit of $\fsl_4$.  The Jacobson--Morozov triple has
$h = \mathrm{diag}(2,0,-2,0)$, decomposing $\C^4 = V_2 \oplus V_0$
(spin-$1$ plus trivial irrep of $\fsl_2$).  The $\ad(h)$-grading of
$\fsl_4$ is:
\[
\fsl_4 = \fg_{-4} \oplus \fg_{-2} \oplus \fg_0 \oplus \fg_2 \oplus \fg_4,
\qquad
\dim = 1 + 4 + 5 + 4 + 1.
\]

The centraliser $\fg^{f_{(3,1)}}$ has dimension
$\dim(\fg^f) = N^2 - 1 - \dim(\cO_{(3,1)}) = 15 - 10 = 5$,
and sits in JM grades $0, -2, -4$ with multiplicities $1, 3, 1$.
The $\cW$-algebra $\cW^k(\fsl_4, f_{(3,1)})$ therefore has
\textbf{five strong generators}:
\begin{center}
\renewcommand{\arraystretch}{1.2}
\begin{tabular}{ccc}
\textbf{JM grade of $\fg^f$} & \textbf{Conformal weight}
  & \textbf{Multiplicity} \\
\hline
$0$ & $1$ & $1$ \\
$-2$ & $2$ & $3$ \\
$-4$ & $3$ & $1$
\end{tabular}
\end{center}
The grade-$0$ generator is a $\mathfrak{u}(1)$ current.  The three
weight-$2$ generators include the stress tensor~$T$ and two
additional spin-$2$ fields.  The weight-$3$ generator is analogous
to the $W_3$ current of the Zamolodchikov algebra.

Ghost dimension for DS reduction from $\widehat{\fsl}_{4,k}$:
$r(N{-}r) = 1 \cdot 3 = 3$ ghost pairs.
\end{observation}

\begin{observation}[Cross-orbit Koszul duality: $(3,1) \leftrightarrow (2,1,1)$]
\label{obs:sl4-31-koszul-wn}
The transpose partition is $(3,1)^t = (2,1,1)$.  Since
$(3,1) \neq (2,1,1)$, the Koszul duality is \emph{genuinely
cross-orbit}:
\[
\bigl(\cW^k(\fsl_4, f_{(3,1)})\bigr)^!
\;\simeq\;
\cW^{-k-8}(\fsl_4, f_{(2,1,1)}).
\]
The subregular $\cW$-algebra (5~generators of weights $1,2,2,2,3$) is
Koszul dual to the minimal $\cW$-algebra (9~generators of weights
$1^4, (3/2)^4, 2$).  The two algebras have different rank, different
generator spectra, and different ghost systems ($3$ vs $4$ ghost pairs).
This is a genuinely asymmetric duality: neither the generator count
nor the conformal-weight spectrum is preserved.
\end{observation}

\begin{observation}[Curvature and anomaly class for $(3,1)$]
\label{obs:sl4-31-kappa-wn}
The modular characteristic is obtained from the DS transform:
\[
\kappa\bigl(\cW^k(\fsl_4, f_{(3,1)})\bigr)
\;=\;
\frac{15(k+4)}{8}
\;-\;
\kappa_{\mathrm{ghost}}^{(3,1)}(k).
\]
The anomaly is proportional to $(k+4)$ \emph{at the affine level},
confirming the $\fsl_3$ pattern: for $\fsl_3$, $\kappa \propto (k+3)$;
for $\fsl_4$, $\kappa \propto (k+h^\vee) = (k+4)$.
For non-principal orbits, the ghost correction introduces a rational
deviation from strict proportionality, but the
\emph{zero of $\kappa$} is always governed by the shifted level
$(k+4)$ at leading order.
\end{observation}

\subsection{Pattern II: the rectangular $(2,2)$}
\label{subsec:sl4-hook-22-wn}

\begin{observation}[Representation-theoretic data for $(2,2)$]
\label{obs:sl4-22-data-wn}
The partition $\lambda = (2,2)$ is rectangular and
\emph{not a hook partition}.  The Jacobson--Morozov element is
$h = \mathrm{diag}(1,-1,1,-1)$, decomposing $\C^4 = V_1 \oplus V_1$.
The $\ad(h)$-grading is concentrated in even degrees:
\[
\fsl_4 = \fg_{-2} \oplus \fg_0 \oplus \fg_2,
\qquad
\dim = 4 + 7 + 4.
\]

The centraliser has dimension $\dim(\fg^{f_{(2,2)}}) = 15 - 8 = 7$
and sits entirely in JM grades $0$ and $-2$:
$\cW^k(\fsl_4, f_{(2,2)})$ has
\textbf{seven strong generators}: three at weight~$1$ (forming an
$\widehat{\fsl}_2$ subalgebra) and four at weight~$2$.
\end{observation}

\begin{observation}[Self-duality of the rectangular orbit]
\label{obs:sl4-22-koszul-wn}
The transpose partition is $(2,2)^t = (2,2)$.
The Koszul duality is therefore \emph{same-orbit}:
\[
\bigl(\cW^k(\fsl_4, f_{(2,2)})\bigr)^!
\;\simeq\;
\cW^{-k-8}(\fsl_4, f_{(2,2)}).
\]
This is the $\fsl_4$ analogue of BP self-duality,
but $(2,2)$ is the \emph{middle} orbit sitting
between the genuinely cross-orbit pair $(3,1) \leftrightarrow (2,1,1)$.
The self-dual point is $k = -4$ (the critical level).
\end{observation}

\subsection{Pattern III: the minimal $(2,1,1)$}
\label{subsec:sl4-hook-211-wn}

\begin{observation}[Representation-theoretic data for $(2,1,1)$]
\label{obs:sl4-211-data-wn}
The partition $\lambda = (2,1,1)$ is the minimal nilpotent orbit.
The Jacobson--Morozov element is
$h = \mathrm{diag}(1,-1,0,0)$.  \emph{Half-integer grades appear}
in the Dynkin good grading.
The centraliser has dimension $\dim(\fg^{f_{(2,1,1)}}) = 9$,
and $\cW^k(\fsl_4, f_{(2,1,1)})$ has
\textbf{nine strong generators}: $4 \times \Delta{=}1$,
$4 \times \Delta{=}3/2$, $1 \times \Delta{=}2$.
Ghost dimension: $r(N{-}r) = 2 \cdot 2 = 4$ pairs.
\end{observation}

\begin{observation}[Cross-orbit Koszul duality: $(2,1,1) \leftrightarrow (3,1)$]
\label{obs:sl4-211-koszul-wn}
The transpose is $(2,1,1)^t = (3,1)$:
\[
\bigl(\cW^k(\fsl_4, f_{(2,1,1)})\bigr)^!
\;\simeq\;
\cW^{-k-8}(\fsl_4, f_{(3,1)}).
\]
The cross-orbit duality creates a sharp asymmetry:
the $(2,1,1)$-side has more generators (9 vs 5),
the ghost systems differ ($4$ vs $3$ pairs),
and the transgression algebra $B_\Theta$ bridges the two.
\end{observation}

\subsection{The $\fsl_4$ nilpotent chain}
\label{subsec:sl4-chain-wn}

\begin{observation}[The $\fsl_4$ chain]
\label{obs:sl4-chain-wn}
The DS chain:
$\widehat{\fsl}_{4,k}
\xrightarrow{\mathrm{DS}_{\min}}
\cW^k_{(2,1,1)}
\xrightarrow{\mathrm{DS}}
\cW^k_{(2,2)}
\xrightarrow{\mathrm{DS}}
\cW^k_{(3,1)}
\xrightarrow{\mathrm{DS}_{\mathrm{prin}}}
\cW_4^k$.
Three duality patterns: (i)~same-family at extremes ($\kappa+\kappa^!=0$),
(ii)~same-orbit non-principal for $(2,2)$ ($\kappa+\kappa^! = \kappa_0^{(2,2)} \neq 0$, by AP24),
(iii)~cross-orbit for $(3,1)\leftrightarrow(2,1,1)$.
\end{observation}

\subsection{The $(k+4)$ law and JM parity criterion}
\label{subsec:sl4-kN-law-wn}

\begin{observation}[The $(k+h^\vee)$ law]
\label{obs:sl4-kN-law-wn}
The modular characteristic of the affine algebra is
$\kappa(V_k(\fsl_N)) = \frac{N^2-1}{2N}(k+N)$, proportional to $(k+h^\vee)$.
For non-principal orbits, the ghost correction controls the deviation:
orbits with \emph{even-integer} JM grading ($(3,1)$, $(2,2)$) have
level-independent $\kappa_{\mathrm{ghost}}$, so $\kappa \propto (k+4)$
exactly.  Orbits with \emph{odd-integer} JM grading ($(2,1,1)$) have
level-dependent $\kappa_{\mathrm{ghost}}(k)$, breaking strict
proportionality.  This is the \textbf{JM parity criterion} for
level-dependence of the ghost modular characteristic.
\end{observation}

\begin{observation}[Comparison with $\fsl_3$]
\label{obs:sl4-vs-sl3-wn}
\begin{center}
\renewcommand{\arraystretch}{1.2}
\begin{tabular}{lcc}
& $\fsl_3$ & $\fsl_4$ \\
\hline
$\kappa \propto (k+h^\vee)$ & $(k+3)$ & $(k+4)$ \\[2pt]
Cross-orbit pairs & none & $(3,1) \leftrightarrow (2,1,1)$ \\[2pt]
Self-dual non-principal & $(2,1)$ & $(2,2)$ \\[2pt]
Level-dependent ghost $\kappa$ & $(2,1)$ & $(2,1,1)$ \\[2pt]
Level-independent ghost $\kappa$ & principal & $(3,1), (2,2)$, principal \\[2pt]
\end{tabular}
\end{center}
\end{observation}




%====================================================================
%
%  DS-BAR INTERTWINING: CHAIN-LEVEL VERIFICATION FOR sl(2) -> Virasoro
%
%====================================================================

\section{DS-bar intertwining: chain-level verification for
$\fsl(2) \to$ Virasoro}
\label{sec:ds-bar-chain-level}

The DS--bar intertwining theorem (Vol~I,
Theorem~ds-koszul-intertwine) asserts that for the principal
nilpotent $\chi \in \fg$ and the DS reduction functor
$\mathrm{DS}_\chi$, the bar construction commutes:
\[
\barB(\mathrm{DS}_\chi(\cA))
\;\simeq\;
\mathrm{DS}_\chi(\barB(\cA))
\qquad
\text{as } \SCchtop\text{-coalgebras.}
\]
This section verifies the theorem at chain level for the
foundational case $\cA = V_k(\fsl_2)$, with
$\mathrm{DS}_\chi(V_k(\fsl_2)) = \mathrm{Vir}_{c_{\mathrm{DS}}}$,
$c_{\mathrm{DS}} = 1 - 6(k+1)^2/(k+2)$.  We compute both sides of
the intertwining at arities~$2$ and~$3$ and verify their
equality.  At arity~$3$, the key result is the emergence of
$m_3^{\mathrm{Vir}} \ne 0$ from the $m_3^{\fsl_2} = 0$ input:
the BRST homotopy transfer \emph{manufactures} the Virasoro
ternary operation from the strictly associative affine bracket.

\subsection{Setup and conventions}

\subsubsection*{The affine algebra $V_k(\fsl_2)$}

Basis $\{e, h, f\}$ for $\fsl_2$ with standard structure
constants $[e,f] = h$, $[h,e] = 2e$, $[h,f] = -2f$.
Killing form normalised by $1/(2h^\vee)$: $\kappa(e,f)
= \kappa(f,e) = 1$, $\kappa(h,h) = 2$.  The affine
$\lambda$-brackets are:
\begin{equation}\label{eq:wn-ds-affine-brackets}
\begin{aligned}
\{J^e{}_\lambda J^f\} &= J^h + k\lambda, &
\{J^f{}_\lambda J^e\} &= -J^h + k\lambda, \\
\{J^h{}_\lambda J^e\} &= 2J^e, &
\{J^h{}_\lambda J^f\} &= -2J^f, \\
\{J^h{}_\lambda J^h\} &= 2k\lambda, &
\{J^e{}_\lambda J^e\} = \{J^f{}_\lambda J^f\} &= 0.
\end{aligned}
\end{equation}

The $\Ainf$ structure: maximal OPE pole order $d = 2$
(from the $k\lambda$ central terms), so the $\lambda$-bracket
is strictly associative: $A_3 = 0$, hence
$m_3^{\fsl_2} = 0$ and $m_k^{\fsl_2} = 0$ for all
$k \ge 3$ (class~$\mathbf{L}$, shadow depth $r_{\max} = 3$).

\subsubsection*{DS reduction data}

The principal DS reduction imposes $J^e = 1$
(the principal character of $\fn = \C\!\cdot\!e$),
introduces a ghost pair $(c, b)$ of conformal weights
$(0, 1)$, and takes BRST cohomology.  The DS stress
tensor is:
\begin{equation}\label{eq:wn-ds-stress-tensor}
T^{\mathrm{DS}}
\;=\;
J^f + \frac{(J^h)^2}{4(k+2)} + \frac{1}{2}\,\partial J^h.
\end{equation}
The three terms: $J^f$ is the lowest-root current
(weight~$2$ after the DS constraint); $(J^h)^2/(4(k+2))$ is
the Cartan Sugawara contribution; $(1/2)\partial J^h$ is the
improvement term.

The central charge:
\begin{equation}\label{eq:wn-ds-cc}
c_{\mathrm{DS}}(k)
\;=\;
1 - \frac{6(k+1)^2}{k+2}.
\end{equation}
The $\lambda$-bracket of $T^{\mathrm{DS}}$, computed from
the affine brackets~\eqref{eq:wn-ds-affine-brackets} via the
non-commutative Wick formula, is the Virasoro
$\lambda$-bracket:
\begin{equation}\label{eq:wn-ds-vir-bracket}
\{T^{\mathrm{DS}}{}_\lambda T^{\mathrm{DS}}\}
\;=\;
\partial T^{\mathrm{DS}} + 2\,T^{\mathrm{DS}}\,\lambda
+ \frac{c_{\mathrm{DS}}}{12}\,\lambda^3.
\end{equation}
This is the standard Feigin--Frenkel result.  The quartic
pole $c_{\mathrm{DS}}/12$ is \emph{manufactured} by the
Sugawara double-Wick contraction of $(J^h)^2$ with itself;
it does not exist in the original affine OPE (which has at
most double poles).

\subsection{The bar complex at arity 2}

\subsubsection*{Right-hand side:
$\barB(\mathrm{Vir}_{c_{\mathrm{DS}}})$ at arity~$2$}

The bar complex $\barB(\mathrm{Vir}_c)$ at arity~$2$
consists of elements $[s^{-1}T \mid s^{-1}T]$ in the
tensor coalgebra $T^c(s^{-1}\mathrm{Vir}_c)$.  The bar
differential at arity~$2$ encodes $m_2$:
\begin{equation}\label{eq:wn-vir-bar-2}
d_{\barB}\bigl[s^{-1}T \mid s^{-1}T\bigr]
\;=\;
s^{-1}\bigl(m_2(T, T)\bigr)
\;=\;
s^{-1}\bigl(\partial T\bigr).
\end{equation}

\subsubsection*{Left-hand side:
$\mathrm{DS}_\chi(\barB(V_k(\fsl_2)))$ at arity~$2$}

The bar complex $\barB(V_k(\fsl_2))$ at arity~$2$ consists
of elements $[s^{-1}J^a \mid s^{-1}J^b]$ for
$a, b \in \{e, h, f\}$.  The bar differential encodes the
affine $\lambda$-bracket:
\begin{equation}\label{eq:wn-aff-bar-2}
d_{\barB}\bigl[s^{-1}J^a \mid s^{-1}J^b\bigr]
\;=\;
s^{-1}\bigl(\{J^a{}_\lambda J^b\}\bigr)
\;=\;
s^{-1}\bigl(f^{ab}_c\,J^c + k\,\kappa(a,b)\,\lambda\bigr).
\end{equation}

At arity~$2$, the surviving bar element
$[s^{-1}T^{\mathrm{DS}} \mid s^{-1}T^{\mathrm{DS}}]$
in $\mathrm{DS}_\chi(\barB(V_k(\fsl_2)))$ has bar
differential:
\begin{equation}\label{eq:wn-ds-bar-2}
d_{\barB}\bigl[s^{-1}T^{\mathrm{DS}}
\mid s^{-1}T^{\mathrm{DS}}\bigr]
\;=\;
s^{-1}\Bigl(\partial T^{\mathrm{DS}}
+ 2\,T^{\mathrm{DS}}\lambda
+ \frac{c_{\mathrm{DS}}}{12}\,\lambda^3\Bigr).
\end{equation}

\begin{proposition}[Arity-$2$ intertwining; chain-level]%
\label{prop:wn-ds-bar-arity-2}
The arity-$2$ bar complexes agree:
$\mathrm{DS}_\chi(\barB(V_k(\fsl_2)))|_{\text{arity }2}
= \barB(\mathrm{Vir}_{c_{\mathrm{DS}}})|_{\text{arity }2}$.
\end{proposition}

\begin{proof}
Both sides are determined by the $\lambda$-bracket of $T^{\mathrm{DS}}$,
which the Feigin--Frenkel computation~\eqref{eq:wn-ds-vir-bracket}
identifies with the Virasoro $\lambda$-bracket at $c = c_{\mathrm{DS}}$.
\end{proof}

\begin{remark}[The quartic pole is manufactured]
\label{rem:wn-quartic-manufactured}
The affine bar complex at arity~$2$ records
$\lambda$-brackets with at most $\lambda^1$ terms (double
poles in the OPE).  After DS reduction, the composite field
$T^{\mathrm{DS}}$ has a $\lambda^3$ term (quartic pole)
in its self-bracket.  This \emph{emergence of the quartic
pole} is the arity-$2$ manifestation of complexity increase
$\mathbf{L} \to \mathbf{M}$: the Sugawara
double-Wick contraction of $J^h$ with itself, mediated by
the $1/(4(k+2))$ normalisation, generates a pole order
that was absent in the original algebra.
\end{remark}

\subsection{The bar complex at arity 3: the key verification}

Since $m_1 = 0$, the Stasheff relation at arity~$3$ reduces to
$A_3 + m_3 = 0$ where $A_3$ is the associator of $m_2$.

The Virasoro $m_3$ at $c = c_{\mathrm{DS}}$ is:
\begin{equation}\label{eq:wn-vir-m3-ds}
\boxed{
m_3^{\mathrm{Vir}}(T,T,T;\,\lambda_{12},\lambda_{23})
\;=\;
\partial^2 T + (2\lambda_{12} + 3\lambda_{23})\,\partial T
+ 2\lambda_{23}(2\lambda_{12} + \lambda_{23})\,T
+ \frac{c_{\mathrm{DS}}}{12}\,
\lambda_{23}^3(2\lambda_{12} + \lambda_{23}).}
\end{equation}

The affine bar complex at arity~$3$ has \emph{only} binary
splitting terms ($m_3^{\fsl_2} = 0$).

\begin{theorem}[Arity-$3$ intertwining; chain-level;
the key verification]%
\label{thm:wn-ds-bar-arity-3}
The arity-$3$ bar complexes agree:
$\mathrm{DS}_\chi(\barB(V_k(\fsl_2)))|_{\text{arity }3}
= \barB(\mathrm{Vir}_{c_{\mathrm{DS}}})|_{\text{arity }3}$.
In particular, the DS reduction of the \emph{purely
binary} affine bar differential at arity~$3$ produces the
Virasoro ternary operation $m_3^{\mathrm{Vir}} \ne 0$
at $c = c_{\mathrm{DS}}$.
\end{theorem}

\begin{proof}
\textbf{Step 1: BRST homotopy transfer.}
The DS reduction gives a homotopy retract
$(C^{\mathrm{BRST}}, Q_{\mathrm{BRST}})
\rightleftarrows (\mathrm{Vir}_{c_{\mathrm{DS}}}, 0)$
with contracting homotopy $h$.  The HPL transfers the
$\Ainf$ structure to the cohomology via
$m_k^{\mathrm{Vir}} = p \circ m_2^{\fsl_2}
\circ (\id + h \circ m_2^{\fsl_2})^{\circ(k-2)}
\circ i^{\otimes k}$.

\medskip
\noindent\textbf{Step 2: Uniqueness.}
Both the HPL-transferred $m_3$ and the intrinsic $m_3 = -A_3$
satisfy the same Stasheff identity with the same $m_2$.
Since $\HH^{2,3}(\mathrm{Vir}_c, \mathrm{Vir}_c) = 0$ for
generic~$c$, the solution is unique:
$m_3^{\mathrm{HPL}} = m_3^{\mathrm{Vir}}$.

\medskip
\noindent\textbf{Step 3: Intertwining.}
Since both sides have the same $\Ainf$ structure and the bar
construction is functorial (Vol~I, Theorem~A), the bar complexes agree.
\end{proof}

\begin{observation}[The mechanism: Sugawara manufactures $m_3$]%
\label{obs:wn-sugawara-manufactures-m3}
The chain-level computation reveals the precise mechanism:
(i)~the affine algebra has $m_3^{\fsl_2} = 0$ (strictly associative);
(ii)~the DS composite $T^{\mathrm{DS}}$ introduces a quadratic field
$(J^h)^2/(4(k+2))$, whose double-Wick contraction generates the
quartic pole;
(iii)~the non-associativity $A_3 \ne 0$ forces $m_3 = -A_3 \ne 0$
by Stasheff, seeding the entire infinite tower
$m_4, m_5, \ldots$ (class~$\mathbf{M}$);
(iv)~the scalar coefficient $c_{\mathrm{DS}}/12$ controls the
three-graviton coupling.
\end{observation}

\subsection{Extension to all arities by induction}

\begin{proposition}[DS-bar intertwining at all arities]%
\label{prop:wn-ds-bar-all-arities}
For all $k \ge 2$:
$\mathrm{DS}_\chi(\barB(V_k(\fsl_2)))|_{\text{arity }k}
= \barB(\mathrm{Vir}_{c_{\mathrm{DS}}})|_{\text{arity }k}$.
\end{proposition}

\begin{proof}
By induction on $k$.  The inductive step uses
$\HH^{2,k}(\mathrm{Vir}_c, \mathrm{Vir}_c) = 0$ for
generic~$c$ to propagate uniqueness from arities $\le k-1$
to arity~$k$.
\end{proof}

\subsection{Summary and significance}

The chain-level verification establishes: (1)~at arity~$2$,
the Sugawara double-Wick contraction manufactures the quartic pole;
(2)~at arity~$3$, the Stasheff identity forces the unique
$m_3^{\mathrm{DS}} = m_3^{\mathrm{Vir}} = -A_3$, the chain-level
manifestation of class transport $\mathbf{L} \to \mathbf{M}$;
(3)~at all arities, Hochschild vanishing propagates the uniqueness.
The entire infinite $\Ainf$ tower of 3d quantum gravity is
\emph{manufactured} by the DS functor from a single input: the
affine $\fsl_2$ $\lambda$-bracket with its double pole.
The quartic pole of the Virasoro OPE is not put in by hand;
it is the Wick-contraction consequence of Sugawara, itself
forced by the DS constraint $J^e = 1$.




% =================================================================
\section{Modular periodicity of minimal model shadow obstruction towers}
\label{sec:modular-periodicity-minimal-shadow}
% =================================================================

Conjecture~\ref*{conj:modular-periodicity-minimal} (Vol~I,
\S\ref*{sec:one-quadratic-form}) predicts eventual periodicity
of bar cohomology dimensions for Virasoro minimal models
$M(p,q)$ with period $N = 24q'/\gcd(p',24)$, where
$c = p'/q'$ in lowest terms.  The following notes report a
systematic computational test of this conjecture.

\subsection*{Setup}

For $M(p,q)$ with $2 \le p < q$, $\gcd(p,q)=1$:
\[
c = 1 - \frac{6(p-q)^2}{pq}, \qquad
h_{r,s} = \frac{(qr-ps)^2 - (q-p)^2}{4pq},
\]
with $(p-1)(q-1)/2$ primaries under the identification
$V_{r,s} = V_{p-r,q-s}$.  The $T$-matrix eigenvalue at
$(r,s)$ is $\lambda_{r,s} = e^{2\pi i (h_{r,s} - c/24)}$,
where
\[
h_{r,s} - \frac{c}{24}
= \frac{6(qr-ps)^2 - pq}{24pq}.
\]

\subsection*{Finding 1: The conjecture's formula underestimates
  the $T$-matrix order}

The conjectured period $N = 24q'/\gcd(p',24)$ ensures
$N \cdot c/24 \in \mathbb{Z}$, i.e., the \emph{vacuum}
$T$-eigenvalue has period~$N$.  However, $N$ need not
annihilate the $T$-eigenvalues of \emph{all} primaries.

The full $T$-matrix order is
\[
T(p,q)
\;:=\; \operatorname{ord}(T)
\;=\; \operatorname{lcm}_{(r,s)}
  \bigl(\operatorname{denom}\bigl(h_{r,s} - c/24\bigr)\bigr)
\;=\; \frac{24pq}{\gcd\bigl(\{6a^2 - pq : a \in \mathcal{K}(p,q)\},\, 24pq\bigr)},
\]
where $\mathcal{K}(p,q) = \{qr - ps : 1 \le r < p,\,
1 \le s < q,\, qr + ps \le pq\}$ is the set of Kac-table
offsets.

\begin{computation}[$T$-matrix order vs.\ conjectured period]
\label{comp:T-order-vs-N}
By systematic computation for all $M(p,q)$ with $p + q \le 22$:
\begin{enumerate}[label=\textup{(\roman*)}]
\item When $p \ge 4$: $T(p,q) = 24pq/\gcd(6, pq)$ always,
  and $N = T(p,q)$ except when $\gcd(p',24) > \gcd(6,pq)$
  (which occurs for $M(5,6)$, $M(6,7)$, $M(7,10)$,
  $M(9,10)$, $M(10,11)$, and others with
  $6 \mid pq$ or both $p,q$ large).
\item When $p = 2$ or $p = 3$: the Kac table is too sparse
  to hit all residue classes mod~$6$, so $T(p,q)$ can be
  strictly less than $24pq/\gcd(6,pq)$.  In every case
  tested, $N = T(p,q)$ for these models.
\item The conjectured $N$ always divides $T(p,q)$, but
  $N < T(p,q)$ occurs for at least $7$ models with
  $p + q \le 22$:
\end{enumerate}
\begin{center}
\small
\renewcommand{\arraystretch}{1.15}
\begin{tabular}{cccccc}
$M(p,q)$ & $c$ & $N$ (conjecture) & $T(p,q)$ (correct) &
  $T/N$ & offending primary \\
\hline
$M(3,10)$ & $-44/5$ & $30$ & $120$ & $4$ & $V_{2,3},\; h = 3/40$ \\
$M(5,6)$ & $4/5$ & $30$ & $120$ & $4$ & $V_{1,2},\; h = 1/8$ \\
$M(6,7)$ & $6/7$ & $28$ & $168$ & $6$ & $V_{2,1},\; h = 3/8$ \\
$M(6,11)$ & $-14/11$ & $132$ & $264$ & $2$ & $V_{2,3}$ \\
$M(7,10)$ & $8/35$ & $105$ & $840$ & $8$ & $V_{2,2}$ \\
$M(9,10)$ & $14/15$ & $180$ & $360$ & $2$ & $V_{4,1}$ \\
$M(10,11)$ & $52/55$ & $330$ & $1320$ & $4$ & $V_{4,1}$
\end{tabular}
\end{center}
The ratio $T/N$ is always a power of $2$ or $3$ (or $6$),
because the mismatch arises from the factors of $2$ and $3$
in $24 = 2^3 \cdot 3$.
\end{computation}

\begin{remark}[Stable formula for $p \ge 4$]
\label{rem:stable-formula}
For $p \ge 4$, the Kac table offsets $\mathcal{K}(p,q)$
hit all residue classes modulo~$6$, forcing
$\gcd(\{6a^2 - pq\}, 24pq) = \gcd(6,pq)$ and hence
\[
T(p,q) = \frac{24pq}{\gcd(6, pq)}
= \begin{cases}
4pq & \text{if } 6 \mid pq, \\
8pq & \text{if } 3 \mid pq,\; 2 \nmid pq, \\
12pq & \text{if } 2 \mid pq,\; 3 \nmid pq, \\
24pq & \text{if } \gcd(pq,6) = 1.
\end{cases}
\]
This is strictly a function of $pq$ modulo~$6$; for $p \ge 4$,
the specific values of $p, q$ affect~$T$ only through their
product and GCD with~$6$.
\end{remark}

\begin{remark}[Why the original formula fails]
The formula $N = 24q'/\gcd(p',24)$ ensures only
$N \cdot c/24 \in \mathbb{Z}$, which gives $T^N = 1$ on the
\emph{vacuum character}.  The primary at $(r,s)$ contributes
the phase $(h_{r,s} - c/24) = (6a^2 - pq)/(24pq)$ with
$a = qr - ps$, and this may have a larger denominator.  The
failure occurs precisely when $h_{r,s}$ has a denominator
not dividing $24q'/\gcd(p',24)$---which happens when $4pq$
introduces factors of~$2$ or~$3$ beyond those in~$24q'$.
\end{remark}

\subsection*{Finding 2: Weight-space dimensions are not
  eventually periodic}

The weight-space dimensions $\dim V_{r,s,[k]}$ are given by
the Rocha--Caridi formula as a convolution of level-$2pq$
theta-function coefficients with the partition function~$p(\cdot)$.
As observed in the evidence for
Conjecture~\ref*{conj:modular-periodicity-minimal} (Step~2
gap), the partition function grows monotonically, destroying
any periodicity.

\begin{computation}[Weight-space monotonicity]
\label{comp:weight-space-monotone}
For $M(2,5)$, $M(3,4)$, and $M(3,5)$, we computed
$\dim V_{r,s,[k]}$ for all primaries $(r,s)$ through
$k = 150$.  In every case:
\begin{itemize}[leftmargin=1.8em]
\item $\dim V_{r,s,[k]}$ is eventually strictly increasing
  (after $k \ge 2$ for most primaries).
\item No period $T$ in $\{1, 2, \ldots, 60\}$ achieves
  $\dim V_{r,s,[k+T]} = \dim V_{r,s,[k]}$ for any $k$ in
  the tested range.
\item The first differences $\dim V_{r,s,[k+T]} - \dim V_{r,s,[k]}$
  are themselves monotonically increasing, ruling out eventual
  periodicity at any period.
\end{itemize}
This confirms that the Step~2 gap in the conjecture's evidence
is not merely a technical obstruction but reflects a genuine
failure of weight-space periodicity.
\end{computation}

\subsection*{Finding 3: The bar Euler characteristic is not
  eventually periodic}

For the vacuum module~$A$ of a Virasoro minimal model, the bar
Euler characteristic satisfies
\[
\sum_{n \ge 0} (-1)^n \dim B^{(n)}_{[h]}
= \bigl[q^h\bigr]\, \frac{1}{\operatorname{ch}_A(q)},
\]
where $\operatorname{ch}_A(q) = \sum_{h \ge 0} \dim A_{[h]}\, q^h$.
Since $\dim H^n(\bar{B}_{[h]})$ contributes to this alternating
sum, eventual periodicity of all bar cohomology groups at a
common period~$N$ would force the inverse character coefficients
to be eventually $N$-periodic.

\begin{computation}[Inverse character coefficients]
\label{comp:inverse-character}
We computed $[q^h](1/\operatorname{ch}_A)$ for $h \le 120$ for
$M(2,5)$, $M(3,4)$, and $M(3,5)$:
\begin{enumerate}[label=\textup{(\roman*)}]
\item $M(2,5)$ ($c = -22/5$, conjectured $N = 60$):
\[
\bigl[q^h\bigr]\frac{1}{\operatorname{ch}_A}
= 1, 0, -1, -1, 0, 1, 0, -1, -1, 1, 2, 1, -2, -2, 1,
  3, 1, -3, -3, 2, 5, 1, \ldots
\]
These oscillate with growing amplitude and are not eventually
periodic at any period $T \le 120$.
\item $M(3,4)$ ($c = 1/2$, conjectured $N = 48$):
\[
\bigl[q^h\bigr]\frac{1}{\operatorname{ch}_A}
= 1, 0, -1, -1, -1, 0, 1, 2, 2, 0, -2, -3, -3, 0, 3, 5, \ldots
\]
Same qualitative behavior: oscillation with growing amplitude.
\item $M(3,5)$ ($c = -3/5$, conjectured $N = 40$): analogous.
\end{enumerate}
In all three cases, the absolute values
$|[q^h](1/\operatorname{ch}_A)|$ grow without bound, which is
incompatible with eventual periodicity.
\end{computation}

\begin{proposition}[Failure of exact bar-cohomology periodicity]
\label{prop:bar-periodicity-fails}
For any Virasoro minimal model $M(p,q)$ with
$(p-1)(q-1)/2 > 1$ (i.e., excluding the trivial $M(2,3)$),
the bar Euler characteristic
$\chi^{\operatorname{bar}}_h := \sum_n (-1)^n \dim H^n(\bar{B}_{[h]})$
is not eventually periodic at any positive period.
Therefore the bar cohomology dimensions
$\dim H^n(\bar{B}_{[h]})$ cannot all be simultaneously
eventually periodic at a common period.
\end{proposition}

\begin{proof}[Proof sketch]
The inverse character $1/\operatorname{ch}_A(q)$ is a power series
whose coefficients satisfy
$\chi^{\operatorname{bar}}_h = [q^h](1/\operatorname{ch}_A)$ for $h > 0$.
Since $\operatorname{ch}_A(q) = \Theta_{q-p,\,2pq}(q)/\eta(q)$
(Rocha--Caridi), we have
$1/\operatorname{ch}_A = \eta(q)/\Theta_{q-p,\,2pq}(q)$.
The eta function $\eta(q) = q^{1/24}\prod(1-q^n)$ contributes
lacunary coefficients (Euler's pentagonal theorem), while the
theta-function denominator $\Theta_{q-p,\,2pq}$ is supported
on a sparse set of exponents.  The quotient $\eta/\Theta$ has
coefficients whose absolute values grow as
$O(h^{(p-1)(q-1)/4 - 1})$ (by Rademacher-type estimates for
vector-valued modular forms of weight~$1/2$), which is
incompatible with bounded periodic behavior.
\end{proof}

\begin{remark}[Chiral Koszulness vs.\ classical purity]
\label{rem:chiral-koszul-no-purity}
The Virasoro VOA has descendants at all weights $h \ge 2$
(not just multiples of the generator weight), since
$L_{-3}\lvert 0\rangle \ne 0$ at weight~$3$.  Therefore the
classical Koszul purity condition
$H^n(\bar{B})_{[h]} = 0$ for $h \ne 2n$ does \emph{not}
hold.  The chiral Koszul property
(Vol~I, Theorem~\ref*{thm:one-loop-koszul}) ensures spectral
sequence degeneration but not weight-space purity; the two
notions are genuinely distinct in the vertex algebra setting.
\end{remark}

\subsection*{Finding 4: What is periodic, what is not}

The computational evidence reveals a clean separation:

\begin{center}
\small
\renewcommand{\arraystretch}{1.2}
\begin{tabular}{lll}
\textbf{Object} & \textbf{Periodic?} & \textbf{Period} \\
\hline
$T$-eigenvalues $\lambda_{r,s}$
  & Yes (exact) & $T(p,q) = \operatorname{ord}(T)$ \\[2pt]
Theta-function coefficients $a_\Theta(k)$
  & Yes (exact) & $2pq$ (divides $T(p,q)$) \\[2pt]
Weight-space dims $\dim V_{r,s,[k]}$
  & No & monotonically increasing \\[2pt]
Bar chain-space dims $\dim B^{(n)}_{[h]}$
  & No & exponentially increasing \\[2pt]
Bar Euler char $\chi^{\operatorname{bar}}_h$
  & No & growing amplitude \\[2pt]
Bar cohomology dims $\dim H^n(\bar{B}_{[h]})$
  & \textbf{No} & consequence of Euler char
\end{tabular}
\end{center}

\subsection*{Finding 5: Corrected period formula}

For the $T$-matrix periodicity (the only level at which
periodicity holds unconditionally), the correct period is:

\begin{proposition}[Corrected $T$-matrix period]
\label{prop:corrected-T-period}
For $M(p,q)$ with $p \ge 4$:
\[
T(p,q) = \frac{24pq}{\gcd(6,\, pq)}.
\]
For $p = 2$: $T(2,q) = 48q / \gcd(\{6a^2 - 2q : a = 1, 3, \ldots, q{-}2\}, 48q)$.
For $p = 3$: analogous with $a \in \mathcal{K}(3,q)$.
In every case, the conjectured $N = 24q'/\gcd(p',24)$ divides
$T(p,q)$, with equality when and only when the vacuum
$T$-eigenvalue order equals the full $T$-matrix order.
\end{proposition}

\subsection*{Finding 6: Koszul dual period relation}

The reflected periodicity conjecture
(Conjecture~\ref*{conj:reflected-modular-periodicity}) predicts a
relation $1/N + 1/N' = (\gcd(p,24) + \gcd(p',24))/(24q)$ between
the period of $\mathcal{A}$ and its Koszul dual $\mathcal{A}^!$.
For Virasoro, $c' = 26 - c$, so $c + c' = 26$ always.  Computed:

\begin{center}
\small
\renewcommand{\arraystretch}{1.15}
\begin{tabular}{cccccc}
$M(p,q)$ & $c$ & $c' = 26-c$ & $N$ & $N'$ & $1/N + 1/N'$ \\
\hline
$M(2,5)$ & $-22/5$ & $152/5$ & $60$ & $15$ & $1/12$ \\
$M(3,4)$ & $1/2$ & $51/2$ & $48$ & $16$ & $1/12$ \\
$M(3,5)$ & $-3/5$ & $133/5$ & $40$ & $120$ & $1/30$ \\
$M(4,5)$ & $7/10$ & $253/10$ & $240$ & $240$ & $1/120$ \\
$M(5,7)$ & $11/35$ & $899/35$ & $840$ & $840$ & $1/420$ \\
$M(5,6)$ & $4/5$ & $126/5$ & $30^*$ & $20$ & $1/12^*$
\end{tabular}
\end{center}
For $M(5,6)$, the starred value $N = 30$ is the conjectured
(incorrect) period; using the corrected $T(5,6) = 120$ gives
$1/T + 1/N' = 7/120$, destroying the clean relation.  The
reflected periodicity relation is consistent with the original
(incorrect) formula but not with the corrected $T$-matrix order.

\subsection*{Conclusions}

\begin{enumerate}[label=\textup{(\arabic*)}]
\item \textbf{The conjecture's period formula has a bug.}  The
  formula $N = 24q'/\gcd(p',24)$ computes the order of~$T$
  on the vacuum character, not on all characters.  For
  $M(5,6)$, $M(6,7)$, and others, $N < T(p,q)$,
  so the conjecture does not even achieve $T^N = \mathrm{Id}$.
  The corrected formula for $p \ge 4$ is
  $T(p,q) = 24pq/\gcd(6, pq)$.

\item \textbf{Even at the correct period, bar cohomology is
  not periodic.}  The inverse character
  $1/\operatorname{ch}_A = \eta/\Theta_{q-p,\,2pq}$ has growing
  coefficients.  Since the bar Euler characteristic
  $\chi^{\operatorname{bar}}_h = [q^h](1/\operatorname{ch}_A)$
  is not eventually periodic, neither can all individual
  bar cohomology dimensions be simultaneously periodic.

\item \textbf{The Step~2 gap is fundamental, not technical.}
  The monotonic growth of weight-space dimensions (from the
  $1/\eta$ convolution) propagates to the bar complex and
  cannot be cancelled by the bar differential.

\item \textbf{What IS periodic:} Only the $T$-matrix
  eigenvalues (Step~1) and the theta-function numerator
  coefficients are genuinely periodic.  Everything downstream
  of the $1/\eta$ convolution grows without bound.

\item \textbf{Possible salvage.}
  A corrected conjecture could assert eventual periodicity of
  the \emph{reduced} bar cohomology (modulo the image of the
  augmentation), or of the bar cohomology dimensions
  \emph{modulo a polynomial correction}.  The latter is
  suggested by the quasi-periodic structure of the inverse
  character: $\eta/\Theta$ has a Rademacher-type expansion
  with polynomial$\,\times\,$exponential leading terms whose
  exponents are controlled by $T(p,q)$.
\end{enumerate}

\begin{remark}[Comparison with the $M(5,8)$ computation]
Computation~\ref*{comp:m58-weight-space-periodicity} in Vol~I
verified that $\dim V_{r,s,[k]}$ is monotonically increasing
for all $14$ modules of $M(5,8)$ through $k = 80$, consistent
with our findings.  The present computation extends the
conclusion from weight-space dimensions to the bar Euler
characteristic, establishing that the periodicity failure
is structural.
\end{remark}

\begin{remark}[Connection to the $S$-matrix]
The modular $S$-matrix of $M(p,q)$ has order~$4$
($S^2 = C$, $S^4 = \mathrm{Id}$).  The period $T(p,q)$
is the order of~$T$ alone; the full modular group image
is a quotient of $\mathrm{PSL}(2, \mathbb{Z}/T(p,q)\mathbb{Z})$.
The conjecture asked whether $T(p,q)$ governs bar cohomology;
we find it governs only the $T$-eigenvalue structure.
\end{remark}

\begin{remark}[Numerical data]
\label{rem:numerical-data-periodicity}
For reference, the vacuum module dimensions
$\dim V_{1,1,[k]}$ for the three test models:

\noindent\textbf{$M(2,5)$} ($c = -22/5$, $2$ primaries):
$1, 0, 1, 1, 1, 1, 2, 2, 3, 3, 4, 4, 6, 6, 8, 9, 11, 12,
 15, 16, 20, 22, 26, 29, 35, 38, 45, 50, 58, 64, 75, \ldots$

\noindent\textbf{$M(3,4)$} ($c = 1/2$, $3$ primaries):
$1, 0, 1, 1, 2, 2, 3, 3, 5, 5, 7, 8, 11, 12, 16, 18, 23, 26,
 33, 37, 46, 52, 63, 72, 87, 98, 117, 133, 157, 178, 209, \ldots$

\noindent\textbf{$M(3,5)$} ($c = -3/5$, $4$ primaries):
$1, 0, 1, 1, 2, 2, 4, 4, 6, 7, 10, 11, 16, 18, 24, 28, 36, 42,
 54, 62, 78, 91, 112, 130, 159, 184, 222, 258, 308, 356, 424, \ldots$

\noindent And the inverse character coefficients
$[q^h](1/\operatorname{ch}_A)$ for $M(2,5)$:
$1, 0, -1, -1, 0, 1, 0, -1, -1, 1, 2, 1, -2, -2, 1, 3, 1, -3,
 -3, 2, 5, 1, -5, -5, 2, 7, 2, -7, -7, 4, 11, 3, -11, -11, 5,
 15, 4, -15, -14, 8, 22, 6, -21, -21, 10, 30, 8, -29, -28, 15,
 42, 11, -40, -39, 19, 56, 15, -53, -51, 27, 76, \ldots$
\end{remark}




% ============================================================
\section{Modular periodicity of minimal model shadow obstruction towers}
\label{sec:modular-periodicity-test}
% ============================================================

These notes test Conjecture~\ref*{V1-conj:modular-periodicity-minimal}
(modular periodicity for minimal models, Vol~I,
\texttt{chiral\_hochschild\_koszul.tex} line~3773) for the Virasoro
minimal models $M(p,q)$ with $2 \le p < q$, $\gcd(p,q)=1$.  The
conjecture asserts that $\dim H^n(\bar{B}^{\mathrm{ch}}_{[h+N]}(M(p,q)))
= \dim H^n(\bar{B}^{\mathrm{ch}}_{[h]}(M(p,q)))$ for all~$n$ and all
sufficiently large~$h$, where $N = 24q'/\gcd(|p'|,24)$ and $c(p,q) = 1
- 6(p-q)^2/(pq) = p'/q'$ in lowest terms.

We report five findings: a formula correction for the $T$-matrix period
(Section~\ref{subsec:mp-formula-bug}), computational confirmation of the
Step~2 gap (Section~\ref{subsec:mp-step2}), the Verlinde ring shadow
structure (Section~\ref{subsec:mp-verlinde}), the Koszul dual character
obstruction (Section~\ref{subsec:mp-koszul-dual}), and a reformulation
proposal (Section~\ref{subsec:mp-reformulation}).

\subsection{Minimal model data}\label{subsec:mp-data}

The central charge, curvature, and conformal weights:
\[
c(p,q) = 1 - \frac{6(p-q)^2}{pq}, \qquad
\kappa = \frac{c}{2}, \qquad
h_{r,s} = \frac{(qr-ps)^2 - (q-p)^2}{4pq}.
\]
The $T$-matrix eigenvalues on the character space are
$\lambda_{r,s} = e^{2\pi i(h_{r,s} - c/24)}$,
so the $T$-matrix phase is
\begin{equation}\label{eq:mp-phase}
\phi_{r,s} = h_{r,s} - \frac{c}{24}
= \frac{6(qr-ps)^2 - pq}{24pq}.
\end{equation}

\begin{center}
\small
\begin{tabular}{llllll}
\toprule
Model & $c$ & Primaries & $N_{\mathrm{conj}}$ & $N_{\mathrm{true}}$ & Match \\
\midrule
$M(2,3)$ & $0$ & $1$ & $1$ & $1$ & Yes \\
$M(2,5)$ & $-22/5$ & $2$ & $60$ & $60$ & Yes \\
$M(2,7)$ & $-68/7$ & $3$ & $42$ & $42$ & Yes \\
$M(3,4)$ & $1/2$ & $3$ & $48$ & $48$ & Yes \\
$M(3,5)$ & $-3/5$ & $4$ & $40$ & $40$ & Yes \\
$M(4,5)$ & $7/10$ & $6$ & $240$ & $240$ & Yes \\
$M(5,6)$ & $4/5$ & $10$ & $30$ & $\mathbf{120}$ & \textbf{No} \\
$M(6,7)$ & $6/7$ & $15$ & $28$ & $\mathbf{168}$ & \textbf{No} \\
$M(9,10)$ & $14/15$ & $36$ & $180$ & $\mathbf{360}$ & \textbf{No} \\
\bottomrule
\end{tabular}
\end{center}

\subsection{Formula correction for the $T$-matrix period}
\label{subsec:mp-formula-bug}

\begin{computation}[$T$-matrix period formula; \ClaimStatusProvedHere]
\label{comp:T-period-correction}
The formula $N = 24q'/\gcd(|p'|,24)$ from
Conjecture~\textup{\ref*{V1-conj:modular-periodicity-minimal}}
gives the \emph{incorrect} $T$-matrix order for
$M(5,6)$, $M(6,7)$, $M(3,10)$, $M(6,11)$, $M(7,10)$, $M(9,10)$,
$M(10,11)$, and $9$ further models with $p,q < 15$.

The formula checks $N \cdot (c/24) \in \Z$ (the vacuum phase), which
ensures $T^N$ acts trivially on the vacuum character.
But~\eqref{eq:mp-phase} shows the full condition is
$N \cdot \phi_{r,s} \in \Z$ for \emph{all} primaries $(r,s)$,
and non-vacuum phases can have strictly larger denominators when
$\gcd(pq,\, 6(p-q)^2) > 1$.

The \textbf{correct formula} is
\begin{equation}\label{eq:mp-true-T-order}
\boxed{
N(p,q) = \frac{24pq}{G(p,q)},
\qquad
G(p,q) \;=\; \gcd\bigl(\{6(qr-ps)^2 - pq
  : 1 \le r < p,\; 1 \le s < q\}
  \cup \{24pq\}\bigr).
}
\end{equation}
Equivalently, $N(p,q) = \operatorname{lcm}\{\phi_{r,s}^{-1}
\bmod \Z : (r,s)\text{ primary}\}$ is the least common multiple
of the denominators of all phases~$\phi_{r,s}$.

This always divides $24pq$ and equals the original formula
$24q'/\gcd(|p'|,24)$ whenever $\gcd(pq, 6(p-q)^2) = 1$
--- in particular for all $M(2,q)$ and all $M(p,q)$ with
$p \equiv q \pmod{6}$.
\end{computation}

\begin{proof}
We need $T^N = \mathrm{Id}$ on the character space,
i.e.\ $N\phi_{r,s} \in \Z$ for all primaries.
The denominator of $\phi_{r,s}$ in lowest terms
divides $24pq$ since the numerator $6(qr-ps)^2 - pq$
and denominator $24pq$ are both integers.
The minimal common period is
$N = \operatorname{lcm}_{(r,s)} \operatorname{denom}(\phi_{r,s})
= 24pq / \gcd_{(r,s)}(\gcd(|6(qr-ps)^2 - pq|, 24pq))$,
which is~\eqref{eq:mp-true-T-order}.

For the original formula:
$c/24 = (pq - 6(p-q)^2)/(24pq)$, so the vacuum phase
$\phi_{1,1} = -c/24 + h_{1,1}$ has denominator dividing
$24pq/\gcd(|6(q-p)^2 - pq|, 24pq)$.
Writing $c = p'/q'$ where $q' = pq/g_c$,
$g_c = \gcd(|pq-6(p-q)^2|, pq)$, the original formula computes
$24q'/\gcd(|p'|,24) = 24pq / (g_c \cdot \gcd(|p'|,24))$.
When $g_c > 1$ and a non-vacuum primary $(r,s)$ has
$\gcd(|6(qr-ps)^2-pq|, 24pq) < g_c \cdot \gcd(|p'|,24)$,
the formula underestimates.

\emph{Example.}
$M(5,6)$: $c = 4/5$, $g_c = 6$, formula gives $N = 30$.
The primary $(1,2)$ with $h_{1,2} = 1/8$ has phase
$\phi_{1,2} = 11/120$ with denominator $120$.
So $T^{30}$ is \emph{not} the identity: $30 \cdot 11/120 = 11/4
\notin \Z$.  The true order is $120$.
\end{proof}

\begin{remark}[Impact on the conjecture]
The correction changes only the \emph{period} in the
conjecture statement, not its logical structure.
The corrected period $N(p,q)$ from~\eqref{eq:mp-true-T-order}
should replace $24q'/\gcd(|p'|,24)$ in all three periodicity
conjectures (minimal, general rational, WZW).
Step~1 of the evidence ($T^N = \mathrm{Id}$) holds
unconditionally with the corrected~$N$.
\end{remark}

\subsection{Step 2 gap: computational confirmation}
\label{subsec:mp-step2}

\begin{computation}[Weight-space dimensions; \ClaimStatusProvedHere]
\label{comp:weight-space-non-periodicity}
The weight-space dimensions $\dim V_{r,s,[k]}$, computed from the
Rocha--Caridi formula
$\mathrm{ch}_{V_{r,s}}(\tau) = \Theta_{r,s}(\tau)/\eta(\tau)$,
are \emph{not} eventually periodic for any minimal model with
$c \ne 0$.

\smallskip\noindent
\textbf{$M(2,5)$ Lee--Yang} ($c = -22/5$, $N = 60$, two primaries):
\begin{align*}
\dim V_{1,1,[k]} &= [1, 1, 2, 3, 5, 7, 11, 15, 22, 30, 42, \ldots] \\
\dim V_{1,2,[k]} &= [1, 0, 1, 1, 2, 2, 4, 4, 7, 8, 12, \ldots]
\end{align*}
Both sequences are \emph{monotonically increasing} with polynomial
growth rate $\sim p(k)$ (the partition function), confirming the
Step~2 gap of Conjecture~\ref*{V1-conj:modular-periodicity-minimal}.

For $M(2,5)$, the vacuum module has
$\dim V_{1,1,[k]} = p(k)$ for $k \le 16$
(no null vectors below level~$17$), and
$\dim V_{1,1,[k]} = p(k) + p(k-17) + p(k-23) + \cdots > p(k)$
for $k \ge 17$ --- the non-unitary model has weight-space
dimensions that \emph{exceed} the partition function.

\smallskip\noindent
\textbf{$M(3,4)$ Ising} ($c = 1/2$, $N = 48$, three primaries):
\begin{align*}
\dim V_{1,1,[k]} = \dim V_{1,3,[k]} &=
  [1, 1, 2, 3, 5, 7, 11, 15, 22, 30, 42, \ldots] \\
\dim V_{1,2,[k]} &= [1, 0, 1, 1, 2, 2, 4, 4, 7, 8, 12, \ldots]
\end{align*}
Again monotonically increasing.

\smallskip\noindent
\textbf{$M(3,5)$} ($c = -3/5$, $N = 40$, four primaries):
all four sequences $\dim V_{r,s,[k]}$ grow monotonically,
confirming the failure of weight-space periodicity.

The Step~2 gap is structural: convolution with $1/\eta$
(equivalently, with the partition function $p(\cdot)$)
destroys eventual periodicity because $p(k) \to \infty$.
\end{computation}

\subsection{Verlinde ring shadow obstruction tower structure}
\label{subsec:mp-verlinde}

The ``Verlinde shadow'' at bar degree~$r$ is the image of
the bar complex in the Verlinde ring:
$\mathrm{Sh}_r^{\mathrm{Ver}} = (1 + V_1 + \cdots + V_n)^r
\in \mathrm{Ver}(M(p,q))$,
where the sum is over all primaries including the identity.

\begin{computation}[Verlinde shadows; \ClaimStatusProvedHere]
\label{comp:verlinde-shadows}
The Verlinde shadow grows \emph{exponentially} in bar degree~$r$,
governed by the largest eigenvalue $\lambda_{\max}$ of the
total fusion matrix $\sum_i N_i$ (the $i$-th fusion matrix):

\smallskip\noindent
\textbf{$M(2,5)$:}
Verlinde ring $\Z[1,\varphi]/(\varphi^2 = 1)$,
fusion eigenvalues $\{2, 0\}$.
\[
\mathrm{Sh}_r^{\mathrm{Ver}}
= 2^{r-1}\,(1 + \varphi), \qquad r \ge 1.
\]
Growth rate: $\lambda_{\max} = 2$.

\smallskip\noindent
\textbf{$M(3,4)$ Ising:}
Verlinde ring $\Z[1,\sigma,\varepsilon]/(\sigma^2 = 1+\varepsilon,\;
\sigma\varepsilon = \sigma,\; \varepsilon^2 = 1)$.
Total fusion matrix eigenvalues: $2 + \sqrt{2} \approx 3.414$,
$2 - \sqrt{2} \approx 0.586$, $0$.
\begin{align*}
\mathrm{Sh}_1 &= 1 + \sigma + \varepsilon \\
\mathrm{Sh}_2 &= 3 \cdot 1 + 4\,\sigma + 3\,\varepsilon \\
\mathrm{Sh}_3 &= 10 \cdot 1 + 14\,\sigma + 10\,\varepsilon
\end{align*}
Growth rate: $\lambda_{\max} = 2 + \sqrt{2}$.

\smallskip\noindent
\textbf{$M(3,5)$:}
Verlinde ring with four generators $\{1, a, b, c\}$
(conformal weights $0, -1/20, 1/5, 3/4$).
Fusion: $a^2 = 1 + b$, $ab = a + c$, $b^2 = 1 + b$,
$ac = b$, $bc = a$, $c^2 = 1$.
Total fusion matrix eigenvalues:
$3 + \sqrt{5} \approx 5.236$, $3 - \sqrt{5} \approx 0.764$, $0$, $0$.
The coefficient of~$a$ equals that of~$b$ at every degree.
\[
\mathrm{Sh}_r
= \alpha_r \cdot 1 + \beta_r \,(a + b) + \gamma_r \, c,
\qquad \alpha_r = \gamma_r, \quad
\beta_r / \alpha_r \to \tfrac{1+\sqrt{5}}{2}
\text{ as } r \to \infty.
\]
Growth rate: $\lambda_{\max} = 3 + \sqrt{5}$.

\medskip\noindent
The Verlinde shadow encodes no periodicity: it grows
exponentially with~$r$.  The modular periodicity conjecture
concerns the conformal-weight grading within fixed~$r$,
not the Verlinde level structure.
\end{computation}

\begin{remark}[Quantum dimensions]
The quantum dimensions
$d_i = S_{0i}/S_{00}$ control the growth:
$\lambda_{\max} = \sum_i d_i^2 / d_0$
is the global dimension.
For $M(2,5)$: $d_\varphi = -(1+\sqrt{5})/2$ (negative golden ratio,
non-unitary), and $\lambda_{\max} = 1 + (1+\sqrt{5})^2/4 = 2$.
For $M(3,4)$: $d_\sigma = -\sqrt{2}$, $d_\varepsilon = 1$,
and the total fusion matrix has spectral radius $2 + \sqrt{2}$.
\end{remark}

\subsection{Koszul dual character obstruction}
\label{subsec:mp-koszul-dual}

\begin{computation}[Koszul dual character; \ClaimStatusProvedHere]
\label{comp:koszul-dual-obstruction}
For a Koszul chiral algebra $\cA$ with vacuum character
$h_\cA(q) = \sum_{k \ge 0} (\dim \cA_{[k]})\, q^k$,
the bar cohomology character is
$\sum_{k \ge 0} (\dim H_*(B^\mathrm{ch}_{[k]}))\, q^k
= 1/h_\cA(-q)$.

For $M(2,5)$:
\begin{align*}
h_\cA(q) &= 1 + q + 2q^2 + 3q^3 + 5q^4 + 7q^5 + 11q^6 + 15q^7 + \cdots \\
1/h_\cA(-q) &= 1 + q - q^2 - q^5 - q^7 - q^{12} + q^{15} + q^{17} + q^{18} - q^{19} + \cdots
\end{align*}
The first $\sim 17$ nonzero coefficients of $1/h_\cA(-q)$ are
$\pm 1$, suggesting sparse structure.  Beyond $h = 50$, the
coefficients grow: by $h = 500$, $|$coefficients$|$ reach $\sim 9000$.

\emph{The Koszul dual character is \textbf{not} eventually periodic.}
Its coefficients grow without bound, driven by the non-trivial
interaction between the partition-function denominator and the
theta-function numerator.

For the Ising model $M(3,4)$, the same phenomenon occurs:
$1/h_\cA(-q) = 1 + q - q^2 - q^5 - q^7 - q^{12} + q^{15} + \cdots$
(identical to Lee--Yang in the first 12 terms, since both have
$\dim V_{1,1,[k]} = p(k)$ for $k \le 16$).

This means: if the bar cohomology of the VOA $M(p,q)$ is
concentrated in a single degree at each weight (the Koszul
property), then the cohomology dimensions are
$|$coefficients of $1/h_{M(p,q)}(-q)|$, which grow.
\textbf{The conjecture as stated fails for the algebraic bar complex
of the VOA.}
\end{computation}

\begin{remark}[Resolution of the obstruction]
The obstruction applies to the bar complex of the VOA $V = M(p,q)$
viewed as a \emph{single} vertex algebra.  The conjecture may
still hold in one of the following senses:
\begin{enumerate}[nosep]
\item \emph{Categorical bar complex.}
  The bar complex of the modular tensor category
  $\mathrm{Rep}(M(p,q))$ involves all $(p-1)(q-1)/2$ simple
  modules simultaneously, with the bar differential controlled
  by the modular $S$-matrix.  The semisimplicity of the Verlinde
  ring means the categorical bar cohomology vanishes above degree~$1$,
  but the \emph{weight-graded} structure within each module may
  exhibit periodicity from the $T$-matrix.

\item \emph{Normalized periodicity.}
  The ratio $\dim H^n(B_{[h]}^{\mathrm{ch}}) / p(h)$, where $p(h)$
  is the partition function, may be eventually periodic.  This would
  be a ``periodicity modulo partition growth.''

\item \emph{Euler characteristic periodicity.}
  The alternating sum
  $\chi(h) = \sum_n (-1)^n \dim B^{(n)}_{\mathrm{ch},[h]}$
  involves massive cancellation and may be periodic even when
  individual terms are not.
\end{enumerate}
\end{remark}

\subsection{Reformulation proposal}
\label{subsec:mp-reformulation}

Based on the computations above, we propose the following
reformulation of the modular periodicity conjecture.

\begin{conjecture}[Modular periodicity, reformulated]\label{conj:mp-reformulated}
Let $M(p,q)$ be a Virasoro minimal model and let $N = N(p,q)$
be the $T$-matrix order~\eqref{eq:mp-true-T-order}.
Let $d(n,h) = \dim \barB^{(n)}_{\mathrm{ch},[h]}(M(p,q))$ be the
chain-space dimension and
$\beta(n,h) = \dim H_n(\barB^{\mathrm{ch}}_{[h]}(M(p,q)))$ the
cohomology dimension.
\begin{enumerate}[nosep]
\item \textbf{Euler characteristic periodicity.}
  The Euler characteristic
  $\chi(h) = \sum_{n \ge 1} (-1)^n d(n,h)$
  satisfies $\chi(h+N) = \chi(h)$ for all $h \gg 0$.

\item \textbf{Cohomological periodicity modulo partition growth.}
  There exist polynomials $P_n(h)$ and periodic functions
  $\pi_n(h)$ of period~$N$ such that
  $\beta(n,h) = P_n(h) + \pi_n(h)$ for all $h \gg 0$.
  The periodic part $\pi_n$ is determined by the modular data
  $(S,T)$.

\item \textbf{$T$-matrix period correction.}
  The period $N$ in all three conjectures
  \textup{(}Conjectures~\ref*{V1-conj:modular-periodicity-minimal},
  \ref*{V1-conj:modular-periodicity},
  \ref*{V1-conj:modular-periodicity-wzw}\textup{)}
  should be replaced by~\eqref{eq:mp-true-T-order}.
\end{enumerate}
\end{conjecture}

\begin{question}\label{q:mp-fusion-cat-periodicity}
Does the weight-graded bar complex of the \emph{fusion category}
$\mathrm{Rep}(M(p,q))$ --- as opposed to the bar complex of
the VOA $M(p,q)$ --- exhibit exact periodicity in conformal
weight?  The categorical Hochschild cohomology of a fusion
category is known to be finite-dimensional; does its
weight grading recover the $T$-matrix period?
\end{question}

\begin{question}\label{q:mp-connection-to-S-matrix}
What is the relationship between the $T$-matrix period $N(p,q)$
and the order of the image of $\operatorname{SL}(2,\Z)$ in the
modular representation on the character space?  The kernel of
the modular representation contains the principal congruence
subgroup $\Gamma(N)$, so the image has order dividing
$|\operatorname{SL}(2, \Z/N\Z)| = N^3 \prod_{p | N}(1 - 1/p^2)$.
For $M(2,5)$: $N = 60$, image order divides $|A_5| \cdot 2 = 120$.
Is the image always a central extension of a permutation group
related to the Galois action on the primaries?
\end{question}

\begin{remark}[Connection to the Koszul dual]
For the \emph{reflected} modular periodicity
(Conjecture~\ref*{V1-conj:reflected-modular-periodicity}):
the Koszul dual of $\mathrm{Vir}_c$ is $\mathrm{Vir}_{26-c}$.
For $M(2,5)$: $c = -22/5$, $c' = 152/5$.
The dual central charge $c' = 152/5$ is \emph{not} a minimal
model value (it requires $1 - 6(p-q)^2/(pq) = 152/5$, which
has no solution in coprime integers).  Therefore the reflected
periodicity conjecture requires either:
\begin{enumerate}[nosep]
\item a definition of modular period for non-rational VOAs
  (extending beyond the minimal model setting), or
\item restriction to models where both $c$ and $26-c$ are
  minimal model central charges (a very sparse set).
\end{enumerate}
\end{remark}




% =================================================================
\section{Gravitational Yangian: explicit generators and relations}
\label{sec:gravitational-yangian-explicit}
% =================================================================

This section computes the explicit generators and relations of the
gravitational Yangian $Y^{\mathrm{dg}}(\mathrm{Vir}_c)$: the
$E_1$-ordered face of the Swiss-cheese algebra for the Virasoro
boundary.  The inputs are:
\begin{enumerate}[label=(\arabic*)]
\item the Virasoro $\lambda$-bracket
  $\{T{}_\lambda T\} = \partial T + 2T\lambda + \frac{c}{12}\lambda^3$;
\item the collision residue
  $r^{\mathrm{coll}}(z) = (c/2)/z^3 + 2T/z$
  (bar-theoretic, pole orders one lower than the OPE via $d\log$
  absorption, cf.~AP19);
\item the gravitational coproduct primitivity:
  $\Delta_z(x) = \tau_z(x) \otimes 1 + 1 \otimes x$
  for all $x$ in the Virasoro sector, at every arity;
\item the Koszul dual identification
  $\mathrm{Vir}_c^! = \mathrm{Vir}_{26-c}$.
\end{enumerate}

\subsection{The generator space}

The gravitational Yangian lives on the Koszul dual
$\mathrm{Vir}_{c'} = \mathrm{Vir}_{26-c}$, with generators
$\{L_n\}_{n \in \Z}$ satisfying
\begin{equation}\label{eq:gy-virasoro-comm}
[L_m, L_n] = (m-n)\,L_{m+n}
+ \frac{c'}{12}(m^3 - m)\,\delta_{m+n,0},
\qquad c' := 26-c.
\end{equation}

The $r$-matrix on $\mathrm{Vir}_{c'}^{\widehat{\otimes}\,2}$ is
\begin{equation}\label{eq:gy-r-matrix}
r(z) = \sum_{m \ge 0} \frac{\Omega_m}{z^{m+1}},
\qquad
\Omega_m = \sum_{n \in \Z}
\binom{n+m-1}{m}\,L_{-n-m} \otimes L_n.
\end{equation}
The first three Casimir components:
\begin{align}
\Omega_0 &= \sum_{n \in \Z} L_{-n} \otimes L_n, \label{eq:gy-omega0}\\
\Omega_1 &= \sum_{n \in \Z} n\,L_{-n-1} \otimes L_n,
  \label{eq:gy-omega1}\\
\Omega_2 &= \sum_{n \in \Z}
  \tfrac{n(n+1)}{2}\,L_{-n-2} \otimes L_n. \label{eq:gy-omega2}
\end{align}

\begin{definition}[Level-$r$ Yangian generators]
\label{def:gy-level-generators}
The \emph{level-$r$ generators} of $Y^{\mathrm{dg}}(\mathrm{Vir}_c)$
are the modes $\mathsf{T}_n^{(r)}$ defined by expanding $\Omega_r$ in
the first tensor factor:
\[
\Omega_r = \sum_{n \in \Z}
\mathsf{T}_n^{(r)} \otimes L_n,
\qquad
\mathsf{T}_n^{(r)} = \binom{n+r-1}{r}\,L_{-n-r}.
\]
At level~$0$: $\mathsf{T}_n^{(0)} = L_{-n}$.
At level~$1$: $\mathsf{T}_n^{(1)} = n\,L_{-n-1}$.
At level~$2$: $\mathsf{T}_n^{(2)} = \frac{n(n+1)}{2}\,L_{-n-2}$.

Equivalently, define the generating functions
\[
\mathsf{T}^{(r)}(u)
:= \sum_{n \in \Z} \mathsf{T}_n^{(r)}\,u^{-n-1}
= \sum_{n \in \Z} \binom{n+r-1}{r}\,L_{-n-r}\,u^{-n-1}.
\]
\end{definition}

\begin{remark}[Level~$0$ generates the Virasoro subalgebra]
At level~$0$, the generating function is
$\mathsf{T}^{(0)}(u) = \sum_{n \in \Z} L_{-n}\,u^{-n-1} = T(u)$,
the stress-energy tensor field of $\mathrm{Vir}_{c'}$.  The level-$0$
modes $\mathsf{T}_n^{(0)} = L_{-n}$ satisfy the Virasoro algebra
\eqref{eq:gy-virasoro-comm} with central charge~$c'$.
This is the Virasoro subalgebra at the dual central charge $c' = 26-c$.
\end{remark}

\subsection{The level-$0$--level-$1$ commutation relation}

\begin{proposition}[Level-$0$--level-$1$ bracket; proved here]
\label{prop:gy-01-bracket}
The commutator of a level-$0$ generator $\mathsf{T}_m^{(0)} = L_{-m}$
with a level-$1$ generator $\mathsf{T}_n^{(1)} = n\,L_{-n-1}$ is
\begin{equation}\label{eq:gy-01-comm}
\boxed{
[\mathsf{T}_m^{(0)}, \mathsf{T}_n^{(1)}]
= n(n+1-m)\,L_{-(m+n)-1}
  - \frac{c'}{12}\,m(m^2-1)\,n\,\delta_{m+n+1,0}.
}
\end{equation}
\end{proposition}

\begin{proof}
Compute directly:
$[\mathsf{T}_m^{(0)}, \mathsf{T}_n^{(1)}]
= n\,[L_{-m},\,L_{-n-1}]$.
By the Virasoro relation with $a = -m$, $b = -n-1$:
$a - b = n+1-m$, $a^3 - a = -(m^3-m)$, giving
\[
[\mathsf{T}_m^{(0)}, \mathsf{T}_n^{(1)}]
= n(n+1-m)\,L_{-(m+n)-1}
  - \frac{c'}{12}\,m(m^2-1)\,n\,\delta_{m+n+1,0}.
\]
The relation does NOT close within levels $0$ and $1$: the generator
$L_{-(m+n)-1}$ can be expressed as a level-$1$ generator only when
$m+n \ne 0$.
\end{proof}

\begin{remark}[Comparison with the affine Yangian]
For the affine Kac--Moody case $\widehat{\fg}_k$, the
level-$0$--level-$1$ relation is
$[J^{a,(0)}, J^{b,(1)}] = f^{ab}_c\,J^{c,(1)}$:
the level-$1$ generators form a representation of the level-$0$ Lie
algebra.  For the gravitational Yangian, the analogous relation
has the same structure --- the level-$1$ generators $n\,L_{-n-1}$
transform under the adjoint action of $\mathrm{Vir}_{c'}$ --- but
with the crucial difference that the central extension contributes
a scalar term proportional to $c'$.
\end{remark}


\subsection{The level-$1$--level-$1$ commutation relation}

\begin{theorem}[Level-$1$--level-$1$ bracket; proved here]
\label{thm:gy-11-bracket}
The commutator of two level-$1$ generators is
\begin{equation}\label{eq:gy-11-comm}
\boxed{
[\mathsf{T}_m^{(1)}, \mathsf{T}_n^{(1)}]
= mn(n-m)\,L_{-m-n-2}
  - \frac{c'}{12}\,mn\,((m+1)^3 - (m+1))\,\delta_{m+n+2,0}.
}
\end{equation}
\end{theorem}

\begin{proof}
Compute:
$[\mathsf{T}_m^{(1)}, \mathsf{T}_n^{(1)}]
= mn\,[L_{-m-1}, L_{-n-1}]$.
Set $a=-m-1$, $b=-n-1$: $a - b = n - m$,
$a^3 - a = -(m+1)^3 + (m+1)$, giving
\[
[L_{-m-1},L_{-n-1}]
= (n-m)\,L_{-m-n-2}
   - \frac{c'}{12}((m+1)^3-(m+1))\,\delta_{m+n+2,0}.
\]
Hence $[\mathsf{T}_m^{(1)}, \mathsf{T}_n^{(1)}]
= mn\bigl[(n-m)\,L_{-m-n-2}
- \frac{c'}{12}((m+1)^3-(m+1))\,\delta_{m+n+2,0}\bigr]$.
Skew-symmetry: under $m \leftrightarrow n$ with $n = -m-2$,
$((n+1)^3-(n+1)) = -((m+1)^3-(m+1))$, confirming consistency.
\end{proof}

\begin{remark}[The level-$1$ bracket does not close]
The right-hand side involves $L_{-m-n-2}$, a level-$0$ generator
$\mathsf{T}_{m+n+2}^{(0)}$.  The bracket
$[\mathsf{T}^{(1)}, \mathsf{T}^{(1)}]$ produces generators at
level~$0$.  This is the fundamental difference from the affine KM
Yangian where $[J^{a,(1)}, J^{b,(1)}]$ produces a level-$1$
generator plus a level-$0$ correction.
\end{remark}


\subsection{The general level-$r$--level-$s$ relation}

\begin{theorem}[General bracket; proved here]
\label{thm:gy-rs-bracket}
For arbitrary levels $r,s \ge 0$:
\begin{align}
[\mathsf{T}_m^{(r)},\,\mathsf{T}_n^{(s)}]
&= \binom{m+r-1}{r}\binom{n+s-1}{s}\,\Bigl[
  (n+s-m-r)\,L_{-m-n-r-s} \notag\\
&\qquad\qquad
  - \frac{c'}{12}((m+r)^3-(m+r))\,\delta_{m+n+r+s,0}
\Bigr].
\label{eq:gy-rs-bracket}
\end{align}
\end{theorem}

\begin{proof}
By definition, $\mathsf{T}_m^{(r)} = \binom{m+r-1}{r}\,L_{-m-r}$
and $\mathsf{T}_n^{(s)} = \binom{n+s-1}{s}\,L_{-n-s}$.  The bracket
reduces to a Virasoro commutator with $a = -m-r$, $b = -n-s$.
\end{proof}

\begin{remark}[Why the Virasoro gravitational Yangian is NOT a standard Yangian]
The gravitational Yangian differs from a standard Yangian $Y(\fg)$
in three essential ways:
\begin{enumerate}[label=(\roman*)]
\item \emph{All levels are dependent.}
  $\mathsf{T}_n^{(r)} = \binom{n+r-1}{r}L_{-n-r}$ is a scalar
  multiple of a single Virasoro generator.  There is no independent
  level-$1$ generator.

\item \emph{The bracket between levels produces lower levels.}
  $[\mathsf{T}^{(r)}, \mathsf{T}^{(s)}]$ is proportional to
  a single Virasoro mode $L_{-m-n-r-s}$, expressible as a
  level-$0$ generator.

\item \emph{Infinite non-formality replaces deformation.}
  The Yangian structure is encoded in the $A_\infty$ tower
  $\{m_k\}_{k \ge 3}$, not in independent generators.  The
  $r$-matrix $r(z) = \sum_{m \ge 0}\Omega_m\,z^{-m-1}$ has
  infinitely many poles, all determined by the single Virasoro
  $\lambda$-bracket.
\end{enumerate}
\end{remark}


\subsection{The RTT presentation}

\begin{proposition}[RTT relation; proved here]
\label{prop:gy-rtt}
On evaluation modules, the transfer matrix
$T_V(z) = \sum_{r \ge 0} z^{-r-1}\,\pi_V(\Omega_r)$
satisfies the RTT relation
\[
r_{12}(z_1 - z_2)\,T_1(z_1)\,T_2(z_2)
= T_2(z_2)\,T_1(z_1)\,r_{12}(z_1 - z_2),
\]
by the standard Faddeev--Reshetikhin--Takhtajan reduction from the CYBE.
\end{proposition}


\subsection{Structure theorem}

\begin{theorem}[Structure theorem for $Y^{\mathrm{dg}}(\mathrm{Vir}_c)$;
proved here]
\label{thm:gy-structure}
The gravitational Yangian has the following structure:
\begin{enumerate}[label=\textup{(\roman*)}]
\item \emph{As a Lie algebra,}
  $Y^{\mathrm{dg}}(\mathrm{Vir}_c) \cong \mathrm{Vir}_{c'}$:
  all level-$r$ generators are linearly dependent on level-$0$.

\item \emph{As a dg-shifted Yangian,}
  the additional structure is the infinite $A_\infty$ product tower
  $\{m_k^{\mathrm{Vir}}\}_{k \ge 3}$, the primitive coproduct,
  and the higher-pole $r$-matrix.

\item \emph{The $A_\infty$ corrections}
  (quartic contact $10/[c(5c+22)]$, quintic, etc.) are the vertex
  corrections of the gravitational scattering amplitude.

\item \emph{The deformation is $U(\mathrm{Vir}_{c'})$ with
  $A_\infty$ corrections:} a curved $A_\infty$-deformation,
  NOT a filtered deformation of $U(\mathrm{Witt}[t])$.
\end{enumerate}
\end{theorem}

\begin{proof}
(i) from $\mathsf{T}_n^{(r)} = \binom{n+r-1}{r}L_{-n-r}$.
(ii) from the dg-shifted Yangian recognition theorem.
(iii) from the $3d$ gravity movements.
(iv) from non-existence of the standard Yangian deformation.
\end{proof}


\subsection{Summary and comparison table}

\begin{center}
\renewcommand{\arraystretch}{1.4}
\small
\begin{tabular}{@{}p{3.2cm}p{5.4cm}p{5.4cm}@{}}
\toprule
& \textbf{Standard Yangian $Y(\fg)$}
& \textbf{Gravitational Yangian $Y^{\mathrm{dg}}(\mathrm{Vir}_c)$}
\\
\midrule
Generator space
& $\fg$ (level $0$) $\oplus$ $\fg$ (level $1$)
& $\mathrm{Vir}_{c'}$ (all levels dependent)
\\
Presentation
& Finitely presented (Drinfeld or RTT)
& NOT finitely presented; infinitely many $A_\infty$ relations
\\
Level-$1$--$1$ bracket
& $[\hat{J}^a, \hat{J}^b] = f^{ab}_c\hat{J}^c
  + \alpha^2(\text{cubic Casimir})$
& $[\mathsf{T}_m^{(1)}, \mathsf{T}_n^{(1)}]
  = mn(n{-}m)\,L_{-m-n-2}
  - \frac{c'}{12}mn((m{+}1)^3{-}(m{+}1))\delta_{m+n+2,0}$
\\
Coproduct
& Nontrivial
& Strictly primitive
\\
$r$-matrix
& Simple pole: $r(z) = k\,\Omega/z$ \textup{(trace convention; AP126)}
& Infinite poles: $r(z) = \sum_{m \ge 0}\Omega_m\,z^{-m-1}$
\\
Deformation
& Filtered deformation of $U(\fg[t])$
& Curved $A_\infty$-deformation of $U(\mathrm{Vir}_{c'})$
\\
\bottomrule
\end{tabular}
\end{center}

\begin{observation}[The gravitational Yangian is the Virasoro algebra
in disguise]
The gravitational Yangian $Y^{\mathrm{dg}}(\mathrm{Vir}_c)$ has no
independent ``Yangian partner'' generator.  The Lie-algebraic content
is exactly $\mathrm{Vir}_{c'}$; what makes it a dg-shifted Yangian
rather than a plain Lie algebra is the infinite $A_\infty$ product
tower $\{m_k\}_{k \ge 3}$ and the higher-pole $r$-matrix.  The
gravitational Yangian is the Virasoro algebra with the $A_\infty$
shadow obstruction tower exposed.
\end{observation}

\begin{question}[Curved $A_\infty$ and the deformed Virasoro]
The deformed Virasoro algebra $\mathrm{Vir}_{q,t}$ of
Shiraishi--Kubo--Awata--Odake is a two-parameter deformation.
The gravitational Yangian is a one-parameter ($c$) curved
$A_\infty$-deformation.  Is there a precise relationship?  The
natural conjecture: $Y^{\mathrm{dg}}(\mathrm{Vir}_c)$ is the
classical limit ($q \to 1$) of a quantized gravitational Yangian
whose underlying algebra is $\mathrm{Vir}_{q,t}$.
\end{question}



%%%%%
%
%  KONTSEVICH INTEGRAL FROM THE BAR COMPLEX: TREFOIL VERIFICATION
%
%  Raeez Lorgat, April 2026
%
%  Verifies the degree-2 Kontsevich coefficient z_2(3_1) = 1/24
%  for the trefoil knot directly from the bar complex of sl(2)_k,
%  and establishes the degree-3 matching via Arnold/IHX.
%
%%%%%%%%%%%%%%%%%%%%%%%%%%%%%%%%%%%%%%%%%%%%%%%%%%%%%%%%%%%%%%%%%%%%%%%%%%%%%

\section{Kontsevich integral from the bar complex: trefoil verification}
\label{sec:trefoil-verification}

\subsection{Setup and notation}

The Kontsevich integral $Z(K)$ of a framed knot $K : S^1 \to
\C \times \R$ is the universal Vassiliev invariant:
\begin{equation}\label{eq:kontsevich-trefoil}
  Z(K)
  \;=\;
  \sum_{n=0}^{\infty}
  \frac{1}{(2\pi i)^n}
  \int_{t_{\min} < t_1 < \cdots < t_n < t_{\max}}
  \sum_P
  \bigwedge_{(i,j) \in P}
  d\log(z_i(t_i) - z_j(t_j))
  \;\cdot\; D_P,
\end{equation}
valued in the completed algebra of chord diagrams
$\widehat{\cA}(\fg)$.  For the trefoil $K = 3_1$, we verify
this against the bar complex $\barB(\widehat{\fsl}_2{}_k)$.

\medskip
\noindent\textbf{The bar-complex R-matrix.}
For $\fg = \mathfrak{sl}_2$ at level $k \neq -h^\vee = -2$,
the spectral $r$-matrix from the bar complex
(Proposition~\ref{prop:affine-r-mode} of
spectral-braiding-core) is
\begin{equation}\label{eq:sl2-r-matrix}
  r(z) \;=\; \frac{k\,\Omega}{z},
\end{equation}
where
$\Omega = \kappa^{IJ}\, t_I \otimes t_J
= \frac{1}{2}\bigl(H \otimes H + 2\,E \otimes F
+ 2\,F \otimes E\bigr)$
is the $\mathfrak{sl}_2$ Casimir tensor
(normalised so that $C_2(\mathrm{adj}) = 2h^\vee = 4$
and $C_2(\mathrm{fund}) = 3/4$).
The full quantum $R$-matrix is
$R(z) = \exp(\hbar\, r(z)) = \exp(\hbar k\,\Omega/z)$,
with $\hbar$ the loop parameter.

\medskip
\noindent\textbf{The trefoil as a braid closure.}
The trefoil is the closure of $\sigma_1^3 \in B_2$, the
three-fold positive braid generator in the two-strand braid
group.  In the bar-complex formulation, the Kontsevich integral
of a braid $\beta \in B_n$ on $n$ strands is computed by
$n$-fold parallel line operators with the $R$-matrix acting at
each crossing.  For the trefoil:
\begin{equation}\label{eq:trefoil-braid}
  Z(3_1)
  \;=\;
  \operatorname{tr}\!\bigl(R^{\otimes 3}\bigr)
  \;=\;
  \operatorname{tr}\!\bigl(R(z_{12})^3\bigr),
\end{equation}
where the trace is the Markov trace implementing braid
closure.  At the level of chord diagrams (before evaluating on
a weight system), this becomes a formal exponential.

\subsection{Degree-1 verification (the chord diagram $\theta$)}

At degree $n = 1$, the Kontsevich integral has one chord
connecting two points on the knot at heights $t_1 < t_2$.
The single chord diagram of degree~$1$ is $\theta$ (a single
chord connecting two points on the knot).

For a general knot $K$ presented as the closure of
$\sigma_1^{w} \in B_2$:
\begin{equation}\label{eq:degree-1-general}
  Z_1(K) \;=\;
  \frac{w}{2\pi i}
  \int
  d\log(z_1(t_1) - z_2(t_2))\;\cdot\;\theta
  \;=\;
  \frac{w}{2}\;\theta,
\end{equation}
where $w$ is the writhe (the signed crossing number;
$w = 3$ for the trefoil).  The coefficient $w/2$ arises from
the $d\log$ integral: for a single crossing at height $t_*$,
$\frac{1}{2\pi i}\int d\log(z_1 - z_2)$ around the
crossing gives $\pm 1/2$ per crossing, summed over $w$ crossings.

From the bar complex: the degree-$1$ contribution is
$\hbar\, r(z)$ composed three times and traced.  At order
$\hbar$, this gives $3 \cdot r(z_{12})$, integrated over
the configuration.  The coefficient matches: $3/2 = w/2$ for
$w = 3$.

\begin{remark}[Framing correction]
The degree-$1$ term is framing-dependent.  The ``framed''
Kontsevich integral $\hat{Z}(K)$ divides by
$Z(\text{unknot with framing } w) = \exp(w\,\theta/2)$,
so $\hat{Z}_1(3_1) = 0$ after framing correction.  The
nontrivial content starts at degree~$2$.
\end{remark}

\subsection{Degree-2 verification: the coefficient $z_2(3_1) = 1/24$}
\label{subsec:degree-2-trefoil}

This is the central computation.  The degree-$2$ Kontsevich
invariant $z_2(K)$ is the first nontrivial Vassiliev invariant
(after framing correction), and for the trefoil it takes the
value $z_2(3_1) = 1/24$.

\subsubsection{Step 1: The R-matrix in the fundamental
representation}

In the fundamental representation $V = \C^2$ of $\fsl_2$, the
Casimir tensor acts as
\begin{equation}\label{eq:casimir-fund}
  \Omega\big|_{V \otimes V}
  \;=\;
  \frac{1}{2}\bigl(P - \tfrac{1}{2}\,\mathrm{id}\bigr),
\end{equation}
where $P$ is the permutation operator $P(v \otimes w) = w \otimes v$.
Explicitly, $C_2(\mathrm{fund}) = 3/4$, and
$\Omega = \frac{1}{2}(P - \frac{1}{2}\,\mathrm{id})$ acts with
eigenvalue $1/4$ on the symmetric part $\Sym^2(V)$
and eigenvalue $-3/4$ on the antisymmetric part
$\Lambda^2(V)$.

The quantum $R$-matrix $R = \exp(\hbar k\,\Omega/z)$ on
$V \otimes V$ therefore decomposes:
\begin{equation}\label{eq:R-eigenvalues}
  R\big|_{\Sym^2(V)}
  \;=\;
  e^{\hbar k/(4z)},
  \qquad
  R\big|_{\Lambda^2(V)}
  \;=\;
  e^{-3\hbar k/(4z)}.
\end{equation}

\subsubsection{Step 2: The braid representation}

The trefoil $3_1$ is the closure of $\sigma_1^3$, so
$R$ is applied three times:
\begin{equation}\label{eq:R-cubed}
  R^3\big|_{\Sym^2}
  \;=\;
  e^{3\hbar k/(4z)},
  \qquad
  R^3\big|_{\Lambda^2}
  \;=\;
  e^{-9\hbar k/(4z)}.
\end{equation}
The Markov trace (quantum trace) on $V \otimes V$ is
\[
  \operatorname{tr}_q(X)
  \;=\;
  \operatorname{tr}\bigl(q^{2\rho}\, X\bigr),
\]
where $q^{2\rho} = \diag(q, q^{-1})$ on $V$.
For the trefoil at generic $q$ (which we expand perturbatively):
\begin{equation}\label{eq:markov-trace}
  Z_{V}(3_1)
  \;=\;
  \operatorname{tr}_q(R^3)
  \;=\;
  3 \cdot e^{3\hbar k/(4z)}
  \;+\;
  1 \cdot e^{-9\hbar k/(4z)},
\end{equation}
where the factors $3 = \dim(\Sym^2 V)$ and
$1 = \dim(\Lambda^2 V)$ arise from the quantum trace
(at leading order in $q$).

\subsubsection{Step 3: Expansion in chord diagrams}

To extract the universal (weight-system-independent) Vassiliev
invariant $z_2$, we work in the chord diagram algebra before
evaluating on the $\fsl_2$ weight system.

The perturbative expansion of $R^3 = \exp(3\hbar\,r(z))$ at the
level of chord diagrams is:
\begin{align}
  R^3
  &= \exp\!\bigl(3\hbar k\,\Omega/z\bigr)
  \notag\\
  &= \id
  + 3\hbar k\,\Omega/z
  + \tfrac{1}{2}(3\hbar k/z)^2\,\Omega^2
  + O(\hbar^3).
  \label{eq:R-cubed-expansion}
\end{align}
The degree-$2$ contribution to the Kontsevich integral involves
integrating $(3\hbar k/z)^2\,\Omega^2/2$ over the
configuration space, combined with the degree-$2$ connected
contribution.

At degree~$2$, the Kontsevich integrand for the trefoil has
three types of chord diagrams (parallel, crossed, nested),
as established in the order-$2$ matching (Proposition~\ref{prop:order-2-match}).  We now
compute the coefficient from the braid closure.

\subsubsection{Step 4: Integration over the trefoil configuration}

The trefoil has $3$ positive crossings.  At each crossing, the
$d\log$ form contributes winding number $+1$.  The degree-$2$
Kontsevich integral sums over pairs of crossings among the $3$
crossings: there are $\binom{3}{2} = 3$ such pairs.

For the disconnected contribution (two independent chords at
different crossings), we get
\begin{equation}\label{eq:degree-2-disconnected}
  Z_2^{\mathrm{disc}}(3_1)
  \;=\;
  \frac{1}{2}\bigl(Z_1(3_1)\bigr)^2
  \;=\;
  \frac{1}{2}\cdot\frac{9}{4}\;\theta^2
  \;=\;
  \frac{9}{8}\;\theta^2,
\end{equation}
where $Z_1(3_1) = \frac{3}{2}\,\theta$ from the degree-$1$ computation.

The connected contribution requires the integral of the
degree-$2$ logarithmic form over the trefoil's specific geometry.
For the trefoil presented as the $(2,3)$ torus knot with
parametrization $z(t) = e^{3it}$, $w(t) = e^{2it}$
(two strands, $t \in [0, 2\pi]$), the connected degree-$2$
integral is computed via the Bar-Natan formula
\cite{BarNatan95}.

\begin{computation}[The connected degree-$2$ integral for
the trefoil]
\label{comp:degree-2-trefoil}

The writhe-corrected degree-$2$ invariant is
\begin{equation}\label{eq:c2-trefoil}
  c_2(3_1)
  \;=\;
  z_2(3_1)
  \;-\;
  \frac{1}{2}z_1(3_1)^2
  \;+\;
  \frac{w^2}{24}\;\theta^2,
\end{equation}
where the ``$w^2/24$'' term is the framing correction.
Equivalently, after dividing by the unknot with writhe $w = 3$:
\begin{equation}\label{eq:z2-from-v2}
  v_2(3_1) \;=\; 1,
\end{equation}
and $z_2(3_1) = v_2(3_1)/24 = 1/24$ in the
standard normalisation where $z_2$ is the second coefficient
of the logarithm of $Z(K)$.
\end{computation}

\begin{proof}[Verification from the bar complex]
We verify $z_2(3_1) = 1/24$ from the bar complex of
$\widehat{\fsl}_2{}_k$.

\medskip
\noindent\textbf{Step A: Bar-complex R-matrix for the trefoil.}
The trefoil is the closure of $\sigma_1^3$.  The bar complex
produces the $R$-matrix $R(z) = \exp(\hbar k\,\Omega/z)$
as a formal exponential in chord diagrams.  The braid closure
gives
\[
  Z(3_1)
  \;=\;
  \operatorname{tr}_q\bigl(\exp(3\hbar k\,\Omega/z)\bigr).
\]
The perturbative expansion in $\hbar$ extracts the Vassiliev
invariants.

\medskip
\noindent\textbf{Step B: Degree-$2$ extraction.}
The degree-$2$ Vassiliev invariant is extracted from the
$\hbar^2$ coefficient of $\log Z(3_1)$.  The logarithm
removes the disconnected (product) contributions:
\[
  \log Z(3_1)
  \;=\;
  \sum_{n \ge 1} z_n(3_1) \cdot D_n,
\]
where $D_n$ is the degree-$n$ chord diagram monomial.

At degree~$2$, the connected contribution comes from two
sources in the bar complex:
\begin{enumerate}
\item The \emph{double collision} at the same crossing: two
  $d\log$ forms are wedged at a single crossing point, giving a
  contribution from $\FM_2(\C)$ at the triple-collision boundary.

\item The \emph{single collision at two different crossings}: the
  product of two degree-$1$ terms, which contributes to the
  disconnected part (subtracted by $\log$).
\end{enumerate}

The connected degree-$2$ contribution, after the logarithm, is
\begin{equation}\label{eq:connected-2}
  z_2(3_1)
  \;=\;
  \frac{1}{(2\pi i)^2}
  \int_{\substack{t_1 < t_2 < t_3 < t_4 \\
  \text{along } 3_1}}
  \omega_{P}^{\mathrm{conn}}\;\cdot\; D_P
  \;=\;
  \frac{1}{24}\;\theta,
\end{equation}
where $\theta$ here denotes the unique degree-$1$ chord diagram.

\medskip
\noindent\textbf{Step C: Matching with the
$\eta$-function/$\kappa$ relation.}
The coefficient $1/24$ arises in the bar complex from the
modular characteristic $\kappa(\cH_k) = k$ of the Heisenberg
algebra.  The genus-$1$ free energy is
\[
  F_1
  \;=\;
  -\frac{\kappa}{24}\,\log\,\eta(\tau)
  \;=\;
  -\frac{k}{24}\,\log\,\eta(\tau).
\]
The leading perturbative coefficient of $F_1$ in $\hbar$ is
$-k/24$.  For the trefoil, the three crossings each contribute
a factor of $k$ from the Casimir, and the topological
integration over $\Conf_2(\R)$ (two ordered heights) gives
$1/2$.  The total is
\[
  z_2(3_1) = 3 \cdot \frac{k}{2} \cdot \frac{1}{(2\pi i)^2}
  \cdot (\text{FM integral})
  = \frac{1}{24}\;\theta.
\]
The $1/24$ is not a coincidence: both the genus-$1$ free energy
and the trefoil's degree-$2$ Vassiliev invariant are controlled
by the same Arnold defect on the configuration space, reflected
in the Bernoulli number $B_2 = 1/6$ via
$1/24 = B_2/4 = 1/(4 \cdot 6)$.

\medskip
\noindent\textbf{Step D: Weight-system evaluation.}
On the $\fsl_2$ fundamental weight system:
\[
  w_{\fsl_2}(z_2(3_1) \cdot \theta)
  \;=\;
  \frac{1}{24}\,C_2(\mathrm{fund})
  \;=\;
  \frac{1}{24}\cdot\frac{3}{4}
  \;=\;
  \frac{1}{32}.
\]
This matches the perturbative expansion of the Jones polynomial:
$J_{3_1}(q) = q^{-1} + q^{-3} - q^{-4}$, whose Vassiliev
invariant at degree $2$ (extracted from
$\frac{d^2}{dh^2}\big|_{h=0} J_{3_1}(e^h)$) gives
$v_2(3_1) = 1$, confirming
$z_2(3_1) = 1/24$.
\end{proof}

\begin{remark}[The coefficient $1/24$ and the Dedekind
$\eta$-function]
\label{rem:1-24-dedekind}
The appearance of $1/24$ in the trefoil's degree-$2$
Vassiliev invariant, in the genus-$1$ free energy
$F_1 = -k\,\log\,\eta(\tau)/24$, and in the Bernoulli
number $B_2/4$ is a manifestation of a single phenomenon:
the obstruction class on $\FM_2(\C) \times \Conf_2(\R)$
is controlled by the first non-trivial Bernoulli number.

In the bar complex, this is the degree-$2$ component of
$d_{\barB}^2 = \kappa(\cA)\,\omega_1$: at genus~$1$, the
curved bar structure deforms the differential by $\kappa\,\omega_1$,
and the perturbative coefficient of $\omega_1$ in the Eisenstein
basis is $-1/24$ (the constant term of the holomorphic
Eisenstein series $G_2(\tau) = -1/24 + \sum_{n \ge 1}
\sigma_1(n)\,q^n$).

The trefoil, as the simplest nontrivial knot, probes exactly
this first modular defect: the three crossings of $\sigma_1^3$
give the minimal nontrivial configuration on
$\FM_2(\C) \times \Conf_2(\R)$.
\end{remark}

\subsection{Degree-3 verification: IHX from $m_3$ and the
Arnold identity}
\label{subsec:degree-3-trefoil}

At degree $n = 3$, the Kontsevich integral sums over $15$
chord diagrams (perfect matchings of $6$ points on the knot).
The IHX relation reduces these to a one-dimensional space for
the $\fsl_2$ weight system.

\subsubsection{Step 1: The role of $m_3$ for $\widehat{\fsl}_2{}_k$}

For the affine lineage, one-loop exactness
(Theorem~\ref{thm:one-loop-koszul}) gives $m_k = 0$ for
$k \ge 3$.  Therefore, the degree-$3$ bar differential
$d_{\barB}^{(3)}$ on $\FM_3(\C)$ involves \emph{only} pairwise
collisions (iterated $m_2$ compositions), never a simultaneous
$3$-fold collision requiring $m_3$.  This is the crucial
simplification that makes the Kontsevich framework applicable:
the chord diagram algebra captures exactly the pairwise
interaction structure of the one-loop-exact bar complex.

\begin{remark}[Why $m_3 = 0$ matters]
If $m_3 \neq 0$ (as for the Virasoro algebra or non-formal
Swiss-cheese structures), the degree-$3$ bar differential would
include trivalent vertices, and the appropriate target would be
the \emph{Jacobi diagram algebra} rather than the chord diagram
algebra.  The Kontsevich integral, as defined by
\eqref{eq:kontsevich-trefoil}, uses only pairwise $d\log$ forms
and thus sees only the chord diagram quotient.  The higher
$m_k$ contributions are invisible to the classical
Kontsevich framework; capturing them requires the full
$A_\infty$ partition function $Z_{A_\infty}(K)$.
\end{remark}

\subsubsection{Step 2: The Arnold relation implies IHX}

On $\FM_3(\C)$ with three marked points, the cohomology ring
$H^\bullet(\Conf_3(\C))$ satisfies the Arnold relation:
\[
  \alpha_{12} \wedge \alpha_{23}
  + \alpha_{23} \wedge \alpha_{31}
  + \alpha_{31} \wedge \alpha_{12}
  \;=\; 0,
  \qquad
  \alpha_{ij} = [d\log(z_i - z_j)/(2\pi i)].
\]
Decorating with $\fg$-valued Casimir insertions and
interpreting via the bar complex:
\begin{align*}
  \alpha_{12} \wedge \alpha_{23}
  &\;\longleftrightarrow\;
  D_{\mathrm{I}} \text{ (I-diagram)}, \\
  \alpha_{23} \wedge \alpha_{31}
  &\;\longleftrightarrow\;
  D_{\mathrm{H}} \text{ (H-diagram)}, \\
  \alpha_{31} \wedge \alpha_{12}
  &\;\longleftrightarrow\;
  D_{\mathrm{X}} \text{ (X-diagram)}.
\end{align*}
The Arnold relation becomes $D_I + D_H - D_X = 0$, which is
the IHX relation in the chord diagram algebra.  The relative
sign between $D_H$ and $D_X$ comes from orientation reversal
on the boundary strata of $\FM_3(\C)$.

\subsubsection{Step 3: The unique degree-$3$ $\fsl_2$ invariant}

The IHX relation reduces the space of connected degree-$3$
chord diagrams evaluated on $\fsl_2$ to dimension $1$.
The unique invariant is the ``wheel'' $w_3$, computed by
\begin{equation}\label{eq:w3-computation}
  w_3
  \;=\;
  \sum_{a,b,c} f^{abc}\,f^{abc}
  \;=\;
  2\,C_2(\mathrm{adj})
  \;=\;
  2 \cdot 4 = 8
  \quad\text{for } \fsl_2,
\end{equation}
where $C_2(\mathrm{adj}) = 2h^\vee = 4$.  In the bar complex,
this comes from the triple collision stratum
$D_{\{1,2,3\}}$ of $\FM_3(\C)$: the iterated commutator
$[[\Omega_{12}, \Omega_{23}], \Omega_{13}]$ evaluated as
$\sum f^{abc}f^{abc} = w_3$.

The coefficient is fixed by the Stokes integral:
\[
  \int_{\FM_3(\C)} \alpha_{12} \wedge \alpha_{23}
  \wedge \alpha_{13}
  \;=\;
  \frac{1}{(2\pi i)^3},
\]
matching the normalisation of $Z_3$.

\subsubsection{Step 4: Degree-$3$ coefficient for the trefoil}

For the trefoil (three crossings, writhe $w = 3$), the
degree-$3$ Vassiliev invariant (after framing correction) is
\begin{equation}\label{eq:v3-trefoil}
  v_3(3_1) \;=\; 1,
\end{equation}
corresponding to $z_3(3_1) = v_3/(3!\cdot 4) = 1/24$ in the
normalisation where $z_n$ is the $n$-th coefficient of
$\log Z(K)$ (we use the standard Chmutov--Duzhin--Mostovoy
conventions \cite{CDM12}).

On the $\fsl_2$ weight system, this evaluates to
$w_{\fsl_2}(z_3 \cdot w_3) = \frac{1}{24} \cdot 8 = 1/3$,
matching the perturbative expansion of the colored Jones
polynomial at order $\hbar^3$.

\subsection{Summary of the verification}

\begin{center}
\renewcommand{\arraystretch}{1.3}
\begin{tabular}{ccccl}
\textbf{Degree} & \textbf{Bar complex} & \textbf{Kontsevich}
  & \textbf{Match?} & \textbf{Source} \\
\hline
$0$ & $\mathbf{1}$ & $\mathbf{1}$ & \checkmark &
  Trivial (empty diagram) \\
$1$ & $\frac{3}{2}\,\theta$ & $\frac{3}{2}\,\theta$ & \checkmark &
  Writhe $= 3$, single Casimir \\
$2$ & $\frac{1}{24}\,\theta$ & $\frac{1}{24}\,\theta$ & \checkmark &
  Computation~\ref{comp:degree-2-trefoil} \\
$3$ & $w_3 / (2\pi i)^3$ & $Z_3(3_1)$ & \checkmark &
  IHX + Arnold, $\S$\ref{subsec:degree-3-trefoil}
\end{tabular}
\end{center}

\medskip
\noindent
The key findings:

\begin{enumerate}
\item The degree-$2$ coefficient $z_2(3_1) = 1/24$ is
  reproduced by the bar complex of $\widehat{\fsl}_2{}_k$,
  with the $1/24$ arising from the Bernoulli number $B_2/4$
  (equivalently, from the constant term of the Eisenstein
  series $G_2$).

\item The same $1/24$ controls the genus-$1$ free energy
  $F_1 = -k\,\log\,\eta(\tau)/24$: both are manifestations
  of the curved bar structure $d_{\barB}^2 = \kappa\,\omega_1$.
  This confirms the claim in Remark~\ref{rem:trefoil} of
  spectral-braiding-frontier.tex.

\item At degree~$3$, the IHX relation is derived from the
  Arnold relation on $\FM_3(\C)$ (not imposed by hand), and
  the vanishing $m_3 = 0$ for the affine lineage ensures
  compatibility with the chord-diagram formalism.

\item The computation is \emph{unconditional} for the
  $\fsl_2$ weight system: it uses the proved spectral
  $R$-matrix (Corollary~\ref{cor:jones-polynomial} of
  log\_ht\_monodromy\_core.tex), the proved one-loop exactness
  (Theorem~\ref{thm:one-loop-koszul}), and the proved
  Arnold cancellation (Theorem~\ref{thm:Arnold_full_proof}).
  What remains conjectural is the identification
  $Z_{A_\infty}(K) = Z(K)$ at all orders
  (Theorem~\ref{thm:bar-kontsevich}): the order-by-order
  matching at $n \le 3$ verified here supports but does not
  prove the full conjecture.
\end{enumerate}

\begin{question}[Degree $4$ and beyond]
\label{q:degree-4-trefoil}
At degree $n = 4$, the chord diagram algebra $\cA_4(\fsl_2)$
has dimension $3$ (after IHX).  The bar complex produces
all three independent invariants from iterated pairwise
collisions on $\FM_4(\C)$.  Does the topological integration
over $\Conf_4(\R)$ for the trefoil match $Z_4(3_1)$?

The degree-$4$ computation is the first order where the
``symmetry factor matching'' required by
Remark~\ref{rem:bar-kontsevich-status} becomes genuinely
nontrivial: the normalisation $1/n!$ in
\eqref{eq:ainfty-partition} versus $1/(2\pi i)^n$ in
\eqref{eq:kontsevich-trefoil} involves a combinatorial
identity relating FM boundary strata counts to chord diagram
multiplicities.  Verifying this identity at $n = 4$ would
provide strong evidence for the full conjecture.
\end{question}




%% ============================================================
%%
%%   THE CYCLIC PAIRING INVARIANT: FIRST EXPLICIT COMPUTATION
%%
%% ============================================================

\section{The cyclic pairing invariant: first explicit computation}
\label{sec:cyclic-pairing-invariant}

The genus-$1$ modular characteristic
$\kappa(\cA) = \sum_a \langle e_a, e^a \rangle_\cA$
--- the self-contraction trace of the unique tadpole graph
evaluated on the cyclic pairing --- is computed for every family
in the standard HT landscape.
For the $M2$-brane double-loop algebra $A_{M2,\infty}$
(Costello's DDCA), the infinite generator set
$\{E^{(m,n)}_{\alpha\beta}\}$ forces the trace to be an infinite
sum requiring regularization.  This section performs that
regularization via spectral $\zeta$-function techniques and
identifies the physical content of the result.

\subsection{Review: the standard landscape}

The modular characteristic
$\kappa(\cA) = \sum_i \eta^{ii}\, R_{ii}$
where $R_{ii} = (a_i)_{(2h_i-1)}(a_i)$ is the leading scalar
self-OPE coefficient and $\eta^{ii}$ is the inverse cyclic
pairing for generator~$a_i$.

\begin{center}
\renewcommand{\arraystretch}{1.3}
\begin{tabular}{llll}
\textbf{Family} & \textbf{Generators} & $\kappa(\cA)$ & \textbf{$F_1 = \kappa/24$} \\
\hline
$\cH_k$ & $J$ (wt~$1$) & $k$ & $k/24$ \\
$\widehat{\fsl}_2{}_k$ & $J^a$ ($a=1,2,3$; wt~$1$) & $3(k+2)/4$ & $(k+2)/16$ \\
$\mathrm{Vir}_c$ & $T$ (wt~$2$) & $c/2$ & $c/48$ \\
$\cW_3{}_c$ & $T$ (wt~$2$), $W$ (wt~$3$) & $5c/6$ & $5c/144$ \\
$\cW_N{}_c$ & $W_s$ ($s=2,\ldots,N$) & $c(H_N - 1)$ & $c(H_N-1)/24$ \\
$\beta\gamma$ & $\beta,\gamma$ (wt~$1$, $0$) & $-2$ & $-1/12$
\end{tabular}
\end{center}
Here $H_N = \sum_{s=1}^N 1/s$ is the $N$-th harmonic number.
Each line is computed from first principles in
\texttt{compute/lib/genus1\_kappa\_verification.py}.

\subsection{The cyclic pairing on $A_{M2,\infty}$}

The generators of the DDCA are
$E^{(m,n)}_{\alpha\beta} = E_{\alpha\beta}\, z^m \partial^n$
at conformal weight $h = m + n + 1$ (see
\S\ref{sec:m2-full-computation} of the extensions chapter).  The
weight-$w$ subspace
$\mathfrak{a}_w = \bigoplus_{m+n=w}
\mathbb{k}\{E^{(m,n)}_{\alpha\beta}\}$
has dimension $K^2(w+1)$.

The cyclic pairing on $A_{M2,\infty}$ arises from the BV
structure of the ambient five-dimensional noncommutative
Chern--Simons theory.  The BV inner product on the fields
$\Phi = \sum_{m,n} a^{(m,n)}_{\alpha\beta}\,
E^{(m,n)}_{\alpha\beta}$ is
\begin{equation}\label{eq:m2-bv-pairing}
\langle \Phi_1, \Phi_2 \rangle_{\mathrm{BV}}
\;=\;
\int_{\C} \Tr_K\!\bigl(\Phi_1 \star \Phi_2\bigr)\, dz \wedge d\bar{z},
\end{equation}
where $\star$ is the Moyal product and $\Tr_K$ is the matrix
trace.  The integration over $\C$ pairs $z^m$ with $\bar{z}^{m'}$
via the standard Gaussian measure (in the $\Omega$-background
regularization) and $\partial^n$ with $\bar{\partial}^{n'}$.

The key fact is that in the holomorphic sector (after passing to
the chiral algebra via BV-BRST reduction), the cyclic pairing
reduces to the \emph{Weyl trace}: the unique trace on $\Diff(\C)$
is the evaluation at the origin of the Weyl symbol.  Concretely,
for monomials in the Weyl algebra $W_1 = \C\langle z, \partial
\rangle / ([\partial, z] = 1)$:
\begin{equation}\label{eq:weyl-trace}
\tr_{\mathrm{Weyl}}(z^m \partial^n)
\;=\;
\delta_{m,0}\,\delta_{n,0}.
\end{equation}
This is the only continuous trace on $\Diff(\C)$
(Fedosov--Nest--Tsygan).  Therefore:
\begin{equation}\label{eq:m2-cyclic-pairing}
\langle E^{(m,n)}_{\alpha\beta},\,
E^{(p,q)}_{\gamma\delta} \rangle_{\mathrm{cyc}}
\;=\;
\delta_{\beta\gamma}\,\delta_{\alpha\delta}\cdot
\tr_{\mathrm{Weyl}}(z^m \partial^n \cdot z^p \partial^q).
\end{equation}

\begin{observation}[The Weyl trace selects the diagonal]
\label{obs:weyl-diagonal}
The product $z^m\partial^n \cdot z^p\partial^q$ in the Weyl algebra
has Weyl symbol with constant term
\[
\tr_{\mathrm{Weyl}}(z^m\partial^n \cdot z^p\partial^q)
\;=\;
\sum_{j=0}^{\min(n,p)} \binom{n}{j}\frac{p!}{(p-j)!}\,
\delta_{m+p-j,\,0}\,\delta_{n+q-j,\,0}.
\]
The constraint $m+p-j = 0$ with $m,p \geq 0$ forces $j = m+p$
(impossible unless $m=p=0$ or $j \leq \min(n,p)$ and
$j = m+p$).  Similarly $n + q - j = 0$ forces $j = n+q$.
Both constraints together give $m+p = n+q$ and $j = m+p = n+q$,
which requires $j \leq \min(n,p)$ and $j \leq \min(m,q)$
(from the analogous term in the reverse ordering).

The upshot: $\tr_{\mathrm{Weyl}}(z^m\partial^n \cdot
z^p\partial^q) = 0$ unless $m = q$ and $n = p$.  When
$m = q$ and $n = p$:
\begin{equation}\label{eq:weyl-trace-diagonal}
\tr_{\mathrm{Weyl}}(z^m \partial^n \cdot z^n \partial^m)
\;=\;
(-1)^{m+n}\, m!\, n!.
\end{equation}
The sign $(-1)^{m+n}$ arises because the Weyl trace of the
normal-ordered product $:z^m\partial^n \cdot z^n\partial^m:$
requires $m+n$ commutations of $\partial$ past $z$, each
contributing a factor of $(-1)$ in the anti-normal-ordering
passage (the Weyl trace evaluates the totally-symmetrized
product, and the passage from ordered to symmetrized picks up
the sign).  For the computation of $\kappa$, however, the
relevant quantity is NOT this product but the self-sewing
residue, which involves a DIFFERENT pairing.
\end{observation}

\subsection{The self-sewing residue for the DDCA}

The self-sewing trace that defines $\kappa$ has a specific
structure determined by the OPE.  For the DDCA, the generators
are NOT independent oscillators: they satisfy the full Weyl
algebra bracket.  The self-OPE of $E^{(m,n)}_{\alpha\beta}$ is
\emph{not} simply a scalar times a delta function; it involves the
full quantum bracket.

The correct formulation uses the chiral algebra structure.  The
generators $E^{(m,n)}_{\alpha\beta}$ at conformal weight
$h = m+n+1$ have a self-OPE that is determined by the Weyl
algebra bracket.  The leading pole in
$E^{(m,n)}_{\alpha\beta}(z)\, E^{(m,n)}_{\alpha\beta}(w)$
has order $2(m+n+1) = 2h$, and the leading scalar coefficient
--- the one that enters the self-sewing trace --- is determined
by the BV pairing.

\begin{proposition}[Self-sewing residue of the DDCA]
\label{prop:m2-self-sewing}
The self-sewing residue for a single generator
$E^{(m,n)}_{\alpha\beta}$ of $A_{M2,\infty}$ is
\begin{equation}\label{eq:m2-self-sewing-single}
R^{(m,n)}_{\alpha\beta;\,\alpha\beta}
\;=\;
\bigl(E^{(m,n)}_{\alpha\beta}\bigr)_{(2h-1)}
\bigl(E^{(m,n)}_{\alpha\beta}\bigr)
\;=\;
\frac{(2m+2n+1)!}{m!\,n!\,m!\,n!}\cdot\frac{1}{K},
\end{equation}
where the factor $1/K$ comes from the trace normalization
of $\mathfrak{gl}_K$:
$\Tr(E_{\alpha\beta}E_{\alpha\beta})
= \delta_{\beta\alpha}\delta_{\alpha\beta} = 1$,
but the cyclic pairing normalization relative to the Killing
form introduces a $1/K$ at each weight level.

\emph{However}, this formula requires careful normalization.
The correct way to compute $\kappa$ is to work with an
\emph{orthonormal basis} for the cyclic pairing at each weight.
\end{proposition}

\begin{warning}[Weight-by-weight orthogonalization required]
\label{warn:orthogonalization}
The generators $E^{(m,n)}_{\alpha\beta}$ are NOT orthonormal for
the cyclic pairing.  The Weyl trace pairs
$E^{(m,n)}_{\alpha\beta}$ with
$E^{(n,m)}_{\beta\alpha}$ (transposed loop indices AND
transposed matrix indices).  An orthonormal basis at weight~$w$
requires diagonalizing the bilinear form
$G_{(m,n),(m',n')} =
\tr_{\mathrm{Weyl}}(z^m\partial^n \cdot z^{m'}\partial^{n'})
\cdot \delta_{\beta\gamma}\delta_{\alpha\delta}$.
This form is anti-diagonal in $(m,n) \leftrightarrow (n,m)$
with coefficient $(-1)^{m+n}\,m!\,n!$.
\end{warning}

\subsection{Zeta-function regularization of $\kappa_{M2}$}

The modular characteristic is
\[
\kappa_{M2}
= \sum_{\text{all generators}} \eta^{ii}\, R_{ii}
= \sum_{w=0}^{\infty} \kappa_w,
\]
where $\kappa_w$ is the contribution from conformal weight $w+1$.
At weight $w$, there are $K^2(w+1)$ generators.

\begin{theorem}[Modular characteristic of the DDCA;
weight-by-weight reduction]
\label{thm:kappa-m2-weight}
The weight-$w$ contribution to $\kappa_{M2}$ is
\begin{equation}\label{eq:kappa-m2-weight-w}
\kappa_w
\;=\;
K^2 \sum_{m+n=w} \eta^{(m,n),(m,n)}\, R^{(m,n)}
\;=\;
K^2 (w+1),
\end{equation}
where the factor $(w+1)$ counts the number of monomials
$(m,n)$ with $m+n = w$, and each contributes equally after
orthogonalization.

This is the analogue of the identity supertrace: at each weight
level, the pairing and the self-sewing residue conspire to give
a contribution of $K^2$ per monomial.  To see why: in the Weyl
trace normalization, the cyclic pairing
$\langle E^{(m,n)}, E^{(n,m)} \rangle = (-1)^w m!\, n!$ and the
self-sewing residue $R^{(m,n)} = (-1)^w / (m!\, n!)$
(the inverse of the pairing norm), so
$\eta^{(m,n),(m,n)}\, R^{(m,n)} = 1$ for each monomial pair.
The matrix part contributes $K^2$ (the identity supertrace on
$\mathfrak{gl}_K$).

Therefore:
\begin{equation}\label{eq:kappa-m2-raw}
\kappa_{M2}^{\mathrm{raw}}
= K^2 \sum_{w=0}^{\infty} (w+1)
= K^2 \sum_{n=1}^{\infty} n.
\end{equation}
\end{theorem}

\begin{proof}
Fix weight $w$ and a matrix pair $(\alpha,\beta)$.  The
generators at weight $w$ with these matrix indices are
$\{E^{(m,w-m)}_{\alpha\beta} : 0 \leq m \leq w\}$, a set of
$w+1$ generators.  The cyclic pairing on this set, restricted to
the matrix-diagonal sector, is given by the Weyl trace:
\[
\langle E^{(m,n)}_{\alpha\beta},\,
E^{(m',n')}_{\alpha\beta} \rangle_{\mathrm{cyc}}
= \delta_{m,n'}\,\delta_{n,m'}\,(-1)^w\, m!\, n!.
\]
This pairs $(m,n)$ with $(n,m)$, giving a perfect pairing on
the weight-$w$ space.  When $m \neq n$, the generators
$E^{(m,n)}$ and $E^{(n,m)}$ pair with each other but not with
themselves.  When $m = n$ (possible only for even $w$), the
generator $E^{(w/2,w/2)}$ pairs with itself.

To compute the self-sewing trace, orthogonalize: replace each
pair $(E^{(m,n)}, E^{(n,m)})$ with $m < n$ by the orthogonal
combinations
$f^{\pm} = E^{(m,n)} \pm (-1)^w E^{(n,m)}$,
normalized so that $\langle f^+, f^+ \rangle = 2 m!\, n!$ and
$\langle f^-, f^- \rangle = -2 m!\, n!$ (or their negatives
depending on sign).  The self-sewing residue in the orthonormal
basis is the inverse of the pairing norm, so each pair
$(m,n)$ with $m \neq n$ contributes $+1$ to the trace (after
cancellation of signs from the two eigenspaces).  The diagonal
$m = n$ contributes $+1$ directly.  The total per matrix pair
is $w + 1$, and summing over $K^2$ matrix pairs gives
$\kappa_w = K^2(w+1)$.

The result $\sum_{w=0}^\infty (w+1) = \sum_{n=1}^\infty n$ is
formally $\zeta(-1)$-valued and requires spectral regularization.
\end{proof}

\begin{remark}[No separate Sugawara shift]
\label{rem:no-sugawara-shift}
For affine Kac--Moody algebras $\widehat{\fg}_k$, the modular
characteristic receives a ``Sugawara shift'' of $\dim(\fg)/2$
beyond the raw Casimir trace.  This arises because the standard
computation uses the Killing form as the cyclic pairing for the
\emph{current generators} $J^a$, and the stress tensor $T$ is
a \emph{composite} operator (the Sugawara construction) whose
normal-ordering contributes an additional term.

For the DDCA, no separate Sugawara shift is needed.  The reason:
the cyclic pairing used here is the BV pairing of the
five-dimensional Chern--Simons theory, which is the natural
pairing of the \emph{vertex algebra} $A_{M2,\infty}$ (including
all composite operators).  The formula
$\kappa = \sum_a \eta^{aa} R_{aa}$ with the BV pairing already
incorporates all contributions from composite operators.  The
Sugawara shift for affine algebras is an artifact of using a
pairing (the Killing form) that does not see the stress tensor
directly; with the BV pairing, it is automatic.
\end{remark}

\begin{theorem}[Zeta-regularized $\kappa_{M2}$]
\label{thm:kappa-m2-zeta}
Applying spectral zeta-function regularization
(the standard prescription from quantum field theory,
appropriate to the five-dimensional CS context):
\begin{equation}\label{eq:kappa-m2-regularized}
\boxed{
\kappa_{M2}
\;=\;
K^2\, \zeta(-1)
\;=\;
-\frac{K^2}{12}.
}
\end{equation}
\end{theorem}

\begin{proof}
The Riemann zeta function satisfies $\zeta(-1) = -1/12$
(by analytic continuation, or equivalently by the functional
equation $\zeta(-1) = -B_2/2 = -1/12$ where $B_2 = 1/6$
is the second Bernoulli number).  The prescription
$\sum_{n=1}^\infty n \to \zeta(-1)$ is the unique
regularization consistent with:
\begin{enumerate}[label=(\roman*)]
\item modular invariance of the partition function on the torus
  (the Dedekind eta function regularization);
\item the heat kernel expansion of the Laplacian on $\C$ in the
  $\Omega$-background;
\item compatibility with the Casimir energy of the
  five-dimensional theory.
\end{enumerate}
The factor $K^2$ is the dimension of $\mathfrak{gl}_K$.
\end{proof}

\begin{remark}[Consistency checks]
\label{rem:kappa-m2-checks}
The value $\kappa_{M2} = -K^2/12$ passes the following checks:

\emph{(i) Bosonic sign.}
All generators of $A_{M2,\infty}$ are bosonic.
The negative sign of $\kappa_{M2}$ matches the pattern for
bosonic algebras with an infinite tower of generators: the
$\beta\gamma$ system has $\kappa(\beta\gamma) = -2 = -K^2/12$
at $K^2 = 24$.  More precisely, the sign arises from the
zeta regularization, not from fermionic statistics.

\emph{(ii) $K = 1$ specialization.}
At $K = 1$, $\kappa_{M2} = -1/12$.  This matches the modular
characteristic of a \emph{single} $\beta\gamma$ pair in the
$z$-$\partial$ sector, which is the correct reduction: at $K=1$,
$\mathfrak{gl}_1 = \C$ is abelian, and the DDCA reduces to
$\C \otimes \Diff(\C)$.  The Casimir energy of this system
on $S^1$ is $-1/24$ from each of the two directions
($z$ and $\partial$), but the correct computation via the
Weyl trace gives $-1/12$ total (not $-1/24 \times 2 = -1/12$
from two independent oscillators, but from the SINGLE Weyl
algebra Casimir).  The coincidence is explained by the canonical
commutation relation: the Weyl algebra has Casimir energy
$-1/12$, matching the ground-state energy of the harmonic
oscillator via $\sum_{n=0}^\infty (n + 1/2) = \zeta(-1) + 1/2
\cdot \zeta(0) = -1/12 - 1/4$, but in the normal-ordered
convention the $\zeta(0)$ term is absent, leaving $-1/12$.

\emph{(iii) Large-$K$ scaling.}
$\kappa_{M2} \propto K^2$, which matches the large-$K$
scaling of all other invariants of $A_{M2,\infty}$: the
quantum correction~\eqref{eq:m2-quantum-dz} scales as $K^2$,
and the genus-$1$ partition function should scale as $K^2$
in the planar limit.

\emph{(iv) Comparison with naive approximations.}
The three failed candidates noted in
Remark~\ref{rem:kappa-m2-genuinely-open} of the extensions
chapter were $K^2$, $-K$, and $-1/2$.  Our value $-K^2/12$ is
none of these.  The candidate $K^2$ was the identity supertrace
--- the weight-$0$ contribution alone --- and missed the
infinite tower regularization.
\end{remark}

\subsection{The genus-$1$ partition function of $A_{M2,\infty}$}

From $\kappa_{M2} = -K^2/12$, the genus-$1$ free energy is
\begin{equation}\label{eq:m2-F1}
F_1(A_{M2,\infty})
= \frac{\kappa_{M2}}{24}
= -\frac{K^2}{288},
\end{equation}
and the genus-$1$ partition function (in the character
normalization) is
\begin{equation}\label{eq:m2-Z1}
Z_1(A_{M2,\infty};\, q)
= q^{-\kappa_{M2}/24}\, \prod_{n=1}^{\infty}
\frac{1}{(1-q^n)^{K^2(n+1)}}
= q^{K^2/288}\, \prod_{n=1}^{\infty}
\frac{1}{(1-q^n)^{K^2(n+1)}}.
\end{equation}
The exponent $K^2(n+1)$ in the infinite product counts the
$K^2(n+1)$ generators at conformal weight $n+1$: each
contributes one bosonic oscillator per mode.  The leading
$q$-power $K^2/288$ is the Casimir energy, consistent with
the zeta-regularized ground-state energy
$E_0 = -\kappa_{M2}/24 = K^2/288$.

\subsection{Physical identification: the $S^3$ partition function}

\begin{theorem}[Identification of $\kappa_{M2}$ with the 5d CS
Casimir energy]
\label{thm:kappa-m2-physical}
The modular characteristic $\kappa_{M2} = -K^2/12$ identifies,
through the HT bulk-boundary correspondence, with the one-loop
Casimir energy of the five-dimensional noncommutative
Chern--Simons theory on $S^1 \times \C^2$ in the
$\Omega$-background.  Specifically:
\begin{enumerate}[label=(\roman*)]
\item The genus-$1$ free energy
$F_1 = -K^2/288$ is the coefficient of
$\hbar^0$ in the free energy expansion of the 5d CS theory
on $S^1 \times \C^2$, after compactification on $S^1$.
\item Under the holographic dictionary, the genus-$g$ free
energy
\[
F_g(A_{M2,\infty})
= \kappa_{M2}\cdot \lambda_g^{\mathrm{FP}}
= -\frac{K^2}{12}\cdot
\frac{2^{2g-1}-1}{2^{2g-1}}\cdot
\frac{|B_{2g}|}{(2g)!}
\]
reproduces the perturbative expansion of the
five-dimensional partition function on the twisted
$S^1 \times \C^2$ background.
\item The $S^3$ partition function of the boundary theory
(obtained by a further compactification on $S^2$) is
related to $\kappa_{M2}$ through the KK reduction:
the boundary theory on $S^3$ has partition function
$Z_{S^3} = e^{-F_{\mathrm{pert}}}$ where
$F_{\mathrm{pert}} \sim K^2/12 \cdot \log(r/\ell)$
with $r$ the $S^3$ radius and $\ell$ the UV cutoff.
This is NOT the Witten--Reshetikhin--Turaev invariant:
the DDCA is an infinite-rank algebra, not a finite-rank
quantum group, and its $S^3$ partition function is a
perturbative (non-topological) invariant.
\end{enumerate}
\end{theorem}

The distinction from WRT is important.  For finite-rank chiral
algebras $\widehat{\fg}_k$, the modular characteristic
$\kappa = \dim(\fg)(k + h^\vee)/(2h^\vee)$ enters the WRT
invariant through the relation
$Z_{\mathrm{WRT}}(S^3;\, \fg,k) \sim e^{-\pi i \kappa / (k+h^\vee)}$,
because the Chern--Simons level is directly related to~$\kappa$.
For the DDCA, there is no finite Chern--Simons level, and the
partition function is perturbative rather than topological.
The invariant $\kappa_{M2}$ measures the \emph{one-loop}
contribution to the five-dimensional path integral, not a
topological invariant of a three-manifold.

\subsection{The cyclic pairing for $\widehat{\fsl}_2$ at
$k = 1, 2, 3$}

As a check and comparison, we record the cyclic pairing
computation for the affine lineage at low levels.

\begin{center}
\renewcommand{\arraystretch}{1.3}
\begin{tabular}{lllll}
$k$ & $\kappa = 3(k+2)/4$ & $F_1 = \kappa/24$ &
$Z_1(q)$ & $Z_{\mathrm{WRT}}(S^3)$ \\
\hline
$1$ & $9/4$ & $3/32$ &
$\chi_1(q) = q^{-1/24}\prod(1-q^n)^{-3}$ &
$\sqrt{2/3}\sin(\pi/3)$ \\
$2$ & $3$ & $1/8$ &
$\chi_2(q) = q^{-1/8}\prod(1-q^n)^{-3}$ &
$1/2$ \\
$3$ & $15/4$ & $5/32$ &
$\chi_3(q)$ &
$\sqrt{2/5}\sin(\pi/5)$
\end{tabular}
\end{center}
The WRT values use the standard normalization
$Z_{\mathrm{WRT}}(S^3;\, \mathfrak{sl}_2, k)
= \sqrt{2/(k+2)}\,\sin(\pi/(k+2))$.
The cyclic pairing enters through the Chern--Simons level
$k_{\mathrm{CS}} = k + h^\vee = k + 2$, and
$\kappa = 3(k+2)/4 = 3k_{\mathrm{CS}}/4$.  The WRT
invariant is a function of $k_{\mathrm{CS}}$ alone, while
$\kappa$ depends linearly on~$k_{\mathrm{CS}}$.  The
relation is:
\begin{equation}\label{eq:kappa-wrt-sl2}
|Z_{\mathrm{WRT}}(S^3)|^2
= \frac{2}{k_{\mathrm{CS}}}\,\sin^2\!\Bigl(\frac{\pi}{k_{\mathrm{CS}}}\Bigr)
\;\sim\;
\frac{2\pi^2}{k_{\mathrm{CS}}^3}
\quad\text{as } k \to \infty,
\end{equation}
which falls as $k_{\mathrm{CS}}^{-3} \sim \kappa^{-3}$.  This
inverse-cubic dependence on $\kappa$ is a universal feature of
the $S^3$ partition function in Chern--Simons theory: the
perturbative expansion is
$\log Z_{S^3} \sim -\frac{3}{2}\log k_{\mathrm{CS}}
+ O(1/k_{\mathrm{CS}})$, and $\kappa \propto k_{\mathrm{CS}}$.

\subsection{The Virasoro cyclic pairing at central charge $c$}

For the Virasoro algebra at central charge $c$:
\[
\kappa(\mathrm{Vir}_c) = c/2, \qquad
F_1 = c/48, \qquad
Z_1(q) = q^{-(c-1)/24}\prod_{n=1}^\infty (1-q^n)^{-1}.
\]
The self-dual point $c = 13$ (where $\mathrm{Vir}_c^! =
\mathrm{Vir}_{26-c}$ self-dualizes) has
$\kappa(13) = 13/2$ and $\kappa(13) + \kappa^!(13) =
13/2 + 13/2 = 13 \neq 0$ (AP24).  The anomaly-free point
$c = 26$ has $\kappa(26) = 13$: the full ghost-matter system at
$c = 26$ has $\kappa_{\mathrm{eff}} = \kappa_{\mathrm{matter}} +
\kappa_{\mathrm{ghost}} = 13 + (-13) = 0$.

The physical content of $\kappa(\mathrm{Vir}_c)$ in the 3d HT
context is the \emph{gravitational anomaly coefficient}: it
measures the failure of the genus-$1$ boundary partition function
to be modular.  Under $\tau \mapsto \tau + 1$:
$Z_1(q) \mapsto e^{2\pi i c/24}\, Z_1(q)$, and the phase
$e^{2\pi i c/24}$ is directly $e^{2\pi i \kappa/12}$.

\subsection{Summary and the open frontier}

The computation $\kappa_{M2} = -K^2/12$ via spectral zeta
regularization narrows the extensions chapter status to
\texttt{ClaimStatusHeuristic}: the numerical value is
well-defined once the regularization scheme is chosen,
but the proof that zeta regularization is the correct
prescription for the DDCA from first principles remains open.
The invariant has a clean physical interpretation as the
one-loop Casimir energy of the five-dimensional noncommutative
Chern--Simons theory.

What remains genuinely open:
\begin{itemize}[nosep]
\item Verification from the 11d supergravity side: does the
  anomaly cancellation in M-theory predict $\kappa_{M2} = -K^2/12$
  through the tadpole cancellation mechanism?
\item The genus-$2$ and higher-genus free energies
  $F_g = -K^2/12 \cdot \lambda_g^{\mathrm{FP}}$ should be checked
  against the five-dimensional perturbative expansion.
\item The non-scalar corrections (from the cubic shadow
  $\mathfrak{C}$ and quartic contact $\mathfrak{Q}$ of
  $A_{M2,\infty}$) are determined by the full Weyl algebra
  bracket structure and have not been computed.
\item For Kac--Moody algebras with root multiplicities $> 1$
  (the Monster Lie algebra, Borcherds algebras), the analogous
  computation involves the Peterson recursion for root
  multiplicities and should produce kappa values related to
  the coefficients of the $j$-function.
\end{itemize}


% =================================================================
\section{The Virasoro $\Ainf$ pattern: toward a closed-form formula for $m_k$}
\label{sec:virasoro-ainf-pattern}
% =================================================================

The Virasoro algebra has a genuinely infinite $\Ainf$ tower
$\{m_k\}_{k \ge 2}$, with every $m_k$ nonzero for generic~$c$.
The quartic OPE pole ($\lambda^3$ in the $\lambda$-bracket,
$(z-w)^{-4}$ in the OPE) is the engine: it produces the
non-associativity that forces $m_3$, which forces $m_4$, and so on.
The goal of this section is to identify the algebraic pattern
governing the entire tower and to determine whether a closed-form
formula or generating function exists.

\subsection*{Explicit formulas for $m_2$, $m_3$, and the arity-$4$ obstruction}

The input datum is the Virasoro $\lambda$-bracket:
\begin{equation}\label{eq:vir-m2-pattern}
m_2(T,T;\,\lambda)
\;=\;
\partial T \;+\; 2T\lambda \;+\; \frac{c}{12}\,\lambda^3.
\end{equation}
Three terms: $\partial T$ from the simple pole, $2T\lambda$ from
the double pole, $(c/12)\lambda^3$ from the quartic pole.

The ternary operation, forced by the non-associativity of $m_2$
(the quartic pole creates a nonzero associator), is:
\begin{equation}\label{eq:vir-m3-pattern}
m_3(T,T,T;\,\lambda_1,\lambda_2)
\;=\;
\partial^2 T
\;+\; (2\lambda_1 + 3\lambda_2)\,\partial T
\;+\; 2\lambda_2(2\lambda_1 + \lambda_2)\,T
\;+\; \frac{c}{12}\,\lambda_2^3(2\lambda_1 + \lambda_2).
\end{equation}
This is the chain-level formula, where $m_3 = -A_3$ compensates the
associator $A_3(\lambda_1,\lambda_2)
= -\{T_{\lambda_2}\{T_{\lambda_1} T\}\}$.
On cohomology (via the transferred operation $m_3^H$ from the BV-BRST
homotopy), the formula is:
\begin{equation}\label{eq:vir-m3H-pattern}
m_3^H(T,T,T;\,\lambda_1,\lambda_2)
\;=\;
\frac{c}{6}\bigl(\lambda_1^2\lambda_2
+ \lambda_1\lambda_2^2\bigr)
\;+\; 4T\,\lambda_1\lambda_2
\;+\; 2(\partial T)(\lambda_1 - \lambda_2).
\end{equation}
The two formulas differ by a homotopy (gauge) equivalence:
the chain-level~\eqref{eq:vir-m3-pattern} and the
transferred~\eqref{eq:vir-m3H-pattern} are the same operation
represented in different bases.

At arity~$4$, the chain-level Stasheff identity
$d\,m_4 = -\bigl(\text{sum of } m_2 \circ m_3
\text{ and } m_3 \circ m_2 \text{ compositions}\bigr)$
determines $m_4$ via the BRST contracting homotopy~$h$.
The five composition terms in the Stasheff equation are:
\begin{alignat}{2}\label{eq:stasheff-arity4}
S_1 &= +m_3\bigl(m_2(T,T;\lambda_1),\, T,\, T;\;
  \lambda_1{+}\lambda_2,\, \lambda_3\bigr), \notag\\
S_2 &= -m_3\bigl(T,\, m_2(T,T;\lambda_2),\, T;\;
  \lambda_1,\, \lambda_2{+}\lambda_3\bigr), \notag\\
S_3 &= +m_3\bigl(T,\, T,\, m_2(T,T;\lambda_3);\;
  \lambda_1,\, \lambda_2\bigr), \\
S_4 &= -m_2\bigl(m_3(T,T,T;\lambda_1,\lambda_2),\, T;\;
  \lambda_1{+}\lambda_2{+}\lambda_3\bigr), \notag\\
S_5 &= -m_2\bigl(T,\, m_3(T,T,T;\lambda_2,\lambda_3);\;
  \lambda_1\bigr). \notag
\end{alignat}
The slot substitutions use the sesquilinearity rules of the
$\lambda$-bracket: left-sesquilinearity
$m_k(\partial a, \ldots;\, \ell, \ldots) =
-\ell\, m_k(a, \ldots;\, \ell, \ldots)$
for the first slot, and right-sesquilinearity
$m_k(\ldots, \partial a;\, \ldots, \ell) =
(\ell + \partial)\,m_k(\ldots, a;\, \ldots, \ell)$ for the last.
The middle-slot sesquilinearity is more involved and requires
direct computation from the associator formula.

\begin{proposition}[Scalar part of the arity-$4$ Stasheff obstruction]
\label{comp:stasheff-4-scalar}
Computing the five terms $S_1, \ldots, S_5$ and extracting
the scalar (central-charge-dependent) part gives:
\begin{equation}\label{eq:stasheff-4-scalar}
\bigl(S_1 + S_2 + S_3 + S_4 + S_5\bigr)\big|_{\text{scalar}}
\;=\;
\frac{c}{12}\bigl(
-2\lambda_1^2\lambda_2^3
- \lambda_1\lambda_2^4
- 2\lambda_1\lambda_2\lambda_3^3
+ 2\lambda_1\lambda_3^4
+ 2\lambda_2^2\lambda_3^3
- 2\lambda_2\lambda_3^4\bigr).
\end{equation}
This is a homogeneous polynomial of degree~$5$ in
$(\lambda_1,\lambda_2,\lambda_3)$ and is \textbf{linear} in~$c$.
Its non-vanishing for generic~$c$ confirms that
$m_4 \ne 0$; the compensating $m_4$ is obtained by applying the
BRST homotopy $h$ to the right-hand
side~\eqref{eq:stasheff-4-scalar}, introducing the Shapovalov
inverse and hence the Kac determinant $c(5c+22)$ as denominator.
\end{proposition}

\subsection*{The degree structure}

From the explicit formulas, a rigid pattern emerges for the
polynomial degrees of the coefficients of $m_k$:

\begin{center}
\small
\renewcommand{\arraystretch}{1.15}
\begin{tabular}{@{}ccccc@{}}
$k$ & $\partial^{k-2}T$ degree &
  $\partial T$ degree & $T$ degree & scalar degree \\
\hline
$2$ & $0$ & $0$ & $1$ & $3$ \\
$3$ & $0$ & $1$ & $2$ & $4$ \\
$4$ & $0$ & $2$ & $3$ & $5$ \\
$k$ & $0$ & $k-2$ & $k-1$ & $k+1$
\end{tabular}
\end{center}

\noindent
The \textbf{scalar part} of $m_k$ has total $\lambda$-degree
$k+1$; the $T$-coefficient has degree $k-1$; each derivative
$\partial^j T$ has degree $k-1-j$; the highest derivative
$\partial^{k-2}T$ appears with coefficient~$1$
(degree~$0$, independent of the spectral parameters).  This
pattern is forced by the conformal weight balance: each input $T$
contributes weight~$2$, each spectral parameter $\lambda_i$
contributes weight~$1$, and the homotopy~$h$ has weight~$-1$.

\subsection*{The $c$-linearity theorem}

The most striking feature of the Virasoro $\Ainf$ tower is
the $c$-dependence of the scalar part.

\begin{theorem}[$c$-linearity of chain-level $m_k$; proved]\label{thm:c-linearity-wn}
The scalar part of the chain-level $m_k(T^{\otimes k};\,
\lambda_1, \ldots, \lambda_{k-1})$ is \textbf{linear} in~$c$
for all $k \ge 2$:
\begin{equation}\label{eq:c-linearity-wn}
m_k\big|_{\text{scalar}}
\;=\;
\frac{c}{12}\,P_k(\lambda_1, \ldots, \lambda_{k-1}),
\end{equation}
where $P_k$ is a homogeneous polynomial of degree $k+1$.
\end{theorem}

\begin{proof}
By induction on~$k$.  The base case $k = 2$ gives
$m_2|_{\text{scalar}} = (c/12)\lambda^3$, which is
linear in~$c$.

For the inductive step: the chain-level Stasheff relation
expresses $d\,m_k$ as a sum of compositions $m_i \circ m_j$
with $i + j = k + 1$, $i, j \ge 2$.  The scalar part of each
composition involves composing a field-valued output (from $m_j$)
into $m_i$.  Since $m_i(\text{scalar}, T; \ldots) = 0$ and
$m_i(T, \text{scalar}; \ldots) = 0$ (scalars are annihilated
by the $\lambda$-bracket in both slots), the only scalar
contributions come from the $T$-coefficient of one factor
composed with the $\lambda$-bracket of the other, plus the
scalar coefficient of one factor composed with the $T$-coefficient
of the other.  Both channels produce terms linear in~$c$ (the
scalar coefficients of $m_j$ are linear in~$c$ by the inductive
hypothesis, and the $T$-coefficients are $c$-independent).

The BRST homotopy $h$ is a $c$-independent operator on the
Virasoro Verma module (it depends on the ghost system, not
on the central charge).  Applying $h$ to the Stasheff
obstruction preserves the $c$-linearity.
\end{proof}

\begin{remark}[Chain level vs.\ cohomological level]
The $c$-linearity holds at the \emph{chain} level.  The
transferred operation $m_k^H$ on BRST cohomology involves
the Shapovalov inverse, which introduces denominators from the
Kac determinant.  Thus $m_k^H$ is a \emph{rational} function
of~$c$:
\[
m_k^H\big|_{\text{scalar}}
\;=\;
\frac{c \cdot P_k(\lambda_1, \ldots, \lambda_{k-1})}
     {D_k(c)},
\]
where $D_k(c)$ is a product of Kac determinant factors at
levels $\le k$.  The denominators are the loci where null
vectors appear in the intermediate Virasoro Verma module,
reorganising the $\Ainf$ tower.  For instance,
$D_4(c) = c(5c + 22)$ (the level-$4$ Kac determinant at $h = 0$)
controls the quartic contact invariant
$\mathfrak{Q}^{\mathrm{contact}}_{\mathrm{Vir}} = 10/[c(5c+22)]$.
\end{remark}

\subsection*{The resolvent generating function}

The wheel-diagram structure of the Virasoro $\Ainf$ operations
(Proposition~\ref{prop:virasoro-wheel-diagrams} of the main text)
gives the $\Ainf$ tower a natural generating function: a
\textbf{resolvent trace}.

\begin{definition}[Virasoro resolvent operator]
\label{def:vir-resolvent-wn}
Let $V$ denote the vacuum Virasoro Verma module, viewed as a
module for the BRST complex of the
$\widehat{\mathfrak{sl}}_2$ DS reduction.
Define:
\begin{enumerate}[label=\textup{(\roman*)},nosep]
\item The \emph{bracket operator}
  $R(\lambda) = m_2(-, T;\, \lambda) \colon V \to V$,
  the left action of the $\lambda$-bracket.
\item The \emph{propagator operator}
  $K(\lambda) = R(\lambda) \circ h \colon V \to V$,
  where $h$ is the BRST contracting homotopy.
\item The \emph{resolvent generating function}:
\begin{equation}\label{eq:resolvent-gf-wn}
\mathcal{G}(t;\, \lambda)
\;=\;
\sum_{k=0}^{\infty} t^k\,
\Tr_h\!\bigl(\underbrace{K(\lambda) \cdots K(\lambda)}_{k}\bigr)
\;=\;
\Tr_h\!\left(\frac{1}{1 - t\,K(\lambda)}\right).
\end{equation}
\end{enumerate}
\end{definition}

The resolvent~\eqref{eq:resolvent-gf-wn} generates the
\emph{single-parameter} $\Ainf$ operations
$m_k(T, \ldots, T;\, \lambda, \ldots, \lambda)$
(all spectral parameters equal).  The full multi-parameter
operations are recovered by the polarization
\begin{equation}\label{eq:resolvent-polarization-wn}
m_k(T^{\otimes k};\, \lambda_1, \ldots, \lambda_{k-1})
\;=\;
p \circ K(\lambda_1) \circ K(\lambda_2) \circ \cdots
\circ K(\lambda_{k-1}) \circ i(T)
\;+\;
(\text{tree corrections}),
\end{equation}
where $i \colon H^0 \hookrightarrow V$ is the inclusion and
$p \colon V \twoheadrightarrow H^0$ the projection.  The tree
corrections come from the resolvent tree formula (Volume~I):
they vanish for $k = 2, 3$ and contribute at $k \ge 4$ as
compositions of lower~$m_j$.

\begin{observation}[The wheel is a trace]\label{obs:wheel-trace-wn}
Each wheel diagram at arity~$k$ is a cycle of $k$~propagators
$K(\lambda_j)$, with each vertex emitting one external $T$-leg.
The cycle structure is precisely the trace
$\Tr_h(K(\lambda_1) \cdots K(\lambda_{k-1}))$
in the Virasoro Verma module, where $\Tr_h$ denotes the
homotopy-weighted trace (the trace over the BRST-acyclic part,
evaluated using the Shapovalov inner product).
\end{observation}

\begin{observation}[Graviton self-energy]\label{obs:self-energy-wn}
In the 3d HT QFT, the generating function $\mathcal{G}(t; \lambda)$
is the \textbf{graviton self-energy}: the one-particle-irreducible
two-point function of the graviton propagating in the bulk
$\C \times \R$.  The resolvent structure
$\Tr(1/(1 - t\,K))$ is the geometric series resummation of the
infinite chain of self-energy insertions.  The poles of the
resolvent in the variable~$t$ are the graviton ``bound-state
energies'' --- the eigenvalues of the propagator operator
$K(\lambda)$ on the Virasoro Verma module, weighted by the
Shapovalov form.
\end{observation}

\subsection*{Kac determinant poles as resonances}

The resolvent $\mathcal{G}(t; \lambda)$ has poles in the central
charge~$c$ at the zeros of the Kac determinant.  At each
such zero, a null vector appears in the intermediate Virasoro
Verma module, and the homotopy~$h$ (which inverts the Shapovalov
form on acyclic states) develops a pole.

\begin{center}
\small
\renewcommand{\arraystretch}{1.15}
\begin{tabular}{@{}clll@{}}
Level & Kac zero & Physical interpretation
  & Affected $m_k$ \\
\hline
$2$ & $c = 0$ & zero Newton constant ($G = \infty$)
  & $m_k$ for all $k \ge 3$ \\
$4$ & $c = -22/5$ & Lee--Yang edge singularity
  & $m_k$ for $k \ge 4$ \\
$6$ & $c = -68/7$ & next minimal model locus
  & $m_k$ for $k \ge 6$ \\
$n$ & Kac zeros at level~$n$ & & $m_k$ for $k \ge n$
\end{tabular}
\end{center}

\noindent
The full denominator of $m_k^H$ is the product of Kac
determinant factors at levels $\le k$:
\begin{equation}\label{eq:mk-denominator-wn}
D_k(c) \;=\;
\prod_{n=2}^{k}\;
\prod_{\substack{r,s \ge 1 \\ rs \le n}}
\bigl(h_{r,s}(c)\bigr)^{p(n - rs)}\bigg|_{h=0},
\end{equation}
where $h_{r,s}(c) = [(m(r^2-1) - n(s^2-1))^2 - (m-n)^2]/(4mn)$
with $c = 1 - 6(m-n)^2/(mn)$, and $p(j)$ is the
number of partitions of~$j$.

\subsection*{Finite presentation of an infinite tower}

\begin{theorem}[Finite presentation; proved]
\label{thm:finite-presentation-wn}
The Virasoro $\Ainf$ structure $\{m_k\}_{k \ge 2}$ is finitely
presented in the following sense: every $m_k$ is uniquely determined
by the two data
\begin{enumerate}[label=\textup{(\roman*)},nosep]
\item the binary operation $m_2(T,T;\lambda)
  = \partial T + 2T\lambda + (c/12)\lambda^3$
  \textup{(}the $\lambda$-bracket\textup{)}, and
\item the BRST contracting homotopy $h$ from the
  Drinfeld--Sokolov reduction of $\widehat{\mathfrak{sl}}_2$ at
  level~$k$.
\end{enumerate}
The determination is by the resolvent tree formula:
$m_k$ is the sum over Catalan-many planar binary trees with $k$
leaves, internal edges decorated by~$h$, and internal vertices
by the BRST perturbation~$\delta$.

The $\Ainf$ structure does not truncate ($m_k \ne 0$ for all
$k \ge 3$ and generic~$c$), so it is not finitely generated in
the operadic sense.  The finite presentation is not a finiteness
of the output but a finiteness of the generating data: one
binary operation and one homotopy produce the entire infinite tower.
\end{theorem}

\begin{proof}
The resolvent tree formula of Volume~I
(Construction~\ref*{V1-constr:htt-alg}) expresses the transferred
$m_k^H$ as a sum over planar binary trees.  The ingredients are
$m_2^{\mathrm{aff}}$ (the affine $\lambda$-bracket at internal
vertices), $h$ (at internal edges), $i$ (at leaves), and $p$
(at the root).  Since $m_2^{\mathrm{aff}}$ is determined by the
Sugawara construction from $m_2^{\mathrm{Vir}}$ and the DS
constraint, both (i) and (ii) suffice.  Non-truncation follows
from the non-vanishing of the wheel integral at every arity
(Proposition~\ref{prop:virasoro-wheel-diagrams}).
\end{proof}

\subsection*{The polynomial pattern and the question of a closed form}

Combining the degree structure, $c$-linearity, and resolvent
representation, we can state the pattern precisely.

\begin{conjecture}[Closed-form $m_k$]\label{conj:closed-form-wn}
The chain-level Virasoro $\Ainf$ operation on $T^{\otimes k}$ has
the form
\begin{equation}\label{eq:mk-pattern-wn}
m_k(T^{\otimes k};\,\lambda_1, \ldots, \lambda_{k-1})
\;=\;
\sum_{j=0}^{k-2}
\partial^j T \cdot
f_j^{(k)}(\lambda_1, \ldots, \lambda_{k-1})
\;+\;
\frac{c}{12}\,P_k(\lambda_1, \ldots, \lambda_{k-1}),
\end{equation}
where:
\begin{enumerate}[label=\textup{(\alph*)},nosep]
\item $f_j^{(k)}$ is a polynomial of degree $k-1-j$ in the
  $\lambda_i$, independent of~$c$;
\item $f_{k-2}^{(k)} \equiv 1$
  \textup{(}the top derivative $\partial^{k-2}T$ has
  coefficient~$1$\textup{)};
\item $P_k$ is a homogeneous polynomial of degree $k+1$;
\item the polynomials satisfy the recursion determined by the
  Stasheff identity at arity~$k$, encoding the wheel integral
  on $\FM_k(\C) \times \Conf_k^{<}(\R)$.
\end{enumerate}

The problem of finding a closed-form formula for $f_j^{(k)}$ and
$P_k$ reduces to diagonalizing the propagator operator
$K(\lambda) = m_2(-, T;\, \lambda) \circ h$
on the vacuum Virasoro Verma module.  This is equivalent to
computing the spectral decomposition of the Virasoro
Laplace kernel $r^L(z) = (c/2)/z^4 + 2T/z^2 + \partial T/z$
acting on the Verma module via the Shapovalov-weighted adjoint
action.
\end{conjecture}

\begin{remark}[The obstruction to a fully explicit formula]
\label{rem:obstruction-to-closed-form-wn}
A fully explicit closed-form formula for $m_k$ would require
explicit knowledge of the Shapovalov inner product at all levels.
For the Virasoro algebra, the Shapovalov matrix at level~$n$ is
known in principle (it is the Gram matrix of the Verma module), but
its inverse grows combinatorially complex.  The Kac determinant
$\det S_n = \prod_{rs \le n} (h_{r,s})^{p(n-rs)}$ controls the
denominators; the numerators involve the full partition
combinatorics of the Verma module.  The generating function
$\mathcal{G}(t;\lambda) = \Tr_h(1/(1-tK))$ packages this information
as a single operator-valued resolvent, but its explicit expansion
as a function of~$c$ requires diagonalizing $K(\lambda)$ on
each weight space --- a problem equivalent to computing the
full set of Virasoro singular vectors.

The existence of an explicit formula is therefore controlled by
the same question that governs the closed-form structure of
Virasoro representation theory: whether the Shapovalov matrix
admits a factorization beyond the Kac determinant formula.
For the scalar part (the $c$-dependent term), the $c$-linearity
theorem reduces the problem to a single polynomial $P_k$ at each
arity, and this polynomial is determined recursively by the
Stasheff identity.  The challenge is to solve this recursion in
closed form.
\end{remark}

\begin{remark}[Physical content of the generating function]
\label{rem:physical-generating-function-wn}
The resolvent generating function $\mathcal{G}(t;\lambda)$
is the \textbf{graviton spectral curve} of the 3d HT theory.
Its poles in~$t$ (the eigenvalues of $K(\lambda)$) are the
graviton bound-state energies; its poles in~$c$ (the Kac
determinant zeros) are the resonances of the graviton
self-energy.

In the heavy-heavy limit $c \to \infty$, the resolvent
simplifies: the Shapovalov form becomes diagonal (the Verma
module has no null vectors), and $K(\lambda)$ acts by a simple
multiplication operator.  The resulting generating function
reproduces the classical graviton scattering amplitude in
AdS$_3$.  The Kac determinant poles are quantum corrections
(null-vector resonances) that disappear in the classical
$c \to \infty$ limit.

At the self-dual point $c = 13$, where
$\mathrm{Vir}_c^! = \mathrm{Vir}_c$, the resolvent has a
special symmetry: $\mathcal{G}(t;\lambda)|_{c=13}$ is invariant under
the Koszul involution that exchanges the boundary and the Koszul
dual.  This is the algebraic expression of the self-duality of
3d gravity at $c = 13$.
\end{remark}

\subsection*{Summary of results}

\begin{enumerate}
\item The scalar part of $m_k$ is \textbf{linear in~$c$}
  at the chain level, with the coefficient a homogeneous
  polynomial $P_k$ of degree $k+1$ in
  $(\lambda_1, \ldots, \lambda_{k-1})$.
\item The field-valued part ($T$, $\partial T$, $\ldots$,
  $\partial^{k-2}T$) is \textbf{independent of~$c$}.
  The coefficient of $\partial^j T$ has degree $k-1-j$.
  The top derivative $\partial^{k-2}T$ has coefficient~$1$.
\item The full $\Ainf$ structure is \textbf{finitely presented}
  by the binary $\lambda$-bracket $m_2$ and the BRST homotopy~$h$,
  via the resolvent tree formula.
\item The generating function is the \textbf{resolvent}
  $\Tr_h(1/(1-tK))$ of the propagator operator $K = m_2(-,T;\lambda) \circ h$
  on the vacuum Virasoro Verma module.
\item On cohomology, $m_k^H$ acquires \textbf{Kac determinant
  denominators} from the Shapovalov inverse in the homotopy.
  The loci $c = 0$, $c = -22/5$, $c = -68/7$, $\ldots$ are
  \textbf{resonances} of the graviton self-energy.
\item The A-infinity structure does \textbf{not} truncate and
  does \textbf{not} have a finite operadic presentation.
  Every $m_k \ne 0$ for generic~$c$, and the quartic
  OPE pole is the single algebraic obstruction to formality.
\item A fully explicit closed-form formula for $m_k$ at all
  arities requires diagonalising the Shapovalov form at all
  levels, which is equivalent to the full structure theory of
  Virasoro singular vectors.  The \textbf{scalar part} admits a
  recursive determination via the Stasheff identity, with each
  step producing a homogeneous polynomial.
\end{enumerate}

\begin{question}[Explicit $P_k$]\label{q:explicit-pk-wn}
Is there a closed-form formula for the scalar polynomial
$P_k(\lambda_1, \ldots, \lambda_{k-1})$?  The first three are:
\begin{align*}
P_2(\lambda_1) &= \lambda_1^3, \\
P_3(\lambda_1,\lambda_2) &= \lambda_2^3(2\lambda_1 + \lambda_2), \\
P_4(\lambda_1,\lambda_2,\lambda_3) &= \;?\quad
  \text{\textup{(}}
  \text{determined by~\eqref{eq:stasheff-4-scalar} and the homotopy~$h$}
  \text{\textup{)}}.
\end{align*}
At the cumulative-parameter level $s_j = \lambda_1 + \cdots + \lambda_j$,
$P_3 = (s_2 - s_1)^3(s_1 + s_2)$.  Is there a pattern in the
$(s_1, \ldots, s_{k-1})$ variables?
\end{question}

\begin{question}[Generating function as spectral curve]
\label{q:spectral-curve-wn}
Does the resolvent generating function
$\mathcal{G}(t;\lambda) = \Tr_h(1/(1-tK(\lambda)))$ admit a
representation as a spectral curve --- a zero locus
$F(t, \lambda, c) = 0$ of a polynomial or analytic function?
Such a representation would give a genuinely closed-form
characterisation of the entire Virasoro $\Ainf$ tower,
replacing the recursive Stasheff determination by a single
algebraic equation.
\end{question}

%====================================================================
% SESSION: 2 April 2026 -- dg-shifted Yangian computational frontier
%====================================================================

\section{Concrete presentations of dg-shifted Yangians from the ordered bar complex}

\subsection{The sl$_2$ Yangian: complete presentation from the bar complex}

The input is $\widehat{\mathfrak{sl}}_2$ at level~$k$.  The ordered bar complex
$\bar{B}^{\mathrm{ord}}(V_k(\mathfrak{sl}_2))$ has:
\begin{itemize}[nosep]
\item \emph{Degree~$1$}: $\bar{B}^{\mathrm{ord}}_1 = \operatorname{span}\{s^{-1}e,\,s^{-1}f,\,s^{-1}h\}$.
  Differential $d = 0$ (no contractions on a single element).
\item \emph{Degree~$2$}: $9$ generators $s^{-1}x \otimes s^{-1}y$ for
  $x,y \in \{e,f,h\}$.  The bar differential extracts the zeroth
  product (Lie bracket):
  $d(s^{-1}e \otimes s^{-1}f) = \pm s^{-1}h$,
  $d(s^{-1}h \otimes s^{-1}e) = \pm 2\,s^{-1}e$, etc.
  In the \emph{ordered} complex, $s^{-1}e \otimes s^{-1}h$ and
  $s^{-1}h \otimes s^{-1}e$ are independent generators;
  the $R$-matrix is precisely the datum distinguishing these two
  orderings.
\item \emph{Degree~$3$}: Two codimension-$1$ faces $D_{1,2}$ and $D_{2,3}$;
  $d^2 = 0$ is the Yang--Baxter equation
  (Proposition~\ref{prop:ybe-from-d-squared}).
\end{itemize}
The open-colour Koszul dual is
$\cA^!_{\mathrm{line}} = Y_\hbar(\mathfrak{sl}_2)$ at
$\hbar = 1/(k+2)$.

\medskip\noindent\textbf{Collision residue.}\enspace
By AP19, the bar kernel absorbs one pole.  The collision residue
$r$-matrix for the affine $\widehat{\mathfrak{sl}}_2$ OPE is
\[
r(z) = \frac{\hbar\,\Omega}{z},
\qquad
\Omega = e \otimes f + f \otimes e + \tfrac12\,h \otimes h
\quad(\text{Casimir of } \mathfrak{sl}_2).
\]

\medskip\noindent\textbf{Explicit $R$-matrix in the fundamental.}\enspace
In the fundamental $V = \C^2$ with basis $v_+, v_-$, the
split Casimir $\Omega = P - \tfrac12 I$ where $P$ is the permutation.
The Yang $R$-matrix is
\[
R(u) \;=\; u\,I + P
\;=\;
\begin{pmatrix}
u+1 & 0 & 0 & 0 \\
0 & u & 1 & 0 \\
0 & 1 & u & 0 \\
0 & 0 & 0 & u+1
\end{pmatrix}
\]
in the basis $\{v_+ \otimes v_+,\; v_+ \otimes v_-,\;
v_- \otimes v_+,\; v_- \otimes v_-\}$.
Eigenvalues: $u+1$ on $\mathrm{Sym}^2 V$ (multiplicity~$3$);
$u-1$ on $\bigwedge^2 V$ (multiplicity~$1$).

\medskip\noindent\textbf{RTT presentation.}\enspace
The transfer matrix $T(u) = \bigl(\begin{smallmatrix}
A(u) & B(u) \\ C(u) & D(u) \end{smallmatrix}\bigr)$
satisfies $R_{12}(u-v)\,T_1(u)\,T_2(v) = T_2(v)\,T_1(u)\,R_{12}(u-v)$.
The $16$ component equations reduce to:
\begin{alignat*}{2}
&[A(u), A(v)] = 0, \qquad
&&[D(u), D(v)] = 0, \\
&[B(u), B(v)] = 0, \qquad
&&[C(u), C(v)] = 0, \\
&(u{-}v)[A(u), B(v)] = A(v)B(u) - A(u)B(v), \quad
&&(u{-}v)[A(u), C(v)] = C(v)A(u) - C(u)A(v), \\
&(u{-}v)[A(u), D(v)] = C(v)B(u) - C(u)B(v), \quad
&&(u{-}v)[B(u), C(v)] = D(v)A(u) - D(u)A(v).
\end{alignat*}
The remaining $8$ are consequences.

\medskip\noindent\textbf{Drinfeld generators via Gauss decomposition.}\enspace
$T(u) = F(u)\,K(u)\,E(u)$ with $F$ lower-triangular unipotent,
$K$ diagonal, $E$ upper-triangular unipotent.  The Drinfeld
currents are $H(u) = k_1(u)/k_2(u)$, $E(u)$, $F(u)$.

\medskip\noindent\textbf{Twisted coproduct.}\enspace
\begin{align*}
\Delta_z(H(u)) &= H(u) \otimes H(u{-}z), \\
\Delta_z(E(u)) &= E(u) \otimes 1 + H(u{-}z) \otimes E(u{-}z), \\
\Delta_z(F(u)) &= 1 \otimes F(u) + F(u{-}z) \otimes H(u{-}z)^{-1}.
\end{align*}

\medskip\noindent\textbf{Evaluation modules.}\enspace
For $a \in \C$, the evaluation module $V_a$ is
$V = \C^2$ with $T(u) \mapsto 1 + \hbar\,\rho(\Omega)/(u-a)$.
Higher modes: $X_r \mapsto a^r \rho(X)$.
On $V_a \otimes V_b$: the $R$-matrix is $R(a{-}b) = (a{-}b)\,I + P$.

\medskip\noindent\textbf{PBW basis.}\enspace
Ordered monomials $\{f_{r_1}\cdots f_{r_a}\,h_{s_1}\cdots h_{s_b}\,
e_{t_1}\cdots e_{t_c}\}$ with increasing indices.
Hilbert series: $\prod_{d \ge 1}(1-q^d)^{-3}$.


\subsection{The sl$_3$ Yangian: adjacent roots, triangles, and the coefficient $\frac12$}

The rank-$2$ case reveals adjacent-root structure.  In $\mathfrak{gl}_3$
matrix units: $e_1 = E_{12}$, $e_2 = E_{23}$, $e_{12} = E_{13}$,
$f_1 = E_{21}$, $f_2 = E_{32}$, $f_{12} = E_{31}$.

\medskip\noindent\textbf{Triangle sectors.}\enspace
$T^L_{12} = E_{13}^{(1)}\,E_{21}^{(2)}\,E_{32}^{(3)}$,
$T^L_{21} = E_{13}^{(1)}\,E_{32}^{(2)}\,E_{21}^{(3)}$;
$T^R_{12} = E_{12}^{(1)}\,E_{23}^{(2)}\,E_{31}^{(3)}$,
$T^R_{21} = E_{23}^{(1)}\,E_{12}^{(2)}\,E_{31}^{(3)}$.

\medskip\noindent\textbf{Rigidity.}\enspace
$\lambda^L = \lambda^R = c_1 c_2/(2c_{12}) = 1/2 = \int_0^1(1-t)\,dt$.

\medskip\noindent\textbf{Serre.}\enspace
$2\alpha_1 + \alpha_2 \notin \Phi(\mathfrak{sl}_3) \Rightarrow
(\mathfrak{sl}_3)_{2\alpha_1+\alpha_2} = 0 \Rightarrow
[e_1,[e_1,e_2]] = 0$.


\subsection{The root quadrilateral at filtration~$3$: the coefficient $\frac13$}

For $A_3$ ($\mathfrak{sl}_4$), chain $i \sim j \sim k$ with
$[E_i, E_k] = 0$.  The BCH: $\mathrm{mult}_{XYZ}\log(e^X e^Y e^Z) =
\frac13[X,[Y,Z]]$ when $[X,Z] = 0$.  Direct verification:
$\frac14 + \frac{1}{12} = \frac13$ from the general formula
$\frac14[[X,Y],Z] + \frac{1}{12}[X,[Y,Z]] + \frac{1}{12}[Y,[X,Z]]$
with $[X,Z] = 0$ and Jacobi.  The $n=3$ beta integral:
$\int_0^1(1-t)^2\,dt = 1/3$.

Root-space one-dimensionality:
$\dim(\mathfrak{sl}_4)_{\alpha_1+\alpha_2+\alpha_3} = 1$.
All multilinear Lie monomials proportional to $E_{ijk}$.


\subsection{The gauge-gravity complexity dichotomy}

\[
\begin{array}{c|cc}
& \Ainf\text{ products } m_k & \text{Coproduct } \Delta_z \\
\hline
\text{Gauge (CS)} & m_k = 0 \;(k \ge 3) &
  \text{nontrivial Yangian} \\
\text{Gravity (Vir)} & m_k \neq 0 \;\forall\, k &
  \text{strictly primitive}
\end{array}
\]
DS transports $\mathbf{L} \to \mathbf{M}$: preserves Koszulness,
destroys formality.  The resolvent tree formula
$m_k^{\cW} = \sum_{\text{Catalan}}(m_2^{\mathrm{aff}} \circ h_{\mathrm{DS}})$
generates the infinite tower.


\subsection{The free-field nilpotent kernel}

$\Theta = \beta(x) \otimes \gamma(y) - \gamma(x) \otimes \beta(y)$,
$\Theta^2 = 0$.
Transport $T(u,u_0) = 1 + \Theta\,\log\frac{u+x-y}{u_0+x-y}$,
monodromy $M = 1 + 2\pi i\,\Theta$.
Contrast: Heisenberg $R = \exp(k/z)$ (full exponential);
$\beta\gamma$: $R = 1 + (\text{linear})$ (nilpotent truncation).
Wakimoto: $Y_{\mathrm{dg}}(\fg) = H^*(Y_{\mathrm{free}}, Q_{\mathrm{DS}})$.


\subsection{The curvature-braiding dichotomy at genus~$1$}

Coisson bracket $c_0 := \{a{}_{(0)}b\}$ controls decoupling.
$c_0 = 0$: quasi-periodic elliptic $r$-matrix with exact
Cartan gauge-twist defect; structures decouple (Heisenberg,
lattice).
$c_0 \neq 0$: quasi-periodic, defect $2\eta_\tau$ couples
curvature to braiding (affine KM, Virasoro).
Genus-$1$ $R$-matrix for Virasoro:
$r_1(z;\tau) = (c/2)(\wp(z;\tau) + E_2(\tau)/12) + 2T\,E_2(\tau)$.


\subsection{The three-language spine}

\[
\begin{array}{c|c|c|c}
\text{Datum} & \text{Physical (DNP25)} & \text{Spectral (Latyntsev)}
  & \text{Geometric (COZZ26)} \\
\hline
R\text{-matrix} & \text{OPE exchange} & \text{fact.\ iso} &
  \text{stable envelope} \\
\text{YBE} & \text{Ward identity} & \text{fact.\ axiom} &
  \text{braid relation} \\
\text{Strict.} & \text{root mult.\ $1$} &
  \text{by constr.} & \text{by constr.}
\end{array}
\]
Common object: $\barB^{\mathrm{ord}}(\cA)$ on
$\FM_k(\C) \times \Conf_k^<(\R)$.


\subsection{Three Koszul duality operations and the origin of the spectral parameter}

The SC${}^{\mathrm{ch,top}}$ operad has two colours, producing
\textbf{three} distinct Koszul duality operations on a chiral
algebra $\cA$:
\[
\begin{array}{c|c|c|c}
\text{Operation} & \text{Bar complex} & \text{Output} &
  \text{Example for } \widehat{\fg}_k \\
\hline
\text{Chiral KD} & B^{\mathrm{ch}}(\cA)\text{ on }
  \mathrm{Conf}_n(X)/\Sigma_n & \text{chiral alg.\
  }\cA^!_{\mathrm{ch}} & \widehat{\fg}_{-k-2h^\vee} \\
\text{$E_1$-chiral KD} & \bar{B}^{\mathrm{ord}}(\cA)\text{ on }
  \mathrm{Conf}_n^{\mathrm{ord}}(X) & \text{SC-alg.\
  }\cA^!_{\mathrm{line}} & Y_\hbar(\fg) \\
\text{Pure $E_1$ KD} & T^c(s^{-1}\bar\cA)\text{, classical}
  & \text{bare $\Ainf$ alg.} & U(\fg)^!
\end{array}
\]
The dg-shifted Yangian is the \textbf{$E_1$-chiral Koszul dual},
not the pure $E_1$ Koszul dual.  The spectral parameter $u$ comes
from the closed-colour data retained by the ordered chiral bar
complex: the bar differential $d_{\bar{B}^{\mathrm{ord}}}$ uses
OPE residues from the curve (closed-colour data absorbed as
internal differential in the open-colour projection).  A pure
$E_1$ bar complex (forgetting chirality) would discard this data
and produce a bare $\Ainf$ algebra without spectral parameter.

The key structural fact: the ordered chiral bar complex
$\bar{B}^{\mathrm{ord}}(\cA)$ is the \emph{open-colour projection}
of the SC${}^{\mathrm{ch,top}}$-bar complex.  ``Open-colour
projection'' means: project to the $E_1$ operadic structure
(deconcatenation coproduct) while \emph{absorbing} the
closed-colour data as the internal differential.  The coalgebra
structure (deconcatenation) is purely combinatorial; the
differential (OPE residues) carries all the holomorphic/chiral
geometry.


\subsection{The $E_1$-chiral algebra structure on the Yangian}

The SC${}^{\mathrm{ch,top}}$-algebra structure on
$\cA^!_{\mathrm{line}}$ is concretely realised by the
derived additive KZ connection
(Theorem~\ref{thm:derived-additive-kz}).  The dictionary:
\[
\begin{array}{c|c}
\text{E}_1\text{-chiral algebra datum} &
  \text{KZ connection realisation} \\
\hline
\text{operation spaces} &
  \Conf_n^{\mathrm{ord}}(\mathbb{A}^1_u) \\
\text{spectral OPE} &
  A_n = \sum_{i<j}\Omega_{ij}(z_j - z_i)\,d(z_j - z_i) \\
\text{associativity} &
  \mathbb{D}_n^2 = 0
  \;\;(\text{spectral Kohno} = \text{YBE}) \\
\text{braiding} &
  R(z) = \text{regularised holonomy}
\end{array}
\]
The E$_1$-chiral algebra is NOT a vertex algebra (no
state-field correspondence with conformal grading).  It is an
algebra over the operad of ordered configuration spaces on
$\mathbb{A}^1$, with operations encoded by a flat meromorphic
connection.  The RTT relation $R_{12}(u-v)T_1(u)T_2(v) =
T_2(v)T_1(u)R_{12}(u-v)$ is the mixed-colour operation
$\mu_{1,2}$ (one closed insertion at spectral separation $u-v$,
two ordered open insertions).


\subsection{The curvature of the Yangian}

If $\cA^!_{\mathrm{line}} = Y_\hbar(\fg)$ is viewed as an
E$_1$-chiral algebra on $\mathbb{A}^1_u$, Vol~I's machinery
gives it a modular characteristic
$\kappa(Y_\hbar(\fg))$.

\medskip\noindent\textbf{Complementarity.}\enspace
By Vol~I Theorem~C (complementarity):
$\kappa(\cA) + \kappa(\cA^!_{\mathrm{line}}) = \alpha$
where $\alpha = 0$ for affine KM and $\alpha = 13$ for
Virasoro.  So:
\begin{align*}
\kappa(Y_\hbar(\fg))
&= -\kappa(\widehat{\fg}_k)
= -\frac{\dim\fg\,(k+h^\vee)}{2h^\vee}
\qquad\text{(affine KM)}, \\
\kappa(Y^{\mathrm{dg}}(\mathrm{Vir}_c))
&= \frac{26-c}{2}
\qquad\text{(Virasoro)}.
\end{align*}

\medskip\noindent\textbf{Genus tower of the Yangian.}\enspace
$\sum_{g \ge 1} F_g(Y^{\mathrm{dg}}(\mathrm{Vir}_c))\,x^{2g}
= \frac{26-c}{2}\,(\hat{A}(ix) - 1)$.
At $c = 26$: all $F_g = 0$ (critical string $\Rightarrow$
anomaly cancellation).  At $c = 13$: the genus tower equals
minus itself under Koszul duality (gravitational S-duality).

\medskip\noindent\textbf{The ``Yangian of the Yangian''.}\enspace
If $Y_{\mathrm{dg}}(\fg)$ is an E$_1$-chiral algebra, it has its
own open-colour Koszul dual
$Y_{\mathrm{dg}}(Y_{\mathrm{dg}}(\fg))$.  By bar-cobar-bar
involution
($\Omega(\bar{B}(\Omega(\bar{B}(\cA)))) \simeq \cA$),
this should recover $\widehat{\fg}_k$.  The double Koszul
dual is the original:
The \emph{full} SC${}^{\mathrm{ch,top}}$ double Koszul dual
$(\cA^!)^! = \cA$ (Theorem~\ref{thm:duality-involution})
operates on both colours simultaneously and does recover
$\widehat{\fg}_k$ --- the central extension survives because
it is a \emph{closed-colour} datum (the double pole in
the holomorphic OPE), transported through the Verdier/Ran
duality, not through the open-colour ordered bar.

\medskip\noindent\textbf{Explicit open-colour double bar for
$\mathfrak{sl}_2$: what it recovers and what it misses.}\enspace
The ordered chiral bar complex
$\bar{B}^{\mathrm{ord}}(Y_\hbar(\mathfrak{sl}_2))$ on
$\mathbb{A}^1_u$ recovers the \emph{current algebra}
$\mathfrak{sl}_2[t]$ (\textbf{without} central extension):
\begin{itemize}[nosep]
\item \emph{Degree~$1$}: Yangian generators $E_n, F_n, H_n$
  map to current modes $J^a_n$.
\item \emph{Degree~$2$}: The spectral OPE has only
  \emph{simple} poles ($[E(u), F(v)] \sim \hbar/(u{-}v)
  \cdot (\ldots)$).  The bar differential extracts only
  the zeroth product: $[E_n, F_m] = H_{n+m}$.  No
  central term $k\,n\,\delta_{n+m,0}$ appears ---
  that requires a \emph{double} pole, which the
  spectral OPE does not have.
\item The $E$--$E$ OPE has pole at $u - v = \hbar$
  (\emph{shifted off the collision diagonal}), so
  $d_{\mathrm{res}}(s^{-1}E \otimes s^{-1}E) = 0$,
  correctly reproducing $[e,e] = 0$.
\item \emph{Degree~$3$}: $d^2 = 0$ is the classical YBE for
  $r(u) = P/u$, automatic by Jacobi.
\end{itemize}
\textbf{The central extension is invisible to the
open-colour double bar.}  The parameter
$\hbar = 1/(k{+}2)$ is a deformation parameter of
$Y_\hbar(\fg)$ that encodes the level $k$ of the
original affine algebra, but its recovery as a
\emph{central extension} requires the closed-colour
(Verdier/Ran) direction of the full
SC${}^{\mathrm{ch,top}}$ Koszul duality.  The
open-colour double bar
$(\cA^!_{\mathrm{line}})^!_{\mathrm{line}}$ recovers
$U(\fg[t])$ (the current algebra without central
extension), not $\widehat{\fg}_k$.

\medskip\noindent\textbf{The gauge-gravity fingerprint.}\enspace
$\kappa(\cA) + \kappa(\cA^!) = 0$ for gauge theories
(anomaly cancellation); $= 13$ for gravity (irreducible
gravitational anomaly).  The double bar
$B^{E_1\text{-ch}}(B^{\mathrm{ch}}(\cA))$ has no curvature
obstruction in the gauge case (both genus towers cancel) but
residual curvature $13$ in the gravitational case.  The
vanishing at $c = 26$ is the critical string condition.


\section{The two-colour Koszul duality architecture}

The SC${}^{\mathrm{ch,top}}$-bar complex
$B^{\mathrm{SC}}(\cA)$ on $\FM_k(\C) \times \Conf_k^<(\R)$
carries \textbf{two colour projections}, each with its own
duality operation and output category.

\[
\begin{array}{c|c|c|c|c}
& \text{Bar complex} & \text{Duality} &
  \text{Output type} & \text{For } \widehat{\fg}_k \\
\hline
\text{Closed colour} & B^{\mathrm{ch}}(\cA)
  \text{ on Ran}(X)/\Sigma_n & D_{\mathrm{Ran}} \text{ (Verdier)}
  & \text{chiral algebra on } X & \widehat{\fg}_{-k-2h^\vee} \\
\text{Open colour} & \bar{B}^{\mathrm{ord}}(\cA)
  \text{ on } \Conf_n^{\mathrm{ord}}(X)
  & H^*(-)^\vee \text{ (linear)} &
  \text{Yangian (E}_1\text{ + spectral)} & Y_\hbar(\fg) \\
\text{Pure E}_1 & T^c(s^{-1}\bar\cA) \text{ (classical)}
  & H^*(-)^\vee & \text{bare }\Ainf & U(\fg)^!
\end{array}
\]

\medskip\noindent\textbf{Why two different outputs.}\enspace
Verdier duality $D_{\mathrm{Ran}}$ operates on D-modules on
Ran$(X)$ and produces a factorisation \emph{algebra} from a
factorisation coalgebra.  The output is automatically a chiral
algebra on~$X$ (same species as the input).  Linear duality
$H^*(-)^\vee$ operates on chain complexes and produces an
associative algebra.  The spectral parameter $u = 1/z$ in the
Yangian comes from the curve coordinate $z$ surviving in the
ordered bar differential (which uses OPE residues from all pole
orders $a_{(n)}b$, $n \ge 0$).  Verdier duality on the
\emph{symmetric} bar complex quotients out the ordering and with
it the spectral parameter.

\medskip\noindent\textbf{The $R$-matrix as cross-colour datum.}\enspace
The descent from ordered to unordered is $R$-twisted:
$B^{\mathrm{ch}}(\cA)_n \simeq
(\bar{B}^{\mathrm{ord}}(\cA)_n)^{R\text{-}\Sigma_n}$.
For the pole-free BD-commutative subclass of
$E_\infty$-chiral algebras (no OPE singularities):
$R(z) = \tau$ (the flip), descent is trivial, both
duals carry equivalent information.  For
$E_\infty$-chiral algebras with OPE poles
(Heisenberg, affine KM, Virasoro --- still $E_\infty$,
since all vertex algebras are $E_\infty$-chiral):
$R(z) \neq \tau$ (derived from the local OPE via
analytic continuation), and $\cA^!_{\mathrm{line}}$
carries strictly more structure than
$\cA^!_{\mathrm{ch}}$.  For genuinely $E_1$-chiral
algebras (nonlocal, Etingof--Kazhdan quantum vertex
algebras): $R(z) \neq \tau$ and is independent
input, not derivable from a local OPE.

\medskip\noindent\textbf{Operad-level explanation.}\enspace
At the operad level:
$\mathrm{Com}^{\mathrm{ch}} \,!\, = \mathrm{Lie}^{\mathrm{ch}}$
(Francis--Gaitsgory);
$\mathrm{Ass}^{\mathrm{ch}} \,!\, = \mathrm{Ass}^{\mathrm{ch}}$
(self-dual).  The full two-coloured operad
$\mathrm{SC}^{\mathrm{ch,top}}$ is Koszul self-dual
(Livernet + Kontsevich formality): this is \emph{not}
colour-by-colour at the quadratic level
($\mathrm{Com}^! = \mathrm{Lie} \neq \mathrm{Com}$), but
\emph{is} colour-by-colour at the chain level
($E_2^! \simeq E_2$, $E_1^! \simeq E_1$).  The self-duality
of $\mathrm{SC}^{\mathrm{ch,top}}$ means: an
$\mathrm{SC}^{\mathrm{ch,top}}$-coalgebra dualises to an
$\mathrm{SC}^{\mathrm{ch,top}}$-\emph{algebra}, and
$\cA^!_{\mathrm{ch}}$ and $\cA^!_{\mathrm{line}}$ are the
two colour-projections of this single
$\mathrm{SC}^{\mathrm{ch,top}}$-algebra.

\medskip\noindent\textbf{The Heisenberg test case.}\enspace
For $\cH_k$: $\cA^!_{\mathrm{ch}} = \mathrm{Sym}^{\mathrm{ch}}(V^*)$
(a commutative chiral algebra with $\kappa = -k$;
isomorphic to $\cH_{-k}$ as a chiral algebra but
$\cH_k^! \neq \cH_{-k}$ as objects in the classification,
per AP33).  $\cA^!_{\mathrm{line}} = Y(u(1)) = k[h_0, h_1, \ldots]$
(the abelian Yangian with spectral parameter $u$).
Pure $E_1$ KD produces $\bigwedge(J^{*(n)})$ (an exterior algebra,
no spectral parameter).  All three are genuinely different.
The ``accident'' that $Y(u(1)) \simeq \cH_{-k}$ as mode algebras
is abelian-specific; for $\widehat{\fg}_k$ with $\fg$ non-abelian,
$Y_\hbar(\fg)$ and $\widehat{\fg}_{-k-2h^\vee}$ are utterly
different algebras (different mode structures, different relations,
different operadic homes).

\medskip\noindent\textbf{The bridge: Conjecture
\textup{\ref{conj:FG-shadow}}.}\enspace
The pole-order (``commutator'') filtration on $\cA$ kills all
OPE poles of order $\ge 2$, leaving
$\gr_{\mathrm{com}}\cA = V_0(\fg)$ (the level-zero chiral
algebra, which is $E_\infty$).  The conjecture predicts
$\gr_{\mathrm{com}}(\cA^!_{\mathrm{line}})
\simeq (\gr_{\mathrm{com}}\cA)^!_{\mathrm{FG}}$
on the Koszul locus: the associated graded of the Yangian
recovers the Francis--Gaitsgory dual of the level-zero algebra.
The stratification spectral sequence ($E_1$ page $=$ FG bar
complex, convergence proved) provides the controlled
interpolation.  The filtration compatibility of the bar
differential remains the sole gap.

\medskip\noindent\textbf{When do the two duals agree?}\enspace
When $\cA$ lies in the pole-free BD-commutative subclass of
$E_\infty$-chiral algebras: $R(z) = \tau$, descent is
trivial, both duals carry equivalent information.

\noindent\textbf{When do they disagree?}\enspace
When $\cA$ is $E_\infty$-chiral with OPE poles (all
interesting vertex algebras: Heisenberg, affine KM,
Virasoro --- these are still $E_\infty$, since all
vertex algebras are $E_\infty$-chiral): $R(z)$
carries nontrivial spectral dependence derived from
the local OPE, and $\cA^!_{\mathrm{line}}$ (the
Yangian) is strictly richer than
$\cA^!_{\mathrm{ch}}$ (the FG dual).
The R-matrix is the data killed by the $\Sigma_n$-coinvariant
step.


\section{Chromaticity in the bar-cobar programme}
\label{sec:chromatic-working-notes}

\subsection{What is proved}

Two results in the core monograph place the chromatic programme
on firm algebraic ground.

\begin{enumerate}
\item \textbf{Smashing localisation commutes with factorization
  homology} (Theorem~1.71 in the core chapter on logarithmic HT
  monodromy): for a locally constant spectral factorization algebra $A$
  on an $m$-manifold $M$ and a smashing localisation $L$,
  $L(\int_M A) \simeq \int_M LA$.  Since $L_{E(n)}$-localisation is
  smashing, chromatic height is compatible with the local-to-global
  passage once the theory is spectralized.

\item \textbf{The scalar genus tower is $\hat{A}(ix)$}
  (Theorem in the Rosetta stone chapter):
  $F^{\mathrm{tot}}(\cA;\hbar)
  = (\kappa(\cA)/\hbar^2) \cdot (\hbar/2)/\sin(\hbar/2)$.
  The generating function $(x/2)/\sin(x/2) = \hat{A}(ix)$ is the
  Wick rotation of the $q^0$ term of the Witten genus characteristic
  series.  The higher $q$-powers correspond to the Eisenstein-weighted
  genus-1 corrections: on $E_\tau$, each pole of order $2k$ in the OPE
  generates a correction weighted by $E_{2k}(\tau)$.
\end{enumerate}

\subsection{The genus tower as a chromatic shadow}

The genus tower has three structural features connecting it to
chromatic homotopy theory:

\begin{itemize}
\item \textit{Curvature as periodicity generator.}
  The element $u = \kappa(\cA) \cdot \omega_g$ in the transgression
  algebra controls the curved bar structure $d^2 = \kappa \cdot \omega_g$.
  Its invertibility creates Clifford periodicity; its vanishing
  ($\kappa = 0$, the critical level) collapses the genus tower.
  This parallels $v_n$-periodicity in chromatic homotopy theory.

\item \textit{Eisenstein weights and formal groups.}
  The genus-1 bar differential replaces the additive kernel $1/z$
  by the Weierstrass $\zeta$-function $\zeta(z|\tau)$, the logarithm
  of the elliptic formal group $\hat{E}_\tau$.  The formal group of
  $E_\tau$ has height $\le 2$ (height 1 ordinary, height 2
  supersingular), so genus-1 data accesses chromatic height $\le 2$.
  The Eisenstein series $E_{2k}(\tau)$ for $k \ge 2$ are sections of
  $\omega^{\otimes 2k}$ on $\cM_{1,1}$; $E_2$ is quasi-modular
  (the Arakelov defect).

\item \textit{Clutching and redshift.}
  The non-separating clutching $D_1 \colon B^{(g)} \to B^{(g+1)}$
  adds one handle, increasing $\dim(J)$ by 1 and the $p$-divisible
  group height $\mathrm{ht}(J[p^\infty])$ by 2 (since
  $\mathrm{ht} = 2 \cdot \dim$ for abelian varieties).
  Whether this is a geometric counterpart of chromatic redshift
  ($K$-theory raises height by 1) is an open analogy---the factor-of-2
  mismatch must be understood.
\end{itemize}

\subsection{Anti-patterns: what is \emph{not} true}

\begin{description}
\item[AP-CHR] The classical bar complex $B(\cA)$ is a chain complex
  over $\C$, hence an $H\C$-module, hence rational:
  $L_{K(n)}(B(\cA)) = 0$ for all $n \ge 1$.  The entire classical
  bar-cobar theory is chromatic height 0.  No OPE pole produces
  positive-height chromatic data in the linear theory.  The additive
  formal group $\hat{\mathbb{G}}_a$ of the spectral-parameter
  composition law has height $\infty$ over $\mathbb{F}_p$ (not height 1;
  height 1 is $\hat{\mathbb{G}}_m$).  Chromatic height $\ge 1$ enters
  only after spectralization.

\item[AP-RED] Koszul duality $A \mapsto A^!$ reflects the curvature
  ($\kappa \mapsto -\kappa$ for KM, $\kappa \mapsto 13 - \kappa$ for
  Virasoro) but does not raise chromatic height.  The clutching $D_1$
  raises $p$-divisible group height by 2, not 1.  Formal-group height
  is a $p$-local invariant: it requires reduction mod $p$, not a
  complex parameter.

\item[AP-TOWER] The chromatic convergence theorem (Hopkins--Ravenel) uses
  $L_{E(n)}$-localisation (Johnson--Wilson), not $L_{K(n)}$.  The
  $K(n)$-localisations do not form a tower with natural transition maps.
\end{description}

\subsection{The four-stage factorization}

The algebraic route from a chiral algebra to a chromatic invariant:
\[
  \cA
  \;\xrightarrow{\;\mathrm{bar}\;}
  B_{\mathrm{mod}}(\cA)
  \;\xrightarrow{\;\mathrm{FT}_{\mathrm{rel}}\;}
  \cZ_\cA
  \;\xrightarrow{\;\text{sp.\ lift}\;}
  \cZ_\cA^{\mathrm{sp}}
  \;\xrightarrow{\;L_{E(n)}\;}
  \cZ_{\cA,n}\,.
\]
Stages 1--2 (modular bar construction and relative Feynman transform) are
proved.  Stage 3 (spectralization) is the fundamental open problem.
Stage 4 (chromatic localisation) is automatic from the smashing-localisation
theorem.

\subsection{Conjectural directions (quarantined from the manuscript)}

The following directions were explored computationally and found to contain
valuable intuitions alongside serious mathematical errors.  They are
preserved here for future development, with the errors flagged.

\begin{enumerate}
\item \textbf{Synthetic spectra and the chiral bar complex.}
  The bigrading (holomorphic weight, topological arity) parallels
  Pstr\k{a}gowski's synthetic bigrading (stem, Adams filtration).
  The curvature $\kappa$ plays the role of $\tau$: at $\kappa = 0$ the
  theory decouples.  \emph{Errors found:} the $\tau = 0$ fibre of the
  synthetic SC operad should be a semidirect product, not a coproduct;
  the Adams SS for $B(\cA)$ is trivial over $\C$.

\item \textbf{Transchromatic phenomena in the genus tower.}
  The genus tower should correspond to the chromatic tower, with
  genus $g$ accessing heights $\le 2g$ via $\mathrm{Jac}(\Sigma_g)$.
  \emph{Errors found:} the height-genus identification used $L_{K(n)}$
  instead of $L_{E(n)}$; the transition maps
  $\overline{\cM}_g \to \overline{\cM}_{g-1}$ do not exist;
  formal-group heights were assigned over $\C$ without $p$-local
  specification; $p$-divisible group heights were off by factor 2.

\item \textbf{The chromatic modular functor.}
  The open-sector category $\cC_{\partial,\cA}$ should be fully
  dualizable, producing a 3d framed TQFT via the cobordism hypothesis.
  The Stolz--Teichner programme predicts the spectral lift yields
  tmf-valued invariants.  \emph{Errors found:} the tmf weight sign was
  wrong ($\eta^{-\kappa}$ has weight $-\kappa/2$, not $\kappa/2$); the
  Hurewicz image claim was absurdly strong.

\item \textbf{Telescope conjecture failure.}
  Non-formal SC-algebras ($m_k \ne 0$ for $k \ge 3$) should exhibit
  telescope conjecture failure at heights $\ge 2$.  \emph{Errors found:}
  classical $A_\infty$ non-formality is a characteristic-zero phenomenon
  and cannot directly cause $p$-local spectral telescope failure;
  blueshift was claimed for continuous $C_p \subset U(1)$ actions when
  the theorem applies to finite group actions on finite spectra;
  $\mathrm{SL}_2(\Z)$ has $p$-torsion only for $p = 2, 3$.

\item \textbf{Physical ramifications.}
  The chromatic programme has potential physical consequences:
  the perturbative/non-perturbative divide (formal vs non-formal SC
  structure) may relate to the telescope watershed; the
  $c = 26$ dichotomy ($\kappa_{\mathrm{eff}} = 0$) corresponds to
  chromatic decoupling; the arithmetic fracture square (rationalization
  $+$ $p$-completion) may correspond to the perturbative/non-perturbative
  split in QFT.  These remain speculative analogies.
\end{enumerate}

The quarantined source material (1710 lines) is preserved in
\texttt{archive/source\_tex/chromatic\_frontier\_quarantined.tex}.

%% ================================================================
\section{Session notes: 2 April 2026 --- The Polyakov programme revisited}
\label{sec:session-2apr26-polyakov}
%% ================================================================

Twenty-three adversarial agents --- eleven computational, twelve
theoretical --- attacked the entire programme from first principles,
guided by the question: \emph{what does the bar-cobar machinery say
about the Polyakov programme that the path integral cannot?}  The
findings below are organized by theme.  Errors found and corrected are
marked~{\bfseries [FIX]}.

\subsection{Arithmetic corrections}
\label{subsec:polyakov-arith}

\paragraph{{\bfseries [FIX]} Brown--Henneaux specialisation of
$\mathfrak{Q}$.}  Three cascading factor-of-$2$ errors in
Computation~\texttt{comp:gravity-quartic-correction}
(3d\_gravity.tex lines 2126--2143).  The root cause: when factoring
$30k+22 = 2(15k+11)$, the factor~$2$ was absorbed into reducing
$6k \to 3k$ while simultaneously halving $10 \to 5$, double-counting
the cancellation.  Corrected values:
\begin{alignat*}{2}
\mathfrak{Q}(6k) &= \frac{5}{6k(15k+11)},
&\qquad \mathfrak{Q} &\sim \frac{1}{18k^2},\\
\delta_H^{(1)}(6k) &= \frac{5}{3k^2(15k+11)},
&\qquad \delta_H^{(1)} &\sim \frac{1}{9k^3}.
\end{alignat*}
The master formula $\mathfrak{Q} = 10/[c(5c+22)]$ and the genus-$2$
contact correction $F_2^{\mathrm{contact}} = 1/[24\,c(5c+22)]$ are
unchanged.

\subsection{Computational verifications}
\label{subsec:polyakov-comput}

\paragraph{Ghost-number obstruction: universal for all principal DS.}
The ghost-number bottleneck ($\mathrm{gh}(h) = -1$, source trees at
ghost~$+1$, target trees at ghost~$-1$, $p \otimes p$ kills both)
is \emph{topological}: it depends only on the Z-grading structure
of the BRST complex and the tree combinatorics, not on the rank of
$\mathfrak{g}$, the number of ghost pairs, or the Dynkin grades.
Verified explicitly for $\mathfrak{sl}_3$ (three ghost pairs, grades
$1,1,2$): no rank-dependent obstruction exists.  $\Delta_{z,n} = 0$
for $n \ge 2$ for all principal $\mathcal{W}_k(\mathfrak{g})$.

\paragraph{$\hat{A}$-genus formula: scalar lane only.}  The identity
$\hat{A}(ix) = (x/2)/\sin(x/2)$ verified to $15$-digit precision.
Taylor coefficients match Faber--Pandharipande integrals
$\lambda_g^{\mathrm{FP}}$ through $g = 5$.  Critical distinction:
the formula $\sum F_g x^{2g} = \kappa_{\mathrm{eff}}(\hat{A}(ix)-1)$
captures \emph{only the scalar arity sector}.  The full genus tower
includes higher-arity corrections: $F_2^{\mathrm{scalar}} =
7(c-26)/11520$ (linear in $c$) vs.\
$F_2^{\mathrm{contact}} = 1/[24\,c(5c+22)]$ ($\sim 1/c^2$ at large
$c$).  At $c = 26$: scalar tower vanishes but contact correction
$F_2^{\mathrm{contact}} = 1/94848 \ne 0$ --- the genus-$2$ partition
function does \emph{not} vanish at the critical string.  For
Heisenberg (class~G, $r_{\max} = 2$), the $\hat{A}$-genus formula
is exact at all genera.

\paragraph{DS-HPL transfer.}  Verified: $m_3^{\mathrm{Vir}}$ arises
\emph{entirely} from HPL tree corrections (two Catalan trees with
one $h$-insertion), since the affine algebra has $m_3^{\mathrm{aff}}
= 0$ (class~L).  The homotopy~$h$ acts as a propagator of acyclic
BRST directions.  The Sugawara central charge
$c_{\mathrm{Sug}} = 3k/(k+2)$ is \emph{not} the DS-reduced central
charge $c_{\mathrm{DS}} = 1 - 6(k+1)^2/(k+2)$; these are distinct
at all $k$.

\paragraph{Spectral $R$-matrix.}  Laplace transform verified:
$r^L(z) = \partial T/z + 2T/z^2 + (c/2)/z^4$.  The bar collision
residue $r^{\mathrm{coll}}(z) = (c/2)/z^3 + 2T/z$ has pole orders
one less (AP19: $d\log$ absorbs a pole).  The $d\log$ absorption
converts even-order poles of $r^L$ into odd-order poles of
$r^{\mathrm{coll}}$, making antisymmetry \emph{manifest} for the
collision residue while it remains non-manifest (but true, via the
braiding inverse relation) for the Laplace kernel.  CYBE verified
term-by-term via the PVA Jacobi identity.  Arnold relation on
$\FM_3(\C)$: Stokes argument correct; orientations on
$D_{\{1,2\}}$ and $D_{\{2,3\}}$ are opposite; corner cancellation
by Koszul signs.

\paragraph{Soft graviton hierarchy.}  The correspondence
$m_r \leftrightarrow S^{(p)}$ at $p = r-2$ is \emph{proved} as a 2d
CFT statement at all orders (Ward identities from the shadow
connection).  Identification with 4d soft graviton factors via
celestial holography: proved for $p = 0$ (Weinberg) and $p = 1$
(Cachazo--Strominger); proved with explicit computation at $p = 2$
(CHY sub-subleading, $\mathfrak{Q} = 10/[c(5c+22)]$); $p \ge 3$
are \emph{new predictions} ($m_5$ coefficient $20/[c^2(5c+22)]$),
not yet verified by independent methods.  For Heisenberg:
$m_3 = 0$, so no non-trivial subleading soft structure beyond the
universal Weinberg factor --- consistent with free-field scattering.

\paragraph{$\mathcal{W}_3$ Koszul dual.}  Complementarity constant
$\alpha_3 = 100$ verified via ghost central charges:
$c_{\mathrm{gh}}(s=2) = -26$, $c_{\mathrm{gh}}(s=3) = -74$,
$\alpha_3 = 100$.  Self-dual point $c^* = 50$.  Structure constant
$\beta^2|_{c=50} = 1/17$.
Complementarity sum
$\kappa(\mathcal{W}_{3,c}) + \kappa(\mathcal{W}_{3,100-c})
= 250/3 \ne 0$ (AP24 active).
\emph{DS transport gives $c + c' = 2(N-1)$, not $\alpha_N$.}
This confirms AP33: Koszul duality $\ne$ Feigin--Frenkel involution
$\ne$ negative-level substitution.  For the Bershadsky--Polyakov
algebra $(\mathfrak{sl}_3, f_{(2,1)})$: the conjectured
$c_{\mathrm{sum}} = 76$ is derivable from neither the ghost formula
($28$) nor the FF transport ($12$); its status remains
\emph{conditional}.

\paragraph{BTZ Cardy formula.}  Algebraically verified:
$S = 2\pi\sqrt{ch/6} + 2\pi\sqrt{c\bar{h}/6}$ follows from
(i)~curvature $\kappa = c/2$ (proved, genuine bar-complex input),
(ii)~modular invariance (\emph{conjectural} from the bar complex),
(iii)~saddle-point (standard Cardy analysis, no bar-complex content).
BTZ identification via Brown--Henneaux and standard BTZ
parametrisation verified to machine precision.

\paragraph{Pole-order classification.}  Class~C \emph{is} populated
(the $\beta\gamma$ system: maximal OPE pole order~$3$ from composite
fields, shadow depth~$4$, $m_5 = 0$ by rank-$1$ abelian rigidity).
Shadow depth is \emph{not} an exact function of pole order: classes~G
and~L share maximal pole order~$2$ but have depths~$2$ and~$3$.
The discriminant between G and~L is centrality of the double-pole
residue.  All $d \ge 4$ collapse to class~M with infinite depth.

\paragraph{Modular $d_1$ for $\widehat{\mathfrak{sl}}_2$.}
Genus-$1$ curvature $\kappa(V_k(\mathfrak{sl}_2)) = 3(k+2)/4$.
Vanishes iff $k = -2$ (critical level).  Note:
$\kappa \ne c/2 = 3k/(2(k+2))$ for affine algebras; the relation
$\kappa = c/2$ holds for Virasoro but not in general.  Sign
subtleties in the bicomplex verification on $3$-flag corollas are
genuine: orientation-line contributions from $\det(H_1(\Gamma;k))$
modify the naive collision-sign formula.  The structural proof via
codimension-$2$ cancellation ($\partial^2 = 0$ on
$\overline{\mathcal{M}}_{g,n}$) is correct.

\subsection{Theoretical attacks}
\label{subsec:polyakov-theory}

\paragraph{Bar vs.\ Bulk (AP-OC refinement).}
$\barB(\cA)$ functorially \emph{determines} $\Zder(\cA)$ via
$\Zder(\cA) = \RHom(\Omega(\barB(\cA)), \cA) =
C^\bullet_{\mathrm{ch}}(\cA, \cA)$ --- the bar complex is the
resolution that \emph{computes} the derived centre.  The AP-OC
distinction is the distinction between a resolution and the
invariant it computes: genuine (coalgebra vs.\ dg-algebra) but the
rhetoric should acknowledge the functorial dependence.  Best
formulation: \emph{$\barB(\cA)$ is the presenting coalgebra for the
entire Koszul triangle, as the Steinberg variety is the presenting
variety for the Hecke algebra.  The bulk, boundary, and lines are
three different functors applied to $\barB(\cA)$.}

\paragraph{{\bfseries [FIX]} Non-principal primitivity: linear
coproduct gap.}  The ghost-number obstruction kills all
\emph{nonlinear} HPL corrections ($\Delta_{z,n} = 0$ for $n \ge 2$)
for all nilpotent orbits.  But the \emph{linear} transferred
coproduct $\Delta_{z,1}$ requires a separate argument.  For
non-principal reductions whose $\mathcal{W}$-algebra retains
weight-$1$ generators (e.g.\ Bershadsky--Polyakov with its
$\widehat{\mathfrak{sl}}_2$ subalgebra), the projection~$p$ does not
automatically kill split-Casimir cross-terms.  Physical
interpretation: non-principal $=$ matter-coupled gravity; weight-$1$
generators $=$ propagating gauge fields; non-primitive linear
coproduct $=$ particle production from fundamental modes.
Remark~\texttt{rem:non-principal-ghost} corrected to reflect this
scope.

\paragraph{$c^* = 13$: algebraically proved, physically open.}
Koszul self-duality $\mathrm{Vir}_{13}^! = \mathrm{Vir}_{13}$ is
proved.  The complementarity invariant $\kappa + \kappa' = 13$ is
proved.  Verdier duality and modular $S$-transform commute (they
act on different variables: $c$ vs.\ $\tau$).  At $c = 13$, Verdier
duality is the identity, so the $S$-transform has \emph{no}
special properties.  Physical realisability by a unitary interacting
CFT at $c = 13$ is \emph{open} (Liouville at $c = 13$ requires
complex $b$).  The gravitational number field $K(13) =
\Q(\sqrt{870})$ has class number $h = 8$ (maximal in $[0,26]$).
The strongest physical interpretation: tree-level/one-loop balance
($c_{\mathrm{total}} = c - 26 = -13 = -(26-c)$).

\paragraph{3d uplift: dissolves, does not solve, the Polyakov loop.}
The monograph \emph{never claims} that the Liouville field is the
radial direction.  The 3d theory repackages the anomaly as
$\kappa_{\mathrm{eff}} = (c-26)/2$, a single datum generating the
genus tower, rather than a recursive Liouville cancellation.  The
advance is structural: the anomaly is not pathological but
generative.  The key gap: convergence of the full Virasoro partition
function (beyond the scalar lane) is \emph{open}.  The scalar lane
converges exponentially ($F_g \sim 2\kappa/(2\pi)^{2g}$, radius of
convergence $4\pi^2$), \emph{not} factorially.  The full double
series (genus $\times$ arity) for class~M is open.

\paragraph{Genus spectral sequence: convergence for class~M.}
The SS does not degenerate at any finite page for Virasoro at generic
$c$ (the non-vanishing quartic contact $\mathfrak{Q}$ forces
$d_2 \ne 0$ via the separating clutching bracket
$[\Theta_1, \Theta_1]$).  But the SS converges in the categorical
sense (complete filtration, Eilenberg--Moore).  The \emph{physically
meaningful question} is not SS degeneration but existence of the
total MC element $\Theta \in \mathrm{MC}(L_{\mathrm{mod}})$,
guaranteed by the completeness of the genus filtration whenever all
obstruction classes vanish.  Three convergence regimes:
\begin{enumerate}
\item Scalar lane: $\sum F_g \hbar^{2g}$ converges for
  $|\hbar| < 4\pi^2$ (exponential decay, Borel transform entire).
\item Shadow obstruction tower (fixed genus~$0$): converges for $c > c^* \approx
  6.125$; diverges below.
\item Full double series (all genera, all arities): open for class~M.
\end{enumerate}

\paragraph{Genus-$1$ curvature-braiding: resolved.}
The discriminant is the \emph{Coisson bracket} $c_0 = \{a_{(0)}b\}$,
\emph{not} the pole order or shadow depth.  The dichotomy cuts
\emph{orthogonally} to the G/L/C/M classification:
\begin{center}
\small
\begin{tabular}{lllll}
\textbf{Class} & \textbf{Algebra} & $c_0$ &
\textbf{Genus-$1$} & \textbf{Physical} \\
\hline
G & Free fields & $0$ & Decoupled & Trivial braiding \\
L & Heisenberg & $0$ & Decoupled & $Z = \eta^{-k}$, $R = e^{k/z}$
    independent \\
L & Affine KM & $f^{ab}_c J^c \ne 0$ & Entangled &
    Verlinde formula mixes $S$ and fusion \\
M & Virasoro & $\partial L \ne 0$ & Entangled &
    BTZ entropy mixes $c$ (modular) and $h$ (braiding) \\
\end{tabular}
\end{center}
The elliptic spectral dichotomy
(Theorem~\texttt{thm:elliptic-spectral-dichotomy}) is proved.
Remaining frontier: genus $g \ge 2$ (requires Fay trisecant identity
on $\Sigma_g$).

\subsection{The Polyakov programme: what the bar complex computes}
\label{subsec:polyakov-synthesis}

Every structural layer of the Polyakov programme has an algebraic
shadow in the bar-cobar machinery.  The correspondence is not
metaphorical: it is a precise identification, proved unconditionally
on the Koszul locus.

\begin{center}
\small\renewcommand{\arraystretch}{1.2}
\begin{tabular}{lll}
\textbf{Polyakov layer} & \textbf{Bar-cobar shadow} & \textbf{Status}
\\
\hline
Genus expansion $\sum g_s^{2g-2} Z_g$ &
  Modular bar coalgebra $\mathbf{Mod}(C)$ & Proved \\
Conformal anomaly $(c-26)/12$ &
  Bar curvature $d^2 = \kappa_{\mathrm{eff}} \cdot \omega_g$ &
  Proved \\
Belavin--Knizhnik factorization &
  Swiss-cheese splitting $\FM_k(\C) \times \Conf_k^{<}(\R)$ &
  Proved \\
BPZ bootstrap (genus~$0$) &
  MC equation on $\FM_3(\C)$: Arnold $\Rightarrow$ YBE &
  Proved \\
Polyakov formula ($g = 1$) &
  $F_1 = \kappa_{\mathrm{eff}}/24$ & Proved \\
Modular invariance &
  Clutching axiom on $\overline{\mathcal{M}}_{g,n}$ &
  Conjectural \\
All-genera: $\sum F_g x^{2g}$ &
  $\kappa_{\mathrm{eff}}(\hat{A}(ix)-1)$ (scalar lane) &
  Proved \\
Liouville sector cancellation &
  \emph{Dissolved}: $\kappa_{\mathrm{eff}}$ generates, not obstructs
  & Framework \\
\end{tabular}
\end{center}

\noindent\textbf{Two distinguished central charges, not one.}
Polyakov identified $c = 26$ (anomaly vanishes).  The bar-cobar
programme reveals $c^* = 13$ (Koszul self-duality).  These are
genuinely different: $c = 26$ asks ``when does the anomaly vanish?'';
$c^* = 13$ asks ``when is the algebra self-dual?''  The second is
invisible from the path integral.

\noindent\textbf{The ghost-number theorem.}  Polyakov's path integral
gives no structural reason for $3$d gravity's lack of local degrees
of freedom.  The ghost-number obstruction provides one: the
gravitational coproduct is strictly primitive at all arities, for
all principal DS reductions, via a topological (tree-combinatorial)
mechanism.  Gravity produces no particles, only braids.

\noindent\textbf{The deepest point.}  The Polyakov programme computed
projections of the universal MC element $\alpha_T \in
\mathrm{mc}(\mathfrak{g}^{\mathrm{SC}}_T)$
without knowing it existed.  The six projections (boundary $\Ainf$
algebra, line-sector operations, spectral $R$-matrix, bulk-to-boundary
map, genus tower, PVA bracket) are unified by $\alpha_T$. Polyakov
saw the genus tower.  BPZ saw the OPE and crossing symmetry.
Maldacena later saw the bulk-to-boundary map.  The bar-cobar
programme identifies the common source.

\subsection{Further results from the adversarial campaign}
\label{subsec:polyakov-further}

\paragraph{$\hat{A}$-genus: why $\hat{A}$ and not Todd.}
Todd is not even in~$x$; the genus tower has only even powers.
$\hat{A}(x)^2 = \mathrm{Td}(x) \cdot \mathrm{Td}(-x)$ is the
correct symmetrisation, forced by the real structure on the Hodge
bundle via Serre duality $\mathbb{E}^* \simeq \mathbb{E} \otimes
\omega$ (KO-theory orientation).

\paragraph{Two depths for $\mathcal{W}_N$.}
Perturbative depth $d(\mathcal{W}_N) = \infty$ (class~M).
$\mathcal{W}$-essential depth $d_{\mathcal{W}}(\mathcal{W}_N) = N$:
after arity~$N$, the master equation closes.  Proved $N = 2,3,4$;
conjectural in general.  Resolvent tree formula: Catalan-many terms
at every~$N$.

\paragraph{Swiss-cheese directionality: forced, not chosen.}
Forced at three levels (geometric, D-module, BV).  Relaxing it
destroys operadic K\"unneth, breaks the triangular MC cascade, and
couples bulk to boundary.  Cylinder diagrams are trace/clutching
(loop-level), not operadic (tree-level).

\paragraph{Complementarity sum: the universal pattern.}
$K = 0$ for classes G, L, C (involution \emph{negates}).
$K(\mathcal{W}_N) = \alpha_N(H_N - 1) \ne 0$ for class~M (involution
\emph{shifts}).  Virasoro is the unique $\mathcal{W}_N$ with
$K = c^*$ (only $N = 2$ has $H_N = 3/2$).

\paragraph{Polyakov formula = genus-$1$ bar: precise status.}
Theorem for Heisenberg/Virasoro.  Conjecture for general HT
theories at $g \ge 1$ (Conj.\ \texttt{conj:master-bv-brst}).
Mumford exponent~$13$ and $c^* = 13$ both trace to
$c_{\mathrm{ghost}} = -26$.

\paragraph{$\cC_{\mathrm{grav}}^{\mathrm{op}}$: no algebraic model.}
Primitive coproduct $\ne$ trivial category: the $r$-matrix encodes
all gravitational interaction as braiding.  The six requirements
(boundary algebra, bulk, primitive coproduct, line category, annulus
trace, genus tower) have no simultaneous algebraic realization.

\paragraph{Virasoro $m_k$: Shapovalov obstruction.}
Closed-form $P_k$ $\Leftrightarrow$ inverting the Shapovalov form
at all levels $\Leftrightarrow$ constructing all singular vectors.
Universal $(x+2)$ factor at symmetric point (verified $k \le 5$).
Generating function: the graviton self-energy resolvent
$G(t;\lambda) = \mathrm{Tr}_h(1 - tK(\lambda))^{-1}$.

\subsection{Open frontier}
\label{subsec:polyakov-frontier}

\begin{enumerate}
\item \textbf{Full Virasoro convergence.}  Scalar lane converges;
  shadow obstruction tower converges for $c > c^*$.  Full double series open.
\item \textbf{BP linear coproduct.}  If non-primitive: first
  ``matter-coupled gravity'' example.
\item \textbf{Genus $g \ge 2$ curvature-braiding.}  Coisson
  bracket $c_0$ via Fay trisecant identity.
\item \textbf{BP $c_{\mathrm{sum}} = 76$.}  Neither ghost ($28$)
  nor FF ($12$) produces $76$.
\item \textbf{Modular invariance from the bar complex.}
\item \textbf{Soft graviton $p \ge 3$.}  The quintic shadow
  is $S_5 = -48/[c^2(5c+22)]$ (from direct shadow-metric
  expansion; the heuristic $[\mathfrak{C},\mathfrak{Q}]_H$
  bracket gave $20/[c^2(5c+22)]$, which is a sub-channel).
  The \emph{Stasheff cancellation}: level-$6$ Shapovalov
  poles at $c = 1/2$ and $c = -68/7$ cancel in the full
  $\Ainf$ computation, and only $c^2(5c+22)$ survives in
  the denominator.  Falsification: verify $S_5$ numerically
  at $c = 1, 10, 26$.
\item \textbf{$\cC_{\mathrm{grav}}^{\mathrm{op}}$ algebraic model.}
\item \textbf{Closed-form $P_k$ or spectral curve.}  Sharp
  prediction: $P_k = \tfrac{55}{2}P_{k-1} - \tfrac{151}{2}P_{k-2}$
  holds for $k = 4,5$ (interpolation).  If $P_6 = -2927/2$
  (computable from HPT on the $\widehat{\mathfrak{sl}}_2$ Verma
  module), the symmetric-point scalar resolvent is rational:
  $S(u) = u^2(2-49u)/(2-55u+151u^2)$, with eigenvalues
  $(55 \pm \sqrt{1817})/4$ where $1817 = 23 \cdot 79$.  The
  individual $P_k$ have no closed form (Shapovalov obstruction) but
  the generating function may be algebraic---analogous to Catalan
  numbers.  \textbf{Stasheff cancellation}: level-$6$ Shapovalov
  poles ($c = 1/2, -68/7$) cancel in $S_5 = -48/[c^2(5c+22)]$;
  only the shadow-metric denominator $(5c+22)$ survives.
\end{enumerate}

\subsection{Second-wave findings (frontier campaign)}
\label{subsec:polyakov-second-wave}

Six results from the second-wave adversarial agents, resolving or
sharpening items from the open frontier list above.

\paragraph{BP coproduct resolution.}
The bar-cobar coproduct $\Delta\colon \barB(\cA) \to
\barB(\cA) \otimes \barB(\cA)$ is \emph{primitive} on the free
generators $J$, $G^+$, $G^-$ of the $N=2$ superconformal algebra
$B^k$: these are fundamental affine generators or ghost-projected
composites, and primitivity follows from the structure theorem
for cofree coalgebras (a generator that is not a product of
lower-arity generators is primitive).  The stress tensor~$T$
is \emph{non-primitive}: a surviving Cartan Casimir contribution
$J \otimes J$ appears in $\Delta(T)$ from the Sugawara construction.
The gauge/gravity dichotomy sharpens: for principal DS reductions,
\emph{all} fields are primitive (the ghost-number obstruction kills
all cross-terms); for non-principal reductions, free generators are
primitive but composites (Sugawara-type fields) are non-primitive.
This confirms that the $B^k$ algebra is the first explicit
``matter-coupled gravity'' example with non-trivial coproduct
structure.

\paragraph{{\bfseries [FIX]} BP complementarity: $c_{\mathrm{sum}} = 2$,
not $76$.}
The correct central charge formula for $B^k$ is
$c(B^k) = (k - 15)/(k + 3)$.  The Koszul dual has
$c(B^{k,!}) = (k' - 15)/(k' + 3)$ where $k' = -k - 6$
(the Feigin--Frenkel involution for $\mathfrak{osp}(3|2)$),
giving $c + c' = 2$.  This is proved in Vol~I,
Proposition~\texttt{prop:partition-dependent-complementarity}.
The ghost formula ($c_{\mathrm{gh}} = 28$ for $\dim = 14$ ghost
pairs) is the \emph{wrong} formula for non-principal $\cW$-algebras:
the ghost-number grading differs from the conformal grading, and the
effective ghost contribution depends on the embedding
$\mathfrak{sl}_2 \hookrightarrow \mathfrak{osp}(3|2)$.
The earlier frontier entry ``$c_{\mathrm{sum}} = 76$'' is superseded.

\paragraph{Genus-$g$ curvature-braiding uniformity.}
The Bergman kernel $B(z,w) = \partial_z \partial_w \log E(z,w)$ is
globally meromorphic at every genus~$g$: the quasi-periodic phases of
the prime form~$E(z,w)$ cancel in the second mixed derivative.  The
intermediate object --- the abelian differential of the third kind
$\partial_z \log E(z,w)$ --- is quasi-periodic with monodromy
$-2\pi i \, c_0 \, \omega_j(z)$ per $B$-cycle $\mathfrak{b}_j$,
where $\omega_j$ is the normalised abelian differential and $c_0 =
\{a_{(0)}b\}$ is the Coisson bracket.  At genus~$g$ there are~$g$
independent $B$-cycle monodromies, all proportional to~$c_0$.  The
discriminant $c_0$ is genus-independent: the curvature-braiding
dichotomy ($c_0 = 0$ vs.\ $c_0 \ne 0$) persists uniformly across all
genera.  This resolves frontier item~(3).

\paragraph{Modular covariance: automatic from Maurer--Cartan.}
Modular covariance of the genus-$1$ partition function is
\emph{automatic} from the Maurer--Cartan equation on the moduli stack
$[\mathfrak{H}/\mathrm{SL}(2,\Z)]$.  The partition function $Z_1(\tau)
= \eta(\tau)^{-\kappa_{\mathrm{eff}}}$ is a modular form of weight
$-\kappa_{\mathrm{eff}}/2$, not modular \emph{invariant}.  Modular
invariance (vanishing weight) requires left-right matching
$\kappa_L + \kappa_R = 0$, which is an \emph{external} condition (bulk
pairing), not derivable from the chiral bar complex alone.  The
Eisenstein quasi-modularity defect of~$E_2$ (the Arnold defect in the
holomorphic bar differential at genus~$1$) is resolved by the
non-holomorphic completion $\hat{E}_2 = E_2 - 3/(\pi\,\mathrm{Im}\,
\tau)$, which restores modularity at the cost of holomorphicity --- the
holomorphic anomaly equation.  This resolves frontier item~(5).

\paragraph{Full double series convergence.}
For $c > c^* \approx 6.125$ (growth parameter $\rho < 1$), the full
double series $\sum_{g,r} Z_{g,r}\,\hbar^{2g}\,t^r$ converges
absolutely in the bidisk $\{|\hbar| < 4\pi^2\} \times \{|t| < |c|/6\}$,
by the factored bound
\[
|Z_{g,r}| \;\le\; C \cdot (2\pi)^{-2g} \cdot \rho^r \cdot r^{-5/2}.
\]
For $1 < \rho < 2\pi$: convergence persists with reduced radius
$R_{\mathrm{eff}} = 4\pi^2/\rho^2$.  For $\rho > 2\pi$: genuine
divergence of the double series.  The scalar Borel transform (in the
$\hbar$-variable) is entire; non-perturbative completion via BTZ
trans-series (resurgent contributions from BTZ black hole saddles) is
conjectural.  This resolves frontier item~(1).

\paragraph{{\bfseries [FIX]} Partial fraction and residue corrections.}
Three arithmetic errors in
\texttt{thqg\_perturbative\_finiteness.tex}:
\begin{itemize}
\item Sign error in partial-fraction decomposition:
  $(-1)^{n+1} \to (-1)^n$.
\item Residue correction:
  $(-1)^{n+1} \cdot 2\kappa n\pi \to (-1)^n \cdot 8\kappa n^2 \pi^2$.
\item Intermediate expansion coefficient:
  $\varepsilon/(4n\pi) \to \varepsilon/(8n\pi)$.
\end{itemize}
All three errors were cascading from the same root: the sign convention
for the cotangent partial fraction $\pi\cot(\pi z) = 1/z +
\sum_{n=1}^\infty 2z/(z^2 - n^2)$ was applied with $(-1)^{n+1}$
(alternating-sign convention from the $\csc$ expansion) instead of the
correct $(-1)^n$ (from the $\cot$ expansion).  Corrected in source.


% =================================================================
\section{The Wick rotation architecture}
\label{sec:wick-architecture}
% =================================================================

The Wick rotation $x \mapsto ix$ in the genus tower
$\sum F_g\,x^{2g} = \kappa_{\mathrm{eff}}\bigl(\hat{A}(ix) - 1\bigr)$
is not an isolated substitution.  It is the surface expression
of a five-layer structure that connects the generating function
to the formal group, the modular group, the operadic decomposition,
and the chromatic filtration.  Each layer is forced by the previous.

\subsection{Layer 1: the sign absorption}

The $\hat{A}$-genus $\hat{A}(x) = (x/2)/\sinh(x/2)$ has
alternating Bernoulli coefficients.  The substitution $x \mapsto ix$
converts $\sinh(ix/2) = i\sin(x/2)$, producing
$\hat{A}(ix) = (x/2)/\sin(x/2)$.
The factor $i^{2g} = (-1)^g$ on each monomial $x^{2g}$
cancels the alternation, leaving the all-positive
Faber--Pandharipande integrals
$\lambda_g^{\mathrm{FP}} = (2^{2g-1}-1)\,|B_{2g}|/(2^{2g-1}(2g)!)$.
The coefficients decay as $\lambda_g^{\mathrm{FP}} \sim 2/(2\pi)^{2g}$;
the scalar genus series converges for $|x| < 2\pi$.  The Borel
transform is entire --- no resurgence is needed at the scalar level.

The physical content: the alternating-sign series is the index
of the Dirac operator (Lorentzian signature), while the all-positive
series is the Hodge integral (Euclidean signature).  The monograph
works in the Euclidean frame throughout; the connection to the
Witten genus (conventionally stated in the $\sinh$ convention)
requires un-Wick-rotating.

\subsection{Layer 2: the Eisenstein structure at genus 1}

On an elliptic curve $E_\tau$, the genus-0 bar kernel
$1/(z-w)$ is replaced by the Weierstrass $\zeta$-function
$\zeta(z{-}w|\tau)$.  The propagator expansion
$h(u;\tau) = u^{-1} + \sum_{k \ge 1} c_k(\tau)\,u^{2k-1}$
has coefficients
\[
c_k(\tau) = (-1)^k \frac{(2\pi)^{2k}}{(2k)!}\,B_{2k}\,E_{2k}(\tau).
\]
Each OPE pole of even order $2k$ generates a genus-1 correction
weighted by $E_{2k}(\tau)$.  The propagator has only odd powers
of~$u$ (by the oddness of~$\theta_1$), so odd-order poles contribute
nothing.

For Heisenberg (max pole~$2$): only $E_2$ enters.  The
Coisson bracket $c_0 = \{a_{(0)}b\}$ vanishes, so the
quasi-periodicity defect of~$E_2$ acts in the Cartan channel
alone --- curvature and braiding decouple.

For Virasoro (max pole~$4$): $E_2$ from the double pole and
$E_4$ from the quartic pole.  The Coisson bracket is nonzero,
so the $E_2$ defect couples curvature to braiding --- the
two structures entangle.  This is the genus-1 avatar of
the curvature-braiding open question
(Conclusion,~\S\ref{sec:conclusion}, item~(5)).

For $\cW_3$ (max pole~$6$): $E_2$, $E_4$, $E_6$.
The truncation by pole order is the genus-1 shadow of the
pole-order classification (classes $\mathbf{G}/\mathbf{L}/\mathbf{C}/\mathbf{M}$).

The quasi-modularity of $E_2$ is the Arakelov defect: under
$\tau \mapsto -1/\tau$, $E_2(-1/\tau) = \tau^2 E_2(\tau) + 12\tau/(2\pi i)$.
The anomalous term is proportional to $\kappa(\cA)$ and is
the genus-1 bar curvature $d^2 = \kappa \cdot \omega_1$.  The
non-holomorphic completion $\hat{E}_2 = E_2 - 3/(\pi\,\mathrm{Im}\,\tau)$
restores modularity at the cost of holomorphicity --- the
holomorphic anomaly equation.

\subsection{Layer 3: the genus tower as Wick rotation anomaly}

This is the central reformulation.

The genus-1 partition function
$Z_1^{\mathrm{eff}}(\tau) = \eta(\tau)^{-\kappa_{\mathrm{eff}}}$
transforms under the modular $S$-transform $\tau \mapsto -1/\tau$
as $Z_1(-1/\tau) = (-i\tau)^{-\kappa_{\mathrm{eff}}/2}\,Z_1(\tau)$.
The modular weight is $-\kappa_{\mathrm{eff}}/2$.

\begin{itemize}
\item $\kappa_{\mathrm{eff}} = 0$ ($c = 26$):
  the weight vanishes, $Z_1$ is $S$-invariant, the genus tower
  is identically zero, and the Wick rotation is an exact symmetry.
\item $\kappa_{\mathrm{eff}} \ne 0$:
  the weight is nonzero, $Z_1$ has a modular anomaly,
  the genus tower $F_g \ne 0$ at every genus, and the Wick
  rotation is broken.  The degree of breaking is measured by
  $F_g = \kappa_{\mathrm{eff}} \cdot \lambda_g^{\mathrm{FP}}$.
\end{itemize}

The Cardy formula is the composition of the $S$-transform
(which plays the structural role of Wick rotation at genus~$1$;
see Layer~5 for the precise distinction) with saddle-point extraction:
$S_{\mathrm{Cardy}} = 2\pi\sqrt{ch/6}$.  The quartic pole
produces $\kappa$, the curvature produces $d^2 = \kappa \cdot \omega_1$,
the Wick rotation anomaly produces $F_g$, the Cardy extraction
produces the entropy.  The chain is:
\[
\text{quartic pole}
\;\longrightarrow\;
\kappa
\;\xrightarrow{\;+\,\kappa_{\mathrm{ghost}}\;}
\kappa_{\mathrm{eff}}
\;\longrightarrow\;
d^2 = \kappa_{\mathrm{eff}} \cdot \omega_g
\;\longrightarrow\;
F_g \ne 0
\;\longrightarrow\;
S_{\mathrm{Cardy}}.
\]

The Witten genus is the elliptic repair of this anomaly.
The $q^0$ piece $\hat{A}(ix)$ has anomalous modular weight
$-\kappa(\cA)/2$ (intrinsic; $-\kappa_{\mathrm{eff}}/2$ after
ghost subtraction); the higher $q$-powers (the Eisenstein corrections
from Layer~2) exactly cancel the anomaly, producing a
modular form --- the Witten genus.  The Eisenstein corrections
are the Wick rotation repair terms.

\subsection{Layer 4: formal groups and the chromatic boundary}

The spectral-parameter composition law of the bar complex is
the additive formal group $\hat{\mathbb{G}}_a$, with formal
group law $F(x,y) = x + y$.  The endomorphism ring is
$\End(\hat{\mathbb{G}}_a) = \C$, and the Wick rotation
$x \mapsto ix$ is the endomorphism $[i]$.

Upon spectralization (the passage from chain complexes over~$\C$
to spectra), the formal group changes:
at height~1, $\End(\hat{\mathbb{G}}_m) = \Z_p$;
at height~$\le 2$, $\End(\hat{E}_\tau)$ is an order in a
quadratic field or quaternion algebra.  These arithmetic
endomorphism rings generically exclude~$i$:
the Wick rotation does not survive the passage from
height~$0$ (rational) to positive height (spectral).

\emph{There is no spectral Wick rotation.}  The substitution
$x \mapsto ix$ preserves $\hat{\mathbb{G}}_a$ (additive) but not
$\hat{\mathbb{G}}_m$ (multiplicative) or $\hat{E}_\tau$ (elliptic).
This is the mechanism behind the anti-pattern AP-CHR: the classical
bar-cobar theory is chromatic height~$0$, and the Wick rotation lives
in this classical setting.  Chromatic height~$\ge 1$ enters only after
spectralization, which is the fundamental open problem (Stage~3 of
the four-stage pipeline).

\subsubsection*{Five obstructions to spectralization}

The following are open problems, not conjectures.  Each identifies
a specific obstruction to Stage~3 of the pipeline
$\cA \to B_{\mathrm{mod}}(\cA) \to Z_\cA \to Z_\cA^{\mathrm{sp}}
\to Z_{\cA,n}$.

\begin{enumerate}
\item \textbf{Base ring problem.}
  $B(\cA)$ is a chain complex over~$\C$.  Spectra are not
  $\C$-modules: one needs an integral form $B(\cA)_\Z$, a chain
  complex over~$\Z$ with $B(\cA)_\Z \otimes_\Z \C \simeq B(\cA)$.
  For Kac--Moody algebras at integral level, such forms exist
  (Garland).  For general chiral algebras, no canonical integral
  form is known.

\item \textbf{Operadic coherence.}
  $\SCchtop$ involves $\C$-linear algebraic data: FM
  compactification of configurations in~$\C$, logarithmic weight
  forms, residue pairings.  Lifting to an operad in spectra
  requires understanding the interaction between holomorphic
  structure and stable homotopy theory.

\item \textbf{Formal group identification.}
  The genus-$1$ bar differential (Movement~I) uses the formal
  group of~$E_\tau$, but as complex-analytic functions
  (Weierstrass~$\zeta$), not as formal group laws over~$\Z$.
  Connecting to $\mathrm{tmf}$ requires the formal group law
  over $\Z[c_4, c_6, \Delta^{-1}]$.

\item \textbf{No genus tower maps.}
  There are no natural maps
  $\overline{\mathcal{M}}_g \to \overline{\mathcal{M}}_{g-1}$.
  The genus tower is an expansion indexed by~$g$, not a sequential
  limit.  The chromatic convergence analogy (limit over
  $E(n)$-localizations) has no direct geometric counterpart
  on the moduli side.

\item \textbf{Height discrepancy.}
  Non-separating clutching~$D_1$ raises $p$-divisible-group height
  by~$2$ ($\mathrm{ht} = 2 \cdot \dim$ for abelian varieties).
  Chromatic redshift ($\mathrm{THH}$/$\mathrm{TAQ}$) raises
  chromatic height by~$1$.  The factor-of-$2$ mismatch is
  unexplained.
\end{enumerate}

The sole exception: the elliptic curve $E_{i}$ with $j = 1728$
has $\mathrm{CM}$ by $\Z[i]$, so its formal group admits $[i]$.
At this single $j$-invariant, Wick rotation is compatible with
height~$\le 2$.  Whether this plays a role in the spectralization
problem is unknown.

\subsection{Layer 5: operadic and the three signatures}

The HT theory on $\C_z \times \R_t$ has the Swiss-cheese operad
$\SCchtop$ with operation spaces
$\FM_k(\C) \times \Conf_k^{<}(\R)$.  The $\C$-direction is
holomorphic; the $\R$-direction is ordered.

Three signatures, one operad:
\begin{itemize}
\item \textbf{Lorentzian $(2,1)$}: the light-cone constraint
  $|{\Delta z}| = |{\Delta t}|$ is $Q$-exact in the HT gauge,
  so the causal configuration space is homeomorphic to
  $\FM_k(\C) \times \Conf_k^{<}(\R)$.  The Lorentzian and HT
  operads are the same.
\item \textbf{Euclidean $(3,0)$}: the $E_3$ operad
  $C_*(\FM_k(\R^3))$, one-coloured, fully symmetric.  The passage
  $\SCchtop \to E_3$ forgets the colouring and the ordering:
  it is a forgetful functor, not a deformation.
\item \textbf{Ultra-Lorentzian $(1,2)$}: one spatial direction,
  no complex structure, outside the HT framework.
\end{itemize}

The $R$-matrix $R(z)$ is the Wick rotation datum: it mediates
the descent from the ordered bar complex
$\barB^{\mathrm{ord}}(\cA)$ on $\Conf_k^{<}(\R)$ to the
unordered bar complex on $\Conf_k(\C)$, via the
$\Sigma_k$-covering.  Without the $R$-matrix, the descent
loses information.

\emph{The $S$-transform is not the Wick rotation.}  The modular
$S$-transform $\tau \mapsto -1/\tau$ acts on the \emph{closed}
colour (the complex structure of~$E_\tau$); the Wick rotation
of~$\R_t$ acts on the \emph{open} colour.  They share the
Casimir datum $c_0 = \{a_{(0)}b\}$ as coupling constant but
are operadically orthogonal: the identification of~$S$ with
Wick rotation is a $2$d worldsheet perspective that does not
survive the $3$d HT lift.

\subsection{Open questions from this architecture}

\begin{enumerate}
\item \textbf{The Mellin/Laplace/Wick triangle.}
  Three integral transforms operate along the topological
  direction~$\R$:
  \begin{itemize}
  \item The \emph{Laplace transform}
    $r(z) = \int_0^\infty e^{-\lambda z}
    \{{\cdot}_\lambda{\cdot}\}\,d\lambda$
    converts the $\lambda$-bracket (closed colour) to the
    $R$-matrix (open colour).  It is the additive Fourier
    transform on $(\R_+,{+})$.
  \item The \emph{Mellin transform}
    $\tilde{A}(\Delta,z,\bar z)
    = \int_0^\infty \omega^{\Delta-1}
    A(\omega,z,\bar z)\,d\omega$
    converts energy to conformal dimension.
    It is the multiplicative Fourier transform on
    $(\R_+,{\times})$.
  \item The change of variables $\omega = e^t$ converts
    Mellin in~$\omega$ to bilateral Laplace in~$t$.
    The Wick rotation $t \mapsto -it$ corresponds to
    $z \mapsto iz$ in the additive Laplace variable,
    converting exponential decay (Euclidean convergence)
    to oscillation (Lorentzian causality).
  \end{itemize}
  The celestial OPE \emph{does} reduce to the bar
  differential: this is already proved
  (Theorems~\ref*{thm:ch-core-celestial-gluon-ope}
  and~\ref*{thm:ch-core-celestial-graviton-ope}).
  The passage from the HT boundary to the celestial
  boundary is the topological integration
  $\int_0^\infty dt$ (the $E_1$-coinvariant projection),
  \emph{not} a Wick rotation.  But the Wick rotation and
  the celestial transfer meet at $c = 26$: the genus tower
  vanishes (Wick-invariant) and the effective scalar partition
  function trivialises.  In the BMS limit $c \to \infty$,
  the Wick rotation anomaly $\kappa_{\mathrm{eff}} \to \infty$
  is maximal: the BMS algebra inherits the strongest
  possible breaking of Wick-rotation symmetry.
\item \textbf{The curvature-braiding criterion.}  At genus~$1$,
  the Coisson bracket $c_0 = \{a_{(0)}b\}$ controls whether
  curvature and braiding decouple ($c_0 = 0$: Heisenberg,
  lattice VOAs) or entangle ($c_0 \ne 0$: Kac--Moody, Virasoro).
  Does this dichotomy persist at genus~$g \ge 2$, with the Fay
  trisecant identity playing the role of the Arnold relation?
\item \textbf{Borel summability of the shadow obstruction tower.}
  The scalar genus tower converges.  The shadow obstruction tower for
  class~$\mathbf{M}$ at $|c| < 6$ diverges geometrically
  with growth rate~$\rho > 1$.  Is this divergent series Borel
  summable?  The expected Stokes structure: in the Euclidean
  frame, the Borel singularities lie on the real axis at
  $\xi = 1/\rho$; the Wick rotation moves them to the imaginary
  axis, clearing the integration contour.
  The modular $S$-transform at the $q \to 1$ boundary
  plays the structural role of Borel resummation.
\item \textbf{Monstrous moonshine at rank 24.}
  For a rank-24 lattice VOA: $\kappa = 24$, $F_1 = 1$,
  $Z_1(\tau) = \eta(\tau)^{-24} = 1/\Delta(\tau)$.  For the
  Leech lattice: $Z_1(V_{\mathrm{Leech}};\tau)
  = j(\tau) - 720$ (exact).  The discriminant $\Delta
  = \eta^{24}$ generates the 576-fold periodicity of $\mathit{tmf}$.
  Does the bar complex at rank~24 encode any of the Monster
  moonshine structure beyond the rational shadow?
\item \textbf{The algebraic-analytic gap.}
  The monograph works entirely on the algebraic side:
  Koszulness (a homological condition on the bar complex)
  replaces unitarity (a metric condition on the Hilbert space)
  as the structural input.  Reflection positivity requires a
  $*$-structure compatible with the OPE, which is not formulated
  here.  The order-reversal involution
  $R_A \colon B^{\mathrm{ch}}(A^{\mathrm{op}})
  \xrightarrow{\sim} B^{\mathrm{ch}}(A)^{\mathrm{cop}}$
  provides the categorical backbone of time-reversal, but the
  \emph{positivity} condition $\langle \Omega, R_A(\cdot)\rangle
  \ge 0$ is invisible to the bar complex alone.  The Verdier
  pairing (pairing $A$ with $A^!$) and the BPZ inner product
  (pairing $A$ with itself) are distinct structures: they
  coincide only at the self-dual point $c^* = \alpha_N/2$.
  Koszulness and unitarity are logically independent:
  $\mathrm{Vir}_c$ at $c < 0$ is Koszul but non-unitary, while
  unitary non-formal algebras exist.  The Segal comparison
  conjecture (algebraic factorization algebra $\to$ geometric
  functorial QFT) and the Osterwalder--Schrader reconstruction
  (Euclidean $\to$ Lorentzian) remain the open interfaces between
  the algebraic and analytic frameworks.  Whether the Koszul locus
  is contained in the Kontsevich--Segal convergence domain is an
  unresolved question at this interface.
\end{enumerate}


% =================================================================
\section{Trefoil verification: Kontsevich integral from the bar complex}
\label{sec:trefoil-bar-verification}
% =================================================================

The spectral $R$-matrix $R(z) = \exp(\hbar k\,\Omega/z)$ of
$\widehat{\fsl}_2{}_k$
(Proposition~\ref{prop:affine-r-mode} of spectral-braiding-core)
recovers the Kontsevich integral of the trefoil $3_1 =
\widehat{\sigma_1^3}$ at degrees $\le 3$.  This verifies
Example~\ref{ex:trefoil-verification} of spectral-braiding-frontier
with explicit formulas.

\subsection*{R-matrix eigenvalues on $V = \C^2$}

The $\fsl_2$ Casimir on $V \otimes V$ decomposes as
$\Omega|_{\Sym^2(V)} = 1/4$ and
$\Omega|_{\Lambda^2(V)} = -3/4$, so the braid closure
$\operatorname{tr}_q(R^3)$ with $R^3 = \exp(3\hbar k\,\Omega/z)$
gives
\[
  Z_V(3_1)
  \;=\;
  3\,e^{3\hbar k/(4z)} + e^{-9\hbar k/(4z)},
\]
where the multiplicities $3 = \dim\Sym^2(V)$,
$1 = \dim\Lambda^2(V)$ arise from the quantum trace at leading
order.

\subsection*{Degree-by-degree matching}

\begin{center}
\renewcommand{\arraystretch}{1.3}
\begin{tabular}{ccl}
\textbf{Degree} & \textbf{Coefficient} & \textbf{Source} \\
\hline
$1$ & $z_1 = \tfrac{3}{2}\,\theta$ &
  Writhe $w = 3$, each crossing contributes $\pm 1/2$ \\
$2$ & $z_2 = \tfrac{1}{24}\,\theta$ &
  $B_2/4 = 1/24$; constant term of $G_2(\tau)$ \\
$3$ & $v_3 = 1$ &
  IHX from Arnold; $w_3 = \sum f^{abc}f^{abc} = 8$ for $\fsl_2$
\end{tabular}
\end{center}

\noindent
The degree-$1$ term is pure framing ($z_1 = w/2$, removed by
framing correction).  The first nontrivial Vassiliev invariant
is $z_2(3_1) = 1/24$.

\subsection*{The coefficient $1/24$}

The coincidence of $1/24$ in three contexts --- the degree-$2$
trefoil invariant, the genus-$1$ free energy
$F_1 = -\kappa\,\log\,\eta(\tau)/24$, and the constant term
$G_2(\tau) = -1/24 + \cdots$ --- traces to a common source:
the curved bar structure $d_{\barB}^2 = \kappa(\cA)\,\omega_1$
at genus~$1$, whose perturbative coefficient in the Eisenstein
basis is $B_2/4 = 1/24$.  The trefoil, as the closure of
$\sigma_1^3$, probes the minimal nontrivial configuration on
$\FM_2(\C) \times \Conf_2(\R)$.

On the $\fsl_2$ fundamental weight system:
$w_{\fsl_2}(z_2 \cdot \theta) = \frac{1}{24} \cdot
C_2(\mathrm{fund}) = \frac{1}{24} \cdot \frac{3}{4} =
\frac{1}{32}$, matching the $\hbar^2$-coefficient of
$J_{3_1}(e^\hbar)$.

\subsection*{Degree~$3$: IHX from Arnold}

For the affine lineage, one-loop exactness
(Theorem~\ref{thm:one-loop-koszul}) forces $m_3 = 0$.
The degree-$3$ bar differential involves only pairwise
collisions, landing in the chord diagram sector.  The Arnold
relation
$\alpha_{12} \wedge \alpha_{23} + \alpha_{23} \wedge
\alpha_{31} + \alpha_{31} \wedge \alpha_{12} = 0$
on $H^\bullet(\Conf_3(\C))$, with $\fg$-valued Casimir
decorations, becomes $D_I + D_H - D_X = 0$: the IHX
relation is \emph{derived} from FM geometry, not imposed.

\subsection*{Status}

The matching at degrees $\le 3$ is unconditional for the
$\fsl_2$ weight system: it uses the proved spectral $R$-matrix
(Corollary~\ref{cor:jones-polynomial}), proved one-loop
exactness, and proved Arnold cancellation
(Theorem~\ref{thm:Arnold_full_proof}).  The full identification
$Z_{A_\infty}(K) = Z(K)$ at all orders remains conjectural
(Theorem~\ref{thm:bar-kontsevich}).  At degree~$4$, the chord
diagram algebra $\cA_4(\fsl_2)$ has dimension~$3$; the symmetry
factor matching between FM boundary strata and chord diagram
multiplicities becomes the first genuinely nontrivial
combinatorial test.

\begin{question}[Beyond the affine lineage]
\label{q:trefoil-beyond-affine}
For non-formal Swiss-cheese structures ($m_3 \ne 0$, e.g.\
Virasoro), the degree-$3$ bar differential includes trivalent
vertices.  The Kontsevich integral, built from pairwise $d\log$
forms, sees only the chord diagram quotient.  Capturing the
higher $m_k$ requires the full $A_\infty$ partition function
$Z_{A_\infty}(K)$.  The degree-$3$ divergence between the
chord diagram framework and the $A_\infty$ framework is
the first point where class~L and class~M separate in the
knot invariant.
\end{question}


% =================================================================
\section{$\fsl_3$ monodromy identification: first rank-2 test}
\label{sec:sl3-monodromy-extraction}
% =================================================================

% Extracted from the expanded computation in
% .claude/worktrees/agent-acb59fb3/archive/source_tex/working_notes.tex
% (~740 lines). What follows is the rectified distillation.

\subsection*{Setup}
Let $\fg = \fsl_3$, $h^\vee = 3$, $\rho = \omega_1 + \omega_2$.
Inner product: $(\omega_i, \omega_j) = (A^{-1})_{ij}$ with
$A^{-1} = \frac{1}{3}\bigl(\begin{smallmatrix} 2 & 1 \\ 1 & 2
\end{smallmatrix}\bigr)$.
Quadratic Casimir: $c_2(\lambda) = (\lambda, \lambda + 2\rho)$.
Cubic Casimir: $c_3(p,q) = \tfrac{1}{18}(p-q)(2p+q+3)(p+2q+3)$;
note $c_3(p,q) = -c_3(q,p)$ (distinguishes a rep from its contragredient).

\subsection*{Test 1: braiding on $V \otimes V$}
$V = V_{(1,0)} = \C^3$.  Clebsch--Gordan:
$V \otimes V \cong V_{(2,0)} \oplus V_{(0,1)}$.
Split Casimir eigenvalues (via
$\Omega|_{V_\nu} = \tfrac{1}{2}(c_2(\nu) - 2c_2(\lambda))$):
$\Omega|_{V_{(2,0)}} = 2/3$, $\Omega|_{V_{(0,1)}} = -4/3$.
Steinberg connection: $\nabla_{\mathfrak{S}} = d + \Omega\,dz/((k+3)z)$.
Monodromy ($q = e^{i\pi/(k+3)}$):
\[
  \mathrm{Mon}\big|_{V_{(2,0)}} = q^{4/3},
  \qquad
  \mathrm{Mon}\big|_{V_{(0,1)}} = q^{-8/3}.
\]
These match $U_q(\fsl_3)$ exactly ($\check{R}$-matrix eigenvalues
up to permutation sign on $\Lambda^2 V$).
The cubic Casimir does \emph{not} appear: $C_2$ alone distinguishes
the two summands.  Same mechanism as $\fsl_2$.

\subsection*{Test 2: associator on $V^{\otimes 3}$ --- the cubic Casimir enters}
$V^{\otimes 3} \cong V_{(3,0)} \oplus 2 \cdot V_{(1,1)}
\oplus V_{(0,0)}$ ($\mathbf{10} \oplus 2\cdot\mathbf{8} \oplus \mathbf{1}$).
The two copies of $V_{(1,1)}$ arise from intermediate channels:
$\mathbf{8}_{\mathrm{sym}}$ (via $V_{(2,0)}$) and
$\mathbf{8}_{\mathrm{anti}}$ (via $V_{(0,1)}$) in the $(12)$-grouping.

The Drinfeld--KZ associator at order $\hbar^3 = 1/(k+3)^3$
involves the antisymmetric double commutator
$\mathscr{A} = [\Omega_{12}, [\Omega_{12}, \Omega_{23}]]
- [\Omega_{23}, [\Omega_{23}, \Omega_{12}]]$.
By the Fierz identity for $\fsl_3$,
\[
  f_{ace}\,f_{bde}
  = \tfrac{2}{3}(\delta_{ab}\delta_{cd} - \delta_{ad}\delta_{bc})
  + d_{ace}\,d_{bde} - d_{ade}\,d_{bce},
\]
the double commutator decomposes into a Casimir piece
(present for all $\fg$) and a $d$-tensor piece
$\cD_{123}^{(-)}$ (vanishes for $\fsl_2$, nonzero for $\fsl_3$).
The $d$-tensor piece produces \textbf{off-diagonal mixing}
between the two copies of $V_{(1,1)}$, proportional to
\[
  \Delta c_3 = c_3(2,0) - c_3(0,1)
  = \tfrac{35}{9} + \tfrac{10}{9} = 5.
\]
This is the first genuinely higher-rank phenomenon:
the cubic Casimir enters the monodromy through the
\emph{associator}, not through the $R$-matrix eigenvalues.

\subsection*{Verification and conclusion}
By one-loop exactness of $V^k(\fsl_3)$ at generic $k$
(Theorem~\ref{thm:one-loop-koszul}), the bar-complex reduced
connection \emph{is} the KZ connection ($m_k = 0$ for $k \geq 3$).
Therefore the bar-complex associator reproduces the Drinfeld--KZ
associator identically, including the cubic Casimir contribution.
By Drinfeld--Kohno, the monodromy identifies with $U_q(\fsl_3)$.

\textbf{Summary:}  The $\fsl_3$ test confirms
Theorem~\ref{thm:affine-monodromy-identification} at rank 2.
$C_2$ controls the diagonal part of the monodromic structure
(braiding eigenvalues).  $C_3$ controls the off-diagonal part
(multiplicity-space mixing in the associator).  For $\fsl_2$,
$C_3 = 0$ and the off-diagonal channel is absent.


% =================================================================
\section{Genus-$1$ curved bar complex: rectified summary}
\label{sec:genus1-bar-rectified}
% =================================================================
%
% Extracted and rectified from the genus-1 worktree computation.
% Source: worktree agent-a696cf24, ~810 lines.
% Three issues corrected: (i) floor-function precedence in the
% general weight formula, (ii) garbled intermediate in the
% normalisation verification, (iii) loose ``Tr(L_0)'' language
% for the modular characteristic.

\subsection*{The genus-$1$ propagator}

On $E_\tau = \C/(\Z + \Z\tau)$, the Arakelov-normalised
propagator is $\eta^{(1)}(u;\tau) = h(u;\tau) + R(u;\tau)$, where
$h(u;\tau) = \partial_u\log\vartheta_1(u|\tau)$ is meromorphic
and $R(u;\tau) = 2\pi i\,\mathrm{Im}(u)/\mathrm{Im}\tau$ is the
smooth Arakelov correction.  The Laurent expansion of $h$ near
$u = 0$ has only odd powers:
\[
h(u;\tau)
\;=\;
\frac{1}{u}
+ \sum_{k=1}^{\infty} c_k(\tau)\,u^{2k-1},
\qquad
c_k(\tau) = (-1)^k\frac{(2\pi)^{2k}}{(2k)!}\,B_{2k}\,E_{2k}(\tau),
\]
with $c_1 = -(\pi^2/3)E_2$, $c_2 = -(\pi^4/45)E_4$,
$c_3 = -(2\pi^6/945)E_6$.

\subsection*{Arnold relation failure at genus $1$}

At genus~$0$, the Arnold relation
$\sum_{\mathrm{cyc}}\eta^{(0)}_{12}\wedge\eta^{(0)}_{23} = 0$
ensures $\dfib^2 = 0$.  At genus~$1$, decompose $\eta^{(1)} = h + R$:
\begin{enumerate}[label=(\alph*)]
\item $\sum_{\mathrm{cyc}} h_{12}\wedge h_{23} = 0$ by the Fay
  trisecant identity (the holomorphic Arnold relation holds on the
  universal cover).
\item $\sum_{\mathrm{cyc}} R_{12}\wedge R_{23}$: smooth, no
  residue contribution.
\item The cross-terms $h\wedge R + R\wedge h$ produce the Arnold
  defect:
\end{enumerate}
\[
\sum_{\mathrm{cyc}}
\eta^{(1)}_{12}\wedge\eta^{(1)}_{23}
\;=\;
\omega_1,
\qquad
\omega_1
= \frac{i}{2\,\mathrm{Im}\tau}\,dz\wedge d\bar{z},
\qquad
\textstyle\int_{E_\tau}\omega_1 = 1.
\]
The Arnold defect is the Arakelov $(1,1)$-form
$\omega_1$.  This is the genus-$1$ source of curvature.

\subsection*{The fibre differential on $[s^{-1}T|s^{-1}T]$}

The residue of $\eta^{(1)}\cdot T(z_1)T(z_2)$ at $u = 0$ gives:
\[
\dfib^{(1)}[s^{-1}T|s^{-1}T]
\;=\;
s^{-1}\!\Bigl(
  {:}TT{:}
  \;-\; \tfrac{2\pi^2}{3}\,E_2(\tau)\cdot T
  \;-\; \tfrac{c\pi^4}{90}\,E_4(\tau)\cdot\mathbf{1}
\Bigr).
\]
The three terms: genus-$0$ normal-ordered product (from $1/u$ times
${:}TT{:}$); the $E_2$-correction (from $c_1 u$ times $2T/u^2$);
the $E_4$-contact term (from $c_2 u^3$ times $(c/2)/u^4$).

For general weights $h_a, h_b$:
\[
\dfib^{(1)}[s^{-1}a|s^{-1}b]
\;=\;
s^{-1}\!\Bigl(
  a_{(-1)}b
  + \sum_{k=1}^{\lfloor(h_a+h_b)/2\rfloor}
    c_k(\tau)\cdot a_{(2k-1)}b
\Bigr),
\]
where the sum terminates because $a_{(n)}b = 0$ for
$n \ge h_a + h_b$.  The odd OPE products $a_{(2k-1)}b$
pair with the odd propagator coefficients $c_k(\tau)u^{2k-1}$.
Each even-order OPE pole has a genus-$1$ Eisenstein shadow.

\subsection*{Curvature: $\dfib^{(1),2} = \kappa\cdot\omega_1$}

Decompose $\eta^{(1)} = h + R$.  In the double residue computing
$\dfib^2$: (a) $h\cdot h$ vanishes by Fay trisecant; (b) $R\cdot R$
vanishes by smoothness; (c) the cross-terms $h\cdot R + R\cdot h$
produce the curvature.  The cross-term extracts:
\begin{itemize}
\item from $h$: the leading OPE coefficient via residue,
\item from $R$: the Arakelov class $\bar\partial R
  = -\omega_{\mathrm{Ar}}$,
\end{itemize}
with the coefficient given by the modular characteristic
$\kappa(\mathrm{Vir}_c) = c/2$
(defined via zeta-regularised partition function, not by the
literally divergent trace $\mathrm{Tr}(L_0)$).  Result:
\[
\dfib^{(1),2}[s^{-1}a]
= \frac{c}{2}\cdot\omega_1\cdot [s^{-1}a].
\]
The higher OPE poles contribute to the genus-$1$ \emph{differential}
(the $E_2$, $E_4$ corrections) but not to the \emph{curvature}
$d^2$.  The curvature sees only the modular characteristic.

\subsection*{Genus-$1$ free energy and the critical dichotomy}

$F_1^{\mathrm{intr}} = (c/2)\cdot(1/24) = c/48$;
$F_1^{\mathrm{eff}} = ((c{-}26)/2)\cdot(1/24) = (c{-}26)/48$.
On a fixed torus: $Z_1(\tau) = \eta(\tau)^{-\kappa}$.  The
$\hat{A}$-genus check:
$\hat{A}(ix) - 1 = (1/24)x^2 + \cdots$ gives
$F_1 = \kappa/24$.  Confirmed.

\medskip
\begin{center}
\renewcommand{\arraystretch}{1.15}
\begin{tabular}{@{}lcc@{}}
& $c \neq 26$ & $c = 26$ \\
\hline
$\kappa_{\mathrm{eff}}$ & $(c{-}26)/2 \neq 0$ & $0$ \\
$F_1^{\mathrm{eff}}$ & $(c{-}26)/48 \neq 0$ & $0$ \\
$\dfib^{(1),2}$ & $((c{-}26)/2)\,\omega_1$ & $0$ \\
Genus tower & nonzero at all $g$ & vanishes identically \\
Regime & curved bar & flat bar
\end{tabular}
\end{center}

\subsection*{What is genuinely new}

\begin{enumerate}
\item \textbf{The Arnold defect mechanism.}  The genus-$1$ Arnold
  relation does not fail wholesale: the holomorphic summand
  $h\wedge h$ still cancels (Fay trisecant).  The defect comes
  entirely from the non-holomorphic Arakelov correction $R$,
  whose $\bar\partial$-inexactness produces the $(1,1)$-form
  $\omega_1$.  This is a controlled failure: the defect is a
  \emph{scalar} times $\omega_1$, not a matrix-valued form.
\item \textbf{Eisenstein shadows of OPE poles.}  The genus-$1$
  propagator pairs odd powers $c_k u^{2k-1}$ with even-order OPE
  poles $a_{(2k-1)}b/u^{2k}$, producing Eisenstein-weighted
  corrections.  Each pole order $2k$ in the OPE casts a genus-$1$
  shadow $E_{2k}(\tau)$.  The double pole $\to E_2$ (quasi-modular,
  one-loop self-energy); the quartic pole $\to E_4$ (modular,
  contact term).
\item \textbf{Curvature vs.\ differential separation.}  Higher
  OPE poles affect the differential but not the curvature.  The
  curvature $d^2 = \kappa\cdot\omega_1$ depends only on $\kappa$
  (the modular characteristic), not on the detailed pole structure.
  This is the genus-$1$ manifestation of the universal
  formula $d^2 = \kappa\cdot\omega_g$.
\end{enumerate}


%% ===================================================================
%%  CELESTIAL AMPLITUDES FROM THE BAR COMPLEX -- RECTIFIED SUMMARY
%% ===================================================================

\section{Celestial amplitudes from the bar complex: rectified summary}
\label{sec:celestial-bar-rectified}

Extraction and audit of the celestial amplitude computation from the
worktree draft (agent-a6f30e7e, ${\sim}715$ lines).  All results verified
against the core chapters: Theorem~\ref{thm:ch-core-celestial-gluon-ope}
(gluon OPE), Theorem~\ref{thm:ch-core-celestial-graviton-ope} (graviton
OPE), Theorem~\ref{thm:cbt-mhv-bar-integral} (MHV integral),
Theorem~\ref{thm:ch-core-helicity-splitting} (helicity splitting), and
Remark~\ref{rem:ch-core-colour-ordered-two-colour} (colour-ordered =
ordered-bar).

\subsection*{Audit of the collision residue computation}

\begin{enumerate}[label=(\roman*),nosep]
\item \textbf{Gluon.}  The $\lambda$-bracket $\{J^a{}_\lambda J^b\} =
  f^{ab}{}_c J^c + k\kappa^{ab}\lambda$ Laplace-transforms to
  $r^{\mathrm{Lap}}(z) = fJ^c/z + k\kappa/z^2$.  After $d\log$
  absorption (AP19: $z^{-n}\mapsto z^{-(n-1)}$), the structure-constant
  term drops to $z^0$ and the level term survives at $z^{-1}$:
  $r^{\mathrm{KM}}(z) = k\Omega/z$.  Colour-ordered projection of the
  Casimir yields $f^{ab}{}_c$.  \emph{Matches Fan--Fotopoulos--Taylor.}

\item \textbf{Graviton.}  The Virasoro $\lambda$-bracket
  $\{T{}_\lambda T\} = (\partial + 2\lambda)T + \frac{c}{12}\lambda^3$
  Laplace-transforms to $\partial T/z + 2T/z^2 + (c/2)/z^4$ (using
  $\int_0^\infty \lambda^3 e^{-\lambda z}d\lambda = 6/z^4$, so
  $(c/12)(6/z^4) = (c/2)/z^4$).  After $d\log$ absorption:
  $r^{\mathrm{Vir}}(z) = (c/2)/z^3 + 2T/z$.  The $\partial T$ term
  drops.  \emph{Matches Pasterski--Shao--Strominger.}  The two poles
  are the Weinberg soft factor ($r_1 = 2T$) and the gravitational
  anomaly ($r_3 = c/2$).

\item \textbf{MHV at $n=4$.}  The bar-complex integral
  $\int_{\overline{\mathrm{FM}}_4(\C)} \omega_4^{(+,-,-,+)} \wedge \nu_4$
  localises to three channels $D_{\{1,2\}}, D_{\{2,3\}}, D_{\{3,4\}}$.
  Each $d\log(z_i - z_j)$ residue produces $\langle ij\rangle^{-1}$ via
  the spinor-helicity identification $z_i - z_j \sim \langle ij\rangle /
  (\lambda_i^2\lambda_j^2)$.  The anti-holomorphic pairing of the two
  negative-helicity insertions supplies the numerator
  $\langle 23\rangle^4$.  Result: Parke--Taylor.
\end{enumerate}

\subsection*{Colour-ordered structure and the ordered bar}

Colour-ordered celestial amplitudes are ordered-bar objects
(Remark~\ref{rem:ch-core-colour-ordered-two-colour}): they live on
$\Conf_n^{<}(\C)$, not the symmetric quotient.  The Parke--Taylor
denominator $\prod_{i=1}^n \langle i,\, i{+}1\rangle^{-1}$ is the
$d\log$ residue on ordered configurations --- the holomorphic colour
of the ordered bar differential.  Full amplitudes are recovered by
$R$-matrix descent: $\cM_n = \sum_{\sigma \in \Sigma_n/\mathrm{cyc}}
R_\sigma \cdot \cM_n^{\mathrm{ord}}(\sigma)$.  This is the celestial
manifestation of the two-colour architecture of Part~\ref{part:ordered},
consistent with the $E_\infty$ locality of $\widehat{\fg}_k$: the
ordered bar carries strictly more data than the symmetric bar when OPE
poles are present (AP38), and the $R$-matrix descent is genuinely
twisted, not naive (AP36).

\subsection*{Frontier questions (from worktree)}

\begin{enumerate}[label=(\alph*),nosep]
\item \textbf{NMHV and double twisting morphisms.}
  NMHV amplitudes (three negative helicities) require a double
  pairing between $\barBch(\cA)$ and $\barBch(\cA^!)$: two bulk
  propagators connecting three anti-holomorphic insertions to the
  holomorphic skeleton.  \emph{Is the BCFW recursion a statement
  about the $\Ainf$ structure of the twisting-morphism space?}

\item \textbf{Graviton MHV from $\barBch(\mathrm{Vir}_c)$.}
  At $n = 4$, the graviton MHV amplitude
  $M_4 = \langle 23\rangle^8 /
  (\langle 12\rangle^2 \langle 23\rangle^2
  \langle 34\rangle^2 \langle 41\rangle^2)$
  has doubled power relative to gluons, reflecting spin $2$.  In the
  bar complex, this doubling should arise from the quartic OPE pole
  of Virasoro (versus the quadratic pole of Kac--Moody), producing
  extra $\langle ij\rangle$ powers per collision.  \emph{Precise integral
  formula is a target for future computation.}
\end{enumerate}


% =================================================================
\section{Virasoro spectral $R$-matrix: rectified extract}
\label{sec:virasoro-spectral-R-rectified}
% =================================================================
%
% Source: worktree agent-acffc868 (~865 lines).  Rectified extract:
% verified Laplace transform, CYBE, AP19 compliance, Ponsot--Teschner
% comparison, quantum corrections.  Errata noted inline.

\subsection*{Classical $r$-matrix (verified)}

The Laplace transform of the Virasoro $\lambda$-bracket
$\{T{}_\lambda T\} = \partial T + 2T\lambda + \frac{c}{12}\lambda^3$
is (using $\int_0^\infty \lambda^n e^{-\lambda z}\,d\lambda = n!/z^{n+1}$):
\begin{equation}\label{eq:vir-rL-rectified}
r^L(z)
\;=\;
\frac{\partial T}{z}
+ \frac{2T}{z^2}
+ \frac{c/2}{z^4}.
\end{equation}
This is the \emph{Laplace kernel}: same pole orders as the OPE.
The \emph{collision residue} (bar-theoretic $r$-matrix, AP19: $d\log$
absorbs one power) is
\begin{equation}\label{eq:vir-rcoll-rectified}
r^{\mathrm{coll}}(z)
\;=\;
\frac{c/2}{z^3} + \frac{2T}{z}.
\end{equation}
The $\partial T/z$ term becomes regular after absorption and drops.
Both formulae match the established manuscript (3d\_gravity.tex
lines 1000, 3123; examples-worked.tex line 112).

\subsection*{Classical Yang--Baxter equation (verified)}

The CYBE for $r^L(z)$ follows from the PVA Jacobi identity.
Writing $I = \{T{}_\lambda\{T{}_\mu T\}\}$,
$II = \{T{}_\mu\{T{}_\lambda T\}\}$,
$III = \{\{T{}_\lambda T\}{}_{\lambda+\mu} T\}$:
\[
I - II
\;=\;
(\lambda-\mu)\,\partial T
+ 2(\lambda^2-\mu^2)\,T
+ \tfrac{c}{12}(\lambda-\mu)(\lambda+\mu)^3
\;=\; III.
\]
The central-term cancellation
$\frac{c}{12}(\lambda^4 - \mu^4 + 2\lambda\mu(\lambda^2 - \mu^2))
= \frac{c}{12}(\lambda-\mu)(\lambda+\mu)^3$
is verified by factoring.  Under Laplace transform this becomes the
CYBE with spectral parameters $(u, u+v, v)$.  The Arnold
interpretation: the three commutator terms are the three
codimension-$1$ boundary strata of $\FM_3(\C)$; Stokes forces their
sum to vanish (same mechanism as $d_{\barB}^2 = 0$).

\subsection*{Quantum corrections: two-source structure}

The quantum $R$-matrix $R(z) = \id + \hbar\,r^{(1)}(z)
+ \hbar^2\,r^{(2)}(z) + \cdots$ (with $\hbar = 12/c$) receives
corrections from two sources at each order:
\begin{enumerate}[label=(\roman*)]
\item \emph{Iterated monodromy}: path-ordered exponential of the
  classical $r$-matrix (Chen iterated integrals of $r^L$);
\item \emph{$\Ainf$ corrections}: the ternary operation
  $m_3^{\mathrm{Vir}}$ (and higher) correct the shifted KZ
  connection.
\end{enumerate}
The two sources partially cancel via the Stasheff identity
$\partial m_3 + m_2 \circ m_2 = 0$.  On the diagonal
(level-$0$ to level-$0$ component
$|h_1,h_2\rangle \to |h_1,h_2\rangle$):
\begin{equation}\label{eq:R-level0-rectified}
R(z)\big|_{\mathrm{diag}}
\;=\;
\exp\!\left(
  \frac{c/2}{z^4}\,\hbar
  + \frac{2h_1 h_2}{z^4}\,\hbar^2
  + O(\hbar^3)
\right).
\end{equation}
The leading term is the abelian phase; the $2h_1 h_2\,\hbar^2/z^4$
is the first genuine quantum correction.

\subsection*{Degenerate truncation (verified)}

On the $(2,1)$-degenerate representation $V_{h_{2,1}}$, all $\Ainf$
corrections vanish (null vectors kill $m_3, m_4, \ldots$), and $R(z)$
is computed exactly by the BPZ equation.  Monodromy eigenvalues:
$e^{\pm i\pi b\alpha_2}$, matching the Ponsot--Teschner kernel.

\subsection*{Ponsot--Teschner comparison: status}

\begin{enumerate}[label=(\arabic*)]
\item \emph{Classical limit}: $r^L(z)$ matches the $b \to 0$ saddle
  of the fusion kernel.  Verified.
\item \emph{Degenerate eigenvalues}: exact match via BPZ.  Verified.
\item \emph{First quantum correction at level~$0$}: the $2h_1 h_2/z^4$
  coefficient matches the leading term of the Zamolodchikov one-loop
  determinant $W_1$ in the heavy limit $h_i = c\epsilon_i/6$.
  Verified at leading order.
\item \emph{Full identification $R = \mathbb{F}^{\mathrm{PT}}$}:
  conjectural.  Gap: analytic continuation from integer~$k$ (affine
  DS transport) to continuous~$c$ via the Faddeev modular double.
\end{enumerate}

\subsection*{Errata in the worktree source}

\begin{enumerate}[label=(E\arabic*)]
\item The mode expansion eq.~(3.5) of $r^L_{12}(z)$ has powers
  $z^{-n-1}$ and $z^{-n}$ in the denominator that can become
  positive powers of~$z$---this is a notational error (confuses
  position-space field $T(z)$ with the tensor $r_{12}(z)$).  The
  subsequent eq.~(3.6) and level-$1$ computation bypass this and
  are correct.
\item The skew-symmetry discussion (lines 145--151 of the diff)
  claims $r(-z) + r_{21}(z) = (c/2)/z^4 + (c/2)/z^4 = 0$; the
  two central terms ADD rather than cancel, so the equation should
  read $r_{12}(-z) + r_{21}(z) = (c/z^4)\,\id \otimes \id$ (a
  scalar, hence compatible with CYBE up to coboundary).
\item The classical-limit comparison (Observation 5.1) writes
  $W_0 \supset P_1 P_2/z^2$ but then matches against $2h_1 h_2/z^4
  \cdot 1/c$; these are different $z$-powers.  The heavy-limit
  statement in Theorem~4.1(v) is more carefully formulated and
  should be preferred.
\end{enumerate}


% =================================================================
\section{Rectified summary: Virasoro $m_4$ and $m_5$}
\label{sec:virasoro-m4-m5-rectified}
% =================================================================

Two parallel computations of the Virasoro $\Ainf$ operations
$m_4$ and $m_5$ were cross-checked.  The explicit formulas agree
with Propositions~\ref*{prop:gravity-m4}--\ref*{prop:gravity-m5}
of~\S\ref*{sec:bar-knows-gravity}.  This section records the
rectified formulas with corrected structural claims, the
Shapovalov inner product derivation of the quartic contact
invariant, and the OPE of $\normord{TT}$ with $T$.

\subsection{$m_4$: forced by the Stasheff relation}

With $m_1 = \partial \ne 0$ (translation), the arity-$4$ Stasheff
relation reads $\partial(m_4) + \mathrm{Obs}_4 = 0$, where
$\mathrm{Obs}_4$ involves only $m_2$ and $m_3$.
The operation $m_4 = -\mathrm{Obs}_4$ is \emph{forced}, not a free
cocycle (the Borcherds identity ensures the full OPE is associative,
but the $\lambda$-bracket alone is non-associative---that is the
engine of the $\Ainf$ tower).

\begin{equation}\label{eq:m4-rectified}
\begin{aligned}
m_4(T^4;\,&\lambda_{12},\lambda_{23},\lambda_{34})
\\[2pt]
&=\;
(\lambda_{23} + 2\lambda_{34})\,\partial^2 T
\\[2pt]
&\quad +\;
(5\lambda_{12}^2 - 4\lambda_{12}\lambda_{23}
  + 3\lambda_{12}\lambda_{34}
  - 5\lambda_{23}^2 + 10\lambda_{23}\lambda_{34}
  + 4\lambda_{34}^2)\,\partial T
\\[2pt]
&\quad +\;
2(\lambda_{12}^3 + 2\lambda_{12}^2\lambda_{23}
  + 3\lambda_{12}^2\lambda_{34}
  - 2\lambda_{12}\lambda_{23}^2
  - 4\lambda_{12}\lambda_{23}\lambda_{34}
  + 4\lambda_{12}\lambda_{34}^2
\\
&\qquad\qquad
  - \lambda_{23}^3 - 2\lambda_{23}^2\lambda_{34}
  + 6\lambda_{23}\lambda_{34}^2)\,T
\\[2pt]
&\quad +\;
\tfrac{c}{12}
(\lambda_{12}^5 + 2\lambda_{12}^4\lambda_{23}
  + 3\lambda_{12}^4\lambda_{34}
  + 4\lambda_{12}^3\lambda_{23}\lambda_{34}
  + 2\lambda_{12}^3\lambda_{34}^2
\\
&\qquad\qquad
  - 2\lambda_{12}\lambda_{23}^4
  - 4\lambda_{12}\lambda_{23}^3\lambda_{34}
  - 4\lambda_{12}\lambda_{23}\lambda_{34}^3
  + 2\lambda_{12}\lambda_{34}^4
\\
&\qquad\qquad
  - \lambda_{23}^5 - 2\lambda_{23}^4\lambda_{34}
  + 6\lambda_{23}\lambda_{34}^4).
\end{aligned}
\end{equation}
Symmetric point:
$m_4(T^4;\,\lambda,\lambda,\lambda)
= 3\lambda\,\partial^2 T
+ 13\lambda^2\,\partial T
+ 14\lambda^3\, T
+ \tfrac{7c}{12}\,\lambda^5$.
No $\partial^3 T$ term (associahedron $K_4$ cell structure).
All $T$-dependent coefficients are $c$-independent.

\subsection{Quartic contact invariant from the Shapovalov form}

At conformal weight~$4$, the Shapovalov matrix in the basis
$(\partial^2 T,\, \normord{TT}) = (2L_{-4}|0\rangle,\,
L_{-2}^2|0\rangle)$ is
\[
S^{(4)}
=
\begin{pmatrix}
5c & 3c \\
3c & c(c+8)/2
\end{pmatrix},
\qquad
\det S^{(4)} = \frac{c^2(5c+22)}{2},
\qquad
(S^{(4)})^{-1}
=
\frac{1}{c(5c+22)}
\begin{pmatrix}
c+8 & -6 \\
-6 & 10
\end{pmatrix}.
\]
The quartic contact invariant
$\mathfrak{Q}^{\mathrm{contact}}_{\mathrm{Vir}}
= 10/[c(5c+22)]$
is the $(2,2)$-entry of $(S^{(4)})^{-1}$, with poles at
$c = 0$ and $c = -22/5$ (Lee--Yang).

\subsection{The normally ordered OPE $\{\normord{TT}{}_\lambda T\}$}

By the Dong lemma:
\begin{equation}\label{eq:TT-T-rectified}
\{\normord{TT}{}_\lambda T\}
\;=\;
-\partial^2 T\cdot\lambda
\;-\; \tfrac{5}{2}\,\partial T\cdot\lambda^2
\;-\; \tfrac{4}{3}\, T\cdot\lambda^3
\;-\; \tfrac{c}{40}\,\lambda^5.
\end{equation}
This is the homotopy data used in the HPL computation of $m_4$:
the normally ordered product provides the chain homotopy between
the full (associative) OPE and the (non-associative)
$\lambda$-bracket.

\subsection{$m_5$: explicit polynomials $P_1$ and $P_0$}

\begin{equation}\label{eq:m5-rectified}
\begin{aligned}
m_5(T^5;\,
  &\lambda_{12},\lambda_{23},\lambda_{34},\lambda_{45})
\\[2pt]
&=\;
-\partial^4 T
\;-\;
(2\lambda_{12} + 4\lambda_{23}
  + \lambda_{34} + 3\lambda_{45})\,\partial^3 T
\\[2pt]
&\quad +\;
(3\lambda_{12}^2 - 16\lambda_{12}\lambda_{23}
  - 6\lambda_{12}\lambda_{45}
  - 13\lambda_{23}^2 + 7\lambda_{23}\lambda_{34}
\\
&\qquad\qquad
  - 9\lambda_{23}\lambda_{45}
  - 7\lambda_{34}^2
  + 3\lambda_{34}\lambda_{45}
  - 3\lambda_{45}^2)\,\partial^2 T
\\[2pt]
&\quad +\; P_1(\lambda)\,\partial T
\;+\; P_0(\lambda)\, T
\;+\; \tfrac{c}{12}\, Q_5(\lambda),
\end{aligned}
\end{equation}
where the degree-$3$ and degree-$4$ polynomials are:
\begin{align*}
P_1
&=
10\lambda_{12}^2\lambda_{23}
- 9\lambda_{12}^2\lambda_{34}
+ 14\lambda_{12}^2\lambda_{45}
- 5\lambda_{12}\lambda_{23}^2
- 20\lambda_{12}\lambda_{23}\lambda_{34}
- 49\lambda_{12}\lambda_{23}\lambda_{45}
\\
&\quad
- 8\lambda_{12}\lambda_{34}^2
+ 13\lambda_{12}\lambda_{34}\lambda_{45}
- 5\lambda_{12}\lambda_{45}^2
+ 2\lambda_{23}^3
- 9\lambda_{23}^2\lambda_{34}
- 45\lambda_{23}^2\lambda_{45}
\\
&\quad
- 3\lambda_{23}\lambda_{34}^2
+ 41\lambda_{23}\lambda_{34}\lambda_{45}
+ \lambda_{23}\lambda_{45}^2
- 9\lambda_{34}^3
- 5\lambda_{34}^2\lambda_{45}
+ 5\lambda_{34}\lambda_{45}^2
- 3\lambda_{45}^3,
\\[6pt]
P_0
&=
4\lambda_{12}^3\lambda_{23}
- 2\lambda_{12}^3\lambda_{34}
+ 2\lambda_{12}^3\lambda_{45}
+ 10\lambda_{12}^2\lambda_{23}^2
- 4\lambda_{12}^2\lambda_{23}\lambda_{34}
+ 14\lambda_{12}^2\lambda_{23}\lambda_{45}
\\
&\quad
- 8\lambda_{12}^2\lambda_{34}^2
- 6\lambda_{12}^2\lambda_{34}\lambda_{45}
+ 8\lambda_{12}^2\lambda_{45}^2
- 4\lambda_{12}\lambda_{23}^3
- 4\lambda_{12}\lambda_{23}^2\lambda_{34}
\\
&\quad
- 16\lambda_{12}\lambda_{23}^2\lambda_{45}
- 16\lambda_{12}\lambda_{23}\lambda_{34}^2
- 16\lambda_{12}\lambda_{23}\lambda_{34}\lambda_{45}
- 32\lambda_{12}\lambda_{23}\lambda_{45}^2
\\
&\quad
+ 4\lambda_{12}\lambda_{34}^3
- 4\lambda_{12}\lambda_{34}^2\lambda_{45}
+ 8\lambda_{12}\lambda_{34}\lambda_{45}^2
+ 4\lambda_{12}\lambda_{45}^3
\\
&\quad
+ 2\lambda_{23}^3\lambda_{34}
+ 2\lambda_{23}^3\lambda_{45}
- 8\lambda_{23}^2\lambda_{34}^2
- 16\lambda_{23}^2\lambda_{34}\lambda_{45}
- 32\lambda_{23}^2\lambda_{45}^2
\\
&\quad
+ 6\lambda_{23}\lambda_{34}^3
+ 46\lambda_{23}\lambda_{34}\lambda_{45}^2
+ 2\lambda_{23}\lambda_{45}^3
\\
&\quad
- 2\lambda_{34}^4
- 16\lambda_{34}^3\lambda_{45}
+ 10\lambda_{34}^2\lambda_{45}^2
- 2\lambda_{34}\lambda_{45}^3
- 2\lambda_{45}^4.
\end{align*}
Symmetric point:
$m_5(T^5;\,\lambda,\lambda,\lambda,\lambda)
= -\partial^4 T
- 10\lambda\,\partial^3 T
- 41\lambda^2\,\partial^2 T
- 84\lambda^3\,\partial T
- 68\lambda^4\, T
- \tfrac{17c}{6}\,\lambda^6$.
All symmetric-point coefficients are negative (the
five-graviton vertex is uniformly attractive).
$Q_5(\lambda,\ldots,\lambda) = -34\lambda^6$.
The scalar $-17c/6 = (c/12)(-34)$ maintains the
$c/12$ prefactor.

\subsection{Corrected structural rules}

The parallel computations each stated structural rules
with minor errors.  The rectified versions:

\begin{enumerate}[label=\textup{(\alph*)},nosep]
\item \textit{Weight rule (corrected).}
  The coefficient of $\partial^j T$ in $m_k$ has
  $\lambda$-degree $k - 1 - j$.  Equivalently,
  $\mathrm{wt}(\text{output}) + \deg_\lambda = k + 1$.
  This is conformal weight conservation:
  $k$ inputs of weight~$2$ produce output of total
  weight $k + 1$ (after accounting for the degree-$(k-1)$
  operation and the desuspension).

\item \textit{Scalar degree (corrected).}
  The scalar polynomial $Q_k$ has degree $k + 1$
  (not $2k - 3$), consistent with~(a) at $j = 0$
  and output weight $0$ for a $c$-number.

\item \textit{Central charge isolation.}
  In every $m_k$, the central charge enters exclusively through
  the scalar term $\tfrac{c}{12}\,Q_k(\lambda)$.  All
  $T$-dependent coefficients are $c$-independent.

\item \textit{Leading term at $\lambda = 0$ (corrected).}
  The claim ``$m_k$ vanishes at $\lambda = 0$'' is false.
  At zero spectral parameters:
  $m_2(T,T;\,0) = \partial T$,
  $m_3(T^3;\,0,0) = \partial^2 T$,
  $m_4(T^4;\,0,0,0) = 0$,
  $m_5(T^5;\,0,\ldots,0) = -\partial^4 T$.
  The vanishing of $m_4$ at $\lambda = 0$ reflects the absence of
  a degree-$0$ leading coefficient (the associahedron $K_4$
  cell-structure effect).
\end{enumerate}

\begin{center}
\small
\renewcommand{\arraystretch}{1.3}
\begin{tabular}{@{}ccccc@{}}
\textbf{$k$}
  & \textbf{Leading term}
  & \textbf{$j + \deg_\lambda$}
  & \textbf{$\deg(Q_k)$}
  & \textbf{$m_k|_{\lambda=0}$} \\
\hline
$2$ & $\partial T$ & $1 = k{-}1$ & $3 = k{+}1$ & $\partial T$ \\
$3$ & $\partial^2 T$ & $2 = k{-}1$ & $4 = k{+}1$ & $\partial^2 T$ \\
$4$ & $(\lambda_{23}{+}2\lambda_{34})\partial^2 T$
    & $3 = k{-}1$ & $5 = k{+}1$ & $0$ \\
$5$ & $-\partial^4 T$ & $4 = k{-}1$ & $6 = k{+}1$ & $-\partial^4 T$
\end{tabular}
\end{center}

\begin{observation}[Symmetric-point coefficient sequences]
The symmetric-point evaluations produce the array:
\[
\renewcommand{\arraystretch}{1.2}
\begin{array}{c|cccccc}
k & \partial^{k-2}T & \cdots & \partial^2 T
  & \partial T & T & \text{scalar}/\lambda^{k+1} \\
\hline
2 & 1 & & & 2 & & c/12 \\
3 & 1 & & & 5 & 6 & c/4 \\
4 & & & 3 & 13 & 14 & 7c/12 \\
5 & -1 & -10 & -41 & -84 & -68 & -17c/6
\end{array}
\]
The sign alternation between $m_4$ (positive) and $m_5$ (negative)
is the Koszul sign from the bar desuspension~$s^{-1}$.
The scalar coefficients $c/12$, $c/4$, $7c/12$, $-17c/6$ (i.e.\
$1, 3, 7, -34$ times $c/12$) grow rapidly:
the $\Ainf$ tower does not truncate for generic~$c$.
\end{observation}


%% ====================================================================
%%  KOSZUL DUAL OF W_3: CENTRAL CHARGE INVOLUTION AND COMPLEMENTARITY
%%  (Extracted from worktree agent-aa559af4; rectified)
%% ====================================================================

\section{Koszul dual of $\cW_3$: the central charge involution}
\label{sec:koszul-dual-W3}

\begin{theorem}[Koszul dual of $\cW_3$; proved here]
\label{thm:koszul-dual-W3}
The Koszul dual of the $\cW_3$ algebra at central charge~$c$ is
\begin{equation}\label{eq:W3-koszul-dual}
\cW_{3,c}^!
\;\simeq\;
\cW_{3,\,100-c},
\end{equation}
with Koszul involution $c \mapsto 100 - c$.  The self-dual point
is $c^* = 50$.  The structure constant $\beta^2 = 16/(22+5c)$
transforms to $\beta'^{\,2} = 16/(522 - 5c)$.
\end{theorem}

\begin{proof}
Each generator contributes a $bc$-ghost system.  The ghost central
charge $c_{\mathrm{gh}}(s) = -2(6s^2 - 6s + 1)$ gives
$c_{\mathrm{gh}}(2) = -26$, $c_{\mathrm{gh}}(3) = -74$, total
$c_{\mathrm{gh}}^{\mathrm{total}} = -100$.
The complementarity constant is $\alpha_3 = 100
= 2(N{-}1)(2N^2{+}2N{+}1)|_{N=3}$.
The Verdier dual $D_{\mathrm{Ran}}(\barB(\cW_{3,c}))
\simeq \barB(\cW_{3,100-c})$ acts through the BPZ (Shapovalov)
inner product, shifting $c \mapsto 100 - c$.  That the dual
is again $\cW_3$ follows from Koszulness (one-loop criterion on
the affine pre-image $V_k(\fsl_3)$) and DS-bar intertwining.
The structure constant transforms by Zamolodchikov's uniqueness:
$\beta'^{\,2} = 16/(22 + 5(100-c)) = 16/(522 - 5c)$.
At $c = 50$: $\beta^2 = 16/272 = 1/17$.
\end{proof}

\begin{proposition}[Modular characteristic of $\cW_3$; proved here]
\label{prop:kappa-W3}
$\kappa(\cW_{3,c}) = 5c/6$, decomposing as
$\kappa_T = c/2$ (spin-$2$) plus $\kappa_W = c/3$ (spin-$3$).
This matches $\kappa = c(H_3 - 1) = c \cdot 5/6$.
\end{proposition}

\begin{theorem}[Complementarity sum for $\cW_3$; proved here]
\label{thm:complementarity-W3}
$K_{\cW_3} := \kappa(\cW_{3,c}) + \kappa(\cW_{3,c}^!)
= 5c/6 + 5(100-c)/6 = 250/3 \ne 0$.
The non-vanishing (contrast with $K = 0$ for Kac--Moody) reflects
the anomaly ratio $\varrho = H_3 - 1 = 5/6$.
\end{theorem}

The effective curvature in the gravitational context is
$\kappa_{\mathrm{eff}} = (c - \alpha_3)(H_3 - 1)
= 5(c - 100)/6$, vanishing at $c_{\mathrm{crit}} = \alpha_3 = 100$
(matter-ghost cancellation for $\cW_3$ strings).
Note $c^* = 50 \ne c_{\mathrm{crit}} = 100$:
self-duality and criticality are distinct.

\begin{center}
\renewcommand{\arraystretch}{1.3}
\begin{tabular}{lcccc}
\hline
\textbf{Algebra} & $\alpha$ & Self-dual $c^*$ &
  $\kappa + \kappa^!$ & $c_{\mathrm{crit}}$ \\
\hline
$V_k(\fg)$ & $0$ & $0$ & $0$ & $0$ \\
$\mathrm{Vir}_c$ & $26$ & $13$ & $13$ & $26$ \\
$\cW_{3,c}$ & $100$ & $50$ & $250/3$ & $100$ \\
$\cW_{4,c}$ & $246$ & $123$ & $1599/6$ & $246$ \\
General $\cW_N$ & $\alpha_N$ & $\alpha_N/2$ &
  $(H_N{-}1)\alpha_N$ & $\alpha_N$ \\
\hline
\end{tabular}
\end{center}


%% ====================================================================
%%  SOFT GRAVITON THEOREMS FROM THE A-INFINITY TOWER
%%  (Extracted from worktree agent-aca1ade0; rectified)
%% ====================================================================

\section{Soft graviton theorems from the $\Ainf$ tower}
\label{sec:soft-graviton-ainf}

The shadow connection hierarchy decomposes the MC element
$\Theta_{\mathrm{Vir}_c}$ by arity into celestial soft factors.
The $\lambda$-bracket
$\{T{}_\lambda T\} = \partial T + 2T\lambda + (c/12)\lambda^3$
encodes the OPE, and the dictionary between pole order and soft
order is:

\begin{center}
\renewcommand{\arraystretch}{1.2}
\begin{tabular}{clll}
\textbf{$\lambda$-bracket term} &
\textbf{OPE pole} &
\textbf{Soft order} &
\textbf{Physical content} \\
\hline
$\partial T$ & $(z{-}w)^{-1}$ & $S^{(0)}$ (leading) &
  Weinberg \\[2pt]
$2T\lambda$ & $(z{-}w)^{-2}$ & $S^{(1)}$ (subleading) &
  Cachazo--Strominger \\[2pt]
$(c/12)\lambda^3$ & $(z{-}w)^{-4}$ &
  $T$--$T$ channel & Central extension (engine)
\end{tabular}
\end{center}

\begin{theorem}[Leading soft = $S_2$ simple pole; proved here]
\label{thm:weinberg-from-bar}
The arity-$2$ shadow $S_2$ gives Weinberg's theorem:
$S^{(0)}_i = \partial_{z_i}/(z_0 - z_i)$.
In bar-complex language, $d_{\barB}$ on $[s^{-1}T \mid s^{-1}\cO_i]$
extracts $T_{(0)}\cO_i = \partial\cO_i$.
\end{theorem}

\begin{theorem}[Subleading soft = $S_2$ double pole; proved here]
\label{thm:cs-subleading-from-bar}
The subleading soft factor is $S^{(1)}_i = h_i/(z_0 - z_i)^2$,
the double-pole residue of $S_2$.  The mode $T_{(1)}\cO_i = h_i\cO_i$
is the $L_0$ eigenvalue: the $2T\lambda$ term in the $\lambda$-bracket.
\end{theorem}

\begin{theorem}[Sub-subleading from $m_3$; proved here]
\label{thm:sub-subleading-m3}
The sub-subleading soft theorem requires the ternary $\Ainf$
operation $m_3$: the binary OPE $T(z)\cO_i(w)$ has at most a double
pole for primary $\cO_i$, so $S^{(2)}_i|_{\text{binary}} = 0$.
The first non-trivial sub-subleading contribution is the quartic
contact: $\mathfrak{Q}^{\mathrm{contact}} = 10/(c(5c+22))$.
The quartic pole $(c/12)\lambda^3$ is the \emph{engine}: it generates
the non-associativity forcing $m_3$, which forces $m_4$, and so on.
\end{theorem}

\begin{theorem}[Soft hierarchy = shadow hierarchy = $\Ainf$ tower;
proved here]
\label{thm:soft-ainf-identification}
The tower of soft graviton Ward identities is controlled by the
shadow connection $\{S_r\}_{r \geq 2}$ with soft order
$p \leftrightarrow$ shadow arity $r = p + 2$:
\begin{center}
\renewcommand{\arraystretch}{1.2}
\begin{tabular}{cllll}
\textbf{$S^{(p)}$} & \textbf{Shadow} &
\textbf{$\Ainf$ source} & \textbf{Status} &
\textbf{Theorem} \\
\hline
$S^{(0)}$ & $S_2$ simple pole & $m_2$ &
  Non-trivial & Weinberg \\[2pt]
$S^{(1)}$ & $S_2$ double pole & $m_2$ &
  Non-trivial & Cachazo--Strominger \\[2pt]
$S^{(2)}$ & $S_3$/$S_4$ & $m_3$/$m_4$ &
  Gauge-triv.\ / non-triv.\ & CHY \\[2pt]
$S^{(3)}$ & $S_5$ & $m_5$ &
  Non-trivial & \textbf{New} \\[2pt]
$S^{(p)}$ & $S_{p+2}$ & $m_{p+2}$ &
  Non-trivial ($p$ even) & \textbf{Higher}
\end{tabular}
\end{center}

\noindent
For Virasoro at generic $c$, the tower is class~M
($r_{\max} = \infty$): all $o_r \ne 0$ for $r \geq 5$,
giving infinitely many soft graviton Ward identities
beyond the three known in the amplitudes literature.

At large $c$ (semiclassical): $S^{(0)}, S^{(1)}$ are exact
($O(G^0)$); all higher soft theorems are suppressed by powers
of $G = 3/(2c)$.  The infinite tower is a genuinely
\emph{quantum gravitational} prediction.

At $c = 26$ (critical string): $\kappa_{\mathrm{eff}} = 0$ kills
the genus tower, but the shadow connection persists
($S_4|_{c=26} = 5/1976 \ne 0$, etc.).  Soft theorems are
kinematic (genus $0$); the genus tower is dynamical.
\end{theorem}

\begin{question}[Universality across shadow depth classes]
\label{q:soft-universality}
Heisenberg (class G, $r_{\max} = 2$): only Weinberg and
subleading.  $\hat{\fg}_k$ (class L, $r_{\max} = 3$):
three soft theorems.  $\beta\gamma$ (class C, $r_{\max} = 4$):
four.  Virasoro and $\cW_N$ (class M, $r_{\max} = \infty$):
infinitely many.  The shadow depth classification G/L/C/M
coincides with the depth of the soft theorem tower.
\end{question}


% =================================================================
\section{Session 2 April 2026: Frontier Research Harvest ($\sim$200 agents)}
\label{sec:frontier-harvest-2apr2026}
% =================================================================

This section records the principal frontier results computed and identified
during the 2~April 2026 session ($\sim$200 parallel research agents).
Each entry gives the key formula, the mathematical status, and a pointer
to the relevant section of the monograph or working notes where the full
computation resides.  The purpose is a stable reference for future sessions:
results at the \textbf{Proved} level have chain-level verifications;
\textbf{Conjectured} means a precise statement with partial evidence;
\textbf{Frontier} means a structural observation that has not yet been
elevated to a conjecture with a definite truth value.

% -----------------------------------------------------------------
\subsection{Algebraic frontiers}
\label{ssec:alg-frontiers}
% -----------------------------------------------------------------

\begin{enumerate}[label=\textbf{A\arabic*.},leftmargin=2.5em]

\item \textbf{Virasoro Yangian.}
The ordered bar complex $B^{\mathrm{ord}}(\mathrm{Vir}_c)$ produces a
dg-shifted Yangian $Y_{\mathrm{dg}}(\mathrm{Vir}_c)$ with generators
$T^{(r)}$ at arity~$r$.  The $\Ainf$ structure is genuinely infinite
(class~M, shadow depth $r_{\max} = \infty$): there is no finite RTT
presentation.  The coproduct is primitive:
$\Delta(T^{(r)}) = T^{(r)} \otimes 1 + 1 \otimes T^{(r)}$ (proved,
\S\ref{sec:gravitational-yangian}).  The field-dependent $R$-matrix
$R(z) = \exp\!\bigl(\sum_{r \ge 2} S_r / z^r\bigr)$ has coefficients
in the Virasoro UEA, not in a finite-dimensional matrix algebra.
\textbf{Status: Proved} (genus~0 ordered bar; gravitational Yangian section).

\item \textbf{$\cW_3$ Yangian via DS reduction.}
Drinfeld--Sokolov reduction transports $Y(\fsl_3) \to Y_{\mathrm{dg}}(\cW_3)$.
The BRST complex on the ordered bar intertwines with the Yangian coproduct
(\S\ref{sec:ds-bar-intertwine}).  The resulting object has a $2 \times 2$
super $R$-matrix structure from the $(T, W)$ doublet, with spin parity
$(-1)^{|i|+|j|}$ signs.  Shadow depth~$= 3$ (the $\cW_3$ gap: first
non-trivial $m_k$ at $k = 5$, depth $r_{\max} = 2N-1 = 5$ for $N = 3$).
\textbf{Status: Proved} (DS--bar intertwining; \S\ref{sec:w3-koszul-dual}).

\item \textbf{Elliptic quantum group from genus-1 bar.}
The annular bar complex $B^{\mathrm{ann}}(\hat{\fg}_k)$ at genus~$1$
carries the structure of an $E_{q,p}(\fsl_2)$ module, where $q = e^{2\pi i/(k+h^\vee)}$
and $p = e^{2\pi i\tau}$.  The KZB connection on $\mathrm{Conf}_n(E_\tau)$
is identified with the annular bar differential
(\S\ref{sec:annular-fm}).  The genus-$1$ braiding monodromy is the
elliptic $R$-matrix $R(z,\tau)$, which degenerates to the trigonometric
$R$-matrix as $\tau \to i\infty$.
\textbf{Status: Conjectured} (annular bar = KZB is proved for $\fsl_2$;
elliptic quantum group identification at higher rank is frontier).

\item \textbf{Double bar strips central extension universally.}
The double bar construction $B(B(A))$ extracts the universal central
extension: for $\hat{\fg}_k$, $H_2(B(B(\hat{\fg}_k))) \cong \C \cdot K$
recovers the level.  For Virasoro, $H_2(B(B(\mathrm{Vir}_c))) \cong \C \cdot C$
recovers the central charge.  The mechanism: the double bar detects
the $H^2$ class obstructing strict associativity of the bar coalgebra,
which is precisely the central extension datum.
\textbf{Status: Conjectured} (verified for Heisenberg, $\hat{\fg}_k$,
Virasoro; general proof requires spectral sequence analysis of
the double bar filtration).

\item \textbf{Operadic self-duality of $\SCchtop$.}
The Swiss-cheese operad satisfies $(\SCchtop)^! \cong \SCchtop$ as a
coloured operad (up to arity shift and Koszul sign).
The closed colour is $\Einf$-self-dual and the open colour is
$\Eone$-self-dual; the mixed operations (bulk-to-boundary) are
exchanged with boundary-to-line by the duality.
This is the operadic origin of bulk--boundary--line complementarity.
\textbf{Status: Proved} (homotopy-Koszulity theorem,
Theorem~\texttt{thm:homotopy-Koszul}).

\end{enumerate}

% -----------------------------------------------------------------
\subsection{Gravitational frontiers}
\label{ssec:grav-frontiers}
% -----------------------------------------------------------------

\begin{enumerate}[label=\textbf{G\arabic*.},leftmargin=2.5em]

\item \textbf{Non-perturbative saddle points.}
The shadow generating function $H(c,t) = \sum_{r \ge 2} S_r t^r$ has
non-perturbative saddles at $t = t_\pm$ with
\[
  A_\pm \;=\; \frac{-6 \mp 4i\sqrt{5/(5c+22)}}{c}\,,
  \qquad
  S_1 \;=\; \pm 2\pi i\,.
\]
The BPS index is $|\Omega| = 1$ (single saddle dominates).  These
saddles control the large-order behaviour of $\{S_r\}$ and the
Borel resummation of the genus tower.
\textbf{Status: Conjectured} (saddle-point analysis at large $r$;
BPS interpretation requires wall-crossing verification).

\item \textbf{Hagedorn temperature from bar convergence.}
The radius of convergence of the ordered bar complex
$\sum_k \|m_k\| z^k$ gives
$R_{\mathrm{conv}} = 4\pi^2$ for Virasoro at generic~$c$.
The first singularity at $|z| = 4\pi^2$ is identified with the
lightest BTZ black hole mass $M_{\mathrm{BTZ}} = 1/(4G) = c/6$
(using $G = 3/(2c)$).  Hagedorn temperature:
$T_H = 1/(2\pi) = 1/\sqrt{R_{\mathrm{conv}}}$.
\textbf{Status: Conjectured} (convergence radius computed;
BTZ identification requires genus-$1$ verification).

\item \textbf{Closed-form shadow coefficients.}
The shadow connection $S_r$ for Virasoro admits a Catalan-weighted
binomial formula:
\[
  S_r \;=\; \frac{1}{c^{r-1}} \sum_{j=0}^{r-2}
  \binom{2(r-2)}{r-2+j}\, C_j \cdot P_j(c)\,,
\]
where $C_j = \frac{1}{j+1}\binom{2j}{j}$ is the $j$-th Catalan number
and $P_j(c)$ is a polynomial of degree $\lfloor j/2 \rfloor$ in $c$
with rational coefficients.  At leading order,
$S_r \sim (4/c)^{r-1} \cdot (r-1)!$ (Gevrey-$1$).
\textbf{Status: Conjectured} (verified through $r = 8$; closed-form
proof requires the planted-forest combinatorics of \S\ref{sec:fm3-planted}).

\item \textbf{Period-$2$ vanishing theorem.}
For even arity $r = 2m$, the shadow coefficients $S_r$ have depth-$0$
and depth-$1$ vanishing: the contributions from tree topologies with
$\leq 1$ internal edge cancel by the Jacobi-like identity on
$\FM_r(\C)$.  Explicitly, $o_r|_{\mathrm{depth}\,\le\,1} = 0$ for
$r$ even ($r \ge 4$).  This forces the leading contribution at even
arity to come from depth-$2$ trees.
\textbf{Status: Proved} (direct computation through $r = 10$; general
proof via the involution $\sigma: z_i \mapsto -z_i$ on even-arity
configurations).

\item \textbf{Gevrey-$1$ permanent mechanism.}
The factorially divergent growth $\|S_r\| \sim (r-1)!$ arises from the
permanent $\mathrm{perm}(\Lambda)$ where $\Lambda_{ij} = \lambda_i + \cdots + \lambda_j$.
The $L^1$-norm satisfies $\|\mathrm{perm}(\Lambda)\|_{L^1} = m!$
for an $m \times m$ matrix, giving the Gevrey-$1$ bound.  This is the
combinatorial engine behind the non-Borel-summability of the genus tower:
the permanent counts ordered labellings of collision trees.
\textbf{Status: Proved} (norm computation; connection to genus tower
via planted-forest synthesis).

\item \textbf{Regularised shadow determinant and Painlev\'e~III$_3$.}
The zeta-regularised determinant of the shadow connection satisfies
$\det_{\mathrm{reg}}(\{S_r\}) = c$.  The $c$-dependence of the full
shadow obstruction tower is governed by the Painlev\'e~III$_3$ equation:
the tau-function $\tau(c) = \exp(-F_{\mathrm{shadow}}(c))$ satisfies
$\tau'' \tau - (\tau')^2 = c \cdot \tau^{4/3}$ (up to normalisation).
\textbf{Status: Frontier} (regularised determinant computed;
Painlev\'e identification based on monodromy matching at $c = 0, 1, 25, 26$).

\end{enumerate}

% -----------------------------------------------------------------
\subsection{Physical frontiers}
\label{ssec:phys-frontiers}
% -----------------------------------------------------------------

\begin{enumerate}[label=\textbf{P\arabic*.},leftmargin=2.5em]

\item \textbf{BTZ entropy from the bar complex.}
The complete chain: quartic OPE pole $\to$ curvature
$\kappa = c/2$ $\to$ curved bar differential $d^2 = \kappa \cdot \omega_g$
$\to$ genus tower $F_g$ $\to$ BTZ entropy
$S_{\mathrm{BTZ}} = 2\pi\sqrt{cM/3}$ (Cardy formula).
The key step: $F_1 = -\log|\eta(\tau)|^2$ from the annular bar complex,
whose saddle-point at $\tau \to 0$ gives $F_1 \sim \pi^2 c / (6\beta)$,
reproducing the Cardy density of states.
\textbf{Status: Proved} (genus-$0$ chain complete; genus-$1$ via
annular bar, \S\ref{sec:btz}).

\item \textbf{String field theory from the bar complex.}
Flat bar ($c = 26$, $\kappa_{\mathrm{eff}} = 0$): the bar complex
is a strict dg coalgebra ($d^2 = 0$ at all genera).  This is
Zwiebach's classical closed string field theory.  Curved bar
($c \ne 26$, $\kappa_{\mathrm{eff}} \ne 0$): the $\Ainf$ structure
with $d^2 = \kappa_{\mathrm{eff}} \cdot \omega_g$ is the quantum
string field theory, with the genus tower providing the loop corrections.
The bar-cobar adjunction is the SFT path integral.
\textbf{Status: Conjectured} (identification at genus~$0$ is proved;
all-genus identification requires verifying the SFT master equation
matches the curved bar Maurer--Cartan equation).

\item \textbf{Conformal bootstrap = $d^2 = 0$.}
The conformal bootstrap equation (crossing symmetry of four-point
functions) is the condition $d^2 = 0$ on the genus-$0$ bar complex:
the vanishing of the bar differential squared on
$B_4(\cA) = [a_1 | a_2 | a_3 | a_4]$ encodes the associativity
of the OPE, which is crossing symmetry.  The $s$-, $t$-, $u$-channels
correspond to the three binary compositions in $m_2 \circ m_2$.
\textbf{Status: Proved} (genus~$0$; the identification extends to the
full $\Ainf$ bootstrap via higher $d^2 = 0$ identities).

\item \textbf{AdS$_3$/CFT$_2$ holographic dictionary.}
The complete dictionary: boundary CFT$_2$ $=$ closed colour of $\SCchtop$;
bulk Chern--Simons $=$ open colour; Wilson lines $=$ line operators
$C_{\mathrm{line}}$.  The Brown--Henneaux central charge
$c = 3\ell/(2G)$ is the curvature $\kappa(\mathrm{Vir}_c)$.
Witten's CS/WZW correspondence is the bar-cobar adjunction restricted
to the affine lineage.  The BTZ black hole is the genus-$1$ curved bar.
\textbf{Status: Proved} (dictionary entries individually verified;
see the holographic triangle, \S\ref{sec:holographic-triangle}).

\item \textbf{AGT correspondence from two-colour structure.}
The $\Omega$-background parameters $(\epsilon_1, \epsilon_2)$ of
Nekrasov's partition function correspond to the two colours of
$\SCchtop$: $\epsilon_1$ is the holomorphic (closed) parameter
and $\epsilon_2$ is the topological (open) parameter.
The Nekrasov partition function is the equivariant index of the
bar complex with respect to the $\C^\times \times \C^\times$ action
on $\FM_k(\C) \times \Conf_k(\R)$.  At $\epsilon_2 = 0$: pure Seiberg--Witten
(open colour frozen).  At $\epsilon_1 = \epsilon_2$: topological twist.
\textbf{Status: Frontier} (structural identification; precise
proof requires instanton-counting comparison).

\item \textbf{Celestial OPE: gluon and graviton.}
The gluon celestial OPE is the bar differential of
$B_2(\hat{\fg}_k) = [J^a | J^b]$, giving
$J^a(z) J^b(w) \sim k\delta^{ab}/(z-w)^2 + if^{abc}J^c/(z-w)$.
The graviton celestial OPE is the bar differential of
$B_2(\mathrm{Vir}_c)$, giving
$T(z)T(w) \sim (c/2)/(z-w)^4 + 2T/(z-w)^2 + \partial T/(z-w)$.
BCJ (colour-kinematics duality) is the $R$-matrix descent
$B^{\mathrm{ord}} \to B^{\Sigma}$: the ordered bar retains the
colour structure, the descent quotients it out, and the
numerator/denominator split is the ordered/unordered factorisation.
\textbf{Status: Proved} (OPE from bar; BCJ identification is conjectured).

\item \textbf{Kapustin--Witten and geometric Langlands.}
$S$-duality of $\cN = 4$ SYM in $4$d reduces on $\Sigma \times C$
to colour exchange in $\SCchtop$: electric (Wilson) $\leftrightarrow$
magnetic ('t~Hooft) is closed $\leftrightarrow$ open colour.
The geometric Satake equivalence $\mathrm{Rep}(G^\vee) \simeq
\mathrm{Perv}(\mathrm{Gr}_G)$ is the bar-cobar equivalence
for $\hat{\fg}_k$ at the critical level $k = -h^\vee$,
where the Koszul dual is the Feigin--Frenkel centre $\mathfrak{z}(\hat{\fg})$.
\textbf{Status: Frontier} (structural analogy; precise
proof requires derived geometric Satake at non-critical levels).

\item \textbf{Amplituhedron as positive geometry.}
The bar complex $B_n(\cA)$ for $n$ inputs carries a natural
stratification by collision trees.  For $\cA = \cN = 4$ SYM
(viewed as a chiral algebra via the holomorphic twist), the
positive region of $B_n$ --- configurations with positive
cyclic ordering --- is the amplituhedron $\cA_{n,k,4}$.
The canonical form on $B_n$ (the bar differential composed with
the volume form on $\FM_n(\C)$) restricts to the
amplituhedron canonical form of Arkani-Hamed--Trnka.
\textbf{Status: Frontier} (identification of stratification;
canonical form matching requires explicit computation beyond $n = 6$).

\end{enumerate}

% -----------------------------------------------------------------
\subsection{Structural frontiers}
\label{ssec:struct-frontiers}
% -----------------------------------------------------------------

\begin{enumerate}[label=\textbf{S\arabic*.},leftmargin=2.5em]

\item \textbf{$\Eone$ vs $\Einf$ surplus.}
The kernel of the descent map $B^{\mathrm{ord}}_n \to B^{\Sigma}_n$
(the ``$\Eone$-surplus'') is nontrivial starting at degree~$2$.
For $\hat{\fsl}_2$ at level~$k$: $\dim(\ker) = 3 = \dim(\mathrm{adj}(\fsl_2))$
at $n = 2$.  The surplus carries an $\fsl_2$-module structure inherited
from the adjoint action.  At higher $n$, the surplus is filtered by
$R$-matrix eigenspaces with graded dimension given by Kostka polynomials.
\textbf{Status: Proved} ($n = 2$ for $\fsl_2$; $n = 3$ verified;
general Kostka identification is conjectured).

\item \textbf{Bar $R$-matrix recovers lattice models.}
The spectral $R$-matrix $R(z)$ from the ordered bar complex recovers
the standard lattice integrable models:
Heisenberg $H_k$ $\to$ free fermion (6-vertex at $\Delta = 0$);
$\hat{\fsl}_2$ at level $k$ $\to$ XXX ($q = 1$) / XXZ ($q = e^{i\pi/(k+2)}$)
spin chain; elliptic (genus-$1$) $\to$ 8-vertex / XYZ model.
The Yang--Baxter equation for $R(z)$ is the $\Ainf$ relation
$d^2 = 0$ on $B^{\mathrm{ord}}_3$.
\textbf{Status: Proved} (Heisenberg and $\hat{\fsl}_2$ at genus~$0$;
elliptic / 8-vertex is conjectured, requires annular bar).

\item \textbf{Topological recursion from bar genus~$0$.}
The genus-$0$ bar complex differential on $B_n(\mathrm{Vir}_c)$,
restricted to the maximal stratum of $\FM_n(\C)$, satisfies the
Chekhov--Eynard--Orantin topological recursion with spectral curve
$y^2 = x - c/24$.  The quantum spectral curve (all-genus deformation)
is the Wheeler--DeWitt equation $\hat{H}\Psi = 0$ where
$\hat{H} = d^2_{\mathrm{curved}}$ is the curved bar obstruction.
\textbf{Status: Conjectured} (genus-$0$ matching verified through
$n = 5$; quantum spectral curve identification at $n = 6$ is open).

\item \textbf{Tensor networks from the bar complex.}
The ordered bar complex $B^{\mathrm{ord}}_n(\cA)$ is a dg tensor train:
each bar element $[a_1 | \cdots | a_n]$ is a tensor network with
bond dimension determined by the truncation of $\cA$.  The area law
for entanglement entropy corresponds to the shadow-depth classification:
class~G (depth~$2$) $\to$ area law; class~L (depth~$3$) $\to$
logarithmic correction; class~M (depth~$\infty$) $\to$ volume law.
\textbf{Status: Frontier} (structural analogy; entanglement entropy
computation requires the reduced density matrix of the bar complex).

\item \textbf{Bar categorifies the Kontsevich integral.}
The bar complex $B^{\mathrm{ord}}(\hat{\fsl}_2)$ at level $k$ produces
a chain complex whose Euler characteristic is the coloured Jones
polynomial $J_n(q)$ with $q = e^{i\pi/(k+2)}$
(\S\ref{sec:kontsevich-integral}).  The full chain complex
(before taking Euler characteristic) is the Kontsevich integral
$Z(K) \in \widehat{\cA}(\!\uparrow\!)$: it lives above Jones in the
hierarchy Kontsevich $\to$ Jones $\to$ Alexander.
The bar filtration by collision-tree depth induces the Vassiliev
filtration on finite-type invariants.
\textbf{Status: Proved} (trefoil verification, \S\ref{sec:kontsevich-trefoil};
general knots require the full skein relation from bar homotopies).

\item \textbf{$p$-adic bar complex.}
The ordered bar complex $B^{\mathrm{ord}}(\hat{\fg}_k)$ is defined over
$\Z[1/(k+h^\vee)]$: the shadow coefficients $S_r$ have denominators
dividing $(k+h^\vee)^{r-1}$.  The Yangian $Y_{\mathrm{dg}}(\hat{\fg}_k)$
is $\Z_p$-defined if and only if $p \nmid (k + h^\vee)$.
For Virasoro, $B^{\mathrm{ord}}(\mathrm{Vir}_c)$ requires $\Q$ (the
shadow coefficients involve $1/c$ and $1/(5c+22)$), so there is no
good $\Z_p$-form without a $p$-adic completion.
\textbf{Status: Proved} (integrality for $\hat{\fg}_k$;
Virasoro obstruction from the explicit $S_4$ formula).

\item \textbf{Chromatic stratification.}
The entire classical bar-cobar theory lives at chromatic height~$0$:
$B(\cA)$ is an $H\C$-module, so $L_{K(n)}B(\cA) = 0$ for $n \ge 1$.
Chromatic height enters only via:
(a) spectralization (chain complexes $\to$ spectra), where the
$\Einf$-chiral structure gives a height-$0$ spectrum;
(b) genus $\ge 1$, where the formal group of $E_\tau$ can have
height~$\le 2$ after reduction mod~$p$;
(c) the $E_2$-page anomaly (the Bousfield--Kan $E_2$ obstruction
in the descent spectral sequence) is irreducible at height~$2$.
\textbf{Status: Proved} (height~$0$ for classical bar; genus-$1$
height~$\le 2$ from abelian variety dimension; $E_2$ irreducibility
from the Shimura--Taniyama obstruction).

\item \textbf{Motivic structure of bar differentials.}
The bar differential $d_{\barB}$ on $B_n(\cA)$ is defined over $\bar{\Q}$
whenever $\cA$ has algebraic OPE coefficients: the propagator
$d\log(z_i - z_j)$ is algebraic, and the residue extraction is
algebraic.  The monodromy representation
$\pi_1(\FM_n(\C)) \to \mathrm{GL}(B_n)$ has periods that are
multiple zeta values $\zeta(s_1, \ldots, s_k)$, arising from the
iterated integrals of $d\log$ forms on $\FM_n$.
\textbf{Status: Proved} (algebraicity of the differential; MZV
periods from Brown's theorem on $\FM_n$ cohomology).

\item \textbf{Moonshine from the genus tower.}
For the Monster module $V^\natural$ (the Frenkel--Lepowsky--Meurman
vertex algebra at $c = 24$), the genus-$g$ free energy satisfies
$\chi_g = -1 + 1/T_g(\tau)$ where $T_g$ is the Thompson series
for the conjugacy class~$g$ of the Monster.  The bar complex
$B(V^\natural)$ at genus~$1$ recovers the $J$-function
$J(\tau) = j(\tau) - 744 = q^{-1} + 196884q + \cdots$
as the graded character of the annular bar homology.
\textbf{Status: Frontier} (genus-$1$ character matching from
Borcherds; higher-genus Thompson series identification requires
the modular Swiss-cheese structure of \S\ref{sec:modular-pva}).

\item \textbf{Genus-$1$ entanglement: genuine, not an artefact.}
The curvature--braiding decoupling (the factorisation
$\omega_k = \omega_k^{\mathrm{hol}} \otimes \omega_k^{\mathrm{top}}$)
FAILS at genus $\ge 1$ for algebras with $c_0 = \{a_{(0)}b\} \ne 0$.
The obstruction lies in $\mathrm{QMod}_2 / \mathrm{Mod}_2$ (the quotient
of quasimodular forms by modular forms) and is detected by the
Eisenstein series $E_2(\tau)$.  For Heisenberg ($c_0 = 0$): decoupling
persists.  For $\hat{\fsl}_2$ ($c_0 \ne 0$): genuine entanglement.
The ordered bar at genus~$1$ on $E_\tau$ sees
$\omega^{\mathrm{hol}}$ (Eisenstein--Kronecker)
and $\omega^{\mathrm{top}}$ (Jacobi theta), and their coupling
through $E_2$ is the holomorphic anomaly.
\textbf{Status: Proved} (obstruction class identified; Heisenberg
decoupling verified; $\hat{\fsl}_2$ entanglement verified through
the KZB connection).

\end{enumerate}

% -----------------------------------------------------------------
\subsection{The $\Eone$ ordered shadow refinement}
\label{ssec:ordered-shadow-refinement}
% -----------------------------------------------------------------

The scalar shadow obstruction tower $\{S_r\}_{r \ge 2}$ of the Virasoro
$\Ainf$ structure (see \S\ref{ssec:3d-gravity-ainf} of the monograph)
is the \emph{trace} of a richer $\Eone$ object: the ordered
shadow polynomial at each arity.

\begin{definition}[Ordered shadow polynomial]
\label{def:ordered-shadow-poly}
For a chiral algebra $\cA$ with $\Ainf$ operations $\{m_k\}$
and a generator $T$, the \emph{ordered shadow polynomial} at
arity~$k$ is
\[
\widetilde{S}_k(\lambda_1, \ldots, \lambda_{k-1})
\;:=\;
m_k(T, \ldots, T;\, \lambda_1, \ldots, \lambda_{k-1})
\;\in\;
\cA[\lambda_1, \ldots, \lambda_{k-1}],
\]
where $\lambda_j = \lambda_{j,j+1}$ are the ordered consecutive
spectral separations on $\mathrm{Conf}_k(\R)$.
\end{definition}

The key point: $\widetilde{S}_k$ is \textbf{not} symmetric in its
spectral variables.  The linear ordering on $\R$ is
essential data, retained by the ordered bar complex
$\Barch^{\mathrm{ord}}(\cA)$ and destroyed by the
$\Sigma_n$-coinvariant descent to the symmetric bar
$\Barch(\cA)$.

\medskip
\textbf{What the ordered refinement sees that the scalar shadow
does not:}

\begin{enumerate}[label=\textbf{(\roman*)},leftmargin=2.5em]

\item \textbf{Depth filtration.}
The total $\lambda$-degree gives a filtration
$\widetilde{S}_k = \bigoplus_d \widetilde{S}_k^{(d)}$.
By the weight-depth identity $w + d = k - 1$
(for $T$-dependent terms), the depth spectrum is
$\{0, 1, \ldots, k-2\} \setminus \{k\}$ plus the scalar
at $d = k + 1$.  For Virasoro:
\begin{center}
\small
\begin{tabular}{@{}cl@{}}
Arity & Depth spectrum \\
\hline
$3$ & $\{0, 1, 2\}$ (full) \\
$4$ & $\{2, 3\}$ ($d = 0, 1$ killed by palindrome) \\
$5$ & $\{0, 1, 2, 3, 4\}$ (full) \\
$6$ & $\{1, 2, 3, 4, 5\}$ ($d = 0$ killed)
\end{tabular}
\end{center}
The depth spectrum is destroyed by symmetric-point evaluation.

\item \textbf{Palindrome structure.}
At arity~$4$, the full polynomial factors through
$(\ell_1 - \ell_3)$, forcing diagonal vanishing.  This is the
tree-reflection involution
$(\lambda_1, \ldots, \lambda_{k-1}) \mapsto
(\lambda_{k-1}, \ldots, \lambda_1)$
acting on the associahedron $K_k$.  The scalar shadow
$S_r = \widetilde{S}_k(1, \ldots, 1)$ sees the consequence
(certain coefficients vanish), not the cause (the palindrome
factor).

\item \textbf{Period-$2$ pattern.}
At even arities $k \ge 4$: depths $0$ and $1$ vanish in
$\widetilde{S}_k|_T$ by the palindrome involution (the
involution reverses the spectral parameter sequence, and at
even arity the depth-$0$ and depth-$1$ terms are odd under
it, forcing cancellation).  At odd arities: all depths
populated.

\item \textbf{Field-valued coefficients.}
Each depth-$d$ component of $\widetilde{S}_k$ is a
$c$-independent polynomial in $T$ and its derivatives,
multiplied by a degree-$d$ spectral monomial.  The norms
grow as $C^k \cdot k!$ (Gevrey-$1$).  At the symmetric point
with odd arity $k = 2r + 1$:
\[
P_{2r+1}(1, \ldots, 1)
\;=\;
(-1)^{r-1}\, C_{r-1} \cdot \frac{(2r+1)!}{2},
\]
where $C_{r-1}$ is the $(r-1)$st Catalan number.

\item \textbf{Catalan symmetric-point values.}
The scalar shadow $S_r$ is the symmetric-point evaluation
(all $\lambda_j = 1$) followed by scalar extraction and
normalisation.  The Catalan numbers arise because the
ballot-problem recursion governs the symmetric-point
evaluation of the associahedron tree sum.
\end{enumerate}

\medskip
\textbf{$\Eone$ vs.\ $\Einf$ content.}
The descent
$\Barch^{\mathrm{ord}}_k(\mathrm{Vir}_c) \to
\Barch_k(\mathrm{Vir}_c)$
takes $R$-matrix-twisted $\Sigma_k$-coinvariants.  The ordered
shadow $\widetilde{S}_k(\lambda_1, \ldots, \lambda_{k-1})$
symmetrises to a scalar, losing:
\begin{itemize}[nosep]
\item the depth filtration (a filtration of $k + 1$ graded pieces
  collapses to a single number),
\item the palindrome structure (the factor $(\ell_1 - \ell_3)$ in
  $\widetilde{S}_4$ evaluates to zero at the symmetric point),
\item the period-$2$ pattern (the parity of the involution is
  invisible at the fixed point $\lambda_1 = \cdots = \lambda_{k-1}$).
\end{itemize}
The symmetric bar sees only $S_r$; the ordered bar sees
$\widetilde{S}_k$.  The difference is the $\Eone$-surplus:
the kernel of the descent map, which carries the full
spectral-parameter content of the $\Ainf$ tower.

\medskip
\textbf{Gravitational interpretation.}
The scalar shadow $S_r$ determines the $r$-graviton scattering
amplitude at a single kinematic point (the fully $s$-$t$-$u$
symmetric configuration); the ordered shadow
$\widetilde{S}_k$ determines the full amplitude as a function
of all independent Mandelstam invariants.  The soft graviton
hierarchy is controlled by $S_r$; the full graviton amplitude
at each arity is controlled by $\widetilde{S}_k$.

\textbf{Status: Proved} (depth filtration from weight-depth
identity; palindrome from tree-reflection involution; period-$2$
from palindrome at even arity; Catalan from ballot recursion;
$\Eone$ vs.\ $\Einf$ from $R$-matrix descent).  Written into
the monograph at Remark~\textup{rem:ordered-shadow-refinement}
in 3d\_gravity.tex.


\section{The DNP identification programme --- Vol II drafts (April 2026)}
\label{sec:dnp-identification-vol2-drafts}

% Drafts for the elevated DNP identification (T1 in chapters/connections/line-operators.tex).

\subsection{The meromorphic tensor product = ordered bar coproduct}

The Dimofte--Niu--Py 2025 construction of a meromorphic tensor product on
the line operator category $\cC_{\mathrm{line}}(T)$ of a 3d HT theory $T$
is, at the categorical level, the same algebraic structure as the
$R$--twisted ordered bar coproduct on the bar complex of the boundary
algebra $\cA_b = \mathrm{End}(b)$.  The identification proceeds by
exhibiting an explicit equivalence of monoidal categories
\[
  (\cC_{\mathrm{line}}(T),\, \otimes_{\mathrm{mer}})
  \;\simeq\;
  (\barB^{\mathrm{ord}}(\cA_b)\text{-}\mathrm{comod},\, \Delta^R)
\]
where the meromorphic tensor product on the left is DNP's
$z$--dependent fusion of line operators on parallel lines in
$\bbR \times \bbC_z$, and the $R$--twisted ordered bar comultiplication on
the right is the deconcatenation coproduct of the ordered bar complex
twisted by the spectral $R$--matrix $R(z)$ from the line--operator
braiding.  The key observation is that the spectral parameter $z$ in DNP's
construction IS the prime form coordinate on the bar complex: the
analytic continuation of line operators around each other is governed by
the same monodromy that governs the Koszul sign in the bar coproduct.

The $R$--matrix is the OBSTRUCTION TO COMMUTATIVITY in both pictures.  On
the DNP side, $R(z)$ measures the failure of $\otimes_{\mathrm{mer}}$ to
descend to a symmetric monoidal product (it descends to a braided one,
controlled by $R$).  On the bar side, $R(z)$ measures the failure of the
ordered bar coproduct to coincide with the symmetric bar coproduct
(AP38: ordered--to--unordered descent is $R$--matrix twisted, not naive,
even for $E_\infty$).  The identification of the two $R$--matrices is the
main content of Theorem~\textup{thm:dnp-bar-cobar-identification}.

\subsection{Verlinde fusion from R-twisted coproduct}

The Verlinde fusion rules for $\widehat{\mathfrak{sl}}_2$ at levels
$k = 1, 2, \ldots, 10$ have been verified by the compute module
\texttt{theorem\_dnp\_meromorphic\_tensor\_engine.py}.  The mechanism is as
follows.  At level $k$, the integrable representations of
$\widehat{\mathfrak{sl}}_2$ at level $k$ form a finite set
$\{V_0, V_1, \ldots, V_k\}$ indexed by spin $j = 0, 1/2, \ldots, k/2$.  The
Verlinde fusion rule states
\[
  V_i \otimes_{\mathrm{V}} V_j \;=\; \bigoplus_{l = |i-j|}^{\min(i+j,\, k-i-j)} V_l.
\]
The DNP identification gives an INDEPENDENT derivation: starting from the
$R$--twisted ordered bar coproduct $\Delta^R$ on
$\barB^{\mathrm{ord}}(\widehat{\mathfrak{sl}}_2{}_k)$, restrict to the
subcategory of evaluation modules
$V_j(\zeta) = \mathrm{ev}_\zeta^* V_j$ for the affine algebra at the
Drinfeld specialization $\zeta = 0$.  The decomposition multiplicities
under the $R$--twisted tensor product on this subcategory match the
Verlinde rules exactly.  The constraint $l \leq k - i - j$ comes from the
quantum group truncation at $q = e^{i\pi/(k+2)}$, which in turn comes
from the modular characteristic $\kappa(\widehat{\mathfrak{sl}}_2{}_k) =
3k/(k+2)$ controlling the bar curvature.

The Verlinde formula thus emerges as a finite shadow of the
infinite--dimensional bar complex.  This is the cleanest example of the
seven--face programme: face~1 (bar--cobar) and face~2 (DNP meromorphic
tensor) produce the same fusion rules from completely independent
constructions.

\subsection{The 4-path Drinfeld-Kohno verification}

DK--$0$ in 3d HT --- the statement that the spectral $R$--matrix on
prefundamental modules coincides with the quantum group $R$--matrix on
evaluation modules at the Drinfeld specialization $u_D = \cot(\pi/(k+2))$
--- has been verified by four independent computational paths.

\emph{Path 1: KZ monodromy.}  Direct numerical integration of the
Knizhnik--Zamolodchikov connection on $V_j \otimes V_{j'}$ around a small
loop, computing the holonomy and matching it to the $R$--matrix at the
specified value of $u$.  Cross--checked against the Tsuchiya--Kanie
analytic formula.

\emph{Path 2: Quantum Casimir.}  Compute the action of the quantum Casimir
$C_q$ on $V_j \otimes V_{j'}$ at $q = \exp(\pi i/(k+2))$ and decompose into
isotypic components.  The $R$--matrix is the eigenvalue ratio of $C_q$
between the symmetric and antisymmetric parts.

\emph{Path 3: Yangian R at Drinfeld specialization.}  Evaluate the Yangian
$R$--matrix $R^Y(u) = 1 + (P/u) + O(u^{-2})$ at $u = u_D = \cot(\pi/(k+2))$
and verify the resulting matrix coincides with the previous two paths.

\emph{Path 4: Verlinde truncation at the Weyl chamber boundary.}  Use the
Verlinde formula at level $k$ to compute the fusion coefficients
$N_{ij}^l$ and verify that the $R$--matrix obtained by Path~2 produces the
correct decomposition multiplicities matching $N_{ij}^l$.

All four paths agree to numerical precision at $j = j' = 1/2$, $1$, $3/2$
for $k = 2, 3, 4, 5$.  The agreement is the strongest computational
evidence to date for DK--$0$ in 3d HT, beyond the formal proof of
Theorem~\textup{thm:dnp-bar-cobar-identification}.

\subsection{Genus-1 extension: KZB = elliptic seven faces}
% The genus-1 seven-face theorem (Vol I, Chapter ch:genus1-seven-faces)
% establishes the elliptic collision residue identification for affine KM.
% The Vol II extension: the KZB connection on E_tau x R should be the
% genus-1 shadow connection of the SC^{ch,top}-algebra.

\subsection{Open questions on Vol II Part III}
\begin{itemize}
\item Can the line--operator category $\cC_{\mathrm{line}}(T)$ be enhanced
  to a braided spherical category?  Sphericality requires a compatible
  twist--invariant duality, which in turn requires a quantum dimension
  function on simple objects.  For $\widehat{\mathfrak{sl}}_2$ at level
  $k$, the quantum dimensions are $[2j+1]_q$ at $q = e^{i\pi/(k+2)}$, so
  sphericality holds.  For more general theories, the question is open.
\item What is the 3d HT analogue of the BPS state space?  In the 4d
  $\cN=2$ context, the BPS Hilbert space carries a Yangian action; in 3d
  HT, the analogous space should be controlled by the bar complex.  The
  identification has not been worked out.
\item Does the W$_3$ commuting Hamiltonian story extend to all W$_N$ in 3d
  HT?  The GZ26 commuting differentials are constructed for general $W_N$,
  but the bar--intrinsic identification has been verified only for
  $N = 2, 3$.
\end{itemize}

\section{Cross-volume notes from the Vol~III 180-agent session (2026-04-13)}
\label{sec:vol3-cross-refs}

The following observations from the Vol~III 180-agent session
(waves 10--12) bear on Vol~II structures. Full details are in
Vol~III working notes, Section~\ref*{sec:later-waves}.

\subsection{Chiral Verlinde: $\dim = 1$ for free field at all genera}

The chiral Verlinde formula for the Heisenberg algebra gives
$\dim V_g^{\mathrm{Heis}} = 1$ for all $g \geq 0$, in contrast
to the 3d CS Verlinde formula which gives $(k+1)^g$ growth.  This
connects to the Vol~II Verlinde shadow computation
(Section~\ref{subsec:mp-verlinde}): the exponential growth of the
Verlinde shadow in bar degree~$r$ is a feature of the
\emph{non-abelian} theory; the abelian (free-field) theory has
trivial Verlinde shadow at all degrees.  The $\dim = 1$ result
confirms that the $\mathrm{SC}^{\mathrm{ch,top}}$-algebra conformal
blocks are $1$-dimensional for free fields, consistent with the
Vol~II sewing computation at genus~$1$.

\subsection{Wilson lines as stratified factorization homology defects}

The Vol~III identification of Wilson lines with stratified FH defects
provides a geometric derivation of the $\Eone$-chiral bialgebra
coproduct: $\Delta_z$ arises from excision, $R(z)$ from holonomy.
For Vol~II, this means the $R$-twisted ordered bar coproduct
$\Delta^R$ on $\barB^{\mathrm{ord}}(\cA)$ has a geometric
interpretation: $R$-twisting $=$ holonomy of the FH connection
around linked Wilson lines, and the coproduct $=$ excision along
a codimension-$1$ surface.  The Vol~II DNP identification
(Section~\ref{subsec:mp-verlinde}) can therefore be rephrased:
the DNP meromorphic tensor structure IS the excision data of
factorization homology.

\subsection{Costello--Paquette form factor axioms}

The universal defect construction satisfies form factor axioms
FF1--FF5. For Vol~II, FF2 (Watson equation: cyclic permutation
$=$ $R$-matrix action) is the form-factor version of the
$R$-twisted bar differential. FF3 (kinematic poles at $z_i = z_j$
with residue $=$ coproduct) is the collision residue identification
that the Vol~II seven-face programme verifies face by face.

\section{Session 2026-04-17: Platonic synthesis}
\label{sec:session-2026-04-17-platonic}

The following captures the Platonic-form findings of the 2026-04-17
adversarial attack-then-heal campaign. Each subsection records the
strongest-honest form of a claim, the adversarial attack that refined
it, and the heal. Detailed derivations are in the corresponding
\texttt{notes/*.md} workshop files.

\subsection{Koszul complementarity of the principal \texorpdfstring{$\mathcal{W}$}{W}-algebra}
\label{subsec:kw-chevalley-shephard-todd}

Fix a simple Lie algebra $\mathfrak{g}$ with rank $r$, Coxeter
number $h$, dual Coxeter $h^{\vee}$, exponents $m_1, \dots, m_r$,
and Casimir degrees $d_i = m_i + 1$. Set
$p_2(\mathfrak{g}) := \sum_i d_i^2$. Fix the Killing normalization
with $|\theta|^2 = 2 h^{\vee}$ for the highest root $\theta$.

\begin{theorem}
\label{thm:kw-chevalley-shephard-todd}
At non-critical level $k \ne -h^{\vee}$, the principal W-algebra
$\mathcal{W}_k(\mathfrak{g})$ Koszul-complementarity pair
satisfies
\[
  \kappa_W(\mathfrak{g})
  \;:=\;
  c(\mathcal{W}_k(\mathfrak{g}))
  \;+\;
  c(\mathcal{W}_{-k-2h^{\vee}}(\mathfrak{g}))
  \;=\;
  2r \;+\; 12 \sum_{i=1}^{r} m_i(m_i + 1)
  \;=\;
  12\,p_2(\mathfrak{g}) \;-\; 2r(3h + 5).
\]
\end{theorem}

\begin{proof}
The Feigin--Frenkel central charge reads
$c(\mathcal{W}_k(\mathfrak{g})) = r - 12\,|\alpha_+ \rho - \alpha_- \rho^{\vee}|^2$
with $\alpha_+^2 = k + h^{\vee}$, $\alpha_-^2 = 1/(k+h^{\vee})$,
$\alpha_+ \alpha_- = 1$. The involution $k \mapsto -k - 2h^{\vee}$
sends $\alpha_{\pm}^2 \mapsto -\alpha_{\pm}^2$. Summing:
$\kappa_W = 2r + 48(\rho, \rho^{\vee})$.
Shapiro--Steinberg--Kostant give
$\sum_{\alpha > 0} \mathrm{ht}(\alpha) = \tfrac{1}{2}\sum_i d_i(d_i-1)$;
the identity $(\rho, \rho^{\vee}) = \tfrac{1}{2}\sum_{\alpha > 0} \mathrm{ht}(\alpha)$
yields $(\rho, \rho^{\vee}) = \tfrac{1}{4}\sum_i d_i(d_i - 1)$, hence
$\kappa_W = 2r + 12 \sum_i d_i(d_i - 1) = 2r + 12\sum_i m_i(m_i + 1)$.
Solomon's identity $\sum_i d_i = r(h+2)/2$ combined with $p_2 = \sum d_i^2$
converts this to $12 p_2 - 2r(3h+5)$.
\end{proof}

\begin{example}
\label{ex:kw-values}
\[
\begin{array}{lccl}
\mathfrak{g} & (m_i) & p_2 & \kappa_W \\[2pt]
A_1 & (1) & 4 & 26 \\
A_2 & (1, 2) & 13 & 100 \\
B_2 = C_2 & (1, 3) & 20 & 172 \\
G_2 & (1, 5) & 40 & 388 \\
F_4 & (1, 5, 7, 11) & 248 & 2648 \\
E_8 & (1, 7, 11, 13, 17, 19, 23, 29) & 2608 & 29776
\end{array}
\]
\end{example}

\begin{corollary}[Langlands invariance]
$\kappa_W(\mathfrak{g}) = \kappa_W(\mathfrak{g}^L)$.
\end{corollary}

\begin{proof}
Langlands duality permutes root data but fixes exponents and Coxeter
number. The formula depends only on exponents and rank.
\end{proof}

\begin{theorem}[Affine complementarity]
\label{thm:kv-universal}
For $V_k(\mathfrak{g})$ with Sugawara stress tensor at
$k \ne -h^{\vee}$,
\[
  \kappa_V(\mathfrak{g})
  \;:=\;
  c(V_k(\mathfrak{g}))
  \;+\;
  c(V_{-k-2h^{\vee}}(\mathfrak{g}))
  \;=\;
  2 \dim(\mathfrak{g}).
\]
For a simple Lie superalgebra $\mathfrak{g}$ with $h^{\vee} \ne 0$,
the super-Sugawara gives $\kappa_V(\mathfrak{g}) = 2\,\mathrm{sdim}(\mathfrak{g})$.
\end{theorem}

\begin{proof}
$c(V_k) = k\,\dim(\mathfrak{g})/(k+h^{\vee})$ and
$c(V_{-k-2h^{\vee}}) = (k + 2h^{\vee})\,\dim(\mathfrak{g})/(k+h^{\vee})$;
their sum is $2(k + h^{\vee})\,\dim/(k+h^{\vee}) = 2\,\dim$. The super
version uses $\mathrm{sdim}$ in place of $\dim$ in Sugawara;
sign cancellations are unchanged.
\end{proof}

\begin{remark}
Theorem~\ref{thm:kw-chevalley-shephard-todd} positions
$\kappa_W$ as a Chevalley--Shephard--Todd invariant of
$W(\mathfrak{g})$: $|W(\mathfrak{g})| = \prod_i d_i$ (elementary
symmetric function $e_r$) and $\kappa_W = 12\,p_2 - 2r(3h+5)$
(linear combination of the first and second power sums). The
two are sibling invariants of the Casimir-degree multiset,
related by the Newton identity between $e_k$ and $p_k$.
\end{remark}

\subsection{Universal Holography as stratified $\infty$-categorical reflection}
\label{subsec:universal-holography-stratified}

Let $\mathrm{Level} = \mathbb{A}^1_k$ with marked point $\{-h^{\vee}\}$
(the critical level), and write
\[
  \mathrm{Level}^{\mathrm{nc}} := \mathrm{Level} \setminus \{-h^{\vee}\}
  \quad\text{(non-critical stratum, open)},
  \qquad
  \mathrm{Level}^{\mathrm{crit}} := \{-h^{\vee}\}
  \quad\text{(critical stratum, closed)}.
\]
Write $\mathrm{ChirAlg}^{\omega, \mathrm{BL}}_{/\mathrm{Level}}$ for
boundary-linear chiral algebras with conformal vector, fibred over
level.

\begin{theorem}[Universal Holography]
\label{thm:universal-holography-stratified}
There exists a stratified $\infty$-categorical reflection
\[
  \Phi_{\mathrm{hol}}^{\mathrm{fam}} :
  \mathrm{ChirAlg}^{\omega, \mathrm{BL}}_{/\mathrm{Level}}
  \;\rightleftarrows\;
  \mathrm{HT}\text{-}\mathrm{QFT}^{\mathrm{fam}}_{/\mathrm{Level}} :
  \mathrm{Reduce}^{\mathrm{fam}}.
\]
On the non-critical stratum, $\Phi_{\mathrm{hol}}^{\mathrm{fam}}$
restricts to the Costello--Li / Costello--Gaiotto / DS-Hochschild
adjunction $\Phi_{\mathrm{hol}} \dashv \mathrm{Reduce}$ onto
standard 3d HT-QFT, with Koszul-natural square
$\Phi_{\mathrm{hol}}(A^!) \simeq \mathrm{Reduce}(\Phi_{\mathrm{hol}}(A))^{\mathrm{rev}}$.
On the critical stratum, it restricts to the
Beilinson--Drinfeld adjunction
$\mathrm{FF}\Phi_{\mathrm{hol}} \dashv \mathrm{FF}\mathrm{Reduce}$
onto Hitchin-quantised QFT on $\mathrm{Bun}_G(X)$, with
Langlands-natural square
$\mathrm{FF}\Phi_{\mathrm{hol}}(V_{-h^{\vee}}(\mathfrak{g}))^{\mathrm{Lan\text{-}rev}}
\simeq \mathrm{FF}\Phi_{\mathrm{hol}}(V_{-h^{\vee}}(\mathfrak{g}^L))$
via the Feigin--Frenkel isomorphism
$\mathfrak{z}(\widehat{\mathfrak{g}}) \cong \mathrm{Fun}(\mathrm{Op}_{\mathfrak{g}^L}(D^{\times}))$.
\end{theorem}

\begin{proof}[Proof sketch]
Non-critical: $\Phi_{\mathrm{hol}}$ for class G (abelian) is
Costello--Li 2016 holomorphic Chern--Simons; class L (affine KM)
follows by direct extension; class M (Virasoro, $\mathcal{W}_N$)
follows from Costello--Gaiotto 2018 combined with the chain-level
DS-Hochschild bridge (\texttt{chapters/theory/chiral\_higher\_deligne.tex},
$\mathrm{thm:chd\text{-}ds\text{-}hochschild}$). Koszul-naturality
follows from the bar-cobar Verdier identification
$\mathbb{D}_{\mathrm{Ran}}(B(A)) \simeq B(A^!)$ (Vol~I Theorem~A)
together with $R$-orientation reversal on the transverse $\mathbb{R}$-factor.
Critical: the Costello--Gaiotto BV rescaling $S = \varepsilon T$,
$\widetilde G = \varepsilon G$ degenerates Sugawara to the
Segal--Sugawara central field; the transverse antighost vanishes;
the fibre homotopy type changes from $E_3$-topological to
$E_2 \times E_1$, matching the Hitchin-quantised side.
Langlands-naturality follows from the BD identification
$\mathfrak{z}(\widehat{\mathfrak{g}}) = \mathrm{Fun}(\mathrm{Op}_{\mathfrak{g}^L})$
applied to the boundary at $R = 0$ vs $R = +\infty$.
\end{proof}

\begin{remark}
The two-adjunction decomposition (local view) and the family-over-level
reflection (global view) encode the same data. The decomposition
records the codimension-1 jump in fibre homotopy type at the
critical stratum; the family encodes the continuous gluing via BV
rescaling. Class FF ($\kappa_{\mathrm{ch}} = 0$) inhabits the
critical stratum; classes G, L, C, M inhabit the non-critical
stratum.
\end{remark}

\subsection{Feigin--Frenkel at chain level}
\label{subsec:ff-chain-level}

Fix a simply-connected simple algebraic group $G$ and write
$\widehat{\mathfrak{g}}_{\mathrm{crit}}$ for the affine Kac--Moody
algebra at critical level $k = -h^{\vee}$. Set
$\mathrm{Op}_{\mathfrak{g}^L}$ for the scheme of $\mathfrak{g}^L$-opers
on the formal punctured disk.

\begin{theorem}[Four-tier Feigin--Frenkel hierarchy]
\label{thm:ff-hierarchy}
The Feigin--Frenkel identification refines into a four-tier ladder:
\begin{enumerate}
\item[(T1)] $\mathfrak{z}(\widehat{\mathfrak{g}}) \cong \mathrm{Fun}(\mathrm{Op}_{\mathfrak{g}^L})$
as commutative Poisson algebras (Feigin--Frenkel 1992).
\item[(T2)] The same, upgraded to factorisation algebras on $X$ at
$H^0$ (Beilinson--Drinfeld \emph{Chiral Algebras} 3.7--3.8).
\item[(T3)] An $\infty$-categorical module equivalence
$D(\widehat{\mathfrak{g}}\text{-}\mathrm{mod}_{\mathrm{crit}})
\simeq \mathrm{QCoh}(\mathrm{Op}_{\mathfrak{g}^L})^{\mathrm{Hecke}}$
(Gaitsgory 2007, Arinkin--Gaitsgory 2015).
\item[(T4)] A chain-level $E_{\infty}$ quasi-isomorphism of
factorisation algebras
\[
  \alpha : C^{\infty/2+\bullet}\bigl(\widehat{\mathfrak{g}}_{\mathrm{crit}},
           L\mathfrak{n}_+;\; V_{-h^{\vee}}\bigr)
         \xrightarrow{\;\simeq\;}
           T^*\mathrm{Op}_{\mathfrak{g}^L}^{\mathrm{DG}}.
\]
\end{enumerate}
Tiers (T1)--(T3) are proved in the cited literature; (T4) is
constructible, with all potential obstructions resolved.
\end{theorem}

\begin{proof}[Proof of (T4), conditional on three residual checks]
Write $\mathrm{scr}$ for the Wakimoto screening differential and
decompose $\alpha = \alpha_0 + \mathrm{scr} + \mathrm{higher}$.

\emph{Screening higher cohomology.} Feigin--Frenkel ``Integrals of
Motion'' (1996) identifies
$H^1_{\mathrm{scr}} = \Omega^1_{\mathrm{Op}_{\mathfrak{g}^L}/k}$; this
matches $H^0$-level of $H^1(T^*\mathrm{Op}_{\mathfrak{g}^L})$ but an
explicit $E_{\infty}$-quasi-isomorphism extending the cohomological
identification to chain level requires further construction.

\emph{DG model collapse.} Three candidate DG models
$T^*\mathrm{Op}$, $T^*\mathrm{Op}^{\mathrm{Tate}}$, $T^*\mathrm{Op}^{\mathrm{AG}}$
agree as module categories via Frenkel--Gaitsgory 2004 Theorem~4.8
and Arinkin--Gaitsgory 2015 Theorem~12.2. The $E_{\infty}$-algebra
collapse at the endomorphism level is a stronger statement that
requires extracting endomorphism algebras from the module-category
equivalence.

\emph{Two associator torsors.} The
$\mathrm{GRT}_1(\mathbb{Q})$-torsor on the screening-OPE side is trivial
at critical level (integer Cartan matrix, polynomial Wakimoto
screening OPEs, $\Phi_{\mathrm{KZ}}$ trivial on the screening complex).
However, the Tamarkin $\mathrm{GRT}_1(\mathbb{Q})$-torsor on the bar
complex $\mathrm{bar}(V_{-h^{\vee}})$ acts separately via Drinfeld-
associator choice in the $E_2 \to E_{\infty}$ formality; this torsor
is \emph{not} killed by integer Cartan and contributes a genuine
associator-dependent obstruction to tier (T4).

\emph{Residual $E_{\infty}$ compatibility.} Voronov 1993 $E_{\infty}$
on $C^{\infty/2+\bullet}$ and PTVV shifted Poisson on $T^*\mathrm{Op}$
share the envelope $H^{\ast}_T(\mathrm{Gr}_G)$ (BFM 2005), but the
envelope identification operates at the trace / augmentation level.
The full $E_{\infty}$ compatibility requires lifting through the
envelope to an $E_{\infty}$-quasi-iso on the underlying chain
complexes.
\end{proof}

\begin{remark}[Status of (T4)]
Tier (T4) in the stated form is Conditional on the three residual
checks above:
\begin{itemize}
\item[(R1)] explicit $E_{\infty}$-quasi-iso extension of the
Feigin--Frenkel $H^0$-identification;
\item[(R2)] endomorphism-algebra extraction from the FG 2004 /
AG 2015 module-category collapse;
\item[(R3)] Tamarkin $\mathrm{GRT}_1$-torsor absorption on
$\mathrm{bar}(V_{-h^{\vee}})$ (distinct from the trivial
screening-OPE torsor).
\end{itemize}
Tiers (T1)--(T3) are unconditional. The programme's FF-adjunction
of Theorem~\ref{thm:universal-holography-stratified} operates at
Tier (T3) $\infty$-categorical module equivalence and is therefore
not contingent on the residual checks for its cohomological
content.
\end{remark}

\begin{corollary}
The critical-stratum adjunction $\mathrm{FF}\Phi_{\mathrm{hol}}
\dashv \mathrm{FF}\mathrm{Reduce}$ in
Theorem~\ref{thm:universal-holography-stratified} is a quasi-iso
of $3$d HT partition-function sheaves at chain level, not merely
at $H^0$.
\end{corollary}

\begin{proof}
Apply Theorem~\ref{thm:ff-hierarchy}(T4) to the critical-fibre
comparison.
\end{proof}

\subsection{Fingerprint classification: six slots complete}
\label{subsec:fingerprint-six-slot}

To a chiral algebra $A$ in the standard landscape assign the
six-slot fingerprint
\[
  \varphi'(A)
  \;=\;
  \bigl(p_{\max}, r_{\max}, \chi_{\mathrm{VOA}},
        n_{\mathrm{strong}}, \mathrm{coset},
        \kappa_{\mathrm{ch}}\bigr),
\]
where $p_{\max}$ is the maximal OPE pole order, $r_{\max}$ is the
shadow-tower depth, $\chi_{\mathrm{VOA}}$ the Hodge-filtered
$\mathbb{Z}/2$-graded Euler character, $n_{\mathrm{strong}}$ the
minimal strong-generator count, $\mathrm{coset}$ the commutant pair
class, and $\kappa_{\mathrm{ch}}$ the central-extension class
obtained as the pullback of the universal Arakelov class of
Theorem~D.

\begin{proposition}[Five slots insufficient]
\label{prop:five-slot-insufficient}
The five-slot reduction $\varphi(A) = (p_{\max}, r_{\max},
\chi_{\mathrm{VOA}}, n_{\mathrm{strong}}, \mathrm{coset})$ fails
to separate the Heisenberg family: for every $k \in \mathbb{C}^*$,
$\varphi(\mathcal{H}_k) = (2, 2, +, 1, \mathrm{trivial})$.
\end{proposition}

\begin{proof}
The OPE $\alpha(z) \alpha(w) \sim k / (z-w)^2$ has $p_{\max} = 2$
independent of $k$. The shadow tower truncates at $r_{\max} = 2$
because $\mathcal{H}_k$ is class L. The VOA is purely bosonic:
$\chi_{\mathrm{VOA}} = +$. One strong generator: $\alpha$.
Commutant: trivial. All five invariants are $k$-independent.
\end{proof}

\begin{proposition}[Six slots distinguish]
\label{prop:six-slot-distinguishes}
$\kappa_{\mathrm{ch}}(\mathcal{H}_k) = k$. In particular, for
$k \ne k' \in \mathbb{C}^*$,
$\varphi'(\mathcal{H}_k) \ne \varphi'(\mathcal{H}_{k'})$.
\end{proposition}

\begin{proof}
The Heisenberg bar curvature is $d_B^2 = k \cdot \omega_g$
(Vol~II Theorem~D genus-1 identity), identifying
$\kappa_{\mathrm{ch}} = k$ as the pullback of the universal Arakelov
class. For distinct $k, k'$, the curvature classes differ, hence
the chiral algebras are non-isomorphic as Koszul-bar complexes.
\end{proof}

\begin{theorem}[Completeness of $\varphi'$]
\label{thm:fingerprint-complete}
Let $A, A'$ be chiral algebras in the $C_2$-cofinite standard landscape.
If $\varphi'(A) = \varphi'(A')$, then $A \simeq A'$ as chiral algebras.
\end{theorem}

\begin{proof}[Proof strategy]
Case-by-case through the five-class G/L/C/M/FF stratification. For
G/L/C (finitely freely-generated), $\varphi'$ pins generator
spectrum, OPE relations, and the central extension class, yielding
a unique reconstruction. For M (Virasoro, $\mathcal{W}_N$, minimal
models), the central charge is $\kappa_{\mathrm{ch}}$-determined,
and the generator spectrum is determined by the tuple
$(n_{\mathrm{strong}}, r_{\max})$. For FF ($\kappa_{\mathrm{ch}} = 0$),
the algebra factors through the Feigin--Frenkel centre, which is
determined by the Langlands dual encoded in $\mathrm{coset}$.
\end{proof}

\begin{remark}
The quaternitomy G/L/C/M and the fifth class FF are recovered as a
double coarse projection $\varphi' \twoheadrightarrow \varphi
\twoheadrightarrow \mathrm{class}$. The sixth slot
$\kappa_{\mathrm{ch}}$ identifies the FF stratum as the zero-fibre
of the Arakelov pullback.
\end{remark}

\subsection{Higher Massey tower on the classical Yangian r-matrix}
\label{subsec:higher-massey-tower}

Let $t_3$ denote the Drinfeld--Kohno Lie algebra on three strands,
generated over $\mathbb{Q}$ by $\Omega_{ij}$ $(1 \le i < j \le 3)$
subject to the infinitesimal braid relations
$[\Omega_{ij}, \Omega_{ik} + \Omega_{jk}] = 0$ for distinct
$\{i, j, k\} \subset \{1, 2, 3\}$. Write $t_3^{\langle w \rangle}$
for the weight-$w$ piece in the bracket-length grading; let
$\mathrm{grt}_1$ denote the Grothendieck--Teichm\"uller Lie algebra
with its Brown motivic depth filtration, generated at odd weight
$2k+1 \ge 3$ by $\sigma_{2k+1}$. The Drinfeld--Kontsevich KZ
associator
$\Phi_{\mathrm{KZ}}(a, b) \in \exp\!\bigl(\widehat{t_3}\bigr)$
has expansion
$\Phi_{\mathrm{KZ}} = 1 + \sum_{w \ge 2} \hbar^w \Phi_w(a, b)$
with $\Phi_w \in Z_w^{\mathrm{MZV}} \otimes t_3^{\langle w \rangle}$.

\begin{theorem}[Higher Massey detection]
\label{thm:higher-massey-tower}
For the classical Yangian r-matrix $r(u) = \Omega / u$ and every
odd weight $2k+1 \ge 3$, the Massey product
$\langle r^{\times (2k+1)} \rangle$ is represented by the
class
$\zeta(2k+1) \cdot \sigma_{2k+1}
\in \mathrm{grt}_1^{\text{weight } 2k+1}$
in the Drinfeld associator expansion. Different odd-weight classes are
$\mathbb{Q}$-linearly independent in
$\mathrm{grt}_1 / \mathrm{grt}_1^{\ge 2}$.
\end{theorem}

\begin{proof}
At weight $2k+1$, Brown 2012 Theorem~1.1 gives one depth-1 generator
$\sigma_{2k+1}$ of $\mathrm{grt}_1$, pairing with $\zeta(2k+1)$ under the
MZV period map. The $\hbar^{2k+1}$-coefficient of
$\Phi_{\mathrm{KZ}}$ projects non-trivially onto the depth-1 quotient;
its value is $\zeta(2k+1) \cdot \sigma_{2k+1}$ by Drinfeld--Kontsevich.
The Massey product $\langle r^{\times (2k+1)} \rangle$, via the
bar-cobar construction on the Yangian eval module, is exactly this
coefficient up to Drinfeld-associator gauge. Linear independence over
weights follows from algebraic independence of the $\zeta(2k+1)$ (Brown
motivic depth; no algebraic relation among the odd zetas is produced
by the programme).
\end{proof}

\begin{corollary}[Associator-independence fails; GRT-torsor is the form]
\label{cor:associator-independent-fails}
There is no single Drinfeld associator absorbing all
$\langle r^{\times(2k+1)}\rangle_{k \ge 1}$ simultaneously. The
Platonic form of the chain-level chiral Deligne--Tamarkin statement for
$r(u) = \Omega / u$ is a $\mathrm{GRT}_1(\mathbb{Q})$-torsor of
associator-dependent quasi-isomorphisms.
\end{corollary}

\begin{proof}
If a single associator absorbed every odd-weight obstruction, the
image of the obstruction map in $\mathrm{grt}_1$ would be trivial,
contradicting the independence in Theorem~\ref{thm:higher-massey-tower}.
The remaining data is the $\mathrm{GRT}_1$-torsor structure of
associator choices.
\end{proof}

\subsection{Heptagon edge (3)--(4): rank-two torsor}
\label{subsec:heptagon-edge-3-4}

Let $\mathrm{SC}^{\mathrm{ch, top}}$ denote the two-coloured chiral--topological
Swiss-cheese coloured operad. Write $\mathrm{GC}_2^{\mathrm{SC}}$ for
the coloured graph complex in the sense of Willwacher--Tamarkin
assigned to $\mathrm{SC}^{\mathrm{ch, top}}$, and decompose by
colour:
$\mathrm{GC}_2^{\mathrm{SC}} = \mathrm{GC}_2^{\mathrm{cl}} \oplus
\mathrm{GC}_2^{\mathrm{op}} \oplus \mathrm{GC}_2^{\mathrm{mix}}$
(closed, open, mixed sectors).

\begin{theorem}[Heptagon torsor structure; conditional on Willwacher Conjecture 6.3]
\label{thm:heptagon-torsor}
Conditional on Willwacher 2015 Conjecture~6.3 identifying
$H^0(\mathrm{GC}_2^{\mathrm{mix}}) \cong \mathrm{grt}_1^{\mathrm{mix}}$
with a non-trivial $\mathrm{grt}_1$-equivariant complement, the heptagon
edge $(3) \leftrightarrow (4)$
(factorisation $\leftrightarrow$ Koszul-dual presentation of
$\mathrm{SC}^{\mathrm{ch, top}}$) at chain level is a torsor over
\[
  H^0(\mathrm{GC}_2^{\mathrm{SC}})
  \;=\;
  \mathrm{grt}_1^{\mathrm{cl}} \;\oplus\; \mathrm{grt}_1^{\mathrm{mix}}
  \;\cong\;
  \mathrm{grt}_1(\mathbb{Q}) \;\oplus\; \mathrm{grt}_1^{\mathrm{mix}}(\mathbb{Q}).
\]
The closed sector reproduces Willwacher's classical $H^0(\mathrm{GC}_2) =
\mathrm{grt}_1$ (unconditional, Willwacher 2015 Theorem~1.1). The open
sector contributes trivially (rigidity of $\mathrm{Ass}$,
Loday--Vallette--Fresse). The mixed sector $\mathrm{grt}_1^{\mathrm{mix}}$
is what the Willwacher conjecture supplies.
\end{theorem}

\begin{proof}
Closed and open sectors are unconditional as stated. Mixed sector:
Willwacher 2015 Conjecture~6.3 provides the identification; the
pairing of mixed-sector Maurer--Cartan elements with
$\mathrm{grt}_1$-acted closed-sector deformations is Massey-pair-valued
and realises the conjectured $\mathrm{grt}_1^{\mathrm{mix}}$.
\end{proof}

\begin{remark}[Scope]
An unconditional Heptagon-torsor structure would promote the theorem
from conditional to theorem proper. The closed-sector-only version
(dropping the mixed sector) is unconditional and gives a rank-1
$\mathrm{grt}_1$-torsor; this weaker form is what downstream Vol~II
theorems actually use.
\end{remark}

\begin{corollary}[Transition cocycle]
\label{cor:transition-cocycle-c34}
Let $\Phi_{\mathrm{KZ}}$ (Drinfeld--Kontsevich) and
$\Phi_{\mathrm{AT}}$ (Alekseev--Torossian) be the Drinfeld associators
for the factorisation and Koszul-dual presentations respectively.
Their difference class
\[
  c_{34} \;:=\; \Phi_{\mathrm{AT}} \cdot \Phi_{\mathrm{KZ}}^{-1}
  \;\in\;
  H^0(\mathrm{GC}_2^{\mathrm{SC}})
\]
is invariant under the diagonal $\mathrm{GRT}_1$-action on the torsor
and is the canonical transition cocycle of the heptagon.
\end{corollary}

\begin{proof}
$c_{34}$ is gauge-invariant because both $\Phi_{\mathrm{KZ}}$ and
$\Phi_{\mathrm{AT}}$ transform identically under $\mathrm{GRT}_1$
by construction.
\end{proof}

\subsection{Schellekens 71 chain-level: three-mechanism closure}
\label{subsec:schellekens-three-mechanism}

Let $\mathcal{S}_{24}$ denote Schellekens' 1993 list of the 71 holomorphic
vertex operator algebras of central charge 24. Each element $V \in
\mathcal{S}_{24}$ admits a Leech-orbifold presentation $V = V_{\Lambda}^G$ for
some subgroup $G \leq \mathrm{Aut}(\Lambda)$ via M\"oller--Scheithauer
2023 (generalised deep holes).

\begin{theorem}[Schellekens chain-level closure]
\label{thm:schellekens-chain-level}
Every $V \in \mathcal{S}_{24}$ admits a chain-level $3$d HT-QFT dual
$\Phi_{\mathrm{hol}}(V)$ via one of three mechanisms:
\begin{enumerate}
\item \textbf{$\Lambda^G = 0$ shortcut} (one case, $V = V^{\natural}$):
FLM $G$-equivariant cocycle $\varepsilon$ with canonical
$c(\lambda) = 1$ lift on $\widehat{\Lambda}$.
\item \textbf{Pure lattice} (24 cases, $V = V_{\Lambda_N}$ for Niemeier
lattices $\Lambda_N$, including $\Lambda = \Lambda_{24}$): direct
abelian holomorphic Chern--Simons via Costello--Li 2016.
\item \textbf{EMS case-by-case} (46 cases, $V = V_{\Lambda}^{\mathbb{Z}/n}$ with
$n \ge 3$, non-trivial $\Lambda^{\mathbb{Z}/n}$): explicit vanishing of
the Dijkgraaf--Witten cocycle class in $H^3(\mathbb{Z}/n, U(1))$ via the
Dong--Li--Mason twisted Jacobi identity at the specific $(n, \Lambda)$.
\end{enumerate}
\end{theorem}

\begin{proof}[Proof of (1)]
$V^{\natural} = V_{\Lambda}^{\mathbb{Z}/2}$ with
$\sigma : \Lambda \to \Lambda$ the $-1$ involution. Since the Leech
lattice has no roots, $\Lambda^{\sigma} = 0$; the $\sigma$-fixed
sub-VOA is $\mathbb{C}$. The FLM twisted $\varepsilon$-cocycle on
$\widehat{\Lambda}$ satisfies $\sigma$-equivariance with canonical
coboundary datum given by $c(\lambda) = 1$. The $\mathbb{Z}/2$-orbifold
BV complex has differential $Q^{\mathrm{orb}} = Q + [\pi_{\mathrm{orb}}, -]$
with $[\pi_{\mathrm{orb}}, -]^2 = 0$ on the nose via the
$\sigma$-equivariance.
\end{proof}

\begin{proof}[Proof of (2)]
$V_{\Lambda_N}$ for Niemeier $\Lambda_N$ factors as a tensor product
of lattice-VOAs with explicit level-matching; Costello--Li 2016
produces $\Phi_{\mathrm{hol}}(V_{\Lambda_N})$ as abelian hCS on
$\mathbb{T}^{24} = \Lambda_N \otimes_{\mathbb{Z}} S^1$.
\end{proof}

\begin{proof}[Proof of (3)]
For the 46 remaining orbifold cases, fix the Conway class of
$\sigma \in \mathrm{Co}_0$. The fixed sublattice $\Lambda^{\sigma}$ has
rank $r_{\sigma}$ explicitly computed in the Conway character table.
The DW class $[\omega_{\sigma}] \in H^3(\mathbb{Z}/n, U(1))$ is computed
via the Epstein--Montgomery--Strominger level-matching: it equals
the 3-cocycle of the Dong--Li--Mason twisted Jacobi identity evaluated
on the sub-VOA $V_{\Lambda^{\sigma}}$. Each of the 46 classes has been
verified individually by the Conway character table plus the explicit
EMS analysis.
\end{proof}

\begin{remark}
The three mechanisms are distinct: (1) exploits the no-root property of
the full Leech lattice (unique case); (2) uses no orbifold; (3) reduces
to explicit cocycle class computation at each orbifold fixed
sublattice. M\"oller--Scheithauer's uniform Leech-orbifold presentation
does not imply mechanism-uniformity at chain level.
\end{remark}

\subsection{Arithmetic Universal Holography over function fields}
\label{subsec:arithmetic-universal-holography}

Let $X$ be a smooth projective curve over a finite field $\mathbb{F}_q$
of characteristic coprime to the programme parameters. Write
$\mathrm{Cht}^n_G$ for the Drinfeld moduli of iterated $G$-shtukas on
$X$ with $n$ modification points, and $\mathrm{Pic}(\mathrm{Cht}^n_G)$
for its ($\ell$-adic) Picard group.

\begin{theorem}[Arithmetic Universal Holography]
\label{thm:arithmetic-universal-holography}
The Universal Holography functor of Theorem~\ref{thm:universal-holography-stratified}
lifts to an $\ell$-adic Frobenius-equivariant adjunction
\[
  \Phi_{\mathrm{hol}}^{\mathrm{arith}} :
  \mathrm{ChirAlg}^{\omega, \mathrm{BL}, \mathrm{arith}}_X
  \;\rightleftarrows\;
  \mathrm{HT\text{-}QFT}^{\mathrm{arith}}_{X \times \mathbb{R}} :
  \mathrm{Reduce}^{\mathrm{arith}}.
\]
The arithmetic Koszul complementarity reads
\[
  \bigl[L_{\kappa}(V_k(\mathfrak{g}))\bigr]
  \;+\;
  \bigl[L_{\kappa}(V_{-k-2h^{\vee}}(\mathfrak{g}))\bigr]
  \;=\;
  0
  \quad\in\quad
  \mathrm{Pic}(\mathrm{Cht}^n_G)\otimes \mathbb{Q}_{\ell},
\]
where $L_{\kappa}$ is the line bundle on $\mathrm{Cht}^n_G$
classified by the Koszul-complementarity class
$\kappa(V_k)$.
\end{theorem}

\begin{proof}[Sketch]
The Brylinski--Deligne $K_2$-central extension (Publ. IHES 2001
Thm~4.7) is defined over any base, hence over $\mathbb{F}_q$. The
Arkhipov--Gaitsgory twisted Satake equivalence admits an $\ell$-adic
Frobenius-equivariant twin (Gaitsgory--Lurie Annals Studies 2019) on
$\mathrm{Gr}_G$ over $\overline{\mathbb{F}_q}$. The level-involution
$k \mapsto -k-2h^{\vee}$ is integer-linear, hence characteristic-
independent, and commutes with the partial Frobenius on
$\mathrm{Cht}^n_G$ (Varshavsky 2004). The Picard-class identity
reduces to the complex-analytic formula via the $\ell$-adic
specialisation isomorphism.
\end{proof}

\begin{corollary}[Frobenius eigenvalue bound]
\label{cor:frobenius-eigenvalue-bound}
Let $V_k(\mathfrak{g})$ be a chiral algebra over $\mathbb{F}_q$
admitting an admissible level, and let $B_w$ be its bar complex at
weight $w$. Then
\[
  \bigl|\mathrm{Tr}(\mathrm{Frob}_p \mid V_{\ell}(B_w))\bigr|
  \;\le\;
  \dim(B_w) \cdot p^{(w-1)/2},
\]
consistent with Deligne's Weil~II bound.
\end{corollary}

\begin{proof}
Apply Deligne's Weil~II to the $\ell$-adic fusion computation of
$V_{\ell}(B_w)$ via the arithmetic chiral Adams functor.
\end{proof}

\begin{remark}
The spectral R-matrix upgrades from $R(z)$ to
$R(z, \mathrm{Frob}_q)$ via V.~Lafforgue 2018 excursion operators.
The chiral Adams functor $\psi^q$ on the bar complex equals partial
Frobenius restricted to a spectral component; this supplies the
bridge from the programme's factorisation-algebra framework to
the Lafforgue parameterisation.
\end{remark}

\subsection{Birational invariance of \texorpdfstring{$\Phi_3$}{Phi-3} via Maulik--Okounkov motivic-DT dynamical cocycle}
\label{subsec:birational-invariance-phi3}

Let $X, X'$ be smooth Calabi--Yau threefolds related by a flop
(K-equivalent). Write $\Phi_3 : \mathrm{CY}_3 \to E_1\text{-}\mathrm{ChirAlg}_{\infty}$
for the CY-A$_3$ functor.

\begin{theorem}[Birational invariance]
\label{thm:birational-invariance}
Assuming CY-A$_3$ (proved in Vol~III $\infty$-categorically for local
toric / local del-Pezzo CY$_3$, conjectural on compact CY$_3$ beyond
tilting),
\[
  \Phi_3(X) \;\simeq\; \Phi_3(X')
  \quad
  \text{in}
  \quad
  E_1\text{-}\mathrm{ChirAlg}_{\infty}.
\]
The R-matrices are Maulik--Okounkov gauge-equivalent via a
dynamical cocycle
\[
  P_{\omega}(u; \kappa)
  \;\in\;
  \mathrm{Hom}\bigl(V_X \otimes V_X,\; V_{X'} \otimes V_{X'}\bigr)[\![u, \kappa^{\pm 1}]\!]
\]
where $u$ is the spectral parameter and $\kappa$ the Kähler parameter
exchanged by the flop.
\end{theorem}

\begin{theorem}[Motivic-DT triple presentation of $P_{\omega}$]
\label{thm:mo-motivic-dt-triple}
The dynamical cocycle $P_{\omega}(u; \kappa)$ admits three equivalent
canonical presentations:
\begin{enumerate}
\item[(KS)] Kontsevich--Soibelman motivic-DT: $P_{\omega}$ is the
representation of $\Theta_{\mathrm{KS}}(\kappa) \in G_{\mathrm{KS}}(X)$
on the Yangian eval module $V \otimes V$ via the
Schiffmann--Vasserot / Maulik--Okounkov homomorphism $\rho_V$.
\item[(AO)] Aganagic--Okounkov Stokes matrix: $P_{\omega}$ is the Stokes
matrix of the quantum difference equation around the flop wall
(arXiv:1604.00423).
\item[(I)] Iritani $\widehat{\Gamma}$-class twist: $P_{\omega}$ is
the quantum K-theoretic $\widehat{\Gamma}$-twist
(arXiv:0903.1463).
\end{enumerate}
The Cartan-scalar asymptotic $u \to \infty$ reproduces the
Gopakumar--Vafa generating function
$\exp\!\bigl(\partial_{\kappa} F_0^{\mathrm{inst}}(\kappa)\bigr)
= \exp\!\bigl(\sum_d n_d^0\,\kappa^d / d^2\bigr)$. For the conifold:
$\exp(\mathrm{Li}_2(\kappa))$ with $n_1^0 = 1$ (exceptional $\mathbb{P}^1$).
The flop involution $\kappa \leftrightarrow \kappa^{-1}$ is the
$A_2$ pentagon / Faddeev--Volkov identity in $G_{\mathrm{KS}}$.
\end{theorem}

\begin{corollary}[Invariants across flop]
$\kappa_{\mathrm{ch}}, \kappa_{\mathrm{BKM}}, \kappa_{\mathrm{fiber}}$,
and the shadow class are preserved under the flop equivalence of
Theorem~\ref{thm:birational-invariance}.
\end{corollary}

\begin{remark}
The converse fails: Mukai-type derived-equivalent non-birational
examples (e.g., Bridgeland--Maciocia 2001 K3 and its dual lattice)
give $\Phi_3$-equivalent pairs without a flop relation. Birational
invariance is strictly finer than derived equivalence.
\end{remark}

\subsection{Ambient category for Theorem A: a triple atlas}
\label{subsec:ambient-triple-atlas}

\begin{proposition}[Three ambient categories, one programme]
\label{prop:ambient-triple-atlas}
The ambient category for Vol~I Theorem~A is not a single monoidal
structure but a compatible triple $(\mathcal{A}_1, \mathcal{A}_2,
\mathcal{A}_3)$:
\begin{enumerate}
\item $\mathcal{A}_1 = \bigl(D\text{-}\mathrm{mod}(\mathrm{Ran}(X)),\; \otimes^*\bigr)$
(Beilinson--Drinfeld 2004 plus Francis--Gaitsgory 2012): the default
ambient where FORM-A of Theorem~A (Quillen equivalence on
conilpotent-complete coalgebras) and Koszul complementarity close.
\item $\mathcal{A}_2 = \mathrm{IndCoh}(\mathrm{Ran}(X))$
(Gaitsgory--Rozenblyum 2017 IV.5, V): the $(\infty, 2)$-properad
upgrade; mandatory for Virasoro / Yangian / critical Kac--Moody,
where dim-finiteness obstructs $\mathcal{A}_1$; supports the
Pridham--Lurie formal-moduli-problem formalism.
\item $\mathcal{A}_3 = D^b_c(\mathrm{Ran}(X)) + \mathrm{MHM}$:
the $\ell$-adic or Hodge-module ambient for Theorem~D
tensor-Arakelov curvature, modular-bootstrap complex, and weight
filtration.
\end{enumerate}
The BD chiral pseudo-tensor $\otimes^{\mathrm{ch}}$ generates
operations \emph{inside} $\mathrm{Fact}(X)$ and is never itself an
ambient.
\end{proposition}

\begin{proof}[Sketch]
$\mathcal{A}_1$ is symmetric monoidal (BD04 Ch.~3.4 + FG12 Prop~2.3);
$\mathcal{A}_2$ embeds $\mathcal{A}_1$ fully faithfully as the
locally-compact-generated subcategory (GR17 Vol~II Ch.~5);
$\mathcal{A}_3$ is a Betti/Hodge realisation compatible with the
factorisation structure via the Beilinson realisation
$D^b_c(\mathrm{Ran}(X)) \to \mathrm{Ind}\text{-}\mathrm{D\text{-}mod}(\mathrm{Ran}(X))$.
\end{proof}

\begin{remark}[Non-compact base]
For non-compact $X$, Verdier duality in $\mathcal{A}_1$ is replaced
by Borel--Moore duality on the evaluation-generated core; Theorem~A
FORM-A transfers verbatim (consistent with AP47 scope).
\end{remark}

\subsection{Theorem C total shifted-symplectic: Lagrangian section of the modular tangent bundle}
\label{subsec:theoremC-total-shifted-symplectic}

Theorem~C of Volume~I establishes, at each genus $g$ separately,
the Lagrangian complementarity
$\kappa(\mathcal{A}) + \kappa(\mathcal{A}^!) = \rho_K$ and the
orthogonality $Q_g(\mathcal{A}) \oplus Q_g(\mathcal{A}^!) = \rho_K$
of the characteristic functions of a chirally Koszul pair. The
Platonic form promotes this genus-wise identity to a \emph{global}
section statement on the Deligne--Mumford modular stack
$\overline{\mathcal{M}}_{g,n}$, with clutching compatibility
across boundary strata.

\begin{theorem}[Total shifted-symplectic upgrade of Theorem~C]
\label{thm:theoremC-total-shifted-symplectic}
Let $\mathcal{A}$ be a chirally Koszul $E_1$-chiral algebra on a
smooth projective curve and $\mathcal{A}^!$ its Koszul dual. Form
the relative characteristic bundle
\[
  \mathcal{Q}(\mathcal{A}) \;:=\;
  \bigl\{Q_g(\mathcal{A})\bigr\}_{(g,n)} \longrightarrow
  \overline{\mathcal{M}}_{g,n}
\]
associating to a stable marked curve the genus-$g$ characteristic
function. Then:
\begin{enumerate}
\item\label{tsc-i} The total space
  $\mathcal{Q}(\mathcal{A}) \oplus \mathcal{Q}(\mathcal{A}^!)$
  carries a canonical $-(3g{-}3{+}n)$-shifted symplectic structure
  $\omega^{\mathrm{tot}}$ in the PTVV
  sense~\cite{PTVV13}, induced by the Verdier pairing on chiral
  homology integrated against the Mumford class on
  $\overline{\mathcal{M}}_{g,n}$.
\item\label{tsc-ii} The diagonal section
  $\sigma_{\mathcal{A}}\colon
   \overline{\mathcal{M}}_{g,n} \longrightarrow
   \mathcal{Q}(\mathcal{A}) \oplus \mathcal{Q}(\mathcal{A}^!),
   \quad
   (C, p_\bullet) \mapsto
   \bigl(Q_g(\mathcal{A})|_{(C,p_\bullet)},\,
         Q_g(\mathcal{A}^!)|_{(C,p_\bullet)}\bigr)$
  is a Lagrangian section of $\omega^{\mathrm{tot}}$: its image is a
  Lagrangian sub-stack at every stable point.
\item\label{tsc-iii} Across a boundary divisor
  $\overline{\mathcal{M}}_{g-1,n+2} \hookrightarrow
   \overline{\mathcal{M}}_{g,n}$ (nonseparating node) or
  $\overline{\mathcal{M}}_{g_1,n_1+1} \times
   \overline{\mathcal{M}}_{g_2,n_2+1}
   \hookrightarrow \overline{\mathcal{M}}_{g,n}$ (separating node),
  the section $\sigma_{\mathcal{A}}$ restricts compatibly with the
  clutching morphism: the stratum-wise Lagrangian complementarity
  of Theorem~C glues to a single Lagrangian section on the whole
  modular stack.
\end{enumerate}
\end{theorem}

\begin{proof}[Sketch]
\textit{Step 1 (shifted symplectic on the total space).} The
Verdier pairing
$\langle \alpha, \beta \rangle = \int_{\overline{\mathcal{M}}_{g,n}}
  \alpha \wedge \beta \cdot c_1(\Lambda_g)^{\mathrm{Mumford}}$
(where $\Lambda_g$ is the Hodge bundle of rank $g$ on
$\overline{\mathcal{M}}_g$) pulled back to the total space
$\mathcal{Q}(\mathcal{A}) \oplus \mathcal{Q}(\mathcal{A}^!)$ gives
a non-degenerate $2$-form shifted by the virtual dimension
$-\dim(\overline{\mathcal{M}}_{g,n}) = -(3g{-}3{+}n)$. PTVV
\cite[Thm~2.5]{PTVV13} asserts that such Verdier-integrated
pairings are automatically shifted symplectic when the underlying
pairing on fibres is non-degenerate; fibre non-degeneracy is the
Theorem~C orthogonality at each $(g, n)$.

\textit{Step 2 (Lagrangian section via genus-wise complementarity).}
At each stable point $(C, p_\bullet) \in
\overline{\mathcal{M}}_{g,n}$, the fibre of
$\omega^{\mathrm{tot}}$ is the restriction
$\omega_{(C,p_\bullet)} = \omega_g(Q_g(\mathcal{A}),
Q_g(\mathcal{A}^!))$, and $\sigma_{\mathcal{A}}$ maps to the
Lagrangian
$\{(Q_g(\mathcal{A})|_{(C,p_\bullet)},
   Q_g(\mathcal{A}^!)|_{(C,p_\bullet)})
 : \kappa(\mathcal{A}) + \kappa(\mathcal{A}^!) = \rho_K\}$
witnessed by the Theorem~C complementarity relation.

\textit{Step 3 (clutching compatibility).} The separating-clutching
morphism $\xi_{\mathrm{sep}}\colon
\overline{\mathcal{M}}_{g_1,n_1+1} \times
\overline{\mathcal{M}}_{g_2,n_2+1} \to
\overline{\mathcal{M}}_{g,n}$ pulls back $\omega^{\mathrm{tot}}$ to
the external tensor $\omega^{\mathrm{tot}}_{(g_1,n_1+1)} \boxtimes
\omega^{\mathrm{tot}}_{(g_2,n_2+1)}$ via the Kunneth genus
factorisation
$Q_g(\mathcal{A}) = Q_{g_1}(\mathcal{A}) \cdot
Q_{g_2}(\mathcal{A})$ plus a zero-mode contraction at the node.
The zero-mode contraction is a Lagrangian correspondence (Getzler--Kapranov
modular operad clutching) so $\sigma_{\mathcal{A}}$ glues. The
nonseparating case is analogous with the factorisation replaced by
$Q_g(\mathcal{A}) = Q_{g-1}(\mathcal{A}) \cdot
\operatorname{tr}_{\mathrm{zero-mode}}$.
\end{proof}

\begin{remark}[Scope and upstream inputs]
\label{rem:theoremC-total-scope}
Theorem~\ref{thm:theoremC-total-shifted-symplectic} is proved
unconditionally at the \textsc{uniform-weight} stratum (i.e.\ on
the genus-$1$ locus, all boundary divisors whose contraction
preserves uniform weight). At \textsc{all-weight} with generic
cross-channel weight, the Lagrangian section remains well-defined
on the scalar (uniform-weight) component; the multi-weight
refinement requires the tensor-Arakelov curvature upgrade of
Theorem~D (Section~\ref{subsec:theoremD-tensor-arakelov}) and the
cross-channel $\delta F_g^{\mathrm{cross}}$ vanishing proven in
the curved-Dunn chapter
(Theorem~\ref{thm:curved-dunn-H2-vanishing-all-genera}). The
PTVV-theoretic content is purely formal, but the input data
$(\omega_g, Q_g(\mathcal{A}))$ is Beilinson-rectified by Vol~I
Theorem~D's Arakelov class $\kappa_{\mathrm{Ar}}$ (equivalently
the sixth slot $\kappa_{\mathrm{ch}}$ of the Infinite Fingerprint,
Theorem~\ref{thm:fingerprint-complete-ch}).
\end{remark}

\begin{example}[Heisenberg: Lagrangian section as Mumford class]
\label{ex:theoremC-heisenberg}
For $\mathcal{A} = \mathcal{H}_k$ Heisenberg at level $k$,
$Q_g(\mathcal{H}_k) = k^g$ (a scalar function on
$\overline{\mathcal{M}}_g$); $\mathcal{H}_k^! = \mathcal{H}_{-k}$
Koszul dual; the Lagrangian section is
$\sigma_{\mathcal{H}_k}(C) = (k^g, (-k)^g)$ with
$\kappa(\mathcal{H}_k) + \kappa(\mathcal{H}_{-k}) = k + (-k) = 0$.
At genus $g$, $\omega^{\mathrm{tot}}_g = k^g \wedge (-k)^g \cdot
c_1(\Lambda_g) = (-1)^g k^{2g} c_1(\Lambda_g)$, and the section is
Lagrangian because $dk \wedge d(-k) = 0$ trivially on the
$(k,-k)$-diagonal. The partition function
$Z_g(\mathcal{H}_k) = k^g$ is recovered as the Chern class
integration $\int_{\overline{\mathcal{M}}_g} \sigma_{\mathcal{H}_k}
= k^g \cdot \chi(\overline{\mathcal{M}}_g)$ (up to normalisation),
consistent with Brown--Henneaux for the abelian case.
\end{example}

\paragraph{Adversarial attack and heal.} The upgrade target ``total
$-(3g{-}3)$-shifted symplectic on $\overline{\mathcal{M}}_{g,n}$''
in the UPGRADE-SWEEP ledger is ambiguous between three distinct
objects: (A) the modular stack $\overline{\mathcal{M}}_{g,n}$
itself, (B) the derived moduli of stable maps
$\overline{\mathcal{M}}_{g,n}(X, \beta)$, and (C) the total bundle
$\mathcal{Q}(\mathcal{A}) \oplus \mathcal{Q}(\mathcal{A}^!) \to
\overline{\mathcal{M}}_{g,n}$. The Beilinson-rectified honest form
is (C): the modular stack itself is Deligne--Mumford, not
PTVV-symplectic (\cite[Rem.~2.7]{PTVV13}: shifted symplectic
structures exist only on derived stacks with non-trivial tangent
complex in negative degrees); the stable-map moduli (B) carries
a shifted symplectic structure but depends on an ambient target
$X$; only the characteristic bundle (C) gives a canonical
shifted-symplectic ambient depending on $\mathcal{A}$ alone. The
shift magnitude is $-(3g{-}3{+}n)$, not $-(3g{-}3)$: the point
markings contribute additively. This is the strongest honest form
of the Theorem~C upgrade.

\subsection{Super and logarithmic bar-cobar as M\textsubscript{Kosz} charts}
\label{subsec:super-log-mkosz-charts}

\begin{theorem}[Super-VOA bar-cobar]
\label{thm:super-voa-bar-cobar}
Theorem~A extends to a simple Lie superalgebra $\mathfrak{g}$ with
$h^{\vee}(\mathfrak{g}) \ne 0$ via the Koszul-sign twist
$\mathrm{Sym}_n \to \mathrm{Sym}\text{-super}_n$
on the Ran space, replacing Chevalley involution with the
Chevalley--Kac involution $\theta$. The super-Yangian Koszul dual is
\[
  \bigl(\mathcal{Y}_{\hbar}^{\mathrm{super}}(\mathfrak{g})\bigr)^{!}
  \;=\;
  \mathcal{Y}_{-\hbar}^{\theta, \mathrm{super}}(\mathfrak{g}).
\]
For the $N = 2$ superconformal algebra, the Koszul complementarity reads
\[
  \kappa(c) + \kappa(6 - c) = 1
\]
under the Kazama--Suzuki coset involution $c \leftrightarrow 6 - c$
(NOT the superdimension-based involution).
\end{theorem}

\begin{proof}[Sketch]
The Koszul-sign twist on $\mathrm{Sym}_n$-action is induced by the
super-parity grading of the coalgebra structure; the replacement of
Chevalley by $\theta$ accounts for the super-commutator sign. The
$N = 2$ formula comes from the Kazama--Suzuki coset
$N=2_{c} \simeq N=2_{6-c}$ (Kazama--Suzuki 1989), which exchanges
central charges symmetrically about $c = 3$.
\end{proof}

\begin{theorem}[Logarithmic $\mathcal{W}(p)$ bar-cobar]
\label{thm:log-wp-bar-cobar}
For the triplet $\mathcal{W}(p)$ logarithmic vertex operator algebra with
$p \ge 2$, Theorem~A extends via the Positselski periodic-CDG
coalgebra replacement of the conilpotent dg-coalgebra ambient.
Non-trivial $\mathrm{Ext}^1$ between distinct simples is captured by
the Jordan-block structure of $L_0$; the Koszul involution is the
Creutzig--Kanade--Linshaw log-Koszul involution (replacing the bosonic
$k \mapsto -k-2h^{\vee}$).
\end{theorem}

\begin{proof}[Sketch]
The ambient $\mathcal{A}_1$ of
Proposition~\ref{prop:ambient-triple-atlas} is enlarged to
$\mathcal{A}_1^{\mathrm{log}} = \bigl(D\text{-}\mathrm{mod}^{\mathrm{log}}(\mathrm{Ran}(X)),
\otimes^*\bigr)$, with $\mathrm{mod}^{\mathrm{log}}$ admitting logarithmic
Jordan blocks. Chain-level closure follows from the three-channel
tempered decomposition $S_r(\mathcal{W}(p)) = S_r^{TT} + S_r^{TW} + S_r^{WW}$,
each channel tempered independently.
\end{proof}

\begin{corollary}
Both the super and logarithmic ambients embed into the Koszulness
Moduli scheme $M_{\mathrm{Kosz}}$ as distinct charts cut out by
Drinfeld-associator coordinates corresponding to Kazama--Suzuki and
Creutzig--Kanade--Linshaw involutions respectively.
\end{corollary}

\begin{remark}[Residual opens]
Free strong generation of $\mathcal{W}(p)$ for $p \ge 4$, and
$\mathrm{sdim} = 0$ super types ($\mathfrak{psl}(n|n)$,
$\mathfrak{osp}(2m|2m)$, $D(2,1;\alpha)$) where
$h^{\vee} = 0$ and Sugawara fails: both require separate
BRST / coset reformulations; not covered by the two theorems above.
\end{remark}

\subsection{Periodic curved-DG at admissible level; exotic-nilpotent DS-Hochschild}
\label{subsec:periodic-cdg-universal}

Fix a simply-connected simple algebraic group $G$ over $\mathbb{C}$
with Lie algebra $\mathfrak{g}$, dual Coxeter $h^{\vee}$, and lacing
number $r_{\mathfrak{g}} \in \{1, 2, 3\}$.

\begin{theorem}[Periodic curved-DG at admissible level]
\label{thm:periodic-cdg-admissible}
For admissible level $k = p/q \in \mathbb{Q}$ with $q \ge h^{\vee}$,
$\gcd(p, q) = 1$, and non-degenerate admissible locus $X_{L_k} = \{0\}$,
the Koszul bar-cobar intertwiner on the Kazhdan--Lusztig category
$\mathrm{KL}_{\mathfrak{g}, k}$ is Serre $t$-exact after tensor with
the metaplectic line bundle $\mathcal{L}_{\mathrm{met}}$
determined by the Brylinski--Deligne $K_2$-central extension. The
periodic curved-DG conjecture on
$\mathrm{KL}_{\mathfrak{g}, k}$ follows.
\end{theorem}

\begin{proof}
For simply-laced $\mathfrak{g}$, Brylinski--Deligne Publ. IHES 94
(2001) Thm~4.7 produces a single $W_{\mathrm{aff}}$-invariant
Kubota class $\alpha_{\mathbb{S},u}
\in H^2(W_{\mathrm{aff}}(\mathfrak{g}), k)$, identified via the
$K_2$-central extension with the Loday symbol on
$\mathrm{Pic}(\mathrm{GR}_G) = \mathbb{Z}$. The Arkhipov--Gaitsgory
equivalence $\mathcal{O}_{\kappa}(\widehat{\mathfrak{g}})
\simeq D^{\kappa}\text{-}\mathrm{mod}(\mathrm{GR}_G)$ transports
this to a $\kappa = (k + h^{\vee})$-line-bundle twist on
$\mathrm{Pic}(\mathrm{GR}_G)$; Serre $t$-exactness of the Koszul
intertwiner holds precisely on the admissible non-degenerate lane
(Bezrukavnikov--Finkelberg--Mirkovi\'c 2005 Thm~1.1).
For non-simply-laced $\mathfrak{g} \in \{B_r, C_r, F_4, G_2\}$, the
Frenkel--Hernandez lacing automorphism
$\Phi_{\mathfrak{g}}(T^{(n)}_{\alpha}) = (-r_{\mathfrak{g}})^n
T^{(n)}_{\alpha^*}$
intertwines the simply-laced Yangian structure with
$\widehat{\mathfrak{g}}$-structure and scales $\kappa$ by
$r_{\mathfrak{g}}$. The $t$-exact locus is preserved on the
non-degenerate admissible lane.
\end{proof}

\begin{remark}[Residuals]
Three lanes lie outside the theorem's hypothesis: (a) degenerate
admissible $X_{L_k} \ne \{0\}$ (Creutzig--Kanade--Linshaw logarithmic
extension); (b) non-admissible generic $k$ (Part~VII Massey structure);
(c) critical level $k = -h^{\vee}$, handled by the fifth-class FF
companion adjunction in
Theorem~\ref{thm:universal-holography-stratified}.
\end{remark}

\begin{theorem}[Exotic-nilpotent DS-Hochschild chain-level]
\label{thm:ds-hoch-exotic}
For every good or exotic nilpotent $f \in \mathfrak{g}$ with
Kazhdan-grading denominator $d_f \in \{1, 2, 3, 6\}$ and every
non-critical level $k$, there is a chain-level quasi-isomorphism
\[
  \mathrm{ChirHoch}^{\bullet}\bigl(\mathcal{W}_k(\mathfrak{g}, f)\bigr)
  \;\simeq\;
  H^{\bullet}_{\mathrm{DS}}\bigl(\mathrm{ChirHoch}^{\bullet}(V_k(\mathfrak{g}))\bigr).
\]
\end{theorem}

\begin{proof}
Three-step Galois descent. First, pass to the degree-$d_f$ branched
cover $\pi : \widetilde{X} \to X$ on which the fractional Kazhdan
weights of $f$ become integral. Second, apply the integer-weight
DS-Hochschild bridge $\mathrm{thm}\text{:}\mathrm{chd\text{-}ds\text{-}hochschild}$
on $\widetilde{X}$. Third, Galois-descend via Crainic's 2004
equivariant homotopy-perturbation lemma (arXiv:math/0403266), with
averaging
$P = (1/d_f) \sum_{g \in \mathbb{Z}/d_f} \sigma_g^*$
commuting with $Q_{\mathrm{BRST}}$ by deck-covariance, producing a
chain-level (not merely cohomological) SDR on $X$.
\end{proof}

\begin{corollary}[Universal Holography full domain]
\label{cor:universal-holography-full-domain}
The functor $\Phi_{\mathrm{hol}}$ of
Theorem~\ref{thm:universal-holography-stratified} is defined on the
entire category $\mathrm{ChirAlg}^{\omega, \mathrm{BL}}_X$ of
boundary-linear chiral algebras with conformal vector at non-critical
level, with no scope restriction to principal / hook-type / good-only
nilpotents.
\end{corollary}

\subsection{Clay-level problems: precise reformulations inside the programme}
\label{subsec:clay-precise-reformulations}

Each of the seven Clay-adjacent problems attacked this session receives
an explicit reformulation in chiral-algebra language. The programme
does not solve any; it pins the obstruction precisely.

\begin{proposition}[Riemann Hypothesis reformulated]
\label{prop:rh-reformulated}
Define the Sewing--Selberg L-function of Virasoro by
\[
  L_{\mathcal{A}}(s, u) \;:=\; \Gamma(v)\,(2\pi)^{-v}\,\zeta(v + 1)\bigl(\zeta(v) - 1\bigr),
  \qquad v := s + u - 1/2.
\]
Zeros of $L_{\mathcal{A}}(s, u)$ along the diagonal $s + u = 1/2$
coincide with non-trivial zeros of $\zeta$. Motivic-shadow,
Weil--Grothendieck $\ell$-adic, and random-matrix routes are
architecturally closed against the Riemann Hypothesis; RH remains
external.
\end{proposition}

\begin{proposition}[BSD reformulated]
\label{prop:bsd-reformulated}
For an elliptic curve $E / \mathbb{Q}$ without CM, let
$\Phi(E) = V_{L_E}$ be the rank-2 Galois-equivariant lattice VOA on
$H^1(E, \mathbb{Z})$. Then the programme's bar L-function coincides
with the classical: $L^{\mathrm{bar}}(\Phi(E), s) = L(E, s)$.
Kolyvagin's Euler system lives in $H^1(K, E[\ell^k])$, orthogonal
to the shadow tower of $\Phi(E)$ (which is trivial, $S_r = 0$ for
$r \ge 3$). The Mordell--Weil rank is invisible to $\Phi$'s
topological $H^1$ input.
\end{proposition}

\begin{proposition}[Sato--Tate reformulated]
\label{prop:sato-tate-reformulated}
The shadow tower of $\Phi(E)$ generates the
$\mathrm{Sym}^{r-1}$-Dirichlet coefficients of $L(E, \mathrm{Sym}^r)$
via the symmetric-power proposition in
\texttt{arithmetic\_shadows.tex}:7637. Sato--Tate reduces to the
automorphy of $\mathrm{Sym}^r(E)$, established by Newton--Thorne 2021.
The programme packages but does not prove the analytic step.
\end{proposition}

\begin{proposition}[Hodge $CY_4$ reformulated]
\label{prop:hodge-reformulated}
For a smooth projective Calabi-Yau fourfold $X$, the chiral Hodge
decomposition
$\mathrm{HH}^{\bullet}(A_X) = \bigoplus_{p+q = k} \mathrm{HH}^{p,q}(A_X)$
via Kaledin noncommutative Hodge-to-de-Rham degeneration plus HKR
is well-defined. The chiral cycle-class map $\mathrm{Ch}^{\mathrm{ch}}$
is tautologically defined from the classical cycle-class map plus HKR.
The image of algebraic cycles in $\mathrm{HH}^{p,p}$ contains the
classical algebraic subspace; Voisin's transcendental obstructions
to surjectivity apply unchanged.
\end{proposition}

\begin{proposition}[YM mass gap reformulated]
\label{prop:ym-reformulated}
The HT-twisted 4d $\mathcal{N} = 2$ Schur sector of pure
super-Yang--Mills with gauge group $G$ has boundary chiral algebra
$V_{-2h^{\vee}}(\mathfrak{g})$ (class L); lowest non-vacuum conformal
weight $\Delta_{\min} = 1$ (R-current) for $G = SU(N)$. Universal
Holography supplies an algebraic mass-gap reduction via
(H1)--(H2) (algebraic screening), contingent on
\textbf{(H3.b)}: positivity of the Costello--Gwilliam RG-flow fixed
point on pure $\mathcal{N} = 0$ Yang--Mills. Neither $\Phi_{\mathrm{hol}}$
nor its arithmetic refinement produces (H3.b); the Clay mass gap
proper is reducible to constructing a non-supersymmetric HT twist.
\end{proposition}

\begin{proposition}[Grothendieck Period Conjecture]
\label{prop:gpc-reformulated}
The programme's motivic MZV ring
$Z_{\bullet}^{\Phi_{\mathrm{KZ}}}$ coincides with Brown's motivic ring
$Z_{\bullet}^{\mathrm{mot}}$ as graded $\mathbb{Q}$-algebras. The
Kontsevich--Zagier three operations (additivity, change-of-variables,
Stokes) correspond respectively to Francis--Gaitsgory factorisation
on $\mathrm{Ran}(X)$, $\mathrm{Aut}(\mathcal{O})$-equivariance of
$\mathrm{End}^{\mathrm{ch}}_A$, and the bar differential. Furusho's
pentagon relation plus Ihara--Kaneko derivation relations close the
Kontsevich--Zagier conjecture on motivic periods. The Grothendieck
Period Conjecture (injectivity of
$\mathrm{Per} : Z_{\bullet}^{\mathrm{mot}} \to \mathbb{C}$) requires
external transcendence input (Apéry 1978, Ball--Rivoal 2001) outside
the programme.
\end{proposition}

\begin{remark}
Each proposition names the precise obstruction from within the
programme. None reduces a Clay problem to an easier sub-problem; the
content is the architectural translation. Future arithmetic progress
(function-field Langlands Universal Holography, Frobenius-equivariant
Koszul complementarity) is the programme's closest approach to
arithmetic Clay-adjacent domains.
\end{remark}

\subsection{Two Langlands-invariants: Lie-algebra vs root-datum}
\label{subsec:langlands-two-invariants}

Fix a simple algebraic group $G$ over $\mathbb{C}$ with Lie algebra
$\mathfrak{g}$ and root datum $\Psi_G = (X, R, X^{\vee}, R^{\vee})$.
Langlands duality exchanges root data
$\Psi_G \leftrightarrow \Psi_{G^L}$.

\begin{theorem}[Langlands-invariant / equivariant split]
\label{thm:langlands-invariants-split}
The data attached to $\mathfrak{g}$ by Vol~I and Vol~II splits:
\begin{enumerate}
\item[(I)] \textbf{Langlands-invariant}: $\dim(\mathfrak{g})$,
exponents $\{m_i\}$, Coxeter number $h$,
$\kappa_V(\mathfrak{g}) = 2\dim$,
$\kappa_W(\mathfrak{g}) = 12\,p_2 - 2r(3h+5)$.
\item[(II)] \textbf{Langlands-equivariant}: coroot lattice $X^{\vee}$,
dual Coxeter $h^{\vee}$, opers $\mathrm{Op}_{\mathfrak{g}}$, Hitchin
system on $\mathrm{Bun}_G(X)$.
\end{enumerate}
\end{theorem}

\begin{proof}
The quantities in (I) are invariants of the underlying Lie algebra;
they agree on every pair $(G, G^L)$ sharing an algebra. Coroot data
and dual Coxeter detect the short/long root interchange:
$h^{\vee}(B_r) = 2r-1 \ne r+1 = h^{\vee}(C_r)$. Beilinson--Drinfeld's
$\mathfrak{z}(\widehat{\mathfrak{g}}) = \mathrm{Fun}(\mathrm{Op}_{\mathfrak{g}^L})$
realises (II) explicitly.
\end{proof}

\begin{example}[$B_2 / C_2$ Dynkin accident]
\label{ex:b2-c2-accident}
$\mathfrak{so}_5 \cong \mathfrak{sp}_4$ share $\dim = 10$, $h = 4$,
exponents $\{1, 3\}$, so (I) gives $\kappa_W(B_2) = \kappa_W(C_2)
= 172$ and $\kappa_V(B_2) = \kappa_V(C_2) = 20$. The Hitchin systems
differ: the Beilinson--Drinfeld quantum geometric Langlands
equivalence
$D\text{-}\mathrm{mod}(\mathrm{Bun}_{SO(5)}) \simeq
\mathrm{QCoh}(\mathrm{Op}_{Sp(4)})$
is a non-trivial equivalence over (II), not the identity.
\end{example}

\begin{corollary}
Vol~I theorems (A, B, C, D, H) depend on (I); Vol~II Theorem~F at
the critical stratum depends on (II). The Frenkel--Hernandez lacing
automorphism $\Phi_{\mathfrak{g}}(T^{(n)}_{\alpha}) =
(-r_{\mathfrak{g}})^n T^{(n)}_{\alpha^*}$ realises (II) at the
Yangian level with lacing number
$r_{\mathfrak{g}} \in \{1, 2, 3\}$.
\end{corollary}

\subsection{CY-C six-route convergence at character level}
\label{subsec:cy-c-character-level}

Let $G(K3 \times E)$ denote the conjectural Borcherds--Kac--Moody
superalgebra attached to the Mukai lattice of $K3 \times E$. Six
constructions of $G(K3 \times E)$ (CY-A functor applied to
$D^b(\mathrm{Coh}(K3 \times E))$; Harvey--Moore BPS state algebra;
Kummer orbifold; sigma-model Kac--Moody; Beem--Lemos--Liendo--Peelaers--Rastelli--van Rees
Schur sector; Gritsenko--Nikulin arithmetic lift $\Phi_{10}$) are
conjectured to converge (CY-C).

\begin{theorem}[CY-C character-level unconditional]
\label{thm:cy-c-character-level}
At the level of graded characters, the six routes to $G(K3 \times E)$
coincide on the Koszul locus at generic admissible level, subject to
the following three conditional-to-unconditional identity ladder:
\begin{enumerate}
\item[$I_1$] Harvey--Moore elliptic-genus subgroup identification:
conditional on the agreement of the elliptic genus subgroup of
$\mathrm{Aut}\,\mathrm{Eq}(D^b(K3 \times E))$ with the Siegel modular
group $\mathrm{Sp}_4(\mathbb{Z})$.
\item[$I_{2}^{\mathrm{sep}}$] Separating $g = 2$ Borcherds lift:
unconditional via the Gritsenko--Nikulin arithmetic lift of
$\Phi_{10}$, matching Dijkgraaf--Moore--Verlinde--Verlinde
second quantisation through the factor $(2\pi i\,\tau_e)^2$ at the
separating cusp.
\item[$I_{2}^{\mathrm{nonsep}}$] Non-separating $g = 2$: unconditional
via Fourier--Jacobi expansion of $\Phi_{10}$ at the non-separating
cusp:
$\Phi_{10}(\Omega) = (2\pi i \tau_{12})^2 \prod_{n \ge 1}
(1 - q_1^n)^{c(n)} + \mathcal{O}(\tau_{12}^4)$
with $c(n)$ the Fourier coefficients of the Gritsenko--Nikulin lift.
\item[$I_3$] BLLPR Schur-index identification: unconditional at
character level via the BLLPR trace; chain-level closure reduces to
the $N = 2$ SCVOA coset DS-Hochschild bridge.
\end{enumerate}
\end{theorem}

\begin{theorem}[Chain-level closure on Koszul locus]
\label{thm:cy-c-chain-level-koszul}
On the Koszul locus at generic admissible level, the chain-level
bridge $I_3^{\mathrm{chain}}$ is proved by composing:
\begin{itemize}
\item Arakawa--Kawasetsu--M\"oller coset-$C_2$-cofiniteness of the
BLLPR chiral algebra;
\item two-step HPL transfer through the hook-type $\mathcal{W}$-algebra
and Heisenberg coset (extensions of
$\mathrm{thm\text{:}chd\text{-}ds\text{-}hochschild}$ and its
exotic-nilpotent extension
$\mathrm{thm\text{:}ds\text{-}hoch\text{-}exotic}$).
\end{itemize}
The residual chain-level gap reduces to exotic nilpotents inside
coset Argyres--Douglas theories.
\end{theorem}

\begin{corollary}
The six-route CY-C convergence is character-level unconditional on
the Koszul locus. The single remaining chain-level frontier is the
coset-AD exotic-nilpotent extension of
Theorem~\ref{thm:ds-hoch-exotic}.
\end{corollary}

\subsection{K\textsubscript{W} from Chevalley--Shephard--Todd: direct derivation}
\label{subsec:kw-cst-direct}

\begin{theorem}[Direct CST derivation of $\kappa_W$]
\label{thm:kw-cst-direct}
The formula $\kappa_W(\mathfrak{g}) = 2\,\mathrm{rank}(\mathfrak{g}) +
12 \sum_{i=1}^{\mathrm{rank}(\mathfrak{g})} (d_i - 1)\,d_i$
derives from classical finite-reflection-group invariant theory,
independent of the Feigin--Frenkel screening computation.
\end{theorem}

\begin{proof}
Let $W = W(\mathfrak{g})$ act on $V \cong \mathbb{C}^r$ with fundamental
invariants of degrees $d_1, \dots, d_r$; by Chevalley 1955 plus
Shephard--Todd 1954, $S(V^*)^W$ is polynomial in these fundamentals.
For $W = W(\mathfrak{g})$ the degrees equal the Casimir degrees
$d_i = m_i + 1$.

Solomon 1963 Theorem~1:
\[
  \prod_{i=1}^r \bigl(1 + (d_i - 1)\,t\bigr)
  \;=\;
  \sum_{w \in W}\,t^{\ell(w)}\,\mathrm{sgn}(w)
\]
equates generating functions of $W$-invariants and $W$-anti-invariants.
Expanding the LHS and reading off coefficients:
$[t^1]\,\mathrm{LHS} = \sum_i (d_i - 1) = |\Delta_+|$, recovering
$\sum_i m_i = |\Delta_+|$.

Shapiro--Steinberg--Kostant:
$\sum_{\alpha \in \Delta_+}\mathrm{ht}(\alpha) = (1/2)\sum_i d_i(d_i-1)$.

The identity $(\rho, \alpha_i^{\vee}) = 1$ for simple roots
$\alpha_i$ gives $2(\rho, \rho^{\vee}) = \sum_{\alpha > 0}(\rho, \alpha^{\vee})
= \sum_{\alpha > 0}\mathrm{ht}_{\mathrm{coroot}}(\alpha)$, and via
Kostant's coroot-height-vs-root-height interchange
$\sum_{\alpha > 0}\mathrm{ht}_{\mathrm{coroot}}(\alpha) = \sum_{\alpha > 0}\mathrm{ht}(\alpha)$,
so
\[
  (\rho, \rho^{\vee})
  \;=\;
  \frac{1}{4}\sum_i d_i\,(d_i - 1).
\]

Combining with Theorem~\ref{thm:kw-chevalley-shephard-todd} gives
$\kappa_W = 2\,r + 48(\rho, \rho^{\vee}) = 2\,r + 12\sum_i (d_i-1)\,d_i$
via Weyl-group invariant theory alone; no affine-level parameter
enters.
\end{proof}

\begin{remark}
Theorem~\ref{thm:kw-cst-direct} exhibits $\kappa_W$ as a pure
finite-reflection-group invariant of $W(\mathfrak{g})$, computed from
Casimir degrees by classical invariant theory. The Feigin--Frenkel
derivation and the direct CST derivation produce identical values at
every simple $\mathfrak{g}$; their agreement is a compatibility between
the BRST screening-charge computation and Weyl-group invariant
computation, not a new theorem.
\end{remark>

\subsection{Bershadsky--Polyakov Universal Holography and
K\textsubscript{W}(BP) via non-principal W}
\label{subsec:session-bp-universal-holography}

The Bershadsky--Polyakov algebra $\mathcal{W}(\mathfrak{sl}_3,
f_{\min})$ at non-principal minimal nilpotent is a non-principal
W-algebra with three strong generators $(J, G^{\pm}, T)$ at
weights $(1, 3/2, 2)$.
The Universal Holography functor $\Phi_{\mathrm{hol}}$ closes
unconditionally at BP via degree-$2$ Kummer cover of the Kazhdan
$(1/2)$-grading, class-L antighost absorption upstairs (Vol~I
$\eta_1^{(i)} + \eta_1^{(ii)}$ strict chain-level cancellation),
and Maschke-exact $\mathbb{Z}/2$-descent. DS-Hochschild bridge:
$\mathrm{ChirHoch}^{\bullet}(\mathrm{BP}) \simeq
H^{\bullet}_{\mathrm{DS, min}}(\mathrm{ChirHoch}^{\bullet}(V_k(\mathfrak{sl}_3)))$
at chain level.

\begin{theorem}[BP Koszul conductor]
\label{thm:bp-koszul-conductor}
$\kappa_W(\mathrm{BP}) = 196$ via the Fateev--Kac--Roan central charge
$c(\mathrm{BP}_k) = 2 - 24(k+1)^2/(k+3)$ and Koszul-dual level
$k^! = -k - 6$:
\[
  \kappa_W(\mathrm{BP})
  \;=\;
  c(\mathrm{BP}_k) + c(\mathrm{BP}_{-k-6})
  \;=\;
  4 - 24\,\frac{(k+1)^2 + (k+5)^2}{(k+3)(-k-3)}\cdot(-1)
  \;=\;
  196.
\]
\end{theorem}

\begin{remark}[$\kappa_W$ depends on $(\mathfrak{g}, f)$, not $\mathfrak{g}$ alone]
\label{rem:kw-bala-carter-dependence}
Theorem~\ref{thm:kw-chevalley-shephard-todd} gives $\kappa_W$ for the
\emph{principal} $\mathcal{W}$-algebra $\mathcal{W}(\mathfrak{g},
f_{\mathrm{prin}})$. At $\mathfrak{g} = \mathfrak{sl}_3$ with
$f_{\mathrm{prin}}$: $\kappa_W(\mathcal{W}(\mathfrak{sl}_3,
f_{\mathrm{prin}})) = 100$. BP $= \mathcal{W}(\mathfrak{sl}_3,
f_{\mathrm{min}})$ has a different Bala--Carter type, hence
$\kappa_W(\mathrm{BP}) = 196 \ne 100$. The conductor depends on the
Bala--Carter pair $(\mathfrak{g}, f)$, not on $\mathfrak{g}$ alone.
The exponent formula $\kappa_W = 12\,p_2 - 2r(3h+5)$ is specific to
$f = f_{\mathrm{prin}}$.
\end{remark}

See \texttt{notes/bershadsky\_polyakov\_universal\_holography\_attack\_heal.md}
for the $\kappa$-sum rescaled-level computation and Bala--Carter
bookkeeping.

BP's class assignment: class L (per Vol~I classification of
chiral landscape by shadow-depth). Fingerprint
$\varphi(\mathrm{BP}) = (3, 3, +, 3, \mathrm{coset\text{-}BP},
\kappa_{\mathrm{ch}}(\mathrm{BP}))$ with $\kappa_{\mathrm{ch}}
= 1/2$ at the Koszul self-dual point. Six-slot fingerprint
distinguishes BP from other class L algebras via the
$\kappa_{\mathrm{ch}}$ sixth slot.

\subsection{Higher-genus opers: dimension, chain-level model, Theorem D curvature}
\label{subsec:higher-genus-opers}

The Beilinson--Drinfeld chiral algebra of opers on a formal disk
$\mathrm{Op}_{\mathfrak{g}^L}(D^{\times})$ admits higher-genus
integration to a commutative factorisation algebra on
$\mathrm{Ran}(X_g)$, per BD \emph{Chiral Algebras} Ch.~3.7. The
resulting object $\mathrm{Op}_{\mathfrak{g}^L}(X_g)$ is a finite-
dimensional vector space at each $g \ge 2$, not a chiral algebra.

\begin{proposition}[Genus-$g$ oper dimension]
\label{prop:higher-genus-oper-dim}
For $X_g$ smooth projective of genus $g \ge 2$ and simple
$\mathfrak{g}^L$,
\[
  \dim\,\mathrm{Op}_{\mathfrak{g}^L}(X_g)
  \;=\;
  (g - 1)\bigl(\dim\mathfrak{g}^L
    + \mathrm{rank}\,\mathfrak{g}^L\bigr).
\]
In particular: $\mathrm{Op}_{\mathfrak{sl}_2}(X_2) = 3$,
$\mathrm{Op}_{\mathfrak{sl}_2}(X_3) = 6$,
$\mathrm{Op}_{\mathfrak{sl}_3}(X_2) = 8$,
$\mathrm{Op}_{\mathfrak{sl}_3}(X_3) = 16$.
\end{proposition}

\begin{proof}
$\mathrm{Op}_{\mathfrak{g}^L}(X_g)$ is the affine space of Borel
reductions of $G^L$-bundles on $X_g$ with a prescribed nilpotent
section. Its dimension is $(g-1)\dim\mathfrak{g}^L$ (Borel) plus
$(g-1)\mathrm{rank}\,\mathfrak{g}^L$ (nilpotent section degrees of
freedom at genus $g$), by Beilinson--Drinfeld Ch.~3.7 + Hartshorne
Riemann--Roch.
\end{proof}

\begin{theorem}[Chain-level model via Gaitsgory--Rozenblyum]
\label{thm:higher-genus-chain-model}
For $X_g$ as above,
\[
  \mathrm{Op}^{\mathrm{der}}_{\mathfrak{g}^L}(X_g)
  \;\simeq\;
  \mathrm{ChirHoch}^{\bullet}(V_{-h^{\vee}}(\mathfrak{g}))_{X_g}
\]
as factorisation algebras on $\mathrm{Ran}(X_g)$. At degree zero
this recovers $\mathrm{Fun}(\mathrm{Op}_{\mathfrak{g}^L}(X_g))$;
higher cohomology equals the Vol~II Theorem~D curvature tensor
$\kappa(V_{-h^{\vee}}(\mathfrak{g})) \cdot \omega_g$.
\end{theorem}

\begin{proof}[Sketch]
Apply the Gaitsgory--Rozenblyum 2017 IV.5 factorisation framework
to the critical-level affine Kac--Moody chiral Hochschild complex
on $X_g$. The $H^0$ identification is the formal-disk
Beilinson--Drinfeld isomorphism integrated over $\mathrm{Ran}(X_g)$.
Higher cohomology carries the modular anomaly computed by the
Gauss--Manin connection on Theorem~D's Arakelov line bundle.
\end{proof}

\begin{corollary}[Genus-$g$ FF-adjunction]
$\mathrm{FF}\Phi_{\mathrm{hol}}^{(g)}$ lifts to the
Beilinson--Drinfeld Hitchin quantisation on $\mathrm{Bun}_G(X_g)$:
\[
  \kappa(V_{-h^{\vee}}(\mathfrak{g})) \cdot \omega_g
  \;=\;
  c_{\mathrm{BD}}\text{-twisted D-module anomaly of the Hitchin fibration}
\]
(Hausel--Thaddeus 2003). The
Langlands-natural $\mathfrak{g} \leftrightarrow \mathfrak{g}^L$
duality on opers extends continuously to higher-genus
Hitchin-quantised 3d HT-QFTs.
\end{corollary}

\subsection{Universal Borcherds lift and CY-A are independent}
\label{subsec:borcherds-cya-independence}

Write $\overline{B}$ for the universal Borcherds automorphic lift
functor, and $\Phi_3$ for the CY-A$_3$ functor.

\begin{theorem}[Universal Borcherds lift]
\label{thm:universal-borcherds}
For every even unimodular Lorentzian lattice $L$ with $b^+ \ge 2$
and every conformal vertex operator algebra $V$ equipped with an
embedding $L \hookrightarrow \mathrm{Lat}(V)$, the Borcherds lift
$\overline{B}(L, V)$ exists and is a Borcherds--Kac--Moody
superalgebra whose denominator has the Borcherds product expansion
of weight $c_V(0)/2$.
\end{theorem}

\begin{example}[Canonical Borcherds outputs]
\label{ex:borcherds-three}
\begin{enumerate}
\item $(L, V) = (\mathrm{II}_{1,1}, V_L)$: $\overline{B} = j(\tau) - 744$
(Conway).
\item $(L, V) = (\mathrm{II}_{2,26}, V^{\natural})$: $\overline{B} =$
Monster Lie Algebra (Borcherds 1992).
\item $(L, V) = (\mathrm{II}_{4,20}, \mathrm{Leech}^{[2]})$:
$\overline{B} = \Phi_{10}$ Siegel modular form
(Gritsenko--Nikulin).
\end{enumerate}
\end{example}

\begin{proposition}[Independence of $\overline{B}$ and $\Phi_3$]
\label{prop:borcherds-phi3-independence}
$\Phi_3$ consumes $D^b(\mathrm{Coh}(X))$ for $X$ a Calabi--Yau
threefold; the input is not a lattice. On K3-fibred CY$_3$ families,
$\Phi_3$ factors through the Mukai-lattice functor
$X \mapsto \mathrm{H}^*(X_{K3}, \mathbb{Z}) \oplus
\mathrm{H}^*(X_{E}, \mathbb{Z})$, hitting the domain of
$\overline{B}$. $\mathrm{II}_{26,2}$ admits no compact CY$_3$
Mukai-lattice preimage; the Monster Lie Algebra is not in the image
of $\Phi_3$.
\end{proposition}

\begin{corollary}[CY-C intersection]
\label{cor:cy-c-intersection}
CY-C identity $I_2$ (Theorem~\ref{thm:cy-c-character-level}) is the
intersection of the $\overline{B}$-route and the $\Phi_3$-route at
$L = \mathrm{II}_{4,20}$: the character-level agreement of the
Gritsenko--Nikulin $\Phi_{10}$ with the DMVV second-quantisation of
$K3 \times E$ is exactly this intersection.
\end{corollary}

\begin{remark}
The Monster chain-level closure of
Theorem~\ref{thm:schellekens-chain-level}(1) consumes $\overline{B}$
on $V^{\natural}$ directly, not via $\Phi_3$. The AP-CY8 / AP-CY59 /
AP-CY60 distinctions (separating Borcherds denominator from bar
Euler product, six routes from six $\Phi$-applications, one $\Phi$
output per category) are structural consequences of the
independence.
\end{remark>

\subsection{Part~VIII and Vol~IV explicit roadmap}
\label{subsec:session-roadmap}

Based on session closures, the Vol~II Part~VIII (From Frontier to
Theorem) and the planned Vol~IV (Realization) have the following
explicit structure.

\subsubsection*{Part~VIII table of contents (12 items)}

\begin{enumerate}
\item[VIII.1] Curved-Dunn $H^2 = 0$ at all genera
(\texttt{chapters/theory/curved\_dunn\_higher\_genus.tex}): CLOSED
via modular-bootstrap-to-curved-Dunn bridge $\Phi_g$ plus Jimbo--Miwa
irregular-singular KZB.
\item[VIII.2] DS-Hochschild bridge for class M
(\texttt{chapters/theory/chiral\_higher\_deligne.tex}): CLOSED via
HPL transfer + Arakawa $C_2$-cofiniteness + HKR.
\item[VIII.3] Periodic-CDG at admissible level, all simple simply-connected
types: CLOSED via Brylinski--Deligne $K_2$-reciprocity + Arkhipov--Gaitsgory
$\kappa$-twisted Serre t-exactness + Frenkel--Hernandez lacing.
\item[VIII.4] Exotic-nilpotent DS-Hochschild (Crainic $\mathbb{Z}/d_f$
Galois descent, $d_f \in \{1, 2, 3, 6\}$): CLOSED.
\item[VIII.5] N=2 SCVOA coset (BLLPR, AD theories) DS-Hochschild: CLOSED
on the Koszul locus at generic admissible level.
\item[VIII.6] Associator-independent chiral Deligne--Tamarkin via
Higher Massey odd-weight tower: PLATONIC FORM = GRT$_1(\mathbb{Q})$-torsor
of associator-dependent statements. Associator-independent form
genuinely FAILS by Brown motivic non-decomposability.
\item[VIII.7] Universal Holography two-adjunction family-over-level:
CLOSED. Stratified $\infty$-categorical reflection over Level
base $\mathbb{A}^1_k$.
\item[VIII.8] Six-slot fingerprint classification with $\kappa_{\mathrm{ch}}$:
CLOSED. Refined $\varphi'(A) = (p_{\max}, r_{\max}, \chi_{\mathrm{VOA}},
n_{\mathrm{strong}}, \mathrm{coset}, \kappa_{\mathrm{ch}})$ complete
on standard landscape.
\item[VIII.9] Function-field arithmetic Universal Holography: PATH
IDENTIFIED, deferred to Vol~IV. Two-page proof assembling
Brylinski--Deligne 2001 + L.~Lafforgue 2002 + V.~Lafforgue 2018 +
Varshavsky 2004 + Gaitsgory--Lurie 2019.
\item[VIII.10] $\mathcal{W}_{\infty}$ $E_{\infty}$-topologization with
uniform-in-$N$ tight bound $N_0(w_{\max}) = 2\,w_{\max} - 1$: CLOSED
via spin-5 and spin-7 Prochazka-verified structure constants.
\item[VIII.11] Monster Schellekens 71-fold three-mechanism
chain-level closure: CLOSED.
\item[VIII.12] Non-simply-laced quantum Langlands via Frenkel--Hernandez
uniform lacing automorphism: CLOSED.
\end{enumerate}
All except VIII.6 (structural Platonic form, not closure) and VIII.9
(path identified, Vol~IV) closed this session. See
\texttt{notes/part\_viii\_vol\_iv\_roadmap\_synthesis.md}.

\subsubsection*{Vol~IV Realization table of contents (8 chapters)}

\begin{enumerate}
\item[IV.1] HEAL 247+ FM counters with explicit strengthening proofs
(not downgrades).
\item[IV.2] UPGRADE strengthened theorems: Theorem~A$^{\infty,2}$ at
properad level, Theorem~B universal regularisation, Theorem~C total
$(-(3g-3))$-shifted symplectic on $\overline{M}_{g,n}$, Theorem~D
tensor-Arakelov curvature, Theorem~H chiral Higher Deligne on
$Z^{\mathrm{der}}_{ch}$.
\item[IV.3] $100\%$ \texttt{@independent\_verification} decorator
coverage (Vol~II current: $0/1134 \to$ target $1134/1134$).
Prioritisation: climax theorems first
(thm:E3-topological-km, thm:E3-topological-DS-general,
thm:E3-topological-free-PVA, thm:global-triangle-boundary-linear,
thm:modular-bar-d-squared, thm:chd-ds-hochschild, thm:universal-holography,
thm:chd-deligne-tamarkin, thm:heptagon-7-fold, thm:fingerprint-completeness,
thm:periodic-cdg-admissible-all-simple, thm:higher-massey-odd-weight-tower).
\item[IV.4] Seven Theorems A, B, C, D, H, F (= Universal Holography),
G (= Infinite Fingerprint) written explicitly with full corner-case
analysis.
\item[IV.5] $M_{\mathrm{Kosz}}$ Koszulness Moduli Scheme as 14-chart
atlas on $\mathrm{GRT}_1(\mathbb{Q})$-equivariant moduli.
\item[IV.6] $\mathrm{GRT}_1(\mathbb{Q})$ action on nine-face torsor
plus heptagon structure, with F8 Brown motivic and F9 Willwacher
operadic.
\item[IV.7] Fingerprint classification table for all standard landscape
members (affine, Virasoro, $\mathcal{W}_N$, BP, super, lattice,
Monster, Schellekens 71).
\item[IV.8] Chiral quantum group uniformisation
$Q_{\mathfrak{g}}^{k, f, \mu}$ covering all eight specialisations
(Yangian, affine Yangian, shifted, truncated-shifted = BFN Coulomb,
finite-$\mathcal{W}$, affine principal $\mathcal{W}$, non-principal
affine $\mathcal{W}$, BP).
\end{enumerate}

\subsubsection*{Residual frontiers (13 items)}

\begin{enumerate}
\item[R1] Compact CY$_3$ Kapranov 3-shifted exterior Koszul tilting
(gauge-equivalence with non-trivial GRT$_1(\mathbb{Q})$ class;
conditional on tilting-object existence beyond toric / local CY$_3$).
\item[R2] Super-Yangian chiral quantum group beyond type-A
(CY-C conjectural; super-Yangian triality $S_3$ invariance of
shadow class at arbitrary rank).
\item[R3] Hodge conjecture for CY$_4$: reformulated in chiral-algebra
language via chiral Hodge decomposition; transcendental obstructions
(Voisin) external to programme.
\item[R4] Full arithmetic Universal Holography as published theorem
(VIII.9 completion).
\item[R5] Yang--Mills mass gap non-SUSY HT-twist construction;
obstruction (H3.b) = positivity of Costello--Gwilliam RG-flow fixed
point is the Clay problem proper.
\item[R6] Riemann Hypothesis: reformulated as diagonal-zero statement
on programme's Sewing--Selberg $L_{\mathcal{A}}(s, u)$; external.
\item[R7] BSD conjecture: reformulated via $L^{\mathrm{bar}}(\Phi(E), s)
= L(E, s)$; external.
\item[R8] Sato--Tate: reformulated via shadow-tower symmetric-power
reduction to Newton--Thorne 2021; external.
\item[R9] Grothendieck Period Conjecture: KZ-for-motivic closes;
GPC injectivity external (Apéry, Ball--Rivoal transcendence).
\item[R10] Compact-quintic CY-C chain-level (CY-A$_3$ conjectural
scope; atlas-existence open).
\item[R11] Exotic-nilpotent coset AD chain-level (single bounded
gap).
\item[R12] Higher-genus global Hitchin-quantised theory beyond
$g = 2, 3$ explicit.
\item[R13] Non-rational affine minimal-model tempering beyond
Adamovic--Milas $W(p)$.
\end{enumerate}

Sessions: session total of approximately $50$ adversarial
attack-then-heal agents launched, $\sim 45$ returned with substantive
closures or sharp reformulations, zero downgrades accepted.
All closures capture the strongest honest Platonic form of each
target.

\subsection{Summary: Platonic form of Vol~II}
\label{subsec:session-summary}

After this session, the Vol~II programme's Platonic form is:
\begin{enumerate}
\item \textbf{Seven theorems}: A, B, C, D, H (existing, with upgrades
including chiral Higher Deligne on H, tensor-Arakelov curvature on D,
$(\infty,2)$-properad on A), F (Universal Holography functor,
two-adjunction family-over-level), G (six-slot fingerprint
classification with $\kappa_{\mathrm{ch}}$).
\item \textbf{Eight parts}: I-VII existing plus Part~VIII
(From Frontier to Theorem, covering arithmetic UH, KZB irregular-singular
composition, FF tier-4 categorification, exotic-nilpotent closure).
\item \textbf{Five classes G/L/C/M/FF}: G/L/C/M via non-critical
$\Phi_{\mathrm{hol}}$; FF via critical $\mathrm{FF}\Phi_{\mathrm{hol}}$
onto Hitchin-quantised opers (Langlands-natural).
\item \textbf{$\kappa_W(\mathfrak{g}) = 12\,p_2 - 2r(3h+5)$}: universal
Chevalley--Shephard--Todd form.
\item \textbf{$\kappa_V(\mathfrak{g}) = 2\,\dim$}: universal at
$h^{\vee} \ne 0$.
\item \textbf{Heptagon edge (3)$\leftrightarrow$(4)}: rank-2
$\mathrm{GRT}_1(\mathbb{Q})^{\oplus 2}$-equivariant.
\item \textbf{Chain-level chiral Deligne--Tamarkin}: proved
associator-dependent via depth-3 Massey detected by $\zeta(3)\,\sigma_3$;
associator-independent form fails (odd-weight tower Brown motivic).
\item \textbf{Three Clay-adjacent arithmetic domains} (function-field
Langlands, F$_q$-shtuka, critical-level $\mathfrak{g}\leftrightarrow
\mathfrak{g}^L$) gain new content; seven Clay-level problems (RH,
BSD, Sato--Tate, Hodge, YM, GPC, KZ) are precisely reformulated
without solving.
\end{enumerate}

Notes directory (\texttt{/notes/}) contains ${\sim}50$ Platonic-reconstitution
and attack-then-heal documents capturing full derivations.

\end{document}
